\documentclass[11pt]{article}

\usepackage{fontspec}
\usepackage[margin=1in]{geometry}
\usepackage{microtype}
\usepackage{array}
\usepackage{booktabs}
\usepackage{hyperref}
\usepackage{longtable}
\usepackage{xcolor}
\usepackage{xurl}

\definecolor{ExploratoriumNavy}{HTML}{14273D}
\definecolor{ExploratoriumSlate}{HTML}{365A73}
\definecolor{ExploratoriumMist}{HTML}{F4F1EA}
\definecolor{ExploratoriumGold}{HTML}{BC7F16}

\hypersetup{
  colorlinks=true,
  linkcolor=ExploratoriumNavy,
  urlcolor=ExploratoriumSlate,
  pdftitle={Exploratorium Translation},
  pdfauthor={OpenAI Codex}
}

\defaultfontfeatures{Ligatures=TeX}
\setmainfont{DejaVu Serif}
\renewcommand{\arraystretch}{1.2}
\sloppy

\title{Exploratorium Translation\\\large Repository File, Link, and Fetch-State Atlas / Атлас файлов, ссылок и состояния загрузки}
\author{OpenAI Codex}
\date{March 18, 2026}

\begin{document}

\maketitle
\thispagestyle{empty}

\noindent\colorbox{ExploratoriumNavy}{%
  \parbox{\dimexpr\linewidth-2\fboxsep\relax}{%
    \vspace{10pt}
    {\color{white}\Large\bfseries Exploratorium Translation}\par
    \vspace{4pt}
    {\color{ExploratoriumMist}\normalsize English and Russian repository summaries repeated side-by-side, line-by-line, and page-by-page}\par
    \vspace{4pt}
    {\color{ExploratoriumGold}\normalsize Англо-русский документ, который повторяет одни и те же данные параллельно, построчно и постранично}\par
    \vspace{10pt}
  }%
}

\vspace{1em}

\noindent This document is generated from repository metadata through
\texttt{uv run repo-rag sync-exploratorium-translation --root .}. It inventories the current
state of referenced papers and documentation, summarizes all tracked files except the
exploratorium's own generated outputs, and summarizes every authored explicit HTTP or HTTPS URL.

\par\smallskip\noindent Этот документ генерируется из метаданных репозитория командой
\texttt{uv run repo-rag sync-exploratorium-translation --root .}. Он инвентаризирует текущее
состояние статей и документации, суммирует все отслеживаемые файлы, кроме собственных
сгенерированных выходов exploratorium, и описывает каждый авторский явный HTTP- или HTTPS-URL.

\tableofcontents
\newpage

% Generated by repo-rag sync-exploratorium-translation. Do not edit manually.
\section*{Repository Fetching State / Состояние загрузки источников}
Generated at \texttt{2026-05-13T10:30:07Z}. The repository currently tracks 403 summarized files, 91 authored explicit URL references, and 4 bibliography entries. The exploratorium excludes its own generated outputs (3 tracked path(s)) to avoid recursive drift.
\par\smallskip Сгенерировано в \texttt{2026-05-13T10:30:07Z}. Сейчас репозиторий содержит 403 суммаризированных файлов, 91 авторских явных URL и 4 библиографических записей. Сам exploratorium исключает собственные сгенерированные выходы (3 путей), чтобы не раздувать рекурсию.
\par\smallskip Machine-generated link surfaces excluded from the authored link table: 2632 URL occurrences (2427 unique URLs), mostly from lockfiles or raw logs.
\par\smallskip Машинно-сгенерированные ссылки, исключенные из авторской таблицы: 2632 вхождений URL (2427 уникальных URL), в основном из lockfile или сырых логов.
\section{Side-By-Side / Параллельный показ}
\subsection*{Referenced Papers And Documentation Fetch State / Состояние загрузки статей и документации}
\begingroup
\small
\renewcommand{\arraystretch}{1.16}
\begin{longtable}{>{\raggedright\arraybackslash}p{0.47\linewidth} >{\raggedright\arraybackslash}p{0.47\linewidth}}
\toprule
\textbf{English} & \textbf{Русский} \\
\midrule
\endfirsthead
\toprule
\textbf{English} & \textbf{Русский} \\
\midrule
\endhead
\textbf{lewis2020rag}: Retrieval-Augmented Generation for Knowledge-Intensive NLP Tasks. Kind: inproceedings. Cited in the bibliography, but no local PDF or paper mirror is tracked. Closest local companion paths: docs/architecture/dspy-rag-guide.md, docs/architecture/inspired/implementing-rag-with-dspy-technical-guide.md. Related paths: \path{docs/architecture/dspy-rag-guide.md, docs/architecture/inspired/implementing-rag-with-dspy-technical-guide.md}. & \textbf{lewis2020rag}: Retrieval-Augmented Generation for Knowledge-Intensive NLP Tasks. Тип: inproceedings. Источник есть в библиографии, но локальный PDF или зеркало статьи не отслеживаются. Ближайшие локальные сопроводительные пути: docs/architecture/dspy-rag-guide.md, docs/architecture/inspired/implementing-rag-with-dspy-technical-guide.md. Связанные пути: \path{docs/architecture/dspy-rag-guide.md, docs/architecture/inspired/implementing-rag-with-dspy-technical-guide.md}. \\
\textbf{khattab2024dspy}: DSPy: Compiling Declarative Language Model Calls into Self-Improving Pipelines. Kind: article. Cited in the bibliography, but no local PDF or paper mirror is tracked. Closest local companion paths: docs/architecture/dspy-rag-guide.md, docs/architecture/inspired/dspy-rag-tutorial.md, docs/guides/hushwheel-fixture-rag-guide.md. Related paths: \path{docs/architecture/dspy-rag-guide.md, docs/architecture/inspired/dspy-rag-tutorial.md, docs/guides/hushwheel-fixture-rag-guide.md}. & \textbf{khattab2024dspy}: DSPy: Compiling Declarative Language Model Calls into Self-Improving Pipelines. Тип: article. Источник есть в библиографии, но локальный PDF или зеркало статьи не отслеживаются. Ближайшие локальные сопроводительные пути: docs/architecture/dspy-rag-guide.md, docs/architecture/inspired/dspy-rag-tutorial.md, docs/guides/hushwheel-fixture-rag-guide.md. Связанные пути: \path{docs/architecture/dspy-rag-guide.md, docs/architecture/inspired/dspy-rag-tutorial.md, docs/guides/hushwheel-fixture-rag-guide.md}. \\
\textbf{mcp2024}: Model Context Protocol. Kind: misc. External documentation is cited by URL. No mirrored upstream copy is tracked, but local companion paths exist at docs/architecture/mcp-discovery.md. Related paths: \path{docs/architecture/mcp-discovery.md}. & \textbf{mcp2024}: Model Context Protocol. Тип: misc. Внешняя документация цитируется по URL. Зеркальная копия upstream не отслеживается, но есть локальные сопроводительные пути: docs/architecture/mcp-discovery.md. Связанные пути: \path{docs/architecture/mcp-discovery.md}. \\
\textbf{azureinference2025}: Azure AI Inference Documentation. Kind: misc. External documentation is cited by URL. No mirrored upstream copy is tracked, but local companion paths exist at docs/operations/azure-deployment.md. Related paths: \path{docs/operations/azure-deployment.md}. & \textbf{azureinference2025}: Azure AI Inference Documentation. Тип: misc. Внешняя документация цитируется по URL. Зеркальная копия upstream не отслеживается, но есть локальные сопроводительные пути: docs/operations/azure-deployment.md. Связанные пути: \path{docs/operations/azure-deployment.md}. \\
\bottomrule
\end{longtable}
\endgroup
\subsection*{Summaries Of All Files / Сводки по всем файлам}
\begingroup
\small
\renewcommand{\arraystretch}{1.16}
\begin{longtable}{>{\raggedright\arraybackslash}p{0.47\linewidth} >{\raggedright\arraybackslash}p{0.47\linewidth}}
\toprule
\textbf{English} & \textbf{Русский} \\
\midrule
\endfirsthead
\toprule
\textbf{English} & \textbf{Русский} \\
\midrule
\endhead
\textbf{\path{.codex/skills/exploratorium-translation-sync/SKILL.md}}. Kind: markdown. Size: 1722 bytes. Shape: 32 lines. Markdown document, 32 lines, 0 explicit URLs. Lead heading: Exploratorium Translation Sync. & \textbf{\path{.codex/skills/exploratorium-translation-sync/SKILL.md}}. Тип: markdown. Размер: 1722 байт. Форма: 32 строк. Документ Markdown, 32 строк, явных URL: 0. Главный заголовок: Exploratorium Translation Sync. \\
\textbf{\path{.codex/skills/file-summary-sync/SKILL.md}}. Kind: markdown. Size: 1367 bytes. Shape: 29 lines. Markdown document, 29 lines, 0 explicit URLs. Lead heading: File Summary Sync. & \textbf{\path{.codex/skills/file-summary-sync/SKILL.md}}. Тип: markdown. Размер: 1367 байт. Форма: 29 строк. Документ Markdown, 29 строк, явных URL: 0. Главный заголовок: File Summary Sync. \\
\textbf{\path{.codex/skills/notebook-playbook-sync/SKILL.md}}. Kind: markdown. Size: 2085 bytes. Shape: 48 lines. Markdown document, 48 lines, 0 explicit URLs. Lead heading: Notebook Playbook Sync. & \textbf{\path{.codex/skills/notebook-playbook-sync/SKILL.md}}. Тип: markdown. Размер: 2085 байт. Форма: 48 строк. Документ Markdown, 48 строк, явных URL: 0. Главный заголовок: Notebook Playbook Sync. \\
\textbf{\path{.codex/skills/notebook-playbook-sync/agents/openai.yaml}}. Kind: yaml. Size: 300 bytes. Shape: 5 lines. YAML data or workflow file, 5 lines, 0 explicit URLs. Top keys: interface. & \textbf{\path{.codex/skills/notebook-playbook-sync/agents/openai.yaml}}. Тип: yaml. Размер: 300 байт. Форма: 5 строк. Файл YAML с данными или workflow, 5 строк, явных URL: 0. Верхние ключи: interface. \\
\textbf{\path{.codex/skills/post-push-gh-run-logging/SKILL.md}}. Kind: markdown. Size: 2059 bytes. Shape: 41 lines. Markdown document, 41 lines, 0 explicit URLs. Lead heading: Post Push Gh Run Logging. & \textbf{\path{.codex/skills/post-push-gh-run-logging/SKILL.md}}. Тип: markdown. Размер: 2059 байт. Форма: 41 строк. Документ Markdown, 41 строк, явных URL: 0. Главный заголовок: Post Push Gh Run Logging. \\
\textbf{\path{.codex/skills/post-push-gh-run-logging/agents/openai.yaml}}. Kind: yaml. Size: 281 bytes. Shape: 5 lines. YAML data or workflow file, 5 lines, 0 explicit URLs. Top keys: interface. & \textbf{\path{.codex/skills/post-push-gh-run-logging/agents/openai.yaml}}. Тип: yaml. Размер: 281 байт. Форма: 5 строк. Файл YAML с данными или workflow, 5 строк, явных URL: 0. Верхние ключи: interface. \\
\textbf{\path{.codex/skills/post-push-gh-run-logging/scripts/write_gh_run_log.py}}. Kind: python. Size: 6346 bytes. Shape: 178 lines. Python module, 178 lines, 0 explicit URLs. Module intent: No module docstring detected. & \textbf{\path{.codex/skills/post-push-gh-run-logging/scripts/write_gh_run_log.py}}. Тип: python. Размер: 6346 байт. Форма: 178 строк. Модуль Python, 178 строк, явных URL: 0. Назначение модуля: No module docstring detected. \\
\textbf{\path{.codex/skills/repo-verification-audit-loop/SKILL.md}}. Kind: markdown. Size: 2589 bytes. Shape: 53 lines. Markdown document, 53 lines, 0 explicit URLs. Lead heading: Repo Verification Audit Loop. & \textbf{\path{.codex/skills/repo-verification-audit-loop/SKILL.md}}. Тип: markdown. Размер: 2589 байт. Форма: 53 строк. Документ Markdown, 53 строк, явных URL: 0. Главный заголовок: Repo Verification Audit Loop. \\
\textbf{\path{.codex/skills/repo-verification-audit-loop/agents/openai.yaml}}. Kind: yaml. Size: 272 bytes. Shape: 5 lines. YAML data or workflow file, 5 lines, 0 explicit URLs. Top keys: interface. & \textbf{\path{.codex/skills/repo-verification-audit-loop/agents/openai.yaml}}. Тип: yaml. Размер: 272 байт. Форма: 5 строк. Файл YAML с данными или workflow, 5 строк, явных URL: 0. Верхние ключи: interface. \\
\textbf{\path{.codex/skills/rust-sqlite-lookup/SKILL.md}}. Kind: markdown. Size: 1297 bytes. Shape: 36 lines. Markdown document, 36 lines, 0 explicit URLs. Lead heading: Rust SQLite Lookup. & \textbf{\path{.codex/skills/rust-sqlite-lookup/SKILL.md}}. Тип: markdown. Размер: 1297 байт. Форма: 36 строк. Документ Markdown, 36 строк, явных URL: 0. Главный заголовок: Rust SQLite Lookup. \\
\textbf{\path{.codex/skills/todo-backlog-sync/SKILL.md}}. Kind: markdown. Size: 1289 bytes. Shape: 29 lines. Markdown document, 29 lines, 0 explicit URLs. Lead heading: Todo Backlog Sync. & \textbf{\path{.codex/skills/todo-backlog-sync/SKILL.md}}. Тип: markdown. Размер: 1289 байт. Форма: 29 строк. Документ Markdown, 29 строк, явных URL: 0. Главный заголовок: Todo Backlog Sync. \\
\textbf{\path{.dockerignore}}. Kind: generic. Size: 150 bytes. Shape: 17 lines. Tracked repository file, 17 lines, 0 explicit URLs. Opening line: .git & \textbf{\path{.dockerignore}}. Тип: generic. Размер: 150 байт. Форма: 17 строк. Отслеживаемый файл репозитория, 17 строк, явных URL: 0. Первая строка: .git \\
\textbf{\path{.env.sample}}. Kind: generic. Size: 905 bytes. Shape: 17 lines. Tracked repository file, 17 lines, 0 explicit URLs. Opening line: \# Copy this file to .env and replace placeholder values. & \textbf{\path{.env.sample}}. Тип: generic. Размер: 905 байт. Форма: 17 строк. Отслеживаемый файл репозитория, 17 строк, явных URL: 0. Первая строка: \# Copy this file to .env and replace placeholder values. \\
\textbf{\path{.gitattributes}}. Kind: generic. Size: 126 bytes. Shape: 4 lines. Tracked repository file, 4 lines, 0 explicit URLs. Opening line: *.png filter=lfs diff=lfs merge=lfs -text & \textbf{\path{.gitattributes}}. Тип: generic. Размер: 126 байт. Форма: 4 строк. Отслеживаемый файл репозитория, 4 строк, явных URL: 0. Первая строка: *.png filter=lfs diff=lfs merge=lfs -text \\
\textbf{\path{.github/workflows/ci.yml}}. Kind: yaml. Size: 4185 bytes. Shape: 146 lines. YAML data or workflow file, 146 lines, 0 explicit URLs. Top keys: name, True, jobs. & \textbf{\path{.github/workflows/ci.yml}}. Тип: yaml. Размер: 4185 байт. Форма: 146 строк. Файл YAML с данными или workflow, 146 строк, явных URL: 0. Верхние ключи: name, True, jobs. \\
\textbf{\path{.github/workflows/hushwheel-quality.yml}}. Kind: yaml. Size: 3526 bytes. Shape: 104 lines. YAML data or workflow file, 104 lines, 0 explicit URLs. Top keys: name, True, env. & \textbf{\path{.github/workflows/hushwheel-quality.yml}}. Тип: yaml. Размер: 3526 байт. Форма: 104 строк. Файл YAML с данными или workflow, 104 строк, явных URL: 0. Верхние ключи: name, True, env. \\
\textbf{\path{.github/workflows/pages.yml}}. Kind: yaml. Size: 2424 bytes. Shape: 99 lines. YAML data or workflow file, 99 lines, 0 explicit URLs. Top keys: name, True, env. & \textbf{\path{.github/workflows/pages.yml}}. Тип: yaml. Размер: 2424 байт. Форма: 99 строк. Файл YAML с данными или workflow, 99 строк, явных URL: 0. Верхние ключи: name, True, env. \\
\textbf{\path{.github/workflows/publication-pdf.yml}}. Kind: yaml. Size: 5549 bytes. Shape: 144 lines. YAML data or workflow file, 144 lines, 0 explicit URLs. Top keys: name, True, env. & \textbf{\path{.github/workflows/publication-pdf.yml}}. Тип: yaml. Размер: 5549 байт. Форма: 144 строк. Файл YAML с данными или workflow, 144 строк, явных URL: 0. Верхние ключи: name, True, env. \\
\textbf{\path{.github/workflows/publish.yml}}. Kind: yaml. Size: 847 bytes. Shape: 43 lines. YAML data or workflow file, 43 lines, 0 explicit URLs. Top keys: name, True, jobs. & \textbf{\path{.github/workflows/publish.yml}}. Тип: yaml. Размер: 847 байт. Форма: 43 строк. Файл YAML с данными или workflow, 43 строк, явных URL: 0. Верхние ключи: name, True, jobs. \\
\textbf{\path{.gitignore}}. Kind: generic. Size: 222 bytes. Shape: 21 lines. Tracked repository file, 21 lines, 0 explicit URLs. Opening line: .venv/ & \textbf{\path{.gitignore}}. Тип: generic. Размер: 222 байт. Форма: 21 строк. Отслеживаемый файл репозитория, 21 строк, явных URL: 0. Первая строка: .venv/ \\
\textbf{\path{.pre-commit-config.yaml}}. Kind: yaml. Size: 1689 bytes. Shape: 53 lines. YAML data or workflow file, 53 lines, 0 explicit URLs. Top keys: repos. & \textbf{\path{.pre-commit-config.yaml}}. Тип: yaml. Размер: 1689 байт. Форма: 53 строк. Файл YAML с данными или workflow, 53 строк, явных URL: 0. Верхние ключи: repos. \\
\textbf{\path{.python-version}}. Kind: generic. Size: 5 bytes. Shape: 2 lines. Tracked repository file, 2 lines, 0 explicit URLs. Opening line: 3.11 & \textbf{\path{.python-version}}. Тип: generic. Размер: 5 байт. Форма: 2 строк. Отслеживаемый файл репозитория, 2 строк, явных URL: 0. Первая строка: 3.11 \\
\textbf{\path{.yamllint.yml}}. Kind: yaml. Size: 416 bytes. Shape: 20 lines. YAML data or workflow file, 20 lines, 0 explicit URLs. Top keys: extends, rules. & \textbf{\path{.yamllint.yml}}. Тип: yaml. Размер: 416 байт. Форма: 20 строк. Файл YAML с данными или workflow, 20 строк, явных URL: 0. Верхние ключи: extends, rules. \\
\textbf{\path{AGENTS.md}}. Kind: markdown. Size: 6999 bytes. Shape: 108 lines. Markdown document, 108 lines, 0 explicit URLs. Lead heading: Repository Agent Instructions. & \textbf{\path{AGENTS.md}}. Тип: markdown. Размер: 6999 байт. Форма: 108 строк. Документ Markdown, 108 строк, явных URL: 0. Главный заголовок: Repository Agent Instructions. \\
\textbf{\path{AGENTS.md.d/00-audit-baseline.md}}. Kind: markdown. Size: 592 bytes. Shape: 15 lines. Markdown document, 15 lines, 0 explicit URLs. Lead heading: Audit Baseline Instructions. & \textbf{\path{AGENTS.md.d/00-audit-baseline.md}}. Тип: markdown. Размер: 592 байт. Форма: 15 строк. Документ Markdown, 15 строк, явных URL: 0. Главный заголовок: Audit Baseline Instructions. \\
\textbf{\path{AGENTS.md.d/10-verification-loop.md}}. Kind: markdown. Size: 1175 bytes. Shape: 22 lines. Markdown document, 22 lines, 0 explicit URLs. Lead heading: Verification Loop Instructions. & \textbf{\path{AGENTS.md.d/10-verification-loop.md}}. Тип: markdown. Размер: 1175 байт. Форма: 22 строк. Документ Markdown, 22 строк, явных URL: 0. Главный заголовок: Verification Loop Instructions. \\
\textbf{\path{AGENTS.md.d/FILES.md}}. Kind: markdown. Size: 2247 bytes. Shape: 41 lines. Markdown document, 41 lines, 0 explicit URLs. Lead heading: File Summary Instructions. & \textbf{\path{AGENTS.md.d/FILES.md}}. Тип: markdown. Размер: 2247 байт. Форма: 41 строк. Документ Markdown, 41 строк, явных URL: 0. Главный заголовок: File Summary Instructions. \\
\textbf{\path{AGENTS.md.d/RUST_LOOKUP.md}}. Kind: markdown. Size: 1803 bytes. Shape: 36 lines. Markdown document, 36 lines, 0 explicit URLs. Lead heading: Rust SQLite Lookup Instructions. & \textbf{\path{AGENTS.md.d/RUST_LOOKUP.md}}. Тип: markdown. Размер: 1803 байт. Форма: 36 строк. Документ Markdown, 36 строк, явных URL: 0. Главный заголовок: Rust SQLite Lookup Instructions. \\
\textbf{\path{Dockerfile}}. Kind: generic. Size: 1080 bytes. Shape: 47 lines. Tracked repository file, 47 lines, 0 explicit URLs. Opening line: ARG PYTHON\_BASE\_IMAGE=python:3.11-slim-bookworm & \textbf{\path{Dockerfile}}. Тип: generic. Размер: 1080 байт. Форма: 47 строк. Отслеживаемый файл репозитория, 47 строк, явных URL: 0. Первая строка: ARG PYTHON\_BASE\_IMAGE=python:3.11-slim-bookworm \\
\textbf{\path{FILES.csv}}. Kind: generic. Size: 52899 bytes. Shape: binary surface. Tracked repository file, 52899 bytes of binary data, 0 explicit URLs. Opening line: No non-empty preview line. & \textbf{\path{FILES.csv}}. Тип: generic. Размер: 52899 байт. Форма: бинарная поверхность. Отслеживаемый файл репозитория, 52899 байт бинарных данных, явных URL: 0. Первая строка: No non-empty preview line. \\
\textbf{\path{FILES.md}}. Kind: markdown. Size: 59726 bytes. Shape: 417 lines. Markdown document, 417 lines, 1 explicit URLs. Lead heading: Repository File Inventory. & \textbf{\path{FILES.md}}. Тип: markdown. Размер: 59726 байт. Форма: 417 строк. Документ Markdown, 417 строк, явных URL: 1. Главный заголовок: Repository File Inventory. \\
\textbf{\path{Makefile}}. Kind: make. Size: 23701 bytes. Shape: 506 lines. Repository Makefile, 506 lines, 1 explicit URLs. It exposes the shared automation surfaces. & \textbf{\path{Makefile}}. Тип: make. Размер: 23701 байт. Форма: 506 строк. Makefile репозитория, 506 строк, явных URL: 1. Он открывает общие поверхности автоматизации. \\
\textbf{\path{README.md}}. Kind: markdown. Size: 26222 bytes. Shape: 312 lines. Markdown document, 312 lines, 6 explicit URLs. Lead heading: Repository RAG Lab. & \textbf{\path{README.md}}. Тип: markdown. Размер: 26222 байт. Форма: 312 строк. Документ Markdown, 312 строк, явных URL: 6. Главный заголовок: Repository RAG Lab. \\
\textbf{\path{TODO.MD}}. Kind: markdown. Size: 6104 bytes. Shape: 22 lines. Markdown document, 22 lines, 0 explicit URLs. Lead heading: TODO. & \textbf{\path{TODO.MD}}. Тип: markdown. Размер: 6104 байт. Форма: 22 строк. Документ Markdown, 22 строк, явных URL: 0. Главный заголовок: TODO. \\
\textbf{\path{config/retrieval-profile.json}}. Kind: json. Size: 1791 bytes. Shape: 66 lines. JSON data file, 66 lines, 0 explicit URLs. Top keys: name, retrieval\_mode, rerank\_candidate\_pool\_size. & \textbf{\path{config/retrieval-profile.json}}. Тип: json. Размер: 1791 байт. Форма: 66 строк. Файл JSON, 66 строк, явных URL: 0. Верхние ключи: name, retrieval\_mode, rerank\_candidate\_pool\_size. \\
\textbf{\path{data/questions/repository.yaml}}. Kind: yaml. Size: 283 bytes. Shape: 7 lines. YAML data or workflow file, 7 lines, 0 explicit URLs. Top keys: no top-level mapping keys. & \textbf{\path{data/questions/repository.yaml}}. Тип: yaml. Размер: 283 байт. Форма: 7 строк. Файл YAML с данными или workflow, 7 строк, явных URL: 0. Верхние ключи: no top-level mapping keys. \\
\textbf{\path{docs/architecture/dspy-rag-guide.md}}. Kind: markdown. Size: 50272 bytes. Shape: 911 lines. Markdown document, 911 lines, 1 explicit URLs. Lead heading: DSPy Guide For This Repository. & \textbf{\path{docs/architecture/dspy-rag-guide.md}}. Тип: markdown. Размер: 50272 байт. Форма: 911 строк. Документ Markdown, 911 строк, явных URL: 1. Главный заголовок: DSPy Guide For This Repository. \\
\textbf{\path{docs/architecture/inspired/dspy-rag-tutorial.md}}. Kind: markdown. Size: 1474 bytes. Shape: 35 lines. Markdown document, 35 lines, 1 explicit URLs. Lead heading: DSPy Tutorial: Basic RAG. & \textbf{\path{docs/architecture/inspired/dspy-rag-tutorial.md}}. Тип: markdown. Размер: 1474 байт. Форма: 35 строк. Документ Markdown, 35 строк, явных URL: 1. Главный заголовок: DSPy Tutorial: Basic RAG. \\
\textbf{\path{docs/architecture/inspired/implementing-rag-with-dspy-technical-guide.md}}. Kind: markdown. Size: 1555 bytes. Shape: 38 lines. Markdown document, 38 lines, 1 explicit URLs. Lead heading: Implementing RAG with DSPy: Technical Guide. & \textbf{\path{docs/architecture/inspired/implementing-rag-with-dspy-technical-guide.md}}. Тип: markdown. Размер: 1555 байт. Форма: 38 строк. Документ Markdown, 38 строк, явных URL: 1. Главный заголовок: Implementing RAG with DSPy: Technical Guide. \\
\textbf{\path{docs/architecture/mcp-discovery.md}}. Kind: markdown. Size: 1757 bytes. Shape: 54 lines. Markdown document, 54 lines, 0 explicit URLs. Lead heading: MCP Discovery Notes. & \textbf{\path{docs/architecture/mcp-discovery.md}}. Тип: markdown. Размер: 1757 байт. Форма: 54 строк. Документ Markdown, 54 строк, явных URL: 0. Главный заголовок: MCP Discovery Notes. \\
\textbf{\path{docs/architecture/package-api.md}}. Kind: markdown. Size: 9186 bytes. Shape: 127 lines. Markdown document, 127 lines, 0 explicit URLs. Lead heading: Package API Notes. & \textbf{\path{docs/architecture/package-api.md}}. Тип: markdown. Размер: 9186 байт. Форма: 127 строк. Документ Markdown, 127 строк, явных URL: 0. Главный заголовок: Package API Notes. \\
\textbf{\path{docs/architecture/research-narrative.md}}. Kind: markdown. Size: 82061 bytes. Shape: 1077 lines. Markdown document, 1077 lines, 0 explicit URLs. Lead heading: Repository Research Narrative. & \textbf{\path{docs/architecture/research-narrative.md}}. Тип: markdown. Размер: 82061 байт. Форма: 1077 строк. Документ Markdown, 1077 строк, явных URL: 0. Главный заголовок: Repository Research Narrative. \\
\textbf{\path{docs/archive/repo-completeness-checklist.md}}. Kind: markdown. Size: 7450 bytes. Shape: 238 lines. Markdown document, 238 lines, 0 explicit URLs. Lead heading: Repository Completeness Checklist. & \textbf{\path{docs/archive/repo-completeness-checklist.md}}. Тип: markdown. Размер: 7450 байт. Форма: 238 строк. Документ Markdown, 238 строк, явных URL: 0. Главный заголовок: Repository Completeness Checklist. \\
\textbf{\path{docs/audit/2026-03-15-repo-state.md}}. Kind: markdown. Size: 5989 bytes. Shape: 151 lines. Markdown document, 151 lines, 0 explicit URLs. Lead heading: Repository State Audit. & \textbf{\path{docs/audit/2026-03-15-repo-state.md}}. Тип: markdown. Размер: 5989 байт. Форма: 151 строк. Документ Markdown, 151 строк, явных URL: 0. Главный заголовок: Repository State Audit. \\
\textbf{\path{docs/audit/2026-03-17-docs-and-uv-refresh.md}}. Kind: markdown. Size: 3977 bytes. Shape: 79 lines. Markdown document, 79 lines, 0 explicit URLs. Lead heading: Documentation And UV Workflow Refresh Audit. & \textbf{\path{docs/audit/2026-03-17-docs-and-uv-refresh.md}}. Тип: markdown. Размер: 3977 байт. Форма: 79 строк. Документ Markdown, 79 строк, явных URL: 0. Главный заголовок: Documentation And UV Workflow Refresh Audit. \\
\textbf{\path{docs/audit/2026-03-17-full-build.md}}. Kind: markdown. Size: 6715 bytes. Shape: 45 lines. Markdown document, 45 lines, 0 explicit URLs. Lead heading: Full Build Audit. & \textbf{\path{docs/audit/2026-03-17-full-build.md}}. Тип: markdown. Размер: 6715 байт. Форма: 45 строк. Документ Markdown, 45 строк, явных URL: 0. Главный заголовок: Full Build Audit. \\
\textbf{\path{docs/audit/2026-03-17-gh-actions-watch-loop.md}}. Kind: markdown. Size: 3797 bytes. Shape: 78 lines. Markdown document, 78 lines, 0 explicit URLs. Lead heading: GitHub Actions Watch Loop Audit. & \textbf{\path{docs/audit/2026-03-17-gh-actions-watch-loop.md}}. Тип: markdown. Размер: 3797 байт. Форма: 78 строк. Документ Markdown, 78 строк, явных URL: 0. Главный заголовок: GitHub Actions Watch Loop Audit. \\
\textbf{\path{docs/audit/2026-03-17-hushwheel-fixture.md}}. Kind: markdown. Size: 5090 bytes. Shape: 80 lines. Markdown document, 80 lines, 0 explicit URLs. Lead heading: Hushwheel Fixture Audit. & \textbf{\path{docs/audit/2026-03-17-hushwheel-fixture.md}}. Тип: markdown. Размер: 5090 байт. Форма: 80 строк. Документ Markdown, 80 строк, явных URL: 0. Главный заголовок: Hushwheel Fixture Audit. \\
\textbf{\path{docs/audit/2026-03-17-hushwheel-rag-playbook.md}}. Kind: markdown. Size: 4797 bytes. Shape: 83 lines. Markdown document, 83 lines, 0 explicit URLs. Lead heading: Hushwheel RAG Playbook Audit. & \textbf{\path{docs/audit/2026-03-17-hushwheel-rag-playbook.md}}. Тип: markdown. Размер: 4797 байт. Форма: 83 строк. Документ Markdown, 83 строк, явных URL: 0. Главный заголовок: Hushwheel RAG Playbook Audit. \\
\textbf{\path{docs/audit/2026-03-17-notebook-scaffolding.md}}. Kind: markdown. Size: 3641 bytes. Shape: 70 lines. Markdown document, 70 lines, 0 explicit URLs. Lead heading: Notebook Scaffolding Audit. & \textbf{\path{docs/audit/2026-03-17-notebook-scaffolding.md}}. Тип: markdown. Размер: 3641 байт. Форма: 70 строк. Документ Markdown, 70 строк, явных URL: 0. Главный заголовок: Notebook Scaffolding Audit. \\
\textbf{\path{docs/audit/2026-03-17-publication-article.md}}. Kind: markdown. Size: 2321 bytes. Shape: 54 lines. Markdown document, 54 lines, 0 explicit URLs. Lead heading: Publication Article Audit. & \textbf{\path{docs/audit/2026-03-17-publication-article.md}}. Тип: markdown. Размер: 2321 байт. Форма: 54 строк. Документ Markdown, 54 строк, явных URL: 0. Главный заголовок: Publication Article Audit. \\
\textbf{\path{docs/audit/2026-03-17-repo-completeness-checklist.md}}. Kind: markdown. Size: 3114 bytes. Shape: 59 lines. Markdown document, 59 lines, 0 explicit URLs. Lead heading: Repository Completeness Checklist Audit. & \textbf{\path{docs/audit/2026-03-17-repo-completeness-checklist.md}}. Тип: markdown. Размер: 3114 байт. Форма: 59 строк. Документ Markdown, 59 строк, явных URL: 0. Главный заголовок: Repository Completeness Checklist Audit. \\
\textbf{\path{docs/audit/2026-03-18-actions-caching.md}}. Kind: markdown. Size: 4189 bytes. Shape: 87 lines. Markdown document, 87 lines, 0 explicit URLs. Lead heading: GitHub Actions Caching Audit. & \textbf{\path{docs/audit/2026-03-18-actions-caching.md}}. Тип: markdown. Размер: 4189 байт. Форма: 87 строк. Документ Markdown, 87 строк, явных URL: 0. Главный заголовок: GitHub Actions Caching Audit. \\
\textbf{\path{docs/audit/2026-03-18-azure-inference-endpoint-probe.md}}. Kind: markdown. Size: 2648 bytes. Shape: 63 lines. Markdown document, 63 lines, 0 explicit URLs. Lead heading: Azure Inference Endpoint Probe Audit. & \textbf{\path{docs/audit/2026-03-18-azure-inference-endpoint-probe.md}}. Тип: markdown. Размер: 2648 байт. Форма: 63 строк. Документ Markdown, 63 строк, явных URL: 0. Главный заголовок: Azure Inference Endpoint Probe Audit. \\
\textbf{\path{docs/audit/2026-03-18-azure-runtime-surfaces.md}}. Kind: markdown. Size: 5307 bytes. Shape: 91 lines. Markdown document, 91 lines, 0 explicit URLs. Lead heading: Azure Runtime Surfaces Audit. & \textbf{\path{docs/audit/2026-03-18-azure-runtime-surfaces.md}}. Тип: markdown. Размер: 5307 байт. Форма: 91 строк. Документ Markdown, 91 строк, явных URL: 0. Главный заголовок: Azure Runtime Surfaces Audit. \\
\textbf{\path{docs/audit/2026-03-18-dspy-training-path.md}}. Kind: markdown. Size: 4605 bytes. Shape: 81 lines. Markdown document, 81 lines, 0 explicit URLs. Lead heading: DSPy Training Path Audit. & \textbf{\path{docs/audit/2026-03-18-dspy-training-path.md}}. Тип: markdown. Размер: 4605 байт. Форма: 81 строк. Документ Markdown, 81 строк, явных URL: 0. Главный заголовок: DSPy Training Path Audit. \\
\textbf{\path{docs/audit/2026-03-18-hushwheel-program-harness.md}}. Kind: markdown. Size: 3907 bytes. Shape: 76 lines. Markdown document, 76 lines, 0 explicit URLs. Lead heading: Hushwheel Program Harness Audit. & \textbf{\path{docs/audit/2026-03-18-hushwheel-program-harness.md}}. Тип: markdown. Размер: 3907 байт. Форма: 76 строк. Документ Markdown, 76 строк, явных URL: 0. Главный заголовок: Hushwheel Program Harness Audit. \\
\textbf{\path{docs/audit/2026-03-18-notebook-execution.md}}. Kind: markdown. Size: 2577 bytes. Shape: 66 lines. Markdown document, 66 lines, 0 explicit URLs. Lead heading: Notebook Execution Audit. & \textbf{\path{docs/audit/2026-03-18-notebook-execution.md}}. Тип: markdown. Размер: 2577 байт. Форма: 66 строк. Документ Markdown, 66 строк, явных URL: 0. Главный заголовок: Notebook Execution Audit. \\
\textbf{\path{docs/audit/2026-03-18-notebook-runner-harness.md}}. Kind: markdown. Size: 4121 bytes. Shape: 91 lines. Markdown document, 91 lines, 0 explicit URLs. Lead heading: Notebook Runner Harness Audit. & \textbf{\path{docs/audit/2026-03-18-notebook-runner-harness.md}}. Тип: markdown. Размер: 4121 байт. Форма: 91 строк. Документ Markdown, 91 строк, явных URL: 0. Главный заголовок: Notebook Runner Harness Audit. \\
\textbf{\path{docs/audit/2026-03-18-repository-benchmark-broadening.md}}. Kind: markdown. Size: 3728 bytes. Shape: 74 lines. Markdown document, 74 lines, 0 explicit URLs. Lead heading: Repository Benchmark Broadening Audit. & \textbf{\path{docs/audit/2026-03-18-repository-benchmark-broadening.md}}. Тип: markdown. Размер: 3728 байт. Форма: 74 строк. Документ Markdown, 74 строк, явных URL: 0. Главный заголовок: Repository Benchmark Broadening Audit. \\
\textbf{\path{docs/audit/2026-03-18-retest-with-env-refresh.md}}. Kind: markdown. Size: 5102 bytes. Shape: 105 lines. Markdown document, 105 lines, 0 explicit URLs. Lead heading: Env Refresh Retest Audit. & \textbf{\path{docs/audit/2026-03-18-retest-with-env-refresh.md}}. Тип: markdown. Размер: 5102 байт. Форма: 105 строк. Документ Markdown, 105 строк, явных URL: 0. Главный заголовок: Env Refresh Retest Audit. \\
\textbf{\path{docs/audit/2026-03-18-retrieval-evaluation-suite.md}}. Kind: markdown. Size: 4031 bytes. Shape: 80 lines. Markdown document, 80 lines, 0 explicit URLs. Lead heading: Retrieval Evaluation Suite Audit. & \textbf{\path{docs/audit/2026-03-18-retrieval-evaluation-suite.md}}. Тип: markdown. Размер: 4031 байт. Форма: 80 строк. Документ Markdown, 80 строк, явных URL: 0. Главный заголовок: Retrieval Evaluation Suite Audit. \\
\textbf{\path{docs/audit/2026-03-18-retrieval-ranking-refresh.md}}. Kind: markdown. Size: 6479 bytes. Shape: 109 lines. Markdown document, 109 lines, 0 explicit URLs. Lead heading: Retrieval Ranking Refresh Audit. & \textbf{\path{docs/audit/2026-03-18-retrieval-ranking-refresh.md}}. Тип: markdown. Размер: 6479 байт. Форма: 109 строк. Документ Markdown, 109 строк, явных URL: 0. Главный заголовок: Retrieval Ranking Refresh Audit. \\
\textbf{\path{docs/audit/2026-03-18-rust-and-benchmark-hygiene.md}}. Kind: markdown. Size: 2049 bytes. Shape: 53 lines. Markdown document, 53 lines, 0 explicit URLs. Lead heading: Rust And Benchmark Hygiene Audit. & \textbf{\path{docs/audit/2026-03-18-rust-and-benchmark-hygiene.md}}. Тип: markdown. Размер: 2049 байт. Форма: 53 строк. Документ Markdown, 53 строк, явных URL: 0. Главный заголовок: Rust And Benchmark Hygiene Audit. \\
\textbf{\path{docs/audit/2026-03-18-todo-backlog-sync.md}}. Kind: markdown. Size: 4104 bytes. Shape: 83 lines. Markdown document, 83 lines, 0 explicit URLs. Lead heading: Todo Backlog Sync Audit. & \textbf{\path{docs/audit/2026-03-18-todo-backlog-sync.md}}. Тип: markdown. Размер: 4104 байт. Форма: 83 строк. Документ Markdown, 83 строк, явных URL: 0. Главный заголовок: Todo Backlog Sync Audit. \\
\textbf{\path{docs/audit/2026-03-18-z-notebook-runner-harness.md}}. Kind: markdown. Size: 4153 bytes. Shape: 104 lines. Markdown document, 104 lines, 0 explicit URLs. Lead heading: Notebook Runner Harness Audit. & \textbf{\path{docs/audit/2026-03-18-z-notebook-runner-harness.md}}. Тип: markdown. Размер: 4153 байт. Форма: 104 строк. Документ Markdown, 104 строк, явных URL: 0. Главный заголовок: Notebook Runner Harness Audit. \\
\textbf{\path{docs/audit/2026-03-18-zz-file-summary-inventory.md}}. Kind: markdown. Size: 2596 bytes. Shape: 69 lines. Markdown document, 69 lines, 0 explicit URLs. Lead heading: File Summary Inventory Audit. & \textbf{\path{docs/audit/2026-03-18-zz-file-summary-inventory.md}}. Тип: markdown. Размер: 2596 байт. Форма: 69 строк. Документ Markdown, 69 строк, явных URL: 0. Главный заголовок: File Summary Inventory Audit. \\
\textbf{\path{docs/audit/2026-03-18-zz-research-narrative.md}}. Kind: markdown. Size: 3312 bytes. Shape: 79 lines. Markdown document, 79 lines, 0 explicit URLs. Lead heading: Research Narrative Audit. & \textbf{\path{docs/audit/2026-03-18-zz-research-narrative.md}}. Тип: markdown. Размер: 3312 байт. Форма: 79 строк. Документ Markdown, 79 строк, явных URL: 0. Главный заголовок: Research Narrative Audit. \\
\textbf{\path{docs/audit/2026-03-18-zzz-session-consolidation-on-origin-master.md}}. Kind: markdown. Size: 5061 bytes. Shape: 109 lines. Markdown document, 109 lines, 0 explicit URLs. Lead heading: Session Consolidation On Origin Master. & \textbf{\path{docs/audit/2026-03-18-zzz-session-consolidation-on-origin-master.md}}. Тип: markdown. Размер: 5061 байт. Форма: 109 строк. Документ Markdown, 109 строк, явных URL: 0. Главный заголовок: Session Consolidation On Origin Master. \\
\textbf{\path{docs/audit/2026-03-18-zzzz-rust-sqlite-lookup.md}}. Kind: markdown. Size: 5404 bytes. Shape: 105 lines. Markdown document, 105 lines, 0 explicit URLs. Lead heading: Rust SQLite Lookup And Retrieval Tag Audit. & \textbf{\path{docs/audit/2026-03-18-zzzz-rust-sqlite-lookup.md}}. Тип: markdown. Размер: 5404 байт. Форма: 105 строк. Документ Markdown, 105 строк, явных URL: 0. Главный заголовок: Rust SQLite Lookup And Retrieval Tag Audit. \\
\textbf{\path{docs/audit/2026-03-18-zzzz-session-consolidation-master-ci-repair.md}}. Kind: markdown. Size: 3270 bytes. Shape: 73 lines. Markdown document, 73 lines, 0 explicit URLs. Lead heading: Session Consolidation Master CI Repair. & \textbf{\path{docs/audit/2026-03-18-zzzz-session-consolidation-master-ci-repair.md}}. Тип: markdown. Размер: 3270 байт. Форма: 73 строк. Документ Markdown, 73 строк, явных URL: 0. Главный заголовок: Session Consolidation Master CI Repair. \\
\textbf{\path{docs/audit/2026-03-18-zzzzz-hushwheel-quality-instrumentation.md}}. Kind: markdown. Size: 8775 bytes. Shape: 165 lines. Markdown document, 165 lines, 0 explicit URLs. Lead heading: Hushwheel Quality Instrumentation. & \textbf{\path{docs/audit/2026-03-18-zzzzz-hushwheel-quality-instrumentation.md}}. Тип: markdown. Размер: 8775 байт. Форма: 165 строк. Документ Markdown, 165 строк, явных URL: 0. Главный заголовок: Hushwheel Quality Instrumentation. \\
\textbf{\path{docs/audit/2026-03-18-zzzzzz-dspy-pipeline-and-live-azure-proof.md}}. Kind: markdown. Size: 5992 bytes. Shape: 116 lines. Markdown document, 116 lines, 1 explicit URLs. Lead heading: DSPy Pipeline And Live Azure Proof. & \textbf{\path{docs/audit/2026-03-18-zzzzzz-dspy-pipeline-and-live-azure-proof.md}}. Тип: markdown. Размер: 5992 байт. Форма: 116 строк. Документ Markdown, 116 строк, явных URL: 1. Главный заголовок: DSPy Pipeline And Live Azure Proof. \\
\textbf{\path{docs/audit/2026-03-18-zzzzzz-hushwheel-retrieval-record-wording.md}}. Kind: markdown. Size: 3042 bytes. Shape: 77 lines. Markdown document, 77 lines, 0 explicit URLs. Lead heading: Hushwheel Retrieval Record Wording. & \textbf{\path{docs/audit/2026-03-18-zzzzzz-hushwheel-retrieval-record-wording.md}}. Тип: markdown. Размер: 3042 байт. Форма: 77 строк. Документ Markdown, 77 строк, явных URL: 0. Главный заголовок: Hushwheel Retrieval Record Wording. \\
\textbf{\path{docs/audit/2026-03-18-zzzzzzz-hushwheel-pdf-side-effect-fix.md}}. Kind: markdown. Size: 4044 bytes. Shape: 83 lines. Markdown document, 83 lines, 0 explicit URLs. Lead heading: Hushwheel PDF Side-Effect Fix. & \textbf{\path{docs/audit/2026-03-18-zzzzzzz-hushwheel-pdf-side-effect-fix.md}}. Тип: markdown. Размер: 4044 байт. Форма: 83 строк. Документ Markdown, 83 строк, явных URL: 0. Главный заголовок: Hushwheel PDF Side-Effect Fix. \\
\textbf{\path{docs/audit/2026-03-18-zzzzzzz-hushwheel-quality-salvage.md}}. Kind: markdown. Size: 4554 bytes. Shape: 84 lines. Markdown document, 84 lines, 0 explicit URLs. Lead heading: Hushwheel Quality Salvage Audit. & \textbf{\path{docs/audit/2026-03-18-zzzzzzz-hushwheel-quality-salvage.md}}. Тип: markdown. Размер: 4554 байт. Форма: 84 строк. Документ Markdown, 84 строк, явных URL: 0. Главный заголовок: Hushwheel Quality Salvage Audit. \\
\textbf{\path{docs/audit/2026-03-18-zzzzzzz-land-unlanded-surfaces.md}}. Kind: markdown. Size: 2685 bytes. Shape: 72 lines. Markdown document, 72 lines, 0 explicit URLs. Lead heading: Hushwheel Quality Landing Audit. & \textbf{\path{docs/audit/2026-03-18-zzzzzzz-land-unlanded-surfaces.md}}. Тип: markdown. Размер: 2685 байт. Форма: 72 строк. Документ Markdown, 72 строк, явных URL: 0. Главный заголовок: Hushwheel Quality Landing Audit. \\
\textbf{\path{docs/audit/2026-03-18-zzzzzzz-rebased-master-log-and-azure-typecheck-fix.md}}. Kind: markdown. Size: 3472 bytes. Shape: 83 lines. Markdown document, 83 lines, 0 explicit URLs. Lead heading: Rebased Master Log And Azure Typecheck Fix. & \textbf{\path{docs/audit/2026-03-18-zzzzzzz-rebased-master-log-and-azure-typecheck-fix.md}}. Тип: markdown. Размер: 3472 байт. Форма: 83 строк. Документ Markdown, 83 строк, явных URL: 0. Главный заголовок: Rebased Master Log And Azure Typecheck Fix. \\
\textbf{\path{docs/audit/2026-03-18-zzzzzzz-worktree-consolidation.md}}. Kind: markdown. Size: 5914 bytes. Shape: 125 lines. Markdown document, 125 lines, 0 explicit URLs. Lead heading: Worktree Consolidation. & \textbf{\path{docs/audit/2026-03-18-zzzzzzz-worktree-consolidation.md}}. Тип: markdown. Размер: 5914 байт. Форма: 125 строк. Документ Markdown, 125 строк, явных URL: 0. Главный заголовок: Worktree Consolidation. \\
\textbf{\path{docs/audit/2026-03-18-zzzzzzzzzzzz-retrieval-regression-gate.md}}. Kind: markdown. Size: 2898 bytes. Shape: 66 lines. Markdown document, 66 lines, 0 explicit URLs. Lead heading: Retrieval Regression Gate Audit. & \textbf{\path{docs/audit/2026-03-18-zzzzzzzzzzzz-retrieval-regression-gate.md}}. Тип: markdown. Размер: 2898 байт. Форма: 66 строк. Документ Markdown, 66 строк, явных URL: 0. Главный заголовок: Retrieval Regression Gate Audit. \\
\textbf{\path{docs/audit/2026-03-18-zzzzzzzzzzzzz-land-remaining-unlanded-work.md}}. Kind: markdown. Size: 2908 bytes. Shape: 78 lines. Markdown document, 78 lines, 0 explicit URLs. Lead heading: Land Remaining Unlanded Work. & \textbf{\path{docs/audit/2026-03-18-zzzzzzzzzzzzz-land-remaining-unlanded-work.md}}. Тип: markdown. Размер: 2908 байт. Форма: 78 строк. Документ Markdown, 78 строк, явных URL: 0. Главный заголовок: Land Remaining Unlanded Work. \\
\textbf{\path{docs/audit/2026-03-18-zzzzzzzzzzzzzz-node24-actions-refresh.md}}. Kind: markdown. Size: 2989 bytes. Shape: 91 lines. Markdown document, 91 lines, 0 explicit URLs. Lead heading: Node 24 Actions Refresh. & \textbf{\path{docs/audit/2026-03-18-zzzzzzzzzzzzzz-node24-actions-refresh.md}}. Тип: markdown. Размер: 2989 байт. Форма: 91 строк. Документ Markdown, 91 строк, явных URL: 0. Главный заголовок: Node 24 Actions Refresh. \\
\textbf{\path{docs/audit/2026-03-18-zzzzzzzzzzzzzz-rebase-hushwheel-quality-over-master.md}}. Kind: markdown. Size: 1889 bytes. Shape: 47 lines. Markdown document, 47 lines, 0 explicit URLs. Lead heading: Rebase Hushwheel Quality Over Master. & \textbf{\path{docs/audit/2026-03-18-zzzzzzzzzzzzzz-rebase-hushwheel-quality-over-master.md}}. Тип: markdown. Размер: 1889 байт. Форма: 47 строк. Документ Markdown, 47 строк, явных URL: 0. Главный заголовок: Rebase Hushwheel Quality Over Master. \\
\textbf{\path{docs/audit/2026-03-18-zzzzzzzzzzzzzzzz-dspy-backlog-reconciliation.md}}. Kind: markdown. Size: 2956 bytes. Shape: 75 lines. Markdown document, 75 lines, 0 explicit URLs. Lead heading: DSPy Backlog Reconciliation. & \textbf{\path{docs/audit/2026-03-18-zzzzzzzzzzzzzzzz-dspy-backlog-reconciliation.md}}. Тип: markdown. Размер: 2956 байт. Форма: 75 строк. Документ Markdown, 75 строк, явных URL: 0. Главный заголовок: DSPy Backlog Reconciliation. \\
\textbf{\path{docs/audit/2026-03-18-zzzzzzzzzzzzzzzz-file-summary-test-churn-fix.md}}. Kind: markdown. Size: 2929 bytes. Shape: 85 lines. Markdown document, 85 lines, 0 explicit URLs. Lead heading: File Summary Test Churn Fix. & \textbf{\path{docs/audit/2026-03-18-zzzzzzzzzzzzzzzz-file-summary-test-churn-fix.md}}. Тип: markdown. Размер: 2929 байт. Форма: 85 строк. Документ Markdown, 85 строк, явных URL: 0. Главный заголовок: File Summary Test Churn Fix. \\
\textbf{\path{docs/audit/2026-03-18-zzzzzzzzzzzzzzzz-publication-cache-node24-followup.md}}. Kind: markdown. Size: 2900 bytes. Shape: 89 lines. Markdown document, 89 lines, 0 explicit URLs. Lead heading: Publication Cache Node 24 Follow-Up. & \textbf{\path{docs/audit/2026-03-18-zzzzzzzzzzzzzzzz-publication-cache-node24-followup.md}}. Тип: markdown. Размер: 2900 байт. Форма: 89 строк. Документ Markdown, 89 строк, явных URL: 0. Главный заголовок: Publication Cache Node 24 Follow-Up. \\
\textbf{\path{docs/audit/2026-03-18-zzzzzzzzzzzzzzzz-retrieval-gate-restored-after-rebase.md}}. Kind: markdown. Size: 3788 bytes. Shape: 94 lines. Markdown document, 94 lines, 0 explicit URLs. Lead heading: Retrieval Gate Restored After Rebase. & \textbf{\path{docs/audit/2026-03-18-zzzzzzzzzzzzzzzz-retrieval-gate-restored-after-rebase.md}}. Тип: markdown. Размер: 3788 байт. Форма: 94 строк. Документ Markdown, 94 строк, явных URL: 0. Главный заголовок: Retrieval Gate Restored After Rebase. \\
\textbf{\path{docs/audit/2026-03-18-zzzzzzzzzzzzzzzzzz-dspy-backlog-reconciliation.md}}. Kind: markdown. Size: 2457 bytes. Shape: 62 lines. Markdown document, 62 lines, 0 explicit URLs. Lead heading: DSPy Backlog Reconciliation. & \textbf{\path{docs/audit/2026-03-18-zzzzzzzzzzzzzzzzzz-dspy-backlog-reconciliation.md}}. Тип: markdown. Размер: 2457 байт. Форма: 62 строк. Документ Markdown, 62 строк, явных URL: 0. Главный заголовок: DSPy Backlog Reconciliation. \\
\textbf{\path{docs/audit/2026-03-18-zzzzzzzzzzzzzzzzzz-lookup-first-ask-path.md}}. Kind: markdown. Size: 4604 bytes. Shape: 94 lines. Markdown document, 94 lines, 0 explicit URLs. Lead heading: Lookup First Ask Path. & \textbf{\path{docs/audit/2026-03-18-zzzzzzzzzzzzzzzzzz-lookup-first-ask-path.md}}. Тип: markdown. Размер: 4604 байт. Форма: 94 строк. Документ Markdown, 94 строк, явных URL: 0. Главный заголовок: Lookup First Ask Path. \\
\textbf{\path{docs/audit/2026-03-18-zzzzzzzzzzzzzzzzzzz-answer-synthesis-refresh.md}}. Kind: markdown. Size: 2565 bytes. Shape: 48 lines. Markdown document, 48 lines, 0 explicit URLs. Lead heading: 2026-03-18 Answer Synthesis Refresh. & \textbf{\path{docs/audit/2026-03-18-zzzzzzzzzzzzzzzzzzz-answer-synthesis-refresh.md}}. Тип: markdown. Размер: 2565 байт. Форма: 48 строк. Документ Markdown, 48 строк, явных URL: 0. Главный заголовок: 2026-03-18 Answer Synthesis Refresh. \\
\textbf{\path{docs/audit/2026-03-18-zzzzzzzzzzzzzzzzzzzz-retrieval-quality-doc-priority.md}}. Kind: markdown. Size: 3602 bytes. Shape: 68 lines. Markdown document, 68 lines, 0 explicit URLs. Lead heading: Retrieval Quality Doc Priority Audit. & \textbf{\path{docs/audit/2026-03-18-zzzzzzzzzzzzzzzzzzzz-retrieval-quality-doc-priority.md}}. Тип: markdown. Размер: 3602 байт. Форма: 68 строк. Документ Markdown, 68 строк, явных URL: 0. Главный заголовок: Retrieval Quality Doc Priority Audit. \\
\textbf{\path{docs/audit/2026-03-18-zzzzzzzzzzzzzzzzzzzzzzz-merge-hushwheel-quality-into-master.md}}. Kind: markdown. Size: 4335 bytes. Shape: 112 lines. Markdown document, 112 lines, 0 explicit URLs. Lead heading: Merge Hushwheel Quality Into Master. & \textbf{\path{docs/audit/2026-03-18-zzzzzzzzzzzzzzzzzzzzzzz-merge-hushwheel-quality-into-master.md}}. Тип: markdown. Размер: 4335 байт. Форма: 112 строк. Документ Markdown, 112 строк, явных URL: 0. Главный заголовок: Merge Hushwheel Quality Into Master. \\
\textbf{\path{docs/audit/2026-03-18-zzzzzzzzzzzzzzzzzzzzzzzz-hushwheel-atlas-consolidation.md}}. Kind: markdown. Size: 4762 bytes. Shape: 109 lines. Markdown document, 109 lines, 0 explicit URLs. Lead heading: Hushwheel Atlas Consolidation. & \textbf{\path{docs/audit/2026-03-18-zzzzzzzzzzzzzzzzzzzzzzzz-hushwheel-atlas-consolidation.md}}. Тип: markdown. Размер: 4762 байт. Форма: 109 строк. Документ Markdown, 109 строк, явных URL: 0. Главный заголовок: Hushwheel Atlas Consolidation. \\
\textbf{\path{docs/audit/2026-03-19-github-pr-gate-sync.md}}. Kind: markdown. Size: 6599 bytes. Shape: 136 lines. Markdown document, 136 lines, 0 explicit URLs. Lead heading: Master Consolidation And GitHub PR Gate Sync. & \textbf{\path{docs/audit/2026-03-19-github-pr-gate-sync.md}}. Тип: markdown. Размер: 6599 байт. Форма: 136 строк. Документ Markdown, 136 строк, явных URL: 0. Главный заголовок: Master Consolidation And GitHub PR Gate Sync. \\
\textbf{\path{docs/audit/2026-03-19-node24-actions-opt-in.md}}. Kind: markdown. Size: 2929 bytes. Shape: 80 lines. Markdown document, 80 lines, 0 explicit URLs. Lead heading: Node 24 Actions Opt-In. & \textbf{\path{docs/audit/2026-03-19-node24-actions-opt-in.md}}. Тип: markdown. Размер: 2929 байт. Форма: 80 строк. Документ Markdown, 80 строк, явных URL: 0. Главный заголовок: Node 24 Actions Opt-In. \\
\textbf{\path{docs/audit/2026-03-19-public-pages-and-secret-scan.md}}. Kind: markdown. Size: 5804 bytes. Shape: 113 lines. Markdown document, 113 lines, 2 explicit URLs. Lead heading: Public Pages And Secret Scan. & \textbf{\path{docs/audit/2026-03-19-public-pages-and-secret-scan.md}}. Тип: markdown. Размер: 5804 байт. Форма: 113 строк. Документ Markdown, 113 строк, явных URL: 2. Главный заголовок: Public Pages And Secret Scan. \\
\textbf{\path{docs/audit/2026-03-19-simple-end-to-end-verification-guide.md}}. Kind: markdown. Size: 4236 bytes. Shape: 103 lines. Markdown document, 103 lines, 0 explicit URLs. Lead heading: Simple End-To-End Verification Guide. & \textbf{\path{docs/audit/2026-03-19-simple-end-to-end-verification-guide.md}}. Тип: markdown. Размер: 4236 байт. Форма: 103 строк. Документ Markdown, 103 строк, явных URL: 0. Главный заголовок: Simple End-To-End Verification Guide. \\
\textbf{\path{docs/audit/2026-04-28-cli-envelope-normalization.md}}. Kind: markdown. Size: 4028 bytes. Shape: 95 lines. Markdown document, 95 lines, 0 explicit URLs. Lead heading: CLI Envelope Normalization. & \textbf{\path{docs/audit/2026-04-28-cli-envelope-normalization.md}}. Тип: markdown. Размер: 4028 байт. Форма: 95 строк. Документ Markdown, 95 строк, явных URL: 0. Главный заголовок: CLI Envelope Normalization. \\
\textbf{\path{docs/audit/2026-04-28-doc-tree-restructure-and-planning.md}}. Kind: markdown. Size: 4309 bytes. Shape: 105 lines. Markdown document, 105 lines, 0 explicit URLs. Lead heading: Doc Tree Restructure And Planning. & \textbf{\path{docs/audit/2026-04-28-doc-tree-restructure-and-planning.md}}. Тип: markdown. Размер: 4309 байт. Форма: 105 строк. Документ Markdown, 105 строк, явных URL: 0. Главный заголовок: Doc Tree Restructure And Planning. \\
\textbf{\path{docs/audit/2026-04-28-json-runtime-contract.md}}. Kind: markdown. Size: 5281 bytes. Shape: 131 lines. Markdown document, 131 lines, 0 explicit URLs. Lead heading: JSON Runtime Contract. & \textbf{\path{docs/audit/2026-04-28-json-runtime-contract.md}}. Тип: markdown. Размер: 5281 байт. Форма: 131 строк. Документ Markdown, 131 строк, явных URL: 0. Главный заголовок: JSON Runtime Contract. \\
\textbf{\path{docs/audit/2026-04-28-live-azure-dspy-integration-gating.md}}. Kind: markdown. Size: 3562 bytes. Shape: 79 lines. Markdown document, 79 lines, 0 explicit URLs. Lead heading: Live Azure And DSPy Integration Gating. & \textbf{\path{docs/audit/2026-04-28-live-azure-dspy-integration-gating.md}}. Тип: markdown. Размер: 3562 байт. Форма: 79 строк. Документ Markdown, 79 строк, явных URL: 0. Главный заголовок: Live Azure And DSPy Integration Gating. \\
\textbf{\path{docs/audit/2026-04-28-retrieval-profile-layer.md}}. Kind: markdown. Size: 4123 bytes. Shape: 92 lines. Markdown document, 92 lines, 0 explicit URLs. Lead heading: Retrieval Profile Layer. & \textbf{\path{docs/audit/2026-04-28-retrieval-profile-layer.md}}. Тип: markdown. Размер: 4123 байт. Форма: 92 строк. Документ Markdown, 92 строк, явных URL: 0. Главный заголовок: Retrieval Profile Layer. \\
\textbf{\path{docs/audit/2026-04-28-rust-lookup-arbitrary-repo-roots.md}}. Kind: markdown. Size: 4078 bytes. Shape: 88 lines. Markdown document, 88 lines, 0 explicit URLs. Lead heading: Rust Lookup Arbitrary Repo Roots. & \textbf{\path{docs/audit/2026-04-28-rust-lookup-arbitrary-repo-roots.md}}. Тип: markdown. Размер: 4078 байт. Форма: 88 строк. Документ Markdown, 88 строк, явных URL: 0. Главный заголовок: Rust Lookup Arbitrary Repo Roots. \\
\textbf{\path{docs/audit/2026-04-29-bounded-mcp-server-surface.md}}. Kind: markdown. Size: 6688 bytes. Shape: 154 lines. Markdown document, 154 lines, 0 explicit URLs. Lead heading: 2026-04-29 Bounded MCP Server Surface. & \textbf{\path{docs/audit/2026-04-29-bounded-mcp-server-surface.md}}. Тип: markdown. Размер: 6688 байт. Форма: 154 строк. Документ Markdown, 154 строк, явных URL: 0. Главный заголовок: 2026-04-29 Bounded MCP Server Surface. \\
\textbf{\path{docs/audit/2026-04-29-bundle-channel-publish-promote-rollback.md}}. Kind: markdown. Size: 5633 bytes. Shape: 125 lines. Markdown document, 125 lines, 0 explicit URLs. Lead heading: 2026-04-29 Bundle Channel Publish Promote Rollback. & \textbf{\path{docs/audit/2026-04-29-bundle-channel-publish-promote-rollback.md}}. Тип: markdown. Размер: 5633 байт. Форма: 125 строк. Документ Markdown, 125 строк, явных URL: 0. Главный заголовок: 2026-04-29 Bundle Channel Publish Promote Rollback. \\
\textbf{\path{docs/audit/2026-04-29-bundle-overlay-and-runtime-trace-contract.md}}. Kind: markdown. Size: 4332 bytes. Shape: 96 lines. Markdown document, 96 lines, 0 explicit URLs. Lead heading: 2026-04-29 Bundle Overlay And Runtime Trace Contract. & \textbf{\path{docs/audit/2026-04-29-bundle-overlay-and-runtime-trace-contract.md}}. Тип: markdown. Размер: 4332 байт. Форма: 96 строк. Документ Markdown, 96 строк, явных URL: 0. Главный заголовок: 2026-04-29 Bundle Overlay And Runtime Trace Contract. \\
\textbf{\path{docs/audit/2026-04-29-dataset-accepted-outcome-ingest.md}}. Kind: markdown. Size: 6670 bytes. Shape: 139 lines. Markdown document, 139 lines, 0 explicit URLs. Lead heading: 2026-04-29 Dataset Accepted-Outcome Ingest. & \textbf{\path{docs/audit/2026-04-29-dataset-accepted-outcome-ingest.md}}. Тип: markdown. Размер: 6670 байт. Форма: 139 строк. Документ Markdown, 139 строк, явных URL: 0. Главный заголовок: 2026-04-29 Dataset Accepted-Outcome Ingest. \\
\textbf{\path{docs/audit/2026-04-29-dataset-bundle-fetch-and-trace-import-hooks.md}}. Kind: markdown. Size: 6215 bytes. Shape: 130 lines. Markdown document, 130 lines, 0 explicit URLs. Lead heading: 2026-04-29 Dataset Bundle Fetch And Trace Import Hooks. & \textbf{\path{docs/audit/2026-04-29-dataset-bundle-fetch-and-trace-import-hooks.md}}. Тип: markdown. Размер: 6215 байт. Форма: 130 строк. Документ Markdown, 130 строк, явных URL: 0. Главный заголовок: 2026-04-29 Dataset Bundle Fetch And Trace Import Hooks. \\
\textbf{\path{docs/audit/2026-04-29-dataset-local-repo-rag-cli-backend.md}}. Kind: markdown. Size: 5290 bytes. Shape: 122 lines. Markdown document, 122 lines, 0 explicit URLs. Lead heading: 2026-04-29 Dataset Local `repo\_rag\_cli` Backend. & \textbf{\path{docs/audit/2026-04-29-dataset-local-repo-rag-cli-backend.md}}. Тип: markdown. Размер: 5290 байт. Форма: 122 строк. Документ Markdown, 122 строк, явных URL: 0. Главный заголовок: 2026-04-29 Dataset Local `repo\_rag\_cli` Backend. \\
\textbf{\path{docs/audit/2026-04-29-dataset-repo-rag-auto-detection.md}}. Kind: markdown. Size: 6051 bytes. Shape: 124 lines. Markdown document, 124 lines, 0 explicit URLs. Lead heading: 2026-04-29 Dataset Repo-RAG Auto-Detection. & \textbf{\path{docs/audit/2026-04-29-dataset-repo-rag-auto-detection.md}}. Тип: markdown. Размер: 6051 байт. Форма: 124 строк. Документ Markdown, 124 строк, явных URL: 0. Главный заголовок: 2026-04-29 Dataset Repo-RAG Auto-Detection. \\
\textbf{\path{docs/audit/2026-04-29-dataset-trace-queue-handoff.md}}. Kind: markdown. Size: 6656 bytes. Shape: 141 lines. Markdown document, 141 lines, 0 explicit URLs. Lead heading: 2026-04-29 Dataset Trace Queue Handoff. & \textbf{\path{docs/audit/2026-04-29-dataset-trace-queue-handoff.md}}. Тип: markdown. Размер: 6656 байт. Форма: 141 строк. Документ Markdown, 141 строк, явных URL: 0. Главный заголовок: 2026-04-29 Dataset Trace Queue Handoff. \\
\textbf{\path{docs/audit/2026-04-29-dataset-worker-repo-rag-cli-backend.md}}. Kind: markdown. Size: 6630 bytes. Shape: 128 lines. Markdown document, 128 lines, 0 explicit URLs. Lead heading: 2026-04-29 Dataset Worker `repo\_rag\_cli` Backend. & \textbf{\path{docs/audit/2026-04-29-dataset-worker-repo-rag-cli-backend.md}}. Тип: markdown. Размер: 6630 байт. Форма: 128 строк. Документ Markdown, 128 строк, явных URL: 0. Главный заголовок: 2026-04-29 Dataset Worker `repo\_rag\_cli` Backend. \\
\textbf{\path{docs/audit/2026-04-29-global-blob-queue-trainer-migration.md}}. Kind: markdown. Size: 4807 bytes. Shape: 91 lines. Markdown document, 91 lines, 0 explicit URLs. Lead heading: 2026-04-29 Global Blob And Queue Trainer Migration. & \textbf{\path{docs/audit/2026-04-29-global-blob-queue-trainer-migration.md}}. Тип: markdown. Размер: 4807 байт. Форма: 91 строк. Документ Markdown, 91 строк, явных URL: 0. Главный заголовок: 2026-04-29 Global Blob And Queue Trainer Migration. \\
\textbf{\path{docs/audit/2026-04-29-live-azure-gpt54-validation.md}}. Kind: markdown. Size: 8279 bytes. Shape: 145 lines. Markdown document, 145 lines, 6 explicit URLs. Lead heading: 2026-04-29 Live Azure GPT-5.4 Validation. & \textbf{\path{docs/audit/2026-04-29-live-azure-gpt54-validation.md}}. Тип: markdown. Размер: 8279 байт. Форма: 145 строк. Документ Markdown, 145 строк, явных URL: 6. Главный заголовок: 2026-04-29 Live Azure GPT-5.4 Validation. \\
\textbf{\path{docs/audit/2026-04-29-retrieval-rerank-and-relative-root-normalization.md}}. Kind: markdown. Size: 3070 bytes. Shape: 71 lines. Markdown document, 71 lines, 0 explicit URLs. Lead heading: 2026-04-29 Retrieval Rerank And Relative Root Normalization. & \textbf{\path{docs/audit/2026-04-29-retrieval-rerank-and-relative-root-normalization.md}}. Тип: markdown. Размер: 3070 байт. Форма: 71 строк. Документ Markdown, 71 строк, явных URL: 0. Главный заголовок: 2026-04-29 Retrieval Rerank And Relative Root Normalization. \\
\textbf{\path{docs/audit/2026-04-29-runtime-image-and-dataset-build-deploy-wiring.md}}. Kind: markdown. Size: 5729 bytes. Shape: 140 lines. Markdown document, 140 lines, 0 explicit URLs. Lead heading: 2026-04-29 Runtime Image And Dataset Build/Deploy Wiring. & \textbf{\path{docs/audit/2026-04-29-runtime-image-and-dataset-build-deploy-wiring.md}}. Тип: markdown. Размер: 5729 байт. Форма: 140 строк. Документ Markdown, 140 строк, явных URL: 0. Главный заголовок: 2026-04-29 Runtime Image And Dataset Build/Deploy Wiring. \\
\textbf{\path{docs/audit/2026-04-29-trainer-bundle-benchmark-gates.md}}. Kind: markdown. Size: 4942 bytes. Shape: 111 lines. Markdown document, 111 lines, 0 explicit URLs. Lead heading: 2026-04-29 Trainer Bundle Benchmark Gates. & \textbf{\path{docs/audit/2026-04-29-trainer-bundle-benchmark-gates.md}}. Тип: markdown. Размер: 4942 байт. Форма: 111 строк. Документ Markdown, 111 строк, явных URL: 0. Главный заголовок: 2026-04-29 Trainer Bundle Benchmark Gates. \\
\textbf{\path{docs/audit/2026-04-29-trainer-cycle-background-pass.md}}. Kind: markdown. Size: 4637 bytes. Shape: 105 lines. Markdown document, 105 lines, 0 explicit URLs. Lead heading: 2026-04-29 Trainer Cycle Background Pass. & \textbf{\path{docs/audit/2026-04-29-trainer-cycle-background-pass.md}}. Тип: markdown. Размер: 4637 байт. Форма: 105 строк. Документ Markdown, 105 строк, явных URL: 0. Главный заголовок: 2026-04-29 Trainer Cycle Background Pass. \\
\textbf{\path{docs/audit/2026-04-29-trainer-k8s-deployment-surface.md}}. Kind: markdown. Size: 6040 bytes. Shape: 131 lines. Markdown document, 131 lines, 0 explicit URLs. Lead heading: 2026-04-29 Trainer Kubernetes Deployment Surface. & \textbf{\path{docs/audit/2026-04-29-trainer-k8s-deployment-surface.md}}. Тип: markdown. Размер: 6040 байт. Форма: 131 строк. Документ Markdown, 131 строк, явных URL: 0. Главный заголовок: 2026-04-29 Trainer Kubernetes Deployment Surface. \\
\textbf{\path{docs/audit/2026-04-29-trainer-recompile-from-candidates.md}}. Kind: markdown. Size: 6238 bytes. Shape: 135 lines. Markdown document, 135 lines, 0 explicit URLs. Lead heading: 2026-04-29 Trainer Recompile From Candidates. & \textbf{\path{docs/audit/2026-04-29-trainer-recompile-from-candidates.md}}. Тип: markdown. Размер: 6238 байт. Форма: 135 строк. Документ Markdown, 135 строк, явных URL: 0. Главный заголовок: 2026-04-29 Trainer Recompile From Candidates. \\
\textbf{\path{docs/audit/2026-04-29-trainer-service-loop.md}}. Kind: markdown. Size: 6185 bytes. Shape: 133 lines. Markdown document, 133 lines, 0 explicit URLs. Lead heading: 2026-04-29 Trainer Service Loop. & \textbf{\path{docs/audit/2026-04-29-trainer-service-loop.md}}. Тип: markdown. Размер: 6185 байт. Форма: 133 строк. Документ Markdown, 133 строк, явных URL: 0. Главный заголовок: 2026-04-29 Trainer Service Loop. \\
\textbf{\path{docs/audit/2026-04-30-aks-run-25179828865-stale-image-inspection.md}}. Kind: markdown. Size: 6206 bytes. Shape: 117 lines. Markdown document, 117 lines, 0 explicit URLs. Lead heading: 2026-04-30 AKS Run 25179828865 Stale Image Inspection. & \textbf{\path{docs/audit/2026-04-30-aks-run-25179828865-stale-image-inspection.md}}. Тип: markdown. Размер: 6206 байт. Форма: 117 строк. Документ Markdown, 117 строк, явных URL: 0. Главный заголовок: 2026-04-30 AKS Run 25179828865 Stale Image Inspection. \\
\textbf{\path{docs/audit/2026-04-30-aks-run-25186264345-image-fix-gap-inspection.md}}. Kind: markdown. Size: 8407 bytes. Shape: 165 lines. Markdown document, 165 lines, 0 explicit URLs. Lead heading: 2026-04-30 AKS Run 25186264345 Image/Fix Gap Inspection. & \textbf{\path{docs/audit/2026-04-30-aks-run-25186264345-image-fix-gap-inspection.md}}. Тип: markdown. Размер: 8407 байт. Форма: 165 строк. Документ Markdown, 165 строк, явных URL: 0. Главный заголовок: 2026-04-30 AKS Run 25186264345 Image/Fix Gap Inspection. \\
\textbf{\path{docs/audit/2026-04-30-codex-proxy-mediation-layer.md}}. Kind: markdown. Size: 5279 bytes. Shape: 95 lines. Markdown document, 95 lines, 0 explicit URLs. Lead heading: 2026-04-30 Codex Proxy Mediation Layer. & \textbf{\path{docs/audit/2026-04-30-codex-proxy-mediation-layer.md}}. Тип: markdown. Размер: 5279 байт. Форма: 95 строк. Документ Markdown, 95 строк, явных URL: 0. Главный заголовок: 2026-04-30 Codex Proxy Mediation Layer. \\
\textbf{\path{docs/audit/2026-04-30-codex-proxy-trace-handoff.md}}. Kind: markdown. Size: 6543 bytes. Shape: 122 lines. Markdown document, 122 lines, 0 explicit URLs. Lead heading: 2026-04-30 Codex Proxy Trace Handoff. & \textbf{\path{docs/audit/2026-04-30-codex-proxy-trace-handoff.md}}. Тип: markdown. Размер: 6543 байт. Форма: 122 строк. Документ Markdown, 122 строк, явных URL: 0. Главный заголовок: 2026-04-30 Codex Proxy Trace Handoff. \\
\textbf{\path{docs/audit/2026-05-01-aks-run-25191504814-project-root-resolution-inspection.md}}. Kind: markdown. Size: 6002 bytes. Shape: 133 lines. Markdown document, 133 lines, 0 explicit URLs. Lead heading: 2026-05-01 AKS Run 25191504814 Project Root Resolution Inspection. & \textbf{\path{docs/audit/2026-05-01-aks-run-25191504814-project-root-resolution-inspection.md}}. Тип: markdown. Размер: 6002 байт. Форма: 133 строк. Документ Markdown, 133 строк, явных URL: 0. Главный заголовок: 2026-05-01 AKS Run 25191504814 Project Root Resolution Inspection. \\
\textbf{\path{docs/audit/2026-05-01-aks-run-25209573387-proxy-runtime-env-inspection.md}}. Kind: markdown. Size: 5738 bytes. Shape: 129 lines. Markdown document, 129 lines, 0 explicit URLs. Lead heading: 2026-05-01 AKS Run 25209573387 Proxy Runtime Env Inspection. & \textbf{\path{docs/audit/2026-05-01-aks-run-25209573387-proxy-runtime-env-inspection.md}}. Тип: markdown. Размер: 5738 байт. Форма: 129 строк. Документ Markdown, 129 строк, явных URL: 0. Главный заголовок: 2026-05-01 AKS Run 25209573387 Proxy Runtime Env Inspection. \\
\textbf{\path{docs/audit/2026-05-01-aks-run-25212955759-trace-handoff-secret-gap.md}}. Kind: markdown. Size: 4594 bytes. Shape: 121 lines. Markdown document, 121 lines, 0 explicit URLs. Lead heading: AKS Run `25212955759`: Trace Handoff Skipped Because `repo-rag-storage-config` Lacks Blob Credentials. & \textbf{\path{docs/audit/2026-05-01-aks-run-25212955759-trace-handoff-secret-gap.md}}. Тип: markdown. Размер: 4594 байт. Форма: 121 строк. Документ Markdown, 121 строк, явных URL: 0. Главный заголовок: AKS Run `25212955759`: Trace Handoff Skipped Because `repo-rag-storage-config` Lacks Blob Credentials. \\
\textbf{\path{docs/audit/2026-05-01-aks-run-25226558751-rag-heuristic-trusted-handoff.md}}. Kind: markdown. Size: 5092 bytes. Shape: 130 lines. Markdown document, 130 lines, 0 explicit URLs. Lead heading: AKS Run 25226558751: RAG Success, Heuristic DSPy Fallback, Trusted Handoff Success. & \textbf{\path{docs/audit/2026-05-01-aks-run-25226558751-rag-heuristic-trusted-handoff.md}}. Тип: markdown. Размер: 5092 байт. Форма: 130 строк. Документ Markdown, 130 строк, явных URL: 0. Главный заголовок: AKS Run 25226558751: RAG Success, Heuristic DSPy Fallback, Trusted Handoff Success. \\
\textbf{\path{docs/audit/2026-05-01-dataset-azure-responses-api-version-default-fix.md}}. Kind: markdown. Size: 6700 bytes. Shape: 133 lines. Markdown document, 133 lines, 1 explicit URLs. Lead heading: 2026-05-01 Dataset Azure Responses API-Version Default Fix. & \textbf{\path{docs/audit/2026-05-01-dataset-azure-responses-api-version-default-fix.md}}. Тип: markdown. Размер: 6700 байт. Форма: 133 строк. Документ Markdown, 133 строк, явных URL: 1. Главный заголовок: 2026-05-01 Dataset Azure Responses API-Version Default Fix. \\
\textbf{\path{docs/audit/2026-05-01-dataset-azure-runtime-env-propagation-fix.md}}. Kind: markdown. Size: 4200 bytes. Shape: 81 lines. Markdown document, 81 lines, 0 explicit URLs. Lead heading: 2026-05-01 Dataset Azure Runtime Env Propagation Fix. & \textbf{\path{docs/audit/2026-05-01-dataset-azure-runtime-env-propagation-fix.md}}. Тип: markdown. Размер: 4200 байт. Форма: 81 строк. Документ Markdown, 81 строк, явных URL: 0. Главный заголовок: 2026-05-01 Dataset Azure Runtime Env Propagation Fix. \\
\textbf{\path{docs/audit/2026-05-01-dataset-trusted-trace-handoff-postprocess-fix.md}}. Kind: markdown. Size: 3880 bytes. Shape: 81 lines. Markdown document, 81 lines, 0 explicit URLs. Lead heading: Dataset Trusted Trace Handoff Post-Processing Fix. & \textbf{\path{docs/audit/2026-05-01-dataset-trusted-trace-handoff-postprocess-fix.md}}. Тип: markdown. Размер: 3880 байт. Форма: 81 строк. Документ Markdown, 81 строк, явных URL: 0. Главный заголовок: Dataset Trusted Trace Handoff Post-Processing Fix. \\
\textbf{\path{docs/audit/2026-05-01-dspy-bundle-version-pinning-contract.md}}. Kind: markdown. Size: 5763 bytes. Shape: 118 lines. Markdown document, 118 lines, 0 explicit URLs. Lead heading: DSPY Bundle Version Pinning Contract. & \textbf{\path{docs/audit/2026-05-01-dspy-bundle-version-pinning-contract.md}}. Тип: markdown. Размер: 5763 байт. Форма: 118 строк. Документ Markdown, 118 строк, явных URL: 0. Главный заголовок: DSPY Bundle Version Pinning Contract. \\
\textbf{\path{docs/audit/2026-05-01-durable-trainer-progress-remote-recovery.md}}. Kind: markdown. Size: 5141 bytes. Shape: 115 lines. Markdown document, 115 lines, 0 explicit URLs. Lead heading: Durable Trainer Progress Remote Recovery. & \textbf{\path{docs/audit/2026-05-01-durable-trainer-progress-remote-recovery.md}}. Тип: markdown. Размер: 5141 байт. Форма: 115 строк. Документ Markdown, 115 строк, явных URL: 0. Главный заголовок: Durable Trainer Progress Remote Recovery. \\
\textbf{\path{docs/audit/2026-05-01-live-trainer-no-timestamp-bundle-publish.md}}. Kind: markdown. Size: 7368 bytes. Shape: 172 lines. Markdown document, 172 lines, 0 explicit URLs. Lead heading: Live Trainer No Timestamp Bundle Publish. & \textbf{\path{docs/audit/2026-05-01-live-trainer-no-timestamp-bundle-publish.md}}. Тип: markdown. Размер: 7368 байт. Форма: 172 строк. Документ Markdown, 172 строк, явных URL: 0. Главный заголовок: Live Trainer No Timestamp Bundle Publish. \\
\textbf{\path{docs/audit/2026-05-01-live-trainer-remote-bundle-publish.md}}. Kind: markdown. Size: 7444 bytes. Shape: 148 lines. Markdown document, 148 lines, 0 explicit URLs. Lead heading: Live Trainer Remote Bundle Publish. & \textbf{\path{docs/audit/2026-05-01-live-trainer-remote-bundle-publish.md}}. Тип: markdown. Размер: 7444 байт. Форма: 148 строк. Документ Markdown, 148 строк, явных URL: 0. Главный заголовок: Live Trainer Remote Bundle Publish. \\
\textbf{\path{docs/audit/2026-05-01-live-trainer-service-azure-fallback.md}}. Kind: markdown. Size: 6390 bytes. Shape: 167 lines. Markdown document, 167 lines, 0 explicit URLs. Lead heading: Live Trainer Service Azure Fallback. & \textbf{\path{docs/audit/2026-05-01-live-trainer-service-azure-fallback.md}}. Тип: markdown. Размер: 6390 байт. Форма: 167 строк. Документ Markdown, 167 строк, явных URL: 0. Главный заголовок: Live Trainer Service Azure Fallback. \\
\textbf{\path{docs/audit/2026-05-01-versioned-trainer-lineage-bundles.md}}. Kind: markdown. Size: 5821 bytes. Shape: 123 lines. Markdown document, 123 lines, 0 explicit URLs. Lead heading: Versioned Trainer Lineage Bundles. & \textbf{\path{docs/audit/2026-05-01-versioned-trainer-lineage-bundles.md}}. Тип: markdown. Размер: 5821 байт. Форма: 123 строк. Документ Markdown, 123 строк, явных URL: 0. Главный заголовок: Versioned Trainer Lineage Bundles. \\
\textbf{\path{docs/audit/2026-05-02-hybrid-vector-retrieval-default.md}}. Kind: markdown. Size: 18302 bytes. Shape: 385 lines. Markdown document, 385 lines, 0 explicit URLs. Lead heading: Hybrid Vector Retrieval Default. & \textbf{\path{docs/audit/2026-05-02-hybrid-vector-retrieval-default.md}}. Тип: markdown. Размер: 18302 байт. Форма: 385 строк. Документ Markdown, 385 строк, явных URL: 0. Главный заголовок: Hybrid Vector Retrieval Default. \\
\textbf{\path{docs/audit/2026-05-02-live-trainer-still-not-publishing-bundles.md}}. Kind: markdown. Size: 60556 bytes. Shape: 1281 lines. Markdown document, 1281 lines, 0 explicit URLs. Lead heading: Live Trainer Still Not Publishing Bundles. & \textbf{\path{docs/audit/2026-05-02-live-trainer-still-not-publishing-bundles.md}}. Тип: markdown. Размер: 60556 байт. Форма: 1281 строк. Документ Markdown, 1281 строк, явных URL: 0. Главный заголовок: Live Trainer Still Not Publishing Bundles. \\
\textbf{\path{docs/audit/2026-05-02-trainer-context-group-champion-index-stage1.md}}. Kind: markdown. Size: 18201 bytes. Shape: 330 lines. Markdown document, 330 lines, 0 explicit URLs. Lead heading: Trainer Context-Group Champion Index Stage 1. & \textbf{\path{docs/audit/2026-05-02-trainer-context-group-champion-index-stage1.md}}. Тип: markdown. Размер: 18201 байт. Форма: 330 строк. Документ Markdown, 330 строк, явных URL: 0. Главный заголовок: Trainer Context-Group Champion Index Stage 1. \\
\textbf{\path{docs/audit/2026-05-02-zz-codex-exec-resume-session-state-stage0.md}}. Kind: markdown. Size: 170338 bytes. Shape: 3866 lines. Markdown document, 3866 lines, 2 explicit URLs. Lead heading: Repository audit note for 2026-05-02 codex exec resume session state stage 0. & \textbf{\path{docs/audit/2026-05-02-zz-codex-exec-resume-session-state-stage0.md}}. Тип: markdown. Размер: 170338 байт. Форма: 3866 строк. Документ Markdown, 3866 строк, явных URL: 2. Главный заголовок: Repository audit note for 2026-05-02 codex exec resume session state stage 0. \\
\textbf{\path{docs/audit/2026-05-08-zz-per-turn-dspy-mediation-contract-stage0.md}}. Kind: markdown. Size: 5821 bytes. Shape: 113 lines. Markdown document, 113 lines, 0 explicit URLs. Lead heading: Repository audit note for 2026-05-08 per-turn DSPy mediation contract stage 0. & \textbf{\path{docs/audit/2026-05-08-zz-per-turn-dspy-mediation-contract-stage0.md}}. Тип: markdown. Размер: 5821 байт. Форма: 113 строк. Документ Markdown, 113 строк, явных URL: 0. Главный заголовок: Repository audit note for 2026-05-08 per-turn DSPy mediation contract stage 0. \\
\textbf{\path{docs/audit/2026-05-09-aks-run-25598675829-per-turn-mediation-gap.md}}. Kind: markdown. Size: 7119 bytes. Shape: 215 lines. Markdown document, 215 lines, 0 explicit URLs. Lead heading: Repository audit note for 2026-05-09 AKS run 25598675829 per-turn mediation gap. & \textbf{\path{docs/audit/2026-05-09-aks-run-25598675829-per-turn-mediation-gap.md}}. Тип: markdown. Размер: 7119 байт. Форма: 215 строк. Документ Markdown, 215 строк, явных URL: 0. Главный заголовок: Repository audit note for 2026-05-09 AKS run 25598675829 per-turn mediation gap. \\
\textbf{\path{docs/audit/2026-05-09-family-first-mipro-contract-stage1.md}}. Kind: markdown. Size: 4684 bytes. Shape: 120 lines. Markdown document, 120 lines, 0 explicit URLs. Lead heading: Repository audit note for 2026-05-09 family-first MIPROv2 contract stage 1. & \textbf{\path{docs/audit/2026-05-09-family-first-mipro-contract-stage1.md}}. Тип: markdown. Размер: 4684 байт. Форма: 120 строк. Документ Markdown, 120 строк, явных URL: 0. Главный заголовок: Repository audit note for 2026-05-09 family-first MIPROv2 contract stage 1. \\
\textbf{\path{docs/audit/2026-05-09-family-first-mipro-contract-stage10.md}}. Kind: markdown. Size: 5329 bytes. Shape: 122 lines. Markdown document, 122 lines, 0 explicit URLs. Lead heading: Repository audit note for 2026-05-09 family-first MIPROv2 contract stage 10. & \textbf{\path{docs/audit/2026-05-09-family-first-mipro-contract-stage10.md}}. Тип: markdown. Размер: 5329 байт. Форма: 122 строк. Документ Markdown, 122 строк, явных URL: 0. Главный заголовок: Repository audit note for 2026-05-09 family-first MIPROv2 contract stage 10. \\
\textbf{\path{docs/audit/2026-05-09-family-first-mipro-contract-stage11.md}}. Kind: markdown. Size: 4638 bytes. Shape: 116 lines. Markdown document, 116 lines, 0 explicit URLs. Lead heading: Repository audit note for 2026-05-09 family-first MIPROv2 contract stage 11. & \textbf{\path{docs/audit/2026-05-09-family-first-mipro-contract-stage11.md}}. Тип: markdown. Размер: 4638 байт. Форма: 116 строк. Документ Markdown, 116 строк, явных URL: 0. Главный заголовок: Repository audit note for 2026-05-09 family-first MIPROv2 contract stage 11. \\
\textbf{\path{docs/audit/2026-05-09-family-first-mipro-contract-stage12.md}}. Kind: markdown. Size: 4941 bytes. Shape: 122 lines. Markdown document, 122 lines, 0 explicit URLs. Lead heading: Repository audit note for 2026-05-09 family-first MIPROv2 contract stage 12. & \textbf{\path{docs/audit/2026-05-09-family-first-mipro-contract-stage12.md}}. Тип: markdown. Размер: 4941 байт. Форма: 122 строк. Документ Markdown, 122 строк, явных URL: 0. Главный заголовок: Repository audit note for 2026-05-09 family-first MIPROv2 contract stage 12. \\
\textbf{\path{docs/audit/2026-05-09-family-first-mipro-contract-stage13.md}}. Kind: markdown. Size: 4992 bytes. Shape: 128 lines. Markdown document, 128 lines, 0 explicit URLs. Lead heading: Repository audit note for 2026-05-09 family-first MIPROv2 contract stage 13. & \textbf{\path{docs/audit/2026-05-09-family-first-mipro-contract-stage13.md}}. Тип: markdown. Размер: 4992 байт. Форма: 128 строк. Документ Markdown, 128 строк, явных URL: 0. Главный заголовок: Repository audit note for 2026-05-09 family-first MIPROv2 contract stage 13. \\
\textbf{\path{docs/audit/2026-05-09-family-first-mipro-contract-stage14.md}}. Kind: markdown. Size: 3944 bytes. Shape: 109 lines. Markdown document, 109 lines, 0 explicit URLs. Lead heading: Repository audit note for 2026-05-09 family-first MIPROv2 contract stage 14. & \textbf{\path{docs/audit/2026-05-09-family-first-mipro-contract-stage14.md}}. Тип: markdown. Размер: 3944 байт. Форма: 109 строк. Документ Markdown, 109 строк, явных URL: 0. Главный заголовок: Repository audit note for 2026-05-09 family-first MIPROv2 contract stage 14. \\
\textbf{\path{docs/audit/2026-05-09-family-first-mipro-contract-stage15.md}}. Kind: markdown. Size: 4728 bytes. Shape: 113 lines. Markdown document, 113 lines, 0 explicit URLs. Lead heading: Repository audit note for 2026-05-09 family-first MIPROv2 contract stage 15. & \textbf{\path{docs/audit/2026-05-09-family-first-mipro-contract-stage15.md}}. Тип: markdown. Размер: 4728 байт. Форма: 113 строк. Документ Markdown, 113 строк, явных URL: 0. Главный заголовок: Repository audit note for 2026-05-09 family-first MIPROv2 contract stage 15. \\
\textbf{\path{docs/audit/2026-05-09-family-first-mipro-contract-stage16.md}}. Kind: markdown. Size: 3748 bytes. Shape: 103 lines. Markdown document, 103 lines, 0 explicit URLs. Lead heading: Repository audit note for 2026-05-09 family-first MIPROv2 contract stage 16. & \textbf{\path{docs/audit/2026-05-09-family-first-mipro-contract-stage16.md}}. Тип: markdown. Размер: 3748 байт. Форма: 103 строк. Документ Markdown, 103 строк, явных URL: 0. Главный заголовок: Repository audit note for 2026-05-09 family-first MIPROv2 contract stage 16. \\
\textbf{\path{docs/audit/2026-05-09-family-first-mipro-contract-stage17.md}}. Kind: markdown. Size: 3519 bytes. Shape: 101 lines. Markdown document, 101 lines, 0 explicit URLs. Lead heading: Repository audit note for 2026-05-09 family-first MIPROv2 contract stage 17. & \textbf{\path{docs/audit/2026-05-09-family-first-mipro-contract-stage17.md}}. Тип: markdown. Размер: 3519 байт. Форма: 101 строк. Документ Markdown, 101 строк, явных URL: 0. Главный заголовок: Repository audit note for 2026-05-09 family-first MIPROv2 contract stage 17. \\
\textbf{\path{docs/audit/2026-05-09-family-first-mipro-contract-stage18.md}}. Kind: markdown. Size: 3164 bytes. Shape: 90 lines. Markdown document, 90 lines, 0 explicit URLs. Lead heading: Repository audit note for 2026-05-09 family-first MIPROv2 contract stage 18. & \textbf{\path{docs/audit/2026-05-09-family-first-mipro-contract-stage18.md}}. Тип: markdown. Размер: 3164 байт. Форма: 90 строк. Документ Markdown, 90 строк, явных URL: 0. Главный заголовок: Repository audit note for 2026-05-09 family-first MIPROv2 contract stage 18. \\
\textbf{\path{docs/audit/2026-05-09-family-first-mipro-contract-stage19.md}}. Kind: markdown. Size: 3984 bytes. Shape: 102 lines. Markdown document, 102 lines, 0 explicit URLs. Lead heading: Repository audit note for 2026-05-09 family-first MIPROv2 contract stage 19. & \textbf{\path{docs/audit/2026-05-09-family-first-mipro-contract-stage19.md}}. Тип: markdown. Размер: 3984 байт. Форма: 102 строк. Документ Markdown, 102 строк, явных URL: 0. Главный заголовок: Repository audit note for 2026-05-09 family-first MIPROv2 contract stage 19. \\
\textbf{\path{docs/audit/2026-05-09-family-first-mipro-contract-stage2.md}}. Kind: markdown. Size: 4936 bytes. Shape: 123 lines. Markdown document, 123 lines, 0 explicit URLs. Lead heading: Repository audit note for 2026-05-09 family-first MIPROv2 contract stage 2. & \textbf{\path{docs/audit/2026-05-09-family-first-mipro-contract-stage2.md}}. Тип: markdown. Размер: 4936 байт. Форма: 123 строк. Документ Markdown, 123 строк, явных URL: 0. Главный заголовок: Repository audit note for 2026-05-09 family-first MIPROv2 contract stage 2. \\
\textbf{\path{docs/audit/2026-05-09-family-first-mipro-contract-stage20.md}}. Kind: markdown. Size: 3076 bytes. Shape: 88 lines. Markdown document, 88 lines, 0 explicit URLs. Lead heading: Repository audit note for 2026-05-09 family-first MIPROv2 contract stage 20. & \textbf{\path{docs/audit/2026-05-09-family-first-mipro-contract-stage20.md}}. Тип: markdown. Размер: 3076 байт. Форма: 88 строк. Документ Markdown, 88 строк, явных URL: 0. Главный заголовок: Repository audit note for 2026-05-09 family-first MIPROv2 contract stage 20. \\
\textbf{\path{docs/audit/2026-05-09-family-first-mipro-contract-stage21.md}}. Kind: markdown. Size: 5742 bytes. Shape: 122 lines. Markdown document, 122 lines, 0 explicit URLs. Lead heading: Repository audit note for 2026-05-09 family-first MIPROv2 contract stage 21. & \textbf{\path{docs/audit/2026-05-09-family-first-mipro-contract-stage21.md}}. Тип: markdown. Размер: 5742 байт. Форма: 122 строк. Документ Markdown, 122 строк, явных URL: 0. Главный заголовок: Repository audit note for 2026-05-09 family-first MIPROv2 contract stage 21. \\
\textbf{\path{docs/audit/2026-05-09-family-first-mipro-contract-stage3.md}}. Kind: markdown. Size: 4147 bytes. Shape: 96 lines. Markdown document, 96 lines, 0 explicit URLs. Lead heading: Repository audit note for 2026-05-09 family-first MIPROv2 contract stage 3. & \textbf{\path{docs/audit/2026-05-09-family-first-mipro-contract-stage3.md}}. Тип: markdown. Размер: 4147 байт. Форма: 96 строк. Документ Markdown, 96 строк, явных URL: 0. Главный заголовок: Repository audit note for 2026-05-09 family-first MIPROv2 contract stage 3. \\
\textbf{\path{docs/audit/2026-05-09-family-first-mipro-contract-stage4.md}}. Kind: markdown. Size: 5209 bytes. Shape: 120 lines. Markdown document, 120 lines, 0 explicit URLs. Lead heading: Repository audit note for 2026-05-09 family-first MIPROv2 contract stage 4. & \textbf{\path{docs/audit/2026-05-09-family-first-mipro-contract-stage4.md}}. Тип: markdown. Размер: 5209 байт. Форма: 120 строк. Документ Markdown, 120 строк, явных URL: 0. Главный заголовок: Repository audit note for 2026-05-09 family-first MIPROv2 contract stage 4. \\
\textbf{\path{docs/audit/2026-05-09-family-first-mipro-contract-stage5.md}}. Kind: markdown. Size: 5172 bytes. Shape: 122 lines. Markdown document, 122 lines, 0 explicit URLs. Lead heading: Repository audit note for 2026-05-09 family-first MIPROv2 contract stage 5. & \textbf{\path{docs/audit/2026-05-09-family-first-mipro-contract-stage5.md}}. Тип: markdown. Размер: 5172 байт. Форма: 122 строк. Документ Markdown, 122 строк, явных URL: 0. Главный заголовок: Repository audit note for 2026-05-09 family-first MIPROv2 contract stage 5. \\
\textbf{\path{docs/audit/2026-05-09-family-first-mipro-contract-stage6.md}}. Kind: markdown. Size: 5120 bytes. Shape: 119 lines. Markdown document, 119 lines, 0 explicit URLs. Lead heading: Repository audit note for 2026-05-09 family-first MIPROv2 contract stage 6. & \textbf{\path{docs/audit/2026-05-09-family-first-mipro-contract-stage6.md}}. Тип: markdown. Размер: 5120 байт. Форма: 119 строк. Документ Markdown, 119 строк, явных URL: 0. Главный заголовок: Repository audit note for 2026-05-09 family-first MIPROv2 contract stage 6. \\
\textbf{\path{docs/audit/2026-05-09-family-first-mipro-contract-stage7.md}}. Kind: markdown. Size: 4680 bytes. Shape: 116 lines. Markdown document, 116 lines, 0 explicit URLs. Lead heading: Repository audit note for 2026-05-09 family-first MIPROv2 contract stage 7. & \textbf{\path{docs/audit/2026-05-09-family-first-mipro-contract-stage7.md}}. Тип: markdown. Размер: 4680 байт. Форма: 116 строк. Документ Markdown, 116 строк, явных URL: 0. Главный заголовок: Repository audit note for 2026-05-09 family-first MIPROv2 contract stage 7. \\
\textbf{\path{docs/audit/2026-05-09-family-first-mipro-contract-stage8.md}}. Kind: markdown. Size: 4477 bytes. Shape: 110 lines. Markdown document, 110 lines, 0 explicit URLs. Lead heading: Repository audit note for 2026-05-09 family-first MIPROv2 contract stage 8. & \textbf{\path{docs/audit/2026-05-09-family-first-mipro-contract-stage8.md}}. Тип: markdown. Размер: 4477 байт. Форма: 110 строк. Документ Markdown, 110 строк, явных URL: 0. Главный заголовок: Repository audit note for 2026-05-09 family-first MIPROv2 contract stage 8. \\
\textbf{\path{docs/audit/2026-05-09-family-first-mipro-contract-stage9.md}}. Kind: markdown. Size: 4003 bytes. Shape: 106 lines. Markdown document, 106 lines, 0 explicit URLs. Lead heading: Repository audit note for 2026-05-09 family-first MIPROv2 contract stage 9. & \textbf{\path{docs/audit/2026-05-09-family-first-mipro-contract-stage9.md}}. Тип: markdown. Размер: 4003 байт. Форма: 106 строк. Документ Markdown, 106 строк, явных URL: 0. Главный заголовок: Repository audit note for 2026-05-09 family-first MIPROv2 contract stage 9. \\
\textbf{\path{docs/audit/2026-05-10-aks-run-25625638264-family-first-runtime-gap.md}}. Kind: markdown. Size: 8472 bytes. Shape: 246 lines. Markdown document, 246 lines, 0 explicit URLs. Lead heading: Repository audit note for 2026-05-10 AKS run 25625638264 family-first runtime gap. & \textbf{\path{docs/audit/2026-05-10-aks-run-25625638264-family-first-runtime-gap.md}}. Тип: markdown. Размер: 8472 байт. Форма: 246 строк. Документ Markdown, 246 строк, явных URL: 0. Главный заголовок: Repository audit note for 2026-05-10 AKS run 25625638264 family-first runtime gap. \\
\textbf{\path{docs/audit/2026-05-10-aks-run-25629990035-family-first-runtime-still-heuristic.md}}. Kind: markdown. Size: 6843 bytes. Shape: 211 lines. Markdown document, 211 lines, 0 explicit URLs. Lead heading: Repository audit note for 2026-05-10 AKS run 25629990035 family-first runtime still heuristic. & \textbf{\path{docs/audit/2026-05-10-aks-run-25629990035-family-first-runtime-still-heuristic.md}}. Тип: markdown. Размер: 6843 байт. Форма: 211 строк. Документ Markdown, 211 строк, явных URL: 0. Главный заголовок: Repository audit note for 2026-05-10 AKS run 25629990035 family-first runtime still heuristic. \\
\textbf{\path{docs/audit/2026-05-10-aks-run-25632110510-handoff-fixed-runtime-still-heuristic.md}}. Kind: markdown. Size: 7832 bytes. Shape: 231 lines. Markdown document, 231 lines, 1 explicit URLs. Lead heading: Repository audit note for 2026-05-10 AKS run 25632110510 handoff fixed but runtime still heuristic. & \textbf{\path{docs/audit/2026-05-10-aks-run-25632110510-handoff-fixed-runtime-still-heuristic.md}}. Тип: markdown. Размер: 7832 байт. Форма: 231 строк. Документ Markdown, 231 строк, явных URL: 1. Главный заголовок: Repository audit note for 2026-05-10 AKS run 25632110510 handoff fixed but runtime still heuristic. \\
\textbf{\path{docs/audit/2026-05-10-aks-run-25634202380-reset-lane-batch-handoff-runtime-still-heuristic.md}}. Kind: markdown. Size: 6096 bytes. Shape: 197 lines. Markdown document, 197 lines, 0 explicit URLs. Lead heading: Repository audit note for AKS run 25634202380 on 2026-05-10. & \textbf{\path{docs/audit/2026-05-10-aks-run-25634202380-reset-lane-batch-handoff-runtime-still-heuristic.md}}. Тип: markdown. Размер: 6096 байт. Форма: 197 строк. Документ Markdown, 197 строк, явных URL: 0. Главный заголовок: Repository audit note for AKS run 25634202380 on 2026-05-10. \\
\textbf{\path{docs/audit/2026-05-10-family-first-runtime-fixes-after-run-25625638264.md}}. Kind: markdown. Size: 12143 bytes. Shape: 261 lines. Markdown document, 261 lines, 0 explicit URLs. Lead heading: Repository audit note for 2026-05-10 family-first runtime fixes after run 25625638264. & \textbf{\path{docs/audit/2026-05-10-family-first-runtime-fixes-after-run-25625638264.md}}. Тип: markdown. Размер: 12143 байт. Форма: 261 строк. Документ Markdown, 261 строк, явных URL: 0. Главный заголовок: Repository audit note for 2026-05-10 family-first runtime fixes after run 25625638264. \\
\textbf{\path{docs/audit/2026-05-10-family-first-runtime-fixes-after-run-25629990035.md}}. Kind: markdown. Size: 4170 bytes. Shape: 98 lines. Markdown document, 98 lines, 0 explicit URLs. Lead heading: Repository audit note for 2026-05-10 family-first runtime fixes after AKS run 25629990035. & \textbf{\path{docs/audit/2026-05-10-family-first-runtime-fixes-after-run-25629990035.md}}. Тип: markdown. Размер: 4170 байт. Форма: 98 строк. Документ Markdown, 98 строк, явных URL: 0. Главный заголовок: Repository audit note for 2026-05-10 family-first runtime fixes after AKS run 25629990035. \\
\textbf{\path{docs/audit/2026-05-10-family-first-runtime-fixes-after-run-25632110510.md}}. Kind: markdown. Size: 5332 bytes. Shape: 112 lines. Markdown document, 112 lines, 0 explicit URLs. Lead heading: Repository audit note for 2026-05-10 family-first runtime fixes after AKS run 25632110510. & \textbf{\path{docs/audit/2026-05-10-family-first-runtime-fixes-after-run-25632110510.md}}. Тип: markdown. Размер: 5332 байт. Форма: 112 строк. Документ Markdown, 112 строк, явных URL: 0. Главный заголовок: Repository audit note for 2026-05-10 family-first runtime fixes after AKS run 25632110510. \\
\textbf{\path{docs/audit/2026-05-10-family-state-populates-but-publish-gate-and-dirty-fathers-still-block-runtime.md}}. Kind: markdown. Size: 7106 bytes. Shape: 183 lines. Markdown document, 183 lines, 0 explicit URLs. Lead heading: 2026-05-10 Family State Populates But Publish Gate And Dirty Fathers Still Block Runtime. & \textbf{\path{docs/audit/2026-05-10-family-state-populates-but-publish-gate-and-dirty-fathers-still-block-runtime.md}}. Тип: markdown. Размер: 7106 байт. Форма: 183 строк. Документ Markdown, 183 строк, явных URL: 0. Главный заголовок: 2026-05-10 Family State Populates But Publish Gate And Dirty Fathers Still Block Runtime. \\
\textbf{\path{docs/audit/2026-05-10-family-state-root-alias-and-latest-bundle-staging-fallback.md}}. Kind: markdown. Size: 3390 bytes. Shape: 81 lines. Markdown document, 81 lines, 0 explicit URLs. Lead heading: 2026-05-10 Family-State Root Alias And Latest-Bundle Staging Fallback. & \textbf{\path{docs/audit/2026-05-10-family-state-root-alias-and-latest-bundle-staging-fallback.md}}. Тип: markdown. Размер: 3390 байт. Форма: 81 строк. Документ Markdown, 81 строк, явных URL: 0. Главный заголовок: 2026-05-10 Family-State Root Alias And Latest-Bundle Staging Fallback. \\
\textbf{\path{docs/audit/2026-05-10-hidden-token-spend-and-incremental-trainer-fixes.md}}. Kind: markdown. Size: 6740 bytes. Shape: 154 lines. Markdown document, 154 lines, 0 explicit URLs. Lead heading: 2026-05-10 Hidden Token Spend And Incremental Trainer Fixes. & \textbf{\path{docs/audit/2026-05-10-hidden-token-spend-and-incremental-trainer-fixes.md}}. Тип: markdown. Размер: 6740 байт. Форма: 154 строк. Документ Markdown, 154 строк, явных URL: 0. Главный заголовок: 2026-05-10 Hidden Token Spend And Incremental Trainer Fixes. \\
\textbf{\path{docs/audit/2026-05-10-live-trainer-service-replays-processed-ledger-while-queue-empty.md}}. Kind: markdown. Size: 5891 bytes. Shape: 160 lines. Markdown document, 160 lines, 0 explicit URLs. Lead heading: 2026-05-10 Live Trainer Service Replays Processed Ledger While Queue Is Empty. & \textbf{\path{docs/audit/2026-05-10-live-trainer-service-replays-processed-ledger-while-queue-empty.md}}. Тип: markdown. Размер: 5891 байт. Форма: 160 строк. Документ Markdown, 160 строк, явных URL: 0. Главный заголовок: 2026-05-10 Live Trainer Service Replays Processed Ledger While Queue Is Empty. \\
\textbf{\path{docs/audit/2026-05-10-stale-submodule-runtime-image-root-cause-and-source-preference-fix.md}}. Kind: markdown. Size: 5021 bytes. Shape: 114 lines. Markdown document, 114 lines, 0 explicit URLs. Lead heading: 2026-05-10 Stale Submodule Runtime Image Root Cause And Source-Preference Fix. & \textbf{\path{docs/audit/2026-05-10-stale-submodule-runtime-image-root-cause-and-source-preference-fix.md}}. Тип: markdown. Размер: 5021 байт. Форма: 114 строк. Документ Markdown, 114 строк, явных URL: 0. Главный заголовок: 2026-05-10 Stale Submodule Runtime Image Root Cause And Source-Preference Fix. \\
\textbf{\path{docs/audit/2026-05-11-family-runtime-artifact-fallback-and-staging-path-fix.md}}. Kind: markdown. Size: 5601 bytes. Shape: 139 lines. Markdown document, 139 lines, 0 explicit URLs. Lead heading: 2026-05-11 Family Runtime Artifact Fallback And Bundle Staging Path Fix. & \textbf{\path{docs/audit/2026-05-11-family-runtime-artifact-fallback-and-staging-path-fix.md}}. Тип: markdown. Размер: 5601 байт. Форма: 139 строк. Документ Markdown, 139 строк, явных URL: 0. Главный заголовок: 2026-05-11 Family Runtime Artifact Fallback And Bundle Staging Path Fix. \\
\textbf{\path{docs/audit/2026-05-11-fresh-run-family-match-without-family-artifact-runtime.md}}. Kind: markdown. Size: 6472 bytes. Shape: 211 lines. Markdown document, 211 lines, 0 explicit URLs. Lead heading: 2026-05-11 Fresh Run Family Match Without Family Artifact Runtime. & \textbf{\path{docs/audit/2026-05-11-fresh-run-family-match-without-family-artifact-runtime.md}}. Тип: markdown. Размер: 6472 байт. Форма: 211 строк. Документ Markdown, 211 строк, явных URL: 0. Главный заголовок: 2026-05-11 Fresh Run Family Match Without Family Artifact Runtime. \\
\textbf{\path{docs/audit/2026-05-11-idle-trainer-cycles-no-longer-autorepublish-bundles.md}}. Kind: markdown. Size: 3162 bytes. Shape: 108 lines. Markdown document, 108 lines, 0 explicit URLs. Lead heading: 2026-05-11 Idle Trainer Cycles No Longer Auto-Republish Bundle Versions. & \textbf{\path{docs/audit/2026-05-11-idle-trainer-cycles-no-longer-autorepublish-bundles.md}}. Тип: markdown. Размер: 3162 байт. Форма: 108 строк. Документ Markdown, 108 строк, явных URL: 0. Главный заголовок: 2026-05-11 Idle Trainer Cycles No Longer Auto-Republish Bundle Versions. \\
\textbf{\path{docs/audit/2026-05-11-original-prompt-routing-and-batch-handoff-fixes.md}}. Kind: markdown. Size: 3829 bytes. Shape: 91 lines. Markdown document, 91 lines, 0 explicit URLs. Lead heading: 2026-05-11 Original Prompt Routing And Batch Handoff Fixes. & \textbf{\path{docs/audit/2026-05-11-original-prompt-routing-and-batch-handoff-fixes.md}}. Тип: markdown. Размер: 3829 байт. Форма: 91 строк. Документ Markdown, 91 строк, явных URL: 0. Главный заголовок: 2026-05-11 Original Prompt Routing And Batch Handoff Fixes. \\
\textbf{\path{docs/audit/2026-05-11-trainer-queue-only-no-processed-replay.md}}. Kind: markdown. Size: 6373 bytes. Shape: 166 lines. Markdown document, 166 lines, 0 explicit URLs. Lead heading: 2026-05-11 Trainer Is Now Queue-Only And Will Not Replay Processed History. & \textbf{\path{docs/audit/2026-05-11-trainer-queue-only-no-processed-replay.md}}. Тип: markdown. Размер: 6373 байт. Форма: 166 строк. Документ Markdown, 166 строк, явных URL: 0. Главный заголовок: 2026-05-11 Trainer Is Now Queue-Only And Will Not Replay Processed History. \\
\textbf{\path{docs/audit/2026-05-11-trainer-service-preflights-queue-before-running-cycles.md}}. Kind: markdown. Size: 4211 bytes. Shape: 121 lines. Markdown document, 121 lines, 0 explicit URLs. Lead heading: 2026-05-11 Trainer Service Preflights Queue Before Running Cycles. & \textbf{\path{docs/audit/2026-05-11-trainer-service-preflights-queue-before-running-cycles.md}}. Тип: markdown. Размер: 4211 байт. Форма: 121 строк. Документ Markdown, 121 строк, явных URL: 0. Главный заголовок: 2026-05-11 Trainer Service Preflights Queue Before Running Cycles. \\
\textbf{\path{docs/audit/2026-05-11-trusted-handoff-syntax-error-blocked-execution-artifact-upload.md}}. Kind: markdown. Size: 4031 bytes. Shape: 86 lines. Markdown document, 86 lines, 0 explicit URLs. Lead heading: 2026-05-11 Trusted Handoff Syntax Error Blocked Execution Artifact Upload. & \textbf{\path{docs/audit/2026-05-11-trusted-handoff-syntax-error-blocked-execution-artifact-upload.md}}. Тип: markdown. Размер: 4031 байт. Форма: 86 строк. Документ Markdown, 86 строк, явных URL: 0. Главный заголовок: 2026-05-11 Trusted Handoff Syntax Error Blocked Execution Artifact Upload. \\
\textbf{\path{docs/audit/README.md}}. Kind: markdown. Size: 443 bytes. Shape: 14 lines. Markdown document, 14 lines, 0 explicit URLs. Lead heading: Repository Audit Notes. & \textbf{\path{docs/audit/README.md}}. Тип: markdown. Размер: 443 байт. Форма: 14 строк. Документ Markdown, 14 строк, явных URL: 0. Главный заголовок: Repository Audit Notes. \\
\textbf{\path{docs/guides/hushwheel-fixture-rag-guide.md}}. Kind: markdown. Size: 5169 bytes. Shape: 138 lines. Markdown document, 138 lines, 0 explicit URLs. Lead heading: RAG Guide For The Hushwheel Fixture. & \textbf{\path{docs/guides/hushwheel-fixture-rag-guide.md}}. Тип: markdown. Размер: 5169 байт. Форма: 138 строк. Документ Markdown, 138 строк, явных URL: 0. Главный заголовок: RAG Guide For The Hushwheel Fixture. \\
\textbf{\path{docs/operations/azure-deployment.md}}. Kind: markdown. Size: 4901 bytes. Shape: 111 lines. Markdown document, 111 lines, 5 explicit URLs. Lead heading: Azure Deployment Notes. & \textbf{\path{docs/operations/azure-deployment.md}}. Тип: markdown. Размер: 4901 байт. Форма: 111 строк. Документ Markdown, 111 строк, явных URL: 5. Главный заголовок: Azure Deployment Notes. \\
\textbf{\path{docs/operations/environment.md}}. Kind: markdown. Size: 7589 bytes. Shape: 130 lines. Markdown document, 130 lines, 0 explicit URLs. Lead heading: Environment Variables. & \textbf{\path{docs/operations/environment.md}}. Тип: markdown. Размер: 7589 байт. Форма: 130 строк. Документ Markdown, 130 строк, явных URL: 0. Главный заголовок: Environment Variables. \\
\textbf{\path{docs/operations/runtime-image.md}}. Kind: markdown. Size: 1989 bytes. Shape: 46 lines. Markdown document, 46 lines, 0 explicit URLs. Lead heading: Runtime Image Notes. & \textbf{\path{docs/operations/runtime-image.md}}. Тип: markdown. Размер: 1989 байт. Форма: 46 строк. Документ Markdown, 46 строк, явных URL: 0. Главный заголовок: Runtime Image Notes. \\
\textbf{\path{docs/operations/simple-end-to-end-verification-guide.md}}. Kind: markdown. Size: 4111 bytes. Shape: 105 lines. Markdown document, 105 lines, 12 explicit URLs. Lead heading: Simple End-To-End Verification Guide. & \textbf{\path{docs/operations/simple-end-to-end-verification-guide.md}}. Тип: markdown. Размер: 4111 байт. Форма: 105 строк. Документ Markdown, 105 строк, явных URL: 12. Главный заголовок: Simple End-To-End Verification Guide. \\
\textbf{\path{docs/operations/trainer-deployment.md}}. Kind: markdown. Size: 6047 bytes. Shape: 152 lines. Markdown document, 152 lines, 0 explicit URLs. Lead heading: Trainer Deployment Notes. & \textbf{\path{docs/operations/trainer-deployment.md}}. Тип: markdown. Размер: 6047 байт. Форма: 152 строк. Документ Markdown, 152 строк, явных URL: 0. Главный заголовок: Trainer Deployment Notes. \\
\textbf{\path{docs/planning/codex-exec-resume-plan.md}}. Kind: markdown. Size: 11561 bytes. Shape: 207 lines. Markdown document, 207 lines, 0 explicit URLs. Lead heading: Codex Exec Resume Plan. & \textbf{\path{docs/planning/codex-exec-resume-plan.md}}. Тип: markdown. Размер: 11561 байт. Форма: 207 строк. Документ Markdown, 207 строк, явных URL: 0. Главный заголовок: Codex Exec Resume Plan. \\
\textbf{\path{docs/planning/dataset-integration-plan.md}}. Kind: markdown. Size: 33882 bytes. Shape: 478 lines. Markdown document, 478 lines, 0 explicit URLs. Lead heading: Dataset Integration Plan. & \textbf{\path{docs/planning/dataset-integration-plan.md}}. Тип: markdown. Размер: 33882 байт. Форма: 478 строк. Документ Markdown, 478 строк, явных URL: 0. Главный заголовок: Dataset Integration Plan. \\
\textbf{\path{docs/planning/family-first-mipro-runtime-contract.md}}. Kind: markdown. Size: 29074 bytes. Shape: 439 lines. Markdown document, 439 lines, 0 explicit URLs. Lead heading: Family-First DSPy / MIPROv2 Contract. & \textbf{\path{docs/planning/family-first-mipro-runtime-contract.md}}. Тип: markdown. Размер: 29074 байт. Форма: 439 строк. Документ Markdown, 439 строк, явных URL: 0. Главный заголовок: Family-First DSPy / MIPROv2 Contract. \\
\textbf{\path{docs/planning/per-turn-dspy-mediation-contract.md}}. Kind: markdown. Size: 4966 bytes. Shape: 115 lines. Markdown document, 115 lines, 0 explicit URLs. Lead heading: Per-Turn DSPy Mediation Contract. & \textbf{\path{docs/planning/per-turn-dspy-mediation-contract.md}}. Тип: markdown. Размер: 4966 байт. Форма: 115 строк. Документ Markdown, 115 строк, явных URL: 0. Главный заголовок: Per-Turn DSPy Mediation Contract. \\
\textbf{\path{docs/planning/repo-hardening-plan.md}}. Kind: markdown. Size: 4545 bytes. Shape: 71 lines. Markdown document, 71 lines, 0 explicit URLs. Lead heading: Repository Hardening Plan. & \textbf{\path{docs/planning/repo-hardening-plan.md}}. Тип: markdown. Размер: 4545 байт. Форма: 71 строк. Документ Markdown, 71 строк, явных URL: 0. Главный заголовок: Repository Hardening Plan. \\
\textbf{\path{docs/planning/trainer-context-group-champion-plan.md}}. Kind: markdown. Size: 7736 bytes. Shape: 119 lines. Markdown document, 119 lines, 0 explicit URLs. Lead heading: Trainer Context-Group Champion Plan. & \textbf{\path{docs/planning/trainer-context-group-champion-plan.md}}. Тип: markdown. Размер: 7736 байт. Форма: 119 строк. Документ Markdown, 119 строк, явных URL: 0. Главный заголовок: Trainer Context-Group Champion Plan. \\
\textbf{\path{mkdocs.yml}}. Kind: yaml. Size: 843 bytes. Shape: 37 lines. YAML data or workflow file, 37 lines, 2 explicit URLs. Top keys: site\_name, site\_description, site\_url. & \textbf{\path{mkdocs.yml}}. Тип: yaml. Размер: 843 байт. Форма: 37 строк. Файл YAML с данными или workflow, 37 строк, явных URL: 2. Верхние ключи: site\_name, site\_description, site\_url. \\
\textbf{\path{notebooks/01_repo_rag_research.ipynb}}. Kind: notebook. Size: 9324 bytes. Shape: 281 lines. Notebook, 11 cells, 0 explicit URLs. Opening heading: Repository RAG Research Playbook. & \textbf{\path{notebooks/01_repo_rag_research.ipynb}}. Тип: notebook. Размер: 9324 байт. Форма: 281 строк. Ноутбук, 11 ячеек, явных URL: 0. Первый заголовок: Repository RAG Research Playbook. \\
\textbf{\path{notebooks/02_agent_workflow_checklist.ipynb}}. Kind: notebook. Size: 17632 bytes. Shape: 404 lines. Notebook, 9 cells, 1 explicit URLs. Opening heading: Agent Workflow Checklist. & \textbf{\path{notebooks/02_agent_workflow_checklist.ipynb}}. Тип: notebook. Размер: 17632 байт. Форма: 404 строк. Ноутбук, 9 ячеек, явных URL: 1. Первый заголовок: Agent Workflow Checklist. \\
\textbf{\path{notebooks/03_dspy_training_lab.ipynb}}. Kind: notebook. Size: 12048 bytes. Shape: 314 lines. Notebook, 11 cells, 0 explicit URLs. Opening heading: DSPy Training Sample Lab. & \textbf{\path{notebooks/03_dspy_training_lab.ipynb}}. Тип: notebook. Размер: 12048 байт. Форма: 314 строк. Ноутбук, 11 ячеек, явных URL: 0. Первый заголовок: DSPy Training Sample Lab. \\
\textbf{\path{notebooks/04_sample_population_lab.ipynb}}. Kind: notebook. Size: 7894 bytes. Shape: 250 lines. Notebook, 9 cells, 0 explicit URLs. Opening heading: Sample Population Lab. & \textbf{\path{notebooks/04_sample_population_lab.ipynb}}. Тип: notebook. Размер: 7894 байт. Форма: 250 строк. Ноутбук, 9 ячеек, явных URL: 0. Первый заголовок: Sample Population Lab. \\
\textbf{\path{notebooks/05_hushwheel_fixture_rag_lab.ipynb}}. Kind: notebook. Size: 13622 bytes. Shape: 378 lines. Notebook, 9 cells, 0 explicit URLs. Opening heading: Hushwheel Fixture RAG Lab. & \textbf{\path{notebooks/05_hushwheel_fixture_rag_lab.ipynb}}. Тип: notebook. Размер: 13622 байт. Форма: 378 строк. Ноутбук, 9 ячеек, явных URL: 0. Первый заголовок: Hushwheel Fixture RAG Lab. \\
\textbf{\path{publication/.gitignore}}. Kind: generic. Size: 80 bytes. Shape: 12 lines. Tracked repository file, 12 lines, 0 explicit URLs. Opening line: .build/ & \textbf{\path{publication/.gitignore}}. Тип: generic. Размер: 80 байт. Форма: 12 строк. Отслеживаемый файл репозитория, 12 строк, явных URL: 0. Первая строка: .build/ \\
\textbf{\path{publication/Makefile}}. Kind: make. Size: 1644 bytes. Shape: 48 lines. Repository Makefile, 48 lines, 0 explicit URLs. It exposes the shared automation surfaces. & \textbf{\path{publication/Makefile}}. Тип: make. Размер: 1644 байт. Форма: 48 строк. Makefile репозитория, 48 строк, явных URL: 0. Он открывает общие поверхности автоматизации. \\
\textbf{\path{publication/README.md}}. Kind: markdown. Size: 1499 bytes. Shape: 34 lines. Markdown document, 34 lines, 0 explicit URLs. Lead heading: Publication Article. & \textbf{\path{publication/README.md}}. Тип: markdown. Размер: 1499 байт. Форма: 34 строк. Документ Markdown, 34 строк, явных URL: 0. Главный заголовок: Publication Article. \\
\textbf{\path{publication/article-banner.png}}. Kind: binary. Size: 79074 bytes. Shape: binary surface. Binary asset, 79074 bytes of binary data. It is tracked for stable publication output. & \textbf{\path{publication/article-banner.png}}. Тип: binary. Размер: 79074 байт. Форма: бинарная поверхность. Бинарный артефакт, 79074 байт бинарных данных. Он хранится в Git для стабильного публикационного вывода. \\
\textbf{\path{publication/exploratorium_translation/.gitignore}}. Kind: generic. Size: 32 bytes. Shape: 6 lines. Tracked repository file, 6 lines, 0 explicit URLs. Opening line: .build/ & \textbf{\path{publication/exploratorium_translation/.gitignore}}. Тип: generic. Размер: 32 байт. Форма: 6 строк. Отслеживаемый файл репозитория, 6 строк, явных URL: 0. Первая строка: .build/ \\
\textbf{\path{publication/exploratorium_translation/Makefile}}. Kind: make. Size: 870 bytes. Shape: 30 lines. Repository Makefile, 30 lines, 0 explicit URLs. It exposes the shared automation surfaces. & \textbf{\path{publication/exploratorium_translation/Makefile}}. Тип: make. Размер: 870 байт. Форма: 30 строк. Makefile репозитория, 30 строк, явных URL: 0. Он открывает общие поверхности автоматизации. \\
\textbf{\path{publication/exploratorium_translation/README.md}}. Kind: markdown. Size: 997 bytes. Shape: 28 lines. Markdown document, 28 lines, 0 explicit URLs. Lead heading: Exploratorium Translation. & \textbf{\path{publication/exploratorium_translation/README.md}}. Тип: markdown. Размер: 997 байт. Форма: 28 строк. Документ Markdown, 28 строк, явных URL: 0. Главный заголовок: Exploratorium Translation. \\
\textbf{\path{publication/exploratorium_translation/exploratorium_translation.tex}}. Kind: latex. Size: 2573 bytes. Shape: 70 lines. LaTeX source, 70 lines, 0 explicit URLs. Sections declared in preview: 0. & \textbf{\path{publication/exploratorium_translation/exploratorium_translation.tex}}. Тип: latex. Размер: 2573 байт. Форма: 70 строк. Исходник LaTeX, 70 строк, явных URL: 0. Секций в превью: 0. \\
\textbf{\path{publication/references.bib}}. Kind: bibliography. Size: 1252 bytes. Shape: 28 lines. Bibliography file, 28 lines, 2 explicit URLs. Entries counted in preview: 4. & \textbf{\path{publication/references.bib}}. Тип: bibliography. Размер: 1252 байт. Форма: 28 строк. Библиографический файл, 28 строк, явных URL: 2. Записей в превью: 4. \\
\textbf{\path{publication/repository-rag-lab-article.pdf}}. Kind: binary. Size: 236232 bytes. Shape: binary surface. Binary asset, 236232 bytes of binary data. It is tracked for stable publication output. & \textbf{\path{publication/repository-rag-lab-article.pdf}}. Тип: binary. Размер: 236232 байт. Форма: бинарная поверхность. Бинарный артефакт, 236232 байт бинарных данных. Он хранится в Git для стабильного публикационного вывода. \\
\textbf{\path{publication/repository-rag-lab-article.tex}}. Kind: latex. Size: 9526 bytes. Shape: 208 lines. LaTeX source, 208 lines, 0 explicit URLs. Sections declared in preview: 4. & \textbf{\path{publication/repository-rag-lab-article.tex}}. Тип: latex. Размер: 9526 байт. Форма: 208 строк. Исходник LaTeX, 208 строк, явных URL: 0. Секций в превью: 4. \\
\textbf{\path{publication/todo-backlog-table.tex}}. Kind: latex. Size: 10547 bytes. Shape: 34 lines. LaTeX source, 34 lines, 36 explicit URLs. Sections declared in preview: 0. & \textbf{\path{publication/todo-backlog-table.tex}}. Тип: latex. Размер: 10547 байт. Форма: 34 строк. Исходник LaTeX, 34 строк, явных URL: 36. Секций в превью: 0. \\
\textbf{\path{pyproject.toml}}. Kind: toml. Size: 2285 bytes. Shape: 97 lines. TOML configuration file, 97 lines, 0 explicit URLs. Top keys: build-system, project, dependency-groups. & \textbf{\path{pyproject.toml}}. Тип: toml. Размер: 2285 байт. Форма: 97 строк. Конфигурация TOML, 97 строк, явных URL: 0. Верхние ключи: build-system, project, dependency-groups. \\
\textbf{\path{rust-cli/Cargo.lock}}. Kind: generic. Size: 3213 bytes. Shape: binary surface. Tracked repository file, 3213 bytes of binary data, 0 explicit URLs. Opening line: No non-empty preview line. & \textbf{\path{rust-cli/Cargo.lock}}. Тип: generic. Размер: 3213 байт. Форма: бинарная поверхность. Отслеживаемый файл репозитория, 3213 байт бинарных данных, явных URL: 0. Первая строка: No non-empty preview line. \\
\textbf{\path{rust-cli/Cargo.toml}}. Kind: toml. Size: 144 bytes. Shape: 8 lines. TOML configuration file, 8 lines, 0 explicit URLs. Top keys: package, dependencies. & \textbf{\path{rust-cli/Cargo.toml}}. Тип: toml. Размер: 144 байт. Форма: 8 строк. Конфигурация TOML, 8 строк, явных URL: 0. Верхние ключи: package, dependencies. \\
\textbf{\path{rust-cli/src/main.rs}}. Kind: source. Size: 20422 bytes. Shape: 626 lines. Source file, 626 lines, 0 explicit URLs. Opening line: use rusqlite::\{params, Connection\}; & \textbf{\path{rust-cli/src/main.rs}}. Тип: source. Размер: 20422 байт. Форма: 626 строк. Исходный файл, 626 строк, явных URL: 0. Первая строка: use rusqlite::\{params, Connection\}; \\
\textbf{\path{samples/README.md}}. Kind: markdown. Size: 713 bytes. Shape: 13 lines. Markdown document, 13 lines, 0 explicit URLs. Lead heading: Samples. & \textbf{\path{samples/README.md}}. Тип: markdown. Размер: 713 байт. Форма: 13 строк. Документ Markdown, 13 строк, явных URL: 0. Главный заголовок: Samples. \\
\textbf{\path{samples/logs/20260315T050730Z-gh-runs.md}}. Kind: markdown. Size: 1341 bytes. Shape: 45 lines. Markdown document, 45 lines, 3 explicit URLs. Lead heading: GitHub Actions Run Log. & \textbf{\path{samples/logs/20260315T050730Z-gh-runs.md}}. Тип: markdown. Размер: 1341 байт. Форма: 45 строк. Документ Markdown, 45 строк, явных URL: 3. Главный заголовок: GitHub Actions Run Log. \\
\textbf{\path{samples/logs/20260315T085140Z-gh-red-status-check.md}}. Kind: markdown. Size: 847 bytes. Shape: 25 lines. Markdown document, 25 lines, 1 explicit URLs. Lead heading: GitHub Actions Red Status Check. & \textbf{\path{samples/logs/20260315T085140Z-gh-red-status-check.md}}. Тип: markdown. Размер: 847 байт. Форма: 25 строк. Документ Markdown, 25 строк, явных URL: 1. Главный заголовок: GitHub Actions Red Status Check. \\
\textbf{\path{samples/logs/20260317T064722Z-full-build.raw.log}}. Kind: generic. Size: 2109 bytes. Shape: 62 lines. Tracked repository file, 62 lines, 0 explicit URLs. Opening line: \# Full Build Raw Log & \textbf{\path{samples/logs/20260317T064722Z-full-build.raw.log}}. Тип: generic. Размер: 2109 байт. Форма: 62 строк. Отслеживаемый файл репозитория, 62 строк, явных URL: 0. Первая строка: \# Full Build Raw Log \\
\textbf{\path{samples/logs/20260317T064833Z-full-build-attempt2.raw.log}}. Kind: generic. Size: 14950 bytes. Shape: 408 lines. Tracked repository file, 408 lines, 0 explicit URLs. Opening line: \# Full Build Raw Log Attempt 2 & \textbf{\path{samples/logs/20260317T064833Z-full-build-attempt2.raw.log}}. Тип: generic. Размер: 14950 байт. Форма: 408 строк. Отслеживаемый файл репозитория, 408 строк, явных URL: 0. Первая строка: \# Full Build Raw Log Attempt 2 \\
\textbf{\path{samples/logs/20260317T073401Z-gh-runs-notebook-scaffolding.md}}. Kind: markdown. Size: 2865 bytes. Shape: 35 lines. Markdown document, 35 lines, 13 explicit URLs. Lead heading: GitHub Actions Run Log. & \textbf{\path{samples/logs/20260317T073401Z-gh-runs-notebook-scaffolding.md}}. Тип: markdown. Размер: 2865 байт. Форма: 35 строк. Документ Markdown, 35 строк, явных URL: 13. Главный заголовок: GitHub Actions Run Log. \\
\textbf{\path{samples/logs/20260317T073404Z-gh-runs.md}}. Kind: markdown. Size: 1269 bytes. Shape: 32 lines. Markdown document, 32 lines, 3 explicit URLs. Lead heading: GitHub Actions Run Log. & \textbf{\path{samples/logs/20260317T073404Z-gh-runs.md}}. Тип: markdown. Размер: 1269 байт. Форма: 32 строк. Документ Markdown, 32 строк, явных URL: 3. Главный заголовок: GitHub Actions Run Log. \\
\textbf{\path{samples/logs/20260317T080018Z-gh-runs-all-notebooks.md}}. Kind: markdown. Size: 2128 bytes. Shape: 34 lines. Markdown document, 34 lines, 8 explicit URLs. Lead heading: GitHub Actions Run Log. & \textbf{\path{samples/logs/20260317T080018Z-gh-runs-all-notebooks.md}}. Тип: markdown. Размер: 2128 байт. Форма: 34 строк. Документ Markdown, 34 строк, явных URL: 8. Главный заголовок: GitHub Actions Run Log. \\
\textbf{\path{samples/logs/20260317T083207Z-gh-runs-todo-backlog.md}}. Kind: markdown. Size: 1063 bytes. Shape: 27 lines. Markdown document, 27 lines, 3 explicit URLs. Lead heading: GitHub Actions Run Log. & \textbf{\path{samples/logs/20260317T083207Z-gh-runs-todo-backlog.md}}. Тип: markdown. Размер: 1063 байт. Форма: 27 строк. Документ Markdown, 27 строк, явных URL: 3. Главный заголовок: GitHub Actions Run Log. \\
\textbf{\path{samples/logs/20260317T084723Z-gh-runs-repo-completeness-checklist.md}}. Kind: markdown. Size: 1186 bytes. Shape: 27 lines. Markdown document, 27 lines, 3 explicit URLs. Lead heading: GitHub Actions Run Log. & \textbf{\path{samples/logs/20260317T084723Z-gh-runs-repo-completeness-checklist.md}}. Тип: markdown. Размер: 1186 байт. Форма: 27 строк. Документ Markdown, 27 строк, явных URL: 3. Главный заголовок: GitHub Actions Run Log. \\
\textbf{\path{samples/logs/20260317T085723Z-gh-runs-hushwheel-fixture.md}}. Kind: markdown. Size: 1472 bytes. Shape: 36 lines. Markdown document, 36 lines, 1 explicit URLs. Lead heading: GitHub Actions Run Log. & \textbf{\path{samples/logs/20260317T085723Z-gh-runs-hushwheel-fixture.md}}. Тип: markdown. Размер: 1472 байт. Форма: 36 строк. Документ Markdown, 36 строк, явных URL: 1. Главный заголовок: GitHub Actions Run Log. \\
\textbf{\path{samples/logs/20260317T091641Z-gh-runs-hushwheel-rag-playbook.md}}. Kind: markdown. Size: 1470 bytes. Shape: 36 lines. Markdown document, 36 lines, 1 explicit URLs. Lead heading: GitHub Actions Run Log. & \textbf{\path{samples/logs/20260317T091641Z-gh-runs-hushwheel-rag-playbook.md}}. Тип: markdown. Размер: 1470 байт. Форма: 36 строк. Документ Markdown, 36 строк, явных URL: 1. Главный заголовок: GitHub Actions Run Log. \\
\textbf{\path{samples/logs/20260317T092251Z-gh-runs-publication-article.md}}. Kind: markdown. Size: 1801 bytes. Shape: 33 lines. Markdown document, 33 lines, 5 explicit URLs. Lead heading: GitHub Actions Run Log. & \textbf{\path{samples/logs/20260317T092251Z-gh-runs-publication-article.md}}. Тип: markdown. Размер: 1801 байт. Форма: 33 строк. Документ Markdown, 33 строк, явных URL: 5. Главный заголовок: GitHub Actions Run Log. \\
\textbf{\path{samples/logs/20260318T003027Z-gh-runs-wire-skill-triggers-for-do-repo-verification-aud.md}}. Kind: markdown. Size: 1113 bytes. Shape: 29 lines. Markdown document, 29 lines, 1 explicit URLs. Lead heading: GitHub Actions Run Log. & \textbf{\path{samples/logs/20260318T003027Z-gh-runs-wire-skill-triggers-for-do-repo-verification-aud.md}}. Тип: markdown. Размер: 1113 байт. Форма: 29 строк. Документ Markdown, 29 строк, явных URL: 1. Главный заголовок: GitHub Actions Run Log. \\
\textbf{\path{samples/logs/20260318T003654Z-gh-runs-execute-all-notebooks.md}}. Kind: markdown. Size: 1682 bytes. Shape: 41 lines. Markdown document, 41 lines, 1 explicit URLs. Lead heading: GitHub Actions Run Log. & \textbf{\path{samples/logs/20260318T003654Z-gh-runs-execute-all-notebooks.md}}. Тип: markdown. Размер: 1682 байт. Форма: 41 строк. Документ Markdown, 41 строк, явных URL: 1. Главный заголовок: GitHub Actions Run Log. \\
\textbf{\path{samples/logs/20260318T003927Z-gh-runs-publication-workflow-fix.md}}. Kind: markdown. Size: 1747 bytes. Shape: 35 lines. Markdown document, 35 lines, 2 explicit URLs. Lead heading: GitHub Actions Run Log. & \textbf{\path{samples/logs/20260318T003927Z-gh-runs-publication-workflow-fix.md}}. Тип: markdown. Размер: 1747 байт. Форма: 35 строк. Документ Markdown, 35 строк, явных URL: 2. Главный заголовок: GitHub Actions Run Log. \\
\textbf{\path{samples/logs/20260318T004712Z-gh-runs-hushwheel-program-harness.md}}. Kind: markdown. Size: 1798 bytes. Shape: 49 lines. Markdown document, 49 lines, 0 explicit URLs. Lead heading: GitHub Actions Run Log. & \textbf{\path{samples/logs/20260318T004712Z-gh-runs-hushwheel-program-harness.md}}. Тип: markdown. Размер: 1798 байт. Форма: 49 строк. Документ Markdown, 49 строк, явных URL: 0. Главный заголовок: GitHub Actions Run Log. \\
\textbf{\path{samples/logs/20260318T010254Z-gh-runs-refresh-final-notebook-outputs-for-no-move-those.md}}. Kind: markdown. Size: 1113 bytes. Shape: 29 lines. Markdown document, 29 lines, 1 explicit URLs. Lead heading: GitHub Actions Run Log. & \textbf{\path{samples/logs/20260318T010254Z-gh-runs-refresh-final-notebook-outputs-for-no-move-those.md}}. Тип: markdown. Размер: 1113 байт. Форма: 29 строк. Документ Markdown, 29 строк, явных URL: 1. Главный заголовок: GitHub Actions Run Log. \\
\textbf{\path{samples/logs/20260318T010637Z-gh-runs-env-refresh-retest.md}}. Kind: markdown. Size: 1029 bytes. Shape: 26 lines. Markdown document, 26 lines, 1 explicit URLs. Lead heading: GitHub Actions Run Log. & \textbf{\path{samples/logs/20260318T010637Z-gh-runs-env-refresh-retest.md}}. Тип: markdown. Размер: 1029 байт. Форма: 26 строк. Документ Markdown, 26 строк, явных URL: 1. Главный заголовок: GitHub Actions Run Log. \\
\textbf{\path{samples/logs/20260318T010741Z-gh-runs-consolidate-session-master.md}}. Kind: markdown. Size: 1950 bytes. Shape: 62 lines. Markdown document, 62 lines, 0 explicit URLs. Lead heading: GitHub Actions Run Log. & \textbf{\path{samples/logs/20260318T010741Z-gh-runs-consolidate-session-master.md}}. Тип: markdown. Размер: 1950 байт. Форма: 62 строк. Документ Markdown, 62 строк, явных URL: 0. Главный заголовок: GitHub Actions Run Log. \\
\textbf{\path{samples/logs/20260318T012137Z-gh-runs-azure-inference-endpoint-probe.md}}. Kind: markdown. Size: 1004 bytes. Shape: 26 lines. Markdown document, 26 lines, 1 explicit URLs. Lead heading: GitHub Actions Run Log. & \textbf{\path{samples/logs/20260318T012137Z-gh-runs-azure-inference-endpoint-probe.md}}. Тип: markdown. Размер: 1004 байт. Форма: 26 строк. Документ Markdown, 26 строк, явных URL: 1. Главный заголовок: GitHub Actions Run Log. \\
\textbf{\path{samples/logs/20260318T012714Z-gh-runs-tighten-repo-hygiene-for-other-agnet-is-doing-1-.md}}. Kind: markdown. Size: 1091 bytes. Shape: 29 lines. Markdown document, 29 lines, 1 explicit URLs. Lead heading: GitHub Actions Run Log. & \textbf{\path{samples/logs/20260318T012714Z-gh-runs-tighten-repo-hygiene-for-other-agnet-is-doing-1-.md}}. Тип: markdown. Размер: 1091 байт. Форма: 29 строк. Документ Markdown, 29 строк, явных URL: 1. Главный заголовок: GitHub Actions Run Log. \\
\textbf{\path{samples/logs/20260318T015109Z-gh-runs-improve-notebook-run-observability-and-reporting.md}}. Kind: markdown. Size: 2151 bytes. Shape: 55 lines. Markdown document, 55 lines, 2 explicit URLs. Lead heading: GitHub Actions Run Log. & \textbf{\path{samples/logs/20260318T015109Z-gh-runs-improve-notebook-run-observability-and-reporting.md}}. Тип: markdown. Размер: 2151 байт. Форма: 55 строк. Документ Markdown, 55 строк, явных URL: 2. Главный заголовок: GitHub Actions Run Log. \\
\textbf{\path{samples/logs/20260318T021736Z-gh-runs-sync-todo-backlog-into-markdown-and-publication-surfaces.md}}. Kind: markdown. Size: 1778 bytes. Shape: 55 lines. Markdown document, 55 lines, 2 explicit URLs. Lead heading: GitHub Actions Run Log. & \textbf{\path{samples/logs/20260318T021736Z-gh-runs-sync-todo-backlog-into-markdown-and-publication-surfaces.md}}. Тип: markdown. Размер: 1778 байт. Форма: 55 строк. Документ Markdown, 55 строк, явных URL: 2. Главный заголовок: GitHub Actions Run Log. \\
\textbf{\path{samples/logs/20260318T023922Z-gh-runs-implement-step-1-dspy-training-path-follow-up.md}}. Kind: markdown. Size: 1735 bytes. Shape: 55 lines. Markdown document, 55 lines, 2 explicit URLs. Lead heading: GitHub Actions Run Log. & \textbf{\path{samples/logs/20260318T023922Z-gh-runs-implement-step-1-dspy-training-path-follow-up.md}}. Тип: markdown. Размер: 1735 байт. Форма: 55 строк. Документ Markdown, 55 строк, явных URL: 2. Главный заголовок: GitHub Actions Run Log. \\
\textbf{\path{samples/logs/20260318T044242Z-gh-runs-add-maintained-repository-research-narrative.md}}. Kind: markdown. Size: 2222 bytes. Shape: 56 lines. Markdown document, 56 lines, 2 explicit URLs. Lead heading: GitHub Actions Run Log. & \textbf{\path{samples/logs/20260318T044242Z-gh-runs-add-maintained-repository-research-narrative.md}}. Тип: markdown. Размер: 2222 байт. Форма: 56 строк. Документ Markdown, 56 строк, явных URL: 2. Главный заголовок: GitHub Actions Run Log. \\
\textbf{\path{samples/logs/20260318T050051Z-gh-runs-add-azure-runtime-surfaces-for-do-1-and-2-and-carry-on-to-next.md}}. Kind: markdown. Size: 1828 bytes. Shape: 38 lines. Markdown document, 38 lines, 2 explicit URLs. Lead heading: GitHub Actions Run Log. & \textbf{\path{samples/logs/20260318T050051Z-gh-runs-add-azure-runtime-surfaces-for-do-1-and-2-and-carry-on-to-next.md}}. Тип: markdown. Размер: 1828 байт. Форма: 38 строк. Документ Markdown, 38 строк, явных URL: 2. Главный заголовок: GitHub Actions Run Log. \\
\textbf{\path{samples/logs/20260318T051627Z-gh-runs-develop-retrieval-evaluation-suite.md}}. Kind: markdown. Size: 1768 bytes. Shape: 37 lines. Markdown document, 37 lines, 2 explicit URLs. Lead heading: GitHub Actions Run Log. & \textbf{\path{samples/logs/20260318T051627Z-gh-runs-develop-retrieval-evaluation-suite.md}}. Тип: markdown. Размер: 1768 байт. Форма: 37 строк. Документ Markdown, 37 строк, явных URL: 2. Главный заголовок: GitHub Actions Run Log. \\
\textbf{\path{samples/logs/20260318T060918Z-gh-runs-consolidate-session-master-repair.md}}. Kind: markdown. Size: 2460 bytes. Shape: 70 lines. Markdown document, 70 lines, 3 explicit URLs. Lead heading: GitHub Actions Run Log. & \textbf{\path{samples/logs/20260318T060918Z-gh-runs-consolidate-session-master-repair.md}}. Тип: markdown. Размер: 2460 байт. Форма: 70 строк. Документ Markdown, 70 строк, явных URL: 3. Главный заголовок: GitHub Actions Run Log. \\
\textbf{\path{samples/logs/20260318T063755Z-gh-runs-keep-utility-sync-tests-side-effect-free.md}}. Kind: markdown. Size: 1074 bytes. Shape: 29 lines. Markdown document, 29 lines, 1 explicit URLs. Lead heading: GitHub Actions Run Log. & \textbf{\path{samples/logs/20260318T063755Z-gh-runs-keep-utility-sync-tests-side-effect-free.md}}. Тип: markdown. Размер: 1074 байт. Форма: 29 строк. Документ Markdown, 29 строк, явных URL: 1. Главный заголовок: GitHub Actions Run Log. \\
\textbf{\path{samples/logs/20260318T064924Z-gh-runs-retrieval-refresh-go-ahead.md}}. Kind: markdown. Size: 1858 bytes. Shape: 54 lines. Markdown document, 54 lines, 2 explicit URLs. Lead heading: GitHub Actions Run Log. & \textbf{\path{samples/logs/20260318T064924Z-gh-runs-retrieval-refresh-go-ahead.md}}. Тип: markdown. Размер: 1858 байт. Форма: 54 строк. Документ Markdown, 54 строк, явных URL: 2. Главный заголовок: GitHub Actions Run Log. \\
\textbf{\path{samples/logs/20260318T083955Z-gh-runs-benchmark-broadening-2.md}}. Kind: markdown. Size: 1217 bytes. Shape: 34 lines. Markdown document, 34 lines, 1 explicit URLs. Lead heading: GitHub Actions Run Log. & \textbf{\path{samples/logs/20260318T083955Z-gh-runs-benchmark-broadening-2.md}}. Тип: markdown. Размер: 1217 байт. Форма: 34 строк. Документ Markdown, 34 строк, явных URL: 1. Главный заголовок: GitHub Actions Run Log. \\
\textbf{\path{samples/logs/20260318T090625Z-gh-runs-hushwheel-quality-instrumentation.md}}. Kind: markdown. Size: 1312 bytes. Shape: 29 lines. Markdown document, 29 lines, 1 explicit URLs. Lead heading: GitHub Actions Run Log. & \textbf{\path{samples/logs/20260318T090625Z-gh-runs-hushwheel-quality-instrumentation.md}}. Тип: markdown. Размер: 1312 байт. Форма: 29 строк. Документ Markdown, 29 строк, явных URL: 1. Главный заголовок: GitHub Actions Run Log. \\
\textbf{\path{samples/logs/20260318T091252Z-gh-runs-add-rust-lookup-and-retrieval-tag-summaries-for-.md}}. Kind: markdown. Size: 2223 bytes. Shape: 53 lines. Markdown document, 53 lines, 2 explicit URLs. Lead heading: GitHub Actions Run Log. & \textbf{\path{samples/logs/20260318T091252Z-gh-runs-add-rust-lookup-and-retrieval-tag-summaries-for-.md}}. Тип: markdown. Размер: 2223 байт. Форма: 53 строк. Документ Markdown, 53 строк, явных URL: 2. Главный заголовок: GitHub Actions Run Log. \\
\textbf{\path{samples/logs/20260318T093632Z-gh-runs-retrieval-regression-gate-great-carry-on.md}}. Kind: markdown. Size: 1325 bytes. Shape: 34 lines. Markdown document, 34 lines, 2 explicit URLs. Lead heading: GitHub Actions Run Log. & \textbf{\path{samples/logs/20260318T093632Z-gh-runs-retrieval-regression-gate-great-carry-on.md}}. Тип: markdown. Размер: 1325 байт. Форма: 34 строк. Документ Markdown, 34 строк, явных URL: 2. Главный заголовок: GitHub Actions Run Log. \\
\textbf{\path{samples/logs/20260318T094120Z-gh-runs-dspy-live-proof-and-rust-wrapper-repair.md}}. Kind: markdown. Size: 1469 bytes. Shape: 45 lines. Markdown document, 45 lines, 0 explicit URLs. Lead heading: GitHub Actions Summary. & \textbf{\path{samples/logs/20260318T094120Z-gh-runs-dspy-live-proof-and-rust-wrapper-repair.md}}. Тип: markdown. Размер: 1469 байт. Форма: 45 строк. Документ Markdown, 45 строк, явных URL: 0. Главный заголовок: GitHub Actions Summary. \\
\textbf{\path{samples/logs/20260318T094347Z-gh-runs-hushwheel-hardening-followup.md}}. Kind: markdown. Size: 2739 bytes. Shape: 72 lines. Markdown document, 72 lines, 2 explicit URLs. Lead heading: GitHub Runs Log. & \textbf{\path{samples/logs/20260318T094347Z-gh-runs-hushwheel-hardening-followup.md}}. Тип: markdown. Размер: 2739 байт. Форма: 72 строк. Документ Markdown, 72 строк, явных URL: 2. Главный заголовок: GitHub Runs Log. \\
\textbf{\path{samples/logs/20260318T100336Z-gh-runs-sync-final-rebased-inventories-for-carry-on-comm.md}}. Kind: markdown. Size: 1103 bytes. Shape: 29 lines. Markdown document, 29 lines, 1 explicit URLs. Lead heading: GitHub Actions Run Log. & \textbf{\path{samples/logs/20260318T100336Z-gh-runs-sync-final-rebased-inventories-for-carry-on-comm.md}}. Тип: markdown. Размер: 1103 байт. Форма: 29 строк. Документ Markdown, 29 строк, явных URL: 1. Главный заголовок: GitHub Actions Run Log. \\
\textbf{\path{samples/logs/20260318T101229Z-gh-runs-hushwheel-audit-refresh.md}}. Kind: markdown. Size: 1707 bytes. Shape: 37 lines. Markdown document, 37 lines, 1 explicit URLs. Lead heading: GitHub Actions Summary. & \textbf{\path{samples/logs/20260318T101229Z-gh-runs-hushwheel-audit-refresh.md}}. Тип: markdown. Размер: 1707 байт. Форма: 37 строк. Документ Markdown, 37 строк, явных URL: 1. Главный заголовок: GitHub Actions Summary. \\
\textbf{\path{samples/logs/20260318T101939Z-gh-runs-land-each-unlanded.md}}. Kind: markdown. Size: 1556 bytes. Shape: 28 lines. Markdown document, 28 lines, 2 explicit URLs. Lead heading: GitHub Actions Run Log. & \textbf{\path{samples/logs/20260318T101939Z-gh-runs-land-each-unlanded.md}}. Тип: markdown. Размер: 1556 байт. Форма: 28 строк. Документ Markdown, 28 строк, явных URL: 2. Главный заголовок: GitHub Actions Run Log. \\
\textbf{\path{samples/logs/20260318T123402Z-gh-runs-fix-file-summary-test-churn-for-1.md}}. Kind: markdown. Size: 1061 bytes. Shape: 29 lines. Markdown document, 29 lines, 1 explicit URLs. Lead heading: GitHub Actions Run Log. & \textbf{\path{samples/logs/20260318T123402Z-gh-runs-fix-file-summary-test-churn-for-1.md}}. Тип: markdown. Размер: 1061 байт. Форма: 29 строк. Документ Markdown, 29 строк, явных URL: 1. Главный заголовок: GitHub Actions Run Log. \\
\textbf{\path{samples/logs/20260318T124437Z-gh-runs-reconcile-stale-dspy-backlog-for-1.md}}. Kind: markdown. Size: 1060 bytes. Shape: 21 lines. Markdown document, 21 lines, 3 explicit URLs. Lead heading: GitHub Actions Log. & \textbf{\path{samples/logs/20260318T124437Z-gh-runs-reconcile-stale-dspy-backlog-for-1.md}}. Тип: markdown. Размер: 1060 байт. Форма: 21 строк. Документ Markdown, 21 строк, явных URL: 3. Главный заголовок: GitHub Actions Log. \\
\textbf{\path{samples/logs/20260318T125547Z-gh-runs-wire-lookup-first-ask-for-1.md}}. Kind: markdown. Size: 1058 bytes. Shape: 22 lines. Markdown document, 22 lines, 2 explicit URLs. Lead heading: GitHub Actions Log. & \textbf{\path{samples/logs/20260318T125547Z-gh-runs-wire-lookup-first-ask-for-1.md}}. Тип: markdown. Размер: 1058 байт. Форма: 22 строк. Документ Markdown, 22 строк, явных URL: 2. Главный заголовок: GitHub Actions Log. \\
\textbf{\path{samples/logs/20260318T125829Z-gh-runs-restore-retrieval-gate-after-rebase.md}}. Kind: markdown. Size: 1585 bytes. Shape: 36 lines. Markdown document, 36 lines, 0 explicit URLs. Lead heading: GitHub Actions Run Log. & \textbf{\path{samples/logs/20260318T125829Z-gh-runs-restore-retrieval-gate-after-rebase.md}}. Тип: markdown. Размер: 1585 байт. Форма: 36 строк. Документ Markdown, 36 строк, явных URL: 0. Главный заголовок: GitHub Actions Run Log. \\
\textbf{\path{samples/logs/20260318T130525Z-gh-runs-finish-node24-publication-cache-followup-for-1.md}}. Kind: markdown. Size: 1119 bytes. Shape: 27 lines. Markdown document, 27 lines, 3 explicit URLs. Lead heading: GitHub Actions Status For `Finish Node 24 publication cache follow-up for "1"`. & \textbf{\path{samples/logs/20260318T130525Z-gh-runs-finish-node24-publication-cache-followup-for-1.md}}. Тип: markdown. Размер: 1119 байт. Форма: 27 строк. Документ Markdown, 27 строк, явных URL: 3. Главный заголовок: GitHub Actions Status For `Finish Node 24 publication cache follow-up for "1"`. \\
\textbf{\path{samples/logs/20260318T131709Z-gh-runs-restore-root-readme-weighting-for-1.md}}. Kind: markdown. Size: 1680 bytes. Shape: 45 lines. Markdown document, 45 lines, 2 explicit URLs. Lead heading: GitHub Actions Run Log. & \textbf{\path{samples/logs/20260318T131709Z-gh-runs-restore-root-readme-weighting-for-1.md}}. Тип: markdown. Размер: 1680 байт. Форма: 45 строк. Документ Markdown, 45 строк, явных URL: 2. Главный заголовок: GitHub Actions Run Log. \\
\textbf{\path{samples/logs/20260318T133247Z-gh-runs-improve-retrieval-doc-priority-for-1.md}}. Kind: markdown. Size: 1335 bytes. Shape: 34 lines. Markdown document, 34 lines, 1 explicit URLs. Lead heading: GitHub Actions Run Log. & \textbf{\path{samples/logs/20260318T133247Z-gh-runs-improve-retrieval-doc-priority-for-1.md}}. Тип: markdown. Размер: 1335 байт. Форма: 34 строк. Документ Markdown, 34 строк, явных URL: 1. Главный заголовок: GitHub Actions Run Log. \\
\textbf{\path{samples/logs/20260318T144102Z-gh-runs-merge-into-master.md}}. Kind: markdown. Size: 2285 bytes. Shape: 69 lines. Markdown document, 69 lines, 2 explicit URLs. Lead heading: GitHub Actions Run Log. & \textbf{\path{samples/logs/20260318T144102Z-gh-runs-merge-into-master.md}}. Тип: markdown. Размер: 2285 байт. Форма: 69 строк. Документ Markdown, 69 строк, явных URL: 2. Главный заголовок: GitHub Actions Run Log. \\
\textbf{\path{samples/logs/20260318T145935Z-gh-runs-consolidate-all-unconsoldiated-changes.md}}. Kind: markdown. Size: 2401 bytes. Shape: 71 lines. Markdown document, 71 lines, 2 explicit URLs. Lead heading: GitHub Actions Run Log. & \textbf{\path{samples/logs/20260318T145935Z-gh-runs-consolidate-all-unconsoldiated-changes.md}}. Тип: markdown. Размер: 2401 байт. Форма: 71 строк. Документ Markdown, 71 строк, явных URL: 2. Главный заголовок: GitHub Actions Run Log. \\
\textbf{\path{samples/logs/20260319T035301Z-gh-runs-public-pages-and-workflow-repairs.md}}. Kind: markdown. Size: 1578 bytes. Shape: 37 lines. Markdown document, 37 lines, 1 explicit URLs. Lead heading: GitHub Actions Log. & \textbf{\path{samples/logs/20260319T035301Z-gh-runs-public-pages-and-workflow-repairs.md}}. Тип: markdown. Размер: 1578 байт. Форма: 37 строк. Документ Markdown, 37 строк, явных URL: 1. Главный заголовок: GitHub Actions Log. \\
\textbf{\path{samples/logs/20260319T071544Z-gh-runs-doncsoldiate-consolidation.md}}. Kind: markdown. Size: 1137 bytes. Shape: 27 lines. Markdown document, 27 lines, 0 explicit URLs. Lead heading: GitHub Actions Log. & \textbf{\path{samples/logs/20260319T071544Z-gh-runs-doncsoldiate-consolidation.md}}. Тип: markdown. Размер: 1137 байт. Форма: 27 строк. Документ Markdown, 27 строк, явных URL: 0. Главный заголовок: GitHub Actions Log. \\
\textbf{\path{samples/logs/20260319T114913Z-gh-runs-fix-publication-verification-for-whats-next-step.md}}. Kind: markdown. Size: 2936 bytes. Shape: 67 lines. Markdown document, 67 lines, 3 explicit URLs. Lead heading: GitHub Actions Run Log. & \textbf{\path{samples/logs/20260319T114913Z-gh-runs-fix-publication-verification-for-whats-next-step.md}}. Тип: markdown. Размер: 2936 байт. Форма: 67 строк. Документ Markdown, 67 строк, явных URL: 3. Главный заголовок: GitHub Actions Run Log. \\
\textbf{\path{samples/logs/20260429T172137Z-gh-runs-develop-no-workflow.md}}. Kind: markdown. Size: 866 bytes. Shape: 23 lines. Markdown document, 23 lines, 0 explicit URLs. Lead heading: GitHub Actions Run Log. & \textbf{\path{samples/logs/20260429T172137Z-gh-runs-develop-no-workflow.md}}. Тип: markdown. Размер: 866 байт. Форма: 23 строк. Документ Markdown, 23 строк, явных URL: 0. Главный заголовок: GitHub Actions Run Log. \\
\textbf{\path{samples/logs/20260429T200446Z-gh-runs-develop-no-new-run-after-blob-queue-migration.md}}. Kind: markdown. Size: 1843 bytes. Shape: 34 lines. Markdown document, 34 lines, 0 explicit URLs. Lead heading: GitHub Actions Run Log. & \textbf{\path{samples/logs/20260429T200446Z-gh-runs-develop-no-new-run-after-blob-queue-migration.md}}. Тип: markdown. Размер: 1843 байт. Форма: 34 строк. Документ Markdown, 34 строк, явных URL: 0. Главный заголовок: GitHub Actions Run Log. \\
\textbf{\path{samples/logs/20260430T122222Z-gh-runs-develop-no-new-run-after-codex-proxy-push.md}}. Kind: markdown. Size: 908 bytes. Shape: 25 lines. Markdown document, 25 lines, 0 explicit URLs. Lead heading: GitHub Actions check after pushing `bbdcf4b`. & \textbf{\path{samples/logs/20260430T122222Z-gh-runs-develop-no-new-run-after-codex-proxy-push.md}}. Тип: markdown. Размер: 908 байт. Форма: 25 строк. Документ Markdown, 25 строк, явных URL: 0. Главный заголовок: GitHub Actions check after pushing `bbdcf4b`. \\
\textbf{\path{samples/logs/20260430T172518Z-gh-runs-develop-no-new-run-after-trace-handoff-push.md}}. Kind: markdown. Size: 991 bytes. Shape: 27 lines. Markdown document, 27 lines, 0 explicit URLs. Lead heading: GitHub Actions check after pushing `839f4cb`. & \textbf{\path{samples/logs/20260430T172518Z-gh-runs-develop-no-new-run-after-trace-handoff-push.md}}. Тип: markdown. Размер: 991 байт. Форма: 27 строк. Документ Markdown, 27 строк, явных URL: 0. Главный заголовок: GitHub Actions check after pushing `839f4cb`. \\
\textbf{\path{samples/logs/20260501T103130Z-gh-runs-develop-no-new-run-after-azure-runtime-env-fix-push.md}}. Kind: markdown. Size: 1411 bytes. Shape: 35 lines. Markdown document, 35 lines, 0 explicit URLs. Lead heading: GitHub Actions check after pushing `9e627e1`. & \textbf{\path{samples/logs/20260501T103130Z-gh-runs-develop-no-new-run-after-azure-runtime-env-fix-push.md}}. Тип: markdown. Размер: 1411 байт. Форма: 35 строк. Документ Markdown, 35 строк, явных URL: 0. Главный заголовок: GitHub Actions check after pushing `9e627e1`. \\
\textbf{\path{samples/logs/20260501T112330Z-gh-runs-develop-no-new-run-after-azure-responses-api-version-fix-push.md}}. Kind: markdown. Size: 1522 bytes. Shape: 37 lines. Markdown document, 37 lines, 0 explicit URLs. Lead heading: GitHub Actions check after pushing `128874b`. & \textbf{\path{samples/logs/20260501T112330Z-gh-runs-develop-no-new-run-after-azure-responses-api-version-fix-push.md}}. Тип: markdown. Размер: 1522 байт. Форма: 37 строк. Документ Markdown, 37 строк, явных URL: 0. Главный заголовок: GitHub Actions check after pushing `128874b`. \\
\textbf{\path{samples/logs/20260501T125058Z-gh-runs-develop-no-new-run-after-trusted-trace-handoff-push.md}}. Kind: markdown. Size: 1355 bytes. Shape: 30 lines. Markdown document, 30 lines, 0 explicit URLs. Lead heading: GitHub Actions check after pushing `66a7ad4`. & \textbf{\path{samples/logs/20260501T125058Z-gh-runs-develop-no-new-run-after-trusted-trace-handoff-push.md}}. Тип: markdown. Размер: 1355 байт. Форма: 30 строк. Документ Markdown, 30 строк, явных URL: 0. Главный заголовок: GitHub Actions check after pushing `66a7ad4`. \\
\textbf{\path{samples/logs/20260501T162125Z-gh-runs-develop-no-new-run-after-live-trainer-bundle-publication-push.md}}. Kind: markdown. Size: 1486 bytes. Shape: 30 lines. Markdown document, 30 lines, 0 explicit URLs. Lead heading: GitHub Run Inspection. & \textbf{\path{samples/logs/20260501T162125Z-gh-runs-develop-no-new-run-after-live-trainer-bundle-publication-push.md}}. Тип: markdown. Размер: 1486 байт. Форма: 30 строк. Документ Markdown, 30 строк, явных URL: 0. Главный заголовок: GitHub Run Inspection. \\
\textbf{\path{samples/logs/20260501T174147Z-gh-runs-develop-no-new-run-after-versioned-trainer-recovery-push.md}}. Kind: markdown. Size: 1505 bytes. Shape: 30 lines. Markdown document, 30 lines, 0 explicit URLs. Lead heading: GitHub Run Inspection. & \textbf{\path{samples/logs/20260501T174147Z-gh-runs-develop-no-new-run-after-versioned-trainer-recovery-push.md}}. Тип: markdown. Размер: 1505 байт. Форма: 30 строк. Документ Markdown, 30 строк, явных URL: 0. Главный заголовок: GitHub Run Inspection. \\
\textbf{\path{samples/logs/20260501T201242Z-gh-runs-develop-no-new-run-after-dspy-bundle-version-pinning-push.md}}. Kind: markdown. Size: 2076 bytes. Shape: 41 lines. Markdown document, 41 lines, 0 explicit URLs. Lead heading: GitHub Run Inspection. & \textbf{\path{samples/logs/20260501T201242Z-gh-runs-develop-no-new-run-after-dspy-bundle-version-pinning-push.md}}. Тип: markdown. Размер: 2076 байт. Форма: 41 строк. Документ Markdown, 41 строк, явных URL: 0. Главный заголовок: GitHub Run Inspection. \\
\textbf{\path{samples/logs/20260502T064702Z-gh-runs-develop-no-new-run-after-live-trainer-timestamp-bundle-audit-push.md}}. Kind: markdown. Size: 1604 bytes. Shape: 38 lines. Markdown document, 38 lines, 0 explicit URLs. Lead heading: GitHub Run Inspection. & \textbf{\path{samples/logs/20260502T064702Z-gh-runs-develop-no-new-run-after-live-trainer-timestamp-bundle-audit-push.md}}. Тип: markdown. Размер: 1604 байт. Форма: 38 строк. Документ Markdown, 38 строк, явных URL: 0. Главный заголовок: GitHub Run Inspection. \\
\textbf{\path{samples/logs/20260502T082402Z-gh-runs-develop-no-new-run-after-hybrid-vector-retrieval-push.md}}. Kind: markdown. Size: 1631 bytes. Shape: 30 lines. Markdown document, 30 lines, 0 explicit URLs. Lead heading: GitHub Run Inspection. & \textbf{\path{samples/logs/20260502T082402Z-gh-runs-develop-no-new-run-after-hybrid-vector-retrieval-push.md}}. Тип: markdown. Размер: 1631 байт. Форма: 30 строк. Документ Markdown, 30 строк, явных URL: 0. Главный заголовок: GitHub Run Inspection. \\
\textbf{\path{samples/logs/20260502T121755Z-gh-runs-develop-no-new-run-after-trainer-gating-and-retrieval-isolation-push.md}}. Kind: markdown. Size: 1175 bytes. Shape: 25 lines. Markdown document, 25 lines, 0 explicit URLs. Lead heading: Summary. & \textbf{\path{samples/logs/20260502T121755Z-gh-runs-develop-no-new-run-after-trainer-gating-and-retrieval-isolation-push.md}}. Тип: markdown. Размер: 1175 байт. Форма: 25 строк. Документ Markdown, 25 строк, явных URL: 0. Главный заголовок: Summary. \\
\textbf{\path{samples/logs/20260502T165750Z-gh-runs-develop-no-new-run-after-trainer-champion-batching-push.md}}. Kind: markdown. Size: 1543 bytes. Shape: 35 lines. Markdown document, 35 lines, 0 explicit URLs. Lead heading: GitHub Run Inspection. & \textbf{\path{samples/logs/20260502T165750Z-gh-runs-develop-no-new-run-after-trainer-champion-batching-push.md}}. Тип: markdown. Размер: 1543 байт. Форма: 35 строк. Документ Markdown, 35 строк, явных URL: 0. Главный заголовок: GitHub Run Inspection. \\
\textbf{\path{samples/logs/20260503T080503Z-gh-runs-develop-no-new-run-after-codex-resume-lane-telemetry-push.md}}. Kind: markdown. Size: 1782 bytes. Shape: 37 lines. Markdown document, 37 lines, 0 explicit URLs. Lead heading: GitHub Run Inspection. & \textbf{\path{samples/logs/20260503T080503Z-gh-runs-develop-no-new-run-after-codex-resume-lane-telemetry-push.md}}. Тип: markdown. Размер: 1782 байт. Форма: 37 строк. Документ Markdown, 37 строк, явных URL: 0. Главный заголовок: GitHub Run Inspection. \\
\textbf{\path{samples/logs/20260503T153706Z-gh-runs-develop-no-new-run-after-staged-bundle-mirror-push.md}}. Kind: markdown. Size: 1171 bytes. Shape: 27 lines. Markdown document, 27 lines, 0 explicit URLs. Lead heading: GitHub Runs After `Support staged bundle mirrors and stale queue skips`. & \textbf{\path{samples/logs/20260503T153706Z-gh-runs-develop-no-new-run-after-staged-bundle-mirror-push.md}}. Тип: markdown. Размер: 1171 байт. Форма: 27 строк. Документ Markdown, 27 строк, явных URL: 0. Главный заголовок: GitHub Runs After `Support staged bundle mirrors and stale queue skips`. \\
\textbf{\path{samples/logs/20260503T160232Z-gh-runs-develop-no-new-run-after-no-op-trainer-promotion-fix-push.md}}. Kind: markdown. Size: 1482 bytes. Shape: 32 lines. Markdown document, 32 lines, 0 explicit URLs. Lead heading: GitHub Run Inspection. & \textbf{\path{samples/logs/20260503T160232Z-gh-runs-develop-no-new-run-after-no-op-trainer-promotion-fix-push.md}}. Тип: markdown. Размер: 1482 байт. Форма: 32 строк. Документ Markdown, 32 строк, явных URL: 0. Главный заголовок: GitHub Run Inspection. \\
\textbf{\path{samples/logs/20260504T090704Z-gh-runs-develop-no-new-run-after-fix-codex-resume-pvc-session-path-push.md}}. Kind: markdown. Size: 1689 bytes. Shape: 46 lines. Markdown document, 46 lines, 0 explicit URLs. Lead heading: GitHub Run Check After `Fix Codex resume PVC session path`. & \textbf{\path{samples/logs/20260504T090704Z-gh-runs-develop-no-new-run-after-fix-codex-resume-pvc-session-path-push.md}}. Тип: markdown. Размер: 1689 байт. Форма: 46 строк. Документ Markdown, 46 строк, явных URL: 0. Главный заголовок: GitHub Run Check After `Fix Codex resume PVC session path`. \\
\textbf{\path{samples/logs/20260504T111427Z-gh-runs-develop-no-new-run-after-codex-resume-restore-fallback-debug-push.md}}. Kind: markdown. Size: 3370 bytes. Shape: 72 lines. Markdown document, 72 lines, 0 explicit URLs. Lead heading: GitHub Run Check After `Document Codex resume restore fallback debug`. & \textbf{\path{samples/logs/20260504T111427Z-gh-runs-develop-no-new-run-after-codex-resume-restore-fallback-debug-push.md}}. Тип: markdown. Размер: 3370 байт. Форма: 72 строк. Документ Markdown, 72 строк, явных URL: 0. Главный заголовок: GitHub Run Check After `Document Codex resume restore fallback debug`. \\
\textbf{\path{samples/logs/20260504T141246Z-gh-runs-develop-no-new-run-after-codex-resume-guard-hardening-push.md}}. Kind: markdown. Size: 2837 bytes. Shape: 67 lines. Markdown document, 67 lines, 0 explicit URLs. Lead heading: GitHub Run Check After `Harden Codex resume restore guards`. & \textbf{\path{samples/logs/20260504T141246Z-gh-runs-develop-no-new-run-after-codex-resume-guard-hardening-push.md}}. Тип: markdown. Размер: 2837 байт. Форма: 67 строк. Документ Markdown, 67 строк, явных URL: 0. Главный заголовок: GitHub Run Check After `Harden Codex resume restore guards`. \\
\textbf{\path{samples/logs/20260504T161106Z-gh-runs-develop-no-new-run-after-explicit-codex-resume-targeting-push.md}}. Kind: markdown. Size: 1175 bytes. Shape: 27 lines. Markdown document, 27 lines, 0 explicit URLs. Lead heading: GitHub Actions inspection after `464e16b` push. & \textbf{\path{samples/logs/20260504T161106Z-gh-runs-develop-no-new-run-after-explicit-codex-resume-targeting-push.md}}. Тип: markdown. Размер: 1175 байт. Форма: 27 строк. Документ Markdown, 27 строк, явных URL: 0. Главный заголовок: GitHub Actions inspection after `464e16b` push. \\
\textbf{\path{samples/logs/20260504T180441Z-gh-runs-develop-no-new-run-after-codex-session-continuity-markers-push.md}}. Kind: markdown. Size: 986 bytes. Shape: 28 lines. Markdown document, 28 lines, 0 explicit URLs. Lead heading: GitHub Actions inspection after `490fe41` push. & \textbf{\path{samples/logs/20260504T180441Z-gh-runs-develop-no-new-run-after-codex-session-continuity-markers-push.md}}. Тип: markdown. Размер: 986 байт. Форма: 28 строк. Документ Markdown, 28 строк, явных URL: 0. Главный заголовок: GitHub Actions inspection after `490fe41` push. \\
\textbf{\path{samples/logs/20260505T072441Z-gh-runs-develop-no-new-run-after-mcp-first-codex-discovery-push.md}}. Kind: markdown. Size: 1152 bytes. Shape: 28 lines. Markdown document, 28 lines, 0 explicit URLs. Lead heading: GitHub Actions inspection after `abcf52c` push. & \textbf{\path{samples/logs/20260505T072441Z-gh-runs-develop-no-new-run-after-mcp-first-codex-discovery-push.md}}. Тип: markdown. Размер: 1152 байт. Форма: 28 строк. Документ Markdown, 28 строк, явных URL: 0. Главный заголовок: GitHub Actions inspection after `abcf52c` push. \\
\textbf{\path{samples/logs/README.md}}. Kind: markdown. Size: 331 bytes. Shape: 11 lines. Markdown document, 11 lines, 0 explicit URLs. Lead heading: GitHub Run Logs. & \textbf{\path{samples/logs/README.md}}. Тип: markdown. Размер: 331 байт. Форма: 11 строк. Документ Markdown, 11 строк, явных URL: 0. Главный заголовок: GitHub Run Logs. \\
\textbf{\path{samples/population/hushwheel_fixture_population_candidates.yaml}}. Kind: yaml. Size: 955 bytes. Shape: 22 lines. YAML data or workflow file, 22 lines, 0 explicit URLs. Top keys: no top-level mapping keys. & \textbf{\path{samples/population/hushwheel_fixture_population_candidates.yaml}}. Тип: yaml. Размер: 955 байт. Форма: 22 строк. Файл YAML с данными или workflow, 22 строк, явных URL: 0. Верхние ключи: no top-level mapping keys. \\
\textbf{\path{samples/population/repository_population_candidates.yaml}}. Kind: yaml. Size: 690 bytes. Shape: 16 lines. YAML data or workflow file, 16 lines, 0 explicit URLs. Top keys: no top-level mapping keys. & \textbf{\path{samples/population/repository_population_candidates.yaml}}. Тип: yaml. Размер: 690 байт. Форма: 16 строк. Файл YAML с данными или workflow, 16 строк, явных URL: 0. Верхние ключи: no top-level mapping keys. \\
\textbf{\path{samples/training/hushwheel_fixture_training_examples.yaml}}. Kind: yaml. Size: 1852 bytes. Shape: 68 lines. YAML data or workflow file, 68 lines, 0 explicit URLs. Top keys: no top-level mapping keys. & \textbf{\path{samples/training/hushwheel_fixture_training_examples.yaml}}. Тип: yaml. Размер: 1852 байт. Форма: 68 строк. Файл YAML с данными или workflow, 68 строк, явных URL: 0. Верхние ключи: no top-level mapping keys. \\
\textbf{\path{samples/training/repository_training_examples.yaml}}. Kind: yaml. Size: 2504 bytes. Shape: 80 lines. YAML data or workflow file, 80 lines, 0 explicit URLs. Top keys: no top-level mapping keys. & \textbf{\path{samples/training/repository_training_examples.yaml}}. Тип: yaml. Размер: 2504 байт. Форма: 80 строк. Файл YAML с данными или workflow, 80 строк, явных URL: 0. Верхние ключи: no top-level mapping keys. \\
\textbf{\path{src/repo_rag_lab/__init__.py}}. Kind: python. Size: 1087 bytes. Shape: 47 lines. Python module, 47 lines, 0 explicit URLs. Module intent: Public package surface for the repository-grounded RAG lab. & \textbf{\path{src/repo_rag_lab/__init__.py}}. Тип: python. Размер: 1087 байт. Форма: 47 строк. Модуль Python, 47 строк, явных URL: 0. Назначение модуля: Public package surface for the repository-grounded RAG lab. \\
\textbf{\path{src/repo_rag_lab/azure.py}}. Kind: python. Size: 2592 bytes. Shape: 76 lines. Python module, 76 lines, 0 explicit URLs. Module intent: Azure deployment manifest helpers for downstream release workflows. & \textbf{\path{src/repo_rag_lab/azure.py}}. Тип: python. Размер: 2592 байт. Форма: 76 строк. Модуль Python, 76 строк, явных URL: 0. Назначение модуля: Azure deployment manifest helpers for downstream release workflows. \\
\textbf{\path{src/repo_rag_lab/azure_artifacts.py}}. Kind: python. Size: 16790 bytes. Shape: 436 lines. Python module, 436 lines, 0 explicit URLs. Module intent: Azure Blob + Queue helpers for trainer-side traces and global bundle storage. & \textbf{\path{src/repo_rag_lab/azure_artifacts.py}}. Тип: python. Размер: 16790 байт. Форма: 436 строк. Модуль Python, 436 строк, явных URL: 0. Назначение модуля: Azure Blob + Queue helpers for trainer-side traces and global bundle storage. \\
\textbf{\path{src/repo_rag_lab/azure_runtime.py}}. Kind: python. Size: 20123 bytes. Shape: 577 lines. Python module, 577 lines, 0 explicit URLs. Module intent: Azure runtime normalization, live probes, and cloud-completion helpers. & \textbf{\path{src/repo_rag_lab/azure_runtime.py}}. Тип: python. Размер: 20123 байт. Форма: 577 строк. Модуль Python, 577 строк, явных URL: 0. Назначение модуля: Azure runtime normalization, live probes, and cloud-completion helpers. \\
\textbf{\path{src/repo_rag_lab/benchmarks.py}}. Kind: python. Size: 16524 bytes. Shape: 484 lines. Python module, 484 lines, 0 explicit URLs. Module intent: Retrieval benchmark helpers for notebook assertions and logging. & \textbf{\path{src/repo_rag_lab/benchmarks.py}}. Тип: python. Размер: 16524 байт. Форма: 484 строк. Модуль Python, 484 строк, явных URL: 0. Назначение модуля: Retrieval benchmark helpers for notebook assertions and logging. \\
\textbf{\path{src/repo_rag_lab/cli.py}}. Kind: python. Size: 58870 bytes. Shape: 1339 lines. Python module, 1339 lines, 0 explicit URLs. Module intent: Command-line entrypoints for the shared repository RAG workflows. & \textbf{\path{src/repo_rag_lab/cli.py}}. Тип: python. Размер: 58870 байт. Форма: 1339 строк. Модуль Python, 1339 строк, явных URL: 0. Назначение модуля: Command-line entrypoints for the shared repository RAG workflows. \\
\textbf{\path{src/repo_rag_lab/codex_proxy.py}}. Kind: python. Size: 75460 bytes. Shape: 1946 lines. Python module, 1946 lines, 0 explicit URLs. Module intent: Local streaming proxy that injects repo-RAG + DSPy mediation for Codex. & \textbf{\path{src/repo_rag_lab/codex_proxy.py}}. Тип: python. Размер: 75460 байт. Форма: 1946 строк. Модуль Python, 1946 строк, явных URL: 0. Назначение модуля: Local streaming proxy that injects repo-RAG + DSPy mediation for Codex. \\
\textbf{\path{src/repo_rag_lab/corpus.py}}. Kind: python. Size: 5245 bytes. Shape: 181 lines. Python module, 181 lines, 0 explicit URLs. Module intent: Repository file loading helpers for the baseline text corpus. & \textbf{\path{src/repo_rag_lab/corpus.py}}. Тип: python. Размер: 5245 байт. Форма: 181 строк. Модуль Python, 181 строк, явных URL: 0. Назначение модуля: Repository file loading helpers for the baseline text corpus. \\
\textbf{\path{src/repo_rag_lab/dspy_training.py}}. Kind: python. Size: 61618 bytes. Shape: 1612 lines. Python module, 1612 lines, 0 explicit URLs. Module intent: DSPy training and artifact helpers for repository-grounded RAG. & \textbf{\path{src/repo_rag_lab/dspy_training.py}}. Тип: python. Размер: 61618 байт. Форма: 1612 строк. Модуль Python, 1612 строк, явных URL: 0. Назначение модуля: DSPy training and artifact helpers for repository-grounded RAG. \\
\textbf{\path{src/repo_rag_lab/dspy_workflow.py}}. Kind: python. Size: 6518 bytes. Shape: 180 lines. Python module, 180 lines, 0 explicit URLs. Module intent: DSPy-backed runtime wrappers around the repository retrieval workflow. & \textbf{\path{src/repo_rag_lab/dspy_workflow.py}}. Тип: python. Размер: 6518 байт. Форма: 180 строк. Модуль Python, 180 строк, явных URL: 0. Назначение модуля: DSPy-backed runtime wrappers around the repository retrieval workflow. \\
\textbf{\path{src/repo_rag_lab/exploratorium_translation.py}}. Kind: python. Size: 38049 bytes. Shape: 1009 lines. Python module, 1009 lines, 0 explicit URLs. Module intent: Generate the bilingual exploratorium translation publication surfaces. & \textbf{\path{src/repo_rag_lab/exploratorium_translation.py}}. Тип: python. Размер: 38049 байт. Форма: 1009 строк. Модуль Python, 1009 строк, явных URL: 0. Назначение модуля: Generate the bilingual exploratorium translation publication surfaces. \\
\textbf{\path{src/repo_rag_lab/file_summaries.py}}. Kind: python. Size: 10564 bytes. Shape: 327 lines. Python module, 327 lines, 0 explicit URLs. Module intent: Generate tracked repository file-summary surfaces. & \textbf{\path{src/repo_rag_lab/file_summaries.py}}. Тип: python. Размер: 10564 байт. Форма: 327 строк. Модуль Python, 327 строк, явных URL: 0. Назначение модуля: Generate tracked repository file-summary surfaces. \\
\textbf{\path{src/repo_rag_lab/github_pr_gates.py}}. Kind: python. Size: 3979 bytes. Shape: 140 lines. Python module, 140 lines, 0 explicit URLs. Module intent: GitHub branch-protection helpers for required pull-request gate checks. & \textbf{\path{src/repo_rag_lab/github_pr_gates.py}}. Тип: python. Размер: 3979 байт. Форма: 140 строк. Модуль Python, 140 строк, явных URL: 0. Назначение модуля: GitHub branch-protection helpers for required pull-request gate checks. \\
\textbf{\path{src/repo_rag_lab/mcp.py}}. Kind: python. Size: 2947 bytes. Shape: 84 lines. Python module, 84 lines, 0 explicit URLs. Module intent: Discovery helpers for MCP-related repository artifacts. & \textbf{\path{src/repo_rag_lab/mcp.py}}. Тип: python. Размер: 2947 байт. Форма: 84 строк. Модуль Python, 84 строк, явных URL: 0. Назначение модуля: Discovery helpers for MCP-related repository artifacts. \\
\textbf{\path{src/repo_rag_lab/mcp_server.py}}. Kind: python. Size: 50399 bytes. Shape: 1361 lines. Python module, 1361 lines, 0 explicit URLs. Module intent: Minimal stdio MCP server for bounded repo-RAG operations. & \textbf{\path{src/repo_rag_lab/mcp_server.py}}. Тип: python. Размер: 50399 байт. Форма: 1361 строк. Модуль Python, 1361 строк, явных URL: 0. Назначение модуля: Minimal stdio MCP server for bounded repo-RAG operations. \\
\textbf{\path{src/repo_rag_lab/mcp_stdio.py}}. Kind: python. Size: 953 bytes. Shape: 38 lines. Python module, 38 lines, 0 explicit URLs. Module intent: Lightweight stdio entrypoint for the bounded repo-RAG MCP server. & \textbf{\path{src/repo_rag_lab/mcp_stdio.py}}. Тип: python. Размер: 953 байт. Форма: 38 строк. Модуль Python, 38 строк, явных URL: 0. Назначение модуля: Lightweight stdio entrypoint for the bounded repo-RAG MCP server. \\
\textbf{\path{src/repo_rag_lab/notebook_runner.py}}. Kind: python. Size: 16433 bytes. Shape: 491 lines. Python module, 491 lines, 0 explicit URLs. Module intent: Monitored batch execution helpers for repository notebooks. & \textbf{\path{src/repo_rag_lab/notebook_runner.py}}. Тип: python. Размер: 16433 байт. Форма: 491 строк. Модуль Python, 491 строк, явных URL: 0. Назначение модуля: Monitored batch execution helpers for repository notebooks. \\
\textbf{\path{src/repo_rag_lab/notebook_scaffolding.py}}. Kind: python. Size: 10965 bytes. Shape: 258 lines. Python module, 258 lines, 0 explicit URLs. Module intent: High-level notebook scaffolds built from tested repository helpers. & \textbf{\path{src/repo_rag_lab/notebook_scaffolding.py}}. Тип: python. Размер: 10965 байт. Форма: 258 строк. Модуль Python, 258 строк, явных URL: 0. Назначение модуля: High-level notebook scaffolds built from tested repository helpers. \\
\textbf{\path{src/repo_rag_lab/notebook_support.py}}. Kind: python. Size: 2843 bytes. Shape: 89 lines. Python module, 89 lines, 0 explicit URLs. Module intent: Helpers shared by notebooks so their setup logic stays in tested Python code. & \textbf{\path{src/repo_rag_lab/notebook_support.py}}. Тип: python. Размер: 2843 байт. Форма: 89 строк. Модуль Python, 89 строк, явных URL: 0. Назначение модуля: Helpers shared by notebooks so their setup logic stays in tested Python code. \\
\textbf{\path{src/repo_rag_lab/pages_site.py}}. Kind: python. Size: 13635 bytes. Shape: 390 lines. Python module, 390 lines, 2 explicit URLs. Module intent: Generate a MkDocs-ready catalog of tracked repository Markdown files. & \textbf{\path{src/repo_rag_lab/pages_site.py}}. Тип: python. Размер: 13635 байт. Форма: 390 строк. Модуль Python, 390 строк, явных URL: 2. Назначение модуля: Generate a MkDocs-ready catalog of tracked repository Markdown files. \\
\textbf{\path{src/repo_rag_lab/population_samples.py}}. Kind: python. Size: 8376 bytes. Shape: 241 lines. Python module, 241 lines, 0 explicit URLs. Module intent: Helpers for reviewing and batching starter corpus-population inputs. & \textbf{\path{src/repo_rag_lab/population_samples.py}}. Тип: python. Размер: 8376 байт. Форма: 241 строк. Модуль Python, 241 строк, явных URL: 0. Назначение модуля: Helpers for reviewing and batching starter corpus-population inputs. \\
\textbf{\path{src/repo_rag_lab/retrieval.py}}. Kind: python. Size: 25533 bytes. Shape: 788 lines. Python module, 788 lines, 0 explicit URLs. Module intent: Baseline chunking and lexical retrieval utilities for repository text. & \textbf{\path{src/repo_rag_lab/retrieval.py}}. Тип: python. Размер: 25533 байт. Форма: 788 строк. Модуль Python, 788 строк, явных URL: 0. Назначение модуля: Baseline chunking and lexical retrieval utilities for repository text. \\
\textbf{\path{src/repo_rag_lab/retrieval_profile.py}}. Kind: python. Size: 9527 bytes. Shape: 236 lines. Python module, 236 lines, 0 explicit URLs. Module intent: Profile-driven retrieval weighting, exclusions, and repo-local overrides. & \textbf{\path{src/repo_rag_lab/retrieval_profile.py}}. Тип: python. Размер: 9527 байт. Форма: 236 строк. Модуль Python, 236 строк, явных URL: 0. Назначение модуля: Profile-driven retrieval weighting, exclusions, and repo-local overrides. \\
\textbf{\path{src/repo_rag_lab/runtime_artifacts.py}}. Kind: python. Size: 118360 bytes. Shape: 2671 lines. Python module, 2671 lines, 0 explicit URLs. Module intent: Versioned bundle, overlay, and runtime-trace helpers. & \textbf{\path{src/repo_rag_lab/runtime_artifacts.py}}. Тип: python. Размер: 118360 байт. Форма: 2671 строк. Модуль Python, 2671 строк, явных URL: 0. Назначение модуля: Versioned bundle, overlay, and runtime-trace helpers. \\
\textbf{\path{src/repo_rag_lab/rust_lookup.py}}. Kind: python. Size: 7724 bytes. Shape: 295 lines. Python module, 295 lines, 0 explicit URLs. Module intent: Helpers for the native Rust-backed SQLite lookup index. & \textbf{\path{src/repo_rag_lab/rust_lookup.py}}. Тип: python. Размер: 7724 байт. Форма: 295 строк. Модуль Python, 295 строк, явных URL: 0. Назначение модуля: Helpers for the native Rust-backed SQLite lookup index. \\
\textbf{\path{src/repo_rag_lab/semantic_retrieval.py}}. Kind: python. Size: 7407 bytes. Shape: 214 lines. Python module, 214 lines, 0 explicit URLs. Module intent: Azure OpenAI embedding-backed semantic retrieval helpers. & \textbf{\path{src/repo_rag_lab/semantic_retrieval.py}}. Тип: python. Размер: 7407 байт. Форма: 214 строк. Модуль Python, 214 строк, явных URL: 0. Назначение модуля: Azure OpenAI embedding-backed semantic retrieval helpers. \\
\textbf{\path{src/repo_rag_lab/settings.py}}. Kind: python. Size: 699 bytes. Shape: 28 lines. Python module, 28 lines, 0 explicit URLs. Module intent: Shared repository path settings derived from a repository root. & \textbf{\path{src/repo_rag_lab/settings.py}}. Тип: python. Размер: 699 байт. Форма: 28 строк. Модуль Python, 28 строк, явных URL: 0. Назначение модуля: Shared repository path settings derived from a repository root. \\
\textbf{\path{src/repo_rag_lab/todo_backlog.py}}. Kind: python. Size: 8434 bytes. Shape: 276 lines. Python module, 276 lines, 1 explicit URLs. Module intent: Shared backlog rendering helpers for Markdown and LaTeX surfaces. & \textbf{\path{src/repo_rag_lab/todo_backlog.py}}. Тип: python. Размер: 8434 байт. Форма: 276 строк. Модуль Python, 276 строк, явных URL: 1. Назначение модуля: Shared backlog rendering helpers for Markdown and LaTeX surfaces. \\
\textbf{\path{src/repo_rag_lab/trainer_deployment.py}}. Kind: python. Size: 18027 bytes. Shape: 431 lines. Python module, 431 lines, 0 explicit URLs. Module intent: Kubernetes manifest helpers for trainer-side repo-RAG deployment surfaces. & \textbf{\path{src/repo_rag_lab/trainer_deployment.py}}. Тип: python. Размер: 18027 байт. Форма: 431 строк. Модуль Python, 431 строк, явных URL: 0. Назначение модуля: Kubernetes manifest helpers for trainer-side repo-RAG deployment surfaces. \\
\textbf{\path{src/repo_rag_lab/training_samples.py}}. Kind: python. Size: 110198 bytes. Shape: 2709 lines. Python module, 2709 lines, 0 explicit URLs. Module intent: Helpers for loading and summarizing starter DSPy training examples. & \textbf{\path{src/repo_rag_lab/training_samples.py}}. Тип: python. Размер: 110198 байт. Форма: 2709 строк. Модуль Python, 2709 строк, явных URL: 0. Назначение модуля: Helpers for loading and summarizing starter DSPy training examples. \\
\textbf{\path{src/repo_rag_lab/utilities.py}}. Kind: python. Size: 116367 bytes. Shape: 2809 lines. Python module, 2809 lines, 1 explicit URLs. Module intent: User-facing utility helpers shared by the CLI, tests, and notebooks. & \textbf{\path{src/repo_rag_lab/utilities.py}}. Тип: python. Размер: 116367 байт. Форма: 2809 строк. Модуль Python, 2809 строк, явных URL: 1. Назначение модуля: User-facing utility helpers shared by the CLI, tests, and notebooks. \\
\textbf{\path{src/repo_rag_lab/verification.py}}. Kind: python. Size: 4476 bytes. Shape: 163 lines. Python module, 163 lines, 0 explicit URLs. Module intent: Verification helpers for Makefile and notebook contract checks. & \textbf{\path{src/repo_rag_lab/verification.py}}. Тип: python. Размер: 4476 байт. Форма: 163 строк. Модуль Python, 163 строк, явных URL: 0. Назначение модуля: Verification helpers for Makefile and notebook contract checks. \\
\textbf{\path{src/repo_rag_lab/workflow.py}}. Kind: python. Size: 15900 bytes. Shape: 503 lines. Python module, 503 lines, 0 explicit URLs. Module intent: Baseline repository question-answering workflow built on file retrieval. & \textbf{\path{src/repo_rag_lab/workflow.py}}. Тип: python. Размер: 15900 байт. Форма: 503 строк. Модуль Python, 503 строк, явных URL: 0. Назначение модуля: Baseline repository question-answering workflow built on file retrieval. \\
\textbf{\path{tests/features/repository_rag.feature}}. Kind: generic. Size: 377 bytes. Shape: 12 lines. Tracked repository file, 12 lines, 0 explicit URLs. Opening line: Feature: Repository RAG & \textbf{\path{tests/features/repository_rag.feature}}. Тип: generic. Размер: 377 байт. Форма: 12 строк. Отслеживаемый файл репозитория, 12 строк, явных URL: 0. Первая строка: Feature: Repository RAG \\
\textbf{\path{tests/fixtures/hushwheel_lexiconarium/.gitignore}}. Kind: generic. Size: 23 bytes. Shape: 4 lines. Tracked repository file, 4 lines, 0 explicit URLs. Opening line: hushwheel & \textbf{\path{tests/fixtures/hushwheel_lexiconarium/.gitignore}}. Тип: generic. Размер: 23 байт. Форма: 4 строк. Отслеживаемый файл репозитория, 4 строк, явных URL: 0. Первая строка: hushwheel \\
\textbf{\path{tests/fixtures/hushwheel_lexiconarium/CHANGELOG.md}}. Kind: markdown. Size: 345 bytes. Shape: 9 lines. Markdown document, 9 lines, 0 explicit URLs. Lead heading: Changelog. & \textbf{\path{tests/fixtures/hushwheel_lexiconarium/CHANGELOG.md}}. Тип: markdown. Размер: 345 байт. Форма: 9 строк. Документ Markdown, 9 строк, явных URL: 0. Главный заголовок: Changelog. \\
\textbf{\path{tests/fixtures/hushwheel_lexiconarium/Doxyfile}}. Kind: generic. Size: 1239 bytes. Shape: 35 lines. Tracked repository file, 35 lines, 0 explicit URLs. Opening line: PROJECT\_NAME           = Hushwheel Lexiconarium & \textbf{\path{tests/fixtures/hushwheel_lexiconarium/Doxyfile}}. Тип: generic. Размер: 1239 байт. Форма: 35 строк. Отслеживаемый файл репозитория, 35 строк, явных URL: 0. Первая строка: PROJECT\_NAME           = Hushwheel Lexiconarium \\
\textbf{\path{tests/fixtures/hushwheel_lexiconarium/LICENSE.md}}. Kind: markdown. Size: 1082 bytes. Shape: 22 lines. Markdown document, 22 lines, 0 explicit URLs. Lead heading: LICENSE. & \textbf{\path{tests/fixtures/hushwheel_lexiconarium/LICENSE.md}}. Тип: markdown. Размер: 1082 байт. Форма: 22 строк. Документ Markdown, 22 строк, явных URL: 0. Главный заголовок: LICENSE. \\
\textbf{\path{tests/fixtures/hushwheel_lexiconarium/Makefile}}. Kind: make. Size: 20908 bytes. Shape: 388 lines. Repository Makefile, 388 lines, 0 explicit URLs. It exposes the shared automation surfaces. & \textbf{\path{tests/fixtures/hushwheel_lexiconarium/Makefile}}. Тип: make. Размер: 20908 байт. Форма: 388 строк. Makefile репозитория, 388 строк, явных URL: 0. Он открывает общие поверхности автоматизации. \\
\textbf{\path{tests/fixtures/hushwheel_lexiconarium/README.md}}. Kind: markdown. Size: 8269 bytes. Shape: 162 lines. Markdown document, 162 lines, 0 explicit URLs. Lead heading: The Hushwheel Lexiconarium. & \textbf{\path{tests/fixtures/hushwheel_lexiconarium/README.md}}. Тип: markdown. Размер: 8269 байт. Форма: 162 строк. Документ Markdown, 162 строк, явных URL: 0. Главный заголовок: The Hushwheel Lexiconarium. \\
\textbf{\path{tests/fixtures/hushwheel_lexiconarium/VERSION}}. Kind: generic. Size: 6 bytes. Shape: 2 lines. Tracked repository file, 2 lines, 0 explicit URLs. Opening line: 0.1.0 & \textbf{\path{tests/fixtures/hushwheel_lexiconarium/VERSION}}. Тип: generic. Размер: 6 байт. Форма: 2 строк. Отслеживаемый файл репозитория, 2 строк, явных URL: 0. Первая строка: 0.1.0 \\
\textbf{\path{tests/fixtures/hushwheel_lexiconarium/docs/architecture.md}}. Kind: markdown. Size: 1404 bytes. Shape: 33 lines. Markdown document, 33 lines, 0 explicit URLs. Lead heading: Hushwheel Architecture. & \textbf{\path{tests/fixtures/hushwheel_lexiconarium/docs/architecture.md}}. Тип: markdown. Размер: 1404 байт. Форма: 33 строк. Документ Markdown, 33 строк, явных URL: 0. Главный заголовок: Hushwheel Architecture. \\
\textbf{\path{tests/fixtures/hushwheel_lexiconarium/docs/catalog.md}}. Kind: markdown. Size: 27404 bytes. Shape: 357 lines. Markdown document, 357 lines, 0 explicit URLs. Lead heading: Hushwheel Representative Catalog. & \textbf{\path{tests/fixtures/hushwheel_lexiconarium/docs/catalog.md}}. Тип: markdown. Размер: 27404 байт. Форма: 357 строк. Документ Markdown, 357 строк, явных URL: 0. Главный заголовок: Hushwheel Representative Catalog. \\
\textbf{\path{tests/fixtures/hushwheel_lexiconarium/docs/concepts.md}}. Kind: markdown. Size: 1559 bytes. Shape: 37 lines. Markdown document, 37 lines, 0 explicit URLs. Lead heading: Hushwheel Concepts. & \textbf{\path{tests/fixtures/hushwheel_lexiconarium/docs/concepts.md}}. Тип: markdown. Размер: 1559 байт. Форма: 37 строк. Документ Markdown, 37 строк, явных URL: 0. Главный заголовок: Hushwheel Concepts. \\
\textbf{\path{tests/fixtures/hushwheel_lexiconarium/docs/constellation-atlas.md}}. Kind: markdown. Size: 5111 bytes. Shape: 142 lines. Markdown document, 142 lines, 0 explicit URLs. Lead heading: Hushwheel Constellation Atlas §::<[]>. & \textbf{\path{tests/fixtures/hushwheel_lexiconarium/docs/constellation-atlas.md}}. Тип: markdown. Размер: 5111 байт. Форма: 142 строк. Документ Markdown, 142 строк, явных URL: 0. Главный заголовок: Hushwheel Constellation Atlas §::<[]>. \\
\textbf{\path{tests/fixtures/hushwheel_lexiconarium/docs/districts.md}}. Kind: markdown. Size: 1151 bytes. Shape: 42 lines. Markdown document, 42 lines, 0 explicit URLs. Lead heading: Hushwheel Districts. & \textbf{\path{tests/fixtures/hushwheel_lexiconarium/docs/districts.md}}. Тип: markdown. Размер: 1151 байт. Форма: 42 строк. Документ Markdown, 42 строк, явных URL: 0. Главный заголовок: Hushwheel Districts. \\
\textbf{\path{tests/fixtures/hushwheel_lexiconarium/docs/doxygen-fonts.sty}}. Kind: generic. Size: 212 bytes. Shape: binary surface. Tracked repository file, 212 bytes of binary data, 0 explicit URLs. Opening line: No non-empty preview line. & \textbf{\path{tests/fixtures/hushwheel_lexiconarium/docs/doxygen-fonts.sty}}. Тип: generic. Размер: 212 байт. Форма: бинарная поверхность. Отслеживаемый файл репозитория, 212 байт бинарных данных, явных URL: 0. Первая строка: No non-empty preview line. \\
\textbf{\path{tests/fixtures/hushwheel_lexiconarium/docs/hushwheel-reference.pdf}}. Kind: binary. Size: 132 bytes. Shape: binary surface. Binary asset, 132 bytes of binary data. It is tracked for stable publication output. & \textbf{\path{tests/fixtures/hushwheel_lexiconarium/docs/hushwheel-reference.pdf}}. Тип: binary. Размер: 132 байт. Форма: бинарная поверхность. Бинарный артефакт, 132 байт бинарных данных. Он хранится в Git для стабильного публикационного вывода. \\
\textbf{\path{tests/fixtures/hushwheel_lexiconarium/docs/operations.md}}. Kind: markdown. Size: 1068 bytes. Shape: 35 lines. Markdown document, 35 lines, 0 explicit URLs. Lead heading: Hushwheel Operations. & \textbf{\path{tests/fixtures/hushwheel_lexiconarium/docs/operations.md}}. Тип: markdown. Размер: 1068 байт. Форма: 35 строк. Документ Markdown, 35 строк, явных URL: 0. Главный заголовок: Hushwheel Operations. \\
\textbf{\path{tests/fixtures/hushwheel_lexiconarium/docs/packaging.md}}. Kind: markdown. Size: 1112 bytes. Shape: 33 lines. Markdown document, 33 lines, 0 explicit URLs. Lead heading: Hushwheel Packaging. & \textbf{\path{tests/fixtures/hushwheel_lexiconarium/docs/packaging.md}}. Тип: markdown. Размер: 1112 байт. Форма: 33 строк. Документ Markdown, 33 строк, явных URL: 0. Главный заголовок: Hushwheel Packaging. \\
\textbf{\path{tests/fixtures/hushwheel_lexiconarium/docs/testing.md}}. Kind: markdown. Size: 3507 bytes. Shape: 80 lines. Markdown document, 80 lines, 0 explicit URLs. Lead heading: Hushwheel Testing. & \textbf{\path{tests/fixtures/hushwheel_lexiconarium/docs/testing.md}}. Тип: markdown. Размер: 3507 байт. Форма: 80 строк. Документ Markdown, 80 строк, явных URL: 0. Главный заголовок: Hushwheel Testing. \\
\textbf{\path{tests/fixtures/hushwheel_lexiconarium/fixture-manifest.json}}. Kind: json. Size: 1098 bytes. Shape: 49 lines. JSON data file, 49 lines, 0 explicit URLs. Top keys: entry\_count, source\_size\_bytes, doc\_files. & \textbf{\path{tests/fixtures/hushwheel_lexiconarium/fixture-manifest.json}}. Тип: json. Размер: 1098 байт. Форма: 49 строк. Файл JSON, 49 строк, явных URL: 0. Верхние ключи: entry\_count, source\_size\_bytes, doc\_files. \\
\textbf{\path{tests/fixtures/hushwheel_lexiconarium/include/hushwheel.h}}. Kind: source. Size: 404 bytes. Shape: 19 lines. Source file, 19 lines, 0 explicit URLs. Opening line: \#ifndef HUSHWHEEL\_H & \textbf{\path{tests/fixtures/hushwheel_lexiconarium/include/hushwheel.h}}. Тип: source. Размер: 404 байт. Форма: 19 строк. Исходный файл, 19 строк, явных URL: 0. Первая строка: \#ifndef HUSHWHEEL\_H \\
\textbf{\path{tests/fixtures/hushwheel_lexiconarium/packaging/hushwheel.1}}. Kind: generic. Size: 1303 bytes. Shape: binary surface. Tracked repository file, 1303 bytes of binary data, 0 explicit URLs. Opening line: No non-empty preview line. & \textbf{\path{tests/fixtures/hushwheel_lexiconarium/packaging/hushwheel.1}}. Тип: generic. Размер: 1303 байт. Форма: бинарная поверхность. Отслеживаемый файл репозитория, 1303 байт бинарных данных, явных URL: 0. Первая строка: No non-empty preview line. \\
\textbf{\path{tests/fixtures/hushwheel_lexiconarium/packaging/hushwheel.package.json}}. Kind: json. Size: 591 bytes. Shape: 29 lines. JSON data file, 29 lines, 0 explicit URLs. Top keys: name, version, summary. & \textbf{\path{tests/fixtures/hushwheel_lexiconarium/packaging/hushwheel.package.json}}. Тип: json. Размер: 591 байт. Форма: 29 строк. Файл JSON, 29 строк, явных URL: 0. Верхние ключи: name, version, summary. \\
\textbf{\path{tests/fixtures/hushwheel_lexiconarium/src/hushwheel.c}}. Kind: source. Size: 1694579 bytes. Shape: 4375 lines. Source file, 4375 lines, 0 explicit URLs. Opening line: /** & \textbf{\path{tests/fixtures/hushwheel_lexiconarium/src/hushwheel.c}}. Тип: source. Размер: 1694579 байт. Форма: 4375 строк. Исходный файл, 4375 строк, явных URL: 0. Первая строка: /** \\
\textbf{\path{tests/fixtures/hushwheel_lexiconarium/src/hushwheel_internal.h}}. Kind: source. Size: 603 bytes. Shape: 25 lines. Source file, 25 lines, 0 explicit URLs. Opening line: \#ifndef HUSHWHEEL\_INTERNAL\_H & \textbf{\path{tests/fixtures/hushwheel_lexiconarium/src/hushwheel_internal.h}}. Тип: source. Размер: 603 байт. Форма: 25 строк. Исходный файл, 25 строк, явных URL: 0. Первая строка: \#ifndef HUSHWHEEL\_INTERNAL\_H \\
\textbf{\path{tests/fixtures/hushwheel_lexiconarium/src/hushwheel_spoke_argent.c}}. Kind: source. Size: 324091 bytes. Shape: 525 lines. Source file, 525 lines, 0 explicit URLs. Opening line: \#include "hushwheel\_internal.h" & \textbf{\path{tests/fixtures/hushwheel_lexiconarium/src/hushwheel_spoke_argent.c}}. Тип: source. Размер: 324091 байт. Форма: 525 строк. Исходный файл, 525 строк, явных URL: 0. Первая строка: \#include "hushwheel\_internal.h" \\
\textbf{\path{tests/fixtures/hushwheel_lexiconarium/src/hushwheel_spoke_bhagavatam.c}}. Kind: source. Size: 326160 bytes. Shape: 525 lines. Source file, 525 lines, 0 explicit URLs. Opening line: \#include "hushwheel\_internal.h" & \textbf{\path{tests/fixtures/hushwheel_lexiconarium/src/hushwheel_spoke_bhagavatam.c}}. Тип: source. Размер: 326160 байт. Форма: 525 строк. Исходный файл, 525 строк, явных URL: 0. Первая строка: \#include "hushwheel\_internal.h" \\
\textbf{\path{tests/fixtures/hushwheel_lexiconarium/src/hushwheel_spoke_crosscanon.c}}. Kind: source. Size: 326144 bytes. Shape: 525 lines. Source file, 525 lines, 0 explicit URLs. Opening line: \#include "hushwheel\_internal.h" & \textbf{\path{tests/fixtures/hushwheel_lexiconarium/src/hushwheel_spoke_crosscanon.c}}. Тип: source. Размер: 326144 байт. Форма: 525 строк. Исходный файл, 525 строк, явных URL: 0. Первая строка: \#include "hushwheel\_internal.h" \\
\textbf{\path{tests/fixtures/hushwheel_lexiconarium/src/hushwheel_spoke_editorial.c}}. Kind: source. Size: 325628 bytes. Shape: 525 lines. Source file, 525 lines, 0 explicit URLs. Opening line: \#include "hushwheel\_internal.h" & \textbf{\path{tests/fixtures/hushwheel_lexiconarium/src/hushwheel_spoke_editorial.c}}. Тип: source. Размер: 325628 байт. Форма: 525 строк. Исходный файл, 525 строк, явных URL: 0. Первая строка: \#include "hushwheel\_internal.h" \\
\textbf{\path{tests/fixtures/hushwheel_lexiconarium/src/hushwheel_spoke_gita.c}}. Kind: source. Size: 323055 bytes. Shape: 525 lines. Source file, 525 lines, 0 explicit URLs. Opening line: \#include "hushwheel\_internal.h" & \textbf{\path{tests/fixtures/hushwheel_lexiconarium/src/hushwheel_spoke_gita.c}}. Тип: source. Размер: 323055 байт. Форма: 525 строк. Исходный файл, 525 строк, явных URL: 0. Первая строка: \#include "hushwheel\_internal.h" \\
\textbf{\path{tests/fixtures/hushwheel_lexiconarium/src/hushwheel_spoke_kernel.c}}. Kind: source. Size: 324068 bytes. Shape: 525 lines. Source file, 525 lines, 0 explicit URLs. Opening line: \#include "hushwheel\_internal.h" & \textbf{\path{tests/fixtures/hushwheel_lexiconarium/src/hushwheel_spoke_kernel.c}}. Тип: source. Размер: 324068 байт. Форма: 525 строк. Исходный файл, 525 строк, явных URL: 0. Первая строка: \#include "hushwheel\_internal.h" \\
\textbf{\path{tests/fixtures/hushwheel_lexiconarium/src/hushwheel_spoke_psalter.c}}. Kind: source. Size: 324600 bytes. Shape: 525 lines. Source file, 525 lines, 0 explicit URLs. Opening line: \#include "hushwheel\_internal.h" & \textbf{\path{tests/fixtures/hushwheel_lexiconarium/src/hushwheel_spoke_psalter.c}}. Тип: source. Размер: 324600 байт. Форма: 525 строк. Исходный файл, 525 строк, явных URL: 0. Первая строка: \#include "hushwheel\_internal.h" \\
\textbf{\path{tests/fixtures/hushwheel_lexiconarium/src/hushwheel_spoke_ranger.c}}. Kind: source. Size: 324105 bytes. Shape: 525 lines. Source file, 525 lines, 0 explicit URLs. Opening line: \#include "hushwheel\_internal.h" & \textbf{\path{tests/fixtures/hushwheel_lexiconarium/src/hushwheel_spoke_ranger.c}}. Тип: source. Размер: 324105 байт. Форма: 525 строк. Исходный файл, 525 строк, явных URL: 0. Первая строка: \#include "hushwheel\_internal.h" \\
\textbf{\path{tests/fixtures/hushwheel_lexiconarium/src/hushwheel_spokes.c}}. Kind: source. Size: 1074 bytes. Shape: 25 lines. Source file, 25 lines, 0 explicit URLs. Opening line: \#include "hushwheel\_internal.h" & \textbf{\path{tests/fixtures/hushwheel_lexiconarium/src/hushwheel_spokes.c}}. Тип: source. Размер: 1074 байт. Форма: 25 строк. Исходный файл, 25 строк, явных URL: 0. Первая строка: \#include "hushwheel\_internal.h" \\
\textbf{\path{tests/fixtures/hushwheel_lexiconarium/tests/bdd/hushwheel.feature}}. Kind: generic. Size: 1021 bytes. Shape: 28 lines. Tracked repository file, 28 lines, 0 explicit URLs. Opening line: Feature: Hushwheel CLI & \textbf{\path{tests/fixtures/hushwheel_lexiconarium/tests/bdd/hushwheel.feature}}. Тип: generic. Размер: 1021 байт. Форма: 28 строк. Отслеживаемый файл репозитория, 28 строк, явных URL: 0. Первая строка: Feature: Hushwheel CLI \\
\textbf{\path{tests/fixtures/hushwheel_lexiconarium/tests/bdd/run_bdd.py}}. Kind: python. Size: 3146 bytes. Shape: 94 lines. Python module, 94 lines, 0 explicit URLs. Module intent: No module docstring detected. & \textbf{\path{tests/fixtures/hushwheel_lexiconarium/tests/bdd/run_bdd.py}}. Тип: python. Размер: 3146 байт. Форма: 94 строк. Модуль Python, 94 строк, явных URL: 0. Назначение модуля: No module docstring detected. \\
\textbf{\path{tests/fixtures/hushwheel_lexiconarium/tests/integration/cli_suite.py}}. Kind: python. Size: 4482 bytes. Shape: 122 lines. Python module, 122 lines, 0 explicit URLs. Module intent: No module docstring detected. & \textbf{\path{tests/fixtures/hushwheel_lexiconarium/tests/integration/cli_suite.py}}. Тип: python. Размер: 4482 байт. Форма: 122 строк. Модуль Python, 122 строк, явных URL: 0. Назначение модуля: No module docstring detected. \\
\textbf{\path{tests/fixtures/hushwheel_lexiconarium/tests/unit/test_hushwheel_unit.c}}. Kind: source. Size: 2002 bytes. Shape: 62 lines. Source file, 62 lines, 0 explicit URLs. Opening line: \#include <assert.h> & \textbf{\path{tests/fixtures/hushwheel_lexiconarium/tests/unit/test_hushwheel_unit.c}}. Тип: source. Размер: 2002 байт. Форма: 62 строк. Исходный файл, 62 строк, явных URL: 0. Первая строка: \#include <assert.h> \\
\textbf{\path{tests/fixtures/hushwheel_lexiconarium/tools/lint_hushwheel.py}}. Kind: python. Size: 4678 bytes. Shape: 143 lines. Python module, 143 lines, 0 explicit URLs. Module intent: No module docstring detected. & \textbf{\path{tests/fixtures/hushwheel_lexiconarium/tools/lint_hushwheel.py}}. Тип: python. Размер: 4678 байт. Форма: 143 строк. Модуль Python, 143 строк, явных URL: 0. Назначение модуля: No module docstring detected. \\
\textbf{\path{tests/fixtures/hushwheel_lexiconarium/tools/regenerate_hushwheel_fixture.py}}. Kind: python. Size: 31375 bytes. Shape: 866 lines. Python module, 866 lines, 0 explicit URLs. Module intent: No module docstring detected. & \textbf{\path{tests/fixtures/hushwheel_lexiconarium/tools/regenerate_hushwheel_fixture.py}}. Тип: python. Размер: 31375 байт. Форма: 866 строк. Модуль Python, 866 строк, явных URL: 0. Назначение модуля: No module docstring detected. \\
\textbf{\path{tests/test_azure_runtime.py}}. Kind: python. Size: 7540 bytes. Shape: 208 lines. Python module, 208 lines, 12 explicit URLs. Module intent: No module docstring detected. & \textbf{\path{tests/test_azure_runtime.py}}. Тип: python. Размер: 7540 байт. Форма: 208 строк. Модуль Python, 208 строк, явных URL: 12. Назначение модуля: No module docstring detected. \\
\textbf{\path{tests/test_benchmarks_and_notebook_scaffolding.py}}. Kind: python. Size: 11859 bytes. Shape: 274 lines. Python module, 274 lines, 0 explicit URLs. Module intent: No module docstring detected. & \textbf{\path{tests/test_benchmarks_and_notebook_scaffolding.py}}. Тип: python. Размер: 11859 байт. Форма: 274 строк. Модуль Python, 274 строк, явных URL: 0. Назначение модуля: No module docstring detected. \\
\textbf{\path{tests/test_cli_and_dspy.py}}. Kind: python. Size: 68622 bytes. Shape: 1876 lines. Python module, 1876 lines, 2 explicit URLs. Module intent: No module docstring detected. & \textbf{\path{tests/test_cli_and_dspy.py}}. Тип: python. Размер: 68622 байт. Форма: 1876 строк. Модуль Python, 1876 строк, явных URL: 2. Назначение модуля: No module docstring detected. \\
\textbf{\path{tests/test_codex_proxy.py}}. Kind: python. Size: 57154 bytes. Shape: 1571 lines. Python module, 1571 lines, 5 explicit URLs. Module intent: No module docstring detected. & \textbf{\path{tests/test_codex_proxy.py}}. Тип: python. Размер: 57154 байт. Форма: 1571 строк. Модуль Python, 1571 строк, явных URL: 5. Назначение модуля: No module docstring detected. \\
\textbf{\path{tests/test_dspy_training.py}}. Kind: python. Size: 66652 bytes. Shape: 1653 lines. Python module, 1653 lines, 11 explicit URLs. Module intent: No module docstring detected. & \textbf{\path{tests/test_dspy_training.py}}. Тип: python. Размер: 66652 байт. Форма: 1653 строк. Модуль Python, 1653 строк, явных URL: 11. Назначение модуля: No module docstring detected. \\
\textbf{\path{tests/test_exploratorium_translation.py}}. Kind: python. Size: 5239 bytes. Shape: 138 lines. Python module, 138 lines, 8 explicit URLs. Module intent: No module docstring detected. & \textbf{\path{tests/test_exploratorium_translation.py}}. Тип: python. Размер: 5239 байт. Форма: 138 строк. Модуль Python, 138 строк, явных URL: 8. Назначение модуля: No module docstring detected. \\
\textbf{\path{tests/test_file_summaries.py}}. Kind: python. Size: 9110 bytes. Shape: 230 lines. Python module, 230 lines, 0 explicit URLs. Module intent: No module docstring detected. & \textbf{\path{tests/test_file_summaries.py}}. Тип: python. Размер: 9110 байт. Форма: 230 строк. Модуль Python, 230 строк, явных URL: 0. Назначение модуля: No module docstring detected. \\
\textbf{\path{tests/test_github_pr_gates.py}}. Kind: python. Size: 5678 bytes. Shape: 183 lines. Python module, 183 lines, 1 explicit URLs. Module intent: No module docstring detected. & \textbf{\path{tests/test_github_pr_gates.py}}. Тип: python. Размер: 5678 байт. Форма: 183 строк. Модуль Python, 183 строк, явных URL: 1. Назначение модуля: No module docstring detected. \\
\textbf{\path{tests/test_hushwheel_fixture.py}}. Kind: python. Size: 4617 bytes. Shape: 101 lines. Python module, 101 lines, 0 explicit URLs. Module intent: No module docstring detected. & \textbf{\path{tests/test_hushwheel_fixture.py}}. Тип: python. Размер: 4617 байт. Форма: 101 строк. Модуль Python, 101 строк, явных URL: 0. Назначение модуля: No module docstring detected. \\
\textbf{\path{tests/test_hushwheel_program_surface.py}}. Kind: python. Size: 4626 bytes. Shape: 125 lines. Python module, 125 lines, 0 explicit URLs. Module intent: No module docstring detected. & \textbf{\path{tests/test_hushwheel_program_surface.py}}. Тип: python. Размер: 4626 байт. Форма: 125 строк. Модуль Python, 125 строк, явных URL: 0. Назначение модуля: No module docstring detected. \\
\textbf{\path{tests/test_live_azure_integration.py}}. Kind: python. Size: 2634 bytes. Shape: 83 lines. Python module, 83 lines, 0 explicit URLs. Module intent: No module docstring detected. & \textbf{\path{tests/test_live_azure_integration.py}}. Тип: python. Размер: 2634 байт. Форма: 83 строк. Модуль Python, 83 строк, явных URL: 0. Назначение модуля: No module docstring detected. \\
\textbf{\path{tests/test_lookup_first.py}}. Kind: python. Size: 3559 bytes. Shape: 108 lines. Python module, 108 lines, 0 explicit URLs. Module intent: No module docstring detected. & \textbf{\path{tests/test_lookup_first.py}}. Тип: python. Размер: 3559 байт. Форма: 108 строк. Модуль Python, 108 строк, явных URL: 0. Назначение модуля: No module docstring detected. \\
\textbf{\path{tests/test_mcp_branches.py}}. Kind: python. Size: 1077 bytes. Shape: 31 lines. Python module, 31 lines, 0 explicit URLs. Module intent: No module docstring detected. & \textbf{\path{tests/test_mcp_branches.py}}. Тип: python. Размер: 1077 байт. Форма: 31 строк. Модуль Python, 31 строк, явных URL: 0. Назначение модуля: No module docstring detected. \\
\textbf{\path{tests/test_mcp_server.py}}. Kind: python. Size: 26557 bytes. Shape: 799 lines. Python module, 799 lines, 0 explicit URLs. Module intent: No module docstring detected. & \textbf{\path{tests/test_mcp_server.py}}. Тип: python. Размер: 26557 байт. Форма: 799 строк. Модуль Python, 799 строк, явных URL: 0. Назначение модуля: No module docstring detected. \\
\textbf{\path{tests/test_mcp_stdio.py}}. Kind: python. Size: 988 bytes. Shape: 42 lines. Python module, 42 lines, 0 explicit URLs. Module intent: No module docstring detected. & \textbf{\path{tests/test_mcp_stdio.py}}. Тип: python. Размер: 988 байт. Форма: 42 строк. Модуль Python, 42 строк, явных URL: 0. Назначение модуля: No module docstring detected. \\
\textbf{\path{tests/test_notebook_runner.py}}. Kind: python. Size: 3088 bytes. Shape: 84 lines. Python module, 84 lines, 0 explicit URLs. Module intent: No module docstring detected. & \textbf{\path{tests/test_notebook_runner.py}}. Тип: python. Размер: 3088 байт. Форма: 84 строк. Модуль Python, 84 строк, явных URL: 0. Назначение модуля: No module docstring detected. \\
\textbf{\path{tests/test_pages_site.py}}. Kind: python. Size: 9045 bytes. Shape: 225 lines. Python module, 225 lines, 14 explicit URLs. Module intent: No module docstring detected. & \textbf{\path{tests/test_pages_site.py}}. Тип: python. Размер: 9045 байт. Форма: 225 строк. Модуль Python, 225 строк, явных URL: 14. Назначение модуля: No module docstring detected. \\
\textbf{\path{tests/test_population_samples.py}}. Kind: python. Size: 3144 bytes. Shape: 76 lines. Python module, 76 lines, 0 explicit URLs. Module intent: No module docstring detected. & \textbf{\path{tests/test_population_samples.py}}. Тип: python. Размер: 3144 байт. Форма: 76 строк. Модуль Python, 76 строк, явных URL: 0. Назначение модуля: No module docstring detected. \\
\textbf{\path{tests/test_project_surfaces.py}}. Kind: python. Size: 21471 bytes. Shape: 484 lines. Python module, 484 lines, 1 explicit URLs. Module intent: No module docstring detected. & \textbf{\path{tests/test_project_surfaces.py}}. Тип: python. Размер: 21471 байт. Форма: 484 строк. Модуль Python, 484 строк, явных URL: 1. Назначение модуля: No module docstring detected. \\
\textbf{\path{tests/test_repository_rag_bdd.py}}. Kind: python. Size: 879 bytes. Shape: 30 lines. Python module, 30 lines, 0 explicit URLs. Module intent: No module docstring detected. & \textbf{\path{tests/test_repository_rag_bdd.py}}. Тип: python. Размер: 879 байт. Форма: 30 строк. Модуль Python, 30 строк, явных URL: 0. Назначение модуля: No module docstring detected. \\
\textbf{\path{tests/test_retrieval.py}}. Kind: python. Size: 13362 bytes. Shape: 389 lines. Python module, 389 lines, 0 explicit URLs. Module intent: No module docstring detected. & \textbf{\path{tests/test_retrieval.py}}. Тип: python. Размер: 13362 байт. Форма: 389 строк. Модуль Python, 389 строк, явных URL: 0. Назначение модуля: No module docstring detected. \\
\textbf{\path{tests/test_runtime_artifacts_azure.py}}. Kind: python. Size: 58883 bytes. Shape: 1554 lines. Python module, 1554 lines, 4 explicit URLs. Module intent: No module docstring detected. & \textbf{\path{tests/test_runtime_artifacts_azure.py}}. Тип: python. Размер: 58883 байт. Форма: 1554 строк. Модуль Python, 1554 строк, явных URL: 4. Назначение модуля: No module docstring detected. \\
\textbf{\path{tests/test_rust_lookup.py}}. Kind: python. Size: 3099 bytes. Shape: 95 lines. Python module, 95 lines, 0 explicit URLs. Module intent: No module docstring detected. & \textbf{\path{tests/test_rust_lookup.py}}. Тип: python. Размер: 3099 байт. Форма: 95 строк. Модуль Python, 95 строк, явных URL: 0. Назначение модуля: No module docstring detected. \\
\textbf{\path{tests/test_todo_backlog.py}}. Kind: python. Size: 3505 bytes. Shape: 103 lines. Python module, 103 lines, 1 explicit URLs. Module intent: No module docstring detected. & \textbf{\path{tests/test_todo_backlog.py}}. Тип: python. Размер: 3505 байт. Форма: 103 строк. Модуль Python, 103 строк, явных URL: 1. Назначение модуля: No module docstring detected. \\
\textbf{\path{tests/test_training_samples.py}}. Kind: python. Size: 84870 bytes. Shape: 2152 lines. Python module, 2152 lines, 0 explicit URLs. Module intent: No module docstring detected. & \textbf{\path{tests/test_training_samples.py}}. Тип: python. Размер: 84870 байт. Форма: 2152 строк. Модуль Python, 2152 строк, явных URL: 0. Назначение модуля: No module docstring detected. \\
\textbf{\path{tests/test_utilities.py}}. Kind: python. Size: 121479 bytes. Shape: 3117 lines. Python module, 3117 lines, 6 explicit URLs. Module intent: No module docstring detected. & \textbf{\path{tests/test_utilities.py}}. Тип: python. Размер: 121479 байт. Форма: 3117 строк. Модуль Python, 3117 строк, явных URL: 6. Назначение модуля: No module docstring detected. \\
\textbf{\path{tests/test_verification.py}}. Kind: python. Size: 3156 bytes. Shape: 66 lines. Python module, 66 lines, 0 explicit URLs. Module intent: No module docstring detected. & \textbf{\path{tests/test_verification.py}}. Тип: python. Размер: 3156 байт. Форма: 66 строк. Модуль Python, 66 строк, явных URL: 0. Назначение модуля: No module docstring detected. \\
\textbf{\path{tests/test_workflow.py}}. Kind: python. Size: 3456 bytes. Shape: 99 lines. Python module, 99 lines, 0 explicit URLs. Module intent: No module docstring detected. & \textbf{\path{tests/test_workflow.py}}. Тип: python. Размер: 3456 байт. Форма: 99 строк. Модуль Python, 99 строк, явных URL: 0. Назначение модуля: No module docstring detected. \\
\textbf{\path{tests/test_workflow_live.py}}. Kind: python. Size: 4292 bytes. Shape: 128 lines. Python module, 128 lines, 0 explicit URLs. Module intent: No module docstring detected. & \textbf{\path{tests/test_workflow_live.py}}. Тип: python. Размер: 4292 байт. Форма: 128 строк. Модуль Python, 128 строк, явных URL: 0. Назначение модуля: No module docstring detected. \\
\textbf{\path{todo-backlog.yaml}}. Kind: yaml. Size: 6116 bytes. Shape: 128 lines. YAML data or workflow file, 128 lines, 1 explicit URLs. Top keys: repository\_web\_base, items. & \textbf{\path{todo-backlog.yaml}}. Тип: yaml. Размер: 6116 байт. Форма: 128 строк. Файл YAML с данными или workflow, 128 строк, явных URL: 1. Верхние ключи: repository\_web\_base, items. \\
\textbf{\path{utilities/README.md}}. Kind: markdown. Size: 3495 bytes. Shape: 54 lines. Markdown document, 54 lines, 0 explicit URLs. Lead heading: Repository Utilities. & \textbf{\path{utilities/README.md}}. Тип: markdown. Размер: 3495 байт. Форма: 54 строк. Документ Markdown, 54 строк, явных URL: 0. Главный заголовок: Repository Utilities. \\
\textbf{\path{uv.lock}}. Kind: generic. Size: 784153 bytes. Shape: binary surface. Tracked repository file, 784153 bytes of binary data, 0 explicit URLs. Opening line: No non-empty preview line. & \textbf{\path{uv.lock}}. Тип: generic. Размер: 784153 байт. Форма: бинарная поверхность. Отслеживаемый файл репозитория, 784153 байт бинарных данных, явных URL: 0. Первая строка: No non-empty preview line. \\
\bottomrule
\end{longtable}
\endgroup
\subsection*{Summaries Of All Authored URLs / Сводки по всем авторским URL}
\begingroup
\small
\renewcommand{\arraystretch}{1.16}
\begin{longtable}{>{\raggedright\arraybackslash}p{0.47\linewidth} >{\raggedright\arraybackslash}p{0.47\linewidth}}
\toprule
\textbf{English} & \textbf{Русский} \\
\midrule
\endfirsthead
\toprule
\textbf{English} & \textbf{Русский} \\
\midrule
\endhead
\textbf{\url{http://127.0.0.1}}. Host: 127.0.0.1. Occurrences: 2. Referenced by URL only. No repository-local mirror is currently tracked. Sources: \path{tests/test_codex_proxy.py}. & \textbf{\url{http://127.0.0.1}}. Хост: 127.0.0.1. Упоминаний: 2. Ссылка хранится только как URL. Локальное зеркало в репозитории сейчас не отслеживается. Источники: \path{tests/test_codex_proxy.py}. \\
\textbf{\url{http://127.0.0.1:40111/openai`}}. Host: 127.0.0.1:40111. Occurrences: 1. Referenced by URL only. No repository-local mirror is currently tracked. Sources: \path{docs/audit/2026-05-02-zz-codex-exec-resume-session-state-stage0.md}. & \textbf{\url{http://127.0.0.1:40111/openai`}}. Хост: 127.0.0.1:40111. Упоминаний: 1. Ссылка хранится только как URL. Локальное зеркало в репозитории сейчас не отслеживается. Источники: \path{docs/audit/2026-05-02-zz-codex-exec-resume-session-state-stage0.md}. \\
\textbf{\url{http://127.0.0.1:44973/openai`}}. Host: 127.0.0.1:44973. Occurrences: 1. Referenced by URL only. No repository-local mirror is currently tracked. Sources: \path{docs/audit/2026-05-02-zz-codex-exec-resume-session-state-stage0.md}. & \textbf{\url{http://127.0.0.1:44973/openai`}}. Хост: 127.0.0.1:44973. Упоминаний: 1. Ссылка хранится только как URL. Локальное зеркало в репозитории сейчас не отслеживается. Источники: \path{docs/audit/2026-05-02-zz-codex-exec-resume-session-state-stage0.md}. \\
\textbf{\url{https://api.github.com/repos/example/demo/branches/master/protection}}. Host: api.github.com. Occurrences: 1. External GitHub URL referenced from 1 file(s). Local companion path(s): none. Sources: \path{tests/test_github_pr_gates.py}. & \textbf{\url{https://api.github.com/repos/example/demo/branches/master/protection}}. Хост: api.github.com. Упоминаний: 1. Внешний URL GitHub, упомянутый в 1 файле(ах). Локальные сопроводительные пути: none. Источники: \path{tests/test_github_pr_gates.py}. \\
\textbf{\url{https://astral.sh/}}. Host: astral.sh. Occurrences: 2. Referenced by URL only. No repository-local mirror is currently tracked. Sources: \path{tests/test_exploratorium_translation.py}\newline \path{tests/test_utilities.py}. & \textbf{\url{https://astral.sh/}}. Хост: astral.sh. Упоминаний: 2. Ссылка хранится только как URL. Локальное зеркало в репозитории сейчас не отслеживается. Источники: \path{tests/test_exploratorium_translation.py}\newline \path{tests/test_utilities.py}. \\
\textbf{\url{https://astral.sh/uv/install.sh}}. Host: astral.sh. Occurrences: 1. Referenced by URL only. No repository-local mirror is currently tracked. Sources: \path{README.md}. & \textbf{\url{https://astral.sh/uv/install.sh}}. Хост: astral.sh. Упоминаний: 1. Ссылка хранится только как URL. Локальное зеркало в репозитории сейчас не отслеживается. Источники: \path{README.md}. \\
\textbf{\url{https://dspy-helper.openai.azure.com}}. Host: dspy-helper.openai.azure.com. Occurrences: 1. Referenced by URL only. No repository-local mirror is currently tracked. Sources: \path{tests/test_dspy_training.py}. & \textbf{\url{https://dspy-helper.openai.azure.com}}. Хост: dspy-helper.openai.azure.com. Упоминаний: 1. Ссылка хранится только как URL. Локальное зеркало в репозитории сейчас не отслеживается. Источники: \path{tests/test_dspy_training.py}. \\
\textbf{\url{https://dspy-helper.openai.azure.com/}}. Host: dspy-helper.openai.azure.com. Occurrences: 1. Referenced by URL only. No repository-local mirror is currently tracked. Sources: \path{tests/test_dspy_training.py}. & \textbf{\url{https://dspy-helper.openai.azure.com/}}. Хост: dspy-helper.openai.azure.com. Упоминаний: 1. Ссылка хранится только как URL. Локальное зеркало в репозитории сейчас не отслеживается. Источники: \path{tests/test_dspy_training.py}. \\
\textbf{\url{https://dspy.ai/tutorials/rag/}}. Host: dspy.ai. Occurrences: 1. Referenced by URL only. No repository-local mirror is currently tracked. Sources: \path{docs/architecture/inspired/dspy-rag-tutorial.md}. & \textbf{\url{https://dspy.ai/tutorials/rag/}}. Хост: dspy.ai. Упоминаний: 1. Ссылка хранится только как URL. Локальное зеркало в репозитории сейчас не отслеживается. Источники: \path{docs/architecture/inspired/dspy-rag-tutorial.md}. \\
\textbf{\url{https://example.com}}. Host: example.com. Occurrences: 2. Referenced by URL only. No repository-local mirror is currently tracked. Sources: \path{tests/test_pages_site.py}. & \textbf{\url{https://example.com}}. Хост: example.com. Упоминаний: 2. Ссылка хранится только как URL. Локальное зеркало в репозитории сейчас не отслеживается. Источники: \path{tests/test_pages_site.py}. \\
\textbf{\url{https://example.com/docs}}. Host: example.com. Occurrences: 1. Referenced by URL only. No repository-local mirror is currently tracked. Sources: \path{tests/test_pages_site.py}. & \textbf{\url{https://example.com/docs}}. Хост: example.com. Упоминаний: 1. Ссылка хранится только как URL. Локальное зеркало в репозитории сейчас не отслеживается. Источники: \path{tests/test_pages_site.py}. \\
\textbf{\url{https://example.com/org/repo/blob/main}}. Host: example.com. Occurrences: 1. Referenced by URL only. No repository-local mirror is currently tracked. Sources: \path{tests/test_todo_backlog.py}. & \textbf{\url{https://example.com/org/repo/blob/main}}. Хост: example.com. Упоминаний: 1. Ссылка хранится только как URL. Локальное зеркало в репозитории сейчас не отслеживается. Источники: \path{tests/test_todo_backlog.py}. \\
\textbf{\url{https://example.invalid/v1}}. Host: example.invalid. Occurrences: 2. Referenced by URL only. No repository-local mirror is currently tracked. Sources: \path{tests/test_dspy_training.py}. & \textbf{\url{https://example.invalid/v1}}. Хост: example.invalid. Упоминаний: 2. Ссылка хранится только как URL. Локальное зеркало в репозитории сейчас не отслеживается. Источники: \path{tests/test_dspy_training.py}. \\
\textbf{\url{https://example.openai.azure.com}}. Host: example.openai.azure.com. Occurrences: 5. Referenced by URL only. No repository-local mirror is currently tracked. Sources: \path{tests/test_azure_runtime.py}\newline \path{tests/test_dspy_training.py}. & \textbf{\url{https://example.openai.azure.com}}. Хост: example.openai.azure.com. Упоминаний: 5. Ссылка хранится только как URL. Локальное зеркало в репозитории сейчас не отслеживается. Источники: \path{tests/test_azure_runtime.py}\newline \path{tests/test_dspy_training.py}. \\
\textbf{\url{https://example.openai.azure.com/}}. Host: example.openai.azure.com. Occurrences: 6. Referenced by URL only. No repository-local mirror is currently tracked. Sources: \path{tests/test_azure_runtime.py}\newline \path{tests/test_dspy_training.py}. & \textbf{\url{https://example.openai.azure.com/}}. Хост: example.openai.azure.com. Упоминаний: 6. Ссылка хранится только как URL. Локальное зеркало в репозитории сейчас не отслеживается. Источники: \path{tests/test_azure_runtime.py}\newline \path{tests/test_dspy_training.py}. \\
\textbf{\url{https://example.openai.azure.com/openai/deployments/repo-rag-ft}}. Host: example.openai.azure.com. Occurrences: 3. Referenced by URL only. No repository-local mirror is currently tracked. Sources: \path{tests/test_azure_runtime.py}. & \textbf{\url{https://example.openai.azure.com/openai/deployments/repo-rag-ft}}. Хост: example.openai.azure.com. Упоминаний: 3. Ссылка хранится только как URL. Локальное зеркало в репозитории сейчас не отслеживается. Источники: \path{tests/test_azure_runtime.py}. \\
\textbf{\url{https://example.openai.azure.com/openai/deployments/repo-rag-ft/}}. Host: example.openai.azure.com. Occurrences: 4. Referenced by URL only. No repository-local mirror is currently tracked. Sources: \path{tests/test_azure_runtime.py}. & \textbf{\url{https://example.openai.azure.com/openai/deployments/repo-rag-ft/}}. Хост: example.openai.azure.com. Упоминаний: 4. Ссылка хранится только как URL. Локальное зеркало в репозитории сейчас не отслеживается. Источники: \path{tests/test_azure_runtime.py}. \\
\textbf{\url{https://example.openai.azure.com/openai/deployments/repo-rag-ft/chat/completions}}. Host: example.openai.azure.com. Occurrences: 1. Referenced by URL only. No repository-local mirror is currently tracked. Sources: \path{tests/test_dspy_training.py}. & \textbf{\url{https://example.openai.azure.com/openai/deployments/repo-rag-ft/chat/completions}}. Хост: example.openai.azure.com. Упоминаний: 1. Ссылка хранится только как URL. Локальное зеркало в репозитории сейчас не отслеживается. Источники: \path{tests/test_dspy_training.py}. \\
\textbf{\url{https://example.services.ai.azure.com/models}}. Host: example.services.ai.azure.com. Occurrences: 6. Referenced by URL only. No repository-local mirror is currently tracked. Sources: \path{Makefile}\newline \path{docs/architecture/dspy-rag-guide.md}\newline \path{docs/operations/azure-deployment.md}\newline \path{src/repo_rag_lab/utilities.py}\newline \path{tests/test_cli_and_dspy.py}\newline \path{tests/test_project_surfaces.py}. & \textbf{\url{https://example.services.ai.azure.com/models}}. Хост: example.services.ai.azure.com. Упоминаний: 6. Ссылка хранится только как URL. Локальное зеркало в репозитории сейчас не отслеживается. Источники: \path{Makefile}\newline \path{docs/architecture/dspy-rag-guide.md}\newline \path{docs/operations/azure-deployment.md}\newline \path{src/repo_rag_lab/utilities.py}\newline \path{tests/test_cli_and_dspy.py}\newline \path{tests/test_project_surfaces.py}. \\
\textbf{\url{https://from-connection-string}}. Host: from-connection-string. Occurrences: 2. Referenced by URL only. No repository-local mirror is currently tracked. Sources: \path{tests/test_runtime_artifacts_azure.py}. & \textbf{\url{https://from-connection-string}}. Хост: from-connection-string. Упоминаний: 2. Ссылка хранится только как URL. Локальное зеркало в репозитории сейчас не отслеживается. Источники: \path{tests/test_runtime_artifacts_azure.py}. \\
\textbf{\url{https://github.com/}}. Host: github.com. Occurrences: 2. External GitHub URL referenced from 1 file(s). Local companion path(s): none. Sources: \path{src/repo_rag_lab/pages_site.py}. & \textbf{\url{https://github.com/}}. Хост: github.com. Упоминаний: 2. Внешний URL GitHub, упомянутый в 1 файле(ах). Локальные сопроводительные пути: none. Источники: \path{src/repo_rag_lab/pages_site.py}. \\
\textbf{\url{https://github.com/acme/repo}}. Host: github.com. Occurrences: 3. External GitHub URL referenced from 1 file(s). Local companion path(s): none. Sources: \path{tests/test_codex_proxy.py}. & \textbf{\url{https://github.com/acme/repo}}. Хост: github.com. Упоминаний: 3. Внешний URL GitHub, упомянутый в 1 файле(ах). Локальные сопроводительные пути: none. Источники: \path{tests/test_codex_proxy.py}. \\
\textbf{\url{https://github.com/example/demo}}. Host: github.com. Occurrences: 10. External GitHub URL referenced from 3 file(s). Local companion path(s): none. Sources: \path{tests/test_cli_and_dspy.py}\newline \path{tests/test_pages_site.py}\newline \path{tests/test_utilities.py}. & \textbf{\url{https://github.com/example/demo}}. Хост: github.com. Упоминаний: 10. Внешний URL GitHub, упомянутый в 3 файле(ах). Локальные сопроводительные пути: none. Источники: \path{tests/test_cli_and_dspy.py}\newline \path{tests/test_pages_site.py}\newline \path{tests/test_utilities.py}. \\
\textbf{\url{https://github.com/example/demo/blob/main/README.md}}. Host: github.com. Occurrences: 1. External GitHub URL referenced from 1 file(s). Local companion path(s): none. Sources: \path{tests/test_pages_site.py}. & \textbf{\url{https://github.com/example/demo/blob/main/README.md}}. Хост: github.com. Упоминаний: 1. Внешний URL GitHub, упомянутый в 1 файле(ах). Локальные сопроводительные пути: none. Источники: \path{tests/test_pages_site.py}. \\
\textbf{\url{https://github.com/example/demo/blob/main/assets/logo.png}}. Host: github.com. Occurrences: 1. External GitHub URL referenced from 1 file(s). Local companion path(s): none. Sources: \path{tests/test_pages_site.py}. & \textbf{\url{https://github.com/example/demo/blob/main/assets/logo.png}}. Хост: github.com. Упоминаний: 1. Внешний URL GitHub, упомянутый в 1 файле(ах). Локальные сопроводительные пути: none. Источники: \path{tests/test_pages_site.py}. \\
\textbf{\url{https://github.com/example/demo/blob/main/src/tool.c}}. Host: github.com. Occurrences: 1. External GitHub URL referenced from 1 file(s). Local companion path(s): none. Sources: \path{tests/test_pages_site.py}. & \textbf{\url{https://github.com/example/demo/blob/main/src/tool.c}}. Хост: github.com. Упоминаний: 1. Внешний URL GitHub, упомянутый в 1 файле(ах). Локальные сопроводительные пути: none. Источники: \path{tests/test_pages_site.py}. \\
\textbf{\url{https://github.com/example/demo/tree/main/docs}}. Host: github.com. Occurrences: 1. External GitHub URL referenced from 1 file(s). Local companion path(s): none. Sources: \path{tests/test_pages_site.py}. & \textbf{\url{https://github.com/example/demo/tree/main/docs}}. Хост: github.com. Упоминаний: 1. Внешний URL GitHub, упомянутый в 1 файле(ах). Локальные сопроводительные пути: none. Источники: \path{tests/test_pages_site.py}. \\
\textbf{\url{https://github.com/example/project}}. Host: github.com. Occurrences: 2. External GitHub URL referenced from 2 file(s). Local companion path(s): none. Sources: \path{tests/test_exploratorium_translation.py}\newline \path{tests/test_utilities.py}. & \textbf{\url{https://github.com/example/project}}. Хост: github.com. Упоминаний: 2. Внешний URL GitHub, упомянутый в 2 файле(ах). Локальные сопроводительные пути: none. Источники: \path{tests/test_exploratorium_translation.py}\newline \path{tests/test_utilities.py}. \\
\textbf{\url{https://github.com/example/project/blob/main/README.md}}. Host: github.com. Occurrences: 1. External GitHub URL referenced from 1 file(s). Local companion path(s): none. Sources: \path{tests/test_exploratorium_translation.py}. & \textbf{\url{https://github.com/example/project/blob/main/README.md}}. Хост: github.com. Упоминаний: 1. Внешний URL GitHub, упомянутый в 1 файле(ах). Локальные сопроводительные пути: none. Источники: \path{tests/test_exploratorium_translation.py}. \\
\textbf{\url{https://github.com/example/project/blob/main/README.md\#L1-L2}}. Host: github.com. Occurrences: 1. External GitHub URL referenced from 1 file(s). Local companion path(s): none. Sources: \path{tests/test_exploratorium_translation.py}. & \textbf{\url{https://github.com/example/project/blob/main/README.md\#L1-L2}}. Хост: github.com. Упоминаний: 1. Внешний URL GitHub, упомянутый в 1 файле(ах). Локальные сопроводительные пути: none. Источники: \path{tests/test_exploratorium_translation.py}. \\
\textbf{\url{https://github.com/realagiorganization/dspy_rag_in_repo_docs_and_impl1}}. Host: github.com. Occurrences: 1. External GitHub URL referenced from 1 file(s). Local companion path(s): none. Sources: \path{mkdocs.yml}. & \textbf{\url{https://github.com/realagiorganization/dspy_rag_in_repo_docs_and_impl1}}. Хост: github.com. Упоминаний: 1. Внешний URL GitHub, упомянутый в 1 файле(ах). Локальные сопроводительные пути: none. Источники: \path{mkdocs.yml}. \\
\textbf{\url{https://github.com/realagiorganization/dspy_rag_in_repo_docs_and_impl1/actions/workflows/ci.yml/badge.svg)](https://github.com/realagiorganization/dspy_rag_in_repo_docs_and_impl1/actions/workflows/ci.yml}}. Host: github.com. Occurrences: 2. External GitHub URL referenced from 2 file(s). Local companion path(s): none. Sources: \path{README.md}\newline \path{notebooks/02_agent_workflow_checklist.ipynb}. & \textbf{\url{https://github.com/realagiorganization/dspy_rag_in_repo_docs_and_impl1/actions/workflows/ci.yml/badge.svg)](https://github.com/realagiorganization/dspy_rag_in_repo_docs_and_impl1/actions/workflows/ci.yml}}. Хост: github.com. Упоминаний: 2. Внешний URL GitHub, упомянутый в 2 файле(ах). Локальные сопроводительные пути: none. Источники: \path{README.md}\newline \path{notebooks/02_agent_workflow_checklist.ipynb}. \\
\textbf{\url{https://github.com/realagiorganization/dspy_rag_in_repo_docs_and_impl1/actions/workflows/publish.yml/badge.svg)](https://github.com/realagiorganization/dspy_rag_in_repo_docs_and_impl1/actions/workflows/publish.yml}}. Host: github.com. Occurrences: 1. External GitHub URL referenced from 1 file(s). Local companion path(s): none. Sources: \path{README.md}. & \textbf{\url{https://github.com/realagiorganization/dspy_rag_in_repo_docs_and_impl1/actions/workflows/publish.yml/badge.svg)](https://github.com/realagiorganization/dspy_rag_in_repo_docs_and_impl1/actions/workflows/publish.yml}}. Хост: github.com. Упоминаний: 1. Внешний URL GitHub, упомянутый в 1 файле(ах). Локальные сопроводительные пути: none. Источники: \path{README.md}. \\
\textbf{\url{https://github.com/realagiorganization/dspy_rag_in_repo_docs_and_impl1/blob/master}}. Host: github.com. Occurrences: 3. External GitHub URL referenced from 3 file(s). Local companion path(s): none. Sources: \path{FILES.md}\newline \path{src/repo_rag_lab/todo_backlog.py}\newline \path{todo-backlog.yaml}. & \textbf{\url{https://github.com/realagiorganization/dspy_rag_in_repo_docs_and_impl1/blob/master}}. Хост: github.com. Упоминаний: 3. Внешний URL GitHub, упомянутый в 3 файле(ах). Локальные сопроводительные пути: none. Источники: \path{FILES.md}\newline \path{src/repo_rag_lab/todo_backlog.py}\newline \path{todo-backlog.yaml}. \\
\textbf{\url{https://github.com/realagiorganization/dspy_rag_in_repo_docs_and_impl1/blob/master/.github/workflows/ci.yml}}. Host: github.com. Occurrences: 1. External GitHub URL referenced from 1 file(s). Local companion path(s): none. Sources: \path{publication/todo-backlog-table.tex}. & \textbf{\url{https://github.com/realagiorganization/dspy_rag_in_repo_docs_and_impl1/blob/master/.github/workflows/ci.yml}}. Хост: github.com. Упоминаний: 1. Внешний URL GitHub, упомянутый в 1 файле(ах). Локальные сопроводительные пути: none. Источники: \path{publication/todo-backlog-table.tex}. \\
\textbf{\url{https://github.com/realagiorganization/dspy_rag_in_repo_docs_and_impl1/blob/master/.github/workflows/publish.yml}}. Host: github.com. Occurrences: 1. External GitHub URL referenced from 1 file(s). Local companion path(s): none. Sources: \path{publication/todo-backlog-table.tex}. & \textbf{\url{https://github.com/realagiorganization/dspy_rag_in_repo_docs_and_impl1/blob/master/.github/workflows/publish.yml}}. Хост: github.com. Упоминаний: 1. Внешний URL GitHub, упомянутый в 1 файле(ах). Локальные сопроводительные пути: none. Источники: \path{publication/todo-backlog-table.tex}. \\
\textbf{\url{https://github.com/realagiorganization/dspy_rag_in_repo_docs_and_impl1/blob/master/Makefile\#L142-L143}}. Host: github.com. Occurrences: 1. External GitHub URL referenced from 1 file(s). Local companion path(s): none. Sources: \path{docs/operations/simple-end-to-end-verification-guide.md}. & \textbf{\url{https://github.com/realagiorganization/dspy_rag_in_repo_docs_and_impl1/blob/master/Makefile\#L142-L143}}. Хост: github.com. Упоминаний: 1. Внешний URL GitHub, упомянутый в 1 файле(ах). Локальные сопроводительные пути: none. Источники: \path{docs/operations/simple-end-to-end-verification-guide.md}. \\
\textbf{\url{https://github.com/realagiorganization/dspy_rag_in_repo_docs_and_impl1/blob/master/Makefile\#L153-L154}}. Host: github.com. Occurrences: 1. External GitHub URL referenced from 1 file(s). Local companion path(s): none. Sources: \path{docs/operations/simple-end-to-end-verification-guide.md}. & \textbf{\url{https://github.com/realagiorganization/dspy_rag_in_repo_docs_and_impl1/blob/master/Makefile\#L153-L154}}. Хост: github.com. Упоминаний: 1. Внешний URL GitHub, упомянутый в 1 файле(ах). Локальные сопроводительные пути: none. Источники: \path{docs/operations/simple-end-to-end-verification-guide.md}. \\
\textbf{\url{https://github.com/realagiorganization/dspy_rag_in_repo_docs_and_impl1/blob/master/Makefile\#L248-L251}}. Host: github.com. Occurrences: 1. External GitHub URL referenced from 1 file(s). Local companion path(s): none. Sources: \path{docs/operations/simple-end-to-end-verification-guide.md}. & \textbf{\url{https://github.com/realagiorganization/dspy_rag_in_repo_docs_and_impl1/blob/master/Makefile\#L248-L251}}. Хост: github.com. Упоминаний: 1. Внешний URL GitHub, упомянутый в 1 файле(ах). Локальные сопроводительные пути: none. Источники: \path{docs/operations/simple-end-to-end-verification-guide.md}. \\
\textbf{\url{https://github.com/realagiorganization/dspy_rag_in_repo_docs_and_impl1/blob/master/Makefile\#L65-L66}}. Host: github.com. Occurrences: 1. External GitHub URL referenced from 1 file(s). Local companion path(s): none. Sources: \path{docs/operations/simple-end-to-end-verification-guide.md}. & \textbf{\url{https://github.com/realagiorganization/dspy_rag_in_repo_docs_and_impl1/blob/master/Makefile\#L65-L66}}. Хост: github.com. Упоминаний: 1. Внешний URL GitHub, упомянутый в 1 файле(ах). Локальные сопроводительные пути: none. Источники: \path{docs/operations/simple-end-to-end-verification-guide.md}. \\
\textbf{\url{https://github.com/realagiorganization/dspy_rag_in_repo_docs_and_impl1/blob/master/README.md}}. Host: github.com. Occurrences: 2. External GitHub URL referenced from 1 file(s). Local companion path(s): none. Sources: \path{publication/todo-backlog-table.tex}. & \textbf{\url{https://github.com/realagiorganization/dspy_rag_in_repo_docs_and_impl1/blob/master/README.md}}. Хост: github.com. Упоминаний: 2. Внешний URL GitHub, упомянутый в 1 файле(ах). Локальные сопроводительные пути: none. Источники: \path{publication/todo-backlog-table.tex}. \\
\textbf{\url{https://github.com/realagiorganization/dspy_rag_in_repo_docs_and_impl1/blob/master/TODO.MD}}. Host: github.com. Occurrences: 1. External GitHub URL referenced from 1 file(s). Local companion path(s): none. Sources: \path{publication/todo-backlog-table.tex}. & \textbf{\url{https://github.com/realagiorganization/dspy_rag_in_repo_docs_and_impl1/blob/master/TODO.MD}}. Хост: github.com. Упоминаний: 1. Внешний URL GitHub, упомянутый в 1 файле(ах). Локальные сопроводительные пути: none. Источники: \path{publication/todo-backlog-table.tex}. \\
\textbf{\url{https://github.com/realagiorganization/dspy_rag_in_repo_docs_and_impl1/blob/master/docs/architecture/dspy-rag-guide.md}}. Host: github.com. Occurrences: 1. External GitHub URL referenced from 1 file(s). Local companion path(s): none. Sources: \path{publication/todo-backlog-table.tex}. & \textbf{\url{https://github.com/realagiorganization/dspy_rag_in_repo_docs_and_impl1/blob/master/docs/architecture/dspy-rag-guide.md}}. Хост: github.com. Упоминаний: 1. Внешний URL GitHub, упомянутый в 1 файле(ах). Локальные сопроводительные пути: none. Источники: \path{publication/todo-backlog-table.tex}. \\
\textbf{\url{https://github.com/realagiorganization/dspy_rag_in_repo_docs_and_impl1/blob/master/docs/archive/repo-completeness-checklist.md}}. Host: github.com. Occurrences: 1. External GitHub URL referenced from 1 file(s). Local companion path(s): none. Sources: \path{publication/todo-backlog-table.tex}. & \textbf{\url{https://github.com/realagiorganization/dspy_rag_in_repo_docs_and_impl1/blob/master/docs/archive/repo-completeness-checklist.md}}. Хост: github.com. Упоминаний: 1. Внешний URL GitHub, упомянутый в 1 файле(ах). Локальные сопроводительные пути: none. Источники: \path{publication/todo-backlog-table.tex}. \\
\textbf{\url{https://github.com/realagiorganization/dspy_rag_in_repo_docs_and_impl1/blob/master/docs/audit/README.md}}. Host: github.com. Occurrences: 3. External GitHub URL referenced from 1 file(s). Local companion path(s): none. Sources: \path{publication/todo-backlog-table.tex}. & \textbf{\url{https://github.com/realagiorganization/dspy_rag_in_repo_docs_and_impl1/blob/master/docs/audit/README.md}}. Хост: github.com. Упоминаний: 3. Внешний URL GitHub, упомянутый в 1 файле(ах). Локальные сопроводительные пути: none. Источники: \path{publication/todo-backlog-table.tex}. \\
\textbf{\url{https://github.com/realagiorganization/dspy_rag_in_repo_docs_and_impl1/blob/master/docs/operations/azure-deployment.md}}. Host: github.com. Occurrences: 1. External GitHub URL referenced from 1 file(s). Local companion path(s): none. Sources: \path{publication/todo-backlog-table.tex}. & \textbf{\url{https://github.com/realagiorganization/dspy_rag_in_repo_docs_and_impl1/blob/master/docs/operations/azure-deployment.md}}. Хост: github.com. Упоминаний: 1. Внешний URL GitHub, упомянутый в 1 файле(ах). Локальные сопроводительные пути: none. Источники: \path{publication/todo-backlog-table.tex}. \\
\textbf{\url{https://github.com/realagiorganization/dspy_rag_in_repo_docs_and_impl1/blob/master/docs/planning/dataset-integration-plan.md}}. Host: github.com. Occurrences: 1. External GitHub URL referenced from 1 file(s). Local companion path(s): none. Sources: \path{publication/todo-backlog-table.tex}. & \textbf{\url{https://github.com/realagiorganization/dspy_rag_in_repo_docs_and_impl1/blob/master/docs/planning/dataset-integration-plan.md}}. Хост: github.com. Упоминаний: 1. Внешний URL GitHub, упомянутый в 1 файле(ах). Локальные сопроводительные пути: none. Источники: \path{publication/todo-backlog-table.tex}. \\
\textbf{\url{https://github.com/realagiorganization/dspy_rag_in_repo_docs_and_impl1/blob/master/docs/planning/repo-hardening-plan.md}}. Host: github.com. Occurrences: 1. External GitHub URL referenced from 1 file(s). Local companion path(s): none. Sources: \path{publication/todo-backlog-table.tex}. & \textbf{\url{https://github.com/realagiorganization/dspy_rag_in_repo_docs_and_impl1/blob/master/docs/planning/repo-hardening-plan.md}}. Хост: github.com. Упоминаний: 1. Внешний URL GitHub, упомянутый в 1 файле(ах). Локальные сопроводительные пути: none. Источники: \path{publication/todo-backlog-table.tex}. \\
\textbf{\url{https://github.com/realagiorganization/dspy_rag_in_repo_docs_and_impl1/blob/master/notebooks/01_repo_rag_research.ipynb}}. Host: github.com. Occurrences: 1. External GitHub URL referenced from 1 file(s). Local companion path(s): none. Sources: \path{publication/todo-backlog-table.tex}. & \textbf{\url{https://github.com/realagiorganization/dspy_rag_in_repo_docs_and_impl1/blob/master/notebooks/01_repo_rag_research.ipynb}}. Хост: github.com. Упоминаний: 1. Внешний URL GitHub, упомянутый в 1 файле(ах). Локальные сопроводительные пути: none. Источники: \path{publication/todo-backlog-table.tex}. \\
\textbf{\url{https://github.com/realagiorganization/dspy_rag_in_repo_docs_and_impl1/blob/master/notebooks/03_dspy_training_lab.ipynb}}. Host: github.com. Occurrences: 1. External GitHub URL referenced from 1 file(s). Local companion path(s): none. Sources: \path{publication/todo-backlog-table.tex}. & \textbf{\url{https://github.com/realagiorganization/dspy_rag_in_repo_docs_and_impl1/blob/master/notebooks/03_dspy_training_lab.ipynb}}. Хост: github.com. Упоминаний: 1. Внешний URL GitHub, упомянутый в 1 файле(ах). Локальные сопроводительные пути: none. Источники: \path{publication/todo-backlog-table.tex}. \\
\textbf{\url{https://github.com/realagiorganization/dspy_rag_in_repo_docs_and_impl1/blob/master/notebooks/04_sample_population_lab.ipynb}}. Host: github.com. Occurrences: 1. External GitHub URL referenced from 1 file(s). Local companion path(s): none. Sources: \path{publication/todo-backlog-table.tex}. & \textbf{\url{https://github.com/realagiorganization/dspy_rag_in_repo_docs_and_impl1/blob/master/notebooks/04_sample_population_lab.ipynb}}. Хост: github.com. Упоминаний: 1. Внешний URL GitHub, упомянутый в 1 файле(ах). Локальные сопроводительные пути: none. Источники: \path{publication/todo-backlog-table.tex}. \\
\textbf{\url{https://github.com/realagiorganization/dspy_rag_in_repo_docs_and_impl1/blob/master/publication/README.md}}. Host: github.com. Occurrences: 1. External GitHub URL referenced from 1 file(s). Local companion path(s): none. Sources: \path{publication/todo-backlog-table.tex}. & \textbf{\url{https://github.com/realagiorganization/dspy_rag_in_repo_docs_and_impl1/blob/master/publication/README.md}}. Хост: github.com. Упоминаний: 1. Внешний URL GitHub, упомянутый в 1 файле(ах). Локальные сопроводительные пути: none. Источники: \path{publication/todo-backlog-table.tex}. \\
\textbf{\url{https://github.com/realagiorganization/dspy_rag_in_repo_docs_and_impl1/blob/master/pyproject.toml}}. Host: github.com. Occurrences: 1. External GitHub URL referenced from 1 file(s). Local companion path(s): none. Sources: \path{publication/todo-backlog-table.tex}. & \textbf{\url{https://github.com/realagiorganization/dspy_rag_in_repo_docs_and_impl1/blob/master/pyproject.toml}}. Хост: github.com. Упоминаний: 1. Внешний URL GitHub, упомянутый в 1 файле(ах). Локальные сопроводительные пути: none. Источники: \path{publication/todo-backlog-table.tex}. \\
\textbf{\url{https://github.com/realagiorganization/dspy_rag_in_repo_docs_and_impl1/blob/master/rust-cli/Cargo.lock}}. Host: github.com. Occurrences: 1. External GitHub URL referenced from 1 file(s). Local companion path(s): none. Sources: \path{publication/todo-backlog-table.tex}. & \textbf{\url{https://github.com/realagiorganization/dspy_rag_in_repo_docs_and_impl1/blob/master/rust-cli/Cargo.lock}}. Хост: github.com. Упоминаний: 1. Внешний URL GitHub, упомянутый в 1 файле(ах). Локальные сопроводительные пути: none. Источники: \path{publication/todo-backlog-table.tex}. \\
\textbf{\url{https://github.com/realagiorganization/dspy_rag_in_repo_docs_and_impl1/blob/master/rust-cli/Cargo.toml}}. Host: github.com. Occurrences: 1. External GitHub URL referenced from 1 file(s). Local companion path(s): none. Sources: \path{publication/todo-backlog-table.tex}. & \textbf{\url{https://github.com/realagiorganization/dspy_rag_in_repo_docs_and_impl1/blob/master/rust-cli/Cargo.toml}}. Хост: github.com. Упоминаний: 1. Внешний URL GitHub, упомянутый в 1 файле(ах). Локальные сопроводительные пути: none. Источники: \path{publication/todo-backlog-table.tex}. \\
\textbf{\url{https://github.com/realagiorganization/dspy_rag_in_repo_docs_and_impl1/blob/master/samples/population/repository_population_candidates.yaml}}. Host: github.com. Occurrences: 1. External GitHub URL referenced from 1 file(s). Local companion path(s): none. Sources: \path{publication/todo-backlog-table.tex}. & \textbf{\url{https://github.com/realagiorganization/dspy_rag_in_repo_docs_and_impl1/blob/master/samples/population/repository_population_candidates.yaml}}. Хост: github.com. Упоминаний: 1. Внешний URL GitHub, упомянутый в 1 файле(ах). Локальные сопроводительные пути: none. Источники: \path{publication/todo-backlog-table.tex}. \\
\textbf{\url{https://github.com/realagiorganization/dspy_rag_in_repo_docs_and_impl1/blob/master/samples/training/repository_training_examples.yaml}}. Host: github.com. Occurrences: 1. External GitHub URL referenced from 1 file(s). Local companion path(s): none. Sources: \path{publication/todo-backlog-table.tex}. & \textbf{\url{https://github.com/realagiorganization/dspy_rag_in_repo_docs_and_impl1/blob/master/samples/training/repository_training_examples.yaml}}. Хост: github.com. Упоминаний: 1. Внешний URL GitHub, упомянутый в 1 файле(ах). Локальные сопроводительные пути: none. Источники: \path{publication/todo-backlog-table.tex}. \\
\textbf{\url{https://github.com/realagiorganization/dspy_rag_in_repo_docs_and_impl1/blob/master/src/repo_rag_lab/benchmarks.py}}. Host: github.com. Occurrences: 2. External GitHub URL referenced from 1 file(s). Local companion path(s): none. Sources: \path{publication/todo-backlog-table.tex}. & \textbf{\url{https://github.com/realagiorganization/dspy_rag_in_repo_docs_and_impl1/blob/master/src/repo_rag_lab/benchmarks.py}}. Хост: github.com. Упоминаний: 2. Внешний URL GitHub, упомянутый в 1 файле(ах). Локальные сопроводительные пути: none. Источники: \path{publication/todo-backlog-table.tex}. \\
\textbf{\url{https://github.com/realagiorganization/dspy_rag_in_repo_docs_and_impl1/blob/master/src/repo_rag_lab/cli.py\#L72-L88}}. Host: github.com. Occurrences: 1. External GitHub URL referenced from 1 file(s). Local companion path(s): none. Sources: \path{docs/operations/simple-end-to-end-verification-guide.md}. & \textbf{\url{https://github.com/realagiorganization/dspy_rag_in_repo_docs_and_impl1/blob/master/src/repo_rag_lab/cli.py\#L72-L88}}. Хост: github.com. Упоминаний: 1. Внешний URL GitHub, упомянутый в 1 файле(ах). Локальные сопроводительные пути: none. Источники: \path{docs/operations/simple-end-to-end-verification-guide.md}. \\
\textbf{\url{https://github.com/realagiorganization/dspy_rag_in_repo_docs_and_impl1/blob/master/src/repo_rag_lab/dspy_workflow.py}}. Host: github.com. Occurrences: 1. External GitHub URL referenced from 1 file(s). Local companion path(s): none. Sources: \path{publication/todo-backlog-table.tex}. & \textbf{\url{https://github.com/realagiorganization/dspy_rag_in_repo_docs_and_impl1/blob/master/src/repo_rag_lab/dspy_workflow.py}}. Хост: github.com. Упоминаний: 1. Внешний URL GitHub, упомянутый в 1 файле(ах). Локальные сопроводительные пути: none. Источники: \path{publication/todo-backlog-table.tex}. \\
\textbf{\url{https://github.com/realagiorganization/dspy_rag_in_repo_docs_and_impl1/blob/master/src/repo_rag_lab/notebook_scaffolding.py}}. Host: github.com. Occurrences: 1. External GitHub URL referenced from 1 file(s). Local companion path(s): none. Sources: \path{publication/todo-backlog-table.tex}. & \textbf{\url{https://github.com/realagiorganization/dspy_rag_in_repo_docs_and_impl1/blob/master/src/repo_rag_lab/notebook_scaffolding.py}}. Хост: github.com. Упоминаний: 1. Внешний URL GitHub, упомянутый в 1 файле(ах). Локальные сопроводительные пути: none. Источники: \path{publication/todo-backlog-table.tex}. \\
\textbf{\url{https://github.com/realagiorganization/dspy_rag_in_repo_docs_and_impl1/blob/master/src/repo_rag_lab/population_samples.py}}. Host: github.com. Occurrences: 1. External GitHub URL referenced from 1 file(s). Local companion path(s): none. Sources: \path{publication/todo-backlog-table.tex}. & \textbf{\url{https://github.com/realagiorganization/dspy_rag_in_repo_docs_and_impl1/blob/master/src/repo_rag_lab/population_samples.py}}. Хост: github.com. Упоминаний: 1. Внешний URL GitHub, упомянутый в 1 файле(ах). Локальные сопроводительные пути: none. Источники: \path{publication/todo-backlog-table.tex}. \\
\textbf{\url{https://github.com/realagiorganization/dspy_rag_in_repo_docs_and_impl1/blob/master/src/repo_rag_lab/retrieval.py}}. Host: github.com. Occurrences: 1. External GitHub URL referenced from 1 file(s). Local companion path(s): none. Sources: \path{publication/todo-backlog-table.tex}. & \textbf{\url{https://github.com/realagiorganization/dspy_rag_in_repo_docs_and_impl1/blob/master/src/repo_rag_lab/retrieval.py}}. Хост: github.com. Упоминаний: 1. Внешний URL GitHub, упомянутый в 1 файле(ах). Локальные сопроводительные пути: none. Источники: \path{publication/todo-backlog-table.tex}. \\
\textbf{\url{https://github.com/realagiorganization/dspy_rag_in_repo_docs_and_impl1/blob/master/src/repo_rag_lab/utilities.py\#L108-L124}}. Host: github.com. Occurrences: 2. External GitHub URL referenced from 1 file(s). Local companion path(s): none. Sources: \path{docs/operations/simple-end-to-end-verification-guide.md}. & \textbf{\url{https://github.com/realagiorganization/dspy_rag_in_repo_docs_and_impl1/blob/master/src/repo_rag_lab/utilities.py\#L108-L124}}. Хост: github.com. Упоминаний: 2. Внешний URL GitHub, упомянутый в 1 файле(ах). Локальные сопроводительные пути: none. Источники: \path{docs/operations/simple-end-to-end-verification-guide.md}. \\
\textbf{\url{https://github.com/realagiorganization/dspy_rag_in_repo_docs_and_impl1/blob/master/src/repo_rag_lab/utilities.py\#L190-L193}}. Host: github.com. Occurrences: 1. External GitHub URL referenced from 1 file(s). Local companion path(s): none. Sources: \path{docs/operations/simple-end-to-end-verification-guide.md}. & \textbf{\url{https://github.com/realagiorganization/dspy_rag_in_repo_docs_and_impl1/blob/master/src/repo_rag_lab/utilities.py\#L190-L193}}. Хост: github.com. Упоминаний: 1. Внешний URL GitHub, упомянутый в 1 файле(ах). Локальные сопроводительные пути: none. Источники: \path{docs/operations/simple-end-to-end-verification-guide.md}. \\
\textbf{\url{https://github.com/realagiorganization/dspy_rag_in_repo_docs_and_impl1/blob/master/src/repo_rag_lab/verification.py\#L11-L51}}. Host: github.com. Occurrences: 1. External GitHub URL referenced from 1 file(s). Local companion path(s): none. Sources: \path{docs/operations/simple-end-to-end-verification-guide.md}. & \textbf{\url{https://github.com/realagiorganization/dspy_rag_in_repo_docs_and_impl1/blob/master/src/repo_rag_lab/verification.py\#L11-L51}}. Хост: github.com. Упоминаний: 1. Внешний URL GitHub, упомянутый в 1 файле(ах). Локальные сопроводительные пути: none. Источники: \path{docs/operations/simple-end-to-end-verification-guide.md}. \\
\textbf{\url{https://github.com/realagiorganization/dspy_rag_in_repo_docs_and_impl1/blob/master/src/repo_rag_lab/verification.py\#L133-L146}}. Host: github.com. Occurrences: 1. External GitHub URL referenced from 1 file(s). Local companion path(s): none. Sources: \path{docs/operations/simple-end-to-end-verification-guide.md}. & \textbf{\url{https://github.com/realagiorganization/dspy_rag_in_repo_docs_and_impl1/blob/master/src/repo_rag_lab/verification.py\#L133-L146}}. Хост: github.com. Упоминаний: 1. Внешний URL GitHub, упомянутый в 1 файле(ах). Локальные сопроводительные пути: none. Источники: \path{docs/operations/simple-end-to-end-verification-guide.md}. \\
\textbf{\url{https://github.com/realagiorganization/dspy_rag_in_repo_docs_and_impl1/blob/master/src/repo_rag_lab/workflow.py}}. Host: github.com. Occurrences: 1. External GitHub URL referenced from 1 file(s). Local companion path(s): none. Sources: \path{publication/todo-backlog-table.tex}. & \textbf{\url{https://github.com/realagiorganization/dspy_rag_in_repo_docs_and_impl1/blob/master/src/repo_rag_lab/workflow.py}}. Хост: github.com. Упоминаний: 1. Внешний URL GitHub, упомянутый в 1 файле(ах). Локальные сопроводительные пути: none. Источники: \path{publication/todo-backlog-table.tex}. \\
\textbf{\url{https://github.com/realagiorganization/dspy_rag_in_repo_docs_and_impl1/blob/master/src/repo_rag_lab/workflow.py\#L119-L145}}. Host: github.com. Occurrences: 1. External GitHub URL referenced from 1 file(s). Local companion path(s): none. Sources: \path{docs/operations/simple-end-to-end-verification-guide.md}. & \textbf{\url{https://github.com/realagiorganization/dspy_rag_in_repo_docs_and_impl1/blob/master/src/repo_rag_lab/workflow.py\#L119-L145}}. Хост: github.com. Упоминаний: 1. Внешний URL GitHub, упомянутый в 1 файле(ах). Локальные сопроводительные пути: none. Источники: \path{docs/operations/simple-end-to-end-verification-guide.md}. \\
\textbf{\url{https://github.com/realagiorganization/dspy_rag_in_repo_docs_and_impl1/blob/master/src/repo_rag_lab/workflow.py\#L54-L72}}. Host: github.com. Occurrences: 1. External GitHub URL referenced from 1 file(s). Local companion path(s): none. Sources: \path{docs/operations/simple-end-to-end-verification-guide.md}. & \textbf{\url{https://github.com/realagiorganization/dspy_rag_in_repo_docs_and_impl1/blob/master/src/repo_rag_lab/workflow.py\#L54-L72}}. Хост: github.com. Упоминаний: 1. Внешний URL GitHub, упомянутый в 1 файле(ах). Локальные сопроводительные пути: none. Источники: \path{docs/operations/simple-end-to-end-verification-guide.md}. \\
\textbf{\url{https://github.com/realagiorganization/dspy_rag_in_repo_docs_and_impl1/blob/master/tests/test_benchmarks_and_notebook_scaffolding.py}}. Host: github.com. Occurrences: 1. External GitHub URL referenced from 1 file(s). Local companion path(s): none. Sources: \path{publication/todo-backlog-table.tex}. & \textbf{\url{https://github.com/realagiorganization/dspy_rag_in_repo_docs_and_impl1/blob/master/tests/test_benchmarks_and_notebook_scaffolding.py}}. Хост: github.com. Упоминаний: 1. Внешний URL GitHub, упомянутый в 1 файле(ах). Локальные сопроводительные пути: none. Источники: \path{publication/todo-backlog-table.tex}. \\
\textbf{\url{https://github.com/realagiorganization/dspy_rag_in_repo_docs_and_impl1/blob/master/tests/test_repository_rag_bdd.py}}. Host: github.com. Occurrences: 1. External GitHub URL referenced from 1 file(s). Local companion path(s): none. Sources: \path{publication/todo-backlog-table.tex}. & \textbf{\url{https://github.com/realagiorganization/dspy_rag_in_repo_docs_and_impl1/blob/master/tests/test_repository_rag_bdd.py}}. Хост: github.com. Упоминаний: 1. Внешний URL GitHub, упомянутый в 1 файле(ах). Локальные сопроводительные пути: none. Источники: \path{publication/todo-backlog-table.tex}. \\
\textbf{\url{https://github.com/realagiorganization/dspy_rag_in_repo_docs_and_impl1/blob/master/utilities/README.md}}. Host: github.com. Occurrences: 1. External GitHub URL referenced from 1 file(s). Local companion path(s): none. Sources: \path{publication/todo-backlog-table.tex}. & \textbf{\url{https://github.com/realagiorganization/dspy_rag_in_repo_docs_and_impl1/blob/master/utilities/README.md}}. Хост: github.com. Упоминаний: 1. Внешний URL GitHub, упомянутый в 1 файле(ах). Локальные сопроводительные пути: none. Источники: \path{publication/todo-backlog-table.tex}. \\
\textbf{\url{https://github.com/realagiorganization/dspy_rag_in_repo_docs_and_impl1/tree/master/samples/logs/}}. Host: github.com. Occurrences: 3. External GitHub URL referenced from 1 file(s). Local companion path(s): none. Sources: \path{publication/todo-backlog-table.tex}. & \textbf{\url{https://github.com/realagiorganization/dspy_rag_in_repo_docs_and_impl1/tree/master/samples/logs/}}. Хост: github.com. Упоминаний: 3. Внешний URL GitHub, упомянутый в 1 файле(ах). Локальные сопроводительные пути: none. Источники: \path{publication/todo-backlog-table.tex}. \\
\textbf{\url{https://github.com/realagiorganization/dspy_rag_in_repo_docs_and_impl1/tree/master/tests/fixtures/hushwheel_lexiconarium/}}. Host: github.com. Occurrences: 1. External GitHub URL referenced from 1 file(s). Local companion path(s): none. Sources: \path{publication/todo-backlog-table.tex}. & \textbf{\url{https://github.com/realagiorganization/dspy_rag_in_repo_docs_and_impl1/tree/master/tests/fixtures/hushwheel_lexiconarium/}}. Хост: github.com. Упоминаний: 1. Внешний URL GitHub, упомянутый в 1 файле(ах). Локальные сопроводительные пути: none. Источники: \path{publication/todo-backlog-table.tex}. \\
\textbf{\url{https://github.com/realagiorganization/national-debt-relief}}. Host: github.com. Occurrences: 1. External GitHub URL referenced from 1 file(s). Local companion path(s): none. Sources: \path{docs/audit/2026-05-10-aks-run-25632110510-handoff-fixed-runtime-still-heuristic.md}. & \textbf{\url{https://github.com/realagiorganization/national-debt-relief}}. Хост: github.com. Упоминаний: 1. Внешний URL GitHub, упомянутый в 1 файле(ах). Локальные сопроводительные пути: none. Источники: \path{docs/audit/2026-05-10-aks-run-25632110510-handoff-fixed-runtime-still-heuristic.md}. \\
\textbf{\url{https://gpt45standard.openai.azure.com/}}. Host: gpt45standard.openai.azure.com. Occurrences: 5. Referenced by URL only. No repository-local mirror is currently tracked. Sources: \path{docs/audit/2026-04-29-live-azure-gpt54-validation.md}. & \textbf{\url{https://gpt45standard.openai.azure.com/}}. Хост: gpt45standard.openai.azure.com. Упоминаний: 5. Ссылка хранится только как URL. Локальное зеркало в репозитории сейчас не отслеживается. Источники: \path{docs/audit/2026-04-29-live-azure-gpt54-validation.md}. \\
\textbf{\url{https://gpt45standard.openai.azure.com/`}}. Host: gpt45standard.openai.azure.com. Occurrences: 3. Referenced by URL only. No repository-local mirror is currently tracked. Sources: \path{README.md}\newline \path{docs/audit/2026-04-29-live-azure-gpt54-validation.md}\newline \path{docs/operations/azure-deployment.md}. & \textbf{\url{https://gpt45standard.openai.azure.com/`}}. Хост: gpt45standard.openai.azure.com. Упоминаний: 3. Ссылка хранится только как URL. Локальное зеркало в репозитории сейчас не отслеживается. Источники: \path{README.md}\newline \path{docs/audit/2026-04-29-live-azure-gpt54-validation.md}\newline \path{docs/operations/azure-deployment.md}. \\
\textbf{\url{https://gpt45standard.openai.azure.com/openai/v1}}. Host: gpt45standard.openai.azure.com. Occurrences: 1. Referenced by URL only. No repository-local mirror is currently tracked. Sources: \path{docs/audit/2026-05-01-dataset-azure-responses-api-version-default-fix.md}. & \textbf{\url{https://gpt45standard.openai.azure.com/openai/v1}}. Хост: gpt45standard.openai.azure.com. Упоминаний: 1. Ссылка хранится только как URL. Локальное зеркало в репозитории сейчас не отслеживается. Источники: \path{docs/audit/2026-05-01-dataset-azure-responses-api-version-default-fix.md}. \\
\textbf{\url{https://gpt45standard.openai.azure.com`}}. Host: gpt45standard.openai.azure.com`. Occurrences: 1. Referenced by URL only. No repository-local mirror is currently tracked. Sources: \path{docs/audit/2026-03-18-zzzzzz-dspy-pipeline-and-live-azure-proof.md}. & \textbf{\url{https://gpt45standard.openai.azure.com`}}. Хост: gpt45standard.openai.azure.com`. Упоминаний: 1. Ссылка хранится только как URL. Локальное зеркало в репозитории сейчас не отслеживается. Источники: \path{docs/audit/2026-03-18-zzzzzz-dspy-pipeline-and-live-azure-proof.md}. \\
\textbf{\url{https://img.shields.io/badge/coverage-94.12\%25-brightgreen)](https://github.com/realagiorganization/dspy_rag_in_repo_docs_and_impl1/actions/workflows/ci.yml}}. Host: img.shields.io. Occurrences: 1. Referenced by URL only. No repository-local mirror is currently tracked. Sources: \path{README.md}. & \textbf{\url{https://img.shields.io/badge/coverage-94.12\%25-brightgreen)](https://github.com/realagiorganization/dspy_rag_in_repo_docs_and_impl1/actions/workflows/ci.yml}}. Хост: img.shields.io. Упоминаний: 1. Ссылка хранится только как URL. Локальное зеркало в репозитории сейчас не отслеживается. Источники: \path{README.md}. \\
\textbf{\url{https://learn.microsoft.com/azure/ai-foundry/foundry-models/how-to/inference}}. Host: learn.microsoft.com. Occurrences: 1. External Azure documentation URL. No mirrored upstream copy is tracked, but local companion notes exist at docs/operations/azure-deployment.md. Sources: \path{docs/operations/azure-deployment.md}. & \textbf{\url{https://learn.microsoft.com/azure/ai-foundry/foundry-models/how-to/inference}}. Хост: learn.microsoft.com. Упоминаний: 1. Внешний URL документации Azure. Зеркальная копия upstream не хранится, но есть локальные сопроводительные заметки: docs/operations/azure-deployment.md. Источники: \path{docs/operations/azure-deployment.md}. \\
\textbf{\url{https://learn.microsoft.com/azure/ai-services/}}. Host: learn.microsoft.com. Occurrences: 4. External Azure documentation URL. No mirrored upstream copy is tracked, but local companion notes exist at docs/operations/azure-deployment.md. Sources: \path{publication/references.bib}\newline \path{tests/test_exploratorium_translation.py}\newline \path{tests/test_utilities.py}. & \textbf{\url{https://learn.microsoft.com/azure/ai-services/}}. Хост: learn.microsoft.com. Упоминаний: 4. Внешний URL документации Azure. Зеркальная копия upstream не хранится, но есть локальные сопроводительные заметки: docs/operations/azure-deployment.md. Источники: \path{publication/references.bib}\newline \path{tests/test_exploratorium_translation.py}\newline \path{tests/test_utilities.py}. \\
\textbf{\url{https://learn.microsoft.com/en-us/azure/ai-foundry/openai/how-to/fine-tuning-deploy}}. Host: learn.microsoft.com. Occurrences: 1. External Azure documentation URL. No mirrored upstream copy is tracked, but local companion notes exist at docs/operations/azure-deployment.md. Sources: \path{docs/operations/azure-deployment.md}. & \textbf{\url{https://learn.microsoft.com/en-us/azure/ai-foundry/openai/how-to/fine-tuning-deploy}}. Хост: learn.microsoft.com. Упоминаний: 1. Внешний URL документации Azure. Зеркальная копия upstream не хранится, но есть локальные сопроводительные заметки: docs/operations/azure-deployment.md. Источники: \path{docs/operations/azure-deployment.md}. \\
\textbf{\url{https://learn.microsoft.com/en-us/azure/ai-services/openai/how-to/fine-tuning}}. Host: learn.microsoft.com. Occurrences: 1. External Azure documentation URL. No mirrored upstream copy is tracked, but local companion notes exist at docs/operations/azure-deployment.md. Sources: \path{docs/operations/azure-deployment.md}. & \textbf{\url{https://learn.microsoft.com/en-us/azure/ai-services/openai/how-to/fine-tuning}}. Хост: learn.microsoft.com. Упоминаний: 1. Внешний URL документации Azure. Зеркальная копия upstream не хранится, но есть локальные сопроводительные заметки: docs/operations/azure-deployment.md. Источники: \path{docs/operations/azure-deployment.md}. \\
\textbf{\url{https://medium.com/@arancibia.juan22/implementing-rag-with-dspy-a-technical-guide-a6ae15f6a455}}. Host: medium.com. Occurrences: 1. Referenced by URL only. No repository-local mirror is currently tracked. Sources: \path{docs/architecture/inspired/implementing-rag-with-dspy-technical-guide.md}. & \textbf{\url{https://medium.com/@arancibia.juan22/implementing-rag-with-dspy-a-technical-guide-a6ae15f6a455}}. Хост: medium.com. Упоминаний: 1. Ссылка хранится только как URL. Локальное зеркало в репозитории сейчас не отслеживается. Источники: \path{docs/architecture/inspired/implementing-rag-with-dspy-technical-guide.md}. \\
\textbf{\url{https://modelcontextprotocol.io/}}. Host: modelcontextprotocol.io. Occurrences: 4. External MCP documentation URL. No mirrored upstream copy is tracked, but local companion notes exist at docs/architecture/mcp-discovery.md. Sources: \path{publication/references.bib}\newline \path{tests/test_exploratorium_translation.py}\newline \path{tests/test_utilities.py}. & \textbf{\url{https://modelcontextprotocol.io/}}. Хост: modelcontextprotocol.io. Упоминаний: 4. Внешний URL документации MCP. Зеркальная копия upstream не хранится, но есть локальные сопроводительные заметки: docs/architecture/mcp-discovery.md. Источники: \path{publication/references.bib}\newline \path{tests/test_exploratorium_translation.py}\newline \path{tests/test_utilities.py}. \\
\textbf{\url{https://realagiorganization.github.io/dspy_rag_in_repo_docs_and_impl1/}}. Host: realagiorganization.github.io. Occurrences: 3. Referenced by URL only. No repository-local mirror is currently tracked. Sources: \path{README.md}\newline \path{docs/audit/2026-03-19-public-pages-and-secret-scan.md}\newline \path{mkdocs.yml}. & \textbf{\url{https://realagiorganization.github.io/dspy_rag_in_repo_docs_and_impl1/}}. Хост: realagiorganization.github.io. Упоминаний: 3. Ссылка хранится только как URL. Локальное зеркало в репозитории сейчас не отслеживается. Источники: \path{README.md}\newline \path{docs/audit/2026-03-19-public-pages-and-secret-scan.md}\newline \path{mkdocs.yml}. \\
\textbf{\url{https://realagiorganization.github.io/dspy_rag_in_repo_docs_and_impl1/`}}. Host: realagiorganization.github.io. Occurrences: 1. Referenced by URL only. No repository-local mirror is currently tracked. Sources: \path{docs/audit/2026-03-19-public-pages-and-secret-scan.md}. & \textbf{\url{https://realagiorganization.github.io/dspy_rag_in_repo_docs_and_impl1/`}}. Хост: realagiorganization.github.io. Упоминаний: 1. Ссылка хранится только как URL. Локальное зеркало в репозитории сейчас не отслеживается. Источники: \path{docs/audit/2026-03-19-public-pages-and-secret-scan.md}. \\
\textbf{\url{https://realagistorage.blob.core.windows.net}}. Host: realagistorage.blob.core.windows.net. Occurrences: 1. Referenced by URL only. No repository-local mirror is currently tracked. Sources: \path{tests/test_runtime_artifacts_azure.py}. & \textbf{\url{https://realagistorage.blob.core.windows.net}}. Хост: realagistorage.blob.core.windows.net. Упоминаний: 1. Ссылка хранится только как URL. Локальное зеркало в репозитории сейчас не отслеживается. Источники: \path{tests/test_runtime_artifacts_azure.py}. \\
\textbf{\url{https://realagistorage.queue.core.windows.net}}. Host: realagistorage.queue.core.windows.net. Occurrences: 1. Referenced by URL only. No repository-local mirror is currently tracked. Sources: \path{tests/test_runtime_artifacts_azure.py}. & \textbf{\url{https://realagistorage.queue.core.windows.net}}. Хост: realagistorage.queue.core.windows.net. Упоминаний: 1. Ссылка хранится только как URL. Локальное зеркало в репозитории сейчас не отслеживается. Источники: \path{tests/test_runtime_artifacts_azure.py}. \\
\bottomrule
\end{longtable}
\endgroup
\section{Line-By-Line / Построчный показ}
\subsection*{Referenced Papers And Documentation Fetch State / Состояние загрузки статей и документации}
\textbf{lewis2020rag}: Retrieval-Augmented Generation for Knowledge-Intensive NLP Tasks. Kind: inproceedings. Cited in the bibliography, but no local PDF or paper mirror is tracked. Closest local companion paths: docs/architecture/dspy-rag-guide.md, docs/architecture/inspired/implementing-rag-with-dspy-technical-guide.md. Related paths: \path{docs/architecture/dspy-rag-guide.md, docs/architecture/inspired/implementing-rag-with-dspy-technical-guide.md}.\\
\textbf{lewis2020rag}: Retrieval-Augmented Generation for Knowledge-Intensive NLP Tasks. Тип: inproceedings. Источник есть в библиографии, но локальный PDF или зеркало статьи не отслеживаются. Ближайшие локальные сопроводительные пути: docs/architecture/dspy-rag-guide.md, docs/architecture/inspired/implementing-rag-with-dspy-technical-guide.md. Связанные пути: \path{docs/architecture/dspy-rag-guide.md, docs/architecture/inspired/implementing-rag-with-dspy-technical-guide.md}.\par\medskip
\textbf{khattab2024dspy}: DSPy: Compiling Declarative Language Model Calls into Self-Improving Pipelines. Kind: article. Cited in the bibliography, but no local PDF or paper mirror is tracked. Closest local companion paths: docs/architecture/dspy-rag-guide.md, docs/architecture/inspired/dspy-rag-tutorial.md, docs/guides/hushwheel-fixture-rag-guide.md. Related paths: \path{docs/architecture/dspy-rag-guide.md, docs/architecture/inspired/dspy-rag-tutorial.md, docs/guides/hushwheel-fixture-rag-guide.md}.\\
\textbf{khattab2024dspy}: DSPy: Compiling Declarative Language Model Calls into Self-Improving Pipelines. Тип: article. Источник есть в библиографии, но локальный PDF или зеркало статьи не отслеживаются. Ближайшие локальные сопроводительные пути: docs/architecture/dspy-rag-guide.md, docs/architecture/inspired/dspy-rag-tutorial.md, docs/guides/hushwheel-fixture-rag-guide.md. Связанные пути: \path{docs/architecture/dspy-rag-guide.md, docs/architecture/inspired/dspy-rag-tutorial.md, docs/guides/hushwheel-fixture-rag-guide.md}.\par\medskip
\textbf{mcp2024}: Model Context Protocol. Kind: misc. External documentation is cited by URL. No mirrored upstream copy is tracked, but local companion paths exist at docs/architecture/mcp-discovery.md. Related paths: \path{docs/architecture/mcp-discovery.md}.\\
\textbf{mcp2024}: Model Context Protocol. Тип: misc. Внешняя документация цитируется по URL. Зеркальная копия upstream не отслеживается, но есть локальные сопроводительные пути: docs/architecture/mcp-discovery.md. Связанные пути: \path{docs/architecture/mcp-discovery.md}.\par\medskip
\textbf{azureinference2025}: Azure AI Inference Documentation. Kind: misc. External documentation is cited by URL. No mirrored upstream copy is tracked, but local companion paths exist at docs/operations/azure-deployment.md. Related paths: \path{docs/operations/azure-deployment.md}.\\
\textbf{azureinference2025}: Azure AI Inference Documentation. Тип: misc. Внешняя документация цитируется по URL. Зеркальная копия upstream не отслеживается, но есть локальные сопроводительные пути: docs/operations/azure-deployment.md. Связанные пути: \path{docs/operations/azure-deployment.md}.\par\medskip
\subsection*{Summaries Of All Files / Сводки по всем файлам}
\textbf{\path{.codex/skills/exploratorium-translation-sync/SKILL.md}}. Kind: markdown. Size: 1722 bytes. Shape: 32 lines. Markdown document, 32 lines, 0 explicit URLs. Lead heading: Exploratorium Translation Sync.\\
\textbf{\path{.codex/skills/exploratorium-translation-sync/SKILL.md}}. Тип: markdown. Размер: 1722 байт. Форма: 32 строк. Документ Markdown, 32 строк, явных URL: 0. Главный заголовок: Exploratorium Translation Sync.\par\medskip
\textbf{\path{.codex/skills/file-summary-sync/SKILL.md}}. Kind: markdown. Size: 1367 bytes. Shape: 29 lines. Markdown document, 29 lines, 0 explicit URLs. Lead heading: File Summary Sync.\\
\textbf{\path{.codex/skills/file-summary-sync/SKILL.md}}. Тип: markdown. Размер: 1367 байт. Форма: 29 строк. Документ Markdown, 29 строк, явных URL: 0. Главный заголовок: File Summary Sync.\par\medskip
\textbf{\path{.codex/skills/notebook-playbook-sync/SKILL.md}}. Kind: markdown. Size: 2085 bytes. Shape: 48 lines. Markdown document, 48 lines, 0 explicit URLs. Lead heading: Notebook Playbook Sync.\\
\textbf{\path{.codex/skills/notebook-playbook-sync/SKILL.md}}. Тип: markdown. Размер: 2085 байт. Форма: 48 строк. Документ Markdown, 48 строк, явных URL: 0. Главный заголовок: Notebook Playbook Sync.\par\medskip
\textbf{\path{.codex/skills/notebook-playbook-sync/agents/openai.yaml}}. Kind: yaml. Size: 300 bytes. Shape: 5 lines. YAML data or workflow file, 5 lines, 0 explicit URLs. Top keys: interface.\\
\textbf{\path{.codex/skills/notebook-playbook-sync/agents/openai.yaml}}. Тип: yaml. Размер: 300 байт. Форма: 5 строк. Файл YAML с данными или workflow, 5 строк, явных URL: 0. Верхние ключи: interface.\par\medskip
\textbf{\path{.codex/skills/post-push-gh-run-logging/SKILL.md}}. Kind: markdown. Size: 2059 bytes. Shape: 41 lines. Markdown document, 41 lines, 0 explicit URLs. Lead heading: Post Push Gh Run Logging.\\
\textbf{\path{.codex/skills/post-push-gh-run-logging/SKILL.md}}. Тип: markdown. Размер: 2059 байт. Форма: 41 строк. Документ Markdown, 41 строк, явных URL: 0. Главный заголовок: Post Push Gh Run Logging.\par\medskip
\textbf{\path{.codex/skills/post-push-gh-run-logging/agents/openai.yaml}}. Kind: yaml. Size: 281 bytes. Shape: 5 lines. YAML data or workflow file, 5 lines, 0 explicit URLs. Top keys: interface.\\
\textbf{\path{.codex/skills/post-push-gh-run-logging/agents/openai.yaml}}. Тип: yaml. Размер: 281 байт. Форма: 5 строк. Файл YAML с данными или workflow, 5 строк, явных URL: 0. Верхние ключи: interface.\par\medskip
\textbf{\path{.codex/skills/post-push-gh-run-logging/scripts/write_gh_run_log.py}}. Kind: python. Size: 6346 bytes. Shape: 178 lines. Python module, 178 lines, 0 explicit URLs. Module intent: No module docstring detected.\\
\textbf{\path{.codex/skills/post-push-gh-run-logging/scripts/write_gh_run_log.py}}. Тип: python. Размер: 6346 байт. Форма: 178 строк. Модуль Python, 178 строк, явных URL: 0. Назначение модуля: No module docstring detected.\par\medskip
\textbf{\path{.codex/skills/repo-verification-audit-loop/SKILL.md}}. Kind: markdown. Size: 2589 bytes. Shape: 53 lines. Markdown document, 53 lines, 0 explicit URLs. Lead heading: Repo Verification Audit Loop.\\
\textbf{\path{.codex/skills/repo-verification-audit-loop/SKILL.md}}. Тип: markdown. Размер: 2589 байт. Форма: 53 строк. Документ Markdown, 53 строк, явных URL: 0. Главный заголовок: Repo Verification Audit Loop.\par\medskip
\textbf{\path{.codex/skills/repo-verification-audit-loop/agents/openai.yaml}}. Kind: yaml. Size: 272 bytes. Shape: 5 lines. YAML data or workflow file, 5 lines, 0 explicit URLs. Top keys: interface.\\
\textbf{\path{.codex/skills/repo-verification-audit-loop/agents/openai.yaml}}. Тип: yaml. Размер: 272 байт. Форма: 5 строк. Файл YAML с данными или workflow, 5 строк, явных URL: 0. Верхние ключи: interface.\par\medskip
\textbf{\path{.codex/skills/rust-sqlite-lookup/SKILL.md}}. Kind: markdown. Size: 1297 bytes. Shape: 36 lines. Markdown document, 36 lines, 0 explicit URLs. Lead heading: Rust SQLite Lookup.\\
\textbf{\path{.codex/skills/rust-sqlite-lookup/SKILL.md}}. Тип: markdown. Размер: 1297 байт. Форма: 36 строк. Документ Markdown, 36 строк, явных URL: 0. Главный заголовок: Rust SQLite Lookup.\par\medskip
\textbf{\path{.codex/skills/todo-backlog-sync/SKILL.md}}. Kind: markdown. Size: 1289 bytes. Shape: 29 lines. Markdown document, 29 lines, 0 explicit URLs. Lead heading: Todo Backlog Sync.\\
\textbf{\path{.codex/skills/todo-backlog-sync/SKILL.md}}. Тип: markdown. Размер: 1289 байт. Форма: 29 строк. Документ Markdown, 29 строк, явных URL: 0. Главный заголовок: Todo Backlog Sync.\par\medskip
\textbf{\path{.dockerignore}}. Kind: generic. Size: 150 bytes. Shape: 17 lines. Tracked repository file, 17 lines, 0 explicit URLs. Opening line: .git\\
\textbf{\path{.dockerignore}}. Тип: generic. Размер: 150 байт. Форма: 17 строк. Отслеживаемый файл репозитория, 17 строк, явных URL: 0. Первая строка: .git\par\medskip
\textbf{\path{.env.sample}}. Kind: generic. Size: 905 bytes. Shape: 17 lines. Tracked repository file, 17 lines, 0 explicit URLs. Opening line: \# Copy this file to .env and replace placeholder values.\\
\textbf{\path{.env.sample}}. Тип: generic. Размер: 905 байт. Форма: 17 строк. Отслеживаемый файл репозитория, 17 строк, явных URL: 0. Первая строка: \# Copy this file to .env and replace placeholder values.\par\medskip
\textbf{\path{.gitattributes}}. Kind: generic. Size: 126 bytes. Shape: 4 lines. Tracked repository file, 4 lines, 0 explicit URLs. Opening line: *.png filter=lfs diff=lfs merge=lfs -text\\
\textbf{\path{.gitattributes}}. Тип: generic. Размер: 126 байт. Форма: 4 строк. Отслеживаемый файл репозитория, 4 строк, явных URL: 0. Первая строка: *.png filter=lfs diff=lfs merge=lfs -text\par\medskip
\textbf{\path{.github/workflows/ci.yml}}. Kind: yaml. Size: 4185 bytes. Shape: 146 lines. YAML data or workflow file, 146 lines, 0 explicit URLs. Top keys: name, True, jobs.\\
\textbf{\path{.github/workflows/ci.yml}}. Тип: yaml. Размер: 4185 байт. Форма: 146 строк. Файл YAML с данными или workflow, 146 строк, явных URL: 0. Верхние ключи: name, True, jobs.\par\medskip
\textbf{\path{.github/workflows/hushwheel-quality.yml}}. Kind: yaml. Size: 3526 bytes. Shape: 104 lines. YAML data or workflow file, 104 lines, 0 explicit URLs. Top keys: name, True, env.\\
\textbf{\path{.github/workflows/hushwheel-quality.yml}}. Тип: yaml. Размер: 3526 байт. Форма: 104 строк. Файл YAML с данными или workflow, 104 строк, явных URL: 0. Верхние ключи: name, True, env.\par\medskip
\textbf{\path{.github/workflows/pages.yml}}. Kind: yaml. Size: 2424 bytes. Shape: 99 lines. YAML data or workflow file, 99 lines, 0 explicit URLs. Top keys: name, True, env.\\
\textbf{\path{.github/workflows/pages.yml}}. Тип: yaml. Размер: 2424 байт. Форма: 99 строк. Файл YAML с данными или workflow, 99 строк, явных URL: 0. Верхние ключи: name, True, env.\par\medskip
\textbf{\path{.github/workflows/publication-pdf.yml}}. Kind: yaml. Size: 5549 bytes. Shape: 144 lines. YAML data or workflow file, 144 lines, 0 explicit URLs. Top keys: name, True, env.\\
\textbf{\path{.github/workflows/publication-pdf.yml}}. Тип: yaml. Размер: 5549 байт. Форма: 144 строк. Файл YAML с данными или workflow, 144 строк, явных URL: 0. Верхние ключи: name, True, env.\par\medskip
\textbf{\path{.github/workflows/publish.yml}}. Kind: yaml. Size: 847 bytes. Shape: 43 lines. YAML data or workflow file, 43 lines, 0 explicit URLs. Top keys: name, True, jobs.\\
\textbf{\path{.github/workflows/publish.yml}}. Тип: yaml. Размер: 847 байт. Форма: 43 строк. Файл YAML с данными или workflow, 43 строк, явных URL: 0. Верхние ключи: name, True, jobs.\par\medskip
\textbf{\path{.gitignore}}. Kind: generic. Size: 222 bytes. Shape: 21 lines. Tracked repository file, 21 lines, 0 explicit URLs. Opening line: .venv/\\
\textbf{\path{.gitignore}}. Тип: generic. Размер: 222 байт. Форма: 21 строк. Отслеживаемый файл репозитория, 21 строк, явных URL: 0. Первая строка: .venv/\par\medskip
\textbf{\path{.pre-commit-config.yaml}}. Kind: yaml. Size: 1689 bytes. Shape: 53 lines. YAML data or workflow file, 53 lines, 0 explicit URLs. Top keys: repos.\\
\textbf{\path{.pre-commit-config.yaml}}. Тип: yaml. Размер: 1689 байт. Форма: 53 строк. Файл YAML с данными или workflow, 53 строк, явных URL: 0. Верхние ключи: repos.\par\medskip
\textbf{\path{.python-version}}. Kind: generic. Size: 5 bytes. Shape: 2 lines. Tracked repository file, 2 lines, 0 explicit URLs. Opening line: 3.11\\
\textbf{\path{.python-version}}. Тип: generic. Размер: 5 байт. Форма: 2 строк. Отслеживаемый файл репозитория, 2 строк, явных URL: 0. Первая строка: 3.11\par\medskip
\textbf{\path{.yamllint.yml}}. Kind: yaml. Size: 416 bytes. Shape: 20 lines. YAML data or workflow file, 20 lines, 0 explicit URLs. Top keys: extends, rules.\\
\textbf{\path{.yamllint.yml}}. Тип: yaml. Размер: 416 байт. Форма: 20 строк. Файл YAML с данными или workflow, 20 строк, явных URL: 0. Верхние ключи: extends, rules.\par\medskip
\textbf{\path{AGENTS.md}}. Kind: markdown. Size: 6999 bytes. Shape: 108 lines. Markdown document, 108 lines, 0 explicit URLs. Lead heading: Repository Agent Instructions.\\
\textbf{\path{AGENTS.md}}. Тип: markdown. Размер: 6999 байт. Форма: 108 строк. Документ Markdown, 108 строк, явных URL: 0. Главный заголовок: Repository Agent Instructions.\par\medskip
\textbf{\path{AGENTS.md.d/00-audit-baseline.md}}. Kind: markdown. Size: 592 bytes. Shape: 15 lines. Markdown document, 15 lines, 0 explicit URLs. Lead heading: Audit Baseline Instructions.\\
\textbf{\path{AGENTS.md.d/00-audit-baseline.md}}. Тип: markdown. Размер: 592 байт. Форма: 15 строк. Документ Markdown, 15 строк, явных URL: 0. Главный заголовок: Audit Baseline Instructions.\par\medskip
\textbf{\path{AGENTS.md.d/10-verification-loop.md}}. Kind: markdown. Size: 1175 bytes. Shape: 22 lines. Markdown document, 22 lines, 0 explicit URLs. Lead heading: Verification Loop Instructions.\\
\textbf{\path{AGENTS.md.d/10-verification-loop.md}}. Тип: markdown. Размер: 1175 байт. Форма: 22 строк. Документ Markdown, 22 строк, явных URL: 0. Главный заголовок: Verification Loop Instructions.\par\medskip
\textbf{\path{AGENTS.md.d/FILES.md}}. Kind: markdown. Size: 2247 bytes. Shape: 41 lines. Markdown document, 41 lines, 0 explicit URLs. Lead heading: File Summary Instructions.\\
\textbf{\path{AGENTS.md.d/FILES.md}}. Тип: markdown. Размер: 2247 байт. Форма: 41 строк. Документ Markdown, 41 строк, явных URL: 0. Главный заголовок: File Summary Instructions.\par\medskip
\textbf{\path{AGENTS.md.d/RUST_LOOKUP.md}}. Kind: markdown. Size: 1803 bytes. Shape: 36 lines. Markdown document, 36 lines, 0 explicit URLs. Lead heading: Rust SQLite Lookup Instructions.\\
\textbf{\path{AGENTS.md.d/RUST_LOOKUP.md}}. Тип: markdown. Размер: 1803 байт. Форма: 36 строк. Документ Markdown, 36 строк, явных URL: 0. Главный заголовок: Rust SQLite Lookup Instructions.\par\medskip
\textbf{\path{Dockerfile}}. Kind: generic. Size: 1080 bytes. Shape: 47 lines. Tracked repository file, 47 lines, 0 explicit URLs. Opening line: ARG PYTHON\_BASE\_IMAGE=python:3.11-slim-bookworm\\
\textbf{\path{Dockerfile}}. Тип: generic. Размер: 1080 байт. Форма: 47 строк. Отслеживаемый файл репозитория, 47 строк, явных URL: 0. Первая строка: ARG PYTHON\_BASE\_IMAGE=python:3.11-slim-bookworm\par\medskip
\textbf{\path{FILES.csv}}. Kind: generic. Size: 52899 bytes. Shape: binary surface. Tracked repository file, 52899 bytes of binary data, 0 explicit URLs. Opening line: No non-empty preview line.\\
\textbf{\path{FILES.csv}}. Тип: generic. Размер: 52899 байт. Форма: бинарная поверхность. Отслеживаемый файл репозитория, 52899 байт бинарных данных, явных URL: 0. Первая строка: No non-empty preview line.\par\medskip
\textbf{\path{FILES.md}}. Kind: markdown. Size: 59726 bytes. Shape: 417 lines. Markdown document, 417 lines, 1 explicit URLs. Lead heading: Repository File Inventory.\\
\textbf{\path{FILES.md}}. Тип: markdown. Размер: 59726 байт. Форма: 417 строк. Документ Markdown, 417 строк, явных URL: 1. Главный заголовок: Repository File Inventory.\par\medskip
\textbf{\path{Makefile}}. Kind: make. Size: 23701 bytes. Shape: 506 lines. Repository Makefile, 506 lines, 1 explicit URLs. It exposes the shared automation surfaces.\\
\textbf{\path{Makefile}}. Тип: make. Размер: 23701 байт. Форма: 506 строк. Makefile репозитория, 506 строк, явных URL: 1. Он открывает общие поверхности автоматизации.\par\medskip
\textbf{\path{README.md}}. Kind: markdown. Size: 26222 bytes. Shape: 312 lines. Markdown document, 312 lines, 6 explicit URLs. Lead heading: Repository RAG Lab.\\
\textbf{\path{README.md}}. Тип: markdown. Размер: 26222 байт. Форма: 312 строк. Документ Markdown, 312 строк, явных URL: 6. Главный заголовок: Repository RAG Lab.\par\medskip
\textbf{\path{TODO.MD}}. Kind: markdown. Size: 6104 bytes. Shape: 22 lines. Markdown document, 22 lines, 0 explicit URLs. Lead heading: TODO.\\
\textbf{\path{TODO.MD}}. Тип: markdown. Размер: 6104 байт. Форма: 22 строк. Документ Markdown, 22 строк, явных URL: 0. Главный заголовок: TODO.\par\medskip
\textbf{\path{config/retrieval-profile.json}}. Kind: json. Size: 1791 bytes. Shape: 66 lines. JSON data file, 66 lines, 0 explicit URLs. Top keys: name, retrieval\_mode, rerank\_candidate\_pool\_size.\\
\textbf{\path{config/retrieval-profile.json}}. Тип: json. Размер: 1791 байт. Форма: 66 строк. Файл JSON, 66 строк, явных URL: 0. Верхние ключи: name, retrieval\_mode, rerank\_candidate\_pool\_size.\par\medskip
\textbf{\path{data/questions/repository.yaml}}. Kind: yaml. Size: 283 bytes. Shape: 7 lines. YAML data or workflow file, 7 lines, 0 explicit URLs. Top keys: no top-level mapping keys.\\
\textbf{\path{data/questions/repository.yaml}}. Тип: yaml. Размер: 283 байт. Форма: 7 строк. Файл YAML с данными или workflow, 7 строк, явных URL: 0. Верхние ключи: no top-level mapping keys.\par\medskip
\textbf{\path{docs/architecture/dspy-rag-guide.md}}. Kind: markdown. Size: 50272 bytes. Shape: 911 lines. Markdown document, 911 lines, 1 explicit URLs. Lead heading: DSPy Guide For This Repository.\\
\textbf{\path{docs/architecture/dspy-rag-guide.md}}. Тип: markdown. Размер: 50272 байт. Форма: 911 строк. Документ Markdown, 911 строк, явных URL: 1. Главный заголовок: DSPy Guide For This Repository.\par\medskip
\textbf{\path{docs/architecture/inspired/dspy-rag-tutorial.md}}. Kind: markdown. Size: 1474 bytes. Shape: 35 lines. Markdown document, 35 lines, 1 explicit URLs. Lead heading: DSPy Tutorial: Basic RAG.\\
\textbf{\path{docs/architecture/inspired/dspy-rag-tutorial.md}}. Тип: markdown. Размер: 1474 байт. Форма: 35 строк. Документ Markdown, 35 строк, явных URL: 1. Главный заголовок: DSPy Tutorial: Basic RAG.\par\medskip
\textbf{\path{docs/architecture/inspired/implementing-rag-with-dspy-technical-guide.md}}. Kind: markdown. Size: 1555 bytes. Shape: 38 lines. Markdown document, 38 lines, 1 explicit URLs. Lead heading: Implementing RAG with DSPy: Technical Guide.\\
\textbf{\path{docs/architecture/inspired/implementing-rag-with-dspy-technical-guide.md}}. Тип: markdown. Размер: 1555 байт. Форма: 38 строк. Документ Markdown, 38 строк, явных URL: 1. Главный заголовок: Implementing RAG with DSPy: Technical Guide.\par\medskip
\textbf{\path{docs/architecture/mcp-discovery.md}}. Kind: markdown. Size: 1757 bytes. Shape: 54 lines. Markdown document, 54 lines, 0 explicit URLs. Lead heading: MCP Discovery Notes.\\
\textbf{\path{docs/architecture/mcp-discovery.md}}. Тип: markdown. Размер: 1757 байт. Форма: 54 строк. Документ Markdown, 54 строк, явных URL: 0. Главный заголовок: MCP Discovery Notes.\par\medskip
\textbf{\path{docs/architecture/package-api.md}}. Kind: markdown. Size: 9186 bytes. Shape: 127 lines. Markdown document, 127 lines, 0 explicit URLs. Lead heading: Package API Notes.\\
\textbf{\path{docs/architecture/package-api.md}}. Тип: markdown. Размер: 9186 байт. Форма: 127 строк. Документ Markdown, 127 строк, явных URL: 0. Главный заголовок: Package API Notes.\par\medskip
\textbf{\path{docs/architecture/research-narrative.md}}. Kind: markdown. Size: 82061 bytes. Shape: 1077 lines. Markdown document, 1077 lines, 0 explicit URLs. Lead heading: Repository Research Narrative.\\
\textbf{\path{docs/architecture/research-narrative.md}}. Тип: markdown. Размер: 82061 байт. Форма: 1077 строк. Документ Markdown, 1077 строк, явных URL: 0. Главный заголовок: Repository Research Narrative.\par\medskip
\textbf{\path{docs/archive/repo-completeness-checklist.md}}. Kind: markdown. Size: 7450 bytes. Shape: 238 lines. Markdown document, 238 lines, 0 explicit URLs. Lead heading: Repository Completeness Checklist.\\
\textbf{\path{docs/archive/repo-completeness-checklist.md}}. Тип: markdown. Размер: 7450 байт. Форма: 238 строк. Документ Markdown, 238 строк, явных URL: 0. Главный заголовок: Repository Completeness Checklist.\par\medskip
\textbf{\path{docs/audit/2026-03-15-repo-state.md}}. Kind: markdown. Size: 5989 bytes. Shape: 151 lines. Markdown document, 151 lines, 0 explicit URLs. Lead heading: Repository State Audit.\\
\textbf{\path{docs/audit/2026-03-15-repo-state.md}}. Тип: markdown. Размер: 5989 байт. Форма: 151 строк. Документ Markdown, 151 строк, явных URL: 0. Главный заголовок: Repository State Audit.\par\medskip
\textbf{\path{docs/audit/2026-03-17-docs-and-uv-refresh.md}}. Kind: markdown. Size: 3977 bytes. Shape: 79 lines. Markdown document, 79 lines, 0 explicit URLs. Lead heading: Documentation And UV Workflow Refresh Audit.\\
\textbf{\path{docs/audit/2026-03-17-docs-and-uv-refresh.md}}. Тип: markdown. Размер: 3977 байт. Форма: 79 строк. Документ Markdown, 79 строк, явных URL: 0. Главный заголовок: Documentation And UV Workflow Refresh Audit.\par\medskip
\textbf{\path{docs/audit/2026-03-17-full-build.md}}. Kind: markdown. Size: 6715 bytes. Shape: 45 lines. Markdown document, 45 lines, 0 explicit URLs. Lead heading: Full Build Audit.\\
\textbf{\path{docs/audit/2026-03-17-full-build.md}}. Тип: markdown. Размер: 6715 байт. Форма: 45 строк. Документ Markdown, 45 строк, явных URL: 0. Главный заголовок: Full Build Audit.\par\medskip
\textbf{\path{docs/audit/2026-03-17-gh-actions-watch-loop.md}}. Kind: markdown. Size: 3797 bytes. Shape: 78 lines. Markdown document, 78 lines, 0 explicit URLs. Lead heading: GitHub Actions Watch Loop Audit.\\
\textbf{\path{docs/audit/2026-03-17-gh-actions-watch-loop.md}}. Тип: markdown. Размер: 3797 байт. Форма: 78 строк. Документ Markdown, 78 строк, явных URL: 0. Главный заголовок: GitHub Actions Watch Loop Audit.\par\medskip
\textbf{\path{docs/audit/2026-03-17-hushwheel-fixture.md}}. Kind: markdown. Size: 5090 bytes. Shape: 80 lines. Markdown document, 80 lines, 0 explicit URLs. Lead heading: Hushwheel Fixture Audit.\\
\textbf{\path{docs/audit/2026-03-17-hushwheel-fixture.md}}. Тип: markdown. Размер: 5090 байт. Форма: 80 строк. Документ Markdown, 80 строк, явных URL: 0. Главный заголовок: Hushwheel Fixture Audit.\par\medskip
\textbf{\path{docs/audit/2026-03-17-hushwheel-rag-playbook.md}}. Kind: markdown. Size: 4797 bytes. Shape: 83 lines. Markdown document, 83 lines, 0 explicit URLs. Lead heading: Hushwheel RAG Playbook Audit.\\
\textbf{\path{docs/audit/2026-03-17-hushwheel-rag-playbook.md}}. Тип: markdown. Размер: 4797 байт. Форма: 83 строк. Документ Markdown, 83 строк, явных URL: 0. Главный заголовок: Hushwheel RAG Playbook Audit.\par\medskip
\textbf{\path{docs/audit/2026-03-17-notebook-scaffolding.md}}. Kind: markdown. Size: 3641 bytes. Shape: 70 lines. Markdown document, 70 lines, 0 explicit URLs. Lead heading: Notebook Scaffolding Audit.\\
\textbf{\path{docs/audit/2026-03-17-notebook-scaffolding.md}}. Тип: markdown. Размер: 3641 байт. Форма: 70 строк. Документ Markdown, 70 строк, явных URL: 0. Главный заголовок: Notebook Scaffolding Audit.\par\medskip
\textbf{\path{docs/audit/2026-03-17-publication-article.md}}. Kind: markdown. Size: 2321 bytes. Shape: 54 lines. Markdown document, 54 lines, 0 explicit URLs. Lead heading: Publication Article Audit.\\
\textbf{\path{docs/audit/2026-03-17-publication-article.md}}. Тип: markdown. Размер: 2321 байт. Форма: 54 строк. Документ Markdown, 54 строк, явных URL: 0. Главный заголовок: Publication Article Audit.\par\medskip
\textbf{\path{docs/audit/2026-03-17-repo-completeness-checklist.md}}. Kind: markdown. Size: 3114 bytes. Shape: 59 lines. Markdown document, 59 lines, 0 explicit URLs. Lead heading: Repository Completeness Checklist Audit.\\
\textbf{\path{docs/audit/2026-03-17-repo-completeness-checklist.md}}. Тип: markdown. Размер: 3114 байт. Форма: 59 строк. Документ Markdown, 59 строк, явных URL: 0. Главный заголовок: Repository Completeness Checklist Audit.\par\medskip
\textbf{\path{docs/audit/2026-03-18-actions-caching.md}}. Kind: markdown. Size: 4189 bytes. Shape: 87 lines. Markdown document, 87 lines, 0 explicit URLs. Lead heading: GitHub Actions Caching Audit.\\
\textbf{\path{docs/audit/2026-03-18-actions-caching.md}}. Тип: markdown. Размер: 4189 байт. Форма: 87 строк. Документ Markdown, 87 строк, явных URL: 0. Главный заголовок: GitHub Actions Caching Audit.\par\medskip
\textbf{\path{docs/audit/2026-03-18-azure-inference-endpoint-probe.md}}. Kind: markdown. Size: 2648 bytes. Shape: 63 lines. Markdown document, 63 lines, 0 explicit URLs. Lead heading: Azure Inference Endpoint Probe Audit.\\
\textbf{\path{docs/audit/2026-03-18-azure-inference-endpoint-probe.md}}. Тип: markdown. Размер: 2648 байт. Форма: 63 строк. Документ Markdown, 63 строк, явных URL: 0. Главный заголовок: Azure Inference Endpoint Probe Audit.\par\medskip
\textbf{\path{docs/audit/2026-03-18-azure-runtime-surfaces.md}}. Kind: markdown. Size: 5307 bytes. Shape: 91 lines. Markdown document, 91 lines, 0 explicit URLs. Lead heading: Azure Runtime Surfaces Audit.\\
\textbf{\path{docs/audit/2026-03-18-azure-runtime-surfaces.md}}. Тип: markdown. Размер: 5307 байт. Форма: 91 строк. Документ Markdown, 91 строк, явных URL: 0. Главный заголовок: Azure Runtime Surfaces Audit.\par\medskip
\textbf{\path{docs/audit/2026-03-18-dspy-training-path.md}}. Kind: markdown. Size: 4605 bytes. Shape: 81 lines. Markdown document, 81 lines, 0 explicit URLs. Lead heading: DSPy Training Path Audit.\\
\textbf{\path{docs/audit/2026-03-18-dspy-training-path.md}}. Тип: markdown. Размер: 4605 байт. Форма: 81 строк. Документ Markdown, 81 строк, явных URL: 0. Главный заголовок: DSPy Training Path Audit.\par\medskip
\textbf{\path{docs/audit/2026-03-18-hushwheel-program-harness.md}}. Kind: markdown. Size: 3907 bytes. Shape: 76 lines. Markdown document, 76 lines, 0 explicit URLs. Lead heading: Hushwheel Program Harness Audit.\\
\textbf{\path{docs/audit/2026-03-18-hushwheel-program-harness.md}}. Тип: markdown. Размер: 3907 байт. Форма: 76 строк. Документ Markdown, 76 строк, явных URL: 0. Главный заголовок: Hushwheel Program Harness Audit.\par\medskip
\textbf{\path{docs/audit/2026-03-18-notebook-execution.md}}. Kind: markdown. Size: 2577 bytes. Shape: 66 lines. Markdown document, 66 lines, 0 explicit URLs. Lead heading: Notebook Execution Audit.\\
\textbf{\path{docs/audit/2026-03-18-notebook-execution.md}}. Тип: markdown. Размер: 2577 байт. Форма: 66 строк. Документ Markdown, 66 строк, явных URL: 0. Главный заголовок: Notebook Execution Audit.\par\medskip
\textbf{\path{docs/audit/2026-03-18-notebook-runner-harness.md}}. Kind: markdown. Size: 4121 bytes. Shape: 91 lines. Markdown document, 91 lines, 0 explicit URLs. Lead heading: Notebook Runner Harness Audit.\\
\textbf{\path{docs/audit/2026-03-18-notebook-runner-harness.md}}. Тип: markdown. Размер: 4121 байт. Форма: 91 строк. Документ Markdown, 91 строк, явных URL: 0. Главный заголовок: Notebook Runner Harness Audit.\par\medskip
\textbf{\path{docs/audit/2026-03-18-repository-benchmark-broadening.md}}. Kind: markdown. Size: 3728 bytes. Shape: 74 lines. Markdown document, 74 lines, 0 explicit URLs. Lead heading: Repository Benchmark Broadening Audit.\\
\textbf{\path{docs/audit/2026-03-18-repository-benchmark-broadening.md}}. Тип: markdown. Размер: 3728 байт. Форма: 74 строк. Документ Markdown, 74 строк, явных URL: 0. Главный заголовок: Repository Benchmark Broadening Audit.\par\medskip
\textbf{\path{docs/audit/2026-03-18-retest-with-env-refresh.md}}. Kind: markdown. Size: 5102 bytes. Shape: 105 lines. Markdown document, 105 lines, 0 explicit URLs. Lead heading: Env Refresh Retest Audit.\\
\textbf{\path{docs/audit/2026-03-18-retest-with-env-refresh.md}}. Тип: markdown. Размер: 5102 байт. Форма: 105 строк. Документ Markdown, 105 строк, явных URL: 0. Главный заголовок: Env Refresh Retest Audit.\par\medskip
\textbf{\path{docs/audit/2026-03-18-retrieval-evaluation-suite.md}}. Kind: markdown. Size: 4031 bytes. Shape: 80 lines. Markdown document, 80 lines, 0 explicit URLs. Lead heading: Retrieval Evaluation Suite Audit.\\
\textbf{\path{docs/audit/2026-03-18-retrieval-evaluation-suite.md}}. Тип: markdown. Размер: 4031 байт. Форма: 80 строк. Документ Markdown, 80 строк, явных URL: 0. Главный заголовок: Retrieval Evaluation Suite Audit.\par\medskip
\textbf{\path{docs/audit/2026-03-18-retrieval-ranking-refresh.md}}. Kind: markdown. Size: 6479 bytes. Shape: 109 lines. Markdown document, 109 lines, 0 explicit URLs. Lead heading: Retrieval Ranking Refresh Audit.\\
\textbf{\path{docs/audit/2026-03-18-retrieval-ranking-refresh.md}}. Тип: markdown. Размер: 6479 байт. Форма: 109 строк. Документ Markdown, 109 строк, явных URL: 0. Главный заголовок: Retrieval Ranking Refresh Audit.\par\medskip
\textbf{\path{docs/audit/2026-03-18-rust-and-benchmark-hygiene.md}}. Kind: markdown. Size: 2049 bytes. Shape: 53 lines. Markdown document, 53 lines, 0 explicit URLs. Lead heading: Rust And Benchmark Hygiene Audit.\\
\textbf{\path{docs/audit/2026-03-18-rust-and-benchmark-hygiene.md}}. Тип: markdown. Размер: 2049 байт. Форма: 53 строк. Документ Markdown, 53 строк, явных URL: 0. Главный заголовок: Rust And Benchmark Hygiene Audit.\par\medskip
\textbf{\path{docs/audit/2026-03-18-todo-backlog-sync.md}}. Kind: markdown. Size: 4104 bytes. Shape: 83 lines. Markdown document, 83 lines, 0 explicit URLs. Lead heading: Todo Backlog Sync Audit.\\
\textbf{\path{docs/audit/2026-03-18-todo-backlog-sync.md}}. Тип: markdown. Размер: 4104 байт. Форма: 83 строк. Документ Markdown, 83 строк, явных URL: 0. Главный заголовок: Todo Backlog Sync Audit.\par\medskip
\textbf{\path{docs/audit/2026-03-18-z-notebook-runner-harness.md}}. Kind: markdown. Size: 4153 bytes. Shape: 104 lines. Markdown document, 104 lines, 0 explicit URLs. Lead heading: Notebook Runner Harness Audit.\\
\textbf{\path{docs/audit/2026-03-18-z-notebook-runner-harness.md}}. Тип: markdown. Размер: 4153 байт. Форма: 104 строк. Документ Markdown, 104 строк, явных URL: 0. Главный заголовок: Notebook Runner Harness Audit.\par\medskip
\textbf{\path{docs/audit/2026-03-18-zz-file-summary-inventory.md}}. Kind: markdown. Size: 2596 bytes. Shape: 69 lines. Markdown document, 69 lines, 0 explicit URLs. Lead heading: File Summary Inventory Audit.\\
\textbf{\path{docs/audit/2026-03-18-zz-file-summary-inventory.md}}. Тип: markdown. Размер: 2596 байт. Форма: 69 строк. Документ Markdown, 69 строк, явных URL: 0. Главный заголовок: File Summary Inventory Audit.\par\medskip
\textbf{\path{docs/audit/2026-03-18-zz-research-narrative.md}}. Kind: markdown. Size: 3312 bytes. Shape: 79 lines. Markdown document, 79 lines, 0 explicit URLs. Lead heading: Research Narrative Audit.\\
\textbf{\path{docs/audit/2026-03-18-zz-research-narrative.md}}. Тип: markdown. Размер: 3312 байт. Форма: 79 строк. Документ Markdown, 79 строк, явных URL: 0. Главный заголовок: Research Narrative Audit.\par\medskip
\textbf{\path{docs/audit/2026-03-18-zzz-session-consolidation-on-origin-master.md}}. Kind: markdown. Size: 5061 bytes. Shape: 109 lines. Markdown document, 109 lines, 0 explicit URLs. Lead heading: Session Consolidation On Origin Master.\\
\textbf{\path{docs/audit/2026-03-18-zzz-session-consolidation-on-origin-master.md}}. Тип: markdown. Размер: 5061 байт. Форма: 109 строк. Документ Markdown, 109 строк, явных URL: 0. Главный заголовок: Session Consolidation On Origin Master.\par\medskip
\textbf{\path{docs/audit/2026-03-18-zzzz-rust-sqlite-lookup.md}}. Kind: markdown. Size: 5404 bytes. Shape: 105 lines. Markdown document, 105 lines, 0 explicit URLs. Lead heading: Rust SQLite Lookup And Retrieval Tag Audit.\\
\textbf{\path{docs/audit/2026-03-18-zzzz-rust-sqlite-lookup.md}}. Тип: markdown. Размер: 5404 байт. Форма: 105 строк. Документ Markdown, 105 строк, явных URL: 0. Главный заголовок: Rust SQLite Lookup And Retrieval Tag Audit.\par\medskip
\textbf{\path{docs/audit/2026-03-18-zzzz-session-consolidation-master-ci-repair.md}}. Kind: markdown. Size: 3270 bytes. Shape: 73 lines. Markdown document, 73 lines, 0 explicit URLs. Lead heading: Session Consolidation Master CI Repair.\\
\textbf{\path{docs/audit/2026-03-18-zzzz-session-consolidation-master-ci-repair.md}}. Тип: markdown. Размер: 3270 байт. Форма: 73 строк. Документ Markdown, 73 строк, явных URL: 0. Главный заголовок: Session Consolidation Master CI Repair.\par\medskip
\textbf{\path{docs/audit/2026-03-18-zzzzz-hushwheel-quality-instrumentation.md}}. Kind: markdown. Size: 8775 bytes. Shape: 165 lines. Markdown document, 165 lines, 0 explicit URLs. Lead heading: Hushwheel Quality Instrumentation.\\
\textbf{\path{docs/audit/2026-03-18-zzzzz-hushwheel-quality-instrumentation.md}}. Тип: markdown. Размер: 8775 байт. Форма: 165 строк. Документ Markdown, 165 строк, явных URL: 0. Главный заголовок: Hushwheel Quality Instrumentation.\par\medskip
\textbf{\path{docs/audit/2026-03-18-zzzzzz-dspy-pipeline-and-live-azure-proof.md}}. Kind: markdown. Size: 5992 bytes. Shape: 116 lines. Markdown document, 116 lines, 1 explicit URLs. Lead heading: DSPy Pipeline And Live Azure Proof.\\
\textbf{\path{docs/audit/2026-03-18-zzzzzz-dspy-pipeline-and-live-azure-proof.md}}. Тип: markdown. Размер: 5992 байт. Форма: 116 строк. Документ Markdown, 116 строк, явных URL: 1. Главный заголовок: DSPy Pipeline And Live Azure Proof.\par\medskip
\textbf{\path{docs/audit/2026-03-18-zzzzzz-hushwheel-retrieval-record-wording.md}}. Kind: markdown. Size: 3042 bytes. Shape: 77 lines. Markdown document, 77 lines, 0 explicit URLs. Lead heading: Hushwheel Retrieval Record Wording.\\
\textbf{\path{docs/audit/2026-03-18-zzzzzz-hushwheel-retrieval-record-wording.md}}. Тип: markdown. Размер: 3042 байт. Форма: 77 строк. Документ Markdown, 77 строк, явных URL: 0. Главный заголовок: Hushwheel Retrieval Record Wording.\par\medskip
\textbf{\path{docs/audit/2026-03-18-zzzzzzz-hushwheel-pdf-side-effect-fix.md}}. Kind: markdown. Size: 4044 bytes. Shape: 83 lines. Markdown document, 83 lines, 0 explicit URLs. Lead heading: Hushwheel PDF Side-Effect Fix.\\
\textbf{\path{docs/audit/2026-03-18-zzzzzzz-hushwheel-pdf-side-effect-fix.md}}. Тип: markdown. Размер: 4044 байт. Форма: 83 строк. Документ Markdown, 83 строк, явных URL: 0. Главный заголовок: Hushwheel PDF Side-Effect Fix.\par\medskip
\textbf{\path{docs/audit/2026-03-18-zzzzzzz-hushwheel-quality-salvage.md}}. Kind: markdown. Size: 4554 bytes. Shape: 84 lines. Markdown document, 84 lines, 0 explicit URLs. Lead heading: Hushwheel Quality Salvage Audit.\\
\textbf{\path{docs/audit/2026-03-18-zzzzzzz-hushwheel-quality-salvage.md}}. Тип: markdown. Размер: 4554 байт. Форма: 84 строк. Документ Markdown, 84 строк, явных URL: 0. Главный заголовок: Hushwheel Quality Salvage Audit.\par\medskip
\textbf{\path{docs/audit/2026-03-18-zzzzzzz-land-unlanded-surfaces.md}}. Kind: markdown. Size: 2685 bytes. Shape: 72 lines. Markdown document, 72 lines, 0 explicit URLs. Lead heading: Hushwheel Quality Landing Audit.\\
\textbf{\path{docs/audit/2026-03-18-zzzzzzz-land-unlanded-surfaces.md}}. Тип: markdown. Размер: 2685 байт. Форма: 72 строк. Документ Markdown, 72 строк, явных URL: 0. Главный заголовок: Hushwheel Quality Landing Audit.\par\medskip
\textbf{\path{docs/audit/2026-03-18-zzzzzzz-rebased-master-log-and-azure-typecheck-fix.md}}. Kind: markdown. Size: 3472 bytes. Shape: 83 lines. Markdown document, 83 lines, 0 explicit URLs. Lead heading: Rebased Master Log And Azure Typecheck Fix.\\
\textbf{\path{docs/audit/2026-03-18-zzzzzzz-rebased-master-log-and-azure-typecheck-fix.md}}. Тип: markdown. Размер: 3472 байт. Форма: 83 строк. Документ Markdown, 83 строк, явных URL: 0. Главный заголовок: Rebased Master Log And Azure Typecheck Fix.\par\medskip
\textbf{\path{docs/audit/2026-03-18-zzzzzzz-worktree-consolidation.md}}. Kind: markdown. Size: 5914 bytes. Shape: 125 lines. Markdown document, 125 lines, 0 explicit URLs. Lead heading: Worktree Consolidation.\\
\textbf{\path{docs/audit/2026-03-18-zzzzzzz-worktree-consolidation.md}}. Тип: markdown. Размер: 5914 байт. Форма: 125 строк. Документ Markdown, 125 строк, явных URL: 0. Главный заголовок: Worktree Consolidation.\par\medskip
\textbf{\path{docs/audit/2026-03-18-zzzzzzzzzzzz-retrieval-regression-gate.md}}. Kind: markdown. Size: 2898 bytes. Shape: 66 lines. Markdown document, 66 lines, 0 explicit URLs. Lead heading: Retrieval Regression Gate Audit.\\
\textbf{\path{docs/audit/2026-03-18-zzzzzzzzzzzz-retrieval-regression-gate.md}}. Тип: markdown. Размер: 2898 байт. Форма: 66 строк. Документ Markdown, 66 строк, явных URL: 0. Главный заголовок: Retrieval Regression Gate Audit.\par\medskip
\textbf{\path{docs/audit/2026-03-18-zzzzzzzzzzzzz-land-remaining-unlanded-work.md}}. Kind: markdown. Size: 2908 bytes. Shape: 78 lines. Markdown document, 78 lines, 0 explicit URLs. Lead heading: Land Remaining Unlanded Work.\\
\textbf{\path{docs/audit/2026-03-18-zzzzzzzzzzzzz-land-remaining-unlanded-work.md}}. Тип: markdown. Размер: 2908 байт. Форма: 78 строк. Документ Markdown, 78 строк, явных URL: 0. Главный заголовок: Land Remaining Unlanded Work.\par\medskip
\textbf{\path{docs/audit/2026-03-18-zzzzzzzzzzzzzz-node24-actions-refresh.md}}. Kind: markdown. Size: 2989 bytes. Shape: 91 lines. Markdown document, 91 lines, 0 explicit URLs. Lead heading: Node 24 Actions Refresh.\\
\textbf{\path{docs/audit/2026-03-18-zzzzzzzzzzzzzz-node24-actions-refresh.md}}. Тип: markdown. Размер: 2989 байт. Форма: 91 строк. Документ Markdown, 91 строк, явных URL: 0. Главный заголовок: Node 24 Actions Refresh.\par\medskip
\textbf{\path{docs/audit/2026-03-18-zzzzzzzzzzzzzz-rebase-hushwheel-quality-over-master.md}}. Kind: markdown. Size: 1889 bytes. Shape: 47 lines. Markdown document, 47 lines, 0 explicit URLs. Lead heading: Rebase Hushwheel Quality Over Master.\\
\textbf{\path{docs/audit/2026-03-18-zzzzzzzzzzzzzz-rebase-hushwheel-quality-over-master.md}}. Тип: markdown. Размер: 1889 байт. Форма: 47 строк. Документ Markdown, 47 строк, явных URL: 0. Главный заголовок: Rebase Hushwheel Quality Over Master.\par\medskip
\textbf{\path{docs/audit/2026-03-18-zzzzzzzzzzzzzzzz-dspy-backlog-reconciliation.md}}. Kind: markdown. Size: 2956 bytes. Shape: 75 lines. Markdown document, 75 lines, 0 explicit URLs. Lead heading: DSPy Backlog Reconciliation.\\
\textbf{\path{docs/audit/2026-03-18-zzzzzzzzzzzzzzzz-dspy-backlog-reconciliation.md}}. Тип: markdown. Размер: 2956 байт. Форма: 75 строк. Документ Markdown, 75 строк, явных URL: 0. Главный заголовок: DSPy Backlog Reconciliation.\par\medskip
\textbf{\path{docs/audit/2026-03-18-zzzzzzzzzzzzzzzz-file-summary-test-churn-fix.md}}. Kind: markdown. Size: 2929 bytes. Shape: 85 lines. Markdown document, 85 lines, 0 explicit URLs. Lead heading: File Summary Test Churn Fix.\\
\textbf{\path{docs/audit/2026-03-18-zzzzzzzzzzzzzzzz-file-summary-test-churn-fix.md}}. Тип: markdown. Размер: 2929 байт. Форма: 85 строк. Документ Markdown, 85 строк, явных URL: 0. Главный заголовок: File Summary Test Churn Fix.\par\medskip
\textbf{\path{docs/audit/2026-03-18-zzzzzzzzzzzzzzzz-publication-cache-node24-followup.md}}. Kind: markdown. Size: 2900 bytes. Shape: 89 lines. Markdown document, 89 lines, 0 explicit URLs. Lead heading: Publication Cache Node 24 Follow-Up.\\
\textbf{\path{docs/audit/2026-03-18-zzzzzzzzzzzzzzzz-publication-cache-node24-followup.md}}. Тип: markdown. Размер: 2900 байт. Форма: 89 строк. Документ Markdown, 89 строк, явных URL: 0. Главный заголовок: Publication Cache Node 24 Follow-Up.\par\medskip
\textbf{\path{docs/audit/2026-03-18-zzzzzzzzzzzzzzzz-retrieval-gate-restored-after-rebase.md}}. Kind: markdown. Size: 3788 bytes. Shape: 94 lines. Markdown document, 94 lines, 0 explicit URLs. Lead heading: Retrieval Gate Restored After Rebase.\\
\textbf{\path{docs/audit/2026-03-18-zzzzzzzzzzzzzzzz-retrieval-gate-restored-after-rebase.md}}. Тип: markdown. Размер: 3788 байт. Форма: 94 строк. Документ Markdown, 94 строк, явных URL: 0. Главный заголовок: Retrieval Gate Restored After Rebase.\par\medskip
\textbf{\path{docs/audit/2026-03-18-zzzzzzzzzzzzzzzzzz-dspy-backlog-reconciliation.md}}. Kind: markdown. Size: 2457 bytes. Shape: 62 lines. Markdown document, 62 lines, 0 explicit URLs. Lead heading: DSPy Backlog Reconciliation.\\
\textbf{\path{docs/audit/2026-03-18-zzzzzzzzzzzzzzzzzz-dspy-backlog-reconciliation.md}}. Тип: markdown. Размер: 2457 байт. Форма: 62 строк. Документ Markdown, 62 строк, явных URL: 0. Главный заголовок: DSPy Backlog Reconciliation.\par\medskip
\textbf{\path{docs/audit/2026-03-18-zzzzzzzzzzzzzzzzzz-lookup-first-ask-path.md}}. Kind: markdown. Size: 4604 bytes. Shape: 94 lines. Markdown document, 94 lines, 0 explicit URLs. Lead heading: Lookup First Ask Path.\\
\textbf{\path{docs/audit/2026-03-18-zzzzzzzzzzzzzzzzzz-lookup-first-ask-path.md}}. Тип: markdown. Размер: 4604 байт. Форма: 94 строк. Документ Markdown, 94 строк, явных URL: 0. Главный заголовок: Lookup First Ask Path.\par\medskip
\textbf{\path{docs/audit/2026-03-18-zzzzzzzzzzzzzzzzzzz-answer-synthesis-refresh.md}}. Kind: markdown. Size: 2565 bytes. Shape: 48 lines. Markdown document, 48 lines, 0 explicit URLs. Lead heading: 2026-03-18 Answer Synthesis Refresh.\\
\textbf{\path{docs/audit/2026-03-18-zzzzzzzzzzzzzzzzzzz-answer-synthesis-refresh.md}}. Тип: markdown. Размер: 2565 байт. Форма: 48 строк. Документ Markdown, 48 строк, явных URL: 0. Главный заголовок: 2026-03-18 Answer Synthesis Refresh.\par\medskip
\textbf{\path{docs/audit/2026-03-18-zzzzzzzzzzzzzzzzzzzz-retrieval-quality-doc-priority.md}}. Kind: markdown. Size: 3602 bytes. Shape: 68 lines. Markdown document, 68 lines, 0 explicit URLs. Lead heading: Retrieval Quality Doc Priority Audit.\\
\textbf{\path{docs/audit/2026-03-18-zzzzzzzzzzzzzzzzzzzz-retrieval-quality-doc-priority.md}}. Тип: markdown. Размер: 3602 байт. Форма: 68 строк. Документ Markdown, 68 строк, явных URL: 0. Главный заголовок: Retrieval Quality Doc Priority Audit.\par\medskip
\textbf{\path{docs/audit/2026-03-18-zzzzzzzzzzzzzzzzzzzzzzz-merge-hushwheel-quality-into-master.md}}. Kind: markdown. Size: 4335 bytes. Shape: 112 lines. Markdown document, 112 lines, 0 explicit URLs. Lead heading: Merge Hushwheel Quality Into Master.\\
\textbf{\path{docs/audit/2026-03-18-zzzzzzzzzzzzzzzzzzzzzzz-merge-hushwheel-quality-into-master.md}}. Тип: markdown. Размер: 4335 байт. Форма: 112 строк. Документ Markdown, 112 строк, явных URL: 0. Главный заголовок: Merge Hushwheel Quality Into Master.\par\medskip
\textbf{\path{docs/audit/2026-03-18-zzzzzzzzzzzzzzzzzzzzzzzz-hushwheel-atlas-consolidation.md}}. Kind: markdown. Size: 4762 bytes. Shape: 109 lines. Markdown document, 109 lines, 0 explicit URLs. Lead heading: Hushwheel Atlas Consolidation.\\
\textbf{\path{docs/audit/2026-03-18-zzzzzzzzzzzzzzzzzzzzzzzz-hushwheel-atlas-consolidation.md}}. Тип: markdown. Размер: 4762 байт. Форма: 109 строк. Документ Markdown, 109 строк, явных URL: 0. Главный заголовок: Hushwheel Atlas Consolidation.\par\medskip
\textbf{\path{docs/audit/2026-03-19-github-pr-gate-sync.md}}. Kind: markdown. Size: 6599 bytes. Shape: 136 lines. Markdown document, 136 lines, 0 explicit URLs. Lead heading: Master Consolidation And GitHub PR Gate Sync.\\
\textbf{\path{docs/audit/2026-03-19-github-pr-gate-sync.md}}. Тип: markdown. Размер: 6599 байт. Форма: 136 строк. Документ Markdown, 136 строк, явных URL: 0. Главный заголовок: Master Consolidation And GitHub PR Gate Sync.\par\medskip
\textbf{\path{docs/audit/2026-03-19-node24-actions-opt-in.md}}. Kind: markdown. Size: 2929 bytes. Shape: 80 lines. Markdown document, 80 lines, 0 explicit URLs. Lead heading: Node 24 Actions Opt-In.\\
\textbf{\path{docs/audit/2026-03-19-node24-actions-opt-in.md}}. Тип: markdown. Размер: 2929 байт. Форма: 80 строк. Документ Markdown, 80 строк, явных URL: 0. Главный заголовок: Node 24 Actions Opt-In.\par\medskip
\textbf{\path{docs/audit/2026-03-19-public-pages-and-secret-scan.md}}. Kind: markdown. Size: 5804 bytes. Shape: 113 lines. Markdown document, 113 lines, 2 explicit URLs. Lead heading: Public Pages And Secret Scan.\\
\textbf{\path{docs/audit/2026-03-19-public-pages-and-secret-scan.md}}. Тип: markdown. Размер: 5804 байт. Форма: 113 строк. Документ Markdown, 113 строк, явных URL: 2. Главный заголовок: Public Pages And Secret Scan.\par\medskip
\textbf{\path{docs/audit/2026-03-19-simple-end-to-end-verification-guide.md}}. Kind: markdown. Size: 4236 bytes. Shape: 103 lines. Markdown document, 103 lines, 0 explicit URLs. Lead heading: Simple End-To-End Verification Guide.\\
\textbf{\path{docs/audit/2026-03-19-simple-end-to-end-verification-guide.md}}. Тип: markdown. Размер: 4236 байт. Форма: 103 строк. Документ Markdown, 103 строк, явных URL: 0. Главный заголовок: Simple End-To-End Verification Guide.\par\medskip
\textbf{\path{docs/audit/2026-04-28-cli-envelope-normalization.md}}. Kind: markdown. Size: 4028 bytes. Shape: 95 lines. Markdown document, 95 lines, 0 explicit URLs. Lead heading: CLI Envelope Normalization.\\
\textbf{\path{docs/audit/2026-04-28-cli-envelope-normalization.md}}. Тип: markdown. Размер: 4028 байт. Форма: 95 строк. Документ Markdown, 95 строк, явных URL: 0. Главный заголовок: CLI Envelope Normalization.\par\medskip
\textbf{\path{docs/audit/2026-04-28-doc-tree-restructure-and-planning.md}}. Kind: markdown. Size: 4309 bytes. Shape: 105 lines. Markdown document, 105 lines, 0 explicit URLs. Lead heading: Doc Tree Restructure And Planning.\\
\textbf{\path{docs/audit/2026-04-28-doc-tree-restructure-and-planning.md}}. Тип: markdown. Размер: 4309 байт. Форма: 105 строк. Документ Markdown, 105 строк, явных URL: 0. Главный заголовок: Doc Tree Restructure And Planning.\par\medskip
\textbf{\path{docs/audit/2026-04-28-json-runtime-contract.md}}. Kind: markdown. Size: 5281 bytes. Shape: 131 lines. Markdown document, 131 lines, 0 explicit URLs. Lead heading: JSON Runtime Contract.\\
\textbf{\path{docs/audit/2026-04-28-json-runtime-contract.md}}. Тип: markdown. Размер: 5281 байт. Форма: 131 строк. Документ Markdown, 131 строк, явных URL: 0. Главный заголовок: JSON Runtime Contract.\par\medskip
\textbf{\path{docs/audit/2026-04-28-live-azure-dspy-integration-gating.md}}. Kind: markdown. Size: 3562 bytes. Shape: 79 lines. Markdown document, 79 lines, 0 explicit URLs. Lead heading: Live Azure And DSPy Integration Gating.\\
\textbf{\path{docs/audit/2026-04-28-live-azure-dspy-integration-gating.md}}. Тип: markdown. Размер: 3562 байт. Форма: 79 строк. Документ Markdown, 79 строк, явных URL: 0. Главный заголовок: Live Azure And DSPy Integration Gating.\par\medskip
\textbf{\path{docs/audit/2026-04-28-retrieval-profile-layer.md}}. Kind: markdown. Size: 4123 bytes. Shape: 92 lines. Markdown document, 92 lines, 0 explicit URLs. Lead heading: Retrieval Profile Layer.\\
\textbf{\path{docs/audit/2026-04-28-retrieval-profile-layer.md}}. Тип: markdown. Размер: 4123 байт. Форма: 92 строк. Документ Markdown, 92 строк, явных URL: 0. Главный заголовок: Retrieval Profile Layer.\par\medskip
\textbf{\path{docs/audit/2026-04-28-rust-lookup-arbitrary-repo-roots.md}}. Kind: markdown. Size: 4078 bytes. Shape: 88 lines. Markdown document, 88 lines, 0 explicit URLs. Lead heading: Rust Lookup Arbitrary Repo Roots.\\
\textbf{\path{docs/audit/2026-04-28-rust-lookup-arbitrary-repo-roots.md}}. Тип: markdown. Размер: 4078 байт. Форма: 88 строк. Документ Markdown, 88 строк, явных URL: 0. Главный заголовок: Rust Lookup Arbitrary Repo Roots.\par\medskip
\textbf{\path{docs/audit/2026-04-29-bounded-mcp-server-surface.md}}. Kind: markdown. Size: 6688 bytes. Shape: 154 lines. Markdown document, 154 lines, 0 explicit URLs. Lead heading: 2026-04-29 Bounded MCP Server Surface.\\
\textbf{\path{docs/audit/2026-04-29-bounded-mcp-server-surface.md}}. Тип: markdown. Размер: 6688 байт. Форма: 154 строк. Документ Markdown, 154 строк, явных URL: 0. Главный заголовок: 2026-04-29 Bounded MCP Server Surface.\par\medskip
\textbf{\path{docs/audit/2026-04-29-bundle-channel-publish-promote-rollback.md}}. Kind: markdown. Size: 5633 bytes. Shape: 125 lines. Markdown document, 125 lines, 0 explicit URLs. Lead heading: 2026-04-29 Bundle Channel Publish Promote Rollback.\\
\textbf{\path{docs/audit/2026-04-29-bundle-channel-publish-promote-rollback.md}}. Тип: markdown. Размер: 5633 байт. Форма: 125 строк. Документ Markdown, 125 строк, явных URL: 0. Главный заголовок: 2026-04-29 Bundle Channel Publish Promote Rollback.\par\medskip
\textbf{\path{docs/audit/2026-04-29-bundle-overlay-and-runtime-trace-contract.md}}. Kind: markdown. Size: 4332 bytes. Shape: 96 lines. Markdown document, 96 lines, 0 explicit URLs. Lead heading: 2026-04-29 Bundle Overlay And Runtime Trace Contract.\\
\textbf{\path{docs/audit/2026-04-29-bundle-overlay-and-runtime-trace-contract.md}}. Тип: markdown. Размер: 4332 байт. Форма: 96 строк. Документ Markdown, 96 строк, явных URL: 0. Главный заголовок: 2026-04-29 Bundle Overlay And Runtime Trace Contract.\par\medskip
\textbf{\path{docs/audit/2026-04-29-dataset-accepted-outcome-ingest.md}}. Kind: markdown. Size: 6670 bytes. Shape: 139 lines. Markdown document, 139 lines, 0 explicit URLs. Lead heading: 2026-04-29 Dataset Accepted-Outcome Ingest.\\
\textbf{\path{docs/audit/2026-04-29-dataset-accepted-outcome-ingest.md}}. Тип: markdown. Размер: 6670 байт. Форма: 139 строк. Документ Markdown, 139 строк, явных URL: 0. Главный заголовок: 2026-04-29 Dataset Accepted-Outcome Ingest.\par\medskip
\textbf{\path{docs/audit/2026-04-29-dataset-bundle-fetch-and-trace-import-hooks.md}}. Kind: markdown. Size: 6215 bytes. Shape: 130 lines. Markdown document, 130 lines, 0 explicit URLs. Lead heading: 2026-04-29 Dataset Bundle Fetch And Trace Import Hooks.\\
\textbf{\path{docs/audit/2026-04-29-dataset-bundle-fetch-and-trace-import-hooks.md}}. Тип: markdown. Размер: 6215 байт. Форма: 130 строк. Документ Markdown, 130 строк, явных URL: 0. Главный заголовок: 2026-04-29 Dataset Bundle Fetch And Trace Import Hooks.\par\medskip
\textbf{\path{docs/audit/2026-04-29-dataset-local-repo-rag-cli-backend.md}}. Kind: markdown. Size: 5290 bytes. Shape: 122 lines. Markdown document, 122 lines, 0 explicit URLs. Lead heading: 2026-04-29 Dataset Local `repo\_rag\_cli` Backend.\\
\textbf{\path{docs/audit/2026-04-29-dataset-local-repo-rag-cli-backend.md}}. Тип: markdown. Размер: 5290 байт. Форма: 122 строк. Документ Markdown, 122 строк, явных URL: 0. Главный заголовок: 2026-04-29 Dataset Local `repo\_rag\_cli` Backend.\par\medskip
\textbf{\path{docs/audit/2026-04-29-dataset-repo-rag-auto-detection.md}}. Kind: markdown. Size: 6051 bytes. Shape: 124 lines. Markdown document, 124 lines, 0 explicit URLs. Lead heading: 2026-04-29 Dataset Repo-RAG Auto-Detection.\\
\textbf{\path{docs/audit/2026-04-29-dataset-repo-rag-auto-detection.md}}. Тип: markdown. Размер: 6051 байт. Форма: 124 строк. Документ Markdown, 124 строк, явных URL: 0. Главный заголовок: 2026-04-29 Dataset Repo-RAG Auto-Detection.\par\medskip
\textbf{\path{docs/audit/2026-04-29-dataset-trace-queue-handoff.md}}. Kind: markdown. Size: 6656 bytes. Shape: 141 lines. Markdown document, 141 lines, 0 explicit URLs. Lead heading: 2026-04-29 Dataset Trace Queue Handoff.\\
\textbf{\path{docs/audit/2026-04-29-dataset-trace-queue-handoff.md}}. Тип: markdown. Размер: 6656 байт. Форма: 141 строк. Документ Markdown, 141 строк, явных URL: 0. Главный заголовок: 2026-04-29 Dataset Trace Queue Handoff.\par\medskip
\textbf{\path{docs/audit/2026-04-29-dataset-worker-repo-rag-cli-backend.md}}. Kind: markdown. Size: 6630 bytes. Shape: 128 lines. Markdown document, 128 lines, 0 explicit URLs. Lead heading: 2026-04-29 Dataset Worker `repo\_rag\_cli` Backend.\\
\textbf{\path{docs/audit/2026-04-29-dataset-worker-repo-rag-cli-backend.md}}. Тип: markdown. Размер: 6630 байт. Форма: 128 строк. Документ Markdown, 128 строк, явных URL: 0. Главный заголовок: 2026-04-29 Dataset Worker `repo\_rag\_cli` Backend.\par\medskip
\textbf{\path{docs/audit/2026-04-29-global-blob-queue-trainer-migration.md}}. Kind: markdown. Size: 4807 bytes. Shape: 91 lines. Markdown document, 91 lines, 0 explicit URLs. Lead heading: 2026-04-29 Global Blob And Queue Trainer Migration.\\
\textbf{\path{docs/audit/2026-04-29-global-blob-queue-trainer-migration.md}}. Тип: markdown. Размер: 4807 байт. Форма: 91 строк. Документ Markdown, 91 строк, явных URL: 0. Главный заголовок: 2026-04-29 Global Blob And Queue Trainer Migration.\par\medskip
\textbf{\path{docs/audit/2026-04-29-live-azure-gpt54-validation.md}}. Kind: markdown. Size: 8279 bytes. Shape: 145 lines. Markdown document, 145 lines, 6 explicit URLs. Lead heading: 2026-04-29 Live Azure GPT-5.4 Validation.\\
\textbf{\path{docs/audit/2026-04-29-live-azure-gpt54-validation.md}}. Тип: markdown. Размер: 8279 байт. Форма: 145 строк. Документ Markdown, 145 строк, явных URL: 6. Главный заголовок: 2026-04-29 Live Azure GPT-5.4 Validation.\par\medskip
\textbf{\path{docs/audit/2026-04-29-retrieval-rerank-and-relative-root-normalization.md}}. Kind: markdown. Size: 3070 bytes. Shape: 71 lines. Markdown document, 71 lines, 0 explicit URLs. Lead heading: 2026-04-29 Retrieval Rerank And Relative Root Normalization.\\
\textbf{\path{docs/audit/2026-04-29-retrieval-rerank-and-relative-root-normalization.md}}. Тип: markdown. Размер: 3070 байт. Форма: 71 строк. Документ Markdown, 71 строк, явных URL: 0. Главный заголовок: 2026-04-29 Retrieval Rerank And Relative Root Normalization.\par\medskip
\textbf{\path{docs/audit/2026-04-29-runtime-image-and-dataset-build-deploy-wiring.md}}. Kind: markdown. Size: 5729 bytes. Shape: 140 lines. Markdown document, 140 lines, 0 explicit URLs. Lead heading: 2026-04-29 Runtime Image And Dataset Build/Deploy Wiring.\\
\textbf{\path{docs/audit/2026-04-29-runtime-image-and-dataset-build-deploy-wiring.md}}. Тип: markdown. Размер: 5729 байт. Форма: 140 строк. Документ Markdown, 140 строк, явных URL: 0. Главный заголовок: 2026-04-29 Runtime Image And Dataset Build/Deploy Wiring.\par\medskip
\textbf{\path{docs/audit/2026-04-29-trainer-bundle-benchmark-gates.md}}. Kind: markdown. Size: 4942 bytes. Shape: 111 lines. Markdown document, 111 lines, 0 explicit URLs. Lead heading: 2026-04-29 Trainer Bundle Benchmark Gates.\\
\textbf{\path{docs/audit/2026-04-29-trainer-bundle-benchmark-gates.md}}. Тип: markdown. Размер: 4942 байт. Форма: 111 строк. Документ Markdown, 111 строк, явных URL: 0. Главный заголовок: 2026-04-29 Trainer Bundle Benchmark Gates.\par\medskip
\textbf{\path{docs/audit/2026-04-29-trainer-cycle-background-pass.md}}. Kind: markdown. Size: 4637 bytes. Shape: 105 lines. Markdown document, 105 lines, 0 explicit URLs. Lead heading: 2026-04-29 Trainer Cycle Background Pass.\\
\textbf{\path{docs/audit/2026-04-29-trainer-cycle-background-pass.md}}. Тип: markdown. Размер: 4637 байт. Форма: 105 строк. Документ Markdown, 105 строк, явных URL: 0. Главный заголовок: 2026-04-29 Trainer Cycle Background Pass.\par\medskip
\textbf{\path{docs/audit/2026-04-29-trainer-k8s-deployment-surface.md}}. Kind: markdown. Size: 6040 bytes. Shape: 131 lines. Markdown document, 131 lines, 0 explicit URLs. Lead heading: 2026-04-29 Trainer Kubernetes Deployment Surface.\\
\textbf{\path{docs/audit/2026-04-29-trainer-k8s-deployment-surface.md}}. Тип: markdown. Размер: 6040 байт. Форма: 131 строк. Документ Markdown, 131 строк, явных URL: 0. Главный заголовок: 2026-04-29 Trainer Kubernetes Deployment Surface.\par\medskip
\textbf{\path{docs/audit/2026-04-29-trainer-recompile-from-candidates.md}}. Kind: markdown. Size: 6238 bytes. Shape: 135 lines. Markdown document, 135 lines, 0 explicit URLs. Lead heading: 2026-04-29 Trainer Recompile From Candidates.\\
\textbf{\path{docs/audit/2026-04-29-trainer-recompile-from-candidates.md}}. Тип: markdown. Размер: 6238 байт. Форма: 135 строк. Документ Markdown, 135 строк, явных URL: 0. Главный заголовок: 2026-04-29 Trainer Recompile From Candidates.\par\medskip
\textbf{\path{docs/audit/2026-04-29-trainer-service-loop.md}}. Kind: markdown. Size: 6185 bytes. Shape: 133 lines. Markdown document, 133 lines, 0 explicit URLs. Lead heading: 2026-04-29 Trainer Service Loop.\\
\textbf{\path{docs/audit/2026-04-29-trainer-service-loop.md}}. Тип: markdown. Размер: 6185 байт. Форма: 133 строк. Документ Markdown, 133 строк, явных URL: 0. Главный заголовок: 2026-04-29 Trainer Service Loop.\par\medskip
\textbf{\path{docs/audit/2026-04-30-aks-run-25179828865-stale-image-inspection.md}}. Kind: markdown. Size: 6206 bytes. Shape: 117 lines. Markdown document, 117 lines, 0 explicit URLs. Lead heading: 2026-04-30 AKS Run 25179828865 Stale Image Inspection.\\
\textbf{\path{docs/audit/2026-04-30-aks-run-25179828865-stale-image-inspection.md}}. Тип: markdown. Размер: 6206 байт. Форма: 117 строк. Документ Markdown, 117 строк, явных URL: 0. Главный заголовок: 2026-04-30 AKS Run 25179828865 Stale Image Inspection.\par\medskip
\textbf{\path{docs/audit/2026-04-30-aks-run-25186264345-image-fix-gap-inspection.md}}. Kind: markdown. Size: 8407 bytes. Shape: 165 lines. Markdown document, 165 lines, 0 explicit URLs. Lead heading: 2026-04-30 AKS Run 25186264345 Image/Fix Gap Inspection.\\
\textbf{\path{docs/audit/2026-04-30-aks-run-25186264345-image-fix-gap-inspection.md}}. Тип: markdown. Размер: 8407 байт. Форма: 165 строк. Документ Markdown, 165 строк, явных URL: 0. Главный заголовок: 2026-04-30 AKS Run 25186264345 Image/Fix Gap Inspection.\par\medskip
\textbf{\path{docs/audit/2026-04-30-codex-proxy-mediation-layer.md}}. Kind: markdown. Size: 5279 bytes. Shape: 95 lines. Markdown document, 95 lines, 0 explicit URLs. Lead heading: 2026-04-30 Codex Proxy Mediation Layer.\\
\textbf{\path{docs/audit/2026-04-30-codex-proxy-mediation-layer.md}}. Тип: markdown. Размер: 5279 байт. Форма: 95 строк. Документ Markdown, 95 строк, явных URL: 0. Главный заголовок: 2026-04-30 Codex Proxy Mediation Layer.\par\medskip
\textbf{\path{docs/audit/2026-04-30-codex-proxy-trace-handoff.md}}. Kind: markdown. Size: 6543 bytes. Shape: 122 lines. Markdown document, 122 lines, 0 explicit URLs. Lead heading: 2026-04-30 Codex Proxy Trace Handoff.\\
\textbf{\path{docs/audit/2026-04-30-codex-proxy-trace-handoff.md}}. Тип: markdown. Размер: 6543 байт. Форма: 122 строк. Документ Markdown, 122 строк, явных URL: 0. Главный заголовок: 2026-04-30 Codex Proxy Trace Handoff.\par\medskip
\textbf{\path{docs/audit/2026-05-01-aks-run-25191504814-project-root-resolution-inspection.md}}. Kind: markdown. Size: 6002 bytes. Shape: 133 lines. Markdown document, 133 lines, 0 explicit URLs. Lead heading: 2026-05-01 AKS Run 25191504814 Project Root Resolution Inspection.\\
\textbf{\path{docs/audit/2026-05-01-aks-run-25191504814-project-root-resolution-inspection.md}}. Тип: markdown. Размер: 6002 байт. Форма: 133 строк. Документ Markdown, 133 строк, явных URL: 0. Главный заголовок: 2026-05-01 AKS Run 25191504814 Project Root Resolution Inspection.\par\medskip
\textbf{\path{docs/audit/2026-05-01-aks-run-25209573387-proxy-runtime-env-inspection.md}}. Kind: markdown. Size: 5738 bytes. Shape: 129 lines. Markdown document, 129 lines, 0 explicit URLs. Lead heading: 2026-05-01 AKS Run 25209573387 Proxy Runtime Env Inspection.\\
\textbf{\path{docs/audit/2026-05-01-aks-run-25209573387-proxy-runtime-env-inspection.md}}. Тип: markdown. Размер: 5738 байт. Форма: 129 строк. Документ Markdown, 129 строк, явных URL: 0. Главный заголовок: 2026-05-01 AKS Run 25209573387 Proxy Runtime Env Inspection.\par\medskip
\textbf{\path{docs/audit/2026-05-01-aks-run-25212955759-trace-handoff-secret-gap.md}}. Kind: markdown. Size: 4594 bytes. Shape: 121 lines. Markdown document, 121 lines, 0 explicit URLs. Lead heading: AKS Run `25212955759`: Trace Handoff Skipped Because `repo-rag-storage-config` Lacks Blob Credentials.\\
\textbf{\path{docs/audit/2026-05-01-aks-run-25212955759-trace-handoff-secret-gap.md}}. Тип: markdown. Размер: 4594 байт. Форма: 121 строк. Документ Markdown, 121 строк, явных URL: 0. Главный заголовок: AKS Run `25212955759`: Trace Handoff Skipped Because `repo-rag-storage-config` Lacks Blob Credentials.\par\medskip
\textbf{\path{docs/audit/2026-05-01-aks-run-25226558751-rag-heuristic-trusted-handoff.md}}. Kind: markdown. Size: 5092 bytes. Shape: 130 lines. Markdown document, 130 lines, 0 explicit URLs. Lead heading: AKS Run 25226558751: RAG Success, Heuristic DSPy Fallback, Trusted Handoff Success.\\
\textbf{\path{docs/audit/2026-05-01-aks-run-25226558751-rag-heuristic-trusted-handoff.md}}. Тип: markdown. Размер: 5092 байт. Форма: 130 строк. Документ Markdown, 130 строк, явных URL: 0. Главный заголовок: AKS Run 25226558751: RAG Success, Heuristic DSPy Fallback, Trusted Handoff Success.\par\medskip
\textbf{\path{docs/audit/2026-05-01-dataset-azure-responses-api-version-default-fix.md}}. Kind: markdown. Size: 6700 bytes. Shape: 133 lines. Markdown document, 133 lines, 1 explicit URLs. Lead heading: 2026-05-01 Dataset Azure Responses API-Version Default Fix.\\
\textbf{\path{docs/audit/2026-05-01-dataset-azure-responses-api-version-default-fix.md}}. Тип: markdown. Размер: 6700 байт. Форма: 133 строк. Документ Markdown, 133 строк, явных URL: 1. Главный заголовок: 2026-05-01 Dataset Azure Responses API-Version Default Fix.\par\medskip
\textbf{\path{docs/audit/2026-05-01-dataset-azure-runtime-env-propagation-fix.md}}. Kind: markdown. Size: 4200 bytes. Shape: 81 lines. Markdown document, 81 lines, 0 explicit URLs. Lead heading: 2026-05-01 Dataset Azure Runtime Env Propagation Fix.\\
\textbf{\path{docs/audit/2026-05-01-dataset-azure-runtime-env-propagation-fix.md}}. Тип: markdown. Размер: 4200 байт. Форма: 81 строк. Документ Markdown, 81 строк, явных URL: 0. Главный заголовок: 2026-05-01 Dataset Azure Runtime Env Propagation Fix.\par\medskip
\textbf{\path{docs/audit/2026-05-01-dataset-trusted-trace-handoff-postprocess-fix.md}}. Kind: markdown. Size: 3880 bytes. Shape: 81 lines. Markdown document, 81 lines, 0 explicit URLs. Lead heading: Dataset Trusted Trace Handoff Post-Processing Fix.\\
\textbf{\path{docs/audit/2026-05-01-dataset-trusted-trace-handoff-postprocess-fix.md}}. Тип: markdown. Размер: 3880 байт. Форма: 81 строк. Документ Markdown, 81 строк, явных URL: 0. Главный заголовок: Dataset Trusted Trace Handoff Post-Processing Fix.\par\medskip
\textbf{\path{docs/audit/2026-05-01-dspy-bundle-version-pinning-contract.md}}. Kind: markdown. Size: 5763 bytes. Shape: 118 lines. Markdown document, 118 lines, 0 explicit URLs. Lead heading: DSPY Bundle Version Pinning Contract.\\
\textbf{\path{docs/audit/2026-05-01-dspy-bundle-version-pinning-contract.md}}. Тип: markdown. Размер: 5763 байт. Форма: 118 строк. Документ Markdown, 118 строк, явных URL: 0. Главный заголовок: DSPY Bundle Version Pinning Contract.\par\medskip
\textbf{\path{docs/audit/2026-05-01-durable-trainer-progress-remote-recovery.md}}. Kind: markdown. Size: 5141 bytes. Shape: 115 lines. Markdown document, 115 lines, 0 explicit URLs. Lead heading: Durable Trainer Progress Remote Recovery.\\
\textbf{\path{docs/audit/2026-05-01-durable-trainer-progress-remote-recovery.md}}. Тип: markdown. Размер: 5141 байт. Форма: 115 строк. Документ Markdown, 115 строк, явных URL: 0. Главный заголовок: Durable Trainer Progress Remote Recovery.\par\medskip
\textbf{\path{docs/audit/2026-05-01-live-trainer-no-timestamp-bundle-publish.md}}. Kind: markdown. Size: 7368 bytes. Shape: 172 lines. Markdown document, 172 lines, 0 explicit URLs. Lead heading: Live Trainer No Timestamp Bundle Publish.\\
\textbf{\path{docs/audit/2026-05-01-live-trainer-no-timestamp-bundle-publish.md}}. Тип: markdown. Размер: 7368 байт. Форма: 172 строк. Документ Markdown, 172 строк, явных URL: 0. Главный заголовок: Live Trainer No Timestamp Bundle Publish.\par\medskip
\textbf{\path{docs/audit/2026-05-01-live-trainer-remote-bundle-publish.md}}. Kind: markdown. Size: 7444 bytes. Shape: 148 lines. Markdown document, 148 lines, 0 explicit URLs. Lead heading: Live Trainer Remote Bundle Publish.\\
\textbf{\path{docs/audit/2026-05-01-live-trainer-remote-bundle-publish.md}}. Тип: markdown. Размер: 7444 байт. Форма: 148 строк. Документ Markdown, 148 строк, явных URL: 0. Главный заголовок: Live Trainer Remote Bundle Publish.\par\medskip
\textbf{\path{docs/audit/2026-05-01-live-trainer-service-azure-fallback.md}}. Kind: markdown. Size: 6390 bytes. Shape: 167 lines. Markdown document, 167 lines, 0 explicit URLs. Lead heading: Live Trainer Service Azure Fallback.\\
\textbf{\path{docs/audit/2026-05-01-live-trainer-service-azure-fallback.md}}. Тип: markdown. Размер: 6390 байт. Форма: 167 строк. Документ Markdown, 167 строк, явных URL: 0. Главный заголовок: Live Trainer Service Azure Fallback.\par\medskip
\textbf{\path{docs/audit/2026-05-01-versioned-trainer-lineage-bundles.md}}. Kind: markdown. Size: 5821 bytes. Shape: 123 lines. Markdown document, 123 lines, 0 explicit URLs. Lead heading: Versioned Trainer Lineage Bundles.\\
\textbf{\path{docs/audit/2026-05-01-versioned-trainer-lineage-bundles.md}}. Тип: markdown. Размер: 5821 байт. Форма: 123 строк. Документ Markdown, 123 строк, явных URL: 0. Главный заголовок: Versioned Trainer Lineage Bundles.\par\medskip
\textbf{\path{docs/audit/2026-05-02-hybrid-vector-retrieval-default.md}}. Kind: markdown. Size: 18302 bytes. Shape: 385 lines. Markdown document, 385 lines, 0 explicit URLs. Lead heading: Hybrid Vector Retrieval Default.\\
\textbf{\path{docs/audit/2026-05-02-hybrid-vector-retrieval-default.md}}. Тип: markdown. Размер: 18302 байт. Форма: 385 строк. Документ Markdown, 385 строк, явных URL: 0. Главный заголовок: Hybrid Vector Retrieval Default.\par\medskip
\textbf{\path{docs/audit/2026-05-02-live-trainer-still-not-publishing-bundles.md}}. Kind: markdown. Size: 60556 bytes. Shape: 1281 lines. Markdown document, 1281 lines, 0 explicit URLs. Lead heading: Live Trainer Still Not Publishing Bundles.\\
\textbf{\path{docs/audit/2026-05-02-live-trainer-still-not-publishing-bundles.md}}. Тип: markdown. Размер: 60556 байт. Форма: 1281 строк. Документ Markdown, 1281 строк, явных URL: 0. Главный заголовок: Live Trainer Still Not Publishing Bundles.\par\medskip
\textbf{\path{docs/audit/2026-05-02-trainer-context-group-champion-index-stage1.md}}. Kind: markdown. Size: 18201 bytes. Shape: 330 lines. Markdown document, 330 lines, 0 explicit URLs. Lead heading: Trainer Context-Group Champion Index Stage 1.\\
\textbf{\path{docs/audit/2026-05-02-trainer-context-group-champion-index-stage1.md}}. Тип: markdown. Размер: 18201 байт. Форма: 330 строк. Документ Markdown, 330 строк, явных URL: 0. Главный заголовок: Trainer Context-Group Champion Index Stage 1.\par\medskip
\textbf{\path{docs/audit/2026-05-02-zz-codex-exec-resume-session-state-stage0.md}}. Kind: markdown. Size: 170338 bytes. Shape: 3866 lines. Markdown document, 3866 lines, 2 explicit URLs. Lead heading: Repository audit note for 2026-05-02 codex exec resume session state stage 0.\\
\textbf{\path{docs/audit/2026-05-02-zz-codex-exec-resume-session-state-stage0.md}}. Тип: markdown. Размер: 170338 байт. Форма: 3866 строк. Документ Markdown, 3866 строк, явных URL: 2. Главный заголовок: Repository audit note for 2026-05-02 codex exec resume session state stage 0.\par\medskip
\textbf{\path{docs/audit/2026-05-08-zz-per-turn-dspy-mediation-contract-stage0.md}}. Kind: markdown. Size: 5821 bytes. Shape: 113 lines. Markdown document, 113 lines, 0 explicit URLs. Lead heading: Repository audit note for 2026-05-08 per-turn DSPy mediation contract stage 0.\\
\textbf{\path{docs/audit/2026-05-08-zz-per-turn-dspy-mediation-contract-stage0.md}}. Тип: markdown. Размер: 5821 байт. Форма: 113 строк. Документ Markdown, 113 строк, явных URL: 0. Главный заголовок: Repository audit note for 2026-05-08 per-turn DSPy mediation contract stage 0.\par\medskip
\textbf{\path{docs/audit/2026-05-09-aks-run-25598675829-per-turn-mediation-gap.md}}. Kind: markdown. Size: 7119 bytes. Shape: 215 lines. Markdown document, 215 lines, 0 explicit URLs. Lead heading: Repository audit note for 2026-05-09 AKS run 25598675829 per-turn mediation gap.\\
\textbf{\path{docs/audit/2026-05-09-aks-run-25598675829-per-turn-mediation-gap.md}}. Тип: markdown. Размер: 7119 байт. Форма: 215 строк. Документ Markdown, 215 строк, явных URL: 0. Главный заголовок: Repository audit note for 2026-05-09 AKS run 25598675829 per-turn mediation gap.\par\medskip
\textbf{\path{docs/audit/2026-05-09-family-first-mipro-contract-stage1.md}}. Kind: markdown. Size: 4684 bytes. Shape: 120 lines. Markdown document, 120 lines, 0 explicit URLs. Lead heading: Repository audit note for 2026-05-09 family-first MIPROv2 contract stage 1.\\
\textbf{\path{docs/audit/2026-05-09-family-first-mipro-contract-stage1.md}}. Тип: markdown. Размер: 4684 байт. Форма: 120 строк. Документ Markdown, 120 строк, явных URL: 0. Главный заголовок: Repository audit note for 2026-05-09 family-first MIPROv2 contract stage 1.\par\medskip
\textbf{\path{docs/audit/2026-05-09-family-first-mipro-contract-stage10.md}}. Kind: markdown. Size: 5329 bytes. Shape: 122 lines. Markdown document, 122 lines, 0 explicit URLs. Lead heading: Repository audit note for 2026-05-09 family-first MIPROv2 contract stage 10.\\
\textbf{\path{docs/audit/2026-05-09-family-first-mipro-contract-stage10.md}}. Тип: markdown. Размер: 5329 байт. Форма: 122 строк. Документ Markdown, 122 строк, явных URL: 0. Главный заголовок: Repository audit note for 2026-05-09 family-first MIPROv2 contract stage 10.\par\medskip
\textbf{\path{docs/audit/2026-05-09-family-first-mipro-contract-stage11.md}}. Kind: markdown. Size: 4638 bytes. Shape: 116 lines. Markdown document, 116 lines, 0 explicit URLs. Lead heading: Repository audit note for 2026-05-09 family-first MIPROv2 contract stage 11.\\
\textbf{\path{docs/audit/2026-05-09-family-first-mipro-contract-stage11.md}}. Тип: markdown. Размер: 4638 байт. Форма: 116 строк. Документ Markdown, 116 строк, явных URL: 0. Главный заголовок: Repository audit note for 2026-05-09 family-first MIPROv2 contract stage 11.\par\medskip
\textbf{\path{docs/audit/2026-05-09-family-first-mipro-contract-stage12.md}}. Kind: markdown. Size: 4941 bytes. Shape: 122 lines. Markdown document, 122 lines, 0 explicit URLs. Lead heading: Repository audit note for 2026-05-09 family-first MIPROv2 contract stage 12.\\
\textbf{\path{docs/audit/2026-05-09-family-first-mipro-contract-stage12.md}}. Тип: markdown. Размер: 4941 байт. Форма: 122 строк. Документ Markdown, 122 строк, явных URL: 0. Главный заголовок: Repository audit note for 2026-05-09 family-first MIPROv2 contract stage 12.\par\medskip
\textbf{\path{docs/audit/2026-05-09-family-first-mipro-contract-stage13.md}}. Kind: markdown. Size: 4992 bytes. Shape: 128 lines. Markdown document, 128 lines, 0 explicit URLs. Lead heading: Repository audit note for 2026-05-09 family-first MIPROv2 contract stage 13.\\
\textbf{\path{docs/audit/2026-05-09-family-first-mipro-contract-stage13.md}}. Тип: markdown. Размер: 4992 байт. Форма: 128 строк. Документ Markdown, 128 строк, явных URL: 0. Главный заголовок: Repository audit note for 2026-05-09 family-first MIPROv2 contract stage 13.\par\medskip
\textbf{\path{docs/audit/2026-05-09-family-first-mipro-contract-stage14.md}}. Kind: markdown. Size: 3944 bytes. Shape: 109 lines. Markdown document, 109 lines, 0 explicit URLs. Lead heading: Repository audit note for 2026-05-09 family-first MIPROv2 contract stage 14.\\
\textbf{\path{docs/audit/2026-05-09-family-first-mipro-contract-stage14.md}}. Тип: markdown. Размер: 3944 байт. Форма: 109 строк. Документ Markdown, 109 строк, явных URL: 0. Главный заголовок: Repository audit note for 2026-05-09 family-first MIPROv2 contract stage 14.\par\medskip
\textbf{\path{docs/audit/2026-05-09-family-first-mipro-contract-stage15.md}}. Kind: markdown. Size: 4728 bytes. Shape: 113 lines. Markdown document, 113 lines, 0 explicit URLs. Lead heading: Repository audit note for 2026-05-09 family-first MIPROv2 contract stage 15.\\
\textbf{\path{docs/audit/2026-05-09-family-first-mipro-contract-stage15.md}}. Тип: markdown. Размер: 4728 байт. Форма: 113 строк. Документ Markdown, 113 строк, явных URL: 0. Главный заголовок: Repository audit note for 2026-05-09 family-first MIPROv2 contract stage 15.\par\medskip
\textbf{\path{docs/audit/2026-05-09-family-first-mipro-contract-stage16.md}}. Kind: markdown. Size: 3748 bytes. Shape: 103 lines. Markdown document, 103 lines, 0 explicit URLs. Lead heading: Repository audit note for 2026-05-09 family-first MIPROv2 contract stage 16.\\
\textbf{\path{docs/audit/2026-05-09-family-first-mipro-contract-stage16.md}}. Тип: markdown. Размер: 3748 байт. Форма: 103 строк. Документ Markdown, 103 строк, явных URL: 0. Главный заголовок: Repository audit note for 2026-05-09 family-first MIPROv2 contract stage 16.\par\medskip
\textbf{\path{docs/audit/2026-05-09-family-first-mipro-contract-stage17.md}}. Kind: markdown. Size: 3519 bytes. Shape: 101 lines. Markdown document, 101 lines, 0 explicit URLs. Lead heading: Repository audit note for 2026-05-09 family-first MIPROv2 contract stage 17.\\
\textbf{\path{docs/audit/2026-05-09-family-first-mipro-contract-stage17.md}}. Тип: markdown. Размер: 3519 байт. Форма: 101 строк. Документ Markdown, 101 строк, явных URL: 0. Главный заголовок: Repository audit note for 2026-05-09 family-first MIPROv2 contract stage 17.\par\medskip
\textbf{\path{docs/audit/2026-05-09-family-first-mipro-contract-stage18.md}}. Kind: markdown. Size: 3164 bytes. Shape: 90 lines. Markdown document, 90 lines, 0 explicit URLs. Lead heading: Repository audit note for 2026-05-09 family-first MIPROv2 contract stage 18.\\
\textbf{\path{docs/audit/2026-05-09-family-first-mipro-contract-stage18.md}}. Тип: markdown. Размер: 3164 байт. Форма: 90 строк. Документ Markdown, 90 строк, явных URL: 0. Главный заголовок: Repository audit note for 2026-05-09 family-first MIPROv2 contract stage 18.\par\medskip
\textbf{\path{docs/audit/2026-05-09-family-first-mipro-contract-stage19.md}}. Kind: markdown. Size: 3984 bytes. Shape: 102 lines. Markdown document, 102 lines, 0 explicit URLs. Lead heading: Repository audit note for 2026-05-09 family-first MIPROv2 contract stage 19.\\
\textbf{\path{docs/audit/2026-05-09-family-first-mipro-contract-stage19.md}}. Тип: markdown. Размер: 3984 байт. Форма: 102 строк. Документ Markdown, 102 строк, явных URL: 0. Главный заголовок: Repository audit note for 2026-05-09 family-first MIPROv2 contract stage 19.\par\medskip
\textbf{\path{docs/audit/2026-05-09-family-first-mipro-contract-stage2.md}}. Kind: markdown. Size: 4936 bytes. Shape: 123 lines. Markdown document, 123 lines, 0 explicit URLs. Lead heading: Repository audit note for 2026-05-09 family-first MIPROv2 contract stage 2.\\
\textbf{\path{docs/audit/2026-05-09-family-first-mipro-contract-stage2.md}}. Тип: markdown. Размер: 4936 байт. Форма: 123 строк. Документ Markdown, 123 строк, явных URL: 0. Главный заголовок: Repository audit note for 2026-05-09 family-first MIPROv2 contract stage 2.\par\medskip
\textbf{\path{docs/audit/2026-05-09-family-first-mipro-contract-stage20.md}}. Kind: markdown. Size: 3076 bytes. Shape: 88 lines. Markdown document, 88 lines, 0 explicit URLs. Lead heading: Repository audit note for 2026-05-09 family-first MIPROv2 contract stage 20.\\
\textbf{\path{docs/audit/2026-05-09-family-first-mipro-contract-stage20.md}}. Тип: markdown. Размер: 3076 байт. Форма: 88 строк. Документ Markdown, 88 строк, явных URL: 0. Главный заголовок: Repository audit note for 2026-05-09 family-first MIPROv2 contract stage 20.\par\medskip
\textbf{\path{docs/audit/2026-05-09-family-first-mipro-contract-stage21.md}}. Kind: markdown. Size: 5742 bytes. Shape: 122 lines. Markdown document, 122 lines, 0 explicit URLs. Lead heading: Repository audit note for 2026-05-09 family-first MIPROv2 contract stage 21.\\
\textbf{\path{docs/audit/2026-05-09-family-first-mipro-contract-stage21.md}}. Тип: markdown. Размер: 5742 байт. Форма: 122 строк. Документ Markdown, 122 строк, явных URL: 0. Главный заголовок: Repository audit note for 2026-05-09 family-first MIPROv2 contract stage 21.\par\medskip
\textbf{\path{docs/audit/2026-05-09-family-first-mipro-contract-stage3.md}}. Kind: markdown. Size: 4147 bytes. Shape: 96 lines. Markdown document, 96 lines, 0 explicit URLs. Lead heading: Repository audit note for 2026-05-09 family-first MIPROv2 contract stage 3.\\
\textbf{\path{docs/audit/2026-05-09-family-first-mipro-contract-stage3.md}}. Тип: markdown. Размер: 4147 байт. Форма: 96 строк. Документ Markdown, 96 строк, явных URL: 0. Главный заголовок: Repository audit note for 2026-05-09 family-first MIPROv2 contract stage 3.\par\medskip
\textbf{\path{docs/audit/2026-05-09-family-first-mipro-contract-stage4.md}}. Kind: markdown. Size: 5209 bytes. Shape: 120 lines. Markdown document, 120 lines, 0 explicit URLs. Lead heading: Repository audit note for 2026-05-09 family-first MIPROv2 contract stage 4.\\
\textbf{\path{docs/audit/2026-05-09-family-first-mipro-contract-stage4.md}}. Тип: markdown. Размер: 5209 байт. Форма: 120 строк. Документ Markdown, 120 строк, явных URL: 0. Главный заголовок: Repository audit note for 2026-05-09 family-first MIPROv2 contract stage 4.\par\medskip
\textbf{\path{docs/audit/2026-05-09-family-first-mipro-contract-stage5.md}}. Kind: markdown. Size: 5172 bytes. Shape: 122 lines. Markdown document, 122 lines, 0 explicit URLs. Lead heading: Repository audit note for 2026-05-09 family-first MIPROv2 contract stage 5.\\
\textbf{\path{docs/audit/2026-05-09-family-first-mipro-contract-stage5.md}}. Тип: markdown. Размер: 5172 байт. Форма: 122 строк. Документ Markdown, 122 строк, явных URL: 0. Главный заголовок: Repository audit note for 2026-05-09 family-first MIPROv2 contract stage 5.\par\medskip
\textbf{\path{docs/audit/2026-05-09-family-first-mipro-contract-stage6.md}}. Kind: markdown. Size: 5120 bytes. Shape: 119 lines. Markdown document, 119 lines, 0 explicit URLs. Lead heading: Repository audit note for 2026-05-09 family-first MIPROv2 contract stage 6.\\
\textbf{\path{docs/audit/2026-05-09-family-first-mipro-contract-stage6.md}}. Тип: markdown. Размер: 5120 байт. Форма: 119 строк. Документ Markdown, 119 строк, явных URL: 0. Главный заголовок: Repository audit note for 2026-05-09 family-first MIPROv2 contract stage 6.\par\medskip
\textbf{\path{docs/audit/2026-05-09-family-first-mipro-contract-stage7.md}}. Kind: markdown. Size: 4680 bytes. Shape: 116 lines. Markdown document, 116 lines, 0 explicit URLs. Lead heading: Repository audit note for 2026-05-09 family-first MIPROv2 contract stage 7.\\
\textbf{\path{docs/audit/2026-05-09-family-first-mipro-contract-stage7.md}}. Тип: markdown. Размер: 4680 байт. Форма: 116 строк. Документ Markdown, 116 строк, явных URL: 0. Главный заголовок: Repository audit note for 2026-05-09 family-first MIPROv2 contract stage 7.\par\medskip
\textbf{\path{docs/audit/2026-05-09-family-first-mipro-contract-stage8.md}}. Kind: markdown. Size: 4477 bytes. Shape: 110 lines. Markdown document, 110 lines, 0 explicit URLs. Lead heading: Repository audit note for 2026-05-09 family-first MIPROv2 contract stage 8.\\
\textbf{\path{docs/audit/2026-05-09-family-first-mipro-contract-stage8.md}}. Тип: markdown. Размер: 4477 байт. Форма: 110 строк. Документ Markdown, 110 строк, явных URL: 0. Главный заголовок: Repository audit note for 2026-05-09 family-first MIPROv2 contract stage 8.\par\medskip
\textbf{\path{docs/audit/2026-05-09-family-first-mipro-contract-stage9.md}}. Kind: markdown. Size: 4003 bytes. Shape: 106 lines. Markdown document, 106 lines, 0 explicit URLs. Lead heading: Repository audit note for 2026-05-09 family-first MIPROv2 contract stage 9.\\
\textbf{\path{docs/audit/2026-05-09-family-first-mipro-contract-stage9.md}}. Тип: markdown. Размер: 4003 байт. Форма: 106 строк. Документ Markdown, 106 строк, явных URL: 0. Главный заголовок: Repository audit note for 2026-05-09 family-first MIPROv2 contract stage 9.\par\medskip
\textbf{\path{docs/audit/2026-05-10-aks-run-25625638264-family-first-runtime-gap.md}}. Kind: markdown. Size: 8472 bytes. Shape: 246 lines. Markdown document, 246 lines, 0 explicit URLs. Lead heading: Repository audit note for 2026-05-10 AKS run 25625638264 family-first runtime gap.\\
\textbf{\path{docs/audit/2026-05-10-aks-run-25625638264-family-first-runtime-gap.md}}. Тип: markdown. Размер: 8472 байт. Форма: 246 строк. Документ Markdown, 246 строк, явных URL: 0. Главный заголовок: Repository audit note for 2026-05-10 AKS run 25625638264 family-first runtime gap.\par\medskip
\textbf{\path{docs/audit/2026-05-10-aks-run-25629990035-family-first-runtime-still-heuristic.md}}. Kind: markdown. Size: 6843 bytes. Shape: 211 lines. Markdown document, 211 lines, 0 explicit URLs. Lead heading: Repository audit note for 2026-05-10 AKS run 25629990035 family-first runtime still heuristic.\\
\textbf{\path{docs/audit/2026-05-10-aks-run-25629990035-family-first-runtime-still-heuristic.md}}. Тип: markdown. Размер: 6843 байт. Форма: 211 строк. Документ Markdown, 211 строк, явных URL: 0. Главный заголовок: Repository audit note for 2026-05-10 AKS run 25629990035 family-first runtime still heuristic.\par\medskip
\textbf{\path{docs/audit/2026-05-10-aks-run-25632110510-handoff-fixed-runtime-still-heuristic.md}}. Kind: markdown. Size: 7832 bytes. Shape: 231 lines. Markdown document, 231 lines, 1 explicit URLs. Lead heading: Repository audit note for 2026-05-10 AKS run 25632110510 handoff fixed but runtime still heuristic.\\
\textbf{\path{docs/audit/2026-05-10-aks-run-25632110510-handoff-fixed-runtime-still-heuristic.md}}. Тип: markdown. Размер: 7832 байт. Форма: 231 строк. Документ Markdown, 231 строк, явных URL: 1. Главный заголовок: Repository audit note for 2026-05-10 AKS run 25632110510 handoff fixed but runtime still heuristic.\par\medskip
\textbf{\path{docs/audit/2026-05-10-aks-run-25634202380-reset-lane-batch-handoff-runtime-still-heuristic.md}}. Kind: markdown. Size: 6096 bytes. Shape: 197 lines. Markdown document, 197 lines, 0 explicit URLs. Lead heading: Repository audit note for AKS run 25634202380 on 2026-05-10.\\
\textbf{\path{docs/audit/2026-05-10-aks-run-25634202380-reset-lane-batch-handoff-runtime-still-heuristic.md}}. Тип: markdown. Размер: 6096 байт. Форма: 197 строк. Документ Markdown, 197 строк, явных URL: 0. Главный заголовок: Repository audit note for AKS run 25634202380 on 2026-05-10.\par\medskip
\textbf{\path{docs/audit/2026-05-10-family-first-runtime-fixes-after-run-25625638264.md}}. Kind: markdown. Size: 12143 bytes. Shape: 261 lines. Markdown document, 261 lines, 0 explicit URLs. Lead heading: Repository audit note for 2026-05-10 family-first runtime fixes after run 25625638264.\\
\textbf{\path{docs/audit/2026-05-10-family-first-runtime-fixes-after-run-25625638264.md}}. Тип: markdown. Размер: 12143 байт. Форма: 261 строк. Документ Markdown, 261 строк, явных URL: 0. Главный заголовок: Repository audit note for 2026-05-10 family-first runtime fixes after run 25625638264.\par\medskip
\textbf{\path{docs/audit/2026-05-10-family-first-runtime-fixes-after-run-25629990035.md}}. Kind: markdown. Size: 4170 bytes. Shape: 98 lines. Markdown document, 98 lines, 0 explicit URLs. Lead heading: Repository audit note for 2026-05-10 family-first runtime fixes after AKS run 25629990035.\\
\textbf{\path{docs/audit/2026-05-10-family-first-runtime-fixes-after-run-25629990035.md}}. Тип: markdown. Размер: 4170 байт. Форма: 98 строк. Документ Markdown, 98 строк, явных URL: 0. Главный заголовок: Repository audit note for 2026-05-10 family-first runtime fixes after AKS run 25629990035.\par\medskip
\textbf{\path{docs/audit/2026-05-10-family-first-runtime-fixes-after-run-25632110510.md}}. Kind: markdown. Size: 5332 bytes. Shape: 112 lines. Markdown document, 112 lines, 0 explicit URLs. Lead heading: Repository audit note for 2026-05-10 family-first runtime fixes after AKS run 25632110510.\\
\textbf{\path{docs/audit/2026-05-10-family-first-runtime-fixes-after-run-25632110510.md}}. Тип: markdown. Размер: 5332 байт. Форма: 112 строк. Документ Markdown, 112 строк, явных URL: 0. Главный заголовок: Repository audit note for 2026-05-10 family-first runtime fixes after AKS run 25632110510.\par\medskip
\textbf{\path{docs/audit/2026-05-10-family-state-populates-but-publish-gate-and-dirty-fathers-still-block-runtime.md}}. Kind: markdown. Size: 7106 bytes. Shape: 183 lines. Markdown document, 183 lines, 0 explicit URLs. Lead heading: 2026-05-10 Family State Populates But Publish Gate And Dirty Fathers Still Block Runtime.\\
\textbf{\path{docs/audit/2026-05-10-family-state-populates-but-publish-gate-and-dirty-fathers-still-block-runtime.md}}. Тип: markdown. Размер: 7106 байт. Форма: 183 строк. Документ Markdown, 183 строк, явных URL: 0. Главный заголовок: 2026-05-10 Family State Populates But Publish Gate And Dirty Fathers Still Block Runtime.\par\medskip
\textbf{\path{docs/audit/2026-05-10-family-state-root-alias-and-latest-bundle-staging-fallback.md}}. Kind: markdown. Size: 3390 bytes. Shape: 81 lines. Markdown document, 81 lines, 0 explicit URLs. Lead heading: 2026-05-10 Family-State Root Alias And Latest-Bundle Staging Fallback.\\
\textbf{\path{docs/audit/2026-05-10-family-state-root-alias-and-latest-bundle-staging-fallback.md}}. Тип: markdown. Размер: 3390 байт. Форма: 81 строк. Документ Markdown, 81 строк, явных URL: 0. Главный заголовок: 2026-05-10 Family-State Root Alias And Latest-Bundle Staging Fallback.\par\medskip
\textbf{\path{docs/audit/2026-05-10-hidden-token-spend-and-incremental-trainer-fixes.md}}. Kind: markdown. Size: 6740 bytes. Shape: 154 lines. Markdown document, 154 lines, 0 explicit URLs. Lead heading: 2026-05-10 Hidden Token Spend And Incremental Trainer Fixes.\\
\textbf{\path{docs/audit/2026-05-10-hidden-token-spend-and-incremental-trainer-fixes.md}}. Тип: markdown. Размер: 6740 байт. Форма: 154 строк. Документ Markdown, 154 строк, явных URL: 0. Главный заголовок: 2026-05-10 Hidden Token Spend And Incremental Trainer Fixes.\par\medskip
\textbf{\path{docs/audit/2026-05-10-live-trainer-service-replays-processed-ledger-while-queue-empty.md}}. Kind: markdown. Size: 5891 bytes. Shape: 160 lines. Markdown document, 160 lines, 0 explicit URLs. Lead heading: 2026-05-10 Live Trainer Service Replays Processed Ledger While Queue Is Empty.\\
\textbf{\path{docs/audit/2026-05-10-live-trainer-service-replays-processed-ledger-while-queue-empty.md}}. Тип: markdown. Размер: 5891 байт. Форма: 160 строк. Документ Markdown, 160 строк, явных URL: 0. Главный заголовок: 2026-05-10 Live Trainer Service Replays Processed Ledger While Queue Is Empty.\par\medskip
\textbf{\path{docs/audit/2026-05-10-stale-submodule-runtime-image-root-cause-and-source-preference-fix.md}}. Kind: markdown. Size: 5021 bytes. Shape: 114 lines. Markdown document, 114 lines, 0 explicit URLs. Lead heading: 2026-05-10 Stale Submodule Runtime Image Root Cause And Source-Preference Fix.\\
\textbf{\path{docs/audit/2026-05-10-stale-submodule-runtime-image-root-cause-and-source-preference-fix.md}}. Тип: markdown. Размер: 5021 байт. Форма: 114 строк. Документ Markdown, 114 строк, явных URL: 0. Главный заголовок: 2026-05-10 Stale Submodule Runtime Image Root Cause And Source-Preference Fix.\par\medskip
\textbf{\path{docs/audit/2026-05-11-family-runtime-artifact-fallback-and-staging-path-fix.md}}. Kind: markdown. Size: 5601 bytes. Shape: 139 lines. Markdown document, 139 lines, 0 explicit URLs. Lead heading: 2026-05-11 Family Runtime Artifact Fallback And Bundle Staging Path Fix.\\
\textbf{\path{docs/audit/2026-05-11-family-runtime-artifact-fallback-and-staging-path-fix.md}}. Тип: markdown. Размер: 5601 байт. Форма: 139 строк. Документ Markdown, 139 строк, явных URL: 0. Главный заголовок: 2026-05-11 Family Runtime Artifact Fallback And Bundle Staging Path Fix.\par\medskip
\textbf{\path{docs/audit/2026-05-11-fresh-run-family-match-without-family-artifact-runtime.md}}. Kind: markdown. Size: 6472 bytes. Shape: 211 lines. Markdown document, 211 lines, 0 explicit URLs. Lead heading: 2026-05-11 Fresh Run Family Match Without Family Artifact Runtime.\\
\textbf{\path{docs/audit/2026-05-11-fresh-run-family-match-without-family-artifact-runtime.md}}. Тип: markdown. Размер: 6472 байт. Форма: 211 строк. Документ Markdown, 211 строк, явных URL: 0. Главный заголовок: 2026-05-11 Fresh Run Family Match Without Family Artifact Runtime.\par\medskip
\textbf{\path{docs/audit/2026-05-11-idle-trainer-cycles-no-longer-autorepublish-bundles.md}}. Kind: markdown. Size: 3162 bytes. Shape: 108 lines. Markdown document, 108 lines, 0 explicit URLs. Lead heading: 2026-05-11 Idle Trainer Cycles No Longer Auto-Republish Bundle Versions.\\
\textbf{\path{docs/audit/2026-05-11-idle-trainer-cycles-no-longer-autorepublish-bundles.md}}. Тип: markdown. Размер: 3162 байт. Форма: 108 строк. Документ Markdown, 108 строк, явных URL: 0. Главный заголовок: 2026-05-11 Idle Trainer Cycles No Longer Auto-Republish Bundle Versions.\par\medskip
\textbf{\path{docs/audit/2026-05-11-original-prompt-routing-and-batch-handoff-fixes.md}}. Kind: markdown. Size: 3829 bytes. Shape: 91 lines. Markdown document, 91 lines, 0 explicit URLs. Lead heading: 2026-05-11 Original Prompt Routing And Batch Handoff Fixes.\\
\textbf{\path{docs/audit/2026-05-11-original-prompt-routing-and-batch-handoff-fixes.md}}. Тип: markdown. Размер: 3829 байт. Форма: 91 строк. Документ Markdown, 91 строк, явных URL: 0. Главный заголовок: 2026-05-11 Original Prompt Routing And Batch Handoff Fixes.\par\medskip
\textbf{\path{docs/audit/2026-05-11-trainer-queue-only-no-processed-replay.md}}. Kind: markdown. Size: 6373 bytes. Shape: 166 lines. Markdown document, 166 lines, 0 explicit URLs. Lead heading: 2026-05-11 Trainer Is Now Queue-Only And Will Not Replay Processed History.\\
\textbf{\path{docs/audit/2026-05-11-trainer-queue-only-no-processed-replay.md}}. Тип: markdown. Размер: 6373 байт. Форма: 166 строк. Документ Markdown, 166 строк, явных URL: 0. Главный заголовок: 2026-05-11 Trainer Is Now Queue-Only And Will Not Replay Processed History.\par\medskip
\textbf{\path{docs/audit/2026-05-11-trainer-service-preflights-queue-before-running-cycles.md}}. Kind: markdown. Size: 4211 bytes. Shape: 121 lines. Markdown document, 121 lines, 0 explicit URLs. Lead heading: 2026-05-11 Trainer Service Preflights Queue Before Running Cycles.\\
\textbf{\path{docs/audit/2026-05-11-trainer-service-preflights-queue-before-running-cycles.md}}. Тип: markdown. Размер: 4211 байт. Форма: 121 строк. Документ Markdown, 121 строк, явных URL: 0. Главный заголовок: 2026-05-11 Trainer Service Preflights Queue Before Running Cycles.\par\medskip
\textbf{\path{docs/audit/2026-05-11-trusted-handoff-syntax-error-blocked-execution-artifact-upload.md}}. Kind: markdown. Size: 4031 bytes. Shape: 86 lines. Markdown document, 86 lines, 0 explicit URLs. Lead heading: 2026-05-11 Trusted Handoff Syntax Error Blocked Execution Artifact Upload.\\
\textbf{\path{docs/audit/2026-05-11-trusted-handoff-syntax-error-blocked-execution-artifact-upload.md}}. Тип: markdown. Размер: 4031 байт. Форма: 86 строк. Документ Markdown, 86 строк, явных URL: 0. Главный заголовок: 2026-05-11 Trusted Handoff Syntax Error Blocked Execution Artifact Upload.\par\medskip
\textbf{\path{docs/audit/README.md}}. Kind: markdown. Size: 443 bytes. Shape: 14 lines. Markdown document, 14 lines, 0 explicit URLs. Lead heading: Repository Audit Notes.\\
\textbf{\path{docs/audit/README.md}}. Тип: markdown. Размер: 443 байт. Форма: 14 строк. Документ Markdown, 14 строк, явных URL: 0. Главный заголовок: Repository Audit Notes.\par\medskip
\textbf{\path{docs/guides/hushwheel-fixture-rag-guide.md}}. Kind: markdown. Size: 5169 bytes. Shape: 138 lines. Markdown document, 138 lines, 0 explicit URLs. Lead heading: RAG Guide For The Hushwheel Fixture.\\
\textbf{\path{docs/guides/hushwheel-fixture-rag-guide.md}}. Тип: markdown. Размер: 5169 байт. Форма: 138 строк. Документ Markdown, 138 строк, явных URL: 0. Главный заголовок: RAG Guide For The Hushwheel Fixture.\par\medskip
\textbf{\path{docs/operations/azure-deployment.md}}. Kind: markdown. Size: 4901 bytes. Shape: 111 lines. Markdown document, 111 lines, 5 explicit URLs. Lead heading: Azure Deployment Notes.\\
\textbf{\path{docs/operations/azure-deployment.md}}. Тип: markdown. Размер: 4901 байт. Форма: 111 строк. Документ Markdown, 111 строк, явных URL: 5. Главный заголовок: Azure Deployment Notes.\par\medskip
\textbf{\path{docs/operations/environment.md}}. Kind: markdown. Size: 7589 bytes. Shape: 130 lines. Markdown document, 130 lines, 0 explicit URLs. Lead heading: Environment Variables.\\
\textbf{\path{docs/operations/environment.md}}. Тип: markdown. Размер: 7589 байт. Форма: 130 строк. Документ Markdown, 130 строк, явных URL: 0. Главный заголовок: Environment Variables.\par\medskip
\textbf{\path{docs/operations/runtime-image.md}}. Kind: markdown. Size: 1989 bytes. Shape: 46 lines. Markdown document, 46 lines, 0 explicit URLs. Lead heading: Runtime Image Notes.\\
\textbf{\path{docs/operations/runtime-image.md}}. Тип: markdown. Размер: 1989 байт. Форма: 46 строк. Документ Markdown, 46 строк, явных URL: 0. Главный заголовок: Runtime Image Notes.\par\medskip
\textbf{\path{docs/operations/simple-end-to-end-verification-guide.md}}. Kind: markdown. Size: 4111 bytes. Shape: 105 lines. Markdown document, 105 lines, 12 explicit URLs. Lead heading: Simple End-To-End Verification Guide.\\
\textbf{\path{docs/operations/simple-end-to-end-verification-guide.md}}. Тип: markdown. Размер: 4111 байт. Форма: 105 строк. Документ Markdown, 105 строк, явных URL: 12. Главный заголовок: Simple End-To-End Verification Guide.\par\medskip
\textbf{\path{docs/operations/trainer-deployment.md}}. Kind: markdown. Size: 6047 bytes. Shape: 152 lines. Markdown document, 152 lines, 0 explicit URLs. Lead heading: Trainer Deployment Notes.\\
\textbf{\path{docs/operations/trainer-deployment.md}}. Тип: markdown. Размер: 6047 байт. Форма: 152 строк. Документ Markdown, 152 строк, явных URL: 0. Главный заголовок: Trainer Deployment Notes.\par\medskip
\textbf{\path{docs/planning/codex-exec-resume-plan.md}}. Kind: markdown. Size: 11561 bytes. Shape: 207 lines. Markdown document, 207 lines, 0 explicit URLs. Lead heading: Codex Exec Resume Plan.\\
\textbf{\path{docs/planning/codex-exec-resume-plan.md}}. Тип: markdown. Размер: 11561 байт. Форма: 207 строк. Документ Markdown, 207 строк, явных URL: 0. Главный заголовок: Codex Exec Resume Plan.\par\medskip
\textbf{\path{docs/planning/dataset-integration-plan.md}}. Kind: markdown. Size: 33882 bytes. Shape: 478 lines. Markdown document, 478 lines, 0 explicit URLs. Lead heading: Dataset Integration Plan.\\
\textbf{\path{docs/planning/dataset-integration-plan.md}}. Тип: markdown. Размер: 33882 байт. Форма: 478 строк. Документ Markdown, 478 строк, явных URL: 0. Главный заголовок: Dataset Integration Plan.\par\medskip
\textbf{\path{docs/planning/family-first-mipro-runtime-contract.md}}. Kind: markdown. Size: 29074 bytes. Shape: 439 lines. Markdown document, 439 lines, 0 explicit URLs. Lead heading: Family-First DSPy / MIPROv2 Contract.\\
\textbf{\path{docs/planning/family-first-mipro-runtime-contract.md}}. Тип: markdown. Размер: 29074 байт. Форма: 439 строк. Документ Markdown, 439 строк, явных URL: 0. Главный заголовок: Family-First DSPy / MIPROv2 Contract.\par\medskip
\textbf{\path{docs/planning/per-turn-dspy-mediation-contract.md}}. Kind: markdown. Size: 4966 bytes. Shape: 115 lines. Markdown document, 115 lines, 0 explicit URLs. Lead heading: Per-Turn DSPy Mediation Contract.\\
\textbf{\path{docs/planning/per-turn-dspy-mediation-contract.md}}. Тип: markdown. Размер: 4966 байт. Форма: 115 строк. Документ Markdown, 115 строк, явных URL: 0. Главный заголовок: Per-Turn DSPy Mediation Contract.\par\medskip
\textbf{\path{docs/planning/repo-hardening-plan.md}}. Kind: markdown. Size: 4545 bytes. Shape: 71 lines. Markdown document, 71 lines, 0 explicit URLs. Lead heading: Repository Hardening Plan.\\
\textbf{\path{docs/planning/repo-hardening-plan.md}}. Тип: markdown. Размер: 4545 байт. Форма: 71 строк. Документ Markdown, 71 строк, явных URL: 0. Главный заголовок: Repository Hardening Plan.\par\medskip
\textbf{\path{docs/planning/trainer-context-group-champion-plan.md}}. Kind: markdown. Size: 7736 bytes. Shape: 119 lines. Markdown document, 119 lines, 0 explicit URLs. Lead heading: Trainer Context-Group Champion Plan.\\
\textbf{\path{docs/planning/trainer-context-group-champion-plan.md}}. Тип: markdown. Размер: 7736 байт. Форма: 119 строк. Документ Markdown, 119 строк, явных URL: 0. Главный заголовок: Trainer Context-Group Champion Plan.\par\medskip
\textbf{\path{mkdocs.yml}}. Kind: yaml. Size: 843 bytes. Shape: 37 lines. YAML data or workflow file, 37 lines, 2 explicit URLs. Top keys: site\_name, site\_description, site\_url.\\
\textbf{\path{mkdocs.yml}}. Тип: yaml. Размер: 843 байт. Форма: 37 строк. Файл YAML с данными или workflow, 37 строк, явных URL: 2. Верхние ключи: site\_name, site\_description, site\_url.\par\medskip
\textbf{\path{notebooks/01_repo_rag_research.ipynb}}. Kind: notebook. Size: 9324 bytes. Shape: 281 lines. Notebook, 11 cells, 0 explicit URLs. Opening heading: Repository RAG Research Playbook.\\
\textbf{\path{notebooks/01_repo_rag_research.ipynb}}. Тип: notebook. Размер: 9324 байт. Форма: 281 строк. Ноутбук, 11 ячеек, явных URL: 0. Первый заголовок: Repository RAG Research Playbook.\par\medskip
\textbf{\path{notebooks/02_agent_workflow_checklist.ipynb}}. Kind: notebook. Size: 17632 bytes. Shape: 404 lines. Notebook, 9 cells, 1 explicit URLs. Opening heading: Agent Workflow Checklist.\\
\textbf{\path{notebooks/02_agent_workflow_checklist.ipynb}}. Тип: notebook. Размер: 17632 байт. Форма: 404 строк. Ноутбук, 9 ячеек, явных URL: 1. Первый заголовок: Agent Workflow Checklist.\par\medskip
\textbf{\path{notebooks/03_dspy_training_lab.ipynb}}. Kind: notebook. Size: 12048 bytes. Shape: 314 lines. Notebook, 11 cells, 0 explicit URLs. Opening heading: DSPy Training Sample Lab.\\
\textbf{\path{notebooks/03_dspy_training_lab.ipynb}}. Тип: notebook. Размер: 12048 байт. Форма: 314 строк. Ноутбук, 11 ячеек, явных URL: 0. Первый заголовок: DSPy Training Sample Lab.\par\medskip
\textbf{\path{notebooks/04_sample_population_lab.ipynb}}. Kind: notebook. Size: 7894 bytes. Shape: 250 lines. Notebook, 9 cells, 0 explicit URLs. Opening heading: Sample Population Lab.\\
\textbf{\path{notebooks/04_sample_population_lab.ipynb}}. Тип: notebook. Размер: 7894 байт. Форма: 250 строк. Ноутбук, 9 ячеек, явных URL: 0. Первый заголовок: Sample Population Lab.\par\medskip
\textbf{\path{notebooks/05_hushwheel_fixture_rag_lab.ipynb}}. Kind: notebook. Size: 13622 bytes. Shape: 378 lines. Notebook, 9 cells, 0 explicit URLs. Opening heading: Hushwheel Fixture RAG Lab.\\
\textbf{\path{notebooks/05_hushwheel_fixture_rag_lab.ipynb}}. Тип: notebook. Размер: 13622 байт. Форма: 378 строк. Ноутбук, 9 ячеек, явных URL: 0. Первый заголовок: Hushwheel Fixture RAG Lab.\par\medskip
\textbf{\path{publication/.gitignore}}. Kind: generic. Size: 80 bytes. Shape: 12 lines. Tracked repository file, 12 lines, 0 explicit URLs. Opening line: .build/\\
\textbf{\path{publication/.gitignore}}. Тип: generic. Размер: 80 байт. Форма: 12 строк. Отслеживаемый файл репозитория, 12 строк, явных URL: 0. Первая строка: .build/\par\medskip
\textbf{\path{publication/Makefile}}. Kind: make. Size: 1644 bytes. Shape: 48 lines. Repository Makefile, 48 lines, 0 explicit URLs. It exposes the shared automation surfaces.\\
\textbf{\path{publication/Makefile}}. Тип: make. Размер: 1644 байт. Форма: 48 строк. Makefile репозитория, 48 строк, явных URL: 0. Он открывает общие поверхности автоматизации.\par\medskip
\textbf{\path{publication/README.md}}. Kind: markdown. Size: 1499 bytes. Shape: 34 lines. Markdown document, 34 lines, 0 explicit URLs. Lead heading: Publication Article.\\
\textbf{\path{publication/README.md}}. Тип: markdown. Размер: 1499 байт. Форма: 34 строк. Документ Markdown, 34 строк, явных URL: 0. Главный заголовок: Publication Article.\par\medskip
\textbf{\path{publication/article-banner.png}}. Kind: binary. Size: 79074 bytes. Shape: binary surface. Binary asset, 79074 bytes of binary data. It is tracked for stable publication output.\\
\textbf{\path{publication/article-banner.png}}. Тип: binary. Размер: 79074 байт. Форма: бинарная поверхность. Бинарный артефакт, 79074 байт бинарных данных. Он хранится в Git для стабильного публикационного вывода.\par\medskip
\textbf{\path{publication/exploratorium_translation/.gitignore}}. Kind: generic. Size: 32 bytes. Shape: 6 lines. Tracked repository file, 6 lines, 0 explicit URLs. Opening line: .build/\\
\textbf{\path{publication/exploratorium_translation/.gitignore}}. Тип: generic. Размер: 32 байт. Форма: 6 строк. Отслеживаемый файл репозитория, 6 строк, явных URL: 0. Первая строка: .build/\par\medskip
\textbf{\path{publication/exploratorium_translation/Makefile}}. Kind: make. Size: 870 bytes. Shape: 30 lines. Repository Makefile, 30 lines, 0 explicit URLs. It exposes the shared automation surfaces.\\
\textbf{\path{publication/exploratorium_translation/Makefile}}. Тип: make. Размер: 870 байт. Форма: 30 строк. Makefile репозитория, 30 строк, явных URL: 0. Он открывает общие поверхности автоматизации.\par\medskip
\textbf{\path{publication/exploratorium_translation/README.md}}. Kind: markdown. Size: 997 bytes. Shape: 28 lines. Markdown document, 28 lines, 0 explicit URLs. Lead heading: Exploratorium Translation.\\
\textbf{\path{publication/exploratorium_translation/README.md}}. Тип: markdown. Размер: 997 байт. Форма: 28 строк. Документ Markdown, 28 строк, явных URL: 0. Главный заголовок: Exploratorium Translation.\par\medskip
\textbf{\path{publication/exploratorium_translation/exploratorium_translation.tex}}. Kind: latex. Size: 2573 bytes. Shape: 70 lines. LaTeX source, 70 lines, 0 explicit URLs. Sections declared in preview: 0.\\
\textbf{\path{publication/exploratorium_translation/exploratorium_translation.tex}}. Тип: latex. Размер: 2573 байт. Форма: 70 строк. Исходник LaTeX, 70 строк, явных URL: 0. Секций в превью: 0.\par\medskip
\textbf{\path{publication/references.bib}}. Kind: bibliography. Size: 1252 bytes. Shape: 28 lines. Bibliography file, 28 lines, 2 explicit URLs. Entries counted in preview: 4.\\
\textbf{\path{publication/references.bib}}. Тип: bibliography. Размер: 1252 байт. Форма: 28 строк. Библиографический файл, 28 строк, явных URL: 2. Записей в превью: 4.\par\medskip
\textbf{\path{publication/repository-rag-lab-article.pdf}}. Kind: binary. Size: 236232 bytes. Shape: binary surface. Binary asset, 236232 bytes of binary data. It is tracked for stable publication output.\\
\textbf{\path{publication/repository-rag-lab-article.pdf}}. Тип: binary. Размер: 236232 байт. Форма: бинарная поверхность. Бинарный артефакт, 236232 байт бинарных данных. Он хранится в Git для стабильного публикационного вывода.\par\medskip
\textbf{\path{publication/repository-rag-lab-article.tex}}. Kind: latex. Size: 9526 bytes. Shape: 208 lines. LaTeX source, 208 lines, 0 explicit URLs. Sections declared in preview: 4.\\
\textbf{\path{publication/repository-rag-lab-article.tex}}. Тип: latex. Размер: 9526 байт. Форма: 208 строк. Исходник LaTeX, 208 строк, явных URL: 0. Секций в превью: 4.\par\medskip
\textbf{\path{publication/todo-backlog-table.tex}}. Kind: latex. Size: 10547 bytes. Shape: 34 lines. LaTeX source, 34 lines, 36 explicit URLs. Sections declared in preview: 0.\\
\textbf{\path{publication/todo-backlog-table.tex}}. Тип: latex. Размер: 10547 байт. Форма: 34 строк. Исходник LaTeX, 34 строк, явных URL: 36. Секций в превью: 0.\par\medskip
\textbf{\path{pyproject.toml}}. Kind: toml. Size: 2285 bytes. Shape: 97 lines. TOML configuration file, 97 lines, 0 explicit URLs. Top keys: build-system, project, dependency-groups.\\
\textbf{\path{pyproject.toml}}. Тип: toml. Размер: 2285 байт. Форма: 97 строк. Конфигурация TOML, 97 строк, явных URL: 0. Верхние ключи: build-system, project, dependency-groups.\par\medskip
\textbf{\path{rust-cli/Cargo.lock}}. Kind: generic. Size: 3213 bytes. Shape: binary surface. Tracked repository file, 3213 bytes of binary data, 0 explicit URLs. Opening line: No non-empty preview line.\\
\textbf{\path{rust-cli/Cargo.lock}}. Тип: generic. Размер: 3213 байт. Форма: бинарная поверхность. Отслеживаемый файл репозитория, 3213 байт бинарных данных, явных URL: 0. Первая строка: No non-empty preview line.\par\medskip
\textbf{\path{rust-cli/Cargo.toml}}. Kind: toml. Size: 144 bytes. Shape: 8 lines. TOML configuration file, 8 lines, 0 explicit URLs. Top keys: package, dependencies.\\
\textbf{\path{rust-cli/Cargo.toml}}. Тип: toml. Размер: 144 байт. Форма: 8 строк. Конфигурация TOML, 8 строк, явных URL: 0. Верхние ключи: package, dependencies.\par\medskip
\textbf{\path{rust-cli/src/main.rs}}. Kind: source. Size: 20422 bytes. Shape: 626 lines. Source file, 626 lines, 0 explicit URLs. Opening line: use rusqlite::\{params, Connection\};\\
\textbf{\path{rust-cli/src/main.rs}}. Тип: source. Размер: 20422 байт. Форма: 626 строк. Исходный файл, 626 строк, явных URL: 0. Первая строка: use rusqlite::\{params, Connection\};\par\medskip
\textbf{\path{samples/README.md}}. Kind: markdown. Size: 713 bytes. Shape: 13 lines. Markdown document, 13 lines, 0 explicit URLs. Lead heading: Samples.\\
\textbf{\path{samples/README.md}}. Тип: markdown. Размер: 713 байт. Форма: 13 строк. Документ Markdown, 13 строк, явных URL: 0. Главный заголовок: Samples.\par\medskip
\textbf{\path{samples/logs/20260315T050730Z-gh-runs.md}}. Kind: markdown. Size: 1341 bytes. Shape: 45 lines. Markdown document, 45 lines, 3 explicit URLs. Lead heading: GitHub Actions Run Log.\\
\textbf{\path{samples/logs/20260315T050730Z-gh-runs.md}}. Тип: markdown. Размер: 1341 байт. Форма: 45 строк. Документ Markdown, 45 строк, явных URL: 3. Главный заголовок: GitHub Actions Run Log.\par\medskip
\textbf{\path{samples/logs/20260315T085140Z-gh-red-status-check.md}}. Kind: markdown. Size: 847 bytes. Shape: 25 lines. Markdown document, 25 lines, 1 explicit URLs. Lead heading: GitHub Actions Red Status Check.\\
\textbf{\path{samples/logs/20260315T085140Z-gh-red-status-check.md}}. Тип: markdown. Размер: 847 байт. Форма: 25 строк. Документ Markdown, 25 строк, явных URL: 1. Главный заголовок: GitHub Actions Red Status Check.\par\medskip
\textbf{\path{samples/logs/20260317T064722Z-full-build.raw.log}}. Kind: generic. Size: 2109 bytes. Shape: 62 lines. Tracked repository file, 62 lines, 0 explicit URLs. Opening line: \# Full Build Raw Log\\
\textbf{\path{samples/logs/20260317T064722Z-full-build.raw.log}}. Тип: generic. Размер: 2109 байт. Форма: 62 строк. Отслеживаемый файл репозитория, 62 строк, явных URL: 0. Первая строка: \# Full Build Raw Log\par\medskip
\textbf{\path{samples/logs/20260317T064833Z-full-build-attempt2.raw.log}}. Kind: generic. Size: 14950 bytes. Shape: 408 lines. Tracked repository file, 408 lines, 0 explicit URLs. Opening line: \# Full Build Raw Log Attempt 2\\
\textbf{\path{samples/logs/20260317T064833Z-full-build-attempt2.raw.log}}. Тип: generic. Размер: 14950 байт. Форма: 408 строк. Отслеживаемый файл репозитория, 408 строк, явных URL: 0. Первая строка: \# Full Build Raw Log Attempt 2\par\medskip
\textbf{\path{samples/logs/20260317T073401Z-gh-runs-notebook-scaffolding.md}}. Kind: markdown. Size: 2865 bytes. Shape: 35 lines. Markdown document, 35 lines, 13 explicit URLs. Lead heading: GitHub Actions Run Log.\\
\textbf{\path{samples/logs/20260317T073401Z-gh-runs-notebook-scaffolding.md}}. Тип: markdown. Размер: 2865 байт. Форма: 35 строк. Документ Markdown, 35 строк, явных URL: 13. Главный заголовок: GitHub Actions Run Log.\par\medskip
\textbf{\path{samples/logs/20260317T073404Z-gh-runs.md}}. Kind: markdown. Size: 1269 bytes. Shape: 32 lines. Markdown document, 32 lines, 3 explicit URLs. Lead heading: GitHub Actions Run Log.\\
\textbf{\path{samples/logs/20260317T073404Z-gh-runs.md}}. Тип: markdown. Размер: 1269 байт. Форма: 32 строк. Документ Markdown, 32 строк, явных URL: 3. Главный заголовок: GitHub Actions Run Log.\par\medskip
\textbf{\path{samples/logs/20260317T080018Z-gh-runs-all-notebooks.md}}. Kind: markdown. Size: 2128 bytes. Shape: 34 lines. Markdown document, 34 lines, 8 explicit URLs. Lead heading: GitHub Actions Run Log.\\
\textbf{\path{samples/logs/20260317T080018Z-gh-runs-all-notebooks.md}}. Тип: markdown. Размер: 2128 байт. Форма: 34 строк. Документ Markdown, 34 строк, явных URL: 8. Главный заголовок: GitHub Actions Run Log.\par\medskip
\textbf{\path{samples/logs/20260317T083207Z-gh-runs-todo-backlog.md}}. Kind: markdown. Size: 1063 bytes. Shape: 27 lines. Markdown document, 27 lines, 3 explicit URLs. Lead heading: GitHub Actions Run Log.\\
\textbf{\path{samples/logs/20260317T083207Z-gh-runs-todo-backlog.md}}. Тип: markdown. Размер: 1063 байт. Форма: 27 строк. Документ Markdown, 27 строк, явных URL: 3. Главный заголовок: GitHub Actions Run Log.\par\medskip
\textbf{\path{samples/logs/20260317T084723Z-gh-runs-repo-completeness-checklist.md}}. Kind: markdown. Size: 1186 bytes. Shape: 27 lines. Markdown document, 27 lines, 3 explicit URLs. Lead heading: GitHub Actions Run Log.\\
\textbf{\path{samples/logs/20260317T084723Z-gh-runs-repo-completeness-checklist.md}}. Тип: markdown. Размер: 1186 байт. Форма: 27 строк. Документ Markdown, 27 строк, явных URL: 3. Главный заголовок: GitHub Actions Run Log.\par\medskip
\textbf{\path{samples/logs/20260317T085723Z-gh-runs-hushwheel-fixture.md}}. Kind: markdown. Size: 1472 bytes. Shape: 36 lines. Markdown document, 36 lines, 1 explicit URLs. Lead heading: GitHub Actions Run Log.\\
\textbf{\path{samples/logs/20260317T085723Z-gh-runs-hushwheel-fixture.md}}. Тип: markdown. Размер: 1472 байт. Форма: 36 строк. Документ Markdown, 36 строк, явных URL: 1. Главный заголовок: GitHub Actions Run Log.\par\medskip
\textbf{\path{samples/logs/20260317T091641Z-gh-runs-hushwheel-rag-playbook.md}}. Kind: markdown. Size: 1470 bytes. Shape: 36 lines. Markdown document, 36 lines, 1 explicit URLs. Lead heading: GitHub Actions Run Log.\\
\textbf{\path{samples/logs/20260317T091641Z-gh-runs-hushwheel-rag-playbook.md}}. Тип: markdown. Размер: 1470 байт. Форма: 36 строк. Документ Markdown, 36 строк, явных URL: 1. Главный заголовок: GitHub Actions Run Log.\par\medskip
\textbf{\path{samples/logs/20260317T092251Z-gh-runs-publication-article.md}}. Kind: markdown. Size: 1801 bytes. Shape: 33 lines. Markdown document, 33 lines, 5 explicit URLs. Lead heading: GitHub Actions Run Log.\\
\textbf{\path{samples/logs/20260317T092251Z-gh-runs-publication-article.md}}. Тип: markdown. Размер: 1801 байт. Форма: 33 строк. Документ Markdown, 33 строк, явных URL: 5. Главный заголовок: GitHub Actions Run Log.\par\medskip
\textbf{\path{samples/logs/20260318T003027Z-gh-runs-wire-skill-triggers-for-do-repo-verification-aud.md}}. Kind: markdown. Size: 1113 bytes. Shape: 29 lines. Markdown document, 29 lines, 1 explicit URLs. Lead heading: GitHub Actions Run Log.\\
\textbf{\path{samples/logs/20260318T003027Z-gh-runs-wire-skill-triggers-for-do-repo-verification-aud.md}}. Тип: markdown. Размер: 1113 байт. Форма: 29 строк. Документ Markdown, 29 строк, явных URL: 1. Главный заголовок: GitHub Actions Run Log.\par\medskip
\textbf{\path{samples/logs/20260318T003654Z-gh-runs-execute-all-notebooks.md}}. Kind: markdown. Size: 1682 bytes. Shape: 41 lines. Markdown document, 41 lines, 1 explicit URLs. Lead heading: GitHub Actions Run Log.\\
\textbf{\path{samples/logs/20260318T003654Z-gh-runs-execute-all-notebooks.md}}. Тип: markdown. Размер: 1682 байт. Форма: 41 строк. Документ Markdown, 41 строк, явных URL: 1. Главный заголовок: GitHub Actions Run Log.\par\medskip
\textbf{\path{samples/logs/20260318T003927Z-gh-runs-publication-workflow-fix.md}}. Kind: markdown. Size: 1747 bytes. Shape: 35 lines. Markdown document, 35 lines, 2 explicit URLs. Lead heading: GitHub Actions Run Log.\\
\textbf{\path{samples/logs/20260318T003927Z-gh-runs-publication-workflow-fix.md}}. Тип: markdown. Размер: 1747 байт. Форма: 35 строк. Документ Markdown, 35 строк, явных URL: 2. Главный заголовок: GitHub Actions Run Log.\par\medskip
\textbf{\path{samples/logs/20260318T004712Z-gh-runs-hushwheel-program-harness.md}}. Kind: markdown. Size: 1798 bytes. Shape: 49 lines. Markdown document, 49 lines, 0 explicit URLs. Lead heading: GitHub Actions Run Log.\\
\textbf{\path{samples/logs/20260318T004712Z-gh-runs-hushwheel-program-harness.md}}. Тип: markdown. Размер: 1798 байт. Форма: 49 строк. Документ Markdown, 49 строк, явных URL: 0. Главный заголовок: GitHub Actions Run Log.\par\medskip
\textbf{\path{samples/logs/20260318T010254Z-gh-runs-refresh-final-notebook-outputs-for-no-move-those.md}}. Kind: markdown. Size: 1113 bytes. Shape: 29 lines. Markdown document, 29 lines, 1 explicit URLs. Lead heading: GitHub Actions Run Log.\\
\textbf{\path{samples/logs/20260318T010254Z-gh-runs-refresh-final-notebook-outputs-for-no-move-those.md}}. Тип: markdown. Размер: 1113 байт. Форма: 29 строк. Документ Markdown, 29 строк, явных URL: 1. Главный заголовок: GitHub Actions Run Log.\par\medskip
\textbf{\path{samples/logs/20260318T010637Z-gh-runs-env-refresh-retest.md}}. Kind: markdown. Size: 1029 bytes. Shape: 26 lines. Markdown document, 26 lines, 1 explicit URLs. Lead heading: GitHub Actions Run Log.\\
\textbf{\path{samples/logs/20260318T010637Z-gh-runs-env-refresh-retest.md}}. Тип: markdown. Размер: 1029 байт. Форма: 26 строк. Документ Markdown, 26 строк, явных URL: 1. Главный заголовок: GitHub Actions Run Log.\par\medskip
\textbf{\path{samples/logs/20260318T010741Z-gh-runs-consolidate-session-master.md}}. Kind: markdown. Size: 1950 bytes. Shape: 62 lines. Markdown document, 62 lines, 0 explicit URLs. Lead heading: GitHub Actions Run Log.\\
\textbf{\path{samples/logs/20260318T010741Z-gh-runs-consolidate-session-master.md}}. Тип: markdown. Размер: 1950 байт. Форма: 62 строк. Документ Markdown, 62 строк, явных URL: 0. Главный заголовок: GitHub Actions Run Log.\par\medskip
\textbf{\path{samples/logs/20260318T012137Z-gh-runs-azure-inference-endpoint-probe.md}}. Kind: markdown. Size: 1004 bytes. Shape: 26 lines. Markdown document, 26 lines, 1 explicit URLs. Lead heading: GitHub Actions Run Log.\\
\textbf{\path{samples/logs/20260318T012137Z-gh-runs-azure-inference-endpoint-probe.md}}. Тип: markdown. Размер: 1004 байт. Форма: 26 строк. Документ Markdown, 26 строк, явных URL: 1. Главный заголовок: GitHub Actions Run Log.\par\medskip
\textbf{\path{samples/logs/20260318T012714Z-gh-runs-tighten-repo-hygiene-for-other-agnet-is-doing-1-.md}}. Kind: markdown. Size: 1091 bytes. Shape: 29 lines. Markdown document, 29 lines, 1 explicit URLs. Lead heading: GitHub Actions Run Log.\\
\textbf{\path{samples/logs/20260318T012714Z-gh-runs-tighten-repo-hygiene-for-other-agnet-is-doing-1-.md}}. Тип: markdown. Размер: 1091 байт. Форма: 29 строк. Документ Markdown, 29 строк, явных URL: 1. Главный заголовок: GitHub Actions Run Log.\par\medskip
\textbf{\path{samples/logs/20260318T015109Z-gh-runs-improve-notebook-run-observability-and-reporting.md}}. Kind: markdown. Size: 2151 bytes. Shape: 55 lines. Markdown document, 55 lines, 2 explicit URLs. Lead heading: GitHub Actions Run Log.\\
\textbf{\path{samples/logs/20260318T015109Z-gh-runs-improve-notebook-run-observability-and-reporting.md}}. Тип: markdown. Размер: 2151 байт. Форма: 55 строк. Документ Markdown, 55 строк, явных URL: 2. Главный заголовок: GitHub Actions Run Log.\par\medskip
\textbf{\path{samples/logs/20260318T021736Z-gh-runs-sync-todo-backlog-into-markdown-and-publication-surfaces.md}}. Kind: markdown. Size: 1778 bytes. Shape: 55 lines. Markdown document, 55 lines, 2 explicit URLs. Lead heading: GitHub Actions Run Log.\\
\textbf{\path{samples/logs/20260318T021736Z-gh-runs-sync-todo-backlog-into-markdown-and-publication-surfaces.md}}. Тип: markdown. Размер: 1778 байт. Форма: 55 строк. Документ Markdown, 55 строк, явных URL: 2. Главный заголовок: GitHub Actions Run Log.\par\medskip
\textbf{\path{samples/logs/20260318T023922Z-gh-runs-implement-step-1-dspy-training-path-follow-up.md}}. Kind: markdown. Size: 1735 bytes. Shape: 55 lines. Markdown document, 55 lines, 2 explicit URLs. Lead heading: GitHub Actions Run Log.\\
\textbf{\path{samples/logs/20260318T023922Z-gh-runs-implement-step-1-dspy-training-path-follow-up.md}}. Тип: markdown. Размер: 1735 байт. Форма: 55 строк. Документ Markdown, 55 строк, явных URL: 2. Главный заголовок: GitHub Actions Run Log.\par\medskip
\textbf{\path{samples/logs/20260318T044242Z-gh-runs-add-maintained-repository-research-narrative.md}}. Kind: markdown. Size: 2222 bytes. Shape: 56 lines. Markdown document, 56 lines, 2 explicit URLs. Lead heading: GitHub Actions Run Log.\\
\textbf{\path{samples/logs/20260318T044242Z-gh-runs-add-maintained-repository-research-narrative.md}}. Тип: markdown. Размер: 2222 байт. Форма: 56 строк. Документ Markdown, 56 строк, явных URL: 2. Главный заголовок: GitHub Actions Run Log.\par\medskip
\textbf{\path{samples/logs/20260318T050051Z-gh-runs-add-azure-runtime-surfaces-for-do-1-and-2-and-carry-on-to-next.md}}. Kind: markdown. Size: 1828 bytes. Shape: 38 lines. Markdown document, 38 lines, 2 explicit URLs. Lead heading: GitHub Actions Run Log.\\
\textbf{\path{samples/logs/20260318T050051Z-gh-runs-add-azure-runtime-surfaces-for-do-1-and-2-and-carry-on-to-next.md}}. Тип: markdown. Размер: 1828 байт. Форма: 38 строк. Документ Markdown, 38 строк, явных URL: 2. Главный заголовок: GitHub Actions Run Log.\par\medskip
\textbf{\path{samples/logs/20260318T051627Z-gh-runs-develop-retrieval-evaluation-suite.md}}. Kind: markdown. Size: 1768 bytes. Shape: 37 lines. Markdown document, 37 lines, 2 explicit URLs. Lead heading: GitHub Actions Run Log.\\
\textbf{\path{samples/logs/20260318T051627Z-gh-runs-develop-retrieval-evaluation-suite.md}}. Тип: markdown. Размер: 1768 байт. Форма: 37 строк. Документ Markdown, 37 строк, явных URL: 2. Главный заголовок: GitHub Actions Run Log.\par\medskip
\textbf{\path{samples/logs/20260318T060918Z-gh-runs-consolidate-session-master-repair.md}}. Kind: markdown. Size: 2460 bytes. Shape: 70 lines. Markdown document, 70 lines, 3 explicit URLs. Lead heading: GitHub Actions Run Log.\\
\textbf{\path{samples/logs/20260318T060918Z-gh-runs-consolidate-session-master-repair.md}}. Тип: markdown. Размер: 2460 байт. Форма: 70 строк. Документ Markdown, 70 строк, явных URL: 3. Главный заголовок: GitHub Actions Run Log.\par\medskip
\textbf{\path{samples/logs/20260318T063755Z-gh-runs-keep-utility-sync-tests-side-effect-free.md}}. Kind: markdown. Size: 1074 bytes. Shape: 29 lines. Markdown document, 29 lines, 1 explicit URLs. Lead heading: GitHub Actions Run Log.\\
\textbf{\path{samples/logs/20260318T063755Z-gh-runs-keep-utility-sync-tests-side-effect-free.md}}. Тип: markdown. Размер: 1074 байт. Форма: 29 строк. Документ Markdown, 29 строк, явных URL: 1. Главный заголовок: GitHub Actions Run Log.\par\medskip
\textbf{\path{samples/logs/20260318T064924Z-gh-runs-retrieval-refresh-go-ahead.md}}. Kind: markdown. Size: 1858 bytes. Shape: 54 lines. Markdown document, 54 lines, 2 explicit URLs. Lead heading: GitHub Actions Run Log.\\
\textbf{\path{samples/logs/20260318T064924Z-gh-runs-retrieval-refresh-go-ahead.md}}. Тип: markdown. Размер: 1858 байт. Форма: 54 строк. Документ Markdown, 54 строк, явных URL: 2. Главный заголовок: GitHub Actions Run Log.\par\medskip
\textbf{\path{samples/logs/20260318T083955Z-gh-runs-benchmark-broadening-2.md}}. Kind: markdown. Size: 1217 bytes. Shape: 34 lines. Markdown document, 34 lines, 1 explicit URLs. Lead heading: GitHub Actions Run Log.\\
\textbf{\path{samples/logs/20260318T083955Z-gh-runs-benchmark-broadening-2.md}}. Тип: markdown. Размер: 1217 байт. Форма: 34 строк. Документ Markdown, 34 строк, явных URL: 1. Главный заголовок: GitHub Actions Run Log.\par\medskip
\textbf{\path{samples/logs/20260318T090625Z-gh-runs-hushwheel-quality-instrumentation.md}}. Kind: markdown. Size: 1312 bytes. Shape: 29 lines. Markdown document, 29 lines, 1 explicit URLs. Lead heading: GitHub Actions Run Log.\\
\textbf{\path{samples/logs/20260318T090625Z-gh-runs-hushwheel-quality-instrumentation.md}}. Тип: markdown. Размер: 1312 байт. Форма: 29 строк. Документ Markdown, 29 строк, явных URL: 1. Главный заголовок: GitHub Actions Run Log.\par\medskip
\textbf{\path{samples/logs/20260318T091252Z-gh-runs-add-rust-lookup-and-retrieval-tag-summaries-for-.md}}. Kind: markdown. Size: 2223 bytes. Shape: 53 lines. Markdown document, 53 lines, 2 explicit URLs. Lead heading: GitHub Actions Run Log.\\
\textbf{\path{samples/logs/20260318T091252Z-gh-runs-add-rust-lookup-and-retrieval-tag-summaries-for-.md}}. Тип: markdown. Размер: 2223 байт. Форма: 53 строк. Документ Markdown, 53 строк, явных URL: 2. Главный заголовок: GitHub Actions Run Log.\par\medskip
\textbf{\path{samples/logs/20260318T093632Z-gh-runs-retrieval-regression-gate-great-carry-on.md}}. Kind: markdown. Size: 1325 bytes. Shape: 34 lines. Markdown document, 34 lines, 2 explicit URLs. Lead heading: GitHub Actions Run Log.\\
\textbf{\path{samples/logs/20260318T093632Z-gh-runs-retrieval-regression-gate-great-carry-on.md}}. Тип: markdown. Размер: 1325 байт. Форма: 34 строк. Документ Markdown, 34 строк, явных URL: 2. Главный заголовок: GitHub Actions Run Log.\par\medskip
\textbf{\path{samples/logs/20260318T094120Z-gh-runs-dspy-live-proof-and-rust-wrapper-repair.md}}. Kind: markdown. Size: 1469 bytes. Shape: 45 lines. Markdown document, 45 lines, 0 explicit URLs. Lead heading: GitHub Actions Summary.\\
\textbf{\path{samples/logs/20260318T094120Z-gh-runs-dspy-live-proof-and-rust-wrapper-repair.md}}. Тип: markdown. Размер: 1469 байт. Форма: 45 строк. Документ Markdown, 45 строк, явных URL: 0. Главный заголовок: GitHub Actions Summary.\par\medskip
\textbf{\path{samples/logs/20260318T094347Z-gh-runs-hushwheel-hardening-followup.md}}. Kind: markdown. Size: 2739 bytes. Shape: 72 lines. Markdown document, 72 lines, 2 explicit URLs. Lead heading: GitHub Runs Log.\\
\textbf{\path{samples/logs/20260318T094347Z-gh-runs-hushwheel-hardening-followup.md}}. Тип: markdown. Размер: 2739 байт. Форма: 72 строк. Документ Markdown, 72 строк, явных URL: 2. Главный заголовок: GitHub Runs Log.\par\medskip
\textbf{\path{samples/logs/20260318T100336Z-gh-runs-sync-final-rebased-inventories-for-carry-on-comm.md}}. Kind: markdown. Size: 1103 bytes. Shape: 29 lines. Markdown document, 29 lines, 1 explicit URLs. Lead heading: GitHub Actions Run Log.\\
\textbf{\path{samples/logs/20260318T100336Z-gh-runs-sync-final-rebased-inventories-for-carry-on-comm.md}}. Тип: markdown. Размер: 1103 байт. Форма: 29 строк. Документ Markdown, 29 строк, явных URL: 1. Главный заголовок: GitHub Actions Run Log.\par\medskip
\textbf{\path{samples/logs/20260318T101229Z-gh-runs-hushwheel-audit-refresh.md}}. Kind: markdown. Size: 1707 bytes. Shape: 37 lines. Markdown document, 37 lines, 1 explicit URLs. Lead heading: GitHub Actions Summary.\\
\textbf{\path{samples/logs/20260318T101229Z-gh-runs-hushwheel-audit-refresh.md}}. Тип: markdown. Размер: 1707 байт. Форма: 37 строк. Документ Markdown, 37 строк, явных URL: 1. Главный заголовок: GitHub Actions Summary.\par\medskip
\textbf{\path{samples/logs/20260318T101939Z-gh-runs-land-each-unlanded.md}}. Kind: markdown. Size: 1556 bytes. Shape: 28 lines. Markdown document, 28 lines, 2 explicit URLs. Lead heading: GitHub Actions Run Log.\\
\textbf{\path{samples/logs/20260318T101939Z-gh-runs-land-each-unlanded.md}}. Тип: markdown. Размер: 1556 байт. Форма: 28 строк. Документ Markdown, 28 строк, явных URL: 2. Главный заголовок: GitHub Actions Run Log.\par\medskip
\textbf{\path{samples/logs/20260318T123402Z-gh-runs-fix-file-summary-test-churn-for-1.md}}. Kind: markdown. Size: 1061 bytes. Shape: 29 lines. Markdown document, 29 lines, 1 explicit URLs. Lead heading: GitHub Actions Run Log.\\
\textbf{\path{samples/logs/20260318T123402Z-gh-runs-fix-file-summary-test-churn-for-1.md}}. Тип: markdown. Размер: 1061 байт. Форма: 29 строк. Документ Markdown, 29 строк, явных URL: 1. Главный заголовок: GitHub Actions Run Log.\par\medskip
\textbf{\path{samples/logs/20260318T124437Z-gh-runs-reconcile-stale-dspy-backlog-for-1.md}}. Kind: markdown. Size: 1060 bytes. Shape: 21 lines. Markdown document, 21 lines, 3 explicit URLs. Lead heading: GitHub Actions Log.\\
\textbf{\path{samples/logs/20260318T124437Z-gh-runs-reconcile-stale-dspy-backlog-for-1.md}}. Тип: markdown. Размер: 1060 байт. Форма: 21 строк. Документ Markdown, 21 строк, явных URL: 3. Главный заголовок: GitHub Actions Log.\par\medskip
\textbf{\path{samples/logs/20260318T125547Z-gh-runs-wire-lookup-first-ask-for-1.md}}. Kind: markdown. Size: 1058 bytes. Shape: 22 lines. Markdown document, 22 lines, 2 explicit URLs. Lead heading: GitHub Actions Log.\\
\textbf{\path{samples/logs/20260318T125547Z-gh-runs-wire-lookup-first-ask-for-1.md}}. Тип: markdown. Размер: 1058 байт. Форма: 22 строк. Документ Markdown, 22 строк, явных URL: 2. Главный заголовок: GitHub Actions Log.\par\medskip
\textbf{\path{samples/logs/20260318T125829Z-gh-runs-restore-retrieval-gate-after-rebase.md}}. Kind: markdown. Size: 1585 bytes. Shape: 36 lines. Markdown document, 36 lines, 0 explicit URLs. Lead heading: GitHub Actions Run Log.\\
\textbf{\path{samples/logs/20260318T125829Z-gh-runs-restore-retrieval-gate-after-rebase.md}}. Тип: markdown. Размер: 1585 байт. Форма: 36 строк. Документ Markdown, 36 строк, явных URL: 0. Главный заголовок: GitHub Actions Run Log.\par\medskip
\textbf{\path{samples/logs/20260318T130525Z-gh-runs-finish-node24-publication-cache-followup-for-1.md}}. Kind: markdown. Size: 1119 bytes. Shape: 27 lines. Markdown document, 27 lines, 3 explicit URLs. Lead heading: GitHub Actions Status For `Finish Node 24 publication cache follow-up for "1"`.\\
\textbf{\path{samples/logs/20260318T130525Z-gh-runs-finish-node24-publication-cache-followup-for-1.md}}. Тип: markdown. Размер: 1119 байт. Форма: 27 строк. Документ Markdown, 27 строк, явных URL: 3. Главный заголовок: GitHub Actions Status For `Finish Node 24 publication cache follow-up for "1"`.\par\medskip
\textbf{\path{samples/logs/20260318T131709Z-gh-runs-restore-root-readme-weighting-for-1.md}}. Kind: markdown. Size: 1680 bytes. Shape: 45 lines. Markdown document, 45 lines, 2 explicit URLs. Lead heading: GitHub Actions Run Log.\\
\textbf{\path{samples/logs/20260318T131709Z-gh-runs-restore-root-readme-weighting-for-1.md}}. Тип: markdown. Размер: 1680 байт. Форма: 45 строк. Документ Markdown, 45 строк, явных URL: 2. Главный заголовок: GitHub Actions Run Log.\par\medskip
\textbf{\path{samples/logs/20260318T133247Z-gh-runs-improve-retrieval-doc-priority-for-1.md}}. Kind: markdown. Size: 1335 bytes. Shape: 34 lines. Markdown document, 34 lines, 1 explicit URLs. Lead heading: GitHub Actions Run Log.\\
\textbf{\path{samples/logs/20260318T133247Z-gh-runs-improve-retrieval-doc-priority-for-1.md}}. Тип: markdown. Размер: 1335 байт. Форма: 34 строк. Документ Markdown, 34 строк, явных URL: 1. Главный заголовок: GitHub Actions Run Log.\par\medskip
\textbf{\path{samples/logs/20260318T144102Z-gh-runs-merge-into-master.md}}. Kind: markdown. Size: 2285 bytes. Shape: 69 lines. Markdown document, 69 lines, 2 explicit URLs. Lead heading: GitHub Actions Run Log.\\
\textbf{\path{samples/logs/20260318T144102Z-gh-runs-merge-into-master.md}}. Тип: markdown. Размер: 2285 байт. Форма: 69 строк. Документ Markdown, 69 строк, явных URL: 2. Главный заголовок: GitHub Actions Run Log.\par\medskip
\textbf{\path{samples/logs/20260318T145935Z-gh-runs-consolidate-all-unconsoldiated-changes.md}}. Kind: markdown. Size: 2401 bytes. Shape: 71 lines. Markdown document, 71 lines, 2 explicit URLs. Lead heading: GitHub Actions Run Log.\\
\textbf{\path{samples/logs/20260318T145935Z-gh-runs-consolidate-all-unconsoldiated-changes.md}}. Тип: markdown. Размер: 2401 байт. Форма: 71 строк. Документ Markdown, 71 строк, явных URL: 2. Главный заголовок: GitHub Actions Run Log.\par\medskip
\textbf{\path{samples/logs/20260319T035301Z-gh-runs-public-pages-and-workflow-repairs.md}}. Kind: markdown. Size: 1578 bytes. Shape: 37 lines. Markdown document, 37 lines, 1 explicit URLs. Lead heading: GitHub Actions Log.\\
\textbf{\path{samples/logs/20260319T035301Z-gh-runs-public-pages-and-workflow-repairs.md}}. Тип: markdown. Размер: 1578 байт. Форма: 37 строк. Документ Markdown, 37 строк, явных URL: 1. Главный заголовок: GitHub Actions Log.\par\medskip
\textbf{\path{samples/logs/20260319T071544Z-gh-runs-doncsoldiate-consolidation.md}}. Kind: markdown. Size: 1137 bytes. Shape: 27 lines. Markdown document, 27 lines, 0 explicit URLs. Lead heading: GitHub Actions Log.\\
\textbf{\path{samples/logs/20260319T071544Z-gh-runs-doncsoldiate-consolidation.md}}. Тип: markdown. Размер: 1137 байт. Форма: 27 строк. Документ Markdown, 27 строк, явных URL: 0. Главный заголовок: GitHub Actions Log.\par\medskip
\textbf{\path{samples/logs/20260319T114913Z-gh-runs-fix-publication-verification-for-whats-next-step.md}}. Kind: markdown. Size: 2936 bytes. Shape: 67 lines. Markdown document, 67 lines, 3 explicit URLs. Lead heading: GitHub Actions Run Log.\\
\textbf{\path{samples/logs/20260319T114913Z-gh-runs-fix-publication-verification-for-whats-next-step.md}}. Тип: markdown. Размер: 2936 байт. Форма: 67 строк. Документ Markdown, 67 строк, явных URL: 3. Главный заголовок: GitHub Actions Run Log.\par\medskip
\textbf{\path{samples/logs/20260429T172137Z-gh-runs-develop-no-workflow.md}}. Kind: markdown. Size: 866 bytes. Shape: 23 lines. Markdown document, 23 lines, 0 explicit URLs. Lead heading: GitHub Actions Run Log.\\
\textbf{\path{samples/logs/20260429T172137Z-gh-runs-develop-no-workflow.md}}. Тип: markdown. Размер: 866 байт. Форма: 23 строк. Документ Markdown, 23 строк, явных URL: 0. Главный заголовок: GitHub Actions Run Log.\par\medskip
\textbf{\path{samples/logs/20260429T200446Z-gh-runs-develop-no-new-run-after-blob-queue-migration.md}}. Kind: markdown. Size: 1843 bytes. Shape: 34 lines. Markdown document, 34 lines, 0 explicit URLs. Lead heading: GitHub Actions Run Log.\\
\textbf{\path{samples/logs/20260429T200446Z-gh-runs-develop-no-new-run-after-blob-queue-migration.md}}. Тип: markdown. Размер: 1843 байт. Форма: 34 строк. Документ Markdown, 34 строк, явных URL: 0. Главный заголовок: GitHub Actions Run Log.\par\medskip
\textbf{\path{samples/logs/20260430T122222Z-gh-runs-develop-no-new-run-after-codex-proxy-push.md}}. Kind: markdown. Size: 908 bytes. Shape: 25 lines. Markdown document, 25 lines, 0 explicit URLs. Lead heading: GitHub Actions check after pushing `bbdcf4b`.\\
\textbf{\path{samples/logs/20260430T122222Z-gh-runs-develop-no-new-run-after-codex-proxy-push.md}}. Тип: markdown. Размер: 908 байт. Форма: 25 строк. Документ Markdown, 25 строк, явных URL: 0. Главный заголовок: GitHub Actions check after pushing `bbdcf4b`.\par\medskip
\textbf{\path{samples/logs/20260430T172518Z-gh-runs-develop-no-new-run-after-trace-handoff-push.md}}. Kind: markdown. Size: 991 bytes. Shape: 27 lines. Markdown document, 27 lines, 0 explicit URLs. Lead heading: GitHub Actions check after pushing `839f4cb`.\\
\textbf{\path{samples/logs/20260430T172518Z-gh-runs-develop-no-new-run-after-trace-handoff-push.md}}. Тип: markdown. Размер: 991 байт. Форма: 27 строк. Документ Markdown, 27 строк, явных URL: 0. Главный заголовок: GitHub Actions check after pushing `839f4cb`.\par\medskip
\textbf{\path{samples/logs/20260501T103130Z-gh-runs-develop-no-new-run-after-azure-runtime-env-fix-push.md}}. Kind: markdown. Size: 1411 bytes. Shape: 35 lines. Markdown document, 35 lines, 0 explicit URLs. Lead heading: GitHub Actions check after pushing `9e627e1`.\\
\textbf{\path{samples/logs/20260501T103130Z-gh-runs-develop-no-new-run-after-azure-runtime-env-fix-push.md}}. Тип: markdown. Размер: 1411 байт. Форма: 35 строк. Документ Markdown, 35 строк, явных URL: 0. Главный заголовок: GitHub Actions check after pushing `9e627e1`.\par\medskip
\textbf{\path{samples/logs/20260501T112330Z-gh-runs-develop-no-new-run-after-azure-responses-api-version-fix-push.md}}. Kind: markdown. Size: 1522 bytes. Shape: 37 lines. Markdown document, 37 lines, 0 explicit URLs. Lead heading: GitHub Actions check after pushing `128874b`.\\
\textbf{\path{samples/logs/20260501T112330Z-gh-runs-develop-no-new-run-after-azure-responses-api-version-fix-push.md}}. Тип: markdown. Размер: 1522 байт. Форма: 37 строк. Документ Markdown, 37 строк, явных URL: 0. Главный заголовок: GitHub Actions check after pushing `128874b`.\par\medskip
\textbf{\path{samples/logs/20260501T125058Z-gh-runs-develop-no-new-run-after-trusted-trace-handoff-push.md}}. Kind: markdown. Size: 1355 bytes. Shape: 30 lines. Markdown document, 30 lines, 0 explicit URLs. Lead heading: GitHub Actions check after pushing `66a7ad4`.\\
\textbf{\path{samples/logs/20260501T125058Z-gh-runs-develop-no-new-run-after-trusted-trace-handoff-push.md}}. Тип: markdown. Размер: 1355 байт. Форма: 30 строк. Документ Markdown, 30 строк, явных URL: 0. Главный заголовок: GitHub Actions check after pushing `66a7ad4`.\par\medskip
\textbf{\path{samples/logs/20260501T162125Z-gh-runs-develop-no-new-run-after-live-trainer-bundle-publication-push.md}}. Kind: markdown. Size: 1486 bytes. Shape: 30 lines. Markdown document, 30 lines, 0 explicit URLs. Lead heading: GitHub Run Inspection.\\
\textbf{\path{samples/logs/20260501T162125Z-gh-runs-develop-no-new-run-after-live-trainer-bundle-publication-push.md}}. Тип: markdown. Размер: 1486 байт. Форма: 30 строк. Документ Markdown, 30 строк, явных URL: 0. Главный заголовок: GitHub Run Inspection.\par\medskip
\textbf{\path{samples/logs/20260501T174147Z-gh-runs-develop-no-new-run-after-versioned-trainer-recovery-push.md}}. Kind: markdown. Size: 1505 bytes. Shape: 30 lines. Markdown document, 30 lines, 0 explicit URLs. Lead heading: GitHub Run Inspection.\\
\textbf{\path{samples/logs/20260501T174147Z-gh-runs-develop-no-new-run-after-versioned-trainer-recovery-push.md}}. Тип: markdown. Размер: 1505 байт. Форма: 30 строк. Документ Markdown, 30 строк, явных URL: 0. Главный заголовок: GitHub Run Inspection.\par\medskip
\textbf{\path{samples/logs/20260501T201242Z-gh-runs-develop-no-new-run-after-dspy-bundle-version-pinning-push.md}}. Kind: markdown. Size: 2076 bytes. Shape: 41 lines. Markdown document, 41 lines, 0 explicit URLs. Lead heading: GitHub Run Inspection.\\
\textbf{\path{samples/logs/20260501T201242Z-gh-runs-develop-no-new-run-after-dspy-bundle-version-pinning-push.md}}. Тип: markdown. Размер: 2076 байт. Форма: 41 строк. Документ Markdown, 41 строк, явных URL: 0. Главный заголовок: GitHub Run Inspection.\par\medskip
\textbf{\path{samples/logs/20260502T064702Z-gh-runs-develop-no-new-run-after-live-trainer-timestamp-bundle-audit-push.md}}. Kind: markdown. Size: 1604 bytes. Shape: 38 lines. Markdown document, 38 lines, 0 explicit URLs. Lead heading: GitHub Run Inspection.\\
\textbf{\path{samples/logs/20260502T064702Z-gh-runs-develop-no-new-run-after-live-trainer-timestamp-bundle-audit-push.md}}. Тип: markdown. Размер: 1604 байт. Форма: 38 строк. Документ Markdown, 38 строк, явных URL: 0. Главный заголовок: GitHub Run Inspection.\par\medskip
\textbf{\path{samples/logs/20260502T082402Z-gh-runs-develop-no-new-run-after-hybrid-vector-retrieval-push.md}}. Kind: markdown. Size: 1631 bytes. Shape: 30 lines. Markdown document, 30 lines, 0 explicit URLs. Lead heading: GitHub Run Inspection.\\
\textbf{\path{samples/logs/20260502T082402Z-gh-runs-develop-no-new-run-after-hybrid-vector-retrieval-push.md}}. Тип: markdown. Размер: 1631 байт. Форма: 30 строк. Документ Markdown, 30 строк, явных URL: 0. Главный заголовок: GitHub Run Inspection.\par\medskip
\textbf{\path{samples/logs/20260502T121755Z-gh-runs-develop-no-new-run-after-trainer-gating-and-retrieval-isolation-push.md}}. Kind: markdown. Size: 1175 bytes. Shape: 25 lines. Markdown document, 25 lines, 0 explicit URLs. Lead heading: Summary.\\
\textbf{\path{samples/logs/20260502T121755Z-gh-runs-develop-no-new-run-after-trainer-gating-and-retrieval-isolation-push.md}}. Тип: markdown. Размер: 1175 байт. Форма: 25 строк. Документ Markdown, 25 строк, явных URL: 0. Главный заголовок: Summary.\par\medskip
\textbf{\path{samples/logs/20260502T165750Z-gh-runs-develop-no-new-run-after-trainer-champion-batching-push.md}}. Kind: markdown. Size: 1543 bytes. Shape: 35 lines. Markdown document, 35 lines, 0 explicit URLs. Lead heading: GitHub Run Inspection.\\
\textbf{\path{samples/logs/20260502T165750Z-gh-runs-develop-no-new-run-after-trainer-champion-batching-push.md}}. Тип: markdown. Размер: 1543 байт. Форма: 35 строк. Документ Markdown, 35 строк, явных URL: 0. Главный заголовок: GitHub Run Inspection.\par\medskip
\textbf{\path{samples/logs/20260503T080503Z-gh-runs-develop-no-new-run-after-codex-resume-lane-telemetry-push.md}}. Kind: markdown. Size: 1782 bytes. Shape: 37 lines. Markdown document, 37 lines, 0 explicit URLs. Lead heading: GitHub Run Inspection.\\
\textbf{\path{samples/logs/20260503T080503Z-gh-runs-develop-no-new-run-after-codex-resume-lane-telemetry-push.md}}. Тип: markdown. Размер: 1782 байт. Форма: 37 строк. Документ Markdown, 37 строк, явных URL: 0. Главный заголовок: GitHub Run Inspection.\par\medskip
\textbf{\path{samples/logs/20260503T153706Z-gh-runs-develop-no-new-run-after-staged-bundle-mirror-push.md}}. Kind: markdown. Size: 1171 bytes. Shape: 27 lines. Markdown document, 27 lines, 0 explicit URLs. Lead heading: GitHub Runs After `Support staged bundle mirrors and stale queue skips`.\\
\textbf{\path{samples/logs/20260503T153706Z-gh-runs-develop-no-new-run-after-staged-bundle-mirror-push.md}}. Тип: markdown. Размер: 1171 байт. Форма: 27 строк. Документ Markdown, 27 строк, явных URL: 0. Главный заголовок: GitHub Runs After `Support staged bundle mirrors and stale queue skips`.\par\medskip
\textbf{\path{samples/logs/20260503T160232Z-gh-runs-develop-no-new-run-after-no-op-trainer-promotion-fix-push.md}}. Kind: markdown. Size: 1482 bytes. Shape: 32 lines. Markdown document, 32 lines, 0 explicit URLs. Lead heading: GitHub Run Inspection.\\
\textbf{\path{samples/logs/20260503T160232Z-gh-runs-develop-no-new-run-after-no-op-trainer-promotion-fix-push.md}}. Тип: markdown. Размер: 1482 байт. Форма: 32 строк. Документ Markdown, 32 строк, явных URL: 0. Главный заголовок: GitHub Run Inspection.\par\medskip
\textbf{\path{samples/logs/20260504T090704Z-gh-runs-develop-no-new-run-after-fix-codex-resume-pvc-session-path-push.md}}. Kind: markdown. Size: 1689 bytes. Shape: 46 lines. Markdown document, 46 lines, 0 explicit URLs. Lead heading: GitHub Run Check After `Fix Codex resume PVC session path`.\\
\textbf{\path{samples/logs/20260504T090704Z-gh-runs-develop-no-new-run-after-fix-codex-resume-pvc-session-path-push.md}}. Тип: markdown. Размер: 1689 байт. Форма: 46 строк. Документ Markdown, 46 строк, явных URL: 0. Главный заголовок: GitHub Run Check After `Fix Codex resume PVC session path`.\par\medskip
\textbf{\path{samples/logs/20260504T111427Z-gh-runs-develop-no-new-run-after-codex-resume-restore-fallback-debug-push.md}}. Kind: markdown. Size: 3370 bytes. Shape: 72 lines. Markdown document, 72 lines, 0 explicit URLs. Lead heading: GitHub Run Check After `Document Codex resume restore fallback debug`.\\
\textbf{\path{samples/logs/20260504T111427Z-gh-runs-develop-no-new-run-after-codex-resume-restore-fallback-debug-push.md}}. Тип: markdown. Размер: 3370 байт. Форма: 72 строк. Документ Markdown, 72 строк, явных URL: 0. Главный заголовок: GitHub Run Check After `Document Codex resume restore fallback debug`.\par\medskip
\textbf{\path{samples/logs/20260504T141246Z-gh-runs-develop-no-new-run-after-codex-resume-guard-hardening-push.md}}. Kind: markdown. Size: 2837 bytes. Shape: 67 lines. Markdown document, 67 lines, 0 explicit URLs. Lead heading: GitHub Run Check After `Harden Codex resume restore guards`.\\
\textbf{\path{samples/logs/20260504T141246Z-gh-runs-develop-no-new-run-after-codex-resume-guard-hardening-push.md}}. Тип: markdown. Размер: 2837 байт. Форма: 67 строк. Документ Markdown, 67 строк, явных URL: 0. Главный заголовок: GitHub Run Check After `Harden Codex resume restore guards`.\par\medskip
\textbf{\path{samples/logs/20260504T161106Z-gh-runs-develop-no-new-run-after-explicit-codex-resume-targeting-push.md}}. Kind: markdown. Size: 1175 bytes. Shape: 27 lines. Markdown document, 27 lines, 0 explicit URLs. Lead heading: GitHub Actions inspection after `464e16b` push.\\
\textbf{\path{samples/logs/20260504T161106Z-gh-runs-develop-no-new-run-after-explicit-codex-resume-targeting-push.md}}. Тип: markdown. Размер: 1175 байт. Форма: 27 строк. Документ Markdown, 27 строк, явных URL: 0. Главный заголовок: GitHub Actions inspection after `464e16b` push.\par\medskip
\textbf{\path{samples/logs/20260504T180441Z-gh-runs-develop-no-new-run-after-codex-session-continuity-markers-push.md}}. Kind: markdown. Size: 986 bytes. Shape: 28 lines. Markdown document, 28 lines, 0 explicit URLs. Lead heading: GitHub Actions inspection after `490fe41` push.\\
\textbf{\path{samples/logs/20260504T180441Z-gh-runs-develop-no-new-run-after-codex-session-continuity-markers-push.md}}. Тип: markdown. Размер: 986 байт. Форма: 28 строк. Документ Markdown, 28 строк, явных URL: 0. Главный заголовок: GitHub Actions inspection after `490fe41` push.\par\medskip
\textbf{\path{samples/logs/20260505T072441Z-gh-runs-develop-no-new-run-after-mcp-first-codex-discovery-push.md}}. Kind: markdown. Size: 1152 bytes. Shape: 28 lines. Markdown document, 28 lines, 0 explicit URLs. Lead heading: GitHub Actions inspection after `abcf52c` push.\\
\textbf{\path{samples/logs/20260505T072441Z-gh-runs-develop-no-new-run-after-mcp-first-codex-discovery-push.md}}. Тип: markdown. Размер: 1152 байт. Форма: 28 строк. Документ Markdown, 28 строк, явных URL: 0. Главный заголовок: GitHub Actions inspection after `abcf52c` push.\par\medskip
\textbf{\path{samples/logs/README.md}}. Kind: markdown. Size: 331 bytes. Shape: 11 lines. Markdown document, 11 lines, 0 explicit URLs. Lead heading: GitHub Run Logs.\\
\textbf{\path{samples/logs/README.md}}. Тип: markdown. Размер: 331 байт. Форма: 11 строк. Документ Markdown, 11 строк, явных URL: 0. Главный заголовок: GitHub Run Logs.\par\medskip
\textbf{\path{samples/population/hushwheel_fixture_population_candidates.yaml}}. Kind: yaml. Size: 955 bytes. Shape: 22 lines. YAML data or workflow file, 22 lines, 0 explicit URLs. Top keys: no top-level mapping keys.\\
\textbf{\path{samples/population/hushwheel_fixture_population_candidates.yaml}}. Тип: yaml. Размер: 955 байт. Форма: 22 строк. Файл YAML с данными или workflow, 22 строк, явных URL: 0. Верхние ключи: no top-level mapping keys.\par\medskip
\textbf{\path{samples/population/repository_population_candidates.yaml}}. Kind: yaml. Size: 690 bytes. Shape: 16 lines. YAML data or workflow file, 16 lines, 0 explicit URLs. Top keys: no top-level mapping keys.\\
\textbf{\path{samples/population/repository_population_candidates.yaml}}. Тип: yaml. Размер: 690 байт. Форма: 16 строк. Файл YAML с данными или workflow, 16 строк, явных URL: 0. Верхние ключи: no top-level mapping keys.\par\medskip
\textbf{\path{samples/training/hushwheel_fixture_training_examples.yaml}}. Kind: yaml. Size: 1852 bytes. Shape: 68 lines. YAML data or workflow file, 68 lines, 0 explicit URLs. Top keys: no top-level mapping keys.\\
\textbf{\path{samples/training/hushwheel_fixture_training_examples.yaml}}. Тип: yaml. Размер: 1852 байт. Форма: 68 строк. Файл YAML с данными или workflow, 68 строк, явных URL: 0. Верхние ключи: no top-level mapping keys.\par\medskip
\textbf{\path{samples/training/repository_training_examples.yaml}}. Kind: yaml. Size: 2504 bytes. Shape: 80 lines. YAML data or workflow file, 80 lines, 0 explicit URLs. Top keys: no top-level mapping keys.\\
\textbf{\path{samples/training/repository_training_examples.yaml}}. Тип: yaml. Размер: 2504 байт. Форма: 80 строк. Файл YAML с данными или workflow, 80 строк, явных URL: 0. Верхние ключи: no top-level mapping keys.\par\medskip
\textbf{\path{src/repo_rag_lab/__init__.py}}. Kind: python. Size: 1087 bytes. Shape: 47 lines. Python module, 47 lines, 0 explicit URLs. Module intent: Public package surface for the repository-grounded RAG lab.\\
\textbf{\path{src/repo_rag_lab/__init__.py}}. Тип: python. Размер: 1087 байт. Форма: 47 строк. Модуль Python, 47 строк, явных URL: 0. Назначение модуля: Public package surface for the repository-grounded RAG lab.\par\medskip
\textbf{\path{src/repo_rag_lab/azure.py}}. Kind: python. Size: 2592 bytes. Shape: 76 lines. Python module, 76 lines, 0 explicit URLs. Module intent: Azure deployment manifest helpers for downstream release workflows.\\
\textbf{\path{src/repo_rag_lab/azure.py}}. Тип: python. Размер: 2592 байт. Форма: 76 строк. Модуль Python, 76 строк, явных URL: 0. Назначение модуля: Azure deployment manifest helpers for downstream release workflows.\par\medskip
\textbf{\path{src/repo_rag_lab/azure_artifacts.py}}. Kind: python. Size: 16790 bytes. Shape: 436 lines. Python module, 436 lines, 0 explicit URLs. Module intent: Azure Blob + Queue helpers for trainer-side traces and global bundle storage.\\
\textbf{\path{src/repo_rag_lab/azure_artifacts.py}}. Тип: python. Размер: 16790 байт. Форма: 436 строк. Модуль Python, 436 строк, явных URL: 0. Назначение модуля: Azure Blob + Queue helpers for trainer-side traces and global bundle storage.\par\medskip
\textbf{\path{src/repo_rag_lab/azure_runtime.py}}. Kind: python. Size: 20123 bytes. Shape: 577 lines. Python module, 577 lines, 0 explicit URLs. Module intent: Azure runtime normalization, live probes, and cloud-completion helpers.\\
\textbf{\path{src/repo_rag_lab/azure_runtime.py}}. Тип: python. Размер: 20123 байт. Форма: 577 строк. Модуль Python, 577 строк, явных URL: 0. Назначение модуля: Azure runtime normalization, live probes, and cloud-completion helpers.\par\medskip
\textbf{\path{src/repo_rag_lab/benchmarks.py}}. Kind: python. Size: 16524 bytes. Shape: 484 lines. Python module, 484 lines, 0 explicit URLs. Module intent: Retrieval benchmark helpers for notebook assertions and logging.\\
\textbf{\path{src/repo_rag_lab/benchmarks.py}}. Тип: python. Размер: 16524 байт. Форма: 484 строк. Модуль Python, 484 строк, явных URL: 0. Назначение модуля: Retrieval benchmark helpers for notebook assertions and logging.\par\medskip
\textbf{\path{src/repo_rag_lab/cli.py}}. Kind: python. Size: 58870 bytes. Shape: 1339 lines. Python module, 1339 lines, 0 explicit URLs. Module intent: Command-line entrypoints for the shared repository RAG workflows.\\
\textbf{\path{src/repo_rag_lab/cli.py}}. Тип: python. Размер: 58870 байт. Форма: 1339 строк. Модуль Python, 1339 строк, явных URL: 0. Назначение модуля: Command-line entrypoints for the shared repository RAG workflows.\par\medskip
\textbf{\path{src/repo_rag_lab/codex_proxy.py}}. Kind: python. Size: 75460 bytes. Shape: 1946 lines. Python module, 1946 lines, 0 explicit URLs. Module intent: Local streaming proxy that injects repo-RAG + DSPy mediation for Codex.\\
\textbf{\path{src/repo_rag_lab/codex_proxy.py}}. Тип: python. Размер: 75460 байт. Форма: 1946 строк. Модуль Python, 1946 строк, явных URL: 0. Назначение модуля: Local streaming proxy that injects repo-RAG + DSPy mediation for Codex.\par\medskip
\textbf{\path{src/repo_rag_lab/corpus.py}}. Kind: python. Size: 5245 bytes. Shape: 181 lines. Python module, 181 lines, 0 explicit URLs. Module intent: Repository file loading helpers for the baseline text corpus.\\
\textbf{\path{src/repo_rag_lab/corpus.py}}. Тип: python. Размер: 5245 байт. Форма: 181 строк. Модуль Python, 181 строк, явных URL: 0. Назначение модуля: Repository file loading helpers for the baseline text corpus.\par\medskip
\textbf{\path{src/repo_rag_lab/dspy_training.py}}. Kind: python. Size: 61618 bytes. Shape: 1612 lines. Python module, 1612 lines, 0 explicit URLs. Module intent: DSPy training and artifact helpers for repository-grounded RAG.\\
\textbf{\path{src/repo_rag_lab/dspy_training.py}}. Тип: python. Размер: 61618 байт. Форма: 1612 строк. Модуль Python, 1612 строк, явных URL: 0. Назначение модуля: DSPy training and artifact helpers for repository-grounded RAG.\par\medskip
\textbf{\path{src/repo_rag_lab/dspy_workflow.py}}. Kind: python. Size: 6518 bytes. Shape: 180 lines. Python module, 180 lines, 0 explicit URLs. Module intent: DSPy-backed runtime wrappers around the repository retrieval workflow.\\
\textbf{\path{src/repo_rag_lab/dspy_workflow.py}}. Тип: python. Размер: 6518 байт. Форма: 180 строк. Модуль Python, 180 строк, явных URL: 0. Назначение модуля: DSPy-backed runtime wrappers around the repository retrieval workflow.\par\medskip
\textbf{\path{src/repo_rag_lab/exploratorium_translation.py}}. Kind: python. Size: 38049 bytes. Shape: 1009 lines. Python module, 1009 lines, 0 explicit URLs. Module intent: Generate the bilingual exploratorium translation publication surfaces.\\
\textbf{\path{src/repo_rag_lab/exploratorium_translation.py}}. Тип: python. Размер: 38049 байт. Форма: 1009 строк. Модуль Python, 1009 строк, явных URL: 0. Назначение модуля: Generate the bilingual exploratorium translation publication surfaces.\par\medskip
\textbf{\path{src/repo_rag_lab/file_summaries.py}}. Kind: python. Size: 10564 bytes. Shape: 327 lines. Python module, 327 lines, 0 explicit URLs. Module intent: Generate tracked repository file-summary surfaces.\\
\textbf{\path{src/repo_rag_lab/file_summaries.py}}. Тип: python. Размер: 10564 байт. Форма: 327 строк. Модуль Python, 327 строк, явных URL: 0. Назначение модуля: Generate tracked repository file-summary surfaces.\par\medskip
\textbf{\path{src/repo_rag_lab/github_pr_gates.py}}. Kind: python. Size: 3979 bytes. Shape: 140 lines. Python module, 140 lines, 0 explicit URLs. Module intent: GitHub branch-protection helpers for required pull-request gate checks.\\
\textbf{\path{src/repo_rag_lab/github_pr_gates.py}}. Тип: python. Размер: 3979 байт. Форма: 140 строк. Модуль Python, 140 строк, явных URL: 0. Назначение модуля: GitHub branch-protection helpers for required pull-request gate checks.\par\medskip
\textbf{\path{src/repo_rag_lab/mcp.py}}. Kind: python. Size: 2947 bytes. Shape: 84 lines. Python module, 84 lines, 0 explicit URLs. Module intent: Discovery helpers for MCP-related repository artifacts.\\
\textbf{\path{src/repo_rag_lab/mcp.py}}. Тип: python. Размер: 2947 байт. Форма: 84 строк. Модуль Python, 84 строк, явных URL: 0. Назначение модуля: Discovery helpers for MCP-related repository artifacts.\par\medskip
\textbf{\path{src/repo_rag_lab/mcp_server.py}}. Kind: python. Size: 50399 bytes. Shape: 1361 lines. Python module, 1361 lines, 0 explicit URLs. Module intent: Minimal stdio MCP server for bounded repo-RAG operations.\\
\textbf{\path{src/repo_rag_lab/mcp_server.py}}. Тип: python. Размер: 50399 байт. Форма: 1361 строк. Модуль Python, 1361 строк, явных URL: 0. Назначение модуля: Minimal stdio MCP server for bounded repo-RAG operations.\par\medskip
\textbf{\path{src/repo_rag_lab/mcp_stdio.py}}. Kind: python. Size: 953 bytes. Shape: 38 lines. Python module, 38 lines, 0 explicit URLs. Module intent: Lightweight stdio entrypoint for the bounded repo-RAG MCP server.\\
\textbf{\path{src/repo_rag_lab/mcp_stdio.py}}. Тип: python. Размер: 953 байт. Форма: 38 строк. Модуль Python, 38 строк, явных URL: 0. Назначение модуля: Lightweight stdio entrypoint for the bounded repo-RAG MCP server.\par\medskip
\textbf{\path{src/repo_rag_lab/notebook_runner.py}}. Kind: python. Size: 16433 bytes. Shape: 491 lines. Python module, 491 lines, 0 explicit URLs. Module intent: Monitored batch execution helpers for repository notebooks.\\
\textbf{\path{src/repo_rag_lab/notebook_runner.py}}. Тип: python. Размер: 16433 байт. Форма: 491 строк. Модуль Python, 491 строк, явных URL: 0. Назначение модуля: Monitored batch execution helpers for repository notebooks.\par\medskip
\textbf{\path{src/repo_rag_lab/notebook_scaffolding.py}}. Kind: python. Size: 10965 bytes. Shape: 258 lines. Python module, 258 lines, 0 explicit URLs. Module intent: High-level notebook scaffolds built from tested repository helpers.\\
\textbf{\path{src/repo_rag_lab/notebook_scaffolding.py}}. Тип: python. Размер: 10965 байт. Форма: 258 строк. Модуль Python, 258 строк, явных URL: 0. Назначение модуля: High-level notebook scaffolds built from tested repository helpers.\par\medskip
\textbf{\path{src/repo_rag_lab/notebook_support.py}}. Kind: python. Size: 2843 bytes. Shape: 89 lines. Python module, 89 lines, 0 explicit URLs. Module intent: Helpers shared by notebooks so their setup logic stays in tested Python code.\\
\textbf{\path{src/repo_rag_lab/notebook_support.py}}. Тип: python. Размер: 2843 байт. Форма: 89 строк. Модуль Python, 89 строк, явных URL: 0. Назначение модуля: Helpers shared by notebooks so their setup logic stays in tested Python code.\par\medskip
\textbf{\path{src/repo_rag_lab/pages_site.py}}. Kind: python. Size: 13635 bytes. Shape: 390 lines. Python module, 390 lines, 2 explicit URLs. Module intent: Generate a MkDocs-ready catalog of tracked repository Markdown files.\\
\textbf{\path{src/repo_rag_lab/pages_site.py}}. Тип: python. Размер: 13635 байт. Форма: 390 строк. Модуль Python, 390 строк, явных URL: 2. Назначение модуля: Generate a MkDocs-ready catalog of tracked repository Markdown files.\par\medskip
\textbf{\path{src/repo_rag_lab/population_samples.py}}. Kind: python. Size: 8376 bytes. Shape: 241 lines. Python module, 241 lines, 0 explicit URLs. Module intent: Helpers for reviewing and batching starter corpus-population inputs.\\
\textbf{\path{src/repo_rag_lab/population_samples.py}}. Тип: python. Размер: 8376 байт. Форма: 241 строк. Модуль Python, 241 строк, явных URL: 0. Назначение модуля: Helpers for reviewing and batching starter corpus-population inputs.\par\medskip
\textbf{\path{src/repo_rag_lab/retrieval.py}}. Kind: python. Size: 25533 bytes. Shape: 788 lines. Python module, 788 lines, 0 explicit URLs. Module intent: Baseline chunking and lexical retrieval utilities for repository text.\\
\textbf{\path{src/repo_rag_lab/retrieval.py}}. Тип: python. Размер: 25533 байт. Форма: 788 строк. Модуль Python, 788 строк, явных URL: 0. Назначение модуля: Baseline chunking and lexical retrieval utilities for repository text.\par\medskip
\textbf{\path{src/repo_rag_lab/retrieval_profile.py}}. Kind: python. Size: 9527 bytes. Shape: 236 lines. Python module, 236 lines, 0 explicit URLs. Module intent: Profile-driven retrieval weighting, exclusions, and repo-local overrides.\\
\textbf{\path{src/repo_rag_lab/retrieval_profile.py}}. Тип: python. Размер: 9527 байт. Форма: 236 строк. Модуль Python, 236 строк, явных URL: 0. Назначение модуля: Profile-driven retrieval weighting, exclusions, and repo-local overrides.\par\medskip
\textbf{\path{src/repo_rag_lab/runtime_artifacts.py}}. Kind: python. Size: 118360 bytes. Shape: 2671 lines. Python module, 2671 lines, 0 explicit URLs. Module intent: Versioned bundle, overlay, and runtime-trace helpers.\\
\textbf{\path{src/repo_rag_lab/runtime_artifacts.py}}. Тип: python. Размер: 118360 байт. Форма: 2671 строк. Модуль Python, 2671 строк, явных URL: 0. Назначение модуля: Versioned bundle, overlay, and runtime-trace helpers.\par\medskip
\textbf{\path{src/repo_rag_lab/rust_lookup.py}}. Kind: python. Size: 7724 bytes. Shape: 295 lines. Python module, 295 lines, 0 explicit URLs. Module intent: Helpers for the native Rust-backed SQLite lookup index.\\
\textbf{\path{src/repo_rag_lab/rust_lookup.py}}. Тип: python. Размер: 7724 байт. Форма: 295 строк. Модуль Python, 295 строк, явных URL: 0. Назначение модуля: Helpers for the native Rust-backed SQLite lookup index.\par\medskip
\textbf{\path{src/repo_rag_lab/semantic_retrieval.py}}. Kind: python. Size: 7407 bytes. Shape: 214 lines. Python module, 214 lines, 0 explicit URLs. Module intent: Azure OpenAI embedding-backed semantic retrieval helpers.\\
\textbf{\path{src/repo_rag_lab/semantic_retrieval.py}}. Тип: python. Размер: 7407 байт. Форма: 214 строк. Модуль Python, 214 строк, явных URL: 0. Назначение модуля: Azure OpenAI embedding-backed semantic retrieval helpers.\par\medskip
\textbf{\path{src/repo_rag_lab/settings.py}}. Kind: python. Size: 699 bytes. Shape: 28 lines. Python module, 28 lines, 0 explicit URLs. Module intent: Shared repository path settings derived from a repository root.\\
\textbf{\path{src/repo_rag_lab/settings.py}}. Тип: python. Размер: 699 байт. Форма: 28 строк. Модуль Python, 28 строк, явных URL: 0. Назначение модуля: Shared repository path settings derived from a repository root.\par\medskip
\textbf{\path{src/repo_rag_lab/todo_backlog.py}}. Kind: python. Size: 8434 bytes. Shape: 276 lines. Python module, 276 lines, 1 explicit URLs. Module intent: Shared backlog rendering helpers for Markdown and LaTeX surfaces.\\
\textbf{\path{src/repo_rag_lab/todo_backlog.py}}. Тип: python. Размер: 8434 байт. Форма: 276 строк. Модуль Python, 276 строк, явных URL: 1. Назначение модуля: Shared backlog rendering helpers for Markdown and LaTeX surfaces.\par\medskip
\textbf{\path{src/repo_rag_lab/trainer_deployment.py}}. Kind: python. Size: 18027 bytes. Shape: 431 lines. Python module, 431 lines, 0 explicit URLs. Module intent: Kubernetes manifest helpers for trainer-side repo-RAG deployment surfaces.\\
\textbf{\path{src/repo_rag_lab/trainer_deployment.py}}. Тип: python. Размер: 18027 байт. Форма: 431 строк. Модуль Python, 431 строк, явных URL: 0. Назначение модуля: Kubernetes manifest helpers for trainer-side repo-RAG deployment surfaces.\par\medskip
\textbf{\path{src/repo_rag_lab/training_samples.py}}. Kind: python. Size: 110198 bytes. Shape: 2709 lines. Python module, 2709 lines, 0 explicit URLs. Module intent: Helpers for loading and summarizing starter DSPy training examples.\\
\textbf{\path{src/repo_rag_lab/training_samples.py}}. Тип: python. Размер: 110198 байт. Форма: 2709 строк. Модуль Python, 2709 строк, явных URL: 0. Назначение модуля: Helpers for loading and summarizing starter DSPy training examples.\par\medskip
\textbf{\path{src/repo_rag_lab/utilities.py}}. Kind: python. Size: 116367 bytes. Shape: 2809 lines. Python module, 2809 lines, 1 explicit URLs. Module intent: User-facing utility helpers shared by the CLI, tests, and notebooks.\\
\textbf{\path{src/repo_rag_lab/utilities.py}}. Тип: python. Размер: 116367 байт. Форма: 2809 строк. Модуль Python, 2809 строк, явных URL: 1. Назначение модуля: User-facing utility helpers shared by the CLI, tests, and notebooks.\par\medskip
\textbf{\path{src/repo_rag_lab/verification.py}}. Kind: python. Size: 4476 bytes. Shape: 163 lines. Python module, 163 lines, 0 explicit URLs. Module intent: Verification helpers for Makefile and notebook contract checks.\\
\textbf{\path{src/repo_rag_lab/verification.py}}. Тип: python. Размер: 4476 байт. Форма: 163 строк. Модуль Python, 163 строк, явных URL: 0. Назначение модуля: Verification helpers for Makefile and notebook contract checks.\par\medskip
\textbf{\path{src/repo_rag_lab/workflow.py}}. Kind: python. Size: 15900 bytes. Shape: 503 lines. Python module, 503 lines, 0 explicit URLs. Module intent: Baseline repository question-answering workflow built on file retrieval.\\
\textbf{\path{src/repo_rag_lab/workflow.py}}. Тип: python. Размер: 15900 байт. Форма: 503 строк. Модуль Python, 503 строк, явных URL: 0. Назначение модуля: Baseline repository question-answering workflow built on file retrieval.\par\medskip
\textbf{\path{tests/features/repository_rag.feature}}. Kind: generic. Size: 377 bytes. Shape: 12 lines. Tracked repository file, 12 lines, 0 explicit URLs. Opening line: Feature: Repository RAG\\
\textbf{\path{tests/features/repository_rag.feature}}. Тип: generic. Размер: 377 байт. Форма: 12 строк. Отслеживаемый файл репозитория, 12 строк, явных URL: 0. Первая строка: Feature: Repository RAG\par\medskip
\textbf{\path{tests/fixtures/hushwheel_lexiconarium/.gitignore}}. Kind: generic. Size: 23 bytes. Shape: 4 lines. Tracked repository file, 4 lines, 0 explicit URLs. Opening line: hushwheel\\
\textbf{\path{tests/fixtures/hushwheel_lexiconarium/.gitignore}}. Тип: generic. Размер: 23 байт. Форма: 4 строк. Отслеживаемый файл репозитория, 4 строк, явных URL: 0. Первая строка: hushwheel\par\medskip
\textbf{\path{tests/fixtures/hushwheel_lexiconarium/CHANGELOG.md}}. Kind: markdown. Size: 345 bytes. Shape: 9 lines. Markdown document, 9 lines, 0 explicit URLs. Lead heading: Changelog.\\
\textbf{\path{tests/fixtures/hushwheel_lexiconarium/CHANGELOG.md}}. Тип: markdown. Размер: 345 байт. Форма: 9 строк. Документ Markdown, 9 строк, явных URL: 0. Главный заголовок: Changelog.\par\medskip
\textbf{\path{tests/fixtures/hushwheel_lexiconarium/Doxyfile}}. Kind: generic. Size: 1239 bytes. Shape: 35 lines. Tracked repository file, 35 lines, 0 explicit URLs. Opening line: PROJECT\_NAME           = Hushwheel Lexiconarium\\
\textbf{\path{tests/fixtures/hushwheel_lexiconarium/Doxyfile}}. Тип: generic. Размер: 1239 байт. Форма: 35 строк. Отслеживаемый файл репозитория, 35 строк, явных URL: 0. Первая строка: PROJECT\_NAME           = Hushwheel Lexiconarium\par\medskip
\textbf{\path{tests/fixtures/hushwheel_lexiconarium/LICENSE.md}}. Kind: markdown. Size: 1082 bytes. Shape: 22 lines. Markdown document, 22 lines, 0 explicit URLs. Lead heading: LICENSE.\\
\textbf{\path{tests/fixtures/hushwheel_lexiconarium/LICENSE.md}}. Тип: markdown. Размер: 1082 байт. Форма: 22 строк. Документ Markdown, 22 строк, явных URL: 0. Главный заголовок: LICENSE.\par\medskip
\textbf{\path{tests/fixtures/hushwheel_lexiconarium/Makefile}}. Kind: make. Size: 20908 bytes. Shape: 388 lines. Repository Makefile, 388 lines, 0 explicit URLs. It exposes the shared automation surfaces.\\
\textbf{\path{tests/fixtures/hushwheel_lexiconarium/Makefile}}. Тип: make. Размер: 20908 байт. Форма: 388 строк. Makefile репозитория, 388 строк, явных URL: 0. Он открывает общие поверхности автоматизации.\par\medskip
\textbf{\path{tests/fixtures/hushwheel_lexiconarium/README.md}}. Kind: markdown. Size: 8269 bytes. Shape: 162 lines. Markdown document, 162 lines, 0 explicit URLs. Lead heading: The Hushwheel Lexiconarium.\\
\textbf{\path{tests/fixtures/hushwheel_lexiconarium/README.md}}. Тип: markdown. Размер: 8269 байт. Форма: 162 строк. Документ Markdown, 162 строк, явных URL: 0. Главный заголовок: The Hushwheel Lexiconarium.\par\medskip
\textbf{\path{tests/fixtures/hushwheel_lexiconarium/VERSION}}. Kind: generic. Size: 6 bytes. Shape: 2 lines. Tracked repository file, 2 lines, 0 explicit URLs. Opening line: 0.1.0\\
\textbf{\path{tests/fixtures/hushwheel_lexiconarium/VERSION}}. Тип: generic. Размер: 6 байт. Форма: 2 строк. Отслеживаемый файл репозитория, 2 строк, явных URL: 0. Первая строка: 0.1.0\par\medskip
\textbf{\path{tests/fixtures/hushwheel_lexiconarium/docs/architecture.md}}. Kind: markdown. Size: 1404 bytes. Shape: 33 lines. Markdown document, 33 lines, 0 explicit URLs. Lead heading: Hushwheel Architecture.\\
\textbf{\path{tests/fixtures/hushwheel_lexiconarium/docs/architecture.md}}. Тип: markdown. Размер: 1404 байт. Форма: 33 строк. Документ Markdown, 33 строк, явных URL: 0. Главный заголовок: Hushwheel Architecture.\par\medskip
\textbf{\path{tests/fixtures/hushwheel_lexiconarium/docs/catalog.md}}. Kind: markdown. Size: 27404 bytes. Shape: 357 lines. Markdown document, 357 lines, 0 explicit URLs. Lead heading: Hushwheel Representative Catalog.\\
\textbf{\path{tests/fixtures/hushwheel_lexiconarium/docs/catalog.md}}. Тип: markdown. Размер: 27404 байт. Форма: 357 строк. Документ Markdown, 357 строк, явных URL: 0. Главный заголовок: Hushwheel Representative Catalog.\par\medskip
\textbf{\path{tests/fixtures/hushwheel_lexiconarium/docs/concepts.md}}. Kind: markdown. Size: 1559 bytes. Shape: 37 lines. Markdown document, 37 lines, 0 explicit URLs. Lead heading: Hushwheel Concepts.\\
\textbf{\path{tests/fixtures/hushwheel_lexiconarium/docs/concepts.md}}. Тип: markdown. Размер: 1559 байт. Форма: 37 строк. Документ Markdown, 37 строк, явных URL: 0. Главный заголовок: Hushwheel Concepts.\par\medskip
\textbf{\path{tests/fixtures/hushwheel_lexiconarium/docs/constellation-atlas.md}}. Kind: markdown. Size: 5111 bytes. Shape: 142 lines. Markdown document, 142 lines, 0 explicit URLs. Lead heading: Hushwheel Constellation Atlas §::<[]>.\\
\textbf{\path{tests/fixtures/hushwheel_lexiconarium/docs/constellation-atlas.md}}. Тип: markdown. Размер: 5111 байт. Форма: 142 строк. Документ Markdown, 142 строк, явных URL: 0. Главный заголовок: Hushwheel Constellation Atlas §::<[]>.\par\medskip
\textbf{\path{tests/fixtures/hushwheel_lexiconarium/docs/districts.md}}. Kind: markdown. Size: 1151 bytes. Shape: 42 lines. Markdown document, 42 lines, 0 explicit URLs. Lead heading: Hushwheel Districts.\\
\textbf{\path{tests/fixtures/hushwheel_lexiconarium/docs/districts.md}}. Тип: markdown. Размер: 1151 байт. Форма: 42 строк. Документ Markdown, 42 строк, явных URL: 0. Главный заголовок: Hushwheel Districts.\par\medskip
\textbf{\path{tests/fixtures/hushwheel_lexiconarium/docs/doxygen-fonts.sty}}. Kind: generic. Size: 212 bytes. Shape: binary surface. Tracked repository file, 212 bytes of binary data, 0 explicit URLs. Opening line: No non-empty preview line.\\
\textbf{\path{tests/fixtures/hushwheel_lexiconarium/docs/doxygen-fonts.sty}}. Тип: generic. Размер: 212 байт. Форма: бинарная поверхность. Отслеживаемый файл репозитория, 212 байт бинарных данных, явных URL: 0. Первая строка: No non-empty preview line.\par\medskip
\textbf{\path{tests/fixtures/hushwheel_lexiconarium/docs/hushwheel-reference.pdf}}. Kind: binary. Size: 132 bytes. Shape: binary surface. Binary asset, 132 bytes of binary data. It is tracked for stable publication output.\\
\textbf{\path{tests/fixtures/hushwheel_lexiconarium/docs/hushwheel-reference.pdf}}. Тип: binary. Размер: 132 байт. Форма: бинарная поверхность. Бинарный артефакт, 132 байт бинарных данных. Он хранится в Git для стабильного публикационного вывода.\par\medskip
\textbf{\path{tests/fixtures/hushwheel_lexiconarium/docs/operations.md}}. Kind: markdown. Size: 1068 bytes. Shape: 35 lines. Markdown document, 35 lines, 0 explicit URLs. Lead heading: Hushwheel Operations.\\
\textbf{\path{tests/fixtures/hushwheel_lexiconarium/docs/operations.md}}. Тип: markdown. Размер: 1068 байт. Форма: 35 строк. Документ Markdown, 35 строк, явных URL: 0. Главный заголовок: Hushwheel Operations.\par\medskip
\textbf{\path{tests/fixtures/hushwheel_lexiconarium/docs/packaging.md}}. Kind: markdown. Size: 1112 bytes. Shape: 33 lines. Markdown document, 33 lines, 0 explicit URLs. Lead heading: Hushwheel Packaging.\\
\textbf{\path{tests/fixtures/hushwheel_lexiconarium/docs/packaging.md}}. Тип: markdown. Размер: 1112 байт. Форма: 33 строк. Документ Markdown, 33 строк, явных URL: 0. Главный заголовок: Hushwheel Packaging.\par\medskip
\textbf{\path{tests/fixtures/hushwheel_lexiconarium/docs/testing.md}}. Kind: markdown. Size: 3507 bytes. Shape: 80 lines. Markdown document, 80 lines, 0 explicit URLs. Lead heading: Hushwheel Testing.\\
\textbf{\path{tests/fixtures/hushwheel_lexiconarium/docs/testing.md}}. Тип: markdown. Размер: 3507 байт. Форма: 80 строк. Документ Markdown, 80 строк, явных URL: 0. Главный заголовок: Hushwheel Testing.\par\medskip
\textbf{\path{tests/fixtures/hushwheel_lexiconarium/fixture-manifest.json}}. Kind: json. Size: 1098 bytes. Shape: 49 lines. JSON data file, 49 lines, 0 explicit URLs. Top keys: entry\_count, source\_size\_bytes, doc\_files.\\
\textbf{\path{tests/fixtures/hushwheel_lexiconarium/fixture-manifest.json}}. Тип: json. Размер: 1098 байт. Форма: 49 строк. Файл JSON, 49 строк, явных URL: 0. Верхние ключи: entry\_count, source\_size\_bytes, doc\_files.\par\medskip
\textbf{\path{tests/fixtures/hushwheel_lexiconarium/include/hushwheel.h}}. Kind: source. Size: 404 bytes. Shape: 19 lines. Source file, 19 lines, 0 explicit URLs. Opening line: \#ifndef HUSHWHEEL\_H\\
\textbf{\path{tests/fixtures/hushwheel_lexiconarium/include/hushwheel.h}}. Тип: source. Размер: 404 байт. Форма: 19 строк. Исходный файл, 19 строк, явных URL: 0. Первая строка: \#ifndef HUSHWHEEL\_H\par\medskip
\textbf{\path{tests/fixtures/hushwheel_lexiconarium/packaging/hushwheel.1}}. Kind: generic. Size: 1303 bytes. Shape: binary surface. Tracked repository file, 1303 bytes of binary data, 0 explicit URLs. Opening line: No non-empty preview line.\\
\textbf{\path{tests/fixtures/hushwheel_lexiconarium/packaging/hushwheel.1}}. Тип: generic. Размер: 1303 байт. Форма: бинарная поверхность. Отслеживаемый файл репозитория, 1303 байт бинарных данных, явных URL: 0. Первая строка: No non-empty preview line.\par\medskip
\textbf{\path{tests/fixtures/hushwheel_lexiconarium/packaging/hushwheel.package.json}}. Kind: json. Size: 591 bytes. Shape: 29 lines. JSON data file, 29 lines, 0 explicit URLs. Top keys: name, version, summary.\\
\textbf{\path{tests/fixtures/hushwheel_lexiconarium/packaging/hushwheel.package.json}}. Тип: json. Размер: 591 байт. Форма: 29 строк. Файл JSON, 29 строк, явных URL: 0. Верхние ключи: name, version, summary.\par\medskip
\textbf{\path{tests/fixtures/hushwheel_lexiconarium/src/hushwheel.c}}. Kind: source. Size: 1694579 bytes. Shape: 4375 lines. Source file, 4375 lines, 0 explicit URLs. Opening line: /**\\
\textbf{\path{tests/fixtures/hushwheel_lexiconarium/src/hushwheel.c}}. Тип: source. Размер: 1694579 байт. Форма: 4375 строк. Исходный файл, 4375 строк, явных URL: 0. Первая строка: /**\par\medskip
\textbf{\path{tests/fixtures/hushwheel_lexiconarium/src/hushwheel_internal.h}}. Kind: source. Size: 603 bytes. Shape: 25 lines. Source file, 25 lines, 0 explicit URLs. Opening line: \#ifndef HUSHWHEEL\_INTERNAL\_H\\
\textbf{\path{tests/fixtures/hushwheel_lexiconarium/src/hushwheel_internal.h}}. Тип: source. Размер: 603 байт. Форма: 25 строк. Исходный файл, 25 строк, явных URL: 0. Первая строка: \#ifndef HUSHWHEEL\_INTERNAL\_H\par\medskip
\textbf{\path{tests/fixtures/hushwheel_lexiconarium/src/hushwheel_spoke_argent.c}}. Kind: source. Size: 324091 bytes. Shape: 525 lines. Source file, 525 lines, 0 explicit URLs. Opening line: \#include "hushwheel\_internal.h"\\
\textbf{\path{tests/fixtures/hushwheel_lexiconarium/src/hushwheel_spoke_argent.c}}. Тип: source. Размер: 324091 байт. Форма: 525 строк. Исходный файл, 525 строк, явных URL: 0. Первая строка: \#include "hushwheel\_internal.h"\par\medskip
\textbf{\path{tests/fixtures/hushwheel_lexiconarium/src/hushwheel_spoke_bhagavatam.c}}. Kind: source. Size: 326160 bytes. Shape: 525 lines. Source file, 525 lines, 0 explicit URLs. Opening line: \#include "hushwheel\_internal.h"\\
\textbf{\path{tests/fixtures/hushwheel_lexiconarium/src/hushwheel_spoke_bhagavatam.c}}. Тип: source. Размер: 326160 байт. Форма: 525 строк. Исходный файл, 525 строк, явных URL: 0. Первая строка: \#include "hushwheel\_internal.h"\par\medskip
\textbf{\path{tests/fixtures/hushwheel_lexiconarium/src/hushwheel_spoke_crosscanon.c}}. Kind: source. Size: 326144 bytes. Shape: 525 lines. Source file, 525 lines, 0 explicit URLs. Opening line: \#include "hushwheel\_internal.h"\\
\textbf{\path{tests/fixtures/hushwheel_lexiconarium/src/hushwheel_spoke_crosscanon.c}}. Тип: source. Размер: 326144 байт. Форма: 525 строк. Исходный файл, 525 строк, явных URL: 0. Первая строка: \#include "hushwheel\_internal.h"\par\medskip
\textbf{\path{tests/fixtures/hushwheel_lexiconarium/src/hushwheel_spoke_editorial.c}}. Kind: source. Size: 325628 bytes. Shape: 525 lines. Source file, 525 lines, 0 explicit URLs. Opening line: \#include "hushwheel\_internal.h"\\
\textbf{\path{tests/fixtures/hushwheel_lexiconarium/src/hushwheel_spoke_editorial.c}}. Тип: source. Размер: 325628 байт. Форма: 525 строк. Исходный файл, 525 строк, явных URL: 0. Первая строка: \#include "hushwheel\_internal.h"\par\medskip
\textbf{\path{tests/fixtures/hushwheel_lexiconarium/src/hushwheel_spoke_gita.c}}. Kind: source. Size: 323055 bytes. Shape: 525 lines. Source file, 525 lines, 0 explicit URLs. Opening line: \#include "hushwheel\_internal.h"\\
\textbf{\path{tests/fixtures/hushwheel_lexiconarium/src/hushwheel_spoke_gita.c}}. Тип: source. Размер: 323055 байт. Форма: 525 строк. Исходный файл, 525 строк, явных URL: 0. Первая строка: \#include "hushwheel\_internal.h"\par\medskip
\textbf{\path{tests/fixtures/hushwheel_lexiconarium/src/hushwheel_spoke_kernel.c}}. Kind: source. Size: 324068 bytes. Shape: 525 lines. Source file, 525 lines, 0 explicit URLs. Opening line: \#include "hushwheel\_internal.h"\\
\textbf{\path{tests/fixtures/hushwheel_lexiconarium/src/hushwheel_spoke_kernel.c}}. Тип: source. Размер: 324068 байт. Форма: 525 строк. Исходный файл, 525 строк, явных URL: 0. Первая строка: \#include "hushwheel\_internal.h"\par\medskip
\textbf{\path{tests/fixtures/hushwheel_lexiconarium/src/hushwheel_spoke_psalter.c}}. Kind: source. Size: 324600 bytes. Shape: 525 lines. Source file, 525 lines, 0 explicit URLs. Opening line: \#include "hushwheel\_internal.h"\\
\textbf{\path{tests/fixtures/hushwheel_lexiconarium/src/hushwheel_spoke_psalter.c}}. Тип: source. Размер: 324600 байт. Форма: 525 строк. Исходный файл, 525 строк, явных URL: 0. Первая строка: \#include "hushwheel\_internal.h"\par\medskip
\textbf{\path{tests/fixtures/hushwheel_lexiconarium/src/hushwheel_spoke_ranger.c}}. Kind: source. Size: 324105 bytes. Shape: 525 lines. Source file, 525 lines, 0 explicit URLs. Opening line: \#include "hushwheel\_internal.h"\\
\textbf{\path{tests/fixtures/hushwheel_lexiconarium/src/hushwheel_spoke_ranger.c}}. Тип: source. Размер: 324105 байт. Форма: 525 строк. Исходный файл, 525 строк, явных URL: 0. Первая строка: \#include "hushwheel\_internal.h"\par\medskip
\textbf{\path{tests/fixtures/hushwheel_lexiconarium/src/hushwheel_spokes.c}}. Kind: source. Size: 1074 bytes. Shape: 25 lines. Source file, 25 lines, 0 explicit URLs. Opening line: \#include "hushwheel\_internal.h"\\
\textbf{\path{tests/fixtures/hushwheel_lexiconarium/src/hushwheel_spokes.c}}. Тип: source. Размер: 1074 байт. Форма: 25 строк. Исходный файл, 25 строк, явных URL: 0. Первая строка: \#include "hushwheel\_internal.h"\par\medskip
\textbf{\path{tests/fixtures/hushwheel_lexiconarium/tests/bdd/hushwheel.feature}}. Kind: generic. Size: 1021 bytes. Shape: 28 lines. Tracked repository file, 28 lines, 0 explicit URLs. Opening line: Feature: Hushwheel CLI\\
\textbf{\path{tests/fixtures/hushwheel_lexiconarium/tests/bdd/hushwheel.feature}}. Тип: generic. Размер: 1021 байт. Форма: 28 строк. Отслеживаемый файл репозитория, 28 строк, явных URL: 0. Первая строка: Feature: Hushwheel CLI\par\medskip
\textbf{\path{tests/fixtures/hushwheel_lexiconarium/tests/bdd/run_bdd.py}}. Kind: python. Size: 3146 bytes. Shape: 94 lines. Python module, 94 lines, 0 explicit URLs. Module intent: No module docstring detected.\\
\textbf{\path{tests/fixtures/hushwheel_lexiconarium/tests/bdd/run_bdd.py}}. Тип: python. Размер: 3146 байт. Форма: 94 строк. Модуль Python, 94 строк, явных URL: 0. Назначение модуля: No module docstring detected.\par\medskip
\textbf{\path{tests/fixtures/hushwheel_lexiconarium/tests/integration/cli_suite.py}}. Kind: python. Size: 4482 bytes. Shape: 122 lines. Python module, 122 lines, 0 explicit URLs. Module intent: No module docstring detected.\\
\textbf{\path{tests/fixtures/hushwheel_lexiconarium/tests/integration/cli_suite.py}}. Тип: python. Размер: 4482 байт. Форма: 122 строк. Модуль Python, 122 строк, явных URL: 0. Назначение модуля: No module docstring detected.\par\medskip
\textbf{\path{tests/fixtures/hushwheel_lexiconarium/tests/unit/test_hushwheel_unit.c}}. Kind: source. Size: 2002 bytes. Shape: 62 lines. Source file, 62 lines, 0 explicit URLs. Opening line: \#include <assert.h>\\
\textbf{\path{tests/fixtures/hushwheel_lexiconarium/tests/unit/test_hushwheel_unit.c}}. Тип: source. Размер: 2002 байт. Форма: 62 строк. Исходный файл, 62 строк, явных URL: 0. Первая строка: \#include <assert.h>\par\medskip
\textbf{\path{tests/fixtures/hushwheel_lexiconarium/tools/lint_hushwheel.py}}. Kind: python. Size: 4678 bytes. Shape: 143 lines. Python module, 143 lines, 0 explicit URLs. Module intent: No module docstring detected.\\
\textbf{\path{tests/fixtures/hushwheel_lexiconarium/tools/lint_hushwheel.py}}. Тип: python. Размер: 4678 байт. Форма: 143 строк. Модуль Python, 143 строк, явных URL: 0. Назначение модуля: No module docstring detected.\par\medskip
\textbf{\path{tests/fixtures/hushwheel_lexiconarium/tools/regenerate_hushwheel_fixture.py}}. Kind: python. Size: 31375 bytes. Shape: 866 lines. Python module, 866 lines, 0 explicit URLs. Module intent: No module docstring detected.\\
\textbf{\path{tests/fixtures/hushwheel_lexiconarium/tools/regenerate_hushwheel_fixture.py}}. Тип: python. Размер: 31375 байт. Форма: 866 строк. Модуль Python, 866 строк, явных URL: 0. Назначение модуля: No module docstring detected.\par\medskip
\textbf{\path{tests/test_azure_runtime.py}}. Kind: python. Size: 7540 bytes. Shape: 208 lines. Python module, 208 lines, 12 explicit URLs. Module intent: No module docstring detected.\\
\textbf{\path{tests/test_azure_runtime.py}}. Тип: python. Размер: 7540 байт. Форма: 208 строк. Модуль Python, 208 строк, явных URL: 12. Назначение модуля: No module docstring detected.\par\medskip
\textbf{\path{tests/test_benchmarks_and_notebook_scaffolding.py}}. Kind: python. Size: 11859 bytes. Shape: 274 lines. Python module, 274 lines, 0 explicit URLs. Module intent: No module docstring detected.\\
\textbf{\path{tests/test_benchmarks_and_notebook_scaffolding.py}}. Тип: python. Размер: 11859 байт. Форма: 274 строк. Модуль Python, 274 строк, явных URL: 0. Назначение модуля: No module docstring detected.\par\medskip
\textbf{\path{tests/test_cli_and_dspy.py}}. Kind: python. Size: 68622 bytes. Shape: 1876 lines. Python module, 1876 lines, 2 explicit URLs. Module intent: No module docstring detected.\\
\textbf{\path{tests/test_cli_and_dspy.py}}. Тип: python. Размер: 68622 байт. Форма: 1876 строк. Модуль Python, 1876 строк, явных URL: 2. Назначение модуля: No module docstring detected.\par\medskip
\textbf{\path{tests/test_codex_proxy.py}}. Kind: python. Size: 57154 bytes. Shape: 1571 lines. Python module, 1571 lines, 5 explicit URLs. Module intent: No module docstring detected.\\
\textbf{\path{tests/test_codex_proxy.py}}. Тип: python. Размер: 57154 байт. Форма: 1571 строк. Модуль Python, 1571 строк, явных URL: 5. Назначение модуля: No module docstring detected.\par\medskip
\textbf{\path{tests/test_dspy_training.py}}. Kind: python. Size: 66652 bytes. Shape: 1653 lines. Python module, 1653 lines, 11 explicit URLs. Module intent: No module docstring detected.\\
\textbf{\path{tests/test_dspy_training.py}}. Тип: python. Размер: 66652 байт. Форма: 1653 строк. Модуль Python, 1653 строк, явных URL: 11. Назначение модуля: No module docstring detected.\par\medskip
\textbf{\path{tests/test_exploratorium_translation.py}}. Kind: python. Size: 5239 bytes. Shape: 138 lines. Python module, 138 lines, 8 explicit URLs. Module intent: No module docstring detected.\\
\textbf{\path{tests/test_exploratorium_translation.py}}. Тип: python. Размер: 5239 байт. Форма: 138 строк. Модуль Python, 138 строк, явных URL: 8. Назначение модуля: No module docstring detected.\par\medskip
\textbf{\path{tests/test_file_summaries.py}}. Kind: python. Size: 9110 bytes. Shape: 230 lines. Python module, 230 lines, 0 explicit URLs. Module intent: No module docstring detected.\\
\textbf{\path{tests/test_file_summaries.py}}. Тип: python. Размер: 9110 байт. Форма: 230 строк. Модуль Python, 230 строк, явных URL: 0. Назначение модуля: No module docstring detected.\par\medskip
\textbf{\path{tests/test_github_pr_gates.py}}. Kind: python. Size: 5678 bytes. Shape: 183 lines. Python module, 183 lines, 1 explicit URLs. Module intent: No module docstring detected.\\
\textbf{\path{tests/test_github_pr_gates.py}}. Тип: python. Размер: 5678 байт. Форма: 183 строк. Модуль Python, 183 строк, явных URL: 1. Назначение модуля: No module docstring detected.\par\medskip
\textbf{\path{tests/test_hushwheel_fixture.py}}. Kind: python. Size: 4617 bytes. Shape: 101 lines. Python module, 101 lines, 0 explicit URLs. Module intent: No module docstring detected.\\
\textbf{\path{tests/test_hushwheel_fixture.py}}. Тип: python. Размер: 4617 байт. Форма: 101 строк. Модуль Python, 101 строк, явных URL: 0. Назначение модуля: No module docstring detected.\par\medskip
\textbf{\path{tests/test_hushwheel_program_surface.py}}. Kind: python. Size: 4626 bytes. Shape: 125 lines. Python module, 125 lines, 0 explicit URLs. Module intent: No module docstring detected.\\
\textbf{\path{tests/test_hushwheel_program_surface.py}}. Тип: python. Размер: 4626 байт. Форма: 125 строк. Модуль Python, 125 строк, явных URL: 0. Назначение модуля: No module docstring detected.\par\medskip
\textbf{\path{tests/test_live_azure_integration.py}}. Kind: python. Size: 2634 bytes. Shape: 83 lines. Python module, 83 lines, 0 explicit URLs. Module intent: No module docstring detected.\\
\textbf{\path{tests/test_live_azure_integration.py}}. Тип: python. Размер: 2634 байт. Форма: 83 строк. Модуль Python, 83 строк, явных URL: 0. Назначение модуля: No module docstring detected.\par\medskip
\textbf{\path{tests/test_lookup_first.py}}. Kind: python. Size: 3559 bytes. Shape: 108 lines. Python module, 108 lines, 0 explicit URLs. Module intent: No module docstring detected.\\
\textbf{\path{tests/test_lookup_first.py}}. Тип: python. Размер: 3559 байт. Форма: 108 строк. Модуль Python, 108 строк, явных URL: 0. Назначение модуля: No module docstring detected.\par\medskip
\textbf{\path{tests/test_mcp_branches.py}}. Kind: python. Size: 1077 bytes. Shape: 31 lines. Python module, 31 lines, 0 explicit URLs. Module intent: No module docstring detected.\\
\textbf{\path{tests/test_mcp_branches.py}}. Тип: python. Размер: 1077 байт. Форма: 31 строк. Модуль Python, 31 строк, явных URL: 0. Назначение модуля: No module docstring detected.\par\medskip
\textbf{\path{tests/test_mcp_server.py}}. Kind: python. Size: 26557 bytes. Shape: 799 lines. Python module, 799 lines, 0 explicit URLs. Module intent: No module docstring detected.\\
\textbf{\path{tests/test_mcp_server.py}}. Тип: python. Размер: 26557 байт. Форма: 799 строк. Модуль Python, 799 строк, явных URL: 0. Назначение модуля: No module docstring detected.\par\medskip
\textbf{\path{tests/test_mcp_stdio.py}}. Kind: python. Size: 988 bytes. Shape: 42 lines. Python module, 42 lines, 0 explicit URLs. Module intent: No module docstring detected.\\
\textbf{\path{tests/test_mcp_stdio.py}}. Тип: python. Размер: 988 байт. Форма: 42 строк. Модуль Python, 42 строк, явных URL: 0. Назначение модуля: No module docstring detected.\par\medskip
\textbf{\path{tests/test_notebook_runner.py}}. Kind: python. Size: 3088 bytes. Shape: 84 lines. Python module, 84 lines, 0 explicit URLs. Module intent: No module docstring detected.\\
\textbf{\path{tests/test_notebook_runner.py}}. Тип: python. Размер: 3088 байт. Форма: 84 строк. Модуль Python, 84 строк, явных URL: 0. Назначение модуля: No module docstring detected.\par\medskip
\textbf{\path{tests/test_pages_site.py}}. Kind: python. Size: 9045 bytes. Shape: 225 lines. Python module, 225 lines, 14 explicit URLs. Module intent: No module docstring detected.\\
\textbf{\path{tests/test_pages_site.py}}. Тип: python. Размер: 9045 байт. Форма: 225 строк. Модуль Python, 225 строк, явных URL: 14. Назначение модуля: No module docstring detected.\par\medskip
\textbf{\path{tests/test_population_samples.py}}. Kind: python. Size: 3144 bytes. Shape: 76 lines. Python module, 76 lines, 0 explicit URLs. Module intent: No module docstring detected.\\
\textbf{\path{tests/test_population_samples.py}}. Тип: python. Размер: 3144 байт. Форма: 76 строк. Модуль Python, 76 строк, явных URL: 0. Назначение модуля: No module docstring detected.\par\medskip
\textbf{\path{tests/test_project_surfaces.py}}. Kind: python. Size: 21471 bytes. Shape: 484 lines. Python module, 484 lines, 1 explicit URLs. Module intent: No module docstring detected.\\
\textbf{\path{tests/test_project_surfaces.py}}. Тип: python. Размер: 21471 байт. Форма: 484 строк. Модуль Python, 484 строк, явных URL: 1. Назначение модуля: No module docstring detected.\par\medskip
\textbf{\path{tests/test_repository_rag_bdd.py}}. Kind: python. Size: 879 bytes. Shape: 30 lines. Python module, 30 lines, 0 explicit URLs. Module intent: No module docstring detected.\\
\textbf{\path{tests/test_repository_rag_bdd.py}}. Тип: python. Размер: 879 байт. Форма: 30 строк. Модуль Python, 30 строк, явных URL: 0. Назначение модуля: No module docstring detected.\par\medskip
\textbf{\path{tests/test_retrieval.py}}. Kind: python. Size: 13362 bytes. Shape: 389 lines. Python module, 389 lines, 0 explicit URLs. Module intent: No module docstring detected.\\
\textbf{\path{tests/test_retrieval.py}}. Тип: python. Размер: 13362 байт. Форма: 389 строк. Модуль Python, 389 строк, явных URL: 0. Назначение модуля: No module docstring detected.\par\medskip
\textbf{\path{tests/test_runtime_artifacts_azure.py}}. Kind: python. Size: 58883 bytes. Shape: 1554 lines. Python module, 1554 lines, 4 explicit URLs. Module intent: No module docstring detected.\\
\textbf{\path{tests/test_runtime_artifacts_azure.py}}. Тип: python. Размер: 58883 байт. Форма: 1554 строк. Модуль Python, 1554 строк, явных URL: 4. Назначение модуля: No module docstring detected.\par\medskip
\textbf{\path{tests/test_rust_lookup.py}}. Kind: python. Size: 3099 bytes. Shape: 95 lines. Python module, 95 lines, 0 explicit URLs. Module intent: No module docstring detected.\\
\textbf{\path{tests/test_rust_lookup.py}}. Тип: python. Размер: 3099 байт. Форма: 95 строк. Модуль Python, 95 строк, явных URL: 0. Назначение модуля: No module docstring detected.\par\medskip
\textbf{\path{tests/test_todo_backlog.py}}. Kind: python. Size: 3505 bytes. Shape: 103 lines. Python module, 103 lines, 1 explicit URLs. Module intent: No module docstring detected.\\
\textbf{\path{tests/test_todo_backlog.py}}. Тип: python. Размер: 3505 байт. Форма: 103 строк. Модуль Python, 103 строк, явных URL: 1. Назначение модуля: No module docstring detected.\par\medskip
\textbf{\path{tests/test_training_samples.py}}. Kind: python. Size: 84870 bytes. Shape: 2152 lines. Python module, 2152 lines, 0 explicit URLs. Module intent: No module docstring detected.\\
\textbf{\path{tests/test_training_samples.py}}. Тип: python. Размер: 84870 байт. Форма: 2152 строк. Модуль Python, 2152 строк, явных URL: 0. Назначение модуля: No module docstring detected.\par\medskip
\textbf{\path{tests/test_utilities.py}}. Kind: python. Size: 121479 bytes. Shape: 3117 lines. Python module, 3117 lines, 6 explicit URLs. Module intent: No module docstring detected.\\
\textbf{\path{tests/test_utilities.py}}. Тип: python. Размер: 121479 байт. Форма: 3117 строк. Модуль Python, 3117 строк, явных URL: 6. Назначение модуля: No module docstring detected.\par\medskip
\textbf{\path{tests/test_verification.py}}. Kind: python. Size: 3156 bytes. Shape: 66 lines. Python module, 66 lines, 0 explicit URLs. Module intent: No module docstring detected.\\
\textbf{\path{tests/test_verification.py}}. Тип: python. Размер: 3156 байт. Форма: 66 строк. Модуль Python, 66 строк, явных URL: 0. Назначение модуля: No module docstring detected.\par\medskip
\textbf{\path{tests/test_workflow.py}}. Kind: python. Size: 3456 bytes. Shape: 99 lines. Python module, 99 lines, 0 explicit URLs. Module intent: No module docstring detected.\\
\textbf{\path{tests/test_workflow.py}}. Тип: python. Размер: 3456 байт. Форма: 99 строк. Модуль Python, 99 строк, явных URL: 0. Назначение модуля: No module docstring detected.\par\medskip
\textbf{\path{tests/test_workflow_live.py}}. Kind: python. Size: 4292 bytes. Shape: 128 lines. Python module, 128 lines, 0 explicit URLs. Module intent: No module docstring detected.\\
\textbf{\path{tests/test_workflow_live.py}}. Тип: python. Размер: 4292 байт. Форма: 128 строк. Модуль Python, 128 строк, явных URL: 0. Назначение модуля: No module docstring detected.\par\medskip
\textbf{\path{todo-backlog.yaml}}. Kind: yaml. Size: 6116 bytes. Shape: 128 lines. YAML data or workflow file, 128 lines, 1 explicit URLs. Top keys: repository\_web\_base, items.\\
\textbf{\path{todo-backlog.yaml}}. Тип: yaml. Размер: 6116 байт. Форма: 128 строк. Файл YAML с данными или workflow, 128 строк, явных URL: 1. Верхние ключи: repository\_web\_base, items.\par\medskip
\textbf{\path{utilities/README.md}}. Kind: markdown. Size: 3495 bytes. Shape: 54 lines. Markdown document, 54 lines, 0 explicit URLs. Lead heading: Repository Utilities.\\
\textbf{\path{utilities/README.md}}. Тип: markdown. Размер: 3495 байт. Форма: 54 строк. Документ Markdown, 54 строк, явных URL: 0. Главный заголовок: Repository Utilities.\par\medskip
\textbf{\path{uv.lock}}. Kind: generic. Size: 784153 bytes. Shape: binary surface. Tracked repository file, 784153 bytes of binary data, 0 explicit URLs. Opening line: No non-empty preview line.\\
\textbf{\path{uv.lock}}. Тип: generic. Размер: 784153 байт. Форма: бинарная поверхность. Отслеживаемый файл репозитория, 784153 байт бинарных данных, явных URL: 0. Первая строка: No non-empty preview line.\par\medskip
\subsection*{Summaries Of All Authored URLs / Сводки по всем авторским URL}
\textbf{\url{http://127.0.0.1}}. Host: 127.0.0.1. Occurrences: 2. Referenced by URL only. No repository-local mirror is currently tracked. Sources: \path{tests/test_codex_proxy.py}.\\
\textbf{\url{http://127.0.0.1}}. Хост: 127.0.0.1. Упоминаний: 2. Ссылка хранится только как URL. Локальное зеркало в репозитории сейчас не отслеживается. Источники: \path{tests/test_codex_proxy.py}.\par\medskip
\textbf{\url{http://127.0.0.1:40111/openai`}}. Host: 127.0.0.1:40111. Occurrences: 1. Referenced by URL only. No repository-local mirror is currently tracked. Sources: \path{docs/audit/2026-05-02-zz-codex-exec-resume-session-state-stage0.md}.\\
\textbf{\url{http://127.0.0.1:40111/openai`}}. Хост: 127.0.0.1:40111. Упоминаний: 1. Ссылка хранится только как URL. Локальное зеркало в репозитории сейчас не отслеживается. Источники: \path{docs/audit/2026-05-02-zz-codex-exec-resume-session-state-stage0.md}.\par\medskip
\textbf{\url{http://127.0.0.1:44973/openai`}}. Host: 127.0.0.1:44973. Occurrences: 1. Referenced by URL only. No repository-local mirror is currently tracked. Sources: \path{docs/audit/2026-05-02-zz-codex-exec-resume-session-state-stage0.md}.\\
\textbf{\url{http://127.0.0.1:44973/openai`}}. Хост: 127.0.0.1:44973. Упоминаний: 1. Ссылка хранится только как URL. Локальное зеркало в репозитории сейчас не отслеживается. Источники: \path{docs/audit/2026-05-02-zz-codex-exec-resume-session-state-stage0.md}.\par\medskip
\textbf{\url{https://api.github.com/repos/example/demo/branches/master/protection}}. Host: api.github.com. Occurrences: 1. External GitHub URL referenced from 1 file(s). Local companion path(s): none. Sources: \path{tests/test_github_pr_gates.py}.\\
\textbf{\url{https://api.github.com/repos/example/demo/branches/master/protection}}. Хост: api.github.com. Упоминаний: 1. Внешний URL GitHub, упомянутый в 1 файле(ах). Локальные сопроводительные пути: none. Источники: \path{tests/test_github_pr_gates.py}.\par\medskip
\textbf{\url{https://astral.sh/}}. Host: astral.sh. Occurrences: 2. Referenced by URL only. No repository-local mirror is currently tracked. Sources: \path{tests/test_exploratorium_translation.py}\newline \path{tests/test_utilities.py}.\\
\textbf{\url{https://astral.sh/}}. Хост: astral.sh. Упоминаний: 2. Ссылка хранится только как URL. Локальное зеркало в репозитории сейчас не отслеживается. Источники: \path{tests/test_exploratorium_translation.py}\newline \path{tests/test_utilities.py}.\par\medskip
\textbf{\url{https://astral.sh/uv/install.sh}}. Host: astral.sh. Occurrences: 1. Referenced by URL only. No repository-local mirror is currently tracked. Sources: \path{README.md}.\\
\textbf{\url{https://astral.sh/uv/install.sh}}. Хост: astral.sh. Упоминаний: 1. Ссылка хранится только как URL. Локальное зеркало в репозитории сейчас не отслеживается. Источники: \path{README.md}.\par\medskip
\textbf{\url{https://dspy-helper.openai.azure.com}}. Host: dspy-helper.openai.azure.com. Occurrences: 1. Referenced by URL only. No repository-local mirror is currently tracked. Sources: \path{tests/test_dspy_training.py}.\\
\textbf{\url{https://dspy-helper.openai.azure.com}}. Хост: dspy-helper.openai.azure.com. Упоминаний: 1. Ссылка хранится только как URL. Локальное зеркало в репозитории сейчас не отслеживается. Источники: \path{tests/test_dspy_training.py}.\par\medskip
\textbf{\url{https://dspy-helper.openai.azure.com/}}. Host: dspy-helper.openai.azure.com. Occurrences: 1. Referenced by URL only. No repository-local mirror is currently tracked. Sources: \path{tests/test_dspy_training.py}.\\
\textbf{\url{https://dspy-helper.openai.azure.com/}}. Хост: dspy-helper.openai.azure.com. Упоминаний: 1. Ссылка хранится только как URL. Локальное зеркало в репозитории сейчас не отслеживается. Источники: \path{tests/test_dspy_training.py}.\par\medskip
\textbf{\url{https://dspy.ai/tutorials/rag/}}. Host: dspy.ai. Occurrences: 1. Referenced by URL only. No repository-local mirror is currently tracked. Sources: \path{docs/architecture/inspired/dspy-rag-tutorial.md}.\\
\textbf{\url{https://dspy.ai/tutorials/rag/}}. Хост: dspy.ai. Упоминаний: 1. Ссылка хранится только как URL. Локальное зеркало в репозитории сейчас не отслеживается. Источники: \path{docs/architecture/inspired/dspy-rag-tutorial.md}.\par\medskip
\textbf{\url{https://example.com}}. Host: example.com. Occurrences: 2. Referenced by URL only. No repository-local mirror is currently tracked. Sources: \path{tests/test_pages_site.py}.\\
\textbf{\url{https://example.com}}. Хост: example.com. Упоминаний: 2. Ссылка хранится только как URL. Локальное зеркало в репозитории сейчас не отслеживается. Источники: \path{tests/test_pages_site.py}.\par\medskip
\textbf{\url{https://example.com/docs}}. Host: example.com. Occurrences: 1. Referenced by URL only. No repository-local mirror is currently tracked. Sources: \path{tests/test_pages_site.py}.\\
\textbf{\url{https://example.com/docs}}. Хост: example.com. Упоминаний: 1. Ссылка хранится только как URL. Локальное зеркало в репозитории сейчас не отслеживается. Источники: \path{tests/test_pages_site.py}.\par\medskip
\textbf{\url{https://example.com/org/repo/blob/main}}. Host: example.com. Occurrences: 1. Referenced by URL only. No repository-local mirror is currently tracked. Sources: \path{tests/test_todo_backlog.py}.\\
\textbf{\url{https://example.com/org/repo/blob/main}}. Хост: example.com. Упоминаний: 1. Ссылка хранится только как URL. Локальное зеркало в репозитории сейчас не отслеживается. Источники: \path{tests/test_todo_backlog.py}.\par\medskip
\textbf{\url{https://example.invalid/v1}}. Host: example.invalid. Occurrences: 2. Referenced by URL only. No repository-local mirror is currently tracked. Sources: \path{tests/test_dspy_training.py}.\\
\textbf{\url{https://example.invalid/v1}}. Хост: example.invalid. Упоминаний: 2. Ссылка хранится только как URL. Локальное зеркало в репозитории сейчас не отслеживается. Источники: \path{tests/test_dspy_training.py}.\par\medskip
\textbf{\url{https://example.openai.azure.com}}. Host: example.openai.azure.com. Occurrences: 5. Referenced by URL only. No repository-local mirror is currently tracked. Sources: \path{tests/test_azure_runtime.py}\newline \path{tests/test_dspy_training.py}.\\
\textbf{\url{https://example.openai.azure.com}}. Хост: example.openai.azure.com. Упоминаний: 5. Ссылка хранится только как URL. Локальное зеркало в репозитории сейчас не отслеживается. Источники: \path{tests/test_azure_runtime.py}\newline \path{tests/test_dspy_training.py}.\par\medskip
\textbf{\url{https://example.openai.azure.com/}}. Host: example.openai.azure.com. Occurrences: 6. Referenced by URL only. No repository-local mirror is currently tracked. Sources: \path{tests/test_azure_runtime.py}\newline \path{tests/test_dspy_training.py}.\\
\textbf{\url{https://example.openai.azure.com/}}. Хост: example.openai.azure.com. Упоминаний: 6. Ссылка хранится только как URL. Локальное зеркало в репозитории сейчас не отслеживается. Источники: \path{tests/test_azure_runtime.py}\newline \path{tests/test_dspy_training.py}.\par\medskip
\textbf{\url{https://example.openai.azure.com/openai/deployments/repo-rag-ft}}. Host: example.openai.azure.com. Occurrences: 3. Referenced by URL only. No repository-local mirror is currently tracked. Sources: \path{tests/test_azure_runtime.py}.\\
\textbf{\url{https://example.openai.azure.com/openai/deployments/repo-rag-ft}}. Хост: example.openai.azure.com. Упоминаний: 3. Ссылка хранится только как URL. Локальное зеркало в репозитории сейчас не отслеживается. Источники: \path{tests/test_azure_runtime.py}.\par\medskip
\textbf{\url{https://example.openai.azure.com/openai/deployments/repo-rag-ft/}}. Host: example.openai.azure.com. Occurrences: 4. Referenced by URL only. No repository-local mirror is currently tracked. Sources: \path{tests/test_azure_runtime.py}.\\
\textbf{\url{https://example.openai.azure.com/openai/deployments/repo-rag-ft/}}. Хост: example.openai.azure.com. Упоминаний: 4. Ссылка хранится только как URL. Локальное зеркало в репозитории сейчас не отслеживается. Источники: \path{tests/test_azure_runtime.py}.\par\medskip
\textbf{\url{https://example.openai.azure.com/openai/deployments/repo-rag-ft/chat/completions}}. Host: example.openai.azure.com. Occurrences: 1. Referenced by URL only. No repository-local mirror is currently tracked. Sources: \path{tests/test_dspy_training.py}.\\
\textbf{\url{https://example.openai.azure.com/openai/deployments/repo-rag-ft/chat/completions}}. Хост: example.openai.azure.com. Упоминаний: 1. Ссылка хранится только как URL. Локальное зеркало в репозитории сейчас не отслеживается. Источники: \path{tests/test_dspy_training.py}.\par\medskip
\textbf{\url{https://example.services.ai.azure.com/models}}. Host: example.services.ai.azure.com. Occurrences: 6. Referenced by URL only. No repository-local mirror is currently tracked. Sources: \path{Makefile}\newline \path{docs/architecture/dspy-rag-guide.md}\newline \path{docs/operations/azure-deployment.md}\newline \path{src/repo_rag_lab/utilities.py}\newline \path{tests/test_cli_and_dspy.py}\newline \path{tests/test_project_surfaces.py}.\\
\textbf{\url{https://example.services.ai.azure.com/models}}. Хост: example.services.ai.azure.com. Упоминаний: 6. Ссылка хранится только как URL. Локальное зеркало в репозитории сейчас не отслеживается. Источники: \path{Makefile}\newline \path{docs/architecture/dspy-rag-guide.md}\newline \path{docs/operations/azure-deployment.md}\newline \path{src/repo_rag_lab/utilities.py}\newline \path{tests/test_cli_and_dspy.py}\newline \path{tests/test_project_surfaces.py}.\par\medskip
\textbf{\url{https://from-connection-string}}. Host: from-connection-string. Occurrences: 2. Referenced by URL only. No repository-local mirror is currently tracked. Sources: \path{tests/test_runtime_artifacts_azure.py}.\\
\textbf{\url{https://from-connection-string}}. Хост: from-connection-string. Упоминаний: 2. Ссылка хранится только как URL. Локальное зеркало в репозитории сейчас не отслеживается. Источники: \path{tests/test_runtime_artifacts_azure.py}.\par\medskip
\textbf{\url{https://github.com/}}. Host: github.com. Occurrences: 2. External GitHub URL referenced from 1 file(s). Local companion path(s): none. Sources: \path{src/repo_rag_lab/pages_site.py}.\\
\textbf{\url{https://github.com/}}. Хост: github.com. Упоминаний: 2. Внешний URL GitHub, упомянутый в 1 файле(ах). Локальные сопроводительные пути: none. Источники: \path{src/repo_rag_lab/pages_site.py}.\par\medskip
\textbf{\url{https://github.com/acme/repo}}. Host: github.com. Occurrences: 3. External GitHub URL referenced from 1 file(s). Local companion path(s): none. Sources: \path{tests/test_codex_proxy.py}.\\
\textbf{\url{https://github.com/acme/repo}}. Хост: github.com. Упоминаний: 3. Внешний URL GitHub, упомянутый в 1 файле(ах). Локальные сопроводительные пути: none. Источники: \path{tests/test_codex_proxy.py}.\par\medskip
\textbf{\url{https://github.com/example/demo}}. Host: github.com. Occurrences: 10. External GitHub URL referenced from 3 file(s). Local companion path(s): none. Sources: \path{tests/test_cli_and_dspy.py}\newline \path{tests/test_pages_site.py}\newline \path{tests/test_utilities.py}.\\
\textbf{\url{https://github.com/example/demo}}. Хост: github.com. Упоминаний: 10. Внешний URL GitHub, упомянутый в 3 файле(ах). Локальные сопроводительные пути: none. Источники: \path{tests/test_cli_and_dspy.py}\newline \path{tests/test_pages_site.py}\newline \path{tests/test_utilities.py}.\par\medskip
\textbf{\url{https://github.com/example/demo/blob/main/README.md}}. Host: github.com. Occurrences: 1. External GitHub URL referenced from 1 file(s). Local companion path(s): none. Sources: \path{tests/test_pages_site.py}.\\
\textbf{\url{https://github.com/example/demo/blob/main/README.md}}. Хост: github.com. Упоминаний: 1. Внешний URL GitHub, упомянутый в 1 файле(ах). Локальные сопроводительные пути: none. Источники: \path{tests/test_pages_site.py}.\par\medskip
\textbf{\url{https://github.com/example/demo/blob/main/assets/logo.png}}. Host: github.com. Occurrences: 1. External GitHub URL referenced from 1 file(s). Local companion path(s): none. Sources: \path{tests/test_pages_site.py}.\\
\textbf{\url{https://github.com/example/demo/blob/main/assets/logo.png}}. Хост: github.com. Упоминаний: 1. Внешний URL GitHub, упомянутый в 1 файле(ах). Локальные сопроводительные пути: none. Источники: \path{tests/test_pages_site.py}.\par\medskip
\textbf{\url{https://github.com/example/demo/blob/main/src/tool.c}}. Host: github.com. Occurrences: 1. External GitHub URL referenced from 1 file(s). Local companion path(s): none. Sources: \path{tests/test_pages_site.py}.\\
\textbf{\url{https://github.com/example/demo/blob/main/src/tool.c}}. Хост: github.com. Упоминаний: 1. Внешний URL GitHub, упомянутый в 1 файле(ах). Локальные сопроводительные пути: none. Источники: \path{tests/test_pages_site.py}.\par\medskip
\textbf{\url{https://github.com/example/demo/tree/main/docs}}. Host: github.com. Occurrences: 1. External GitHub URL referenced from 1 file(s). Local companion path(s): none. Sources: \path{tests/test_pages_site.py}.\\
\textbf{\url{https://github.com/example/demo/tree/main/docs}}. Хост: github.com. Упоминаний: 1. Внешний URL GitHub, упомянутый в 1 файле(ах). Локальные сопроводительные пути: none. Источники: \path{tests/test_pages_site.py}.\par\medskip
\textbf{\url{https://github.com/example/project}}. Host: github.com. Occurrences: 2. External GitHub URL referenced from 2 file(s). Local companion path(s): none. Sources: \path{tests/test_exploratorium_translation.py}\newline \path{tests/test_utilities.py}.\\
\textbf{\url{https://github.com/example/project}}. Хост: github.com. Упоминаний: 2. Внешний URL GitHub, упомянутый в 2 файле(ах). Локальные сопроводительные пути: none. Источники: \path{tests/test_exploratorium_translation.py}\newline \path{tests/test_utilities.py}.\par\medskip
\textbf{\url{https://github.com/example/project/blob/main/README.md}}. Host: github.com. Occurrences: 1. External GitHub URL referenced from 1 file(s). Local companion path(s): none. Sources: \path{tests/test_exploratorium_translation.py}.\\
\textbf{\url{https://github.com/example/project/blob/main/README.md}}. Хост: github.com. Упоминаний: 1. Внешний URL GitHub, упомянутый в 1 файле(ах). Локальные сопроводительные пути: none. Источники: \path{tests/test_exploratorium_translation.py}.\par\medskip
\textbf{\url{https://github.com/example/project/blob/main/README.md\#L1-L2}}. Host: github.com. Occurrences: 1. External GitHub URL referenced from 1 file(s). Local companion path(s): none. Sources: \path{tests/test_exploratorium_translation.py}.\\
\textbf{\url{https://github.com/example/project/blob/main/README.md\#L1-L2}}. Хост: github.com. Упоминаний: 1. Внешний URL GitHub, упомянутый в 1 файле(ах). Локальные сопроводительные пути: none. Источники: \path{tests/test_exploratorium_translation.py}.\par\medskip
\textbf{\url{https://github.com/realagiorganization/dspy_rag_in_repo_docs_and_impl1}}. Host: github.com. Occurrences: 1. External GitHub URL referenced from 1 file(s). Local companion path(s): none. Sources: \path{mkdocs.yml}.\\
\textbf{\url{https://github.com/realagiorganization/dspy_rag_in_repo_docs_and_impl1}}. Хост: github.com. Упоминаний: 1. Внешний URL GitHub, упомянутый в 1 файле(ах). Локальные сопроводительные пути: none. Источники: \path{mkdocs.yml}.\par\medskip
\textbf{\url{https://github.com/realagiorganization/dspy_rag_in_repo_docs_and_impl1/actions/workflows/ci.yml/badge.svg)](https://github.com/realagiorganization/dspy_rag_in_repo_docs_and_impl1/actions/workflows/ci.yml}}. Host: github.com. Occurrences: 2. External GitHub URL referenced from 2 file(s). Local companion path(s): none. Sources: \path{README.md}\newline \path{notebooks/02_agent_workflow_checklist.ipynb}.\\
\textbf{\url{https://github.com/realagiorganization/dspy_rag_in_repo_docs_and_impl1/actions/workflows/ci.yml/badge.svg)](https://github.com/realagiorganization/dspy_rag_in_repo_docs_and_impl1/actions/workflows/ci.yml}}. Хост: github.com. Упоминаний: 2. Внешний URL GitHub, упомянутый в 2 файле(ах). Локальные сопроводительные пути: none. Источники: \path{README.md}\newline \path{notebooks/02_agent_workflow_checklist.ipynb}.\par\medskip
\textbf{\url{https://github.com/realagiorganization/dspy_rag_in_repo_docs_and_impl1/actions/workflows/publish.yml/badge.svg)](https://github.com/realagiorganization/dspy_rag_in_repo_docs_and_impl1/actions/workflows/publish.yml}}. Host: github.com. Occurrences: 1. External GitHub URL referenced from 1 file(s). Local companion path(s): none. Sources: \path{README.md}.\\
\textbf{\url{https://github.com/realagiorganization/dspy_rag_in_repo_docs_and_impl1/actions/workflows/publish.yml/badge.svg)](https://github.com/realagiorganization/dspy_rag_in_repo_docs_and_impl1/actions/workflows/publish.yml}}. Хост: github.com. Упоминаний: 1. Внешний URL GitHub, упомянутый в 1 файле(ах). Локальные сопроводительные пути: none. Источники: \path{README.md}.\par\medskip
\textbf{\url{https://github.com/realagiorganization/dspy_rag_in_repo_docs_and_impl1/blob/master}}. Host: github.com. Occurrences: 3. External GitHub URL referenced from 3 file(s). Local companion path(s): none. Sources: \path{FILES.md}\newline \path{src/repo_rag_lab/todo_backlog.py}\newline \path{todo-backlog.yaml}.\\
\textbf{\url{https://github.com/realagiorganization/dspy_rag_in_repo_docs_and_impl1/blob/master}}. Хост: github.com. Упоминаний: 3. Внешний URL GitHub, упомянутый в 3 файле(ах). Локальные сопроводительные пути: none. Источники: \path{FILES.md}\newline \path{src/repo_rag_lab/todo_backlog.py}\newline \path{todo-backlog.yaml}.\par\medskip
\textbf{\url{https://github.com/realagiorganization/dspy_rag_in_repo_docs_and_impl1/blob/master/.github/workflows/ci.yml}}. Host: github.com. Occurrences: 1. External GitHub URL referenced from 1 file(s). Local companion path(s): none. Sources: \path{publication/todo-backlog-table.tex}.\\
\textbf{\url{https://github.com/realagiorganization/dspy_rag_in_repo_docs_and_impl1/blob/master/.github/workflows/ci.yml}}. Хост: github.com. Упоминаний: 1. Внешний URL GitHub, упомянутый в 1 файле(ах). Локальные сопроводительные пути: none. Источники: \path{publication/todo-backlog-table.tex}.\par\medskip
\textbf{\url{https://github.com/realagiorganization/dspy_rag_in_repo_docs_and_impl1/blob/master/.github/workflows/publish.yml}}. Host: github.com. Occurrences: 1. External GitHub URL referenced from 1 file(s). Local companion path(s): none. Sources: \path{publication/todo-backlog-table.tex}.\\
\textbf{\url{https://github.com/realagiorganization/dspy_rag_in_repo_docs_and_impl1/blob/master/.github/workflows/publish.yml}}. Хост: github.com. Упоминаний: 1. Внешний URL GitHub, упомянутый в 1 файле(ах). Локальные сопроводительные пути: none. Источники: \path{publication/todo-backlog-table.tex}.\par\medskip
\textbf{\url{https://github.com/realagiorganization/dspy_rag_in_repo_docs_and_impl1/blob/master/Makefile\#L142-L143}}. Host: github.com. Occurrences: 1. External GitHub URL referenced from 1 file(s). Local companion path(s): none. Sources: \path{docs/operations/simple-end-to-end-verification-guide.md}.\\
\textbf{\url{https://github.com/realagiorganization/dspy_rag_in_repo_docs_and_impl1/blob/master/Makefile\#L142-L143}}. Хост: github.com. Упоминаний: 1. Внешний URL GitHub, упомянутый в 1 файле(ах). Локальные сопроводительные пути: none. Источники: \path{docs/operations/simple-end-to-end-verification-guide.md}.\par\medskip
\textbf{\url{https://github.com/realagiorganization/dspy_rag_in_repo_docs_and_impl1/blob/master/Makefile\#L153-L154}}. Host: github.com. Occurrences: 1. External GitHub URL referenced from 1 file(s). Local companion path(s): none. Sources: \path{docs/operations/simple-end-to-end-verification-guide.md}.\\
\textbf{\url{https://github.com/realagiorganization/dspy_rag_in_repo_docs_and_impl1/blob/master/Makefile\#L153-L154}}. Хост: github.com. Упоминаний: 1. Внешний URL GitHub, упомянутый в 1 файле(ах). Локальные сопроводительные пути: none. Источники: \path{docs/operations/simple-end-to-end-verification-guide.md}.\par\medskip
\textbf{\url{https://github.com/realagiorganization/dspy_rag_in_repo_docs_and_impl1/blob/master/Makefile\#L248-L251}}. Host: github.com. Occurrences: 1. External GitHub URL referenced from 1 file(s). Local companion path(s): none. Sources: \path{docs/operations/simple-end-to-end-verification-guide.md}.\\
\textbf{\url{https://github.com/realagiorganization/dspy_rag_in_repo_docs_and_impl1/blob/master/Makefile\#L248-L251}}. Хост: github.com. Упоминаний: 1. Внешний URL GitHub, упомянутый в 1 файле(ах). Локальные сопроводительные пути: none. Источники: \path{docs/operations/simple-end-to-end-verification-guide.md}.\par\medskip
\textbf{\url{https://github.com/realagiorganization/dspy_rag_in_repo_docs_and_impl1/blob/master/Makefile\#L65-L66}}. Host: github.com. Occurrences: 1. External GitHub URL referenced from 1 file(s). Local companion path(s): none. Sources: \path{docs/operations/simple-end-to-end-verification-guide.md}.\\
\textbf{\url{https://github.com/realagiorganization/dspy_rag_in_repo_docs_and_impl1/blob/master/Makefile\#L65-L66}}. Хост: github.com. Упоминаний: 1. Внешний URL GitHub, упомянутый в 1 файле(ах). Локальные сопроводительные пути: none. Источники: \path{docs/operations/simple-end-to-end-verification-guide.md}.\par\medskip
\textbf{\url{https://github.com/realagiorganization/dspy_rag_in_repo_docs_and_impl1/blob/master/README.md}}. Host: github.com. Occurrences: 2. External GitHub URL referenced from 1 file(s). Local companion path(s): none. Sources: \path{publication/todo-backlog-table.tex}.\\
\textbf{\url{https://github.com/realagiorganization/dspy_rag_in_repo_docs_and_impl1/blob/master/README.md}}. Хост: github.com. Упоминаний: 2. Внешний URL GitHub, упомянутый в 1 файле(ах). Локальные сопроводительные пути: none. Источники: \path{publication/todo-backlog-table.tex}.\par\medskip
\textbf{\url{https://github.com/realagiorganization/dspy_rag_in_repo_docs_and_impl1/blob/master/TODO.MD}}. Host: github.com. Occurrences: 1. External GitHub URL referenced from 1 file(s). Local companion path(s): none. Sources: \path{publication/todo-backlog-table.tex}.\\
\textbf{\url{https://github.com/realagiorganization/dspy_rag_in_repo_docs_and_impl1/blob/master/TODO.MD}}. Хост: github.com. Упоминаний: 1. Внешний URL GitHub, упомянутый в 1 файле(ах). Локальные сопроводительные пути: none. Источники: \path{publication/todo-backlog-table.tex}.\par\medskip
\textbf{\url{https://github.com/realagiorganization/dspy_rag_in_repo_docs_and_impl1/blob/master/docs/architecture/dspy-rag-guide.md}}. Host: github.com. Occurrences: 1. External GitHub URL referenced from 1 file(s). Local companion path(s): none. Sources: \path{publication/todo-backlog-table.tex}.\\
\textbf{\url{https://github.com/realagiorganization/dspy_rag_in_repo_docs_and_impl1/blob/master/docs/architecture/dspy-rag-guide.md}}. Хост: github.com. Упоминаний: 1. Внешний URL GitHub, упомянутый в 1 файле(ах). Локальные сопроводительные пути: none. Источники: \path{publication/todo-backlog-table.tex}.\par\medskip
\textbf{\url{https://github.com/realagiorganization/dspy_rag_in_repo_docs_and_impl1/blob/master/docs/archive/repo-completeness-checklist.md}}. Host: github.com. Occurrences: 1. External GitHub URL referenced from 1 file(s). Local companion path(s): none. Sources: \path{publication/todo-backlog-table.tex}.\\
\textbf{\url{https://github.com/realagiorganization/dspy_rag_in_repo_docs_and_impl1/blob/master/docs/archive/repo-completeness-checklist.md}}. Хост: github.com. Упоминаний: 1. Внешний URL GitHub, упомянутый в 1 файле(ах). Локальные сопроводительные пути: none. Источники: \path{publication/todo-backlog-table.tex}.\par\medskip
\textbf{\url{https://github.com/realagiorganization/dspy_rag_in_repo_docs_and_impl1/blob/master/docs/audit/README.md}}. Host: github.com. Occurrences: 3. External GitHub URL referenced from 1 file(s). Local companion path(s): none. Sources: \path{publication/todo-backlog-table.tex}.\\
\textbf{\url{https://github.com/realagiorganization/dspy_rag_in_repo_docs_and_impl1/blob/master/docs/audit/README.md}}. Хост: github.com. Упоминаний: 3. Внешний URL GitHub, упомянутый в 1 файле(ах). Локальные сопроводительные пути: none. Источники: \path{publication/todo-backlog-table.tex}.\par\medskip
\textbf{\url{https://github.com/realagiorganization/dspy_rag_in_repo_docs_and_impl1/blob/master/docs/operations/azure-deployment.md}}. Host: github.com. Occurrences: 1. External GitHub URL referenced from 1 file(s). Local companion path(s): none. Sources: \path{publication/todo-backlog-table.tex}.\\
\textbf{\url{https://github.com/realagiorganization/dspy_rag_in_repo_docs_and_impl1/blob/master/docs/operations/azure-deployment.md}}. Хост: github.com. Упоминаний: 1. Внешний URL GitHub, упомянутый в 1 файле(ах). Локальные сопроводительные пути: none. Источники: \path{publication/todo-backlog-table.tex}.\par\medskip
\textbf{\url{https://github.com/realagiorganization/dspy_rag_in_repo_docs_and_impl1/blob/master/docs/planning/dataset-integration-plan.md}}. Host: github.com. Occurrences: 1. External GitHub URL referenced from 1 file(s). Local companion path(s): none. Sources: \path{publication/todo-backlog-table.tex}.\\
\textbf{\url{https://github.com/realagiorganization/dspy_rag_in_repo_docs_and_impl1/blob/master/docs/planning/dataset-integration-plan.md}}. Хост: github.com. Упоминаний: 1. Внешний URL GitHub, упомянутый в 1 файле(ах). Локальные сопроводительные пути: none. Источники: \path{publication/todo-backlog-table.tex}.\par\medskip
\textbf{\url{https://github.com/realagiorganization/dspy_rag_in_repo_docs_and_impl1/blob/master/docs/planning/repo-hardening-plan.md}}. Host: github.com. Occurrences: 1. External GitHub URL referenced from 1 file(s). Local companion path(s): none. Sources: \path{publication/todo-backlog-table.tex}.\\
\textbf{\url{https://github.com/realagiorganization/dspy_rag_in_repo_docs_and_impl1/blob/master/docs/planning/repo-hardening-plan.md}}. Хост: github.com. Упоминаний: 1. Внешний URL GitHub, упомянутый в 1 файле(ах). Локальные сопроводительные пути: none. Источники: \path{publication/todo-backlog-table.tex}.\par\medskip
\textbf{\url{https://github.com/realagiorganization/dspy_rag_in_repo_docs_and_impl1/blob/master/notebooks/01_repo_rag_research.ipynb}}. Host: github.com. Occurrences: 1. External GitHub URL referenced from 1 file(s). Local companion path(s): none. Sources: \path{publication/todo-backlog-table.tex}.\\
\textbf{\url{https://github.com/realagiorganization/dspy_rag_in_repo_docs_and_impl1/blob/master/notebooks/01_repo_rag_research.ipynb}}. Хост: github.com. Упоминаний: 1. Внешний URL GitHub, упомянутый в 1 файле(ах). Локальные сопроводительные пути: none. Источники: \path{publication/todo-backlog-table.tex}.\par\medskip
\textbf{\url{https://github.com/realagiorganization/dspy_rag_in_repo_docs_and_impl1/blob/master/notebooks/03_dspy_training_lab.ipynb}}. Host: github.com. Occurrences: 1. External GitHub URL referenced from 1 file(s). Local companion path(s): none. Sources: \path{publication/todo-backlog-table.tex}.\\
\textbf{\url{https://github.com/realagiorganization/dspy_rag_in_repo_docs_and_impl1/blob/master/notebooks/03_dspy_training_lab.ipynb}}. Хост: github.com. Упоминаний: 1. Внешний URL GitHub, упомянутый в 1 файле(ах). Локальные сопроводительные пути: none. Источники: \path{publication/todo-backlog-table.tex}.\par\medskip
\textbf{\url{https://github.com/realagiorganization/dspy_rag_in_repo_docs_and_impl1/blob/master/notebooks/04_sample_population_lab.ipynb}}. Host: github.com. Occurrences: 1. External GitHub URL referenced from 1 file(s). Local companion path(s): none. Sources: \path{publication/todo-backlog-table.tex}.\\
\textbf{\url{https://github.com/realagiorganization/dspy_rag_in_repo_docs_and_impl1/blob/master/notebooks/04_sample_population_lab.ipynb}}. Хост: github.com. Упоминаний: 1. Внешний URL GitHub, упомянутый в 1 файле(ах). Локальные сопроводительные пути: none. Источники: \path{publication/todo-backlog-table.tex}.\par\medskip
\textbf{\url{https://github.com/realagiorganization/dspy_rag_in_repo_docs_and_impl1/blob/master/publication/README.md}}. Host: github.com. Occurrences: 1. External GitHub URL referenced from 1 file(s). Local companion path(s): none. Sources: \path{publication/todo-backlog-table.tex}.\\
\textbf{\url{https://github.com/realagiorganization/dspy_rag_in_repo_docs_and_impl1/blob/master/publication/README.md}}. Хост: github.com. Упоминаний: 1. Внешний URL GitHub, упомянутый в 1 файле(ах). Локальные сопроводительные пути: none. Источники: \path{publication/todo-backlog-table.tex}.\par\medskip
\textbf{\url{https://github.com/realagiorganization/dspy_rag_in_repo_docs_and_impl1/blob/master/pyproject.toml}}. Host: github.com. Occurrences: 1. External GitHub URL referenced from 1 file(s). Local companion path(s): none. Sources: \path{publication/todo-backlog-table.tex}.\\
\textbf{\url{https://github.com/realagiorganization/dspy_rag_in_repo_docs_and_impl1/blob/master/pyproject.toml}}. Хост: github.com. Упоминаний: 1. Внешний URL GitHub, упомянутый в 1 файле(ах). Локальные сопроводительные пути: none. Источники: \path{publication/todo-backlog-table.tex}.\par\medskip
\textbf{\url{https://github.com/realagiorganization/dspy_rag_in_repo_docs_and_impl1/blob/master/rust-cli/Cargo.lock}}. Host: github.com. Occurrences: 1. External GitHub URL referenced from 1 file(s). Local companion path(s): none. Sources: \path{publication/todo-backlog-table.tex}.\\
\textbf{\url{https://github.com/realagiorganization/dspy_rag_in_repo_docs_and_impl1/blob/master/rust-cli/Cargo.lock}}. Хост: github.com. Упоминаний: 1. Внешний URL GitHub, упомянутый в 1 файле(ах). Локальные сопроводительные пути: none. Источники: \path{publication/todo-backlog-table.tex}.\par\medskip
\textbf{\url{https://github.com/realagiorganization/dspy_rag_in_repo_docs_and_impl1/blob/master/rust-cli/Cargo.toml}}. Host: github.com. Occurrences: 1. External GitHub URL referenced from 1 file(s). Local companion path(s): none. Sources: \path{publication/todo-backlog-table.tex}.\\
\textbf{\url{https://github.com/realagiorganization/dspy_rag_in_repo_docs_and_impl1/blob/master/rust-cli/Cargo.toml}}. Хост: github.com. Упоминаний: 1. Внешний URL GitHub, упомянутый в 1 файле(ах). Локальные сопроводительные пути: none. Источники: \path{publication/todo-backlog-table.tex}.\par\medskip
\textbf{\url{https://github.com/realagiorganization/dspy_rag_in_repo_docs_and_impl1/blob/master/samples/population/repository_population_candidates.yaml}}. Host: github.com. Occurrences: 1. External GitHub URL referenced from 1 file(s). Local companion path(s): none. Sources: \path{publication/todo-backlog-table.tex}.\\
\textbf{\url{https://github.com/realagiorganization/dspy_rag_in_repo_docs_and_impl1/blob/master/samples/population/repository_population_candidates.yaml}}. Хост: github.com. Упоминаний: 1. Внешний URL GitHub, упомянутый в 1 файле(ах). Локальные сопроводительные пути: none. Источники: \path{publication/todo-backlog-table.tex}.\par\medskip
\textbf{\url{https://github.com/realagiorganization/dspy_rag_in_repo_docs_and_impl1/blob/master/samples/training/repository_training_examples.yaml}}. Host: github.com. Occurrences: 1. External GitHub URL referenced from 1 file(s). Local companion path(s): none. Sources: \path{publication/todo-backlog-table.tex}.\\
\textbf{\url{https://github.com/realagiorganization/dspy_rag_in_repo_docs_and_impl1/blob/master/samples/training/repository_training_examples.yaml}}. Хост: github.com. Упоминаний: 1. Внешний URL GitHub, упомянутый в 1 файле(ах). Локальные сопроводительные пути: none. Источники: \path{publication/todo-backlog-table.tex}.\par\medskip
\textbf{\url{https://github.com/realagiorganization/dspy_rag_in_repo_docs_and_impl1/blob/master/src/repo_rag_lab/benchmarks.py}}. Host: github.com. Occurrences: 2. External GitHub URL referenced from 1 file(s). Local companion path(s): none. Sources: \path{publication/todo-backlog-table.tex}.\\
\textbf{\url{https://github.com/realagiorganization/dspy_rag_in_repo_docs_and_impl1/blob/master/src/repo_rag_lab/benchmarks.py}}. Хост: github.com. Упоминаний: 2. Внешний URL GitHub, упомянутый в 1 файле(ах). Локальные сопроводительные пути: none. Источники: \path{publication/todo-backlog-table.tex}.\par\medskip
\textbf{\url{https://github.com/realagiorganization/dspy_rag_in_repo_docs_and_impl1/blob/master/src/repo_rag_lab/cli.py\#L72-L88}}. Host: github.com. Occurrences: 1. External GitHub URL referenced from 1 file(s). Local companion path(s): none. Sources: \path{docs/operations/simple-end-to-end-verification-guide.md}.\\
\textbf{\url{https://github.com/realagiorganization/dspy_rag_in_repo_docs_and_impl1/blob/master/src/repo_rag_lab/cli.py\#L72-L88}}. Хост: github.com. Упоминаний: 1. Внешний URL GitHub, упомянутый в 1 файле(ах). Локальные сопроводительные пути: none. Источники: \path{docs/operations/simple-end-to-end-verification-guide.md}.\par\medskip
\textbf{\url{https://github.com/realagiorganization/dspy_rag_in_repo_docs_and_impl1/blob/master/src/repo_rag_lab/dspy_workflow.py}}. Host: github.com. Occurrences: 1. External GitHub URL referenced from 1 file(s). Local companion path(s): none. Sources: \path{publication/todo-backlog-table.tex}.\\
\textbf{\url{https://github.com/realagiorganization/dspy_rag_in_repo_docs_and_impl1/blob/master/src/repo_rag_lab/dspy_workflow.py}}. Хост: github.com. Упоминаний: 1. Внешний URL GitHub, упомянутый в 1 файле(ах). Локальные сопроводительные пути: none. Источники: \path{publication/todo-backlog-table.tex}.\par\medskip
\textbf{\url{https://github.com/realagiorganization/dspy_rag_in_repo_docs_and_impl1/blob/master/src/repo_rag_lab/notebook_scaffolding.py}}. Host: github.com. Occurrences: 1. External GitHub URL referenced from 1 file(s). Local companion path(s): none. Sources: \path{publication/todo-backlog-table.tex}.\\
\textbf{\url{https://github.com/realagiorganization/dspy_rag_in_repo_docs_and_impl1/blob/master/src/repo_rag_lab/notebook_scaffolding.py}}. Хост: github.com. Упоминаний: 1. Внешний URL GitHub, упомянутый в 1 файле(ах). Локальные сопроводительные пути: none. Источники: \path{publication/todo-backlog-table.tex}.\par\medskip
\textbf{\url{https://github.com/realagiorganization/dspy_rag_in_repo_docs_and_impl1/blob/master/src/repo_rag_lab/population_samples.py}}. Host: github.com. Occurrences: 1. External GitHub URL referenced from 1 file(s). Local companion path(s): none. Sources: \path{publication/todo-backlog-table.tex}.\\
\textbf{\url{https://github.com/realagiorganization/dspy_rag_in_repo_docs_and_impl1/blob/master/src/repo_rag_lab/population_samples.py}}. Хост: github.com. Упоминаний: 1. Внешний URL GitHub, упомянутый в 1 файле(ах). Локальные сопроводительные пути: none. Источники: \path{publication/todo-backlog-table.tex}.\par\medskip
\textbf{\url{https://github.com/realagiorganization/dspy_rag_in_repo_docs_and_impl1/blob/master/src/repo_rag_lab/retrieval.py}}. Host: github.com. Occurrences: 1. External GitHub URL referenced from 1 file(s). Local companion path(s): none. Sources: \path{publication/todo-backlog-table.tex}.\\
\textbf{\url{https://github.com/realagiorganization/dspy_rag_in_repo_docs_and_impl1/blob/master/src/repo_rag_lab/retrieval.py}}. Хост: github.com. Упоминаний: 1. Внешний URL GitHub, упомянутый в 1 файле(ах). Локальные сопроводительные пути: none. Источники: \path{publication/todo-backlog-table.tex}.\par\medskip
\textbf{\url{https://github.com/realagiorganization/dspy_rag_in_repo_docs_and_impl1/blob/master/src/repo_rag_lab/utilities.py\#L108-L124}}. Host: github.com. Occurrences: 2. External GitHub URL referenced from 1 file(s). Local companion path(s): none. Sources: \path{docs/operations/simple-end-to-end-verification-guide.md}.\\
\textbf{\url{https://github.com/realagiorganization/dspy_rag_in_repo_docs_and_impl1/blob/master/src/repo_rag_lab/utilities.py\#L108-L124}}. Хост: github.com. Упоминаний: 2. Внешний URL GitHub, упомянутый в 1 файле(ах). Локальные сопроводительные пути: none. Источники: \path{docs/operations/simple-end-to-end-verification-guide.md}.\par\medskip
\textbf{\url{https://github.com/realagiorganization/dspy_rag_in_repo_docs_and_impl1/blob/master/src/repo_rag_lab/utilities.py\#L190-L193}}. Host: github.com. Occurrences: 1. External GitHub URL referenced from 1 file(s). Local companion path(s): none. Sources: \path{docs/operations/simple-end-to-end-verification-guide.md}.\\
\textbf{\url{https://github.com/realagiorganization/dspy_rag_in_repo_docs_and_impl1/blob/master/src/repo_rag_lab/utilities.py\#L190-L193}}. Хост: github.com. Упоминаний: 1. Внешний URL GitHub, упомянутый в 1 файле(ах). Локальные сопроводительные пути: none. Источники: \path{docs/operations/simple-end-to-end-verification-guide.md}.\par\medskip
\textbf{\url{https://github.com/realagiorganization/dspy_rag_in_repo_docs_and_impl1/blob/master/src/repo_rag_lab/verification.py\#L11-L51}}. Host: github.com. Occurrences: 1. External GitHub URL referenced from 1 file(s). Local companion path(s): none. Sources: \path{docs/operations/simple-end-to-end-verification-guide.md}.\\
\textbf{\url{https://github.com/realagiorganization/dspy_rag_in_repo_docs_and_impl1/blob/master/src/repo_rag_lab/verification.py\#L11-L51}}. Хост: github.com. Упоминаний: 1. Внешний URL GitHub, упомянутый в 1 файле(ах). Локальные сопроводительные пути: none. Источники: \path{docs/operations/simple-end-to-end-verification-guide.md}.\par\medskip
\textbf{\url{https://github.com/realagiorganization/dspy_rag_in_repo_docs_and_impl1/blob/master/src/repo_rag_lab/verification.py\#L133-L146}}. Host: github.com. Occurrences: 1. External GitHub URL referenced from 1 file(s). Local companion path(s): none. Sources: \path{docs/operations/simple-end-to-end-verification-guide.md}.\\
\textbf{\url{https://github.com/realagiorganization/dspy_rag_in_repo_docs_and_impl1/blob/master/src/repo_rag_lab/verification.py\#L133-L146}}. Хост: github.com. Упоминаний: 1. Внешний URL GitHub, упомянутый в 1 файле(ах). Локальные сопроводительные пути: none. Источники: \path{docs/operations/simple-end-to-end-verification-guide.md}.\par\medskip
\textbf{\url{https://github.com/realagiorganization/dspy_rag_in_repo_docs_and_impl1/blob/master/src/repo_rag_lab/workflow.py}}. Host: github.com. Occurrences: 1. External GitHub URL referenced from 1 file(s). Local companion path(s): none. Sources: \path{publication/todo-backlog-table.tex}.\\
\textbf{\url{https://github.com/realagiorganization/dspy_rag_in_repo_docs_and_impl1/blob/master/src/repo_rag_lab/workflow.py}}. Хост: github.com. Упоминаний: 1. Внешний URL GitHub, упомянутый в 1 файле(ах). Локальные сопроводительные пути: none. Источники: \path{publication/todo-backlog-table.tex}.\par\medskip
\textbf{\url{https://github.com/realagiorganization/dspy_rag_in_repo_docs_and_impl1/blob/master/src/repo_rag_lab/workflow.py\#L119-L145}}. Host: github.com. Occurrences: 1. External GitHub URL referenced from 1 file(s). Local companion path(s): none. Sources: \path{docs/operations/simple-end-to-end-verification-guide.md}.\\
\textbf{\url{https://github.com/realagiorganization/dspy_rag_in_repo_docs_and_impl1/blob/master/src/repo_rag_lab/workflow.py\#L119-L145}}. Хост: github.com. Упоминаний: 1. Внешний URL GitHub, упомянутый в 1 файле(ах). Локальные сопроводительные пути: none. Источники: \path{docs/operations/simple-end-to-end-verification-guide.md}.\par\medskip
\textbf{\url{https://github.com/realagiorganization/dspy_rag_in_repo_docs_and_impl1/blob/master/src/repo_rag_lab/workflow.py\#L54-L72}}. Host: github.com. Occurrences: 1. External GitHub URL referenced from 1 file(s). Local companion path(s): none. Sources: \path{docs/operations/simple-end-to-end-verification-guide.md}.\\
\textbf{\url{https://github.com/realagiorganization/dspy_rag_in_repo_docs_and_impl1/blob/master/src/repo_rag_lab/workflow.py\#L54-L72}}. Хост: github.com. Упоминаний: 1. Внешний URL GitHub, упомянутый в 1 файле(ах). Локальные сопроводительные пути: none. Источники: \path{docs/operations/simple-end-to-end-verification-guide.md}.\par\medskip
\textbf{\url{https://github.com/realagiorganization/dspy_rag_in_repo_docs_and_impl1/blob/master/tests/test_benchmarks_and_notebook_scaffolding.py}}. Host: github.com. Occurrences: 1. External GitHub URL referenced from 1 file(s). Local companion path(s): none. Sources: \path{publication/todo-backlog-table.tex}.\\
\textbf{\url{https://github.com/realagiorganization/dspy_rag_in_repo_docs_and_impl1/blob/master/tests/test_benchmarks_and_notebook_scaffolding.py}}. Хост: github.com. Упоминаний: 1. Внешний URL GitHub, упомянутый в 1 файле(ах). Локальные сопроводительные пути: none. Источники: \path{publication/todo-backlog-table.tex}.\par\medskip
\textbf{\url{https://github.com/realagiorganization/dspy_rag_in_repo_docs_and_impl1/blob/master/tests/test_repository_rag_bdd.py}}. Host: github.com. Occurrences: 1. External GitHub URL referenced from 1 file(s). Local companion path(s): none. Sources: \path{publication/todo-backlog-table.tex}.\\
\textbf{\url{https://github.com/realagiorganization/dspy_rag_in_repo_docs_and_impl1/blob/master/tests/test_repository_rag_bdd.py}}. Хост: github.com. Упоминаний: 1. Внешний URL GitHub, упомянутый в 1 файле(ах). Локальные сопроводительные пути: none. Источники: \path{publication/todo-backlog-table.tex}.\par\medskip
\textbf{\url{https://github.com/realagiorganization/dspy_rag_in_repo_docs_and_impl1/blob/master/utilities/README.md}}. Host: github.com. Occurrences: 1. External GitHub URL referenced from 1 file(s). Local companion path(s): none. Sources: \path{publication/todo-backlog-table.tex}.\\
\textbf{\url{https://github.com/realagiorganization/dspy_rag_in_repo_docs_and_impl1/blob/master/utilities/README.md}}. Хост: github.com. Упоминаний: 1. Внешний URL GitHub, упомянутый в 1 файле(ах). Локальные сопроводительные пути: none. Источники: \path{publication/todo-backlog-table.tex}.\par\medskip
\textbf{\url{https://github.com/realagiorganization/dspy_rag_in_repo_docs_and_impl1/tree/master/samples/logs/}}. Host: github.com. Occurrences: 3. External GitHub URL referenced from 1 file(s). Local companion path(s): none. Sources: \path{publication/todo-backlog-table.tex}.\\
\textbf{\url{https://github.com/realagiorganization/dspy_rag_in_repo_docs_and_impl1/tree/master/samples/logs/}}. Хост: github.com. Упоминаний: 3. Внешний URL GitHub, упомянутый в 1 файле(ах). Локальные сопроводительные пути: none. Источники: \path{publication/todo-backlog-table.tex}.\par\medskip
\textbf{\url{https://github.com/realagiorganization/dspy_rag_in_repo_docs_and_impl1/tree/master/tests/fixtures/hushwheel_lexiconarium/}}. Host: github.com. Occurrences: 1. External GitHub URL referenced from 1 file(s). Local companion path(s): none. Sources: \path{publication/todo-backlog-table.tex}.\\
\textbf{\url{https://github.com/realagiorganization/dspy_rag_in_repo_docs_and_impl1/tree/master/tests/fixtures/hushwheel_lexiconarium/}}. Хост: github.com. Упоминаний: 1. Внешний URL GitHub, упомянутый в 1 файле(ах). Локальные сопроводительные пути: none. Источники: \path{publication/todo-backlog-table.tex}.\par\medskip
\textbf{\url{https://github.com/realagiorganization/national-debt-relief}}. Host: github.com. Occurrences: 1. External GitHub URL referenced from 1 file(s). Local companion path(s): none. Sources: \path{docs/audit/2026-05-10-aks-run-25632110510-handoff-fixed-runtime-still-heuristic.md}.\\
\textbf{\url{https://github.com/realagiorganization/national-debt-relief}}. Хост: github.com. Упоминаний: 1. Внешний URL GitHub, упомянутый в 1 файле(ах). Локальные сопроводительные пути: none. Источники: \path{docs/audit/2026-05-10-aks-run-25632110510-handoff-fixed-runtime-still-heuristic.md}.\par\medskip
\textbf{\url{https://gpt45standard.openai.azure.com/}}. Host: gpt45standard.openai.azure.com. Occurrences: 5. Referenced by URL only. No repository-local mirror is currently tracked. Sources: \path{docs/audit/2026-04-29-live-azure-gpt54-validation.md}.\\
\textbf{\url{https://gpt45standard.openai.azure.com/}}. Хост: gpt45standard.openai.azure.com. Упоминаний: 5. Ссылка хранится только как URL. Локальное зеркало в репозитории сейчас не отслеживается. Источники: \path{docs/audit/2026-04-29-live-azure-gpt54-validation.md}.\par\medskip
\textbf{\url{https://gpt45standard.openai.azure.com/`}}. Host: gpt45standard.openai.azure.com. Occurrences: 3. Referenced by URL only. No repository-local mirror is currently tracked. Sources: \path{README.md}\newline \path{docs/audit/2026-04-29-live-azure-gpt54-validation.md}\newline \path{docs/operations/azure-deployment.md}.\\
\textbf{\url{https://gpt45standard.openai.azure.com/`}}. Хост: gpt45standard.openai.azure.com. Упоминаний: 3. Ссылка хранится только как URL. Локальное зеркало в репозитории сейчас не отслеживается. Источники: \path{README.md}\newline \path{docs/audit/2026-04-29-live-azure-gpt54-validation.md}\newline \path{docs/operations/azure-deployment.md}.\par\medskip
\textbf{\url{https://gpt45standard.openai.azure.com/openai/v1}}. Host: gpt45standard.openai.azure.com. Occurrences: 1. Referenced by URL only. No repository-local mirror is currently tracked. Sources: \path{docs/audit/2026-05-01-dataset-azure-responses-api-version-default-fix.md}.\\
\textbf{\url{https://gpt45standard.openai.azure.com/openai/v1}}. Хост: gpt45standard.openai.azure.com. Упоминаний: 1. Ссылка хранится только как URL. Локальное зеркало в репозитории сейчас не отслеживается. Источники: \path{docs/audit/2026-05-01-dataset-azure-responses-api-version-default-fix.md}.\par\medskip
\textbf{\url{https://gpt45standard.openai.azure.com`}}. Host: gpt45standard.openai.azure.com`. Occurrences: 1. Referenced by URL only. No repository-local mirror is currently tracked. Sources: \path{docs/audit/2026-03-18-zzzzzz-dspy-pipeline-and-live-azure-proof.md}.\\
\textbf{\url{https://gpt45standard.openai.azure.com`}}. Хост: gpt45standard.openai.azure.com`. Упоминаний: 1. Ссылка хранится только как URL. Локальное зеркало в репозитории сейчас не отслеживается. Источники: \path{docs/audit/2026-03-18-zzzzzz-dspy-pipeline-and-live-azure-proof.md}.\par\medskip
\textbf{\url{https://img.shields.io/badge/coverage-94.12\%25-brightgreen)](https://github.com/realagiorganization/dspy_rag_in_repo_docs_and_impl1/actions/workflows/ci.yml}}. Host: img.shields.io. Occurrences: 1. Referenced by URL only. No repository-local mirror is currently tracked. Sources: \path{README.md}.\\
\textbf{\url{https://img.shields.io/badge/coverage-94.12\%25-brightgreen)](https://github.com/realagiorganization/dspy_rag_in_repo_docs_and_impl1/actions/workflows/ci.yml}}. Хост: img.shields.io. Упоминаний: 1. Ссылка хранится только как URL. Локальное зеркало в репозитории сейчас не отслеживается. Источники: \path{README.md}.\par\medskip
\textbf{\url{https://learn.microsoft.com/azure/ai-foundry/foundry-models/how-to/inference}}. Host: learn.microsoft.com. Occurrences: 1. External Azure documentation URL. No mirrored upstream copy is tracked, but local companion notes exist at docs/operations/azure-deployment.md. Sources: \path{docs/operations/azure-deployment.md}.\\
\textbf{\url{https://learn.microsoft.com/azure/ai-foundry/foundry-models/how-to/inference}}. Хост: learn.microsoft.com. Упоминаний: 1. Внешний URL документации Azure. Зеркальная копия upstream не хранится, но есть локальные сопроводительные заметки: docs/operations/azure-deployment.md. Источники: \path{docs/operations/azure-deployment.md}.\par\medskip
\textbf{\url{https://learn.microsoft.com/azure/ai-services/}}. Host: learn.microsoft.com. Occurrences: 4. External Azure documentation URL. No mirrored upstream copy is tracked, but local companion notes exist at docs/operations/azure-deployment.md. Sources: \path{publication/references.bib}\newline \path{tests/test_exploratorium_translation.py}\newline \path{tests/test_utilities.py}.\\
\textbf{\url{https://learn.microsoft.com/azure/ai-services/}}. Хост: learn.microsoft.com. Упоминаний: 4. Внешний URL документации Azure. Зеркальная копия upstream не хранится, но есть локальные сопроводительные заметки: docs/operations/azure-deployment.md. Источники: \path{publication/references.bib}\newline \path{tests/test_exploratorium_translation.py}\newline \path{tests/test_utilities.py}.\par\medskip
\textbf{\url{https://learn.microsoft.com/en-us/azure/ai-foundry/openai/how-to/fine-tuning-deploy}}. Host: learn.microsoft.com. Occurrences: 1. External Azure documentation URL. No mirrored upstream copy is tracked, but local companion notes exist at docs/operations/azure-deployment.md. Sources: \path{docs/operations/azure-deployment.md}.\\
\textbf{\url{https://learn.microsoft.com/en-us/azure/ai-foundry/openai/how-to/fine-tuning-deploy}}. Хост: learn.microsoft.com. Упоминаний: 1. Внешний URL документации Azure. Зеркальная копия upstream не хранится, но есть локальные сопроводительные заметки: docs/operations/azure-deployment.md. Источники: \path{docs/operations/azure-deployment.md}.\par\medskip
\textbf{\url{https://learn.microsoft.com/en-us/azure/ai-services/openai/how-to/fine-tuning}}. Host: learn.microsoft.com. Occurrences: 1. External Azure documentation URL. No mirrored upstream copy is tracked, but local companion notes exist at docs/operations/azure-deployment.md. Sources: \path{docs/operations/azure-deployment.md}.\\
\textbf{\url{https://learn.microsoft.com/en-us/azure/ai-services/openai/how-to/fine-tuning}}. Хост: learn.microsoft.com. Упоминаний: 1. Внешний URL документации Azure. Зеркальная копия upstream не хранится, но есть локальные сопроводительные заметки: docs/operations/azure-deployment.md. Источники: \path{docs/operations/azure-deployment.md}.\par\medskip
\textbf{\url{https://medium.com/@arancibia.juan22/implementing-rag-with-dspy-a-technical-guide-a6ae15f6a455}}. Host: medium.com. Occurrences: 1. Referenced by URL only. No repository-local mirror is currently tracked. Sources: \path{docs/architecture/inspired/implementing-rag-with-dspy-technical-guide.md}.\\
\textbf{\url{https://medium.com/@arancibia.juan22/implementing-rag-with-dspy-a-technical-guide-a6ae15f6a455}}. Хост: medium.com. Упоминаний: 1. Ссылка хранится только как URL. Локальное зеркало в репозитории сейчас не отслеживается. Источники: \path{docs/architecture/inspired/implementing-rag-with-dspy-technical-guide.md}.\par\medskip
\textbf{\url{https://modelcontextprotocol.io/}}. Host: modelcontextprotocol.io. Occurrences: 4. External MCP documentation URL. No mirrored upstream copy is tracked, but local companion notes exist at docs/architecture/mcp-discovery.md. Sources: \path{publication/references.bib}\newline \path{tests/test_exploratorium_translation.py}\newline \path{tests/test_utilities.py}.\\
\textbf{\url{https://modelcontextprotocol.io/}}. Хост: modelcontextprotocol.io. Упоминаний: 4. Внешний URL документации MCP. Зеркальная копия upstream не хранится, но есть локальные сопроводительные заметки: docs/architecture/mcp-discovery.md. Источники: \path{publication/references.bib}\newline \path{tests/test_exploratorium_translation.py}\newline \path{tests/test_utilities.py}.\par\medskip
\textbf{\url{https://realagiorganization.github.io/dspy_rag_in_repo_docs_and_impl1/}}. Host: realagiorganization.github.io. Occurrences: 3. Referenced by URL only. No repository-local mirror is currently tracked. Sources: \path{README.md}\newline \path{docs/audit/2026-03-19-public-pages-and-secret-scan.md}\newline \path{mkdocs.yml}.\\
\textbf{\url{https://realagiorganization.github.io/dspy_rag_in_repo_docs_and_impl1/}}. Хост: realagiorganization.github.io. Упоминаний: 3. Ссылка хранится только как URL. Локальное зеркало в репозитории сейчас не отслеживается. Источники: \path{README.md}\newline \path{docs/audit/2026-03-19-public-pages-and-secret-scan.md}\newline \path{mkdocs.yml}.\par\medskip
\textbf{\url{https://realagiorganization.github.io/dspy_rag_in_repo_docs_and_impl1/`}}. Host: realagiorganization.github.io. Occurrences: 1. Referenced by URL only. No repository-local mirror is currently tracked. Sources: \path{docs/audit/2026-03-19-public-pages-and-secret-scan.md}.\\
\textbf{\url{https://realagiorganization.github.io/dspy_rag_in_repo_docs_and_impl1/`}}. Хост: realagiorganization.github.io. Упоминаний: 1. Ссылка хранится только как URL. Локальное зеркало в репозитории сейчас не отслеживается. Источники: \path{docs/audit/2026-03-19-public-pages-and-secret-scan.md}.\par\medskip
\textbf{\url{https://realagistorage.blob.core.windows.net}}. Host: realagistorage.blob.core.windows.net. Occurrences: 1. Referenced by URL only. No repository-local mirror is currently tracked. Sources: \path{tests/test_runtime_artifacts_azure.py}.\\
\textbf{\url{https://realagistorage.blob.core.windows.net}}. Хост: realagistorage.blob.core.windows.net. Упоминаний: 1. Ссылка хранится только как URL. Локальное зеркало в репозитории сейчас не отслеживается. Источники: \path{tests/test_runtime_artifacts_azure.py}.\par\medskip
\textbf{\url{https://realagistorage.queue.core.windows.net}}. Host: realagistorage.queue.core.windows.net. Occurrences: 1. Referenced by URL only. No repository-local mirror is currently tracked. Sources: \path{tests/test_runtime_artifacts_azure.py}.\\
\textbf{\url{https://realagistorage.queue.core.windows.net}}. Хост: realagistorage.queue.core.windows.net. Упоминаний: 1. Ссылка хранится только как URL. Локальное зеркало в репозитории сейчас не отслеживается. Источники: \path{tests/test_runtime_artifacts_azure.py}.\par\medskip
\section{Page-By-Page / Постраничный показ}
\subsection*{English: Referenced Papers And Documentation Fetch State}
\textbf{lewis2020rag}: Retrieval-Augmented Generation for Knowledge-Intensive NLP Tasks. Kind: inproceedings. Cited in the bibliography, but no local PDF or paper mirror is tracked. Closest local companion paths: docs/architecture/dspy-rag-guide.md, docs/architecture/inspired/implementing-rag-with-dspy-technical-guide.md. Related paths: \path{docs/architecture/dspy-rag-guide.md, docs/architecture/inspired/implementing-rag-with-dspy-technical-guide.md}.\par\medskip
\textbf{khattab2024dspy}: DSPy: Compiling Declarative Language Model Calls into Self-Improving Pipelines. Kind: article. Cited in the bibliography, but no local PDF or paper mirror is tracked. Closest local companion paths: docs/architecture/dspy-rag-guide.md, docs/architecture/inspired/dspy-rag-tutorial.md, docs/guides/hushwheel-fixture-rag-guide.md. Related paths: \path{docs/architecture/dspy-rag-guide.md, docs/architecture/inspired/dspy-rag-tutorial.md, docs/guides/hushwheel-fixture-rag-guide.md}.\par\medskip
\textbf{mcp2024}: Model Context Protocol. Kind: misc. External documentation is cited by URL. No mirrored upstream copy is tracked, but local companion paths exist at docs/architecture/mcp-discovery.md. Related paths: \path{docs/architecture/mcp-discovery.md}.\par\medskip
\textbf{azureinference2025}: Azure AI Inference Documentation. Kind: misc. External documentation is cited by URL. No mirrored upstream copy is tracked, but local companion paths exist at docs/operations/azure-deployment.md. Related paths: \path{docs/operations/azure-deployment.md}.\par\medskip
\subsection*{English: Summaries Of All Files}
\textbf{\path{.codex/skills/exploratorium-translation-sync/SKILL.md}}. Kind: markdown. Size: 1722 bytes. Shape: 32 lines. Markdown document, 32 lines, 0 explicit URLs. Lead heading: Exploratorium Translation Sync.\par\medskip
\textbf{\path{.codex/skills/file-summary-sync/SKILL.md}}. Kind: markdown. Size: 1367 bytes. Shape: 29 lines. Markdown document, 29 lines, 0 explicit URLs. Lead heading: File Summary Sync.\par\medskip
\textbf{\path{.codex/skills/notebook-playbook-sync/SKILL.md}}. Kind: markdown. Size: 2085 bytes. Shape: 48 lines. Markdown document, 48 lines, 0 explicit URLs. Lead heading: Notebook Playbook Sync.\par\medskip
\textbf{\path{.codex/skills/notebook-playbook-sync/agents/openai.yaml}}. Kind: yaml. Size: 300 bytes. Shape: 5 lines. YAML data or workflow file, 5 lines, 0 explicit URLs. Top keys: interface.\par\medskip
\textbf{\path{.codex/skills/post-push-gh-run-logging/SKILL.md}}. Kind: markdown. Size: 2059 bytes. Shape: 41 lines. Markdown document, 41 lines, 0 explicit URLs. Lead heading: Post Push Gh Run Logging.\par\medskip
\textbf{\path{.codex/skills/post-push-gh-run-logging/agents/openai.yaml}}. Kind: yaml. Size: 281 bytes. Shape: 5 lines. YAML data or workflow file, 5 lines, 0 explicit URLs. Top keys: interface.\par\medskip
\textbf{\path{.codex/skills/post-push-gh-run-logging/scripts/write_gh_run_log.py}}. Kind: python. Size: 6346 bytes. Shape: 178 lines. Python module, 178 lines, 0 explicit URLs. Module intent: No module docstring detected.\par\medskip
\textbf{\path{.codex/skills/repo-verification-audit-loop/SKILL.md}}. Kind: markdown. Size: 2589 bytes. Shape: 53 lines. Markdown document, 53 lines, 0 explicit URLs. Lead heading: Repo Verification Audit Loop.\par\medskip
\textbf{\path{.codex/skills/repo-verification-audit-loop/agents/openai.yaml}}. Kind: yaml. Size: 272 bytes. Shape: 5 lines. YAML data or workflow file, 5 lines, 0 explicit URLs. Top keys: interface.\par\medskip
\textbf{\path{.codex/skills/rust-sqlite-lookup/SKILL.md}}. Kind: markdown. Size: 1297 bytes. Shape: 36 lines. Markdown document, 36 lines, 0 explicit URLs. Lead heading: Rust SQLite Lookup.\par\medskip
\textbf{\path{.codex/skills/todo-backlog-sync/SKILL.md}}. Kind: markdown. Size: 1289 bytes. Shape: 29 lines. Markdown document, 29 lines, 0 explicit URLs. Lead heading: Todo Backlog Sync.\par\medskip
\textbf{\path{.dockerignore}}. Kind: generic. Size: 150 bytes. Shape: 17 lines. Tracked repository file, 17 lines, 0 explicit URLs. Opening line: .git\par\medskip
\textbf{\path{.env.sample}}. Kind: generic. Size: 905 bytes. Shape: 17 lines. Tracked repository file, 17 lines, 0 explicit URLs. Opening line: \# Copy this file to .env and replace placeholder values.\par\medskip
\textbf{\path{.gitattributes}}. Kind: generic. Size: 126 bytes. Shape: 4 lines. Tracked repository file, 4 lines, 0 explicit URLs. Opening line: *.png filter=lfs diff=lfs merge=lfs -text\par\medskip
\textbf{\path{.github/workflows/ci.yml}}. Kind: yaml. Size: 4185 bytes. Shape: 146 lines. YAML data or workflow file, 146 lines, 0 explicit URLs. Top keys: name, True, jobs.\par\medskip
\textbf{\path{.github/workflows/hushwheel-quality.yml}}. Kind: yaml. Size: 3526 bytes. Shape: 104 lines. YAML data or workflow file, 104 lines, 0 explicit URLs. Top keys: name, True, env.\par\medskip
\textbf{\path{.github/workflows/pages.yml}}. Kind: yaml. Size: 2424 bytes. Shape: 99 lines. YAML data or workflow file, 99 lines, 0 explicit URLs. Top keys: name, True, env.\par\medskip
\textbf{\path{.github/workflows/publication-pdf.yml}}. Kind: yaml. Size: 5549 bytes. Shape: 144 lines. YAML data or workflow file, 144 lines, 0 explicit URLs. Top keys: name, True, env.\par\medskip
\textbf{\path{.github/workflows/publish.yml}}. Kind: yaml. Size: 847 bytes. Shape: 43 lines. YAML data or workflow file, 43 lines, 0 explicit URLs. Top keys: name, True, jobs.\par\medskip
\textbf{\path{.gitignore}}. Kind: generic. Size: 222 bytes. Shape: 21 lines. Tracked repository file, 21 lines, 0 explicit URLs. Opening line: .venv/\par\medskip
\textbf{\path{.pre-commit-config.yaml}}. Kind: yaml. Size: 1689 bytes. Shape: 53 lines. YAML data or workflow file, 53 lines, 0 explicit URLs. Top keys: repos.\par\medskip
\textbf{\path{.python-version}}. Kind: generic. Size: 5 bytes. Shape: 2 lines. Tracked repository file, 2 lines, 0 explicit URLs. Opening line: 3.11\par\medskip
\textbf{\path{.yamllint.yml}}. Kind: yaml. Size: 416 bytes. Shape: 20 lines. YAML data or workflow file, 20 lines, 0 explicit URLs. Top keys: extends, rules.\par\medskip
\textbf{\path{AGENTS.md}}. Kind: markdown. Size: 6999 bytes. Shape: 108 lines. Markdown document, 108 lines, 0 explicit URLs. Lead heading: Repository Agent Instructions.\par\medskip
\textbf{\path{AGENTS.md.d/00-audit-baseline.md}}. Kind: markdown. Size: 592 bytes. Shape: 15 lines. Markdown document, 15 lines, 0 explicit URLs. Lead heading: Audit Baseline Instructions.\par\medskip
\textbf{\path{AGENTS.md.d/10-verification-loop.md}}. Kind: markdown. Size: 1175 bytes. Shape: 22 lines. Markdown document, 22 lines, 0 explicit URLs. Lead heading: Verification Loop Instructions.\par\medskip
\textbf{\path{AGENTS.md.d/FILES.md}}. Kind: markdown. Size: 2247 bytes. Shape: 41 lines. Markdown document, 41 lines, 0 explicit URLs. Lead heading: File Summary Instructions.\par\medskip
\textbf{\path{AGENTS.md.d/RUST_LOOKUP.md}}. Kind: markdown. Size: 1803 bytes. Shape: 36 lines. Markdown document, 36 lines, 0 explicit URLs. Lead heading: Rust SQLite Lookup Instructions.\par\medskip
\textbf{\path{Dockerfile}}. Kind: generic. Size: 1080 bytes. Shape: 47 lines. Tracked repository file, 47 lines, 0 explicit URLs. Opening line: ARG PYTHON\_BASE\_IMAGE=python:3.11-slim-bookworm\par\medskip
\textbf{\path{FILES.csv}}. Kind: generic. Size: 52899 bytes. Shape: binary surface. Tracked repository file, 52899 bytes of binary data, 0 explicit URLs. Opening line: No non-empty preview line.\par\medskip
\textbf{\path{FILES.md}}. Kind: markdown. Size: 59726 bytes. Shape: 417 lines. Markdown document, 417 lines, 1 explicit URLs. Lead heading: Repository File Inventory.\par\medskip
\textbf{\path{Makefile}}. Kind: make. Size: 23701 bytes. Shape: 506 lines. Repository Makefile, 506 lines, 1 explicit URLs. It exposes the shared automation surfaces.\par\medskip
\textbf{\path{README.md}}. Kind: markdown. Size: 26222 bytes. Shape: 312 lines. Markdown document, 312 lines, 6 explicit URLs. Lead heading: Repository RAG Lab.\par\medskip
\textbf{\path{TODO.MD}}. Kind: markdown. Size: 6104 bytes. Shape: 22 lines. Markdown document, 22 lines, 0 explicit URLs. Lead heading: TODO.\par\medskip
\textbf{\path{config/retrieval-profile.json}}. Kind: json. Size: 1791 bytes. Shape: 66 lines. JSON data file, 66 lines, 0 explicit URLs. Top keys: name, retrieval\_mode, rerank\_candidate\_pool\_size.\par\medskip
\textbf{\path{data/questions/repository.yaml}}. Kind: yaml. Size: 283 bytes. Shape: 7 lines. YAML data or workflow file, 7 lines, 0 explicit URLs. Top keys: no top-level mapping keys.\par\medskip
\textbf{\path{docs/architecture/dspy-rag-guide.md}}. Kind: markdown. Size: 50272 bytes. Shape: 911 lines. Markdown document, 911 lines, 1 explicit URLs. Lead heading: DSPy Guide For This Repository.\par\medskip
\textbf{\path{docs/architecture/inspired/dspy-rag-tutorial.md}}. Kind: markdown. Size: 1474 bytes. Shape: 35 lines. Markdown document, 35 lines, 1 explicit URLs. Lead heading: DSPy Tutorial: Basic RAG.\par\medskip
\textbf{\path{docs/architecture/inspired/implementing-rag-with-dspy-technical-guide.md}}. Kind: markdown. Size: 1555 bytes. Shape: 38 lines. Markdown document, 38 lines, 1 explicit URLs. Lead heading: Implementing RAG with DSPy: Technical Guide.\par\medskip
\textbf{\path{docs/architecture/mcp-discovery.md}}. Kind: markdown. Size: 1757 bytes. Shape: 54 lines. Markdown document, 54 lines, 0 explicit URLs. Lead heading: MCP Discovery Notes.\par\medskip
\textbf{\path{docs/architecture/package-api.md}}. Kind: markdown. Size: 9186 bytes. Shape: 127 lines. Markdown document, 127 lines, 0 explicit URLs. Lead heading: Package API Notes.\par\medskip
\textbf{\path{docs/architecture/research-narrative.md}}. Kind: markdown. Size: 82061 bytes. Shape: 1077 lines. Markdown document, 1077 lines, 0 explicit URLs. Lead heading: Repository Research Narrative.\par\medskip
\textbf{\path{docs/archive/repo-completeness-checklist.md}}. Kind: markdown. Size: 7450 bytes. Shape: 238 lines. Markdown document, 238 lines, 0 explicit URLs. Lead heading: Repository Completeness Checklist.\par\medskip
\textbf{\path{docs/audit/2026-03-15-repo-state.md}}. Kind: markdown. Size: 5989 bytes. Shape: 151 lines. Markdown document, 151 lines, 0 explicit URLs. Lead heading: Repository State Audit.\par\medskip
\textbf{\path{docs/audit/2026-03-17-docs-and-uv-refresh.md}}. Kind: markdown. Size: 3977 bytes. Shape: 79 lines. Markdown document, 79 lines, 0 explicit URLs. Lead heading: Documentation And UV Workflow Refresh Audit.\par\medskip
\textbf{\path{docs/audit/2026-03-17-full-build.md}}. Kind: markdown. Size: 6715 bytes. Shape: 45 lines. Markdown document, 45 lines, 0 explicit URLs. Lead heading: Full Build Audit.\par\medskip
\textbf{\path{docs/audit/2026-03-17-gh-actions-watch-loop.md}}. Kind: markdown. Size: 3797 bytes. Shape: 78 lines. Markdown document, 78 lines, 0 explicit URLs. Lead heading: GitHub Actions Watch Loop Audit.\par\medskip
\textbf{\path{docs/audit/2026-03-17-hushwheel-fixture.md}}. Kind: markdown. Size: 5090 bytes. Shape: 80 lines. Markdown document, 80 lines, 0 explicit URLs. Lead heading: Hushwheel Fixture Audit.\par\medskip
\textbf{\path{docs/audit/2026-03-17-hushwheel-rag-playbook.md}}. Kind: markdown. Size: 4797 bytes. Shape: 83 lines. Markdown document, 83 lines, 0 explicit URLs. Lead heading: Hushwheel RAG Playbook Audit.\par\medskip
\textbf{\path{docs/audit/2026-03-17-notebook-scaffolding.md}}. Kind: markdown. Size: 3641 bytes. Shape: 70 lines. Markdown document, 70 lines, 0 explicit URLs. Lead heading: Notebook Scaffolding Audit.\par\medskip
\textbf{\path{docs/audit/2026-03-17-publication-article.md}}. Kind: markdown. Size: 2321 bytes. Shape: 54 lines. Markdown document, 54 lines, 0 explicit URLs. Lead heading: Publication Article Audit.\par\medskip
\textbf{\path{docs/audit/2026-03-17-repo-completeness-checklist.md}}. Kind: markdown. Size: 3114 bytes. Shape: 59 lines. Markdown document, 59 lines, 0 explicit URLs. Lead heading: Repository Completeness Checklist Audit.\par\medskip
\textbf{\path{docs/audit/2026-03-18-actions-caching.md}}. Kind: markdown. Size: 4189 bytes. Shape: 87 lines. Markdown document, 87 lines, 0 explicit URLs. Lead heading: GitHub Actions Caching Audit.\par\medskip
\textbf{\path{docs/audit/2026-03-18-azure-inference-endpoint-probe.md}}. Kind: markdown. Size: 2648 bytes. Shape: 63 lines. Markdown document, 63 lines, 0 explicit URLs. Lead heading: Azure Inference Endpoint Probe Audit.\par\medskip
\textbf{\path{docs/audit/2026-03-18-azure-runtime-surfaces.md}}. Kind: markdown. Size: 5307 bytes. Shape: 91 lines. Markdown document, 91 lines, 0 explicit URLs. Lead heading: Azure Runtime Surfaces Audit.\par\medskip
\textbf{\path{docs/audit/2026-03-18-dspy-training-path.md}}. Kind: markdown. Size: 4605 bytes. Shape: 81 lines. Markdown document, 81 lines, 0 explicit URLs. Lead heading: DSPy Training Path Audit.\par\medskip
\textbf{\path{docs/audit/2026-03-18-hushwheel-program-harness.md}}. Kind: markdown. Size: 3907 bytes. Shape: 76 lines. Markdown document, 76 lines, 0 explicit URLs. Lead heading: Hushwheel Program Harness Audit.\par\medskip
\textbf{\path{docs/audit/2026-03-18-notebook-execution.md}}. Kind: markdown. Size: 2577 bytes. Shape: 66 lines. Markdown document, 66 lines, 0 explicit URLs. Lead heading: Notebook Execution Audit.\par\medskip
\textbf{\path{docs/audit/2026-03-18-notebook-runner-harness.md}}. Kind: markdown. Size: 4121 bytes. Shape: 91 lines. Markdown document, 91 lines, 0 explicit URLs. Lead heading: Notebook Runner Harness Audit.\par\medskip
\textbf{\path{docs/audit/2026-03-18-repository-benchmark-broadening.md}}. Kind: markdown. Size: 3728 bytes. Shape: 74 lines. Markdown document, 74 lines, 0 explicit URLs. Lead heading: Repository Benchmark Broadening Audit.\par\medskip
\textbf{\path{docs/audit/2026-03-18-retest-with-env-refresh.md}}. Kind: markdown. Size: 5102 bytes. Shape: 105 lines. Markdown document, 105 lines, 0 explicit URLs. Lead heading: Env Refresh Retest Audit.\par\medskip
\textbf{\path{docs/audit/2026-03-18-retrieval-evaluation-suite.md}}. Kind: markdown. Size: 4031 bytes. Shape: 80 lines. Markdown document, 80 lines, 0 explicit URLs. Lead heading: Retrieval Evaluation Suite Audit.\par\medskip
\textbf{\path{docs/audit/2026-03-18-retrieval-ranking-refresh.md}}. Kind: markdown. Size: 6479 bytes. Shape: 109 lines. Markdown document, 109 lines, 0 explicit URLs. Lead heading: Retrieval Ranking Refresh Audit.\par\medskip
\textbf{\path{docs/audit/2026-03-18-rust-and-benchmark-hygiene.md}}. Kind: markdown. Size: 2049 bytes. Shape: 53 lines. Markdown document, 53 lines, 0 explicit URLs. Lead heading: Rust And Benchmark Hygiene Audit.\par\medskip
\textbf{\path{docs/audit/2026-03-18-todo-backlog-sync.md}}. Kind: markdown. Size: 4104 bytes. Shape: 83 lines. Markdown document, 83 lines, 0 explicit URLs. Lead heading: Todo Backlog Sync Audit.\par\medskip
\textbf{\path{docs/audit/2026-03-18-z-notebook-runner-harness.md}}. Kind: markdown. Size: 4153 bytes. Shape: 104 lines. Markdown document, 104 lines, 0 explicit URLs. Lead heading: Notebook Runner Harness Audit.\par\medskip
\textbf{\path{docs/audit/2026-03-18-zz-file-summary-inventory.md}}. Kind: markdown. Size: 2596 bytes. Shape: 69 lines. Markdown document, 69 lines, 0 explicit URLs. Lead heading: File Summary Inventory Audit.\par\medskip
\textbf{\path{docs/audit/2026-03-18-zz-research-narrative.md}}. Kind: markdown. Size: 3312 bytes. Shape: 79 lines. Markdown document, 79 lines, 0 explicit URLs. Lead heading: Research Narrative Audit.\par\medskip
\textbf{\path{docs/audit/2026-03-18-zzz-session-consolidation-on-origin-master.md}}. Kind: markdown. Size: 5061 bytes. Shape: 109 lines. Markdown document, 109 lines, 0 explicit URLs. Lead heading: Session Consolidation On Origin Master.\par\medskip
\textbf{\path{docs/audit/2026-03-18-zzzz-rust-sqlite-lookup.md}}. Kind: markdown. Size: 5404 bytes. Shape: 105 lines. Markdown document, 105 lines, 0 explicit URLs. Lead heading: Rust SQLite Lookup And Retrieval Tag Audit.\par\medskip
\textbf{\path{docs/audit/2026-03-18-zzzz-session-consolidation-master-ci-repair.md}}. Kind: markdown. Size: 3270 bytes. Shape: 73 lines. Markdown document, 73 lines, 0 explicit URLs. Lead heading: Session Consolidation Master CI Repair.\par\medskip
\textbf{\path{docs/audit/2026-03-18-zzzzz-hushwheel-quality-instrumentation.md}}. Kind: markdown. Size: 8775 bytes. Shape: 165 lines. Markdown document, 165 lines, 0 explicit URLs. Lead heading: Hushwheel Quality Instrumentation.\par\medskip
\textbf{\path{docs/audit/2026-03-18-zzzzzz-dspy-pipeline-and-live-azure-proof.md}}. Kind: markdown. Size: 5992 bytes. Shape: 116 lines. Markdown document, 116 lines, 1 explicit URLs. Lead heading: DSPy Pipeline And Live Azure Proof.\par\medskip
\textbf{\path{docs/audit/2026-03-18-zzzzzz-hushwheel-retrieval-record-wording.md}}. Kind: markdown. Size: 3042 bytes. Shape: 77 lines. Markdown document, 77 lines, 0 explicit URLs. Lead heading: Hushwheel Retrieval Record Wording.\par\medskip
\textbf{\path{docs/audit/2026-03-18-zzzzzzz-hushwheel-pdf-side-effect-fix.md}}. Kind: markdown. Size: 4044 bytes. Shape: 83 lines. Markdown document, 83 lines, 0 explicit URLs. Lead heading: Hushwheel PDF Side-Effect Fix.\par\medskip
\textbf{\path{docs/audit/2026-03-18-zzzzzzz-hushwheel-quality-salvage.md}}. Kind: markdown. Size: 4554 bytes. Shape: 84 lines. Markdown document, 84 lines, 0 explicit URLs. Lead heading: Hushwheel Quality Salvage Audit.\par\medskip
\textbf{\path{docs/audit/2026-03-18-zzzzzzz-land-unlanded-surfaces.md}}. Kind: markdown. Size: 2685 bytes. Shape: 72 lines. Markdown document, 72 lines, 0 explicit URLs. Lead heading: Hushwheel Quality Landing Audit.\par\medskip
\textbf{\path{docs/audit/2026-03-18-zzzzzzz-rebased-master-log-and-azure-typecheck-fix.md}}. Kind: markdown. Size: 3472 bytes. Shape: 83 lines. Markdown document, 83 lines, 0 explicit URLs. Lead heading: Rebased Master Log And Azure Typecheck Fix.\par\medskip
\textbf{\path{docs/audit/2026-03-18-zzzzzzz-worktree-consolidation.md}}. Kind: markdown. Size: 5914 bytes. Shape: 125 lines. Markdown document, 125 lines, 0 explicit URLs. Lead heading: Worktree Consolidation.\par\medskip
\textbf{\path{docs/audit/2026-03-18-zzzzzzzzzzzz-retrieval-regression-gate.md}}. Kind: markdown. Size: 2898 bytes. Shape: 66 lines. Markdown document, 66 lines, 0 explicit URLs. Lead heading: Retrieval Regression Gate Audit.\par\medskip
\textbf{\path{docs/audit/2026-03-18-zzzzzzzzzzzzz-land-remaining-unlanded-work.md}}. Kind: markdown. Size: 2908 bytes. Shape: 78 lines. Markdown document, 78 lines, 0 explicit URLs. Lead heading: Land Remaining Unlanded Work.\par\medskip
\textbf{\path{docs/audit/2026-03-18-zzzzzzzzzzzzzz-node24-actions-refresh.md}}. Kind: markdown. Size: 2989 bytes. Shape: 91 lines. Markdown document, 91 lines, 0 explicit URLs. Lead heading: Node 24 Actions Refresh.\par\medskip
\textbf{\path{docs/audit/2026-03-18-zzzzzzzzzzzzzz-rebase-hushwheel-quality-over-master.md}}. Kind: markdown. Size: 1889 bytes. Shape: 47 lines. Markdown document, 47 lines, 0 explicit URLs. Lead heading: Rebase Hushwheel Quality Over Master.\par\medskip
\textbf{\path{docs/audit/2026-03-18-zzzzzzzzzzzzzzzz-dspy-backlog-reconciliation.md}}. Kind: markdown. Size: 2956 bytes. Shape: 75 lines. Markdown document, 75 lines, 0 explicit URLs. Lead heading: DSPy Backlog Reconciliation.\par\medskip
\textbf{\path{docs/audit/2026-03-18-zzzzzzzzzzzzzzzz-file-summary-test-churn-fix.md}}. Kind: markdown. Size: 2929 bytes. Shape: 85 lines. Markdown document, 85 lines, 0 explicit URLs. Lead heading: File Summary Test Churn Fix.\par\medskip
\textbf{\path{docs/audit/2026-03-18-zzzzzzzzzzzzzzzz-publication-cache-node24-followup.md}}. Kind: markdown. Size: 2900 bytes. Shape: 89 lines. Markdown document, 89 lines, 0 explicit URLs. Lead heading: Publication Cache Node 24 Follow-Up.\par\medskip
\textbf{\path{docs/audit/2026-03-18-zzzzzzzzzzzzzzzz-retrieval-gate-restored-after-rebase.md}}. Kind: markdown. Size: 3788 bytes. Shape: 94 lines. Markdown document, 94 lines, 0 explicit URLs. Lead heading: Retrieval Gate Restored After Rebase.\par\medskip
\textbf{\path{docs/audit/2026-03-18-zzzzzzzzzzzzzzzzzz-dspy-backlog-reconciliation.md}}. Kind: markdown. Size: 2457 bytes. Shape: 62 lines. Markdown document, 62 lines, 0 explicit URLs. Lead heading: DSPy Backlog Reconciliation.\par\medskip
\textbf{\path{docs/audit/2026-03-18-zzzzzzzzzzzzzzzzzz-lookup-first-ask-path.md}}. Kind: markdown. Size: 4604 bytes. Shape: 94 lines. Markdown document, 94 lines, 0 explicit URLs. Lead heading: Lookup First Ask Path.\par\medskip
\textbf{\path{docs/audit/2026-03-18-zzzzzzzzzzzzzzzzzzz-answer-synthesis-refresh.md}}. Kind: markdown. Size: 2565 bytes. Shape: 48 lines. Markdown document, 48 lines, 0 explicit URLs. Lead heading: 2026-03-18 Answer Synthesis Refresh.\par\medskip
\textbf{\path{docs/audit/2026-03-18-zzzzzzzzzzzzzzzzzzzz-retrieval-quality-doc-priority.md}}. Kind: markdown. Size: 3602 bytes. Shape: 68 lines. Markdown document, 68 lines, 0 explicit URLs. Lead heading: Retrieval Quality Doc Priority Audit.\par\medskip
\textbf{\path{docs/audit/2026-03-18-zzzzzzzzzzzzzzzzzzzzzzz-merge-hushwheel-quality-into-master.md}}. Kind: markdown. Size: 4335 bytes. Shape: 112 lines. Markdown document, 112 lines, 0 explicit URLs. Lead heading: Merge Hushwheel Quality Into Master.\par\medskip
\textbf{\path{docs/audit/2026-03-18-zzzzzzzzzzzzzzzzzzzzzzzz-hushwheel-atlas-consolidation.md}}. Kind: markdown. Size: 4762 bytes. Shape: 109 lines. Markdown document, 109 lines, 0 explicit URLs. Lead heading: Hushwheel Atlas Consolidation.\par\medskip
\textbf{\path{docs/audit/2026-03-19-github-pr-gate-sync.md}}. Kind: markdown. Size: 6599 bytes. Shape: 136 lines. Markdown document, 136 lines, 0 explicit URLs. Lead heading: Master Consolidation And GitHub PR Gate Sync.\par\medskip
\textbf{\path{docs/audit/2026-03-19-node24-actions-opt-in.md}}. Kind: markdown. Size: 2929 bytes. Shape: 80 lines. Markdown document, 80 lines, 0 explicit URLs. Lead heading: Node 24 Actions Opt-In.\par\medskip
\textbf{\path{docs/audit/2026-03-19-public-pages-and-secret-scan.md}}. Kind: markdown. Size: 5804 bytes. Shape: 113 lines. Markdown document, 113 lines, 2 explicit URLs. Lead heading: Public Pages And Secret Scan.\par\medskip
\textbf{\path{docs/audit/2026-03-19-simple-end-to-end-verification-guide.md}}. Kind: markdown. Size: 4236 bytes. Shape: 103 lines. Markdown document, 103 lines, 0 explicit URLs. Lead heading: Simple End-To-End Verification Guide.\par\medskip
\textbf{\path{docs/audit/2026-04-28-cli-envelope-normalization.md}}. Kind: markdown. Size: 4028 bytes. Shape: 95 lines. Markdown document, 95 lines, 0 explicit URLs. Lead heading: CLI Envelope Normalization.\par\medskip
\textbf{\path{docs/audit/2026-04-28-doc-tree-restructure-and-planning.md}}. Kind: markdown. Size: 4309 bytes. Shape: 105 lines. Markdown document, 105 lines, 0 explicit URLs. Lead heading: Doc Tree Restructure And Planning.\par\medskip
\textbf{\path{docs/audit/2026-04-28-json-runtime-contract.md}}. Kind: markdown. Size: 5281 bytes. Shape: 131 lines. Markdown document, 131 lines, 0 explicit URLs. Lead heading: JSON Runtime Contract.\par\medskip
\textbf{\path{docs/audit/2026-04-28-live-azure-dspy-integration-gating.md}}. Kind: markdown. Size: 3562 bytes. Shape: 79 lines. Markdown document, 79 lines, 0 explicit URLs. Lead heading: Live Azure And DSPy Integration Gating.\par\medskip
\textbf{\path{docs/audit/2026-04-28-retrieval-profile-layer.md}}. Kind: markdown. Size: 4123 bytes. Shape: 92 lines. Markdown document, 92 lines, 0 explicit URLs. Lead heading: Retrieval Profile Layer.\par\medskip
\textbf{\path{docs/audit/2026-04-28-rust-lookup-arbitrary-repo-roots.md}}. Kind: markdown. Size: 4078 bytes. Shape: 88 lines. Markdown document, 88 lines, 0 explicit URLs. Lead heading: Rust Lookup Arbitrary Repo Roots.\par\medskip
\textbf{\path{docs/audit/2026-04-29-bounded-mcp-server-surface.md}}. Kind: markdown. Size: 6688 bytes. Shape: 154 lines. Markdown document, 154 lines, 0 explicit URLs. Lead heading: 2026-04-29 Bounded MCP Server Surface.\par\medskip
\textbf{\path{docs/audit/2026-04-29-bundle-channel-publish-promote-rollback.md}}. Kind: markdown. Size: 5633 bytes. Shape: 125 lines. Markdown document, 125 lines, 0 explicit URLs. Lead heading: 2026-04-29 Bundle Channel Publish Promote Rollback.\par\medskip
\textbf{\path{docs/audit/2026-04-29-bundle-overlay-and-runtime-trace-contract.md}}. Kind: markdown. Size: 4332 bytes. Shape: 96 lines. Markdown document, 96 lines, 0 explicit URLs. Lead heading: 2026-04-29 Bundle Overlay And Runtime Trace Contract.\par\medskip
\textbf{\path{docs/audit/2026-04-29-dataset-accepted-outcome-ingest.md}}. Kind: markdown. Size: 6670 bytes. Shape: 139 lines. Markdown document, 139 lines, 0 explicit URLs. Lead heading: 2026-04-29 Dataset Accepted-Outcome Ingest.\par\medskip
\textbf{\path{docs/audit/2026-04-29-dataset-bundle-fetch-and-trace-import-hooks.md}}. Kind: markdown. Size: 6215 bytes. Shape: 130 lines. Markdown document, 130 lines, 0 explicit URLs. Lead heading: 2026-04-29 Dataset Bundle Fetch And Trace Import Hooks.\par\medskip
\textbf{\path{docs/audit/2026-04-29-dataset-local-repo-rag-cli-backend.md}}. Kind: markdown. Size: 5290 bytes. Shape: 122 lines. Markdown document, 122 lines, 0 explicit URLs. Lead heading: 2026-04-29 Dataset Local `repo\_rag\_cli` Backend.\par\medskip
\textbf{\path{docs/audit/2026-04-29-dataset-repo-rag-auto-detection.md}}. Kind: markdown. Size: 6051 bytes. Shape: 124 lines. Markdown document, 124 lines, 0 explicit URLs. Lead heading: 2026-04-29 Dataset Repo-RAG Auto-Detection.\par\medskip
\textbf{\path{docs/audit/2026-04-29-dataset-trace-queue-handoff.md}}. Kind: markdown. Size: 6656 bytes. Shape: 141 lines. Markdown document, 141 lines, 0 explicit URLs. Lead heading: 2026-04-29 Dataset Trace Queue Handoff.\par\medskip
\textbf{\path{docs/audit/2026-04-29-dataset-worker-repo-rag-cli-backend.md}}. Kind: markdown. Size: 6630 bytes. Shape: 128 lines. Markdown document, 128 lines, 0 explicit URLs. Lead heading: 2026-04-29 Dataset Worker `repo\_rag\_cli` Backend.\par\medskip
\textbf{\path{docs/audit/2026-04-29-global-blob-queue-trainer-migration.md}}. Kind: markdown. Size: 4807 bytes. Shape: 91 lines. Markdown document, 91 lines, 0 explicit URLs. Lead heading: 2026-04-29 Global Blob And Queue Trainer Migration.\par\medskip
\textbf{\path{docs/audit/2026-04-29-live-azure-gpt54-validation.md}}. Kind: markdown. Size: 8279 bytes. Shape: 145 lines. Markdown document, 145 lines, 6 explicit URLs. Lead heading: 2026-04-29 Live Azure GPT-5.4 Validation.\par\medskip
\textbf{\path{docs/audit/2026-04-29-retrieval-rerank-and-relative-root-normalization.md}}. Kind: markdown. Size: 3070 bytes. Shape: 71 lines. Markdown document, 71 lines, 0 explicit URLs. Lead heading: 2026-04-29 Retrieval Rerank And Relative Root Normalization.\par\medskip
\textbf{\path{docs/audit/2026-04-29-runtime-image-and-dataset-build-deploy-wiring.md}}. Kind: markdown. Size: 5729 bytes. Shape: 140 lines. Markdown document, 140 lines, 0 explicit URLs. Lead heading: 2026-04-29 Runtime Image And Dataset Build/Deploy Wiring.\par\medskip
\textbf{\path{docs/audit/2026-04-29-trainer-bundle-benchmark-gates.md}}. Kind: markdown. Size: 4942 bytes. Shape: 111 lines. Markdown document, 111 lines, 0 explicit URLs. Lead heading: 2026-04-29 Trainer Bundle Benchmark Gates.\par\medskip
\textbf{\path{docs/audit/2026-04-29-trainer-cycle-background-pass.md}}. Kind: markdown. Size: 4637 bytes. Shape: 105 lines. Markdown document, 105 lines, 0 explicit URLs. Lead heading: 2026-04-29 Trainer Cycle Background Pass.\par\medskip
\textbf{\path{docs/audit/2026-04-29-trainer-k8s-deployment-surface.md}}. Kind: markdown. Size: 6040 bytes. Shape: 131 lines. Markdown document, 131 lines, 0 explicit URLs. Lead heading: 2026-04-29 Trainer Kubernetes Deployment Surface.\par\medskip
\textbf{\path{docs/audit/2026-04-29-trainer-recompile-from-candidates.md}}. Kind: markdown. Size: 6238 bytes. Shape: 135 lines. Markdown document, 135 lines, 0 explicit URLs. Lead heading: 2026-04-29 Trainer Recompile From Candidates.\par\medskip
\textbf{\path{docs/audit/2026-04-29-trainer-service-loop.md}}. Kind: markdown. Size: 6185 bytes. Shape: 133 lines. Markdown document, 133 lines, 0 explicit URLs. Lead heading: 2026-04-29 Trainer Service Loop.\par\medskip
\textbf{\path{docs/audit/2026-04-30-aks-run-25179828865-stale-image-inspection.md}}. Kind: markdown. Size: 6206 bytes. Shape: 117 lines. Markdown document, 117 lines, 0 explicit URLs. Lead heading: 2026-04-30 AKS Run 25179828865 Stale Image Inspection.\par\medskip
\textbf{\path{docs/audit/2026-04-30-aks-run-25186264345-image-fix-gap-inspection.md}}. Kind: markdown. Size: 8407 bytes. Shape: 165 lines. Markdown document, 165 lines, 0 explicit URLs. Lead heading: 2026-04-30 AKS Run 25186264345 Image/Fix Gap Inspection.\par\medskip
\textbf{\path{docs/audit/2026-04-30-codex-proxy-mediation-layer.md}}. Kind: markdown. Size: 5279 bytes. Shape: 95 lines. Markdown document, 95 lines, 0 explicit URLs. Lead heading: 2026-04-30 Codex Proxy Mediation Layer.\par\medskip
\textbf{\path{docs/audit/2026-04-30-codex-proxy-trace-handoff.md}}. Kind: markdown. Size: 6543 bytes. Shape: 122 lines. Markdown document, 122 lines, 0 explicit URLs. Lead heading: 2026-04-30 Codex Proxy Trace Handoff.\par\medskip
\textbf{\path{docs/audit/2026-05-01-aks-run-25191504814-project-root-resolution-inspection.md}}. Kind: markdown. Size: 6002 bytes. Shape: 133 lines. Markdown document, 133 lines, 0 explicit URLs. Lead heading: 2026-05-01 AKS Run 25191504814 Project Root Resolution Inspection.\par\medskip
\textbf{\path{docs/audit/2026-05-01-aks-run-25209573387-proxy-runtime-env-inspection.md}}. Kind: markdown. Size: 5738 bytes. Shape: 129 lines. Markdown document, 129 lines, 0 explicit URLs. Lead heading: 2026-05-01 AKS Run 25209573387 Proxy Runtime Env Inspection.\par\medskip
\textbf{\path{docs/audit/2026-05-01-aks-run-25212955759-trace-handoff-secret-gap.md}}. Kind: markdown. Size: 4594 bytes. Shape: 121 lines. Markdown document, 121 lines, 0 explicit URLs. Lead heading: AKS Run `25212955759`: Trace Handoff Skipped Because `repo-rag-storage-config` Lacks Blob Credentials.\par\medskip
\textbf{\path{docs/audit/2026-05-01-aks-run-25226558751-rag-heuristic-trusted-handoff.md}}. Kind: markdown. Size: 5092 bytes. Shape: 130 lines. Markdown document, 130 lines, 0 explicit URLs. Lead heading: AKS Run 25226558751: RAG Success, Heuristic DSPy Fallback, Trusted Handoff Success.\par\medskip
\textbf{\path{docs/audit/2026-05-01-dataset-azure-responses-api-version-default-fix.md}}. Kind: markdown. Size: 6700 bytes. Shape: 133 lines. Markdown document, 133 lines, 1 explicit URLs. Lead heading: 2026-05-01 Dataset Azure Responses API-Version Default Fix.\par\medskip
\textbf{\path{docs/audit/2026-05-01-dataset-azure-runtime-env-propagation-fix.md}}. Kind: markdown. Size: 4200 bytes. Shape: 81 lines. Markdown document, 81 lines, 0 explicit URLs. Lead heading: 2026-05-01 Dataset Azure Runtime Env Propagation Fix.\par\medskip
\textbf{\path{docs/audit/2026-05-01-dataset-trusted-trace-handoff-postprocess-fix.md}}. Kind: markdown. Size: 3880 bytes. Shape: 81 lines. Markdown document, 81 lines, 0 explicit URLs. Lead heading: Dataset Trusted Trace Handoff Post-Processing Fix.\par\medskip
\textbf{\path{docs/audit/2026-05-01-dspy-bundle-version-pinning-contract.md}}. Kind: markdown. Size: 5763 bytes. Shape: 118 lines. Markdown document, 118 lines, 0 explicit URLs. Lead heading: DSPY Bundle Version Pinning Contract.\par\medskip
\textbf{\path{docs/audit/2026-05-01-durable-trainer-progress-remote-recovery.md}}. Kind: markdown. Size: 5141 bytes. Shape: 115 lines. Markdown document, 115 lines, 0 explicit URLs. Lead heading: Durable Trainer Progress Remote Recovery.\par\medskip
\textbf{\path{docs/audit/2026-05-01-live-trainer-no-timestamp-bundle-publish.md}}. Kind: markdown. Size: 7368 bytes. Shape: 172 lines. Markdown document, 172 lines, 0 explicit URLs. Lead heading: Live Trainer No Timestamp Bundle Publish.\par\medskip
\textbf{\path{docs/audit/2026-05-01-live-trainer-remote-bundle-publish.md}}. Kind: markdown. Size: 7444 bytes. Shape: 148 lines. Markdown document, 148 lines, 0 explicit URLs. Lead heading: Live Trainer Remote Bundle Publish.\par\medskip
\textbf{\path{docs/audit/2026-05-01-live-trainer-service-azure-fallback.md}}. Kind: markdown. Size: 6390 bytes. Shape: 167 lines. Markdown document, 167 lines, 0 explicit URLs. Lead heading: Live Trainer Service Azure Fallback.\par\medskip
\textbf{\path{docs/audit/2026-05-01-versioned-trainer-lineage-bundles.md}}. Kind: markdown. Size: 5821 bytes. Shape: 123 lines. Markdown document, 123 lines, 0 explicit URLs. Lead heading: Versioned Trainer Lineage Bundles.\par\medskip
\textbf{\path{docs/audit/2026-05-02-hybrid-vector-retrieval-default.md}}. Kind: markdown. Size: 18302 bytes. Shape: 385 lines. Markdown document, 385 lines, 0 explicit URLs. Lead heading: Hybrid Vector Retrieval Default.\par\medskip
\textbf{\path{docs/audit/2026-05-02-live-trainer-still-not-publishing-bundles.md}}. Kind: markdown. Size: 60556 bytes. Shape: 1281 lines. Markdown document, 1281 lines, 0 explicit URLs. Lead heading: Live Trainer Still Not Publishing Bundles.\par\medskip
\textbf{\path{docs/audit/2026-05-02-trainer-context-group-champion-index-stage1.md}}. Kind: markdown. Size: 18201 bytes. Shape: 330 lines. Markdown document, 330 lines, 0 explicit URLs. Lead heading: Trainer Context-Group Champion Index Stage 1.\par\medskip
\textbf{\path{docs/audit/2026-05-02-zz-codex-exec-resume-session-state-stage0.md}}. Kind: markdown. Size: 170338 bytes. Shape: 3866 lines. Markdown document, 3866 lines, 2 explicit URLs. Lead heading: Repository audit note for 2026-05-02 codex exec resume session state stage 0.\par\medskip
\textbf{\path{docs/audit/2026-05-08-zz-per-turn-dspy-mediation-contract-stage0.md}}. Kind: markdown. Size: 5821 bytes. Shape: 113 lines. Markdown document, 113 lines, 0 explicit URLs. Lead heading: Repository audit note for 2026-05-08 per-turn DSPy mediation contract stage 0.\par\medskip
\textbf{\path{docs/audit/2026-05-09-aks-run-25598675829-per-turn-mediation-gap.md}}. Kind: markdown. Size: 7119 bytes. Shape: 215 lines. Markdown document, 215 lines, 0 explicit URLs. Lead heading: Repository audit note for 2026-05-09 AKS run 25598675829 per-turn mediation gap.\par\medskip
\textbf{\path{docs/audit/2026-05-09-family-first-mipro-contract-stage1.md}}. Kind: markdown. Size: 4684 bytes. Shape: 120 lines. Markdown document, 120 lines, 0 explicit URLs. Lead heading: Repository audit note for 2026-05-09 family-first MIPROv2 contract stage 1.\par\medskip
\textbf{\path{docs/audit/2026-05-09-family-first-mipro-contract-stage10.md}}. Kind: markdown. Size: 5329 bytes. Shape: 122 lines. Markdown document, 122 lines, 0 explicit URLs. Lead heading: Repository audit note for 2026-05-09 family-first MIPROv2 contract stage 10.\par\medskip
\textbf{\path{docs/audit/2026-05-09-family-first-mipro-contract-stage11.md}}. Kind: markdown. Size: 4638 bytes. Shape: 116 lines. Markdown document, 116 lines, 0 explicit URLs. Lead heading: Repository audit note for 2026-05-09 family-first MIPROv2 contract stage 11.\par\medskip
\textbf{\path{docs/audit/2026-05-09-family-first-mipro-contract-stage12.md}}. Kind: markdown. Size: 4941 bytes. Shape: 122 lines. Markdown document, 122 lines, 0 explicit URLs. Lead heading: Repository audit note for 2026-05-09 family-first MIPROv2 contract stage 12.\par\medskip
\textbf{\path{docs/audit/2026-05-09-family-first-mipro-contract-stage13.md}}. Kind: markdown. Size: 4992 bytes. Shape: 128 lines. Markdown document, 128 lines, 0 explicit URLs. Lead heading: Repository audit note for 2026-05-09 family-first MIPROv2 contract stage 13.\par\medskip
\textbf{\path{docs/audit/2026-05-09-family-first-mipro-contract-stage14.md}}. Kind: markdown. Size: 3944 bytes. Shape: 109 lines. Markdown document, 109 lines, 0 explicit URLs. Lead heading: Repository audit note for 2026-05-09 family-first MIPROv2 contract stage 14.\par\medskip
\textbf{\path{docs/audit/2026-05-09-family-first-mipro-contract-stage15.md}}. Kind: markdown. Size: 4728 bytes. Shape: 113 lines. Markdown document, 113 lines, 0 explicit URLs. Lead heading: Repository audit note for 2026-05-09 family-first MIPROv2 contract stage 15.\par\medskip
\textbf{\path{docs/audit/2026-05-09-family-first-mipro-contract-stage16.md}}. Kind: markdown. Size: 3748 bytes. Shape: 103 lines. Markdown document, 103 lines, 0 explicit URLs. Lead heading: Repository audit note for 2026-05-09 family-first MIPROv2 contract stage 16.\par\medskip
\textbf{\path{docs/audit/2026-05-09-family-first-mipro-contract-stage17.md}}. Kind: markdown. Size: 3519 bytes. Shape: 101 lines. Markdown document, 101 lines, 0 explicit URLs. Lead heading: Repository audit note for 2026-05-09 family-first MIPROv2 contract stage 17.\par\medskip
\textbf{\path{docs/audit/2026-05-09-family-first-mipro-contract-stage18.md}}. Kind: markdown. Size: 3164 bytes. Shape: 90 lines. Markdown document, 90 lines, 0 explicit URLs. Lead heading: Repository audit note for 2026-05-09 family-first MIPROv2 contract stage 18.\par\medskip
\textbf{\path{docs/audit/2026-05-09-family-first-mipro-contract-stage19.md}}. Kind: markdown. Size: 3984 bytes. Shape: 102 lines. Markdown document, 102 lines, 0 explicit URLs. Lead heading: Repository audit note for 2026-05-09 family-first MIPROv2 contract stage 19.\par\medskip
\textbf{\path{docs/audit/2026-05-09-family-first-mipro-contract-stage2.md}}. Kind: markdown. Size: 4936 bytes. Shape: 123 lines. Markdown document, 123 lines, 0 explicit URLs. Lead heading: Repository audit note for 2026-05-09 family-first MIPROv2 contract stage 2.\par\medskip
\textbf{\path{docs/audit/2026-05-09-family-first-mipro-contract-stage20.md}}. Kind: markdown. Size: 3076 bytes. Shape: 88 lines. Markdown document, 88 lines, 0 explicit URLs. Lead heading: Repository audit note for 2026-05-09 family-first MIPROv2 contract stage 20.\par\medskip
\textbf{\path{docs/audit/2026-05-09-family-first-mipro-contract-stage21.md}}. Kind: markdown. Size: 5742 bytes. Shape: 122 lines. Markdown document, 122 lines, 0 explicit URLs. Lead heading: Repository audit note for 2026-05-09 family-first MIPROv2 contract stage 21.\par\medskip
\textbf{\path{docs/audit/2026-05-09-family-first-mipro-contract-stage3.md}}. Kind: markdown. Size: 4147 bytes. Shape: 96 lines. Markdown document, 96 lines, 0 explicit URLs. Lead heading: Repository audit note for 2026-05-09 family-first MIPROv2 contract stage 3.\par\medskip
\textbf{\path{docs/audit/2026-05-09-family-first-mipro-contract-stage4.md}}. Kind: markdown. Size: 5209 bytes. Shape: 120 lines. Markdown document, 120 lines, 0 explicit URLs. Lead heading: Repository audit note for 2026-05-09 family-first MIPROv2 contract stage 4.\par\medskip
\textbf{\path{docs/audit/2026-05-09-family-first-mipro-contract-stage5.md}}. Kind: markdown. Size: 5172 bytes. Shape: 122 lines. Markdown document, 122 lines, 0 explicit URLs. Lead heading: Repository audit note for 2026-05-09 family-first MIPROv2 contract stage 5.\par\medskip
\textbf{\path{docs/audit/2026-05-09-family-first-mipro-contract-stage6.md}}. Kind: markdown. Size: 5120 bytes. Shape: 119 lines. Markdown document, 119 lines, 0 explicit URLs. Lead heading: Repository audit note for 2026-05-09 family-first MIPROv2 contract stage 6.\par\medskip
\textbf{\path{docs/audit/2026-05-09-family-first-mipro-contract-stage7.md}}. Kind: markdown. Size: 4680 bytes. Shape: 116 lines. Markdown document, 116 lines, 0 explicit URLs. Lead heading: Repository audit note for 2026-05-09 family-first MIPROv2 contract stage 7.\par\medskip
\textbf{\path{docs/audit/2026-05-09-family-first-mipro-contract-stage8.md}}. Kind: markdown. Size: 4477 bytes. Shape: 110 lines. Markdown document, 110 lines, 0 explicit URLs. Lead heading: Repository audit note for 2026-05-09 family-first MIPROv2 contract stage 8.\par\medskip
\textbf{\path{docs/audit/2026-05-09-family-first-mipro-contract-stage9.md}}. Kind: markdown. Size: 4003 bytes. Shape: 106 lines. Markdown document, 106 lines, 0 explicit URLs. Lead heading: Repository audit note for 2026-05-09 family-first MIPROv2 contract stage 9.\par\medskip
\textbf{\path{docs/audit/2026-05-10-aks-run-25625638264-family-first-runtime-gap.md}}. Kind: markdown. Size: 8472 bytes. Shape: 246 lines. Markdown document, 246 lines, 0 explicit URLs. Lead heading: Repository audit note for 2026-05-10 AKS run 25625638264 family-first runtime gap.\par\medskip
\textbf{\path{docs/audit/2026-05-10-aks-run-25629990035-family-first-runtime-still-heuristic.md}}. Kind: markdown. Size: 6843 bytes. Shape: 211 lines. Markdown document, 211 lines, 0 explicit URLs. Lead heading: Repository audit note for 2026-05-10 AKS run 25629990035 family-first runtime still heuristic.\par\medskip
\textbf{\path{docs/audit/2026-05-10-aks-run-25632110510-handoff-fixed-runtime-still-heuristic.md}}. Kind: markdown. Size: 7832 bytes. Shape: 231 lines. Markdown document, 231 lines, 1 explicit URLs. Lead heading: Repository audit note for 2026-05-10 AKS run 25632110510 handoff fixed but runtime still heuristic.\par\medskip
\textbf{\path{docs/audit/2026-05-10-aks-run-25634202380-reset-lane-batch-handoff-runtime-still-heuristic.md}}. Kind: markdown. Size: 6096 bytes. Shape: 197 lines. Markdown document, 197 lines, 0 explicit URLs. Lead heading: Repository audit note for AKS run 25634202380 on 2026-05-10.\par\medskip
\textbf{\path{docs/audit/2026-05-10-family-first-runtime-fixes-after-run-25625638264.md}}. Kind: markdown. Size: 12143 bytes. Shape: 261 lines. Markdown document, 261 lines, 0 explicit URLs. Lead heading: Repository audit note for 2026-05-10 family-first runtime fixes after run 25625638264.\par\medskip
\textbf{\path{docs/audit/2026-05-10-family-first-runtime-fixes-after-run-25629990035.md}}. Kind: markdown. Size: 4170 bytes. Shape: 98 lines. Markdown document, 98 lines, 0 explicit URLs. Lead heading: Repository audit note for 2026-05-10 family-first runtime fixes after AKS run 25629990035.\par\medskip
\textbf{\path{docs/audit/2026-05-10-family-first-runtime-fixes-after-run-25632110510.md}}. Kind: markdown. Size: 5332 bytes. Shape: 112 lines. Markdown document, 112 lines, 0 explicit URLs. Lead heading: Repository audit note for 2026-05-10 family-first runtime fixes after AKS run 25632110510.\par\medskip
\textbf{\path{docs/audit/2026-05-10-family-state-populates-but-publish-gate-and-dirty-fathers-still-block-runtime.md}}. Kind: markdown. Size: 7106 bytes. Shape: 183 lines. Markdown document, 183 lines, 0 explicit URLs. Lead heading: 2026-05-10 Family State Populates But Publish Gate And Dirty Fathers Still Block Runtime.\par\medskip
\textbf{\path{docs/audit/2026-05-10-family-state-root-alias-and-latest-bundle-staging-fallback.md}}. Kind: markdown. Size: 3390 bytes. Shape: 81 lines. Markdown document, 81 lines, 0 explicit URLs. Lead heading: 2026-05-10 Family-State Root Alias And Latest-Bundle Staging Fallback.\par\medskip
\textbf{\path{docs/audit/2026-05-10-hidden-token-spend-and-incremental-trainer-fixes.md}}. Kind: markdown. Size: 6740 bytes. Shape: 154 lines. Markdown document, 154 lines, 0 explicit URLs. Lead heading: 2026-05-10 Hidden Token Spend And Incremental Trainer Fixes.\par\medskip
\textbf{\path{docs/audit/2026-05-10-live-trainer-service-replays-processed-ledger-while-queue-empty.md}}. Kind: markdown. Size: 5891 bytes. Shape: 160 lines. Markdown document, 160 lines, 0 explicit URLs. Lead heading: 2026-05-10 Live Trainer Service Replays Processed Ledger While Queue Is Empty.\par\medskip
\textbf{\path{docs/audit/2026-05-10-stale-submodule-runtime-image-root-cause-and-source-preference-fix.md}}. Kind: markdown. Size: 5021 bytes. Shape: 114 lines. Markdown document, 114 lines, 0 explicit URLs. Lead heading: 2026-05-10 Stale Submodule Runtime Image Root Cause And Source-Preference Fix.\par\medskip
\textbf{\path{docs/audit/2026-05-11-family-runtime-artifact-fallback-and-staging-path-fix.md}}. Kind: markdown. Size: 5601 bytes. Shape: 139 lines. Markdown document, 139 lines, 0 explicit URLs. Lead heading: 2026-05-11 Family Runtime Artifact Fallback And Bundle Staging Path Fix.\par\medskip
\textbf{\path{docs/audit/2026-05-11-fresh-run-family-match-without-family-artifact-runtime.md}}. Kind: markdown. Size: 6472 bytes. Shape: 211 lines. Markdown document, 211 lines, 0 explicit URLs. Lead heading: 2026-05-11 Fresh Run Family Match Without Family Artifact Runtime.\par\medskip
\textbf{\path{docs/audit/2026-05-11-idle-trainer-cycles-no-longer-autorepublish-bundles.md}}. Kind: markdown. Size: 3162 bytes. Shape: 108 lines. Markdown document, 108 lines, 0 explicit URLs. Lead heading: 2026-05-11 Idle Trainer Cycles No Longer Auto-Republish Bundle Versions.\par\medskip
\textbf{\path{docs/audit/2026-05-11-original-prompt-routing-and-batch-handoff-fixes.md}}. Kind: markdown. Size: 3829 bytes. Shape: 91 lines. Markdown document, 91 lines, 0 explicit URLs. Lead heading: 2026-05-11 Original Prompt Routing And Batch Handoff Fixes.\par\medskip
\textbf{\path{docs/audit/2026-05-11-trainer-queue-only-no-processed-replay.md}}. Kind: markdown. Size: 6373 bytes. Shape: 166 lines. Markdown document, 166 lines, 0 explicit URLs. Lead heading: 2026-05-11 Trainer Is Now Queue-Only And Will Not Replay Processed History.\par\medskip
\textbf{\path{docs/audit/2026-05-11-trainer-service-preflights-queue-before-running-cycles.md}}. Kind: markdown. Size: 4211 bytes. Shape: 121 lines. Markdown document, 121 lines, 0 explicit URLs. Lead heading: 2026-05-11 Trainer Service Preflights Queue Before Running Cycles.\par\medskip
\textbf{\path{docs/audit/2026-05-11-trusted-handoff-syntax-error-blocked-execution-artifact-upload.md}}. Kind: markdown. Size: 4031 bytes. Shape: 86 lines. Markdown document, 86 lines, 0 explicit URLs. Lead heading: 2026-05-11 Trusted Handoff Syntax Error Blocked Execution Artifact Upload.\par\medskip
\textbf{\path{docs/audit/README.md}}. Kind: markdown. Size: 443 bytes. Shape: 14 lines. Markdown document, 14 lines, 0 explicit URLs. Lead heading: Repository Audit Notes.\par\medskip
\textbf{\path{docs/guides/hushwheel-fixture-rag-guide.md}}. Kind: markdown. Size: 5169 bytes. Shape: 138 lines. Markdown document, 138 lines, 0 explicit URLs. Lead heading: RAG Guide For The Hushwheel Fixture.\par\medskip
\textbf{\path{docs/operations/azure-deployment.md}}. Kind: markdown. Size: 4901 bytes. Shape: 111 lines. Markdown document, 111 lines, 5 explicit URLs. Lead heading: Azure Deployment Notes.\par\medskip
\textbf{\path{docs/operations/environment.md}}. Kind: markdown. Size: 7589 bytes. Shape: 130 lines. Markdown document, 130 lines, 0 explicit URLs. Lead heading: Environment Variables.\par\medskip
\textbf{\path{docs/operations/runtime-image.md}}. Kind: markdown. Size: 1989 bytes. Shape: 46 lines. Markdown document, 46 lines, 0 explicit URLs. Lead heading: Runtime Image Notes.\par\medskip
\textbf{\path{docs/operations/simple-end-to-end-verification-guide.md}}. Kind: markdown. Size: 4111 bytes. Shape: 105 lines. Markdown document, 105 lines, 12 explicit URLs. Lead heading: Simple End-To-End Verification Guide.\par\medskip
\textbf{\path{docs/operations/trainer-deployment.md}}. Kind: markdown. Size: 6047 bytes. Shape: 152 lines. Markdown document, 152 lines, 0 explicit URLs. Lead heading: Trainer Deployment Notes.\par\medskip
\textbf{\path{docs/planning/codex-exec-resume-plan.md}}. Kind: markdown. Size: 11561 bytes. Shape: 207 lines. Markdown document, 207 lines, 0 explicit URLs. Lead heading: Codex Exec Resume Plan.\par\medskip
\textbf{\path{docs/planning/dataset-integration-plan.md}}. Kind: markdown. Size: 33882 bytes. Shape: 478 lines. Markdown document, 478 lines, 0 explicit URLs. Lead heading: Dataset Integration Plan.\par\medskip
\textbf{\path{docs/planning/family-first-mipro-runtime-contract.md}}. Kind: markdown. Size: 29074 bytes. Shape: 439 lines. Markdown document, 439 lines, 0 explicit URLs. Lead heading: Family-First DSPy / MIPROv2 Contract.\par\medskip
\textbf{\path{docs/planning/per-turn-dspy-mediation-contract.md}}. Kind: markdown. Size: 4966 bytes. Shape: 115 lines. Markdown document, 115 lines, 0 explicit URLs. Lead heading: Per-Turn DSPy Mediation Contract.\par\medskip
\textbf{\path{docs/planning/repo-hardening-plan.md}}. Kind: markdown. Size: 4545 bytes. Shape: 71 lines. Markdown document, 71 lines, 0 explicit URLs. Lead heading: Repository Hardening Plan.\par\medskip
\textbf{\path{docs/planning/trainer-context-group-champion-plan.md}}. Kind: markdown. Size: 7736 bytes. Shape: 119 lines. Markdown document, 119 lines, 0 explicit URLs. Lead heading: Trainer Context-Group Champion Plan.\par\medskip
\textbf{\path{mkdocs.yml}}. Kind: yaml. Size: 843 bytes. Shape: 37 lines. YAML data or workflow file, 37 lines, 2 explicit URLs. Top keys: site\_name, site\_description, site\_url.\par\medskip
\textbf{\path{notebooks/01_repo_rag_research.ipynb}}. Kind: notebook. Size: 9324 bytes. Shape: 281 lines. Notebook, 11 cells, 0 explicit URLs. Opening heading: Repository RAG Research Playbook.\par\medskip
\textbf{\path{notebooks/02_agent_workflow_checklist.ipynb}}. Kind: notebook. Size: 17632 bytes. Shape: 404 lines. Notebook, 9 cells, 1 explicit URLs. Opening heading: Agent Workflow Checklist.\par\medskip
\textbf{\path{notebooks/03_dspy_training_lab.ipynb}}. Kind: notebook. Size: 12048 bytes. Shape: 314 lines. Notebook, 11 cells, 0 explicit URLs. Opening heading: DSPy Training Sample Lab.\par\medskip
\textbf{\path{notebooks/04_sample_population_lab.ipynb}}. Kind: notebook. Size: 7894 bytes. Shape: 250 lines. Notebook, 9 cells, 0 explicit URLs. Opening heading: Sample Population Lab.\par\medskip
\textbf{\path{notebooks/05_hushwheel_fixture_rag_lab.ipynb}}. Kind: notebook. Size: 13622 bytes. Shape: 378 lines. Notebook, 9 cells, 0 explicit URLs. Opening heading: Hushwheel Fixture RAG Lab.\par\medskip
\textbf{\path{publication/.gitignore}}. Kind: generic. Size: 80 bytes. Shape: 12 lines. Tracked repository file, 12 lines, 0 explicit URLs. Opening line: .build/\par\medskip
\textbf{\path{publication/Makefile}}. Kind: make. Size: 1644 bytes. Shape: 48 lines. Repository Makefile, 48 lines, 0 explicit URLs. It exposes the shared automation surfaces.\par\medskip
\textbf{\path{publication/README.md}}. Kind: markdown. Size: 1499 bytes. Shape: 34 lines. Markdown document, 34 lines, 0 explicit URLs. Lead heading: Publication Article.\par\medskip
\textbf{\path{publication/article-banner.png}}. Kind: binary. Size: 79074 bytes. Shape: binary surface. Binary asset, 79074 bytes of binary data. It is tracked for stable publication output.\par\medskip
\textbf{\path{publication/exploratorium_translation/.gitignore}}. Kind: generic. Size: 32 bytes. Shape: 6 lines. Tracked repository file, 6 lines, 0 explicit URLs. Opening line: .build/\par\medskip
\textbf{\path{publication/exploratorium_translation/Makefile}}. Kind: make. Size: 870 bytes. Shape: 30 lines. Repository Makefile, 30 lines, 0 explicit URLs. It exposes the shared automation surfaces.\par\medskip
\textbf{\path{publication/exploratorium_translation/README.md}}. Kind: markdown. Size: 997 bytes. Shape: 28 lines. Markdown document, 28 lines, 0 explicit URLs. Lead heading: Exploratorium Translation.\par\medskip
\textbf{\path{publication/exploratorium_translation/exploratorium_translation.tex}}. Kind: latex. Size: 2573 bytes. Shape: 70 lines. LaTeX source, 70 lines, 0 explicit URLs. Sections declared in preview: 0.\par\medskip
\textbf{\path{publication/references.bib}}. Kind: bibliography. Size: 1252 bytes. Shape: 28 lines. Bibliography file, 28 lines, 2 explicit URLs. Entries counted in preview: 4.\par\medskip
\textbf{\path{publication/repository-rag-lab-article.pdf}}. Kind: binary. Size: 236232 bytes. Shape: binary surface. Binary asset, 236232 bytes of binary data. It is tracked for stable publication output.\par\medskip
\textbf{\path{publication/repository-rag-lab-article.tex}}. Kind: latex. Size: 9526 bytes. Shape: 208 lines. LaTeX source, 208 lines, 0 explicit URLs. Sections declared in preview: 4.\par\medskip
\textbf{\path{publication/todo-backlog-table.tex}}. Kind: latex. Size: 10547 bytes. Shape: 34 lines. LaTeX source, 34 lines, 36 explicit URLs. Sections declared in preview: 0.\par\medskip
\textbf{\path{pyproject.toml}}. Kind: toml. Size: 2285 bytes. Shape: 97 lines. TOML configuration file, 97 lines, 0 explicit URLs. Top keys: build-system, project, dependency-groups.\par\medskip
\textbf{\path{rust-cli/Cargo.lock}}. Kind: generic. Size: 3213 bytes. Shape: binary surface. Tracked repository file, 3213 bytes of binary data, 0 explicit URLs. Opening line: No non-empty preview line.\par\medskip
\textbf{\path{rust-cli/Cargo.toml}}. Kind: toml. Size: 144 bytes. Shape: 8 lines. TOML configuration file, 8 lines, 0 explicit URLs. Top keys: package, dependencies.\par\medskip
\textbf{\path{rust-cli/src/main.rs}}. Kind: source. Size: 20422 bytes. Shape: 626 lines. Source file, 626 lines, 0 explicit URLs. Opening line: use rusqlite::\{params, Connection\};\par\medskip
\textbf{\path{samples/README.md}}. Kind: markdown. Size: 713 bytes. Shape: 13 lines. Markdown document, 13 lines, 0 explicit URLs. Lead heading: Samples.\par\medskip
\textbf{\path{samples/logs/20260315T050730Z-gh-runs.md}}. Kind: markdown. Size: 1341 bytes. Shape: 45 lines. Markdown document, 45 lines, 3 explicit URLs. Lead heading: GitHub Actions Run Log.\par\medskip
\textbf{\path{samples/logs/20260315T085140Z-gh-red-status-check.md}}. Kind: markdown. Size: 847 bytes. Shape: 25 lines. Markdown document, 25 lines, 1 explicit URLs. Lead heading: GitHub Actions Red Status Check.\par\medskip
\textbf{\path{samples/logs/20260317T064722Z-full-build.raw.log}}. Kind: generic. Size: 2109 bytes. Shape: 62 lines. Tracked repository file, 62 lines, 0 explicit URLs. Opening line: \# Full Build Raw Log\par\medskip
\textbf{\path{samples/logs/20260317T064833Z-full-build-attempt2.raw.log}}. Kind: generic. Size: 14950 bytes. Shape: 408 lines. Tracked repository file, 408 lines, 0 explicit URLs. Opening line: \# Full Build Raw Log Attempt 2\par\medskip
\textbf{\path{samples/logs/20260317T073401Z-gh-runs-notebook-scaffolding.md}}. Kind: markdown. Size: 2865 bytes. Shape: 35 lines. Markdown document, 35 lines, 13 explicit URLs. Lead heading: GitHub Actions Run Log.\par\medskip
\textbf{\path{samples/logs/20260317T073404Z-gh-runs.md}}. Kind: markdown. Size: 1269 bytes. Shape: 32 lines. Markdown document, 32 lines, 3 explicit URLs. Lead heading: GitHub Actions Run Log.\par\medskip
\textbf{\path{samples/logs/20260317T080018Z-gh-runs-all-notebooks.md}}. Kind: markdown. Size: 2128 bytes. Shape: 34 lines. Markdown document, 34 lines, 8 explicit URLs. Lead heading: GitHub Actions Run Log.\par\medskip
\textbf{\path{samples/logs/20260317T083207Z-gh-runs-todo-backlog.md}}. Kind: markdown. Size: 1063 bytes. Shape: 27 lines. Markdown document, 27 lines, 3 explicit URLs. Lead heading: GitHub Actions Run Log.\par\medskip
\textbf{\path{samples/logs/20260317T084723Z-gh-runs-repo-completeness-checklist.md}}. Kind: markdown. Size: 1186 bytes. Shape: 27 lines. Markdown document, 27 lines, 3 explicit URLs. Lead heading: GitHub Actions Run Log.\par\medskip
\textbf{\path{samples/logs/20260317T085723Z-gh-runs-hushwheel-fixture.md}}. Kind: markdown. Size: 1472 bytes. Shape: 36 lines. Markdown document, 36 lines, 1 explicit URLs. Lead heading: GitHub Actions Run Log.\par\medskip
\textbf{\path{samples/logs/20260317T091641Z-gh-runs-hushwheel-rag-playbook.md}}. Kind: markdown. Size: 1470 bytes. Shape: 36 lines. Markdown document, 36 lines, 1 explicit URLs. Lead heading: GitHub Actions Run Log.\par\medskip
\textbf{\path{samples/logs/20260317T092251Z-gh-runs-publication-article.md}}. Kind: markdown. Size: 1801 bytes. Shape: 33 lines. Markdown document, 33 lines, 5 explicit URLs. Lead heading: GitHub Actions Run Log.\par\medskip
\textbf{\path{samples/logs/20260318T003027Z-gh-runs-wire-skill-triggers-for-do-repo-verification-aud.md}}. Kind: markdown. Size: 1113 bytes. Shape: 29 lines. Markdown document, 29 lines, 1 explicit URLs. Lead heading: GitHub Actions Run Log.\par\medskip
\textbf{\path{samples/logs/20260318T003654Z-gh-runs-execute-all-notebooks.md}}. Kind: markdown. Size: 1682 bytes. Shape: 41 lines. Markdown document, 41 lines, 1 explicit URLs. Lead heading: GitHub Actions Run Log.\par\medskip
\textbf{\path{samples/logs/20260318T003927Z-gh-runs-publication-workflow-fix.md}}. Kind: markdown. Size: 1747 bytes. Shape: 35 lines. Markdown document, 35 lines, 2 explicit URLs. Lead heading: GitHub Actions Run Log.\par\medskip
\textbf{\path{samples/logs/20260318T004712Z-gh-runs-hushwheel-program-harness.md}}. Kind: markdown. Size: 1798 bytes. Shape: 49 lines. Markdown document, 49 lines, 0 explicit URLs. Lead heading: GitHub Actions Run Log.\par\medskip
\textbf{\path{samples/logs/20260318T010254Z-gh-runs-refresh-final-notebook-outputs-for-no-move-those.md}}. Kind: markdown. Size: 1113 bytes. Shape: 29 lines. Markdown document, 29 lines, 1 explicit URLs. Lead heading: GitHub Actions Run Log.\par\medskip
\textbf{\path{samples/logs/20260318T010637Z-gh-runs-env-refresh-retest.md}}. Kind: markdown. Size: 1029 bytes. Shape: 26 lines. Markdown document, 26 lines, 1 explicit URLs. Lead heading: GitHub Actions Run Log.\par\medskip
\textbf{\path{samples/logs/20260318T010741Z-gh-runs-consolidate-session-master.md}}. Kind: markdown. Size: 1950 bytes. Shape: 62 lines. Markdown document, 62 lines, 0 explicit URLs. Lead heading: GitHub Actions Run Log.\par\medskip
\textbf{\path{samples/logs/20260318T012137Z-gh-runs-azure-inference-endpoint-probe.md}}. Kind: markdown. Size: 1004 bytes. Shape: 26 lines. Markdown document, 26 lines, 1 explicit URLs. Lead heading: GitHub Actions Run Log.\par\medskip
\textbf{\path{samples/logs/20260318T012714Z-gh-runs-tighten-repo-hygiene-for-other-agnet-is-doing-1-.md}}. Kind: markdown. Size: 1091 bytes. Shape: 29 lines. Markdown document, 29 lines, 1 explicit URLs. Lead heading: GitHub Actions Run Log.\par\medskip
\textbf{\path{samples/logs/20260318T015109Z-gh-runs-improve-notebook-run-observability-and-reporting.md}}. Kind: markdown. Size: 2151 bytes. Shape: 55 lines. Markdown document, 55 lines, 2 explicit URLs. Lead heading: GitHub Actions Run Log.\par\medskip
\textbf{\path{samples/logs/20260318T021736Z-gh-runs-sync-todo-backlog-into-markdown-and-publication-surfaces.md}}. Kind: markdown. Size: 1778 bytes. Shape: 55 lines. Markdown document, 55 lines, 2 explicit URLs. Lead heading: GitHub Actions Run Log.\par\medskip
\textbf{\path{samples/logs/20260318T023922Z-gh-runs-implement-step-1-dspy-training-path-follow-up.md}}. Kind: markdown. Size: 1735 bytes. Shape: 55 lines. Markdown document, 55 lines, 2 explicit URLs. Lead heading: GitHub Actions Run Log.\par\medskip
\textbf{\path{samples/logs/20260318T044242Z-gh-runs-add-maintained-repository-research-narrative.md}}. Kind: markdown. Size: 2222 bytes. Shape: 56 lines. Markdown document, 56 lines, 2 explicit URLs. Lead heading: GitHub Actions Run Log.\par\medskip
\textbf{\path{samples/logs/20260318T050051Z-gh-runs-add-azure-runtime-surfaces-for-do-1-and-2-and-carry-on-to-next.md}}. Kind: markdown. Size: 1828 bytes. Shape: 38 lines. Markdown document, 38 lines, 2 explicit URLs. Lead heading: GitHub Actions Run Log.\par\medskip
\textbf{\path{samples/logs/20260318T051627Z-gh-runs-develop-retrieval-evaluation-suite.md}}. Kind: markdown. Size: 1768 bytes. Shape: 37 lines. Markdown document, 37 lines, 2 explicit URLs. Lead heading: GitHub Actions Run Log.\par\medskip
\textbf{\path{samples/logs/20260318T060918Z-gh-runs-consolidate-session-master-repair.md}}. Kind: markdown. Size: 2460 bytes. Shape: 70 lines. Markdown document, 70 lines, 3 explicit URLs. Lead heading: GitHub Actions Run Log.\par\medskip
\textbf{\path{samples/logs/20260318T063755Z-gh-runs-keep-utility-sync-tests-side-effect-free.md}}. Kind: markdown. Size: 1074 bytes. Shape: 29 lines. Markdown document, 29 lines, 1 explicit URLs. Lead heading: GitHub Actions Run Log.\par\medskip
\textbf{\path{samples/logs/20260318T064924Z-gh-runs-retrieval-refresh-go-ahead.md}}. Kind: markdown. Size: 1858 bytes. Shape: 54 lines. Markdown document, 54 lines, 2 explicit URLs. Lead heading: GitHub Actions Run Log.\par\medskip
\textbf{\path{samples/logs/20260318T083955Z-gh-runs-benchmark-broadening-2.md}}. Kind: markdown. Size: 1217 bytes. Shape: 34 lines. Markdown document, 34 lines, 1 explicit URLs. Lead heading: GitHub Actions Run Log.\par\medskip
\textbf{\path{samples/logs/20260318T090625Z-gh-runs-hushwheel-quality-instrumentation.md}}. Kind: markdown. Size: 1312 bytes. Shape: 29 lines. Markdown document, 29 lines, 1 explicit URLs. Lead heading: GitHub Actions Run Log.\par\medskip
\textbf{\path{samples/logs/20260318T091252Z-gh-runs-add-rust-lookup-and-retrieval-tag-summaries-for-.md}}. Kind: markdown. Size: 2223 bytes. Shape: 53 lines. Markdown document, 53 lines, 2 explicit URLs. Lead heading: GitHub Actions Run Log.\par\medskip
\textbf{\path{samples/logs/20260318T093632Z-gh-runs-retrieval-regression-gate-great-carry-on.md}}. Kind: markdown. Size: 1325 bytes. Shape: 34 lines. Markdown document, 34 lines, 2 explicit URLs. Lead heading: GitHub Actions Run Log.\par\medskip
\textbf{\path{samples/logs/20260318T094120Z-gh-runs-dspy-live-proof-and-rust-wrapper-repair.md}}. Kind: markdown. Size: 1469 bytes. Shape: 45 lines. Markdown document, 45 lines, 0 explicit URLs. Lead heading: GitHub Actions Summary.\par\medskip
\textbf{\path{samples/logs/20260318T094347Z-gh-runs-hushwheel-hardening-followup.md}}. Kind: markdown. Size: 2739 bytes. Shape: 72 lines. Markdown document, 72 lines, 2 explicit URLs. Lead heading: GitHub Runs Log.\par\medskip
\textbf{\path{samples/logs/20260318T100336Z-gh-runs-sync-final-rebased-inventories-for-carry-on-comm.md}}. Kind: markdown. Size: 1103 bytes. Shape: 29 lines. Markdown document, 29 lines, 1 explicit URLs. Lead heading: GitHub Actions Run Log.\par\medskip
\textbf{\path{samples/logs/20260318T101229Z-gh-runs-hushwheel-audit-refresh.md}}. Kind: markdown. Size: 1707 bytes. Shape: 37 lines. Markdown document, 37 lines, 1 explicit URLs. Lead heading: GitHub Actions Summary.\par\medskip
\textbf{\path{samples/logs/20260318T101939Z-gh-runs-land-each-unlanded.md}}. Kind: markdown. Size: 1556 bytes. Shape: 28 lines. Markdown document, 28 lines, 2 explicit URLs. Lead heading: GitHub Actions Run Log.\par\medskip
\textbf{\path{samples/logs/20260318T123402Z-gh-runs-fix-file-summary-test-churn-for-1.md}}. Kind: markdown. Size: 1061 bytes. Shape: 29 lines. Markdown document, 29 lines, 1 explicit URLs. Lead heading: GitHub Actions Run Log.\par\medskip
\textbf{\path{samples/logs/20260318T124437Z-gh-runs-reconcile-stale-dspy-backlog-for-1.md}}. Kind: markdown. Size: 1060 bytes. Shape: 21 lines. Markdown document, 21 lines, 3 explicit URLs. Lead heading: GitHub Actions Log.\par\medskip
\textbf{\path{samples/logs/20260318T125547Z-gh-runs-wire-lookup-first-ask-for-1.md}}. Kind: markdown. Size: 1058 bytes. Shape: 22 lines. Markdown document, 22 lines, 2 explicit URLs. Lead heading: GitHub Actions Log.\par\medskip
\textbf{\path{samples/logs/20260318T125829Z-gh-runs-restore-retrieval-gate-after-rebase.md}}. Kind: markdown. Size: 1585 bytes. Shape: 36 lines. Markdown document, 36 lines, 0 explicit URLs. Lead heading: GitHub Actions Run Log.\par\medskip
\textbf{\path{samples/logs/20260318T130525Z-gh-runs-finish-node24-publication-cache-followup-for-1.md}}. Kind: markdown. Size: 1119 bytes. Shape: 27 lines. Markdown document, 27 lines, 3 explicit URLs. Lead heading: GitHub Actions Status For `Finish Node 24 publication cache follow-up for "1"`.\par\medskip
\textbf{\path{samples/logs/20260318T131709Z-gh-runs-restore-root-readme-weighting-for-1.md}}. Kind: markdown. Size: 1680 bytes. Shape: 45 lines. Markdown document, 45 lines, 2 explicit URLs. Lead heading: GitHub Actions Run Log.\par\medskip
\textbf{\path{samples/logs/20260318T133247Z-gh-runs-improve-retrieval-doc-priority-for-1.md}}. Kind: markdown. Size: 1335 bytes. Shape: 34 lines. Markdown document, 34 lines, 1 explicit URLs. Lead heading: GitHub Actions Run Log.\par\medskip
\textbf{\path{samples/logs/20260318T144102Z-gh-runs-merge-into-master.md}}. Kind: markdown. Size: 2285 bytes. Shape: 69 lines. Markdown document, 69 lines, 2 explicit URLs. Lead heading: GitHub Actions Run Log.\par\medskip
\textbf{\path{samples/logs/20260318T145935Z-gh-runs-consolidate-all-unconsoldiated-changes.md}}. Kind: markdown. Size: 2401 bytes. Shape: 71 lines. Markdown document, 71 lines, 2 explicit URLs. Lead heading: GitHub Actions Run Log.\par\medskip
\textbf{\path{samples/logs/20260319T035301Z-gh-runs-public-pages-and-workflow-repairs.md}}. Kind: markdown. Size: 1578 bytes. Shape: 37 lines. Markdown document, 37 lines, 1 explicit URLs. Lead heading: GitHub Actions Log.\par\medskip
\textbf{\path{samples/logs/20260319T071544Z-gh-runs-doncsoldiate-consolidation.md}}. Kind: markdown. Size: 1137 bytes. Shape: 27 lines. Markdown document, 27 lines, 0 explicit URLs. Lead heading: GitHub Actions Log.\par\medskip
\textbf{\path{samples/logs/20260319T114913Z-gh-runs-fix-publication-verification-for-whats-next-step.md}}. Kind: markdown. Size: 2936 bytes. Shape: 67 lines. Markdown document, 67 lines, 3 explicit URLs. Lead heading: GitHub Actions Run Log.\par\medskip
\textbf{\path{samples/logs/20260429T172137Z-gh-runs-develop-no-workflow.md}}. Kind: markdown. Size: 866 bytes. Shape: 23 lines. Markdown document, 23 lines, 0 explicit URLs. Lead heading: GitHub Actions Run Log.\par\medskip
\textbf{\path{samples/logs/20260429T200446Z-gh-runs-develop-no-new-run-after-blob-queue-migration.md}}. Kind: markdown. Size: 1843 bytes. Shape: 34 lines. Markdown document, 34 lines, 0 explicit URLs. Lead heading: GitHub Actions Run Log.\par\medskip
\textbf{\path{samples/logs/20260430T122222Z-gh-runs-develop-no-new-run-after-codex-proxy-push.md}}. Kind: markdown. Size: 908 bytes. Shape: 25 lines. Markdown document, 25 lines, 0 explicit URLs. Lead heading: GitHub Actions check after pushing `bbdcf4b`.\par\medskip
\textbf{\path{samples/logs/20260430T172518Z-gh-runs-develop-no-new-run-after-trace-handoff-push.md}}. Kind: markdown. Size: 991 bytes. Shape: 27 lines. Markdown document, 27 lines, 0 explicit URLs. Lead heading: GitHub Actions check after pushing `839f4cb`.\par\medskip
\textbf{\path{samples/logs/20260501T103130Z-gh-runs-develop-no-new-run-after-azure-runtime-env-fix-push.md}}. Kind: markdown. Size: 1411 bytes. Shape: 35 lines. Markdown document, 35 lines, 0 explicit URLs. Lead heading: GitHub Actions check after pushing `9e627e1`.\par\medskip
\textbf{\path{samples/logs/20260501T112330Z-gh-runs-develop-no-new-run-after-azure-responses-api-version-fix-push.md}}. Kind: markdown. Size: 1522 bytes. Shape: 37 lines. Markdown document, 37 lines, 0 explicit URLs. Lead heading: GitHub Actions check after pushing `128874b`.\par\medskip
\textbf{\path{samples/logs/20260501T125058Z-gh-runs-develop-no-new-run-after-trusted-trace-handoff-push.md}}. Kind: markdown. Size: 1355 bytes. Shape: 30 lines. Markdown document, 30 lines, 0 explicit URLs. Lead heading: GitHub Actions check after pushing `66a7ad4`.\par\medskip
\textbf{\path{samples/logs/20260501T162125Z-gh-runs-develop-no-new-run-after-live-trainer-bundle-publication-push.md}}. Kind: markdown. Size: 1486 bytes. Shape: 30 lines. Markdown document, 30 lines, 0 explicit URLs. Lead heading: GitHub Run Inspection.\par\medskip
\textbf{\path{samples/logs/20260501T174147Z-gh-runs-develop-no-new-run-after-versioned-trainer-recovery-push.md}}. Kind: markdown. Size: 1505 bytes. Shape: 30 lines. Markdown document, 30 lines, 0 explicit URLs. Lead heading: GitHub Run Inspection.\par\medskip
\textbf{\path{samples/logs/20260501T201242Z-gh-runs-develop-no-new-run-after-dspy-bundle-version-pinning-push.md}}. Kind: markdown. Size: 2076 bytes. Shape: 41 lines. Markdown document, 41 lines, 0 explicit URLs. Lead heading: GitHub Run Inspection.\par\medskip
\textbf{\path{samples/logs/20260502T064702Z-gh-runs-develop-no-new-run-after-live-trainer-timestamp-bundle-audit-push.md}}. Kind: markdown. Size: 1604 bytes. Shape: 38 lines. Markdown document, 38 lines, 0 explicit URLs. Lead heading: GitHub Run Inspection.\par\medskip
\textbf{\path{samples/logs/20260502T082402Z-gh-runs-develop-no-new-run-after-hybrid-vector-retrieval-push.md}}. Kind: markdown. Size: 1631 bytes. Shape: 30 lines. Markdown document, 30 lines, 0 explicit URLs. Lead heading: GitHub Run Inspection.\par\medskip
\textbf{\path{samples/logs/20260502T121755Z-gh-runs-develop-no-new-run-after-trainer-gating-and-retrieval-isolation-push.md}}. Kind: markdown. Size: 1175 bytes. Shape: 25 lines. Markdown document, 25 lines, 0 explicit URLs. Lead heading: Summary.\par\medskip
\textbf{\path{samples/logs/20260502T165750Z-gh-runs-develop-no-new-run-after-trainer-champion-batching-push.md}}. Kind: markdown. Size: 1543 bytes. Shape: 35 lines. Markdown document, 35 lines, 0 explicit URLs. Lead heading: GitHub Run Inspection.\par\medskip
\textbf{\path{samples/logs/20260503T080503Z-gh-runs-develop-no-new-run-after-codex-resume-lane-telemetry-push.md}}. Kind: markdown. Size: 1782 bytes. Shape: 37 lines. Markdown document, 37 lines, 0 explicit URLs. Lead heading: GitHub Run Inspection.\par\medskip
\textbf{\path{samples/logs/20260503T153706Z-gh-runs-develop-no-new-run-after-staged-bundle-mirror-push.md}}. Kind: markdown. Size: 1171 bytes. Shape: 27 lines. Markdown document, 27 lines, 0 explicit URLs. Lead heading: GitHub Runs After `Support staged bundle mirrors and stale queue skips`.\par\medskip
\textbf{\path{samples/logs/20260503T160232Z-gh-runs-develop-no-new-run-after-no-op-trainer-promotion-fix-push.md}}. Kind: markdown. Size: 1482 bytes. Shape: 32 lines. Markdown document, 32 lines, 0 explicit URLs. Lead heading: GitHub Run Inspection.\par\medskip
\textbf{\path{samples/logs/20260504T090704Z-gh-runs-develop-no-new-run-after-fix-codex-resume-pvc-session-path-push.md}}. Kind: markdown. Size: 1689 bytes. Shape: 46 lines. Markdown document, 46 lines, 0 explicit URLs. Lead heading: GitHub Run Check After `Fix Codex resume PVC session path`.\par\medskip
\textbf{\path{samples/logs/20260504T111427Z-gh-runs-develop-no-new-run-after-codex-resume-restore-fallback-debug-push.md}}. Kind: markdown. Size: 3370 bytes. Shape: 72 lines. Markdown document, 72 lines, 0 explicit URLs. Lead heading: GitHub Run Check After `Document Codex resume restore fallback debug`.\par\medskip
\textbf{\path{samples/logs/20260504T141246Z-gh-runs-develop-no-new-run-after-codex-resume-guard-hardening-push.md}}. Kind: markdown. Size: 2837 bytes. Shape: 67 lines. Markdown document, 67 lines, 0 explicit URLs. Lead heading: GitHub Run Check After `Harden Codex resume restore guards`.\par\medskip
\textbf{\path{samples/logs/20260504T161106Z-gh-runs-develop-no-new-run-after-explicit-codex-resume-targeting-push.md}}. Kind: markdown. Size: 1175 bytes. Shape: 27 lines. Markdown document, 27 lines, 0 explicit URLs. Lead heading: GitHub Actions inspection after `464e16b` push.\par\medskip
\textbf{\path{samples/logs/20260504T180441Z-gh-runs-develop-no-new-run-after-codex-session-continuity-markers-push.md}}. Kind: markdown. Size: 986 bytes. Shape: 28 lines. Markdown document, 28 lines, 0 explicit URLs. Lead heading: GitHub Actions inspection after `490fe41` push.\par\medskip
\textbf{\path{samples/logs/20260505T072441Z-gh-runs-develop-no-new-run-after-mcp-first-codex-discovery-push.md}}. Kind: markdown. Size: 1152 bytes. Shape: 28 lines. Markdown document, 28 lines, 0 explicit URLs. Lead heading: GitHub Actions inspection after `abcf52c` push.\par\medskip
\textbf{\path{samples/logs/README.md}}. Kind: markdown. Size: 331 bytes. Shape: 11 lines. Markdown document, 11 lines, 0 explicit URLs. Lead heading: GitHub Run Logs.\par\medskip
\textbf{\path{samples/population/hushwheel_fixture_population_candidates.yaml}}. Kind: yaml. Size: 955 bytes. Shape: 22 lines. YAML data or workflow file, 22 lines, 0 explicit URLs. Top keys: no top-level mapping keys.\par\medskip
\textbf{\path{samples/population/repository_population_candidates.yaml}}. Kind: yaml. Size: 690 bytes. Shape: 16 lines. YAML data or workflow file, 16 lines, 0 explicit URLs. Top keys: no top-level mapping keys.\par\medskip
\textbf{\path{samples/training/hushwheel_fixture_training_examples.yaml}}. Kind: yaml. Size: 1852 bytes. Shape: 68 lines. YAML data or workflow file, 68 lines, 0 explicit URLs. Top keys: no top-level mapping keys.\par\medskip
\textbf{\path{samples/training/repository_training_examples.yaml}}. Kind: yaml. Size: 2504 bytes. Shape: 80 lines. YAML data or workflow file, 80 lines, 0 explicit URLs. Top keys: no top-level mapping keys.\par\medskip
\textbf{\path{src/repo_rag_lab/__init__.py}}. Kind: python. Size: 1087 bytes. Shape: 47 lines. Python module, 47 lines, 0 explicit URLs. Module intent: Public package surface for the repository-grounded RAG lab.\par\medskip
\textbf{\path{src/repo_rag_lab/azure.py}}. Kind: python. Size: 2592 bytes. Shape: 76 lines. Python module, 76 lines, 0 explicit URLs. Module intent: Azure deployment manifest helpers for downstream release workflows.\par\medskip
\textbf{\path{src/repo_rag_lab/azure_artifacts.py}}. Kind: python. Size: 16790 bytes. Shape: 436 lines. Python module, 436 lines, 0 explicit URLs. Module intent: Azure Blob + Queue helpers for trainer-side traces and global bundle storage.\par\medskip
\textbf{\path{src/repo_rag_lab/azure_runtime.py}}. Kind: python. Size: 20123 bytes. Shape: 577 lines. Python module, 577 lines, 0 explicit URLs. Module intent: Azure runtime normalization, live probes, and cloud-completion helpers.\par\medskip
\textbf{\path{src/repo_rag_lab/benchmarks.py}}. Kind: python. Size: 16524 bytes. Shape: 484 lines. Python module, 484 lines, 0 explicit URLs. Module intent: Retrieval benchmark helpers for notebook assertions and logging.\par\medskip
\textbf{\path{src/repo_rag_lab/cli.py}}. Kind: python. Size: 58870 bytes. Shape: 1339 lines. Python module, 1339 lines, 0 explicit URLs. Module intent: Command-line entrypoints for the shared repository RAG workflows.\par\medskip
\textbf{\path{src/repo_rag_lab/codex_proxy.py}}. Kind: python. Size: 75460 bytes. Shape: 1946 lines. Python module, 1946 lines, 0 explicit URLs. Module intent: Local streaming proxy that injects repo-RAG + DSPy mediation for Codex.\par\medskip
\textbf{\path{src/repo_rag_lab/corpus.py}}. Kind: python. Size: 5245 bytes. Shape: 181 lines. Python module, 181 lines, 0 explicit URLs. Module intent: Repository file loading helpers for the baseline text corpus.\par\medskip
\textbf{\path{src/repo_rag_lab/dspy_training.py}}. Kind: python. Size: 61618 bytes. Shape: 1612 lines. Python module, 1612 lines, 0 explicit URLs. Module intent: DSPy training and artifact helpers for repository-grounded RAG.\par\medskip
\textbf{\path{src/repo_rag_lab/dspy_workflow.py}}. Kind: python. Size: 6518 bytes. Shape: 180 lines. Python module, 180 lines, 0 explicit URLs. Module intent: DSPy-backed runtime wrappers around the repository retrieval workflow.\par\medskip
\textbf{\path{src/repo_rag_lab/exploratorium_translation.py}}. Kind: python. Size: 38049 bytes. Shape: 1009 lines. Python module, 1009 lines, 0 explicit URLs. Module intent: Generate the bilingual exploratorium translation publication surfaces.\par\medskip
\textbf{\path{src/repo_rag_lab/file_summaries.py}}. Kind: python. Size: 10564 bytes. Shape: 327 lines. Python module, 327 lines, 0 explicit URLs. Module intent: Generate tracked repository file-summary surfaces.\par\medskip
\textbf{\path{src/repo_rag_lab/github_pr_gates.py}}. Kind: python. Size: 3979 bytes. Shape: 140 lines. Python module, 140 lines, 0 explicit URLs. Module intent: GitHub branch-protection helpers for required pull-request gate checks.\par\medskip
\textbf{\path{src/repo_rag_lab/mcp.py}}. Kind: python. Size: 2947 bytes. Shape: 84 lines. Python module, 84 lines, 0 explicit URLs. Module intent: Discovery helpers for MCP-related repository artifacts.\par\medskip
\textbf{\path{src/repo_rag_lab/mcp_server.py}}. Kind: python. Size: 50399 bytes. Shape: 1361 lines. Python module, 1361 lines, 0 explicit URLs. Module intent: Minimal stdio MCP server for bounded repo-RAG operations.\par\medskip
\textbf{\path{src/repo_rag_lab/mcp_stdio.py}}. Kind: python. Size: 953 bytes. Shape: 38 lines. Python module, 38 lines, 0 explicit URLs. Module intent: Lightweight stdio entrypoint for the bounded repo-RAG MCP server.\par\medskip
\textbf{\path{src/repo_rag_lab/notebook_runner.py}}. Kind: python. Size: 16433 bytes. Shape: 491 lines. Python module, 491 lines, 0 explicit URLs. Module intent: Monitored batch execution helpers for repository notebooks.\par\medskip
\textbf{\path{src/repo_rag_lab/notebook_scaffolding.py}}. Kind: python. Size: 10965 bytes. Shape: 258 lines. Python module, 258 lines, 0 explicit URLs. Module intent: High-level notebook scaffolds built from tested repository helpers.\par\medskip
\textbf{\path{src/repo_rag_lab/notebook_support.py}}. Kind: python. Size: 2843 bytes. Shape: 89 lines. Python module, 89 lines, 0 explicit URLs. Module intent: Helpers shared by notebooks so their setup logic stays in tested Python code.\par\medskip
\textbf{\path{src/repo_rag_lab/pages_site.py}}. Kind: python. Size: 13635 bytes. Shape: 390 lines. Python module, 390 lines, 2 explicit URLs. Module intent: Generate a MkDocs-ready catalog of tracked repository Markdown files.\par\medskip
\textbf{\path{src/repo_rag_lab/population_samples.py}}. Kind: python. Size: 8376 bytes. Shape: 241 lines. Python module, 241 lines, 0 explicit URLs. Module intent: Helpers for reviewing and batching starter corpus-population inputs.\par\medskip
\textbf{\path{src/repo_rag_lab/retrieval.py}}. Kind: python. Size: 25533 bytes. Shape: 788 lines. Python module, 788 lines, 0 explicit URLs. Module intent: Baseline chunking and lexical retrieval utilities for repository text.\par\medskip
\textbf{\path{src/repo_rag_lab/retrieval_profile.py}}. Kind: python. Size: 9527 bytes. Shape: 236 lines. Python module, 236 lines, 0 explicit URLs. Module intent: Profile-driven retrieval weighting, exclusions, and repo-local overrides.\par\medskip
\textbf{\path{src/repo_rag_lab/runtime_artifacts.py}}. Kind: python. Size: 118360 bytes. Shape: 2671 lines. Python module, 2671 lines, 0 explicit URLs. Module intent: Versioned bundle, overlay, and runtime-trace helpers.\par\medskip
\textbf{\path{src/repo_rag_lab/rust_lookup.py}}. Kind: python. Size: 7724 bytes. Shape: 295 lines. Python module, 295 lines, 0 explicit URLs. Module intent: Helpers for the native Rust-backed SQLite lookup index.\par\medskip
\textbf{\path{src/repo_rag_lab/semantic_retrieval.py}}. Kind: python. Size: 7407 bytes. Shape: 214 lines. Python module, 214 lines, 0 explicit URLs. Module intent: Azure OpenAI embedding-backed semantic retrieval helpers.\par\medskip
\textbf{\path{src/repo_rag_lab/settings.py}}. Kind: python. Size: 699 bytes. Shape: 28 lines. Python module, 28 lines, 0 explicit URLs. Module intent: Shared repository path settings derived from a repository root.\par\medskip
\textbf{\path{src/repo_rag_lab/todo_backlog.py}}. Kind: python. Size: 8434 bytes. Shape: 276 lines. Python module, 276 lines, 1 explicit URLs. Module intent: Shared backlog rendering helpers for Markdown and LaTeX surfaces.\par\medskip
\textbf{\path{src/repo_rag_lab/trainer_deployment.py}}. Kind: python. Size: 18027 bytes. Shape: 431 lines. Python module, 431 lines, 0 explicit URLs. Module intent: Kubernetes manifest helpers for trainer-side repo-RAG deployment surfaces.\par\medskip
\textbf{\path{src/repo_rag_lab/training_samples.py}}. Kind: python. Size: 110198 bytes. Shape: 2709 lines. Python module, 2709 lines, 0 explicit URLs. Module intent: Helpers for loading and summarizing starter DSPy training examples.\par\medskip
\textbf{\path{src/repo_rag_lab/utilities.py}}. Kind: python. Size: 116367 bytes. Shape: 2809 lines. Python module, 2809 lines, 1 explicit URLs. Module intent: User-facing utility helpers shared by the CLI, tests, and notebooks.\par\medskip
\textbf{\path{src/repo_rag_lab/verification.py}}. Kind: python. Size: 4476 bytes. Shape: 163 lines. Python module, 163 lines, 0 explicit URLs. Module intent: Verification helpers for Makefile and notebook contract checks.\par\medskip
\textbf{\path{src/repo_rag_lab/workflow.py}}. Kind: python. Size: 15900 bytes. Shape: 503 lines. Python module, 503 lines, 0 explicit URLs. Module intent: Baseline repository question-answering workflow built on file retrieval.\par\medskip
\textbf{\path{tests/features/repository_rag.feature}}. Kind: generic. Size: 377 bytes. Shape: 12 lines. Tracked repository file, 12 lines, 0 explicit URLs. Opening line: Feature: Repository RAG\par\medskip
\textbf{\path{tests/fixtures/hushwheel_lexiconarium/.gitignore}}. Kind: generic. Size: 23 bytes. Shape: 4 lines. Tracked repository file, 4 lines, 0 explicit URLs. Opening line: hushwheel\par\medskip
\textbf{\path{tests/fixtures/hushwheel_lexiconarium/CHANGELOG.md}}. Kind: markdown. Size: 345 bytes. Shape: 9 lines. Markdown document, 9 lines, 0 explicit URLs. Lead heading: Changelog.\par\medskip
\textbf{\path{tests/fixtures/hushwheel_lexiconarium/Doxyfile}}. Kind: generic. Size: 1239 bytes. Shape: 35 lines. Tracked repository file, 35 lines, 0 explicit URLs. Opening line: PROJECT\_NAME           = Hushwheel Lexiconarium\par\medskip
\textbf{\path{tests/fixtures/hushwheel_lexiconarium/LICENSE.md}}. Kind: markdown. Size: 1082 bytes. Shape: 22 lines. Markdown document, 22 lines, 0 explicit URLs. Lead heading: LICENSE.\par\medskip
\textbf{\path{tests/fixtures/hushwheel_lexiconarium/Makefile}}. Kind: make. Size: 20908 bytes. Shape: 388 lines. Repository Makefile, 388 lines, 0 explicit URLs. It exposes the shared automation surfaces.\par\medskip
\textbf{\path{tests/fixtures/hushwheel_lexiconarium/README.md}}. Kind: markdown. Size: 8269 bytes. Shape: 162 lines. Markdown document, 162 lines, 0 explicit URLs. Lead heading: The Hushwheel Lexiconarium.\par\medskip
\textbf{\path{tests/fixtures/hushwheel_lexiconarium/VERSION}}. Kind: generic. Size: 6 bytes. Shape: 2 lines. Tracked repository file, 2 lines, 0 explicit URLs. Opening line: 0.1.0\par\medskip
\textbf{\path{tests/fixtures/hushwheel_lexiconarium/docs/architecture.md}}. Kind: markdown. Size: 1404 bytes. Shape: 33 lines. Markdown document, 33 lines, 0 explicit URLs. Lead heading: Hushwheel Architecture.\par\medskip
\textbf{\path{tests/fixtures/hushwheel_lexiconarium/docs/catalog.md}}. Kind: markdown. Size: 27404 bytes. Shape: 357 lines. Markdown document, 357 lines, 0 explicit URLs. Lead heading: Hushwheel Representative Catalog.\par\medskip
\textbf{\path{tests/fixtures/hushwheel_lexiconarium/docs/concepts.md}}. Kind: markdown. Size: 1559 bytes. Shape: 37 lines. Markdown document, 37 lines, 0 explicit URLs. Lead heading: Hushwheel Concepts.\par\medskip
\textbf{\path{tests/fixtures/hushwheel_lexiconarium/docs/constellation-atlas.md}}. Kind: markdown. Size: 5111 bytes. Shape: 142 lines. Markdown document, 142 lines, 0 explicit URLs. Lead heading: Hushwheel Constellation Atlas §::<[]>.\par\medskip
\textbf{\path{tests/fixtures/hushwheel_lexiconarium/docs/districts.md}}. Kind: markdown. Size: 1151 bytes. Shape: 42 lines. Markdown document, 42 lines, 0 explicit URLs. Lead heading: Hushwheel Districts.\par\medskip
\textbf{\path{tests/fixtures/hushwheel_lexiconarium/docs/doxygen-fonts.sty}}. Kind: generic. Size: 212 bytes. Shape: binary surface. Tracked repository file, 212 bytes of binary data, 0 explicit URLs. Opening line: No non-empty preview line.\par\medskip
\textbf{\path{tests/fixtures/hushwheel_lexiconarium/docs/hushwheel-reference.pdf}}. Kind: binary. Size: 132 bytes. Shape: binary surface. Binary asset, 132 bytes of binary data. It is tracked for stable publication output.\par\medskip
\textbf{\path{tests/fixtures/hushwheel_lexiconarium/docs/operations.md}}. Kind: markdown. Size: 1068 bytes. Shape: 35 lines. Markdown document, 35 lines, 0 explicit URLs. Lead heading: Hushwheel Operations.\par\medskip
\textbf{\path{tests/fixtures/hushwheel_lexiconarium/docs/packaging.md}}. Kind: markdown. Size: 1112 bytes. Shape: 33 lines. Markdown document, 33 lines, 0 explicit URLs. Lead heading: Hushwheel Packaging.\par\medskip
\textbf{\path{tests/fixtures/hushwheel_lexiconarium/docs/testing.md}}. Kind: markdown. Size: 3507 bytes. Shape: 80 lines. Markdown document, 80 lines, 0 explicit URLs. Lead heading: Hushwheel Testing.\par\medskip
\textbf{\path{tests/fixtures/hushwheel_lexiconarium/fixture-manifest.json}}. Kind: json. Size: 1098 bytes. Shape: 49 lines. JSON data file, 49 lines, 0 explicit URLs. Top keys: entry\_count, source\_size\_bytes, doc\_files.\par\medskip
\textbf{\path{tests/fixtures/hushwheel_lexiconarium/include/hushwheel.h}}. Kind: source. Size: 404 bytes. Shape: 19 lines. Source file, 19 lines, 0 explicit URLs. Opening line: \#ifndef HUSHWHEEL\_H\par\medskip
\textbf{\path{tests/fixtures/hushwheel_lexiconarium/packaging/hushwheel.1}}. Kind: generic. Size: 1303 bytes. Shape: binary surface. Tracked repository file, 1303 bytes of binary data, 0 explicit URLs. Opening line: No non-empty preview line.\par\medskip
\textbf{\path{tests/fixtures/hushwheel_lexiconarium/packaging/hushwheel.package.json}}. Kind: json. Size: 591 bytes. Shape: 29 lines. JSON data file, 29 lines, 0 explicit URLs. Top keys: name, version, summary.\par\medskip
\textbf{\path{tests/fixtures/hushwheel_lexiconarium/src/hushwheel.c}}. Kind: source. Size: 1694579 bytes. Shape: 4375 lines. Source file, 4375 lines, 0 explicit URLs. Opening line: /**\par\medskip
\textbf{\path{tests/fixtures/hushwheel_lexiconarium/src/hushwheel_internal.h}}. Kind: source. Size: 603 bytes. Shape: 25 lines. Source file, 25 lines, 0 explicit URLs. Opening line: \#ifndef HUSHWHEEL\_INTERNAL\_H\par\medskip
\textbf{\path{tests/fixtures/hushwheel_lexiconarium/src/hushwheel_spoke_argent.c}}. Kind: source. Size: 324091 bytes. Shape: 525 lines. Source file, 525 lines, 0 explicit URLs. Opening line: \#include "hushwheel\_internal.h"\par\medskip
\textbf{\path{tests/fixtures/hushwheel_lexiconarium/src/hushwheel_spoke_bhagavatam.c}}. Kind: source. Size: 326160 bytes. Shape: 525 lines. Source file, 525 lines, 0 explicit URLs. Opening line: \#include "hushwheel\_internal.h"\par\medskip
\textbf{\path{tests/fixtures/hushwheel_lexiconarium/src/hushwheel_spoke_crosscanon.c}}. Kind: source. Size: 326144 bytes. Shape: 525 lines. Source file, 525 lines, 0 explicit URLs. Opening line: \#include "hushwheel\_internal.h"\par\medskip
\textbf{\path{tests/fixtures/hushwheel_lexiconarium/src/hushwheel_spoke_editorial.c}}. Kind: source. Size: 325628 bytes. Shape: 525 lines. Source file, 525 lines, 0 explicit URLs. Opening line: \#include "hushwheel\_internal.h"\par\medskip
\textbf{\path{tests/fixtures/hushwheel_lexiconarium/src/hushwheel_spoke_gita.c}}. Kind: source. Size: 323055 bytes. Shape: 525 lines. Source file, 525 lines, 0 explicit URLs. Opening line: \#include "hushwheel\_internal.h"\par\medskip
\textbf{\path{tests/fixtures/hushwheel_lexiconarium/src/hushwheel_spoke_kernel.c}}. Kind: source. Size: 324068 bytes. Shape: 525 lines. Source file, 525 lines, 0 explicit URLs. Opening line: \#include "hushwheel\_internal.h"\par\medskip
\textbf{\path{tests/fixtures/hushwheel_lexiconarium/src/hushwheel_spoke_psalter.c}}. Kind: source. Size: 324600 bytes. Shape: 525 lines. Source file, 525 lines, 0 explicit URLs. Opening line: \#include "hushwheel\_internal.h"\par\medskip
\textbf{\path{tests/fixtures/hushwheel_lexiconarium/src/hushwheel_spoke_ranger.c}}. Kind: source. Size: 324105 bytes. Shape: 525 lines. Source file, 525 lines, 0 explicit URLs. Opening line: \#include "hushwheel\_internal.h"\par\medskip
\textbf{\path{tests/fixtures/hushwheel_lexiconarium/src/hushwheel_spokes.c}}. Kind: source. Size: 1074 bytes. Shape: 25 lines. Source file, 25 lines, 0 explicit URLs. Opening line: \#include "hushwheel\_internal.h"\par\medskip
\textbf{\path{tests/fixtures/hushwheel_lexiconarium/tests/bdd/hushwheel.feature}}. Kind: generic. Size: 1021 bytes. Shape: 28 lines. Tracked repository file, 28 lines, 0 explicit URLs. Opening line: Feature: Hushwheel CLI\par\medskip
\textbf{\path{tests/fixtures/hushwheel_lexiconarium/tests/bdd/run_bdd.py}}. Kind: python. Size: 3146 bytes. Shape: 94 lines. Python module, 94 lines, 0 explicit URLs. Module intent: No module docstring detected.\par\medskip
\textbf{\path{tests/fixtures/hushwheel_lexiconarium/tests/integration/cli_suite.py}}. Kind: python. Size: 4482 bytes. Shape: 122 lines. Python module, 122 lines, 0 explicit URLs. Module intent: No module docstring detected.\par\medskip
\textbf{\path{tests/fixtures/hushwheel_lexiconarium/tests/unit/test_hushwheel_unit.c}}. Kind: source. Size: 2002 bytes. Shape: 62 lines. Source file, 62 lines, 0 explicit URLs. Opening line: \#include <assert.h>\par\medskip
\textbf{\path{tests/fixtures/hushwheel_lexiconarium/tools/lint_hushwheel.py}}. Kind: python. Size: 4678 bytes. Shape: 143 lines. Python module, 143 lines, 0 explicit URLs. Module intent: No module docstring detected.\par\medskip
\textbf{\path{tests/fixtures/hushwheel_lexiconarium/tools/regenerate_hushwheel_fixture.py}}. Kind: python. Size: 31375 bytes. Shape: 866 lines. Python module, 866 lines, 0 explicit URLs. Module intent: No module docstring detected.\par\medskip
\textbf{\path{tests/test_azure_runtime.py}}. Kind: python. Size: 7540 bytes. Shape: 208 lines. Python module, 208 lines, 12 explicit URLs. Module intent: No module docstring detected.\par\medskip
\textbf{\path{tests/test_benchmarks_and_notebook_scaffolding.py}}. Kind: python. Size: 11859 bytes. Shape: 274 lines. Python module, 274 lines, 0 explicit URLs. Module intent: No module docstring detected.\par\medskip
\textbf{\path{tests/test_cli_and_dspy.py}}. Kind: python. Size: 68622 bytes. Shape: 1876 lines. Python module, 1876 lines, 2 explicit URLs. Module intent: No module docstring detected.\par\medskip
\textbf{\path{tests/test_codex_proxy.py}}. Kind: python. Size: 57154 bytes. Shape: 1571 lines. Python module, 1571 lines, 5 explicit URLs. Module intent: No module docstring detected.\par\medskip
\textbf{\path{tests/test_dspy_training.py}}. Kind: python. Size: 66652 bytes. Shape: 1653 lines. Python module, 1653 lines, 11 explicit URLs. Module intent: No module docstring detected.\par\medskip
\textbf{\path{tests/test_exploratorium_translation.py}}. Kind: python. Size: 5239 bytes. Shape: 138 lines. Python module, 138 lines, 8 explicit URLs. Module intent: No module docstring detected.\par\medskip
\textbf{\path{tests/test_file_summaries.py}}. Kind: python. Size: 9110 bytes. Shape: 230 lines. Python module, 230 lines, 0 explicit URLs. Module intent: No module docstring detected.\par\medskip
\textbf{\path{tests/test_github_pr_gates.py}}. Kind: python. Size: 5678 bytes. Shape: 183 lines. Python module, 183 lines, 1 explicit URLs. Module intent: No module docstring detected.\par\medskip
\textbf{\path{tests/test_hushwheel_fixture.py}}. Kind: python. Size: 4617 bytes. Shape: 101 lines. Python module, 101 lines, 0 explicit URLs. Module intent: No module docstring detected.\par\medskip
\textbf{\path{tests/test_hushwheel_program_surface.py}}. Kind: python. Size: 4626 bytes. Shape: 125 lines. Python module, 125 lines, 0 explicit URLs. Module intent: No module docstring detected.\par\medskip
\textbf{\path{tests/test_live_azure_integration.py}}. Kind: python. Size: 2634 bytes. Shape: 83 lines. Python module, 83 lines, 0 explicit URLs. Module intent: No module docstring detected.\par\medskip
\textbf{\path{tests/test_lookup_first.py}}. Kind: python. Size: 3559 bytes. Shape: 108 lines. Python module, 108 lines, 0 explicit URLs. Module intent: No module docstring detected.\par\medskip
\textbf{\path{tests/test_mcp_branches.py}}. Kind: python. Size: 1077 bytes. Shape: 31 lines. Python module, 31 lines, 0 explicit URLs. Module intent: No module docstring detected.\par\medskip
\textbf{\path{tests/test_mcp_server.py}}. Kind: python. Size: 26557 bytes. Shape: 799 lines. Python module, 799 lines, 0 explicit URLs. Module intent: No module docstring detected.\par\medskip
\textbf{\path{tests/test_mcp_stdio.py}}. Kind: python. Size: 988 bytes. Shape: 42 lines. Python module, 42 lines, 0 explicit URLs. Module intent: No module docstring detected.\par\medskip
\textbf{\path{tests/test_notebook_runner.py}}. Kind: python. Size: 3088 bytes. Shape: 84 lines. Python module, 84 lines, 0 explicit URLs. Module intent: No module docstring detected.\par\medskip
\textbf{\path{tests/test_pages_site.py}}. Kind: python. Size: 9045 bytes. Shape: 225 lines. Python module, 225 lines, 14 explicit URLs. Module intent: No module docstring detected.\par\medskip
\textbf{\path{tests/test_population_samples.py}}. Kind: python. Size: 3144 bytes. Shape: 76 lines. Python module, 76 lines, 0 explicit URLs. Module intent: No module docstring detected.\par\medskip
\textbf{\path{tests/test_project_surfaces.py}}. Kind: python. Size: 21471 bytes. Shape: 484 lines. Python module, 484 lines, 1 explicit URLs. Module intent: No module docstring detected.\par\medskip
\textbf{\path{tests/test_repository_rag_bdd.py}}. Kind: python. Size: 879 bytes. Shape: 30 lines. Python module, 30 lines, 0 explicit URLs. Module intent: No module docstring detected.\par\medskip
\textbf{\path{tests/test_retrieval.py}}. Kind: python. Size: 13362 bytes. Shape: 389 lines. Python module, 389 lines, 0 explicit URLs. Module intent: No module docstring detected.\par\medskip
\textbf{\path{tests/test_runtime_artifacts_azure.py}}. Kind: python. Size: 58883 bytes. Shape: 1554 lines. Python module, 1554 lines, 4 explicit URLs. Module intent: No module docstring detected.\par\medskip
\textbf{\path{tests/test_rust_lookup.py}}. Kind: python. Size: 3099 bytes. Shape: 95 lines. Python module, 95 lines, 0 explicit URLs. Module intent: No module docstring detected.\par\medskip
\textbf{\path{tests/test_todo_backlog.py}}. Kind: python. Size: 3505 bytes. Shape: 103 lines. Python module, 103 lines, 1 explicit URLs. Module intent: No module docstring detected.\par\medskip
\textbf{\path{tests/test_training_samples.py}}. Kind: python. Size: 84870 bytes. Shape: 2152 lines. Python module, 2152 lines, 0 explicit URLs. Module intent: No module docstring detected.\par\medskip
\textbf{\path{tests/test_utilities.py}}. Kind: python. Size: 121479 bytes. Shape: 3117 lines. Python module, 3117 lines, 6 explicit URLs. Module intent: No module docstring detected.\par\medskip
\textbf{\path{tests/test_verification.py}}. Kind: python. Size: 3156 bytes. Shape: 66 lines. Python module, 66 lines, 0 explicit URLs. Module intent: No module docstring detected.\par\medskip
\textbf{\path{tests/test_workflow.py}}. Kind: python. Size: 3456 bytes. Shape: 99 lines. Python module, 99 lines, 0 explicit URLs. Module intent: No module docstring detected.\par\medskip
\textbf{\path{tests/test_workflow_live.py}}. Kind: python. Size: 4292 bytes. Shape: 128 lines. Python module, 128 lines, 0 explicit URLs. Module intent: No module docstring detected.\par\medskip
\textbf{\path{todo-backlog.yaml}}. Kind: yaml. Size: 6116 bytes. Shape: 128 lines. YAML data or workflow file, 128 lines, 1 explicit URLs. Top keys: repository\_web\_base, items.\par\medskip
\textbf{\path{utilities/README.md}}. Kind: markdown. Size: 3495 bytes. Shape: 54 lines. Markdown document, 54 lines, 0 explicit URLs. Lead heading: Repository Utilities.\par\medskip
\textbf{\path{uv.lock}}. Kind: generic. Size: 784153 bytes. Shape: binary surface. Tracked repository file, 784153 bytes of binary data, 0 explicit URLs. Opening line: No non-empty preview line.\par\medskip
\subsection*{English: Summaries Of All Authored URLs}
\textbf{\url{http://127.0.0.1}}. Host: 127.0.0.1. Occurrences: 2. Referenced by URL only. No repository-local mirror is currently tracked. Sources: \path{tests/test_codex_proxy.py}.\par\medskip
\textbf{\url{http://127.0.0.1:40111/openai`}}. Host: 127.0.0.1:40111. Occurrences: 1. Referenced by URL only. No repository-local mirror is currently tracked. Sources: \path{docs/audit/2026-05-02-zz-codex-exec-resume-session-state-stage0.md}.\par\medskip
\textbf{\url{http://127.0.0.1:44973/openai`}}. Host: 127.0.0.1:44973. Occurrences: 1. Referenced by URL only. No repository-local mirror is currently tracked. Sources: \path{docs/audit/2026-05-02-zz-codex-exec-resume-session-state-stage0.md}.\par\medskip
\textbf{\url{https://api.github.com/repos/example/demo/branches/master/protection}}. Host: api.github.com. Occurrences: 1. External GitHub URL referenced from 1 file(s). Local companion path(s): none. Sources: \path{tests/test_github_pr_gates.py}.\par\medskip
\textbf{\url{https://astral.sh/}}. Host: astral.sh. Occurrences: 2. Referenced by URL only. No repository-local mirror is currently tracked. Sources: \path{tests/test_exploratorium_translation.py}\newline \path{tests/test_utilities.py}.\par\medskip
\textbf{\url{https://astral.sh/uv/install.sh}}. Host: astral.sh. Occurrences: 1. Referenced by URL only. No repository-local mirror is currently tracked. Sources: \path{README.md}.\par\medskip
\textbf{\url{https://dspy-helper.openai.azure.com}}. Host: dspy-helper.openai.azure.com. Occurrences: 1. Referenced by URL only. No repository-local mirror is currently tracked. Sources: \path{tests/test_dspy_training.py}.\par\medskip
\textbf{\url{https://dspy-helper.openai.azure.com/}}. Host: dspy-helper.openai.azure.com. Occurrences: 1. Referenced by URL only. No repository-local mirror is currently tracked. Sources: \path{tests/test_dspy_training.py}.\par\medskip
\textbf{\url{https://dspy.ai/tutorials/rag/}}. Host: dspy.ai. Occurrences: 1. Referenced by URL only. No repository-local mirror is currently tracked. Sources: \path{docs/architecture/inspired/dspy-rag-tutorial.md}.\par\medskip
\textbf{\url{https://example.com}}. Host: example.com. Occurrences: 2. Referenced by URL only. No repository-local mirror is currently tracked. Sources: \path{tests/test_pages_site.py}.\par\medskip
\textbf{\url{https://example.com/docs}}. Host: example.com. Occurrences: 1. Referenced by URL only. No repository-local mirror is currently tracked. Sources: \path{tests/test_pages_site.py}.\par\medskip
\textbf{\url{https://example.com/org/repo/blob/main}}. Host: example.com. Occurrences: 1. Referenced by URL only. No repository-local mirror is currently tracked. Sources: \path{tests/test_todo_backlog.py}.\par\medskip
\textbf{\url{https://example.invalid/v1}}. Host: example.invalid. Occurrences: 2. Referenced by URL only. No repository-local mirror is currently tracked. Sources: \path{tests/test_dspy_training.py}.\par\medskip
\textbf{\url{https://example.openai.azure.com}}. Host: example.openai.azure.com. Occurrences: 5. Referenced by URL only. No repository-local mirror is currently tracked. Sources: \path{tests/test_azure_runtime.py}\newline \path{tests/test_dspy_training.py}.\par\medskip
\textbf{\url{https://example.openai.azure.com/}}. Host: example.openai.azure.com. Occurrences: 6. Referenced by URL only. No repository-local mirror is currently tracked. Sources: \path{tests/test_azure_runtime.py}\newline \path{tests/test_dspy_training.py}.\par\medskip
\textbf{\url{https://example.openai.azure.com/openai/deployments/repo-rag-ft}}. Host: example.openai.azure.com. Occurrences: 3. Referenced by URL only. No repository-local mirror is currently tracked. Sources: \path{tests/test_azure_runtime.py}.\par\medskip
\textbf{\url{https://example.openai.azure.com/openai/deployments/repo-rag-ft/}}. Host: example.openai.azure.com. Occurrences: 4. Referenced by URL only. No repository-local mirror is currently tracked. Sources: \path{tests/test_azure_runtime.py}.\par\medskip
\textbf{\url{https://example.openai.azure.com/openai/deployments/repo-rag-ft/chat/completions}}. Host: example.openai.azure.com. Occurrences: 1. Referenced by URL only. No repository-local mirror is currently tracked. Sources: \path{tests/test_dspy_training.py}.\par\medskip
\textbf{\url{https://example.services.ai.azure.com/models}}. Host: example.services.ai.azure.com. Occurrences: 6. Referenced by URL only. No repository-local mirror is currently tracked. Sources: \path{Makefile}\newline \path{docs/architecture/dspy-rag-guide.md}\newline \path{docs/operations/azure-deployment.md}\newline \path{src/repo_rag_lab/utilities.py}\newline \path{tests/test_cli_and_dspy.py}\newline \path{tests/test_project_surfaces.py}.\par\medskip
\textbf{\url{https://from-connection-string}}. Host: from-connection-string. Occurrences: 2. Referenced by URL only. No repository-local mirror is currently tracked. Sources: \path{tests/test_runtime_artifacts_azure.py}.\par\medskip
\textbf{\url{https://github.com/}}. Host: github.com. Occurrences: 2. External GitHub URL referenced from 1 file(s). Local companion path(s): none. Sources: \path{src/repo_rag_lab/pages_site.py}.\par\medskip
\textbf{\url{https://github.com/acme/repo}}. Host: github.com. Occurrences: 3. External GitHub URL referenced from 1 file(s). Local companion path(s): none. Sources: \path{tests/test_codex_proxy.py}.\par\medskip
\textbf{\url{https://github.com/example/demo}}. Host: github.com. Occurrences: 10. External GitHub URL referenced from 3 file(s). Local companion path(s): none. Sources: \path{tests/test_cli_and_dspy.py}\newline \path{tests/test_pages_site.py}\newline \path{tests/test_utilities.py}.\par\medskip
\textbf{\url{https://github.com/example/demo/blob/main/README.md}}. Host: github.com. Occurrences: 1. External GitHub URL referenced from 1 file(s). Local companion path(s): none. Sources: \path{tests/test_pages_site.py}.\par\medskip
\textbf{\url{https://github.com/example/demo/blob/main/assets/logo.png}}. Host: github.com. Occurrences: 1. External GitHub URL referenced from 1 file(s). Local companion path(s): none. Sources: \path{tests/test_pages_site.py}.\par\medskip
\textbf{\url{https://github.com/example/demo/blob/main/src/tool.c}}. Host: github.com. Occurrences: 1. External GitHub URL referenced from 1 file(s). Local companion path(s): none. Sources: \path{tests/test_pages_site.py}.\par\medskip
\textbf{\url{https://github.com/example/demo/tree/main/docs}}. Host: github.com. Occurrences: 1. External GitHub URL referenced from 1 file(s). Local companion path(s): none. Sources: \path{tests/test_pages_site.py}.\par\medskip
\textbf{\url{https://github.com/example/project}}. Host: github.com. Occurrences: 2. External GitHub URL referenced from 2 file(s). Local companion path(s): none. Sources: \path{tests/test_exploratorium_translation.py}\newline \path{tests/test_utilities.py}.\par\medskip
\textbf{\url{https://github.com/example/project/blob/main/README.md}}. Host: github.com. Occurrences: 1. External GitHub URL referenced from 1 file(s). Local companion path(s): none. Sources: \path{tests/test_exploratorium_translation.py}.\par\medskip
\textbf{\url{https://github.com/example/project/blob/main/README.md\#L1-L2}}. Host: github.com. Occurrences: 1. External GitHub URL referenced from 1 file(s). Local companion path(s): none. Sources: \path{tests/test_exploratorium_translation.py}.\par\medskip
\textbf{\url{https://github.com/realagiorganization/dspy_rag_in_repo_docs_and_impl1}}. Host: github.com. Occurrences: 1. External GitHub URL referenced from 1 file(s). Local companion path(s): none. Sources: \path{mkdocs.yml}.\par\medskip
\textbf{\url{https://github.com/realagiorganization/dspy_rag_in_repo_docs_and_impl1/actions/workflows/ci.yml/badge.svg)](https://github.com/realagiorganization/dspy_rag_in_repo_docs_and_impl1/actions/workflows/ci.yml}}. Host: github.com. Occurrences: 2. External GitHub URL referenced from 2 file(s). Local companion path(s): none. Sources: \path{README.md}\newline \path{notebooks/02_agent_workflow_checklist.ipynb}.\par\medskip
\textbf{\url{https://github.com/realagiorganization/dspy_rag_in_repo_docs_and_impl1/actions/workflows/publish.yml/badge.svg)](https://github.com/realagiorganization/dspy_rag_in_repo_docs_and_impl1/actions/workflows/publish.yml}}. Host: github.com. Occurrences: 1. External GitHub URL referenced from 1 file(s). Local companion path(s): none. Sources: \path{README.md}.\par\medskip
\textbf{\url{https://github.com/realagiorganization/dspy_rag_in_repo_docs_and_impl1/blob/master}}. Host: github.com. Occurrences: 3. External GitHub URL referenced from 3 file(s). Local companion path(s): none. Sources: \path{FILES.md}\newline \path{src/repo_rag_lab/todo_backlog.py}\newline \path{todo-backlog.yaml}.\par\medskip
\textbf{\url{https://github.com/realagiorganization/dspy_rag_in_repo_docs_and_impl1/blob/master/.github/workflows/ci.yml}}. Host: github.com. Occurrences: 1. External GitHub URL referenced from 1 file(s). Local companion path(s): none. Sources: \path{publication/todo-backlog-table.tex}.\par\medskip
\textbf{\url{https://github.com/realagiorganization/dspy_rag_in_repo_docs_and_impl1/blob/master/.github/workflows/publish.yml}}. Host: github.com. Occurrences: 1. External GitHub URL referenced from 1 file(s). Local companion path(s): none. Sources: \path{publication/todo-backlog-table.tex}.\par\medskip
\textbf{\url{https://github.com/realagiorganization/dspy_rag_in_repo_docs_and_impl1/blob/master/Makefile\#L142-L143}}. Host: github.com. Occurrences: 1. External GitHub URL referenced from 1 file(s). Local companion path(s): none. Sources: \path{docs/operations/simple-end-to-end-verification-guide.md}.\par\medskip
\textbf{\url{https://github.com/realagiorganization/dspy_rag_in_repo_docs_and_impl1/blob/master/Makefile\#L153-L154}}. Host: github.com. Occurrences: 1. External GitHub URL referenced from 1 file(s). Local companion path(s): none. Sources: \path{docs/operations/simple-end-to-end-verification-guide.md}.\par\medskip
\textbf{\url{https://github.com/realagiorganization/dspy_rag_in_repo_docs_and_impl1/blob/master/Makefile\#L248-L251}}. Host: github.com. Occurrences: 1. External GitHub URL referenced from 1 file(s). Local companion path(s): none. Sources: \path{docs/operations/simple-end-to-end-verification-guide.md}.\par\medskip
\textbf{\url{https://github.com/realagiorganization/dspy_rag_in_repo_docs_and_impl1/blob/master/Makefile\#L65-L66}}. Host: github.com. Occurrences: 1. External GitHub URL referenced from 1 file(s). Local companion path(s): none. Sources: \path{docs/operations/simple-end-to-end-verification-guide.md}.\par\medskip
\textbf{\url{https://github.com/realagiorganization/dspy_rag_in_repo_docs_and_impl1/blob/master/README.md}}. Host: github.com. Occurrences: 2. External GitHub URL referenced from 1 file(s). Local companion path(s): none. Sources: \path{publication/todo-backlog-table.tex}.\par\medskip
\textbf{\url{https://github.com/realagiorganization/dspy_rag_in_repo_docs_and_impl1/blob/master/TODO.MD}}. Host: github.com. Occurrences: 1. External GitHub URL referenced from 1 file(s). Local companion path(s): none. Sources: \path{publication/todo-backlog-table.tex}.\par\medskip
\textbf{\url{https://github.com/realagiorganization/dspy_rag_in_repo_docs_and_impl1/blob/master/docs/architecture/dspy-rag-guide.md}}. Host: github.com. Occurrences: 1. External GitHub URL referenced from 1 file(s). Local companion path(s): none. Sources: \path{publication/todo-backlog-table.tex}.\par\medskip
\textbf{\url{https://github.com/realagiorganization/dspy_rag_in_repo_docs_and_impl1/blob/master/docs/archive/repo-completeness-checklist.md}}. Host: github.com. Occurrences: 1. External GitHub URL referenced from 1 file(s). Local companion path(s): none. Sources: \path{publication/todo-backlog-table.tex}.\par\medskip
\textbf{\url{https://github.com/realagiorganization/dspy_rag_in_repo_docs_and_impl1/blob/master/docs/audit/README.md}}. Host: github.com. Occurrences: 3. External GitHub URL referenced from 1 file(s). Local companion path(s): none. Sources: \path{publication/todo-backlog-table.tex}.\par\medskip
\textbf{\url{https://github.com/realagiorganization/dspy_rag_in_repo_docs_and_impl1/blob/master/docs/operations/azure-deployment.md}}. Host: github.com. Occurrences: 1. External GitHub URL referenced from 1 file(s). Local companion path(s): none. Sources: \path{publication/todo-backlog-table.tex}.\par\medskip
\textbf{\url{https://github.com/realagiorganization/dspy_rag_in_repo_docs_and_impl1/blob/master/docs/planning/dataset-integration-plan.md}}. Host: github.com. Occurrences: 1. External GitHub URL referenced from 1 file(s). Local companion path(s): none. Sources: \path{publication/todo-backlog-table.tex}.\par\medskip
\textbf{\url{https://github.com/realagiorganization/dspy_rag_in_repo_docs_and_impl1/blob/master/docs/planning/repo-hardening-plan.md}}. Host: github.com. Occurrences: 1. External GitHub URL referenced from 1 file(s). Local companion path(s): none. Sources: \path{publication/todo-backlog-table.tex}.\par\medskip
\textbf{\url{https://github.com/realagiorganization/dspy_rag_in_repo_docs_and_impl1/blob/master/notebooks/01_repo_rag_research.ipynb}}. Host: github.com. Occurrences: 1. External GitHub URL referenced from 1 file(s). Local companion path(s): none. Sources: \path{publication/todo-backlog-table.tex}.\par\medskip
\textbf{\url{https://github.com/realagiorganization/dspy_rag_in_repo_docs_and_impl1/blob/master/notebooks/03_dspy_training_lab.ipynb}}. Host: github.com. Occurrences: 1. External GitHub URL referenced from 1 file(s). Local companion path(s): none. Sources: \path{publication/todo-backlog-table.tex}.\par\medskip
\textbf{\url{https://github.com/realagiorganization/dspy_rag_in_repo_docs_and_impl1/blob/master/notebooks/04_sample_population_lab.ipynb}}. Host: github.com. Occurrences: 1. External GitHub URL referenced from 1 file(s). Local companion path(s): none. Sources: \path{publication/todo-backlog-table.tex}.\par\medskip
\textbf{\url{https://github.com/realagiorganization/dspy_rag_in_repo_docs_and_impl1/blob/master/publication/README.md}}. Host: github.com. Occurrences: 1. External GitHub URL referenced from 1 file(s). Local companion path(s): none. Sources: \path{publication/todo-backlog-table.tex}.\par\medskip
\textbf{\url{https://github.com/realagiorganization/dspy_rag_in_repo_docs_and_impl1/blob/master/pyproject.toml}}. Host: github.com. Occurrences: 1. External GitHub URL referenced from 1 file(s). Local companion path(s): none. Sources: \path{publication/todo-backlog-table.tex}.\par\medskip
\textbf{\url{https://github.com/realagiorganization/dspy_rag_in_repo_docs_and_impl1/blob/master/rust-cli/Cargo.lock}}. Host: github.com. Occurrences: 1. External GitHub URL referenced from 1 file(s). Local companion path(s): none. Sources: \path{publication/todo-backlog-table.tex}.\par\medskip
\textbf{\url{https://github.com/realagiorganization/dspy_rag_in_repo_docs_and_impl1/blob/master/rust-cli/Cargo.toml}}. Host: github.com. Occurrences: 1. External GitHub URL referenced from 1 file(s). Local companion path(s): none. Sources: \path{publication/todo-backlog-table.tex}.\par\medskip
\textbf{\url{https://github.com/realagiorganization/dspy_rag_in_repo_docs_and_impl1/blob/master/samples/population/repository_population_candidates.yaml}}. Host: github.com. Occurrences: 1. External GitHub URL referenced from 1 file(s). Local companion path(s): none. Sources: \path{publication/todo-backlog-table.tex}.\par\medskip
\textbf{\url{https://github.com/realagiorganization/dspy_rag_in_repo_docs_and_impl1/blob/master/samples/training/repository_training_examples.yaml}}. Host: github.com. Occurrences: 1. External GitHub URL referenced from 1 file(s). Local companion path(s): none. Sources: \path{publication/todo-backlog-table.tex}.\par\medskip
\textbf{\url{https://github.com/realagiorganization/dspy_rag_in_repo_docs_and_impl1/blob/master/src/repo_rag_lab/benchmarks.py}}. Host: github.com. Occurrences: 2. External GitHub URL referenced from 1 file(s). Local companion path(s): none. Sources: \path{publication/todo-backlog-table.tex}.\par\medskip
\textbf{\url{https://github.com/realagiorganization/dspy_rag_in_repo_docs_and_impl1/blob/master/src/repo_rag_lab/cli.py\#L72-L88}}. Host: github.com. Occurrences: 1. External GitHub URL referenced from 1 file(s). Local companion path(s): none. Sources: \path{docs/operations/simple-end-to-end-verification-guide.md}.\par\medskip
\textbf{\url{https://github.com/realagiorganization/dspy_rag_in_repo_docs_and_impl1/blob/master/src/repo_rag_lab/dspy_workflow.py}}. Host: github.com. Occurrences: 1. External GitHub URL referenced from 1 file(s). Local companion path(s): none. Sources: \path{publication/todo-backlog-table.tex}.\par\medskip
\textbf{\url{https://github.com/realagiorganization/dspy_rag_in_repo_docs_and_impl1/blob/master/src/repo_rag_lab/notebook_scaffolding.py}}. Host: github.com. Occurrences: 1. External GitHub URL referenced from 1 file(s). Local companion path(s): none. Sources: \path{publication/todo-backlog-table.tex}.\par\medskip
\textbf{\url{https://github.com/realagiorganization/dspy_rag_in_repo_docs_and_impl1/blob/master/src/repo_rag_lab/population_samples.py}}. Host: github.com. Occurrences: 1. External GitHub URL referenced from 1 file(s). Local companion path(s): none. Sources: \path{publication/todo-backlog-table.tex}.\par\medskip
\textbf{\url{https://github.com/realagiorganization/dspy_rag_in_repo_docs_and_impl1/blob/master/src/repo_rag_lab/retrieval.py}}. Host: github.com. Occurrences: 1. External GitHub URL referenced from 1 file(s). Local companion path(s): none. Sources: \path{publication/todo-backlog-table.tex}.\par\medskip
\textbf{\url{https://github.com/realagiorganization/dspy_rag_in_repo_docs_and_impl1/blob/master/src/repo_rag_lab/utilities.py\#L108-L124}}. Host: github.com. Occurrences: 2. External GitHub URL referenced from 1 file(s). Local companion path(s): none. Sources: \path{docs/operations/simple-end-to-end-verification-guide.md}.\par\medskip
\textbf{\url{https://github.com/realagiorganization/dspy_rag_in_repo_docs_and_impl1/blob/master/src/repo_rag_lab/utilities.py\#L190-L193}}. Host: github.com. Occurrences: 1. External GitHub URL referenced from 1 file(s). Local companion path(s): none. Sources: \path{docs/operations/simple-end-to-end-verification-guide.md}.\par\medskip
\textbf{\url{https://github.com/realagiorganization/dspy_rag_in_repo_docs_and_impl1/blob/master/src/repo_rag_lab/verification.py\#L11-L51}}. Host: github.com. Occurrences: 1. External GitHub URL referenced from 1 file(s). Local companion path(s): none. Sources: \path{docs/operations/simple-end-to-end-verification-guide.md}.\par\medskip
\textbf{\url{https://github.com/realagiorganization/dspy_rag_in_repo_docs_and_impl1/blob/master/src/repo_rag_lab/verification.py\#L133-L146}}. Host: github.com. Occurrences: 1. External GitHub URL referenced from 1 file(s). Local companion path(s): none. Sources: \path{docs/operations/simple-end-to-end-verification-guide.md}.\par\medskip
\textbf{\url{https://github.com/realagiorganization/dspy_rag_in_repo_docs_and_impl1/blob/master/src/repo_rag_lab/workflow.py}}. Host: github.com. Occurrences: 1. External GitHub URL referenced from 1 file(s). Local companion path(s): none. Sources: \path{publication/todo-backlog-table.tex}.\par\medskip
\textbf{\url{https://github.com/realagiorganization/dspy_rag_in_repo_docs_and_impl1/blob/master/src/repo_rag_lab/workflow.py\#L119-L145}}. Host: github.com. Occurrences: 1. External GitHub URL referenced from 1 file(s). Local companion path(s): none. Sources: \path{docs/operations/simple-end-to-end-verification-guide.md}.\par\medskip
\textbf{\url{https://github.com/realagiorganization/dspy_rag_in_repo_docs_and_impl1/blob/master/src/repo_rag_lab/workflow.py\#L54-L72}}. Host: github.com. Occurrences: 1. External GitHub URL referenced from 1 file(s). Local companion path(s): none. Sources: \path{docs/operations/simple-end-to-end-verification-guide.md}.\par\medskip
\textbf{\url{https://github.com/realagiorganization/dspy_rag_in_repo_docs_and_impl1/blob/master/tests/test_benchmarks_and_notebook_scaffolding.py}}. Host: github.com. Occurrences: 1. External GitHub URL referenced from 1 file(s). Local companion path(s): none. Sources: \path{publication/todo-backlog-table.tex}.\par\medskip
\textbf{\url{https://github.com/realagiorganization/dspy_rag_in_repo_docs_and_impl1/blob/master/tests/test_repository_rag_bdd.py}}. Host: github.com. Occurrences: 1. External GitHub URL referenced from 1 file(s). Local companion path(s): none. Sources: \path{publication/todo-backlog-table.tex}.\par\medskip
\textbf{\url{https://github.com/realagiorganization/dspy_rag_in_repo_docs_and_impl1/blob/master/utilities/README.md}}. Host: github.com. Occurrences: 1. External GitHub URL referenced from 1 file(s). Local companion path(s): none. Sources: \path{publication/todo-backlog-table.tex}.\par\medskip
\textbf{\url{https://github.com/realagiorganization/dspy_rag_in_repo_docs_and_impl1/tree/master/samples/logs/}}. Host: github.com. Occurrences: 3. External GitHub URL referenced from 1 file(s). Local companion path(s): none. Sources: \path{publication/todo-backlog-table.tex}.\par\medskip
\textbf{\url{https://github.com/realagiorganization/dspy_rag_in_repo_docs_and_impl1/tree/master/tests/fixtures/hushwheel_lexiconarium/}}. Host: github.com. Occurrences: 1. External GitHub URL referenced from 1 file(s). Local companion path(s): none. Sources: \path{publication/todo-backlog-table.tex}.\par\medskip
\textbf{\url{https://github.com/realagiorganization/national-debt-relief}}. Host: github.com. Occurrences: 1. External GitHub URL referenced from 1 file(s). Local companion path(s): none. Sources: \path{docs/audit/2026-05-10-aks-run-25632110510-handoff-fixed-runtime-still-heuristic.md}.\par\medskip
\textbf{\url{https://gpt45standard.openai.azure.com/}}. Host: gpt45standard.openai.azure.com. Occurrences: 5. Referenced by URL only. No repository-local mirror is currently tracked. Sources: \path{docs/audit/2026-04-29-live-azure-gpt54-validation.md}.\par\medskip
\textbf{\url{https://gpt45standard.openai.azure.com/`}}. Host: gpt45standard.openai.azure.com. Occurrences: 3. Referenced by URL only. No repository-local mirror is currently tracked. Sources: \path{README.md}\newline \path{docs/audit/2026-04-29-live-azure-gpt54-validation.md}\newline \path{docs/operations/azure-deployment.md}.\par\medskip
\textbf{\url{https://gpt45standard.openai.azure.com/openai/v1}}. Host: gpt45standard.openai.azure.com. Occurrences: 1. Referenced by URL only. No repository-local mirror is currently tracked. Sources: \path{docs/audit/2026-05-01-dataset-azure-responses-api-version-default-fix.md}.\par\medskip
\textbf{\url{https://gpt45standard.openai.azure.com`}}. Host: gpt45standard.openai.azure.com`. Occurrences: 1. Referenced by URL only. No repository-local mirror is currently tracked. Sources: \path{docs/audit/2026-03-18-zzzzzz-dspy-pipeline-and-live-azure-proof.md}.\par\medskip
\textbf{\url{https://img.shields.io/badge/coverage-94.12\%25-brightgreen)](https://github.com/realagiorganization/dspy_rag_in_repo_docs_and_impl1/actions/workflows/ci.yml}}. Host: img.shields.io. Occurrences: 1. Referenced by URL only. No repository-local mirror is currently tracked. Sources: \path{README.md}.\par\medskip
\textbf{\url{https://learn.microsoft.com/azure/ai-foundry/foundry-models/how-to/inference}}. Host: learn.microsoft.com. Occurrences: 1. External Azure documentation URL. No mirrored upstream copy is tracked, but local companion notes exist at docs/operations/azure-deployment.md. Sources: \path{docs/operations/azure-deployment.md}.\par\medskip
\textbf{\url{https://learn.microsoft.com/azure/ai-services/}}. Host: learn.microsoft.com. Occurrences: 4. External Azure documentation URL. No mirrored upstream copy is tracked, but local companion notes exist at docs/operations/azure-deployment.md. Sources: \path{publication/references.bib}\newline \path{tests/test_exploratorium_translation.py}\newline \path{tests/test_utilities.py}.\par\medskip
\textbf{\url{https://learn.microsoft.com/en-us/azure/ai-foundry/openai/how-to/fine-tuning-deploy}}. Host: learn.microsoft.com. Occurrences: 1. External Azure documentation URL. No mirrored upstream copy is tracked, but local companion notes exist at docs/operations/azure-deployment.md. Sources: \path{docs/operations/azure-deployment.md}.\par\medskip
\textbf{\url{https://learn.microsoft.com/en-us/azure/ai-services/openai/how-to/fine-tuning}}. Host: learn.microsoft.com. Occurrences: 1. External Azure documentation URL. No mirrored upstream copy is tracked, but local companion notes exist at docs/operations/azure-deployment.md. Sources: \path{docs/operations/azure-deployment.md}.\par\medskip
\textbf{\url{https://medium.com/@arancibia.juan22/implementing-rag-with-dspy-a-technical-guide-a6ae15f6a455}}. Host: medium.com. Occurrences: 1. Referenced by URL only. No repository-local mirror is currently tracked. Sources: \path{docs/architecture/inspired/implementing-rag-with-dspy-technical-guide.md}.\par\medskip
\textbf{\url{https://modelcontextprotocol.io/}}. Host: modelcontextprotocol.io. Occurrences: 4. External MCP documentation URL. No mirrored upstream copy is tracked, but local companion notes exist at docs/architecture/mcp-discovery.md. Sources: \path{publication/references.bib}\newline \path{tests/test_exploratorium_translation.py}\newline \path{tests/test_utilities.py}.\par\medskip
\textbf{\url{https://realagiorganization.github.io/dspy_rag_in_repo_docs_and_impl1/}}. Host: realagiorganization.github.io. Occurrences: 3. Referenced by URL only. No repository-local mirror is currently tracked. Sources: \path{README.md}\newline \path{docs/audit/2026-03-19-public-pages-and-secret-scan.md}\newline \path{mkdocs.yml}.\par\medskip
\textbf{\url{https://realagiorganization.github.io/dspy_rag_in_repo_docs_and_impl1/`}}. Host: realagiorganization.github.io. Occurrences: 1. Referenced by URL only. No repository-local mirror is currently tracked. Sources: \path{docs/audit/2026-03-19-public-pages-and-secret-scan.md}.\par\medskip
\textbf{\url{https://realagistorage.blob.core.windows.net}}. Host: realagistorage.blob.core.windows.net. Occurrences: 1. Referenced by URL only. No repository-local mirror is currently tracked. Sources: \path{tests/test_runtime_artifacts_azure.py}.\par\medskip
\textbf{\url{https://realagistorage.queue.core.windows.net}}. Host: realagistorage.queue.core.windows.net. Occurrences: 1. Referenced by URL only. No repository-local mirror is currently tracked. Sources: \path{tests/test_runtime_artifacts_azure.py}.\par\medskip
\newpage
\subsection*{Русский: Состояние загрузки статей и документации}
\textbf{lewis2020rag}: Retrieval-Augmented Generation for Knowledge-Intensive NLP Tasks. Тип: inproceedings. Источник есть в библиографии, но локальный PDF или зеркало статьи не отслеживаются. Ближайшие локальные сопроводительные пути: docs/architecture/dspy-rag-guide.md, docs/architecture/inspired/implementing-rag-with-dspy-technical-guide.md. Связанные пути: \path{docs/architecture/dspy-rag-guide.md, docs/architecture/inspired/implementing-rag-with-dspy-technical-guide.md}.\par\medskip
\textbf{khattab2024dspy}: DSPy: Compiling Declarative Language Model Calls into Self-Improving Pipelines. Тип: article. Источник есть в библиографии, но локальный PDF или зеркало статьи не отслеживаются. Ближайшие локальные сопроводительные пути: docs/architecture/dspy-rag-guide.md, docs/architecture/inspired/dspy-rag-tutorial.md, docs/guides/hushwheel-fixture-rag-guide.md. Связанные пути: \path{docs/architecture/dspy-rag-guide.md, docs/architecture/inspired/dspy-rag-tutorial.md, docs/guides/hushwheel-fixture-rag-guide.md}.\par\medskip
\textbf{mcp2024}: Model Context Protocol. Тип: misc. Внешняя документация цитируется по URL. Зеркальная копия upstream не отслеживается, но есть локальные сопроводительные пути: docs/architecture/mcp-discovery.md. Связанные пути: \path{docs/architecture/mcp-discovery.md}.\par\medskip
\textbf{azureinference2025}: Azure AI Inference Documentation. Тип: misc. Внешняя документация цитируется по URL. Зеркальная копия upstream не отслеживается, но есть локальные сопроводительные пути: docs/operations/azure-deployment.md. Связанные пути: \path{docs/operations/azure-deployment.md}.\par\medskip
\subsection*{Русский: Сводки по всем файлам}
\textbf{\path{.codex/skills/exploratorium-translation-sync/SKILL.md}}. Тип: markdown. Размер: 1722 байт. Форма: 32 строк. Документ Markdown, 32 строк, явных URL: 0. Главный заголовок: Exploratorium Translation Sync.\par\medskip
\textbf{\path{.codex/skills/file-summary-sync/SKILL.md}}. Тип: markdown. Размер: 1367 байт. Форма: 29 строк. Документ Markdown, 29 строк, явных URL: 0. Главный заголовок: File Summary Sync.\par\medskip
\textbf{\path{.codex/skills/notebook-playbook-sync/SKILL.md}}. Тип: markdown. Размер: 2085 байт. Форма: 48 строк. Документ Markdown, 48 строк, явных URL: 0. Главный заголовок: Notebook Playbook Sync.\par\medskip
\textbf{\path{.codex/skills/notebook-playbook-sync/agents/openai.yaml}}. Тип: yaml. Размер: 300 байт. Форма: 5 строк. Файл YAML с данными или workflow, 5 строк, явных URL: 0. Верхние ключи: interface.\par\medskip
\textbf{\path{.codex/skills/post-push-gh-run-logging/SKILL.md}}. Тип: markdown. Размер: 2059 байт. Форма: 41 строк. Документ Markdown, 41 строк, явных URL: 0. Главный заголовок: Post Push Gh Run Logging.\par\medskip
\textbf{\path{.codex/skills/post-push-gh-run-logging/agents/openai.yaml}}. Тип: yaml. Размер: 281 байт. Форма: 5 строк. Файл YAML с данными или workflow, 5 строк, явных URL: 0. Верхние ключи: interface.\par\medskip
\textbf{\path{.codex/skills/post-push-gh-run-logging/scripts/write_gh_run_log.py}}. Тип: python. Размер: 6346 байт. Форма: 178 строк. Модуль Python, 178 строк, явных URL: 0. Назначение модуля: No module docstring detected.\par\medskip
\textbf{\path{.codex/skills/repo-verification-audit-loop/SKILL.md}}. Тип: markdown. Размер: 2589 байт. Форма: 53 строк. Документ Markdown, 53 строк, явных URL: 0. Главный заголовок: Repo Verification Audit Loop.\par\medskip
\textbf{\path{.codex/skills/repo-verification-audit-loop/agents/openai.yaml}}. Тип: yaml. Размер: 272 байт. Форма: 5 строк. Файл YAML с данными или workflow, 5 строк, явных URL: 0. Верхние ключи: interface.\par\medskip
\textbf{\path{.codex/skills/rust-sqlite-lookup/SKILL.md}}. Тип: markdown. Размер: 1297 байт. Форма: 36 строк. Документ Markdown, 36 строк, явных URL: 0. Главный заголовок: Rust SQLite Lookup.\par\medskip
\textbf{\path{.codex/skills/todo-backlog-sync/SKILL.md}}. Тип: markdown. Размер: 1289 байт. Форма: 29 строк. Документ Markdown, 29 строк, явных URL: 0. Главный заголовок: Todo Backlog Sync.\par\medskip
\textbf{\path{.dockerignore}}. Тип: generic. Размер: 150 байт. Форма: 17 строк. Отслеживаемый файл репозитория, 17 строк, явных URL: 0. Первая строка: .git\par\medskip
\textbf{\path{.env.sample}}. Тип: generic. Размер: 905 байт. Форма: 17 строк. Отслеживаемый файл репозитория, 17 строк, явных URL: 0. Первая строка: \# Copy this file to .env and replace placeholder values.\par\medskip
\textbf{\path{.gitattributes}}. Тип: generic. Размер: 126 байт. Форма: 4 строк. Отслеживаемый файл репозитория, 4 строк, явных URL: 0. Первая строка: *.png filter=lfs diff=lfs merge=lfs -text\par\medskip
\textbf{\path{.github/workflows/ci.yml}}. Тип: yaml. Размер: 4185 байт. Форма: 146 строк. Файл YAML с данными или workflow, 146 строк, явных URL: 0. Верхние ключи: name, True, jobs.\par\medskip
\textbf{\path{.github/workflows/hushwheel-quality.yml}}. Тип: yaml. Размер: 3526 байт. Форма: 104 строк. Файл YAML с данными или workflow, 104 строк, явных URL: 0. Верхние ключи: name, True, env.\par\medskip
\textbf{\path{.github/workflows/pages.yml}}. Тип: yaml. Размер: 2424 байт. Форма: 99 строк. Файл YAML с данными или workflow, 99 строк, явных URL: 0. Верхние ключи: name, True, env.\par\medskip
\textbf{\path{.github/workflows/publication-pdf.yml}}. Тип: yaml. Размер: 5549 байт. Форма: 144 строк. Файл YAML с данными или workflow, 144 строк, явных URL: 0. Верхние ключи: name, True, env.\par\medskip
\textbf{\path{.github/workflows/publish.yml}}. Тип: yaml. Размер: 847 байт. Форма: 43 строк. Файл YAML с данными или workflow, 43 строк, явных URL: 0. Верхние ключи: name, True, jobs.\par\medskip
\textbf{\path{.gitignore}}. Тип: generic. Размер: 222 байт. Форма: 21 строк. Отслеживаемый файл репозитория, 21 строк, явных URL: 0. Первая строка: .venv/\par\medskip
\textbf{\path{.pre-commit-config.yaml}}. Тип: yaml. Размер: 1689 байт. Форма: 53 строк. Файл YAML с данными или workflow, 53 строк, явных URL: 0. Верхние ключи: repos.\par\medskip
\textbf{\path{.python-version}}. Тип: generic. Размер: 5 байт. Форма: 2 строк. Отслеживаемый файл репозитория, 2 строк, явных URL: 0. Первая строка: 3.11\par\medskip
\textbf{\path{.yamllint.yml}}. Тип: yaml. Размер: 416 байт. Форма: 20 строк. Файл YAML с данными или workflow, 20 строк, явных URL: 0. Верхние ключи: extends, rules.\par\medskip
\textbf{\path{AGENTS.md}}. Тип: markdown. Размер: 6999 байт. Форма: 108 строк. Документ Markdown, 108 строк, явных URL: 0. Главный заголовок: Repository Agent Instructions.\par\medskip
\textbf{\path{AGENTS.md.d/00-audit-baseline.md}}. Тип: markdown. Размер: 592 байт. Форма: 15 строк. Документ Markdown, 15 строк, явных URL: 0. Главный заголовок: Audit Baseline Instructions.\par\medskip
\textbf{\path{AGENTS.md.d/10-verification-loop.md}}. Тип: markdown. Размер: 1175 байт. Форма: 22 строк. Документ Markdown, 22 строк, явных URL: 0. Главный заголовок: Verification Loop Instructions.\par\medskip
\textbf{\path{AGENTS.md.d/FILES.md}}. Тип: markdown. Размер: 2247 байт. Форма: 41 строк. Документ Markdown, 41 строк, явных URL: 0. Главный заголовок: File Summary Instructions.\par\medskip
\textbf{\path{AGENTS.md.d/RUST_LOOKUP.md}}. Тип: markdown. Размер: 1803 байт. Форма: 36 строк. Документ Markdown, 36 строк, явных URL: 0. Главный заголовок: Rust SQLite Lookup Instructions.\par\medskip
\textbf{\path{Dockerfile}}. Тип: generic. Размер: 1080 байт. Форма: 47 строк. Отслеживаемый файл репозитория, 47 строк, явных URL: 0. Первая строка: ARG PYTHON\_BASE\_IMAGE=python:3.11-slim-bookworm\par\medskip
\textbf{\path{FILES.csv}}. Тип: generic. Размер: 52899 байт. Форма: бинарная поверхность. Отслеживаемый файл репозитория, 52899 байт бинарных данных, явных URL: 0. Первая строка: No non-empty preview line.\par\medskip
\textbf{\path{FILES.md}}. Тип: markdown. Размер: 59726 байт. Форма: 417 строк. Документ Markdown, 417 строк, явных URL: 1. Главный заголовок: Repository File Inventory.\par\medskip
\textbf{\path{Makefile}}. Тип: make. Размер: 23701 байт. Форма: 506 строк. Makefile репозитория, 506 строк, явных URL: 1. Он открывает общие поверхности автоматизации.\par\medskip
\textbf{\path{README.md}}. Тип: markdown. Размер: 26222 байт. Форма: 312 строк. Документ Markdown, 312 строк, явных URL: 6. Главный заголовок: Repository RAG Lab.\par\medskip
\textbf{\path{TODO.MD}}. Тип: markdown. Размер: 6104 байт. Форма: 22 строк. Документ Markdown, 22 строк, явных URL: 0. Главный заголовок: TODO.\par\medskip
\textbf{\path{config/retrieval-profile.json}}. Тип: json. Размер: 1791 байт. Форма: 66 строк. Файл JSON, 66 строк, явных URL: 0. Верхние ключи: name, retrieval\_mode, rerank\_candidate\_pool\_size.\par\medskip
\textbf{\path{data/questions/repository.yaml}}. Тип: yaml. Размер: 283 байт. Форма: 7 строк. Файл YAML с данными или workflow, 7 строк, явных URL: 0. Верхние ключи: no top-level mapping keys.\par\medskip
\textbf{\path{docs/architecture/dspy-rag-guide.md}}. Тип: markdown. Размер: 50272 байт. Форма: 911 строк. Документ Markdown, 911 строк, явных URL: 1. Главный заголовок: DSPy Guide For This Repository.\par\medskip
\textbf{\path{docs/architecture/inspired/dspy-rag-tutorial.md}}. Тип: markdown. Размер: 1474 байт. Форма: 35 строк. Документ Markdown, 35 строк, явных URL: 1. Главный заголовок: DSPy Tutorial: Basic RAG.\par\medskip
\textbf{\path{docs/architecture/inspired/implementing-rag-with-dspy-technical-guide.md}}. Тип: markdown. Размер: 1555 байт. Форма: 38 строк. Документ Markdown, 38 строк, явных URL: 1. Главный заголовок: Implementing RAG with DSPy: Technical Guide.\par\medskip
\textbf{\path{docs/architecture/mcp-discovery.md}}. Тип: markdown. Размер: 1757 байт. Форма: 54 строк. Документ Markdown, 54 строк, явных URL: 0. Главный заголовок: MCP Discovery Notes.\par\medskip
\textbf{\path{docs/architecture/package-api.md}}. Тип: markdown. Размер: 9186 байт. Форма: 127 строк. Документ Markdown, 127 строк, явных URL: 0. Главный заголовок: Package API Notes.\par\medskip
\textbf{\path{docs/architecture/research-narrative.md}}. Тип: markdown. Размер: 82061 байт. Форма: 1077 строк. Документ Markdown, 1077 строк, явных URL: 0. Главный заголовок: Repository Research Narrative.\par\medskip
\textbf{\path{docs/archive/repo-completeness-checklist.md}}. Тип: markdown. Размер: 7450 байт. Форма: 238 строк. Документ Markdown, 238 строк, явных URL: 0. Главный заголовок: Repository Completeness Checklist.\par\medskip
\textbf{\path{docs/audit/2026-03-15-repo-state.md}}. Тип: markdown. Размер: 5989 байт. Форма: 151 строк. Документ Markdown, 151 строк, явных URL: 0. Главный заголовок: Repository State Audit.\par\medskip
\textbf{\path{docs/audit/2026-03-17-docs-and-uv-refresh.md}}. Тип: markdown. Размер: 3977 байт. Форма: 79 строк. Документ Markdown, 79 строк, явных URL: 0. Главный заголовок: Documentation And UV Workflow Refresh Audit.\par\medskip
\textbf{\path{docs/audit/2026-03-17-full-build.md}}. Тип: markdown. Размер: 6715 байт. Форма: 45 строк. Документ Markdown, 45 строк, явных URL: 0. Главный заголовок: Full Build Audit.\par\medskip
\textbf{\path{docs/audit/2026-03-17-gh-actions-watch-loop.md}}. Тип: markdown. Размер: 3797 байт. Форма: 78 строк. Документ Markdown, 78 строк, явных URL: 0. Главный заголовок: GitHub Actions Watch Loop Audit.\par\medskip
\textbf{\path{docs/audit/2026-03-17-hushwheel-fixture.md}}. Тип: markdown. Размер: 5090 байт. Форма: 80 строк. Документ Markdown, 80 строк, явных URL: 0. Главный заголовок: Hushwheel Fixture Audit.\par\medskip
\textbf{\path{docs/audit/2026-03-17-hushwheel-rag-playbook.md}}. Тип: markdown. Размер: 4797 байт. Форма: 83 строк. Документ Markdown, 83 строк, явных URL: 0. Главный заголовок: Hushwheel RAG Playbook Audit.\par\medskip
\textbf{\path{docs/audit/2026-03-17-notebook-scaffolding.md}}. Тип: markdown. Размер: 3641 байт. Форма: 70 строк. Документ Markdown, 70 строк, явных URL: 0. Главный заголовок: Notebook Scaffolding Audit.\par\medskip
\textbf{\path{docs/audit/2026-03-17-publication-article.md}}. Тип: markdown. Размер: 2321 байт. Форма: 54 строк. Документ Markdown, 54 строк, явных URL: 0. Главный заголовок: Publication Article Audit.\par\medskip
\textbf{\path{docs/audit/2026-03-17-repo-completeness-checklist.md}}. Тип: markdown. Размер: 3114 байт. Форма: 59 строк. Документ Markdown, 59 строк, явных URL: 0. Главный заголовок: Repository Completeness Checklist Audit.\par\medskip
\textbf{\path{docs/audit/2026-03-18-actions-caching.md}}. Тип: markdown. Размер: 4189 байт. Форма: 87 строк. Документ Markdown, 87 строк, явных URL: 0. Главный заголовок: GitHub Actions Caching Audit.\par\medskip
\textbf{\path{docs/audit/2026-03-18-azure-inference-endpoint-probe.md}}. Тип: markdown. Размер: 2648 байт. Форма: 63 строк. Документ Markdown, 63 строк, явных URL: 0. Главный заголовок: Azure Inference Endpoint Probe Audit.\par\medskip
\textbf{\path{docs/audit/2026-03-18-azure-runtime-surfaces.md}}. Тип: markdown. Размер: 5307 байт. Форма: 91 строк. Документ Markdown, 91 строк, явных URL: 0. Главный заголовок: Azure Runtime Surfaces Audit.\par\medskip
\textbf{\path{docs/audit/2026-03-18-dspy-training-path.md}}. Тип: markdown. Размер: 4605 байт. Форма: 81 строк. Документ Markdown, 81 строк, явных URL: 0. Главный заголовок: DSPy Training Path Audit.\par\medskip
\textbf{\path{docs/audit/2026-03-18-hushwheel-program-harness.md}}. Тип: markdown. Размер: 3907 байт. Форма: 76 строк. Документ Markdown, 76 строк, явных URL: 0. Главный заголовок: Hushwheel Program Harness Audit.\par\medskip
\textbf{\path{docs/audit/2026-03-18-notebook-execution.md}}. Тип: markdown. Размер: 2577 байт. Форма: 66 строк. Документ Markdown, 66 строк, явных URL: 0. Главный заголовок: Notebook Execution Audit.\par\medskip
\textbf{\path{docs/audit/2026-03-18-notebook-runner-harness.md}}. Тип: markdown. Размер: 4121 байт. Форма: 91 строк. Документ Markdown, 91 строк, явных URL: 0. Главный заголовок: Notebook Runner Harness Audit.\par\medskip
\textbf{\path{docs/audit/2026-03-18-repository-benchmark-broadening.md}}. Тип: markdown. Размер: 3728 байт. Форма: 74 строк. Документ Markdown, 74 строк, явных URL: 0. Главный заголовок: Repository Benchmark Broadening Audit.\par\medskip
\textbf{\path{docs/audit/2026-03-18-retest-with-env-refresh.md}}. Тип: markdown. Размер: 5102 байт. Форма: 105 строк. Документ Markdown, 105 строк, явных URL: 0. Главный заголовок: Env Refresh Retest Audit.\par\medskip
\textbf{\path{docs/audit/2026-03-18-retrieval-evaluation-suite.md}}. Тип: markdown. Размер: 4031 байт. Форма: 80 строк. Документ Markdown, 80 строк, явных URL: 0. Главный заголовок: Retrieval Evaluation Suite Audit.\par\medskip
\textbf{\path{docs/audit/2026-03-18-retrieval-ranking-refresh.md}}. Тип: markdown. Размер: 6479 байт. Форма: 109 строк. Документ Markdown, 109 строк, явных URL: 0. Главный заголовок: Retrieval Ranking Refresh Audit.\par\medskip
\textbf{\path{docs/audit/2026-03-18-rust-and-benchmark-hygiene.md}}. Тип: markdown. Размер: 2049 байт. Форма: 53 строк. Документ Markdown, 53 строк, явных URL: 0. Главный заголовок: Rust And Benchmark Hygiene Audit.\par\medskip
\textbf{\path{docs/audit/2026-03-18-todo-backlog-sync.md}}. Тип: markdown. Размер: 4104 байт. Форма: 83 строк. Документ Markdown, 83 строк, явных URL: 0. Главный заголовок: Todo Backlog Sync Audit.\par\medskip
\textbf{\path{docs/audit/2026-03-18-z-notebook-runner-harness.md}}. Тип: markdown. Размер: 4153 байт. Форма: 104 строк. Документ Markdown, 104 строк, явных URL: 0. Главный заголовок: Notebook Runner Harness Audit.\par\medskip
\textbf{\path{docs/audit/2026-03-18-zz-file-summary-inventory.md}}. Тип: markdown. Размер: 2596 байт. Форма: 69 строк. Документ Markdown, 69 строк, явных URL: 0. Главный заголовок: File Summary Inventory Audit.\par\medskip
\textbf{\path{docs/audit/2026-03-18-zz-research-narrative.md}}. Тип: markdown. Размер: 3312 байт. Форма: 79 строк. Документ Markdown, 79 строк, явных URL: 0. Главный заголовок: Research Narrative Audit.\par\medskip
\textbf{\path{docs/audit/2026-03-18-zzz-session-consolidation-on-origin-master.md}}. Тип: markdown. Размер: 5061 байт. Форма: 109 строк. Документ Markdown, 109 строк, явных URL: 0. Главный заголовок: Session Consolidation On Origin Master.\par\medskip
\textbf{\path{docs/audit/2026-03-18-zzzz-rust-sqlite-lookup.md}}. Тип: markdown. Размер: 5404 байт. Форма: 105 строк. Документ Markdown, 105 строк, явных URL: 0. Главный заголовок: Rust SQLite Lookup And Retrieval Tag Audit.\par\medskip
\textbf{\path{docs/audit/2026-03-18-zzzz-session-consolidation-master-ci-repair.md}}. Тип: markdown. Размер: 3270 байт. Форма: 73 строк. Документ Markdown, 73 строк, явных URL: 0. Главный заголовок: Session Consolidation Master CI Repair.\par\medskip
\textbf{\path{docs/audit/2026-03-18-zzzzz-hushwheel-quality-instrumentation.md}}. Тип: markdown. Размер: 8775 байт. Форма: 165 строк. Документ Markdown, 165 строк, явных URL: 0. Главный заголовок: Hushwheel Quality Instrumentation.\par\medskip
\textbf{\path{docs/audit/2026-03-18-zzzzzz-dspy-pipeline-and-live-azure-proof.md}}. Тип: markdown. Размер: 5992 байт. Форма: 116 строк. Документ Markdown, 116 строк, явных URL: 1. Главный заголовок: DSPy Pipeline And Live Azure Proof.\par\medskip
\textbf{\path{docs/audit/2026-03-18-zzzzzz-hushwheel-retrieval-record-wording.md}}. Тип: markdown. Размер: 3042 байт. Форма: 77 строк. Документ Markdown, 77 строк, явных URL: 0. Главный заголовок: Hushwheel Retrieval Record Wording.\par\medskip
\textbf{\path{docs/audit/2026-03-18-zzzzzzz-hushwheel-pdf-side-effect-fix.md}}. Тип: markdown. Размер: 4044 байт. Форма: 83 строк. Документ Markdown, 83 строк, явных URL: 0. Главный заголовок: Hushwheel PDF Side-Effect Fix.\par\medskip
\textbf{\path{docs/audit/2026-03-18-zzzzzzz-hushwheel-quality-salvage.md}}. Тип: markdown. Размер: 4554 байт. Форма: 84 строк. Документ Markdown, 84 строк, явных URL: 0. Главный заголовок: Hushwheel Quality Salvage Audit.\par\medskip
\textbf{\path{docs/audit/2026-03-18-zzzzzzz-land-unlanded-surfaces.md}}. Тип: markdown. Размер: 2685 байт. Форма: 72 строк. Документ Markdown, 72 строк, явных URL: 0. Главный заголовок: Hushwheel Quality Landing Audit.\par\medskip
\textbf{\path{docs/audit/2026-03-18-zzzzzzz-rebased-master-log-and-azure-typecheck-fix.md}}. Тип: markdown. Размер: 3472 байт. Форма: 83 строк. Документ Markdown, 83 строк, явных URL: 0. Главный заголовок: Rebased Master Log And Azure Typecheck Fix.\par\medskip
\textbf{\path{docs/audit/2026-03-18-zzzzzzz-worktree-consolidation.md}}. Тип: markdown. Размер: 5914 байт. Форма: 125 строк. Документ Markdown, 125 строк, явных URL: 0. Главный заголовок: Worktree Consolidation.\par\medskip
\textbf{\path{docs/audit/2026-03-18-zzzzzzzzzzzz-retrieval-regression-gate.md}}. Тип: markdown. Размер: 2898 байт. Форма: 66 строк. Документ Markdown, 66 строк, явных URL: 0. Главный заголовок: Retrieval Regression Gate Audit.\par\medskip
\textbf{\path{docs/audit/2026-03-18-zzzzzzzzzzzzz-land-remaining-unlanded-work.md}}. Тип: markdown. Размер: 2908 байт. Форма: 78 строк. Документ Markdown, 78 строк, явных URL: 0. Главный заголовок: Land Remaining Unlanded Work.\par\medskip
\textbf{\path{docs/audit/2026-03-18-zzzzzzzzzzzzzz-node24-actions-refresh.md}}. Тип: markdown. Размер: 2989 байт. Форма: 91 строк. Документ Markdown, 91 строк, явных URL: 0. Главный заголовок: Node 24 Actions Refresh.\par\medskip
\textbf{\path{docs/audit/2026-03-18-zzzzzzzzzzzzzz-rebase-hushwheel-quality-over-master.md}}. Тип: markdown. Размер: 1889 байт. Форма: 47 строк. Документ Markdown, 47 строк, явных URL: 0. Главный заголовок: Rebase Hushwheel Quality Over Master.\par\medskip
\textbf{\path{docs/audit/2026-03-18-zzzzzzzzzzzzzzzz-dspy-backlog-reconciliation.md}}. Тип: markdown. Размер: 2956 байт. Форма: 75 строк. Документ Markdown, 75 строк, явных URL: 0. Главный заголовок: DSPy Backlog Reconciliation.\par\medskip
\textbf{\path{docs/audit/2026-03-18-zzzzzzzzzzzzzzzz-file-summary-test-churn-fix.md}}. Тип: markdown. Размер: 2929 байт. Форма: 85 строк. Документ Markdown, 85 строк, явных URL: 0. Главный заголовок: File Summary Test Churn Fix.\par\medskip
\textbf{\path{docs/audit/2026-03-18-zzzzzzzzzzzzzzzz-publication-cache-node24-followup.md}}. Тип: markdown. Размер: 2900 байт. Форма: 89 строк. Документ Markdown, 89 строк, явных URL: 0. Главный заголовок: Publication Cache Node 24 Follow-Up.\par\medskip
\textbf{\path{docs/audit/2026-03-18-zzzzzzzzzzzzzzzz-retrieval-gate-restored-after-rebase.md}}. Тип: markdown. Размер: 3788 байт. Форма: 94 строк. Документ Markdown, 94 строк, явных URL: 0. Главный заголовок: Retrieval Gate Restored After Rebase.\par\medskip
\textbf{\path{docs/audit/2026-03-18-zzzzzzzzzzzzzzzzzz-dspy-backlog-reconciliation.md}}. Тип: markdown. Размер: 2457 байт. Форма: 62 строк. Документ Markdown, 62 строк, явных URL: 0. Главный заголовок: DSPy Backlog Reconciliation.\par\medskip
\textbf{\path{docs/audit/2026-03-18-zzzzzzzzzzzzzzzzzz-lookup-first-ask-path.md}}. Тип: markdown. Размер: 4604 байт. Форма: 94 строк. Документ Markdown, 94 строк, явных URL: 0. Главный заголовок: Lookup First Ask Path.\par\medskip
\textbf{\path{docs/audit/2026-03-18-zzzzzzzzzzzzzzzzzzz-answer-synthesis-refresh.md}}. Тип: markdown. Размер: 2565 байт. Форма: 48 строк. Документ Markdown, 48 строк, явных URL: 0. Главный заголовок: 2026-03-18 Answer Synthesis Refresh.\par\medskip
\textbf{\path{docs/audit/2026-03-18-zzzzzzzzzzzzzzzzzzzz-retrieval-quality-doc-priority.md}}. Тип: markdown. Размер: 3602 байт. Форма: 68 строк. Документ Markdown, 68 строк, явных URL: 0. Главный заголовок: Retrieval Quality Doc Priority Audit.\par\medskip
\textbf{\path{docs/audit/2026-03-18-zzzzzzzzzzzzzzzzzzzzzzz-merge-hushwheel-quality-into-master.md}}. Тип: markdown. Размер: 4335 байт. Форма: 112 строк. Документ Markdown, 112 строк, явных URL: 0. Главный заголовок: Merge Hushwheel Quality Into Master.\par\medskip
\textbf{\path{docs/audit/2026-03-18-zzzzzzzzzzzzzzzzzzzzzzzz-hushwheel-atlas-consolidation.md}}. Тип: markdown. Размер: 4762 байт. Форма: 109 строк. Документ Markdown, 109 строк, явных URL: 0. Главный заголовок: Hushwheel Atlas Consolidation.\par\medskip
\textbf{\path{docs/audit/2026-03-19-github-pr-gate-sync.md}}. Тип: markdown. Размер: 6599 байт. Форма: 136 строк. Документ Markdown, 136 строк, явных URL: 0. Главный заголовок: Master Consolidation And GitHub PR Gate Sync.\par\medskip
\textbf{\path{docs/audit/2026-03-19-node24-actions-opt-in.md}}. Тип: markdown. Размер: 2929 байт. Форма: 80 строк. Документ Markdown, 80 строк, явных URL: 0. Главный заголовок: Node 24 Actions Opt-In.\par\medskip
\textbf{\path{docs/audit/2026-03-19-public-pages-and-secret-scan.md}}. Тип: markdown. Размер: 5804 байт. Форма: 113 строк. Документ Markdown, 113 строк, явных URL: 2. Главный заголовок: Public Pages And Secret Scan.\par\medskip
\textbf{\path{docs/audit/2026-03-19-simple-end-to-end-verification-guide.md}}. Тип: markdown. Размер: 4236 байт. Форма: 103 строк. Документ Markdown, 103 строк, явных URL: 0. Главный заголовок: Simple End-To-End Verification Guide.\par\medskip
\textbf{\path{docs/audit/2026-04-28-cli-envelope-normalization.md}}. Тип: markdown. Размер: 4028 байт. Форма: 95 строк. Документ Markdown, 95 строк, явных URL: 0. Главный заголовок: CLI Envelope Normalization.\par\medskip
\textbf{\path{docs/audit/2026-04-28-doc-tree-restructure-and-planning.md}}. Тип: markdown. Размер: 4309 байт. Форма: 105 строк. Документ Markdown, 105 строк, явных URL: 0. Главный заголовок: Doc Tree Restructure And Planning.\par\medskip
\textbf{\path{docs/audit/2026-04-28-json-runtime-contract.md}}. Тип: markdown. Размер: 5281 байт. Форма: 131 строк. Документ Markdown, 131 строк, явных URL: 0. Главный заголовок: JSON Runtime Contract.\par\medskip
\textbf{\path{docs/audit/2026-04-28-live-azure-dspy-integration-gating.md}}. Тип: markdown. Размер: 3562 байт. Форма: 79 строк. Документ Markdown, 79 строк, явных URL: 0. Главный заголовок: Live Azure And DSPy Integration Gating.\par\medskip
\textbf{\path{docs/audit/2026-04-28-retrieval-profile-layer.md}}. Тип: markdown. Размер: 4123 байт. Форма: 92 строк. Документ Markdown, 92 строк, явных URL: 0. Главный заголовок: Retrieval Profile Layer.\par\medskip
\textbf{\path{docs/audit/2026-04-28-rust-lookup-arbitrary-repo-roots.md}}. Тип: markdown. Размер: 4078 байт. Форма: 88 строк. Документ Markdown, 88 строк, явных URL: 0. Главный заголовок: Rust Lookup Arbitrary Repo Roots.\par\medskip
\textbf{\path{docs/audit/2026-04-29-bounded-mcp-server-surface.md}}. Тип: markdown. Размер: 6688 байт. Форма: 154 строк. Документ Markdown, 154 строк, явных URL: 0. Главный заголовок: 2026-04-29 Bounded MCP Server Surface.\par\medskip
\textbf{\path{docs/audit/2026-04-29-bundle-channel-publish-promote-rollback.md}}. Тип: markdown. Размер: 5633 байт. Форма: 125 строк. Документ Markdown, 125 строк, явных URL: 0. Главный заголовок: 2026-04-29 Bundle Channel Publish Promote Rollback.\par\medskip
\textbf{\path{docs/audit/2026-04-29-bundle-overlay-and-runtime-trace-contract.md}}. Тип: markdown. Размер: 4332 байт. Форма: 96 строк. Документ Markdown, 96 строк, явных URL: 0. Главный заголовок: 2026-04-29 Bundle Overlay And Runtime Trace Contract.\par\medskip
\textbf{\path{docs/audit/2026-04-29-dataset-accepted-outcome-ingest.md}}. Тип: markdown. Размер: 6670 байт. Форма: 139 строк. Документ Markdown, 139 строк, явных URL: 0. Главный заголовок: 2026-04-29 Dataset Accepted-Outcome Ingest.\par\medskip
\textbf{\path{docs/audit/2026-04-29-dataset-bundle-fetch-and-trace-import-hooks.md}}. Тип: markdown. Размер: 6215 байт. Форма: 130 строк. Документ Markdown, 130 строк, явных URL: 0. Главный заголовок: 2026-04-29 Dataset Bundle Fetch And Trace Import Hooks.\par\medskip
\textbf{\path{docs/audit/2026-04-29-dataset-local-repo-rag-cli-backend.md}}. Тип: markdown. Размер: 5290 байт. Форма: 122 строк. Документ Markdown, 122 строк, явных URL: 0. Главный заголовок: 2026-04-29 Dataset Local `repo\_rag\_cli` Backend.\par\medskip
\textbf{\path{docs/audit/2026-04-29-dataset-repo-rag-auto-detection.md}}. Тип: markdown. Размер: 6051 байт. Форма: 124 строк. Документ Markdown, 124 строк, явных URL: 0. Главный заголовок: 2026-04-29 Dataset Repo-RAG Auto-Detection.\par\medskip
\textbf{\path{docs/audit/2026-04-29-dataset-trace-queue-handoff.md}}. Тип: markdown. Размер: 6656 байт. Форма: 141 строк. Документ Markdown, 141 строк, явных URL: 0. Главный заголовок: 2026-04-29 Dataset Trace Queue Handoff.\par\medskip
\textbf{\path{docs/audit/2026-04-29-dataset-worker-repo-rag-cli-backend.md}}. Тип: markdown. Размер: 6630 байт. Форма: 128 строк. Документ Markdown, 128 строк, явных URL: 0. Главный заголовок: 2026-04-29 Dataset Worker `repo\_rag\_cli` Backend.\par\medskip
\textbf{\path{docs/audit/2026-04-29-global-blob-queue-trainer-migration.md}}. Тип: markdown. Размер: 4807 байт. Форма: 91 строк. Документ Markdown, 91 строк, явных URL: 0. Главный заголовок: 2026-04-29 Global Blob And Queue Trainer Migration.\par\medskip
\textbf{\path{docs/audit/2026-04-29-live-azure-gpt54-validation.md}}. Тип: markdown. Размер: 8279 байт. Форма: 145 строк. Документ Markdown, 145 строк, явных URL: 6. Главный заголовок: 2026-04-29 Live Azure GPT-5.4 Validation.\par\medskip
\textbf{\path{docs/audit/2026-04-29-retrieval-rerank-and-relative-root-normalization.md}}. Тип: markdown. Размер: 3070 байт. Форма: 71 строк. Документ Markdown, 71 строк, явных URL: 0. Главный заголовок: 2026-04-29 Retrieval Rerank And Relative Root Normalization.\par\medskip
\textbf{\path{docs/audit/2026-04-29-runtime-image-and-dataset-build-deploy-wiring.md}}. Тип: markdown. Размер: 5729 байт. Форма: 140 строк. Документ Markdown, 140 строк, явных URL: 0. Главный заголовок: 2026-04-29 Runtime Image And Dataset Build/Deploy Wiring.\par\medskip
\textbf{\path{docs/audit/2026-04-29-trainer-bundle-benchmark-gates.md}}. Тип: markdown. Размер: 4942 байт. Форма: 111 строк. Документ Markdown, 111 строк, явных URL: 0. Главный заголовок: 2026-04-29 Trainer Bundle Benchmark Gates.\par\medskip
\textbf{\path{docs/audit/2026-04-29-trainer-cycle-background-pass.md}}. Тип: markdown. Размер: 4637 байт. Форма: 105 строк. Документ Markdown, 105 строк, явных URL: 0. Главный заголовок: 2026-04-29 Trainer Cycle Background Pass.\par\medskip
\textbf{\path{docs/audit/2026-04-29-trainer-k8s-deployment-surface.md}}. Тип: markdown. Размер: 6040 байт. Форма: 131 строк. Документ Markdown, 131 строк, явных URL: 0. Главный заголовок: 2026-04-29 Trainer Kubernetes Deployment Surface.\par\medskip
\textbf{\path{docs/audit/2026-04-29-trainer-recompile-from-candidates.md}}. Тип: markdown. Размер: 6238 байт. Форма: 135 строк. Документ Markdown, 135 строк, явных URL: 0. Главный заголовок: 2026-04-29 Trainer Recompile From Candidates.\par\medskip
\textbf{\path{docs/audit/2026-04-29-trainer-service-loop.md}}. Тип: markdown. Размер: 6185 байт. Форма: 133 строк. Документ Markdown, 133 строк, явных URL: 0. Главный заголовок: 2026-04-29 Trainer Service Loop.\par\medskip
\textbf{\path{docs/audit/2026-04-30-aks-run-25179828865-stale-image-inspection.md}}. Тип: markdown. Размер: 6206 байт. Форма: 117 строк. Документ Markdown, 117 строк, явных URL: 0. Главный заголовок: 2026-04-30 AKS Run 25179828865 Stale Image Inspection.\par\medskip
\textbf{\path{docs/audit/2026-04-30-aks-run-25186264345-image-fix-gap-inspection.md}}. Тип: markdown. Размер: 8407 байт. Форма: 165 строк. Документ Markdown, 165 строк, явных URL: 0. Главный заголовок: 2026-04-30 AKS Run 25186264345 Image/Fix Gap Inspection.\par\medskip
\textbf{\path{docs/audit/2026-04-30-codex-proxy-mediation-layer.md}}. Тип: markdown. Размер: 5279 байт. Форма: 95 строк. Документ Markdown, 95 строк, явных URL: 0. Главный заголовок: 2026-04-30 Codex Proxy Mediation Layer.\par\medskip
\textbf{\path{docs/audit/2026-04-30-codex-proxy-trace-handoff.md}}. Тип: markdown. Размер: 6543 байт. Форма: 122 строк. Документ Markdown, 122 строк, явных URL: 0. Главный заголовок: 2026-04-30 Codex Proxy Trace Handoff.\par\medskip
\textbf{\path{docs/audit/2026-05-01-aks-run-25191504814-project-root-resolution-inspection.md}}. Тип: markdown. Размер: 6002 байт. Форма: 133 строк. Документ Markdown, 133 строк, явных URL: 0. Главный заголовок: 2026-05-01 AKS Run 25191504814 Project Root Resolution Inspection.\par\medskip
\textbf{\path{docs/audit/2026-05-01-aks-run-25209573387-proxy-runtime-env-inspection.md}}. Тип: markdown. Размер: 5738 байт. Форма: 129 строк. Документ Markdown, 129 строк, явных URL: 0. Главный заголовок: 2026-05-01 AKS Run 25209573387 Proxy Runtime Env Inspection.\par\medskip
\textbf{\path{docs/audit/2026-05-01-aks-run-25212955759-trace-handoff-secret-gap.md}}. Тип: markdown. Размер: 4594 байт. Форма: 121 строк. Документ Markdown, 121 строк, явных URL: 0. Главный заголовок: AKS Run `25212955759`: Trace Handoff Skipped Because `repo-rag-storage-config` Lacks Blob Credentials.\par\medskip
\textbf{\path{docs/audit/2026-05-01-aks-run-25226558751-rag-heuristic-trusted-handoff.md}}. Тип: markdown. Размер: 5092 байт. Форма: 130 строк. Документ Markdown, 130 строк, явных URL: 0. Главный заголовок: AKS Run 25226558751: RAG Success, Heuristic DSPy Fallback, Trusted Handoff Success.\par\medskip
\textbf{\path{docs/audit/2026-05-01-dataset-azure-responses-api-version-default-fix.md}}. Тип: markdown. Размер: 6700 байт. Форма: 133 строк. Документ Markdown, 133 строк, явных URL: 1. Главный заголовок: 2026-05-01 Dataset Azure Responses API-Version Default Fix.\par\medskip
\textbf{\path{docs/audit/2026-05-01-dataset-azure-runtime-env-propagation-fix.md}}. Тип: markdown. Размер: 4200 байт. Форма: 81 строк. Документ Markdown, 81 строк, явных URL: 0. Главный заголовок: 2026-05-01 Dataset Azure Runtime Env Propagation Fix.\par\medskip
\textbf{\path{docs/audit/2026-05-01-dataset-trusted-trace-handoff-postprocess-fix.md}}. Тип: markdown. Размер: 3880 байт. Форма: 81 строк. Документ Markdown, 81 строк, явных URL: 0. Главный заголовок: Dataset Trusted Trace Handoff Post-Processing Fix.\par\medskip
\textbf{\path{docs/audit/2026-05-01-dspy-bundle-version-pinning-contract.md}}. Тип: markdown. Размер: 5763 байт. Форма: 118 строк. Документ Markdown, 118 строк, явных URL: 0. Главный заголовок: DSPY Bundle Version Pinning Contract.\par\medskip
\textbf{\path{docs/audit/2026-05-01-durable-trainer-progress-remote-recovery.md}}. Тип: markdown. Размер: 5141 байт. Форма: 115 строк. Документ Markdown, 115 строк, явных URL: 0. Главный заголовок: Durable Trainer Progress Remote Recovery.\par\medskip
\textbf{\path{docs/audit/2026-05-01-live-trainer-no-timestamp-bundle-publish.md}}. Тип: markdown. Размер: 7368 байт. Форма: 172 строк. Документ Markdown, 172 строк, явных URL: 0. Главный заголовок: Live Trainer No Timestamp Bundle Publish.\par\medskip
\textbf{\path{docs/audit/2026-05-01-live-trainer-remote-bundle-publish.md}}. Тип: markdown. Размер: 7444 байт. Форма: 148 строк. Документ Markdown, 148 строк, явных URL: 0. Главный заголовок: Live Trainer Remote Bundle Publish.\par\medskip
\textbf{\path{docs/audit/2026-05-01-live-trainer-service-azure-fallback.md}}. Тип: markdown. Размер: 6390 байт. Форма: 167 строк. Документ Markdown, 167 строк, явных URL: 0. Главный заголовок: Live Trainer Service Azure Fallback.\par\medskip
\textbf{\path{docs/audit/2026-05-01-versioned-trainer-lineage-bundles.md}}. Тип: markdown. Размер: 5821 байт. Форма: 123 строк. Документ Markdown, 123 строк, явных URL: 0. Главный заголовок: Versioned Trainer Lineage Bundles.\par\medskip
\textbf{\path{docs/audit/2026-05-02-hybrid-vector-retrieval-default.md}}. Тип: markdown. Размер: 18302 байт. Форма: 385 строк. Документ Markdown, 385 строк, явных URL: 0. Главный заголовок: Hybrid Vector Retrieval Default.\par\medskip
\textbf{\path{docs/audit/2026-05-02-live-trainer-still-not-publishing-bundles.md}}. Тип: markdown. Размер: 60556 байт. Форма: 1281 строк. Документ Markdown, 1281 строк, явных URL: 0. Главный заголовок: Live Trainer Still Not Publishing Bundles.\par\medskip
\textbf{\path{docs/audit/2026-05-02-trainer-context-group-champion-index-stage1.md}}. Тип: markdown. Размер: 18201 байт. Форма: 330 строк. Документ Markdown, 330 строк, явных URL: 0. Главный заголовок: Trainer Context-Group Champion Index Stage 1.\par\medskip
\textbf{\path{docs/audit/2026-05-02-zz-codex-exec-resume-session-state-stage0.md}}. Тип: markdown. Размер: 170338 байт. Форма: 3866 строк. Документ Markdown, 3866 строк, явных URL: 2. Главный заголовок: Repository audit note for 2026-05-02 codex exec resume session state stage 0.\par\medskip
\textbf{\path{docs/audit/2026-05-08-zz-per-turn-dspy-mediation-contract-stage0.md}}. Тип: markdown. Размер: 5821 байт. Форма: 113 строк. Документ Markdown, 113 строк, явных URL: 0. Главный заголовок: Repository audit note for 2026-05-08 per-turn DSPy mediation contract stage 0.\par\medskip
\textbf{\path{docs/audit/2026-05-09-aks-run-25598675829-per-turn-mediation-gap.md}}. Тип: markdown. Размер: 7119 байт. Форма: 215 строк. Документ Markdown, 215 строк, явных URL: 0. Главный заголовок: Repository audit note for 2026-05-09 AKS run 25598675829 per-turn mediation gap.\par\medskip
\textbf{\path{docs/audit/2026-05-09-family-first-mipro-contract-stage1.md}}. Тип: markdown. Размер: 4684 байт. Форма: 120 строк. Документ Markdown, 120 строк, явных URL: 0. Главный заголовок: Repository audit note for 2026-05-09 family-first MIPROv2 contract stage 1.\par\medskip
\textbf{\path{docs/audit/2026-05-09-family-first-mipro-contract-stage10.md}}. Тип: markdown. Размер: 5329 байт. Форма: 122 строк. Документ Markdown, 122 строк, явных URL: 0. Главный заголовок: Repository audit note for 2026-05-09 family-first MIPROv2 contract stage 10.\par\medskip
\textbf{\path{docs/audit/2026-05-09-family-first-mipro-contract-stage11.md}}. Тип: markdown. Размер: 4638 байт. Форма: 116 строк. Документ Markdown, 116 строк, явных URL: 0. Главный заголовок: Repository audit note for 2026-05-09 family-first MIPROv2 contract stage 11.\par\medskip
\textbf{\path{docs/audit/2026-05-09-family-first-mipro-contract-stage12.md}}. Тип: markdown. Размер: 4941 байт. Форма: 122 строк. Документ Markdown, 122 строк, явных URL: 0. Главный заголовок: Repository audit note for 2026-05-09 family-first MIPROv2 contract stage 12.\par\medskip
\textbf{\path{docs/audit/2026-05-09-family-first-mipro-contract-stage13.md}}. Тип: markdown. Размер: 4992 байт. Форма: 128 строк. Документ Markdown, 128 строк, явных URL: 0. Главный заголовок: Repository audit note for 2026-05-09 family-first MIPROv2 contract stage 13.\par\medskip
\textbf{\path{docs/audit/2026-05-09-family-first-mipro-contract-stage14.md}}. Тип: markdown. Размер: 3944 байт. Форма: 109 строк. Документ Markdown, 109 строк, явных URL: 0. Главный заголовок: Repository audit note for 2026-05-09 family-first MIPROv2 contract stage 14.\par\medskip
\textbf{\path{docs/audit/2026-05-09-family-first-mipro-contract-stage15.md}}. Тип: markdown. Размер: 4728 байт. Форма: 113 строк. Документ Markdown, 113 строк, явных URL: 0. Главный заголовок: Repository audit note for 2026-05-09 family-first MIPROv2 contract stage 15.\par\medskip
\textbf{\path{docs/audit/2026-05-09-family-first-mipro-contract-stage16.md}}. Тип: markdown. Размер: 3748 байт. Форма: 103 строк. Документ Markdown, 103 строк, явных URL: 0. Главный заголовок: Repository audit note for 2026-05-09 family-first MIPROv2 contract stage 16.\par\medskip
\textbf{\path{docs/audit/2026-05-09-family-first-mipro-contract-stage17.md}}. Тип: markdown. Размер: 3519 байт. Форма: 101 строк. Документ Markdown, 101 строк, явных URL: 0. Главный заголовок: Repository audit note for 2026-05-09 family-first MIPROv2 contract stage 17.\par\medskip
\textbf{\path{docs/audit/2026-05-09-family-first-mipro-contract-stage18.md}}. Тип: markdown. Размер: 3164 байт. Форма: 90 строк. Документ Markdown, 90 строк, явных URL: 0. Главный заголовок: Repository audit note for 2026-05-09 family-first MIPROv2 contract stage 18.\par\medskip
\textbf{\path{docs/audit/2026-05-09-family-first-mipro-contract-stage19.md}}. Тип: markdown. Размер: 3984 байт. Форма: 102 строк. Документ Markdown, 102 строк, явных URL: 0. Главный заголовок: Repository audit note for 2026-05-09 family-first MIPROv2 contract stage 19.\par\medskip
\textbf{\path{docs/audit/2026-05-09-family-first-mipro-contract-stage2.md}}. Тип: markdown. Размер: 4936 байт. Форма: 123 строк. Документ Markdown, 123 строк, явных URL: 0. Главный заголовок: Repository audit note for 2026-05-09 family-first MIPROv2 contract stage 2.\par\medskip
\textbf{\path{docs/audit/2026-05-09-family-first-mipro-contract-stage20.md}}. Тип: markdown. Размер: 3076 байт. Форма: 88 строк. Документ Markdown, 88 строк, явных URL: 0. Главный заголовок: Repository audit note for 2026-05-09 family-first MIPROv2 contract stage 20.\par\medskip
\textbf{\path{docs/audit/2026-05-09-family-first-mipro-contract-stage21.md}}. Тип: markdown. Размер: 5742 байт. Форма: 122 строк. Документ Markdown, 122 строк, явных URL: 0. Главный заголовок: Repository audit note for 2026-05-09 family-first MIPROv2 contract stage 21.\par\medskip
\textbf{\path{docs/audit/2026-05-09-family-first-mipro-contract-stage3.md}}. Тип: markdown. Размер: 4147 байт. Форма: 96 строк. Документ Markdown, 96 строк, явных URL: 0. Главный заголовок: Repository audit note for 2026-05-09 family-first MIPROv2 contract stage 3.\par\medskip
\textbf{\path{docs/audit/2026-05-09-family-first-mipro-contract-stage4.md}}. Тип: markdown. Размер: 5209 байт. Форма: 120 строк. Документ Markdown, 120 строк, явных URL: 0. Главный заголовок: Repository audit note for 2026-05-09 family-first MIPROv2 contract stage 4.\par\medskip
\textbf{\path{docs/audit/2026-05-09-family-first-mipro-contract-stage5.md}}. Тип: markdown. Размер: 5172 байт. Форма: 122 строк. Документ Markdown, 122 строк, явных URL: 0. Главный заголовок: Repository audit note for 2026-05-09 family-first MIPROv2 contract stage 5.\par\medskip
\textbf{\path{docs/audit/2026-05-09-family-first-mipro-contract-stage6.md}}. Тип: markdown. Размер: 5120 байт. Форма: 119 строк. Документ Markdown, 119 строк, явных URL: 0. Главный заголовок: Repository audit note for 2026-05-09 family-first MIPROv2 contract stage 6.\par\medskip
\textbf{\path{docs/audit/2026-05-09-family-first-mipro-contract-stage7.md}}. Тип: markdown. Размер: 4680 байт. Форма: 116 строк. Документ Markdown, 116 строк, явных URL: 0. Главный заголовок: Repository audit note for 2026-05-09 family-first MIPROv2 contract stage 7.\par\medskip
\textbf{\path{docs/audit/2026-05-09-family-first-mipro-contract-stage8.md}}. Тип: markdown. Размер: 4477 байт. Форма: 110 строк. Документ Markdown, 110 строк, явных URL: 0. Главный заголовок: Repository audit note for 2026-05-09 family-first MIPROv2 contract stage 8.\par\medskip
\textbf{\path{docs/audit/2026-05-09-family-first-mipro-contract-stage9.md}}. Тип: markdown. Размер: 4003 байт. Форма: 106 строк. Документ Markdown, 106 строк, явных URL: 0. Главный заголовок: Repository audit note for 2026-05-09 family-first MIPROv2 contract stage 9.\par\medskip
\textbf{\path{docs/audit/2026-05-10-aks-run-25625638264-family-first-runtime-gap.md}}. Тип: markdown. Размер: 8472 байт. Форма: 246 строк. Документ Markdown, 246 строк, явных URL: 0. Главный заголовок: Repository audit note for 2026-05-10 AKS run 25625638264 family-first runtime gap.\par\medskip
\textbf{\path{docs/audit/2026-05-10-aks-run-25629990035-family-first-runtime-still-heuristic.md}}. Тип: markdown. Размер: 6843 байт. Форма: 211 строк. Документ Markdown, 211 строк, явных URL: 0. Главный заголовок: Repository audit note for 2026-05-10 AKS run 25629990035 family-first runtime still heuristic.\par\medskip
\textbf{\path{docs/audit/2026-05-10-aks-run-25632110510-handoff-fixed-runtime-still-heuristic.md}}. Тип: markdown. Размер: 7832 байт. Форма: 231 строк. Документ Markdown, 231 строк, явных URL: 1. Главный заголовок: Repository audit note for 2026-05-10 AKS run 25632110510 handoff fixed but runtime still heuristic.\par\medskip
\textbf{\path{docs/audit/2026-05-10-aks-run-25634202380-reset-lane-batch-handoff-runtime-still-heuristic.md}}. Тип: markdown. Размер: 6096 байт. Форма: 197 строк. Документ Markdown, 197 строк, явных URL: 0. Главный заголовок: Repository audit note for AKS run 25634202380 on 2026-05-10.\par\medskip
\textbf{\path{docs/audit/2026-05-10-family-first-runtime-fixes-after-run-25625638264.md}}. Тип: markdown. Размер: 12143 байт. Форма: 261 строк. Документ Markdown, 261 строк, явных URL: 0. Главный заголовок: Repository audit note for 2026-05-10 family-first runtime fixes after run 25625638264.\par\medskip
\textbf{\path{docs/audit/2026-05-10-family-first-runtime-fixes-after-run-25629990035.md}}. Тип: markdown. Размер: 4170 байт. Форма: 98 строк. Документ Markdown, 98 строк, явных URL: 0. Главный заголовок: Repository audit note for 2026-05-10 family-first runtime fixes after AKS run 25629990035.\par\medskip
\textbf{\path{docs/audit/2026-05-10-family-first-runtime-fixes-after-run-25632110510.md}}. Тип: markdown. Размер: 5332 байт. Форма: 112 строк. Документ Markdown, 112 строк, явных URL: 0. Главный заголовок: Repository audit note for 2026-05-10 family-first runtime fixes after AKS run 25632110510.\par\medskip
\textbf{\path{docs/audit/2026-05-10-family-state-populates-but-publish-gate-and-dirty-fathers-still-block-runtime.md}}. Тип: markdown. Размер: 7106 байт. Форма: 183 строк. Документ Markdown, 183 строк, явных URL: 0. Главный заголовок: 2026-05-10 Family State Populates But Publish Gate And Dirty Fathers Still Block Runtime.\par\medskip
\textbf{\path{docs/audit/2026-05-10-family-state-root-alias-and-latest-bundle-staging-fallback.md}}. Тип: markdown. Размер: 3390 байт. Форма: 81 строк. Документ Markdown, 81 строк, явных URL: 0. Главный заголовок: 2026-05-10 Family-State Root Alias And Latest-Bundle Staging Fallback.\par\medskip
\textbf{\path{docs/audit/2026-05-10-hidden-token-spend-and-incremental-trainer-fixes.md}}. Тип: markdown. Размер: 6740 байт. Форма: 154 строк. Документ Markdown, 154 строк, явных URL: 0. Главный заголовок: 2026-05-10 Hidden Token Spend And Incremental Trainer Fixes.\par\medskip
\textbf{\path{docs/audit/2026-05-10-live-trainer-service-replays-processed-ledger-while-queue-empty.md}}. Тип: markdown. Размер: 5891 байт. Форма: 160 строк. Документ Markdown, 160 строк, явных URL: 0. Главный заголовок: 2026-05-10 Live Trainer Service Replays Processed Ledger While Queue Is Empty.\par\medskip
\textbf{\path{docs/audit/2026-05-10-stale-submodule-runtime-image-root-cause-and-source-preference-fix.md}}. Тип: markdown. Размер: 5021 байт. Форма: 114 строк. Документ Markdown, 114 строк, явных URL: 0. Главный заголовок: 2026-05-10 Stale Submodule Runtime Image Root Cause And Source-Preference Fix.\par\medskip
\textbf{\path{docs/audit/2026-05-11-family-runtime-artifact-fallback-and-staging-path-fix.md}}. Тип: markdown. Размер: 5601 байт. Форма: 139 строк. Документ Markdown, 139 строк, явных URL: 0. Главный заголовок: 2026-05-11 Family Runtime Artifact Fallback And Bundle Staging Path Fix.\par\medskip
\textbf{\path{docs/audit/2026-05-11-fresh-run-family-match-without-family-artifact-runtime.md}}. Тип: markdown. Размер: 6472 байт. Форма: 211 строк. Документ Markdown, 211 строк, явных URL: 0. Главный заголовок: 2026-05-11 Fresh Run Family Match Without Family Artifact Runtime.\par\medskip
\textbf{\path{docs/audit/2026-05-11-idle-trainer-cycles-no-longer-autorepublish-bundles.md}}. Тип: markdown. Размер: 3162 байт. Форма: 108 строк. Документ Markdown, 108 строк, явных URL: 0. Главный заголовок: 2026-05-11 Idle Trainer Cycles No Longer Auto-Republish Bundle Versions.\par\medskip
\textbf{\path{docs/audit/2026-05-11-original-prompt-routing-and-batch-handoff-fixes.md}}. Тип: markdown. Размер: 3829 байт. Форма: 91 строк. Документ Markdown, 91 строк, явных URL: 0. Главный заголовок: 2026-05-11 Original Prompt Routing And Batch Handoff Fixes.\par\medskip
\textbf{\path{docs/audit/2026-05-11-trainer-queue-only-no-processed-replay.md}}. Тип: markdown. Размер: 6373 байт. Форма: 166 строк. Документ Markdown, 166 строк, явных URL: 0. Главный заголовок: 2026-05-11 Trainer Is Now Queue-Only And Will Not Replay Processed History.\par\medskip
\textbf{\path{docs/audit/2026-05-11-trainer-service-preflights-queue-before-running-cycles.md}}. Тип: markdown. Размер: 4211 байт. Форма: 121 строк. Документ Markdown, 121 строк, явных URL: 0. Главный заголовок: 2026-05-11 Trainer Service Preflights Queue Before Running Cycles.\par\medskip
\textbf{\path{docs/audit/2026-05-11-trusted-handoff-syntax-error-blocked-execution-artifact-upload.md}}. Тип: markdown. Размер: 4031 байт. Форма: 86 строк. Документ Markdown, 86 строк, явных URL: 0. Главный заголовок: 2026-05-11 Trusted Handoff Syntax Error Blocked Execution Artifact Upload.\par\medskip
\textbf{\path{docs/audit/README.md}}. Тип: markdown. Размер: 443 байт. Форма: 14 строк. Документ Markdown, 14 строк, явных URL: 0. Главный заголовок: Repository Audit Notes.\par\medskip
\textbf{\path{docs/guides/hushwheel-fixture-rag-guide.md}}. Тип: markdown. Размер: 5169 байт. Форма: 138 строк. Документ Markdown, 138 строк, явных URL: 0. Главный заголовок: RAG Guide For The Hushwheel Fixture.\par\medskip
\textbf{\path{docs/operations/azure-deployment.md}}. Тип: markdown. Размер: 4901 байт. Форма: 111 строк. Документ Markdown, 111 строк, явных URL: 5. Главный заголовок: Azure Deployment Notes.\par\medskip
\textbf{\path{docs/operations/environment.md}}. Тип: markdown. Размер: 7589 байт. Форма: 130 строк. Документ Markdown, 130 строк, явных URL: 0. Главный заголовок: Environment Variables.\par\medskip
\textbf{\path{docs/operations/runtime-image.md}}. Тип: markdown. Размер: 1989 байт. Форма: 46 строк. Документ Markdown, 46 строк, явных URL: 0. Главный заголовок: Runtime Image Notes.\par\medskip
\textbf{\path{docs/operations/simple-end-to-end-verification-guide.md}}. Тип: markdown. Размер: 4111 байт. Форма: 105 строк. Документ Markdown, 105 строк, явных URL: 12. Главный заголовок: Simple End-To-End Verification Guide.\par\medskip
\textbf{\path{docs/operations/trainer-deployment.md}}. Тип: markdown. Размер: 6047 байт. Форма: 152 строк. Документ Markdown, 152 строк, явных URL: 0. Главный заголовок: Trainer Deployment Notes.\par\medskip
\textbf{\path{docs/planning/codex-exec-resume-plan.md}}. Тип: markdown. Размер: 11561 байт. Форма: 207 строк. Документ Markdown, 207 строк, явных URL: 0. Главный заголовок: Codex Exec Resume Plan.\par\medskip
\textbf{\path{docs/planning/dataset-integration-plan.md}}. Тип: markdown. Размер: 33882 байт. Форма: 478 строк. Документ Markdown, 478 строк, явных URL: 0. Главный заголовок: Dataset Integration Plan.\par\medskip
\textbf{\path{docs/planning/family-first-mipro-runtime-contract.md}}. Тип: markdown. Размер: 29074 байт. Форма: 439 строк. Документ Markdown, 439 строк, явных URL: 0. Главный заголовок: Family-First DSPy / MIPROv2 Contract.\par\medskip
\textbf{\path{docs/planning/per-turn-dspy-mediation-contract.md}}. Тип: markdown. Размер: 4966 байт. Форма: 115 строк. Документ Markdown, 115 строк, явных URL: 0. Главный заголовок: Per-Turn DSPy Mediation Contract.\par\medskip
\textbf{\path{docs/planning/repo-hardening-plan.md}}. Тип: markdown. Размер: 4545 байт. Форма: 71 строк. Документ Markdown, 71 строк, явных URL: 0. Главный заголовок: Repository Hardening Plan.\par\medskip
\textbf{\path{docs/planning/trainer-context-group-champion-plan.md}}. Тип: markdown. Размер: 7736 байт. Форма: 119 строк. Документ Markdown, 119 строк, явных URL: 0. Главный заголовок: Trainer Context-Group Champion Plan.\par\medskip
\textbf{\path{mkdocs.yml}}. Тип: yaml. Размер: 843 байт. Форма: 37 строк. Файл YAML с данными или workflow, 37 строк, явных URL: 2. Верхние ключи: site\_name, site\_description, site\_url.\par\medskip
\textbf{\path{notebooks/01_repo_rag_research.ipynb}}. Тип: notebook. Размер: 9324 байт. Форма: 281 строк. Ноутбук, 11 ячеек, явных URL: 0. Первый заголовок: Repository RAG Research Playbook.\par\medskip
\textbf{\path{notebooks/02_agent_workflow_checklist.ipynb}}. Тип: notebook. Размер: 17632 байт. Форма: 404 строк. Ноутбук, 9 ячеек, явных URL: 1. Первый заголовок: Agent Workflow Checklist.\par\medskip
\textbf{\path{notebooks/03_dspy_training_lab.ipynb}}. Тип: notebook. Размер: 12048 байт. Форма: 314 строк. Ноутбук, 11 ячеек, явных URL: 0. Первый заголовок: DSPy Training Sample Lab.\par\medskip
\textbf{\path{notebooks/04_sample_population_lab.ipynb}}. Тип: notebook. Размер: 7894 байт. Форма: 250 строк. Ноутбук, 9 ячеек, явных URL: 0. Первый заголовок: Sample Population Lab.\par\medskip
\textbf{\path{notebooks/05_hushwheel_fixture_rag_lab.ipynb}}. Тип: notebook. Размер: 13622 байт. Форма: 378 строк. Ноутбук, 9 ячеек, явных URL: 0. Первый заголовок: Hushwheel Fixture RAG Lab.\par\medskip
\textbf{\path{publication/.gitignore}}. Тип: generic. Размер: 80 байт. Форма: 12 строк. Отслеживаемый файл репозитория, 12 строк, явных URL: 0. Первая строка: .build/\par\medskip
\textbf{\path{publication/Makefile}}. Тип: make. Размер: 1644 байт. Форма: 48 строк. Makefile репозитория, 48 строк, явных URL: 0. Он открывает общие поверхности автоматизации.\par\medskip
\textbf{\path{publication/README.md}}. Тип: markdown. Размер: 1499 байт. Форма: 34 строк. Документ Markdown, 34 строк, явных URL: 0. Главный заголовок: Publication Article.\par\medskip
\textbf{\path{publication/article-banner.png}}. Тип: binary. Размер: 79074 байт. Форма: бинарная поверхность. Бинарный артефакт, 79074 байт бинарных данных. Он хранится в Git для стабильного публикационного вывода.\par\medskip
\textbf{\path{publication/exploratorium_translation/.gitignore}}. Тип: generic. Размер: 32 байт. Форма: 6 строк. Отслеживаемый файл репозитория, 6 строк, явных URL: 0. Первая строка: .build/\par\medskip
\textbf{\path{publication/exploratorium_translation/Makefile}}. Тип: make. Размер: 870 байт. Форма: 30 строк. Makefile репозитория, 30 строк, явных URL: 0. Он открывает общие поверхности автоматизации.\par\medskip
\textbf{\path{publication/exploratorium_translation/README.md}}. Тип: markdown. Размер: 997 байт. Форма: 28 строк. Документ Markdown, 28 строк, явных URL: 0. Главный заголовок: Exploratorium Translation.\par\medskip
\textbf{\path{publication/exploratorium_translation/exploratorium_translation.tex}}. Тип: latex. Размер: 2573 байт. Форма: 70 строк. Исходник LaTeX, 70 строк, явных URL: 0. Секций в превью: 0.\par\medskip
\textbf{\path{publication/references.bib}}. Тип: bibliography. Размер: 1252 байт. Форма: 28 строк. Библиографический файл, 28 строк, явных URL: 2. Записей в превью: 4.\par\medskip
\textbf{\path{publication/repository-rag-lab-article.pdf}}. Тип: binary. Размер: 236232 байт. Форма: бинарная поверхность. Бинарный артефакт, 236232 байт бинарных данных. Он хранится в Git для стабильного публикационного вывода.\par\medskip
\textbf{\path{publication/repository-rag-lab-article.tex}}. Тип: latex. Размер: 9526 байт. Форма: 208 строк. Исходник LaTeX, 208 строк, явных URL: 0. Секций в превью: 4.\par\medskip
\textbf{\path{publication/todo-backlog-table.tex}}. Тип: latex. Размер: 10547 байт. Форма: 34 строк. Исходник LaTeX, 34 строк, явных URL: 36. Секций в превью: 0.\par\medskip
\textbf{\path{pyproject.toml}}. Тип: toml. Размер: 2285 байт. Форма: 97 строк. Конфигурация TOML, 97 строк, явных URL: 0. Верхние ключи: build-system, project, dependency-groups.\par\medskip
\textbf{\path{rust-cli/Cargo.lock}}. Тип: generic. Размер: 3213 байт. Форма: бинарная поверхность. Отслеживаемый файл репозитория, 3213 байт бинарных данных, явных URL: 0. Первая строка: No non-empty preview line.\par\medskip
\textbf{\path{rust-cli/Cargo.toml}}. Тип: toml. Размер: 144 байт. Форма: 8 строк. Конфигурация TOML, 8 строк, явных URL: 0. Верхние ключи: package, dependencies.\par\medskip
\textbf{\path{rust-cli/src/main.rs}}. Тип: source. Размер: 20422 байт. Форма: 626 строк. Исходный файл, 626 строк, явных URL: 0. Первая строка: use rusqlite::\{params, Connection\};\par\medskip
\textbf{\path{samples/README.md}}. Тип: markdown. Размер: 713 байт. Форма: 13 строк. Документ Markdown, 13 строк, явных URL: 0. Главный заголовок: Samples.\par\medskip
\textbf{\path{samples/logs/20260315T050730Z-gh-runs.md}}. Тип: markdown. Размер: 1341 байт. Форма: 45 строк. Документ Markdown, 45 строк, явных URL: 3. Главный заголовок: GitHub Actions Run Log.\par\medskip
\textbf{\path{samples/logs/20260315T085140Z-gh-red-status-check.md}}. Тип: markdown. Размер: 847 байт. Форма: 25 строк. Документ Markdown, 25 строк, явных URL: 1. Главный заголовок: GitHub Actions Red Status Check.\par\medskip
\textbf{\path{samples/logs/20260317T064722Z-full-build.raw.log}}. Тип: generic. Размер: 2109 байт. Форма: 62 строк. Отслеживаемый файл репозитория, 62 строк, явных URL: 0. Первая строка: \# Full Build Raw Log\par\medskip
\textbf{\path{samples/logs/20260317T064833Z-full-build-attempt2.raw.log}}. Тип: generic. Размер: 14950 байт. Форма: 408 строк. Отслеживаемый файл репозитория, 408 строк, явных URL: 0. Первая строка: \# Full Build Raw Log Attempt 2\par\medskip
\textbf{\path{samples/logs/20260317T073401Z-gh-runs-notebook-scaffolding.md}}. Тип: markdown. Размер: 2865 байт. Форма: 35 строк. Документ Markdown, 35 строк, явных URL: 13. Главный заголовок: GitHub Actions Run Log.\par\medskip
\textbf{\path{samples/logs/20260317T073404Z-gh-runs.md}}. Тип: markdown. Размер: 1269 байт. Форма: 32 строк. Документ Markdown, 32 строк, явных URL: 3. Главный заголовок: GitHub Actions Run Log.\par\medskip
\textbf{\path{samples/logs/20260317T080018Z-gh-runs-all-notebooks.md}}. Тип: markdown. Размер: 2128 байт. Форма: 34 строк. Документ Markdown, 34 строк, явных URL: 8. Главный заголовок: GitHub Actions Run Log.\par\medskip
\textbf{\path{samples/logs/20260317T083207Z-gh-runs-todo-backlog.md}}. Тип: markdown. Размер: 1063 байт. Форма: 27 строк. Документ Markdown, 27 строк, явных URL: 3. Главный заголовок: GitHub Actions Run Log.\par\medskip
\textbf{\path{samples/logs/20260317T084723Z-gh-runs-repo-completeness-checklist.md}}. Тип: markdown. Размер: 1186 байт. Форма: 27 строк. Документ Markdown, 27 строк, явных URL: 3. Главный заголовок: GitHub Actions Run Log.\par\medskip
\textbf{\path{samples/logs/20260317T085723Z-gh-runs-hushwheel-fixture.md}}. Тип: markdown. Размер: 1472 байт. Форма: 36 строк. Документ Markdown, 36 строк, явных URL: 1. Главный заголовок: GitHub Actions Run Log.\par\medskip
\textbf{\path{samples/logs/20260317T091641Z-gh-runs-hushwheel-rag-playbook.md}}. Тип: markdown. Размер: 1470 байт. Форма: 36 строк. Документ Markdown, 36 строк, явных URL: 1. Главный заголовок: GitHub Actions Run Log.\par\medskip
\textbf{\path{samples/logs/20260317T092251Z-gh-runs-publication-article.md}}. Тип: markdown. Размер: 1801 байт. Форма: 33 строк. Документ Markdown, 33 строк, явных URL: 5. Главный заголовок: GitHub Actions Run Log.\par\medskip
\textbf{\path{samples/logs/20260318T003027Z-gh-runs-wire-skill-triggers-for-do-repo-verification-aud.md}}. Тип: markdown. Размер: 1113 байт. Форма: 29 строк. Документ Markdown, 29 строк, явных URL: 1. Главный заголовок: GitHub Actions Run Log.\par\medskip
\textbf{\path{samples/logs/20260318T003654Z-gh-runs-execute-all-notebooks.md}}. Тип: markdown. Размер: 1682 байт. Форма: 41 строк. Документ Markdown, 41 строк, явных URL: 1. Главный заголовок: GitHub Actions Run Log.\par\medskip
\textbf{\path{samples/logs/20260318T003927Z-gh-runs-publication-workflow-fix.md}}. Тип: markdown. Размер: 1747 байт. Форма: 35 строк. Документ Markdown, 35 строк, явных URL: 2. Главный заголовок: GitHub Actions Run Log.\par\medskip
\textbf{\path{samples/logs/20260318T004712Z-gh-runs-hushwheel-program-harness.md}}. Тип: markdown. Размер: 1798 байт. Форма: 49 строк. Документ Markdown, 49 строк, явных URL: 0. Главный заголовок: GitHub Actions Run Log.\par\medskip
\textbf{\path{samples/logs/20260318T010254Z-gh-runs-refresh-final-notebook-outputs-for-no-move-those.md}}. Тип: markdown. Размер: 1113 байт. Форма: 29 строк. Документ Markdown, 29 строк, явных URL: 1. Главный заголовок: GitHub Actions Run Log.\par\medskip
\textbf{\path{samples/logs/20260318T010637Z-gh-runs-env-refresh-retest.md}}. Тип: markdown. Размер: 1029 байт. Форма: 26 строк. Документ Markdown, 26 строк, явных URL: 1. Главный заголовок: GitHub Actions Run Log.\par\medskip
\textbf{\path{samples/logs/20260318T010741Z-gh-runs-consolidate-session-master.md}}. Тип: markdown. Размер: 1950 байт. Форма: 62 строк. Документ Markdown, 62 строк, явных URL: 0. Главный заголовок: GitHub Actions Run Log.\par\medskip
\textbf{\path{samples/logs/20260318T012137Z-gh-runs-azure-inference-endpoint-probe.md}}. Тип: markdown. Размер: 1004 байт. Форма: 26 строк. Документ Markdown, 26 строк, явных URL: 1. Главный заголовок: GitHub Actions Run Log.\par\medskip
\textbf{\path{samples/logs/20260318T012714Z-gh-runs-tighten-repo-hygiene-for-other-agnet-is-doing-1-.md}}. Тип: markdown. Размер: 1091 байт. Форма: 29 строк. Документ Markdown, 29 строк, явных URL: 1. Главный заголовок: GitHub Actions Run Log.\par\medskip
\textbf{\path{samples/logs/20260318T015109Z-gh-runs-improve-notebook-run-observability-and-reporting.md}}. Тип: markdown. Размер: 2151 байт. Форма: 55 строк. Документ Markdown, 55 строк, явных URL: 2. Главный заголовок: GitHub Actions Run Log.\par\medskip
\textbf{\path{samples/logs/20260318T021736Z-gh-runs-sync-todo-backlog-into-markdown-and-publication-surfaces.md}}. Тип: markdown. Размер: 1778 байт. Форма: 55 строк. Документ Markdown, 55 строк, явных URL: 2. Главный заголовок: GitHub Actions Run Log.\par\medskip
\textbf{\path{samples/logs/20260318T023922Z-gh-runs-implement-step-1-dspy-training-path-follow-up.md}}. Тип: markdown. Размер: 1735 байт. Форма: 55 строк. Документ Markdown, 55 строк, явных URL: 2. Главный заголовок: GitHub Actions Run Log.\par\medskip
\textbf{\path{samples/logs/20260318T044242Z-gh-runs-add-maintained-repository-research-narrative.md}}. Тип: markdown. Размер: 2222 байт. Форма: 56 строк. Документ Markdown, 56 строк, явных URL: 2. Главный заголовок: GitHub Actions Run Log.\par\medskip
\textbf{\path{samples/logs/20260318T050051Z-gh-runs-add-azure-runtime-surfaces-for-do-1-and-2-and-carry-on-to-next.md}}. Тип: markdown. Размер: 1828 байт. Форма: 38 строк. Документ Markdown, 38 строк, явных URL: 2. Главный заголовок: GitHub Actions Run Log.\par\medskip
\textbf{\path{samples/logs/20260318T051627Z-gh-runs-develop-retrieval-evaluation-suite.md}}. Тип: markdown. Размер: 1768 байт. Форма: 37 строк. Документ Markdown, 37 строк, явных URL: 2. Главный заголовок: GitHub Actions Run Log.\par\medskip
\textbf{\path{samples/logs/20260318T060918Z-gh-runs-consolidate-session-master-repair.md}}. Тип: markdown. Размер: 2460 байт. Форма: 70 строк. Документ Markdown, 70 строк, явных URL: 3. Главный заголовок: GitHub Actions Run Log.\par\medskip
\textbf{\path{samples/logs/20260318T063755Z-gh-runs-keep-utility-sync-tests-side-effect-free.md}}. Тип: markdown. Размер: 1074 байт. Форма: 29 строк. Документ Markdown, 29 строк, явных URL: 1. Главный заголовок: GitHub Actions Run Log.\par\medskip
\textbf{\path{samples/logs/20260318T064924Z-gh-runs-retrieval-refresh-go-ahead.md}}. Тип: markdown. Размер: 1858 байт. Форма: 54 строк. Документ Markdown, 54 строк, явных URL: 2. Главный заголовок: GitHub Actions Run Log.\par\medskip
\textbf{\path{samples/logs/20260318T083955Z-gh-runs-benchmark-broadening-2.md}}. Тип: markdown. Размер: 1217 байт. Форма: 34 строк. Документ Markdown, 34 строк, явных URL: 1. Главный заголовок: GitHub Actions Run Log.\par\medskip
\textbf{\path{samples/logs/20260318T090625Z-gh-runs-hushwheel-quality-instrumentation.md}}. Тип: markdown. Размер: 1312 байт. Форма: 29 строк. Документ Markdown, 29 строк, явных URL: 1. Главный заголовок: GitHub Actions Run Log.\par\medskip
\textbf{\path{samples/logs/20260318T091252Z-gh-runs-add-rust-lookup-and-retrieval-tag-summaries-for-.md}}. Тип: markdown. Размер: 2223 байт. Форма: 53 строк. Документ Markdown, 53 строк, явных URL: 2. Главный заголовок: GitHub Actions Run Log.\par\medskip
\textbf{\path{samples/logs/20260318T093632Z-gh-runs-retrieval-regression-gate-great-carry-on.md}}. Тип: markdown. Размер: 1325 байт. Форма: 34 строк. Документ Markdown, 34 строк, явных URL: 2. Главный заголовок: GitHub Actions Run Log.\par\medskip
\textbf{\path{samples/logs/20260318T094120Z-gh-runs-dspy-live-proof-and-rust-wrapper-repair.md}}. Тип: markdown. Размер: 1469 байт. Форма: 45 строк. Документ Markdown, 45 строк, явных URL: 0. Главный заголовок: GitHub Actions Summary.\par\medskip
\textbf{\path{samples/logs/20260318T094347Z-gh-runs-hushwheel-hardening-followup.md}}. Тип: markdown. Размер: 2739 байт. Форма: 72 строк. Документ Markdown, 72 строк, явных URL: 2. Главный заголовок: GitHub Runs Log.\par\medskip
\textbf{\path{samples/logs/20260318T100336Z-gh-runs-sync-final-rebased-inventories-for-carry-on-comm.md}}. Тип: markdown. Размер: 1103 байт. Форма: 29 строк. Документ Markdown, 29 строк, явных URL: 1. Главный заголовок: GitHub Actions Run Log.\par\medskip
\textbf{\path{samples/logs/20260318T101229Z-gh-runs-hushwheel-audit-refresh.md}}. Тип: markdown. Размер: 1707 байт. Форма: 37 строк. Документ Markdown, 37 строк, явных URL: 1. Главный заголовок: GitHub Actions Summary.\par\medskip
\textbf{\path{samples/logs/20260318T101939Z-gh-runs-land-each-unlanded.md}}. Тип: markdown. Размер: 1556 байт. Форма: 28 строк. Документ Markdown, 28 строк, явных URL: 2. Главный заголовок: GitHub Actions Run Log.\par\medskip
\textbf{\path{samples/logs/20260318T123402Z-gh-runs-fix-file-summary-test-churn-for-1.md}}. Тип: markdown. Размер: 1061 байт. Форма: 29 строк. Документ Markdown, 29 строк, явных URL: 1. Главный заголовок: GitHub Actions Run Log.\par\medskip
\textbf{\path{samples/logs/20260318T124437Z-gh-runs-reconcile-stale-dspy-backlog-for-1.md}}. Тип: markdown. Размер: 1060 байт. Форма: 21 строк. Документ Markdown, 21 строк, явных URL: 3. Главный заголовок: GitHub Actions Log.\par\medskip
\textbf{\path{samples/logs/20260318T125547Z-gh-runs-wire-lookup-first-ask-for-1.md}}. Тип: markdown. Размер: 1058 байт. Форма: 22 строк. Документ Markdown, 22 строк, явных URL: 2. Главный заголовок: GitHub Actions Log.\par\medskip
\textbf{\path{samples/logs/20260318T125829Z-gh-runs-restore-retrieval-gate-after-rebase.md}}. Тип: markdown. Размер: 1585 байт. Форма: 36 строк. Документ Markdown, 36 строк, явных URL: 0. Главный заголовок: GitHub Actions Run Log.\par\medskip
\textbf{\path{samples/logs/20260318T130525Z-gh-runs-finish-node24-publication-cache-followup-for-1.md}}. Тип: markdown. Размер: 1119 байт. Форма: 27 строк. Документ Markdown, 27 строк, явных URL: 3. Главный заголовок: GitHub Actions Status For `Finish Node 24 publication cache follow-up for "1"`.\par\medskip
\textbf{\path{samples/logs/20260318T131709Z-gh-runs-restore-root-readme-weighting-for-1.md}}. Тип: markdown. Размер: 1680 байт. Форма: 45 строк. Документ Markdown, 45 строк, явных URL: 2. Главный заголовок: GitHub Actions Run Log.\par\medskip
\textbf{\path{samples/logs/20260318T133247Z-gh-runs-improve-retrieval-doc-priority-for-1.md}}. Тип: markdown. Размер: 1335 байт. Форма: 34 строк. Документ Markdown, 34 строк, явных URL: 1. Главный заголовок: GitHub Actions Run Log.\par\medskip
\textbf{\path{samples/logs/20260318T144102Z-gh-runs-merge-into-master.md}}. Тип: markdown. Размер: 2285 байт. Форма: 69 строк. Документ Markdown, 69 строк, явных URL: 2. Главный заголовок: GitHub Actions Run Log.\par\medskip
\textbf{\path{samples/logs/20260318T145935Z-gh-runs-consolidate-all-unconsoldiated-changes.md}}. Тип: markdown. Размер: 2401 байт. Форма: 71 строк. Документ Markdown, 71 строк, явных URL: 2. Главный заголовок: GitHub Actions Run Log.\par\medskip
\textbf{\path{samples/logs/20260319T035301Z-gh-runs-public-pages-and-workflow-repairs.md}}. Тип: markdown. Размер: 1578 байт. Форма: 37 строк. Документ Markdown, 37 строк, явных URL: 1. Главный заголовок: GitHub Actions Log.\par\medskip
\textbf{\path{samples/logs/20260319T071544Z-gh-runs-doncsoldiate-consolidation.md}}. Тип: markdown. Размер: 1137 байт. Форма: 27 строк. Документ Markdown, 27 строк, явных URL: 0. Главный заголовок: GitHub Actions Log.\par\medskip
\textbf{\path{samples/logs/20260319T114913Z-gh-runs-fix-publication-verification-for-whats-next-step.md}}. Тип: markdown. Размер: 2936 байт. Форма: 67 строк. Документ Markdown, 67 строк, явных URL: 3. Главный заголовок: GitHub Actions Run Log.\par\medskip
\textbf{\path{samples/logs/20260429T172137Z-gh-runs-develop-no-workflow.md}}. Тип: markdown. Размер: 866 байт. Форма: 23 строк. Документ Markdown, 23 строк, явных URL: 0. Главный заголовок: GitHub Actions Run Log.\par\medskip
\textbf{\path{samples/logs/20260429T200446Z-gh-runs-develop-no-new-run-after-blob-queue-migration.md}}. Тип: markdown. Размер: 1843 байт. Форма: 34 строк. Документ Markdown, 34 строк, явных URL: 0. Главный заголовок: GitHub Actions Run Log.\par\medskip
\textbf{\path{samples/logs/20260430T122222Z-gh-runs-develop-no-new-run-after-codex-proxy-push.md}}. Тип: markdown. Размер: 908 байт. Форма: 25 строк. Документ Markdown, 25 строк, явных URL: 0. Главный заголовок: GitHub Actions check after pushing `bbdcf4b`.\par\medskip
\textbf{\path{samples/logs/20260430T172518Z-gh-runs-develop-no-new-run-after-trace-handoff-push.md}}. Тип: markdown. Размер: 991 байт. Форма: 27 строк. Документ Markdown, 27 строк, явных URL: 0. Главный заголовок: GitHub Actions check after pushing `839f4cb`.\par\medskip
\textbf{\path{samples/logs/20260501T103130Z-gh-runs-develop-no-new-run-after-azure-runtime-env-fix-push.md}}. Тип: markdown. Размер: 1411 байт. Форма: 35 строк. Документ Markdown, 35 строк, явных URL: 0. Главный заголовок: GitHub Actions check after pushing `9e627e1`.\par\medskip
\textbf{\path{samples/logs/20260501T112330Z-gh-runs-develop-no-new-run-after-azure-responses-api-version-fix-push.md}}. Тип: markdown. Размер: 1522 байт. Форма: 37 строк. Документ Markdown, 37 строк, явных URL: 0. Главный заголовок: GitHub Actions check after pushing `128874b`.\par\medskip
\textbf{\path{samples/logs/20260501T125058Z-gh-runs-develop-no-new-run-after-trusted-trace-handoff-push.md}}. Тип: markdown. Размер: 1355 байт. Форма: 30 строк. Документ Markdown, 30 строк, явных URL: 0. Главный заголовок: GitHub Actions check after pushing `66a7ad4`.\par\medskip
\textbf{\path{samples/logs/20260501T162125Z-gh-runs-develop-no-new-run-after-live-trainer-bundle-publication-push.md}}. Тип: markdown. Размер: 1486 байт. Форма: 30 строк. Документ Markdown, 30 строк, явных URL: 0. Главный заголовок: GitHub Run Inspection.\par\medskip
\textbf{\path{samples/logs/20260501T174147Z-gh-runs-develop-no-new-run-after-versioned-trainer-recovery-push.md}}. Тип: markdown. Размер: 1505 байт. Форма: 30 строк. Документ Markdown, 30 строк, явных URL: 0. Главный заголовок: GitHub Run Inspection.\par\medskip
\textbf{\path{samples/logs/20260501T201242Z-gh-runs-develop-no-new-run-after-dspy-bundle-version-pinning-push.md}}. Тип: markdown. Размер: 2076 байт. Форма: 41 строк. Документ Markdown, 41 строк, явных URL: 0. Главный заголовок: GitHub Run Inspection.\par\medskip
\textbf{\path{samples/logs/20260502T064702Z-gh-runs-develop-no-new-run-after-live-trainer-timestamp-bundle-audit-push.md}}. Тип: markdown. Размер: 1604 байт. Форма: 38 строк. Документ Markdown, 38 строк, явных URL: 0. Главный заголовок: GitHub Run Inspection.\par\medskip
\textbf{\path{samples/logs/20260502T082402Z-gh-runs-develop-no-new-run-after-hybrid-vector-retrieval-push.md}}. Тип: markdown. Размер: 1631 байт. Форма: 30 строк. Документ Markdown, 30 строк, явных URL: 0. Главный заголовок: GitHub Run Inspection.\par\medskip
\textbf{\path{samples/logs/20260502T121755Z-gh-runs-develop-no-new-run-after-trainer-gating-and-retrieval-isolation-push.md}}. Тип: markdown. Размер: 1175 байт. Форма: 25 строк. Документ Markdown, 25 строк, явных URL: 0. Главный заголовок: Summary.\par\medskip
\textbf{\path{samples/logs/20260502T165750Z-gh-runs-develop-no-new-run-after-trainer-champion-batching-push.md}}. Тип: markdown. Размер: 1543 байт. Форма: 35 строк. Документ Markdown, 35 строк, явных URL: 0. Главный заголовок: GitHub Run Inspection.\par\medskip
\textbf{\path{samples/logs/20260503T080503Z-gh-runs-develop-no-new-run-after-codex-resume-lane-telemetry-push.md}}. Тип: markdown. Размер: 1782 байт. Форма: 37 строк. Документ Markdown, 37 строк, явных URL: 0. Главный заголовок: GitHub Run Inspection.\par\medskip
\textbf{\path{samples/logs/20260503T153706Z-gh-runs-develop-no-new-run-after-staged-bundle-mirror-push.md}}. Тип: markdown. Размер: 1171 байт. Форма: 27 строк. Документ Markdown, 27 строк, явных URL: 0. Главный заголовок: GitHub Runs After `Support staged bundle mirrors and stale queue skips`.\par\medskip
\textbf{\path{samples/logs/20260503T160232Z-gh-runs-develop-no-new-run-after-no-op-trainer-promotion-fix-push.md}}. Тип: markdown. Размер: 1482 байт. Форма: 32 строк. Документ Markdown, 32 строк, явных URL: 0. Главный заголовок: GitHub Run Inspection.\par\medskip
\textbf{\path{samples/logs/20260504T090704Z-gh-runs-develop-no-new-run-after-fix-codex-resume-pvc-session-path-push.md}}. Тип: markdown. Размер: 1689 байт. Форма: 46 строк. Документ Markdown, 46 строк, явных URL: 0. Главный заголовок: GitHub Run Check After `Fix Codex resume PVC session path`.\par\medskip
\textbf{\path{samples/logs/20260504T111427Z-gh-runs-develop-no-new-run-after-codex-resume-restore-fallback-debug-push.md}}. Тип: markdown. Размер: 3370 байт. Форма: 72 строк. Документ Markdown, 72 строк, явных URL: 0. Главный заголовок: GitHub Run Check After `Document Codex resume restore fallback debug`.\par\medskip
\textbf{\path{samples/logs/20260504T141246Z-gh-runs-develop-no-new-run-after-codex-resume-guard-hardening-push.md}}. Тип: markdown. Размер: 2837 байт. Форма: 67 строк. Документ Markdown, 67 строк, явных URL: 0. Главный заголовок: GitHub Run Check After `Harden Codex resume restore guards`.\par\medskip
\textbf{\path{samples/logs/20260504T161106Z-gh-runs-develop-no-new-run-after-explicit-codex-resume-targeting-push.md}}. Тип: markdown. Размер: 1175 байт. Форма: 27 строк. Документ Markdown, 27 строк, явных URL: 0. Главный заголовок: GitHub Actions inspection after `464e16b` push.\par\medskip
\textbf{\path{samples/logs/20260504T180441Z-gh-runs-develop-no-new-run-after-codex-session-continuity-markers-push.md}}. Тип: markdown. Размер: 986 байт. Форма: 28 строк. Документ Markdown, 28 строк, явных URL: 0. Главный заголовок: GitHub Actions inspection after `490fe41` push.\par\medskip
\textbf{\path{samples/logs/20260505T072441Z-gh-runs-develop-no-new-run-after-mcp-first-codex-discovery-push.md}}. Тип: markdown. Размер: 1152 байт. Форма: 28 строк. Документ Markdown, 28 строк, явных URL: 0. Главный заголовок: GitHub Actions inspection after `abcf52c` push.\par\medskip
\textbf{\path{samples/logs/README.md}}. Тип: markdown. Размер: 331 байт. Форма: 11 строк. Документ Markdown, 11 строк, явных URL: 0. Главный заголовок: GitHub Run Logs.\par\medskip
\textbf{\path{samples/population/hushwheel_fixture_population_candidates.yaml}}. Тип: yaml. Размер: 955 байт. Форма: 22 строк. Файл YAML с данными или workflow, 22 строк, явных URL: 0. Верхние ключи: no top-level mapping keys.\par\medskip
\textbf{\path{samples/population/repository_population_candidates.yaml}}. Тип: yaml. Размер: 690 байт. Форма: 16 строк. Файл YAML с данными или workflow, 16 строк, явных URL: 0. Верхние ключи: no top-level mapping keys.\par\medskip
\textbf{\path{samples/training/hushwheel_fixture_training_examples.yaml}}. Тип: yaml. Размер: 1852 байт. Форма: 68 строк. Файл YAML с данными или workflow, 68 строк, явных URL: 0. Верхние ключи: no top-level mapping keys.\par\medskip
\textbf{\path{samples/training/repository_training_examples.yaml}}. Тип: yaml. Размер: 2504 байт. Форма: 80 строк. Файл YAML с данными или workflow, 80 строк, явных URL: 0. Верхние ключи: no top-level mapping keys.\par\medskip
\textbf{\path{src/repo_rag_lab/__init__.py}}. Тип: python. Размер: 1087 байт. Форма: 47 строк. Модуль Python, 47 строк, явных URL: 0. Назначение модуля: Public package surface for the repository-grounded RAG lab.\par\medskip
\textbf{\path{src/repo_rag_lab/azure.py}}. Тип: python. Размер: 2592 байт. Форма: 76 строк. Модуль Python, 76 строк, явных URL: 0. Назначение модуля: Azure deployment manifest helpers for downstream release workflows.\par\medskip
\textbf{\path{src/repo_rag_lab/azure_artifacts.py}}. Тип: python. Размер: 16790 байт. Форма: 436 строк. Модуль Python, 436 строк, явных URL: 0. Назначение модуля: Azure Blob + Queue helpers for trainer-side traces and global bundle storage.\par\medskip
\textbf{\path{src/repo_rag_lab/azure_runtime.py}}. Тип: python. Размер: 20123 байт. Форма: 577 строк. Модуль Python, 577 строк, явных URL: 0. Назначение модуля: Azure runtime normalization, live probes, and cloud-completion helpers.\par\medskip
\textbf{\path{src/repo_rag_lab/benchmarks.py}}. Тип: python. Размер: 16524 байт. Форма: 484 строк. Модуль Python, 484 строк, явных URL: 0. Назначение модуля: Retrieval benchmark helpers for notebook assertions and logging.\par\medskip
\textbf{\path{src/repo_rag_lab/cli.py}}. Тип: python. Размер: 58870 байт. Форма: 1339 строк. Модуль Python, 1339 строк, явных URL: 0. Назначение модуля: Command-line entrypoints for the shared repository RAG workflows.\par\medskip
\textbf{\path{src/repo_rag_lab/codex_proxy.py}}. Тип: python. Размер: 75460 байт. Форма: 1946 строк. Модуль Python, 1946 строк, явных URL: 0. Назначение модуля: Local streaming proxy that injects repo-RAG + DSPy mediation for Codex.\par\medskip
\textbf{\path{src/repo_rag_lab/corpus.py}}. Тип: python. Размер: 5245 байт. Форма: 181 строк. Модуль Python, 181 строк, явных URL: 0. Назначение модуля: Repository file loading helpers for the baseline text corpus.\par\medskip
\textbf{\path{src/repo_rag_lab/dspy_training.py}}. Тип: python. Размер: 61618 байт. Форма: 1612 строк. Модуль Python, 1612 строк, явных URL: 0. Назначение модуля: DSPy training and artifact helpers for repository-grounded RAG.\par\medskip
\textbf{\path{src/repo_rag_lab/dspy_workflow.py}}. Тип: python. Размер: 6518 байт. Форма: 180 строк. Модуль Python, 180 строк, явных URL: 0. Назначение модуля: DSPy-backed runtime wrappers around the repository retrieval workflow.\par\medskip
\textbf{\path{src/repo_rag_lab/exploratorium_translation.py}}. Тип: python. Размер: 38049 байт. Форма: 1009 строк. Модуль Python, 1009 строк, явных URL: 0. Назначение модуля: Generate the bilingual exploratorium translation publication surfaces.\par\medskip
\textbf{\path{src/repo_rag_lab/file_summaries.py}}. Тип: python. Размер: 10564 байт. Форма: 327 строк. Модуль Python, 327 строк, явных URL: 0. Назначение модуля: Generate tracked repository file-summary surfaces.\par\medskip
\textbf{\path{src/repo_rag_lab/github_pr_gates.py}}. Тип: python. Размер: 3979 байт. Форма: 140 строк. Модуль Python, 140 строк, явных URL: 0. Назначение модуля: GitHub branch-protection helpers for required pull-request gate checks.\par\medskip
\textbf{\path{src/repo_rag_lab/mcp.py}}. Тип: python. Размер: 2947 байт. Форма: 84 строк. Модуль Python, 84 строк, явных URL: 0. Назначение модуля: Discovery helpers for MCP-related repository artifacts.\par\medskip
\textbf{\path{src/repo_rag_lab/mcp_server.py}}. Тип: python. Размер: 50399 байт. Форма: 1361 строк. Модуль Python, 1361 строк, явных URL: 0. Назначение модуля: Minimal stdio MCP server for bounded repo-RAG operations.\par\medskip
\textbf{\path{src/repo_rag_lab/mcp_stdio.py}}. Тип: python. Размер: 953 байт. Форма: 38 строк. Модуль Python, 38 строк, явных URL: 0. Назначение модуля: Lightweight stdio entrypoint for the bounded repo-RAG MCP server.\par\medskip
\textbf{\path{src/repo_rag_lab/notebook_runner.py}}. Тип: python. Размер: 16433 байт. Форма: 491 строк. Модуль Python, 491 строк, явных URL: 0. Назначение модуля: Monitored batch execution helpers for repository notebooks.\par\medskip
\textbf{\path{src/repo_rag_lab/notebook_scaffolding.py}}. Тип: python. Размер: 10965 байт. Форма: 258 строк. Модуль Python, 258 строк, явных URL: 0. Назначение модуля: High-level notebook scaffolds built from tested repository helpers.\par\medskip
\textbf{\path{src/repo_rag_lab/notebook_support.py}}. Тип: python. Размер: 2843 байт. Форма: 89 строк. Модуль Python, 89 строк, явных URL: 0. Назначение модуля: Helpers shared by notebooks so their setup logic stays in tested Python code.\par\medskip
\textbf{\path{src/repo_rag_lab/pages_site.py}}. Тип: python. Размер: 13635 байт. Форма: 390 строк. Модуль Python, 390 строк, явных URL: 2. Назначение модуля: Generate a MkDocs-ready catalog of tracked repository Markdown files.\par\medskip
\textbf{\path{src/repo_rag_lab/population_samples.py}}. Тип: python. Размер: 8376 байт. Форма: 241 строк. Модуль Python, 241 строк, явных URL: 0. Назначение модуля: Helpers for reviewing and batching starter corpus-population inputs.\par\medskip
\textbf{\path{src/repo_rag_lab/retrieval.py}}. Тип: python. Размер: 25533 байт. Форма: 788 строк. Модуль Python, 788 строк, явных URL: 0. Назначение модуля: Baseline chunking and lexical retrieval utilities for repository text.\par\medskip
\textbf{\path{src/repo_rag_lab/retrieval_profile.py}}. Тип: python. Размер: 9527 байт. Форма: 236 строк. Модуль Python, 236 строк, явных URL: 0. Назначение модуля: Profile-driven retrieval weighting, exclusions, and repo-local overrides.\par\medskip
\textbf{\path{src/repo_rag_lab/runtime_artifacts.py}}. Тип: python. Размер: 118360 байт. Форма: 2671 строк. Модуль Python, 2671 строк, явных URL: 0. Назначение модуля: Versioned bundle, overlay, and runtime-trace helpers.\par\medskip
\textbf{\path{src/repo_rag_lab/rust_lookup.py}}. Тип: python. Размер: 7724 байт. Форма: 295 строк. Модуль Python, 295 строк, явных URL: 0. Назначение модуля: Helpers for the native Rust-backed SQLite lookup index.\par\medskip
\textbf{\path{src/repo_rag_lab/semantic_retrieval.py}}. Тип: python. Размер: 7407 байт. Форма: 214 строк. Модуль Python, 214 строк, явных URL: 0. Назначение модуля: Azure OpenAI embedding-backed semantic retrieval helpers.\par\medskip
\textbf{\path{src/repo_rag_lab/settings.py}}. Тип: python. Размер: 699 байт. Форма: 28 строк. Модуль Python, 28 строк, явных URL: 0. Назначение модуля: Shared repository path settings derived from a repository root.\par\medskip
\textbf{\path{src/repo_rag_lab/todo_backlog.py}}. Тип: python. Размер: 8434 байт. Форма: 276 строк. Модуль Python, 276 строк, явных URL: 1. Назначение модуля: Shared backlog rendering helpers for Markdown and LaTeX surfaces.\par\medskip
\textbf{\path{src/repo_rag_lab/trainer_deployment.py}}. Тип: python. Размер: 18027 байт. Форма: 431 строк. Модуль Python, 431 строк, явных URL: 0. Назначение модуля: Kubernetes manifest helpers for trainer-side repo-RAG deployment surfaces.\par\medskip
\textbf{\path{src/repo_rag_lab/training_samples.py}}. Тип: python. Размер: 110198 байт. Форма: 2709 строк. Модуль Python, 2709 строк, явных URL: 0. Назначение модуля: Helpers for loading and summarizing starter DSPy training examples.\par\medskip
\textbf{\path{src/repo_rag_lab/utilities.py}}. Тип: python. Размер: 116367 байт. Форма: 2809 строк. Модуль Python, 2809 строк, явных URL: 1. Назначение модуля: User-facing utility helpers shared by the CLI, tests, and notebooks.\par\medskip
\textbf{\path{src/repo_rag_lab/verification.py}}. Тип: python. Размер: 4476 байт. Форма: 163 строк. Модуль Python, 163 строк, явных URL: 0. Назначение модуля: Verification helpers for Makefile and notebook contract checks.\par\medskip
\textbf{\path{src/repo_rag_lab/workflow.py}}. Тип: python. Размер: 15900 байт. Форма: 503 строк. Модуль Python, 503 строк, явных URL: 0. Назначение модуля: Baseline repository question-answering workflow built on file retrieval.\par\medskip
\textbf{\path{tests/features/repository_rag.feature}}. Тип: generic. Размер: 377 байт. Форма: 12 строк. Отслеживаемый файл репозитория, 12 строк, явных URL: 0. Первая строка: Feature: Repository RAG\par\medskip
\textbf{\path{tests/fixtures/hushwheel_lexiconarium/.gitignore}}. Тип: generic. Размер: 23 байт. Форма: 4 строк. Отслеживаемый файл репозитория, 4 строк, явных URL: 0. Первая строка: hushwheel\par\medskip
\textbf{\path{tests/fixtures/hushwheel_lexiconarium/CHANGELOG.md}}. Тип: markdown. Размер: 345 байт. Форма: 9 строк. Документ Markdown, 9 строк, явных URL: 0. Главный заголовок: Changelog.\par\medskip
\textbf{\path{tests/fixtures/hushwheel_lexiconarium/Doxyfile}}. Тип: generic. Размер: 1239 байт. Форма: 35 строк. Отслеживаемый файл репозитория, 35 строк, явных URL: 0. Первая строка: PROJECT\_NAME           = Hushwheel Lexiconarium\par\medskip
\textbf{\path{tests/fixtures/hushwheel_lexiconarium/LICENSE.md}}. Тип: markdown. Размер: 1082 байт. Форма: 22 строк. Документ Markdown, 22 строк, явных URL: 0. Главный заголовок: LICENSE.\par\medskip
\textbf{\path{tests/fixtures/hushwheel_lexiconarium/Makefile}}. Тип: make. Размер: 20908 байт. Форма: 388 строк. Makefile репозитория, 388 строк, явных URL: 0. Он открывает общие поверхности автоматизации.\par\medskip
\textbf{\path{tests/fixtures/hushwheel_lexiconarium/README.md}}. Тип: markdown. Размер: 8269 байт. Форма: 162 строк. Документ Markdown, 162 строк, явных URL: 0. Главный заголовок: The Hushwheel Lexiconarium.\par\medskip
\textbf{\path{tests/fixtures/hushwheel_lexiconarium/VERSION}}. Тип: generic. Размер: 6 байт. Форма: 2 строк. Отслеживаемый файл репозитория, 2 строк, явных URL: 0. Первая строка: 0.1.0\par\medskip
\textbf{\path{tests/fixtures/hushwheel_lexiconarium/docs/architecture.md}}. Тип: markdown. Размер: 1404 байт. Форма: 33 строк. Документ Markdown, 33 строк, явных URL: 0. Главный заголовок: Hushwheel Architecture.\par\medskip
\textbf{\path{tests/fixtures/hushwheel_lexiconarium/docs/catalog.md}}. Тип: markdown. Размер: 27404 байт. Форма: 357 строк. Документ Markdown, 357 строк, явных URL: 0. Главный заголовок: Hushwheel Representative Catalog.\par\medskip
\textbf{\path{tests/fixtures/hushwheel_lexiconarium/docs/concepts.md}}. Тип: markdown. Размер: 1559 байт. Форма: 37 строк. Документ Markdown, 37 строк, явных URL: 0. Главный заголовок: Hushwheel Concepts.\par\medskip
\textbf{\path{tests/fixtures/hushwheel_lexiconarium/docs/constellation-atlas.md}}. Тип: markdown. Размер: 5111 байт. Форма: 142 строк. Документ Markdown, 142 строк, явных URL: 0. Главный заголовок: Hushwheel Constellation Atlas §::<[]>.\par\medskip
\textbf{\path{tests/fixtures/hushwheel_lexiconarium/docs/districts.md}}. Тип: markdown. Размер: 1151 байт. Форма: 42 строк. Документ Markdown, 42 строк, явных URL: 0. Главный заголовок: Hushwheel Districts.\par\medskip
\textbf{\path{tests/fixtures/hushwheel_lexiconarium/docs/doxygen-fonts.sty}}. Тип: generic. Размер: 212 байт. Форма: бинарная поверхность. Отслеживаемый файл репозитория, 212 байт бинарных данных, явных URL: 0. Первая строка: No non-empty preview line.\par\medskip
\textbf{\path{tests/fixtures/hushwheel_lexiconarium/docs/hushwheel-reference.pdf}}. Тип: binary. Размер: 132 байт. Форма: бинарная поверхность. Бинарный артефакт, 132 байт бинарных данных. Он хранится в Git для стабильного публикационного вывода.\par\medskip
\textbf{\path{tests/fixtures/hushwheel_lexiconarium/docs/operations.md}}. Тип: markdown. Размер: 1068 байт. Форма: 35 строк. Документ Markdown, 35 строк, явных URL: 0. Главный заголовок: Hushwheel Operations.\par\medskip
\textbf{\path{tests/fixtures/hushwheel_lexiconarium/docs/packaging.md}}. Тип: markdown. Размер: 1112 байт. Форма: 33 строк. Документ Markdown, 33 строк, явных URL: 0. Главный заголовок: Hushwheel Packaging.\par\medskip
\textbf{\path{tests/fixtures/hushwheel_lexiconarium/docs/testing.md}}. Тип: markdown. Размер: 3507 байт. Форма: 80 строк. Документ Markdown, 80 строк, явных URL: 0. Главный заголовок: Hushwheel Testing.\par\medskip
\textbf{\path{tests/fixtures/hushwheel_lexiconarium/fixture-manifest.json}}. Тип: json. Размер: 1098 байт. Форма: 49 строк. Файл JSON, 49 строк, явных URL: 0. Верхние ключи: entry\_count, source\_size\_bytes, doc\_files.\par\medskip
\textbf{\path{tests/fixtures/hushwheel_lexiconarium/include/hushwheel.h}}. Тип: source. Размер: 404 байт. Форма: 19 строк. Исходный файл, 19 строк, явных URL: 0. Первая строка: \#ifndef HUSHWHEEL\_H\par\medskip
\textbf{\path{tests/fixtures/hushwheel_lexiconarium/packaging/hushwheel.1}}. Тип: generic. Размер: 1303 байт. Форма: бинарная поверхность. Отслеживаемый файл репозитория, 1303 байт бинарных данных, явных URL: 0. Первая строка: No non-empty preview line.\par\medskip
\textbf{\path{tests/fixtures/hushwheel_lexiconarium/packaging/hushwheel.package.json}}. Тип: json. Размер: 591 байт. Форма: 29 строк. Файл JSON, 29 строк, явных URL: 0. Верхние ключи: name, version, summary.\par\medskip
\textbf{\path{tests/fixtures/hushwheel_lexiconarium/src/hushwheel.c}}. Тип: source. Размер: 1694579 байт. Форма: 4375 строк. Исходный файл, 4375 строк, явных URL: 0. Первая строка: /**\par\medskip
\textbf{\path{tests/fixtures/hushwheel_lexiconarium/src/hushwheel_internal.h}}. Тип: source. Размер: 603 байт. Форма: 25 строк. Исходный файл, 25 строк, явных URL: 0. Первая строка: \#ifndef HUSHWHEEL\_INTERNAL\_H\par\medskip
\textbf{\path{tests/fixtures/hushwheel_lexiconarium/src/hushwheel_spoke_argent.c}}. Тип: source. Размер: 324091 байт. Форма: 525 строк. Исходный файл, 525 строк, явных URL: 0. Первая строка: \#include "hushwheel\_internal.h"\par\medskip
\textbf{\path{tests/fixtures/hushwheel_lexiconarium/src/hushwheel_spoke_bhagavatam.c}}. Тип: source. Размер: 326160 байт. Форма: 525 строк. Исходный файл, 525 строк, явных URL: 0. Первая строка: \#include "hushwheel\_internal.h"\par\medskip
\textbf{\path{tests/fixtures/hushwheel_lexiconarium/src/hushwheel_spoke_crosscanon.c}}. Тип: source. Размер: 326144 байт. Форма: 525 строк. Исходный файл, 525 строк, явных URL: 0. Первая строка: \#include "hushwheel\_internal.h"\par\medskip
\textbf{\path{tests/fixtures/hushwheel_lexiconarium/src/hushwheel_spoke_editorial.c}}. Тип: source. Размер: 325628 байт. Форма: 525 строк. Исходный файл, 525 строк, явных URL: 0. Первая строка: \#include "hushwheel\_internal.h"\par\medskip
\textbf{\path{tests/fixtures/hushwheel_lexiconarium/src/hushwheel_spoke_gita.c}}. Тип: source. Размер: 323055 байт. Форма: 525 строк. Исходный файл, 525 строк, явных URL: 0. Первая строка: \#include "hushwheel\_internal.h"\par\medskip
\textbf{\path{tests/fixtures/hushwheel_lexiconarium/src/hushwheel_spoke_kernel.c}}. Тип: source. Размер: 324068 байт. Форма: 525 строк. Исходный файл, 525 строк, явных URL: 0. Первая строка: \#include "hushwheel\_internal.h"\par\medskip
\textbf{\path{tests/fixtures/hushwheel_lexiconarium/src/hushwheel_spoke_psalter.c}}. Тип: source. Размер: 324600 байт. Форма: 525 строк. Исходный файл, 525 строк, явных URL: 0. Первая строка: \#include "hushwheel\_internal.h"\par\medskip
\textbf{\path{tests/fixtures/hushwheel_lexiconarium/src/hushwheel_spoke_ranger.c}}. Тип: source. Размер: 324105 байт. Форма: 525 строк. Исходный файл, 525 строк, явных URL: 0. Первая строка: \#include "hushwheel\_internal.h"\par\medskip
\textbf{\path{tests/fixtures/hushwheel_lexiconarium/src/hushwheel_spokes.c}}. Тип: source. Размер: 1074 байт. Форма: 25 строк. Исходный файл, 25 строк, явных URL: 0. Первая строка: \#include "hushwheel\_internal.h"\par\medskip
\textbf{\path{tests/fixtures/hushwheel_lexiconarium/tests/bdd/hushwheel.feature}}. Тип: generic. Размер: 1021 байт. Форма: 28 строк. Отслеживаемый файл репозитория, 28 строк, явных URL: 0. Первая строка: Feature: Hushwheel CLI\par\medskip
\textbf{\path{tests/fixtures/hushwheel_lexiconarium/tests/bdd/run_bdd.py}}. Тип: python. Размер: 3146 байт. Форма: 94 строк. Модуль Python, 94 строк, явных URL: 0. Назначение модуля: No module docstring detected.\par\medskip
\textbf{\path{tests/fixtures/hushwheel_lexiconarium/tests/integration/cli_suite.py}}. Тип: python. Размер: 4482 байт. Форма: 122 строк. Модуль Python, 122 строк, явных URL: 0. Назначение модуля: No module docstring detected.\par\medskip
\textbf{\path{tests/fixtures/hushwheel_lexiconarium/tests/unit/test_hushwheel_unit.c}}. Тип: source. Размер: 2002 байт. Форма: 62 строк. Исходный файл, 62 строк, явных URL: 0. Первая строка: \#include <assert.h>\par\medskip
\textbf{\path{tests/fixtures/hushwheel_lexiconarium/tools/lint_hushwheel.py}}. Тип: python. Размер: 4678 байт. Форма: 143 строк. Модуль Python, 143 строк, явных URL: 0. Назначение модуля: No module docstring detected.\par\medskip
\textbf{\path{tests/fixtures/hushwheel_lexiconarium/tools/regenerate_hushwheel_fixture.py}}. Тип: python. Размер: 31375 байт. Форма: 866 строк. Модуль Python, 866 строк, явных URL: 0. Назначение модуля: No module docstring detected.\par\medskip
\textbf{\path{tests/test_azure_runtime.py}}. Тип: python. Размер: 7540 байт. Форма: 208 строк. Модуль Python, 208 строк, явных URL: 12. Назначение модуля: No module docstring detected.\par\medskip
\textbf{\path{tests/test_benchmarks_and_notebook_scaffolding.py}}. Тип: python. Размер: 11859 байт. Форма: 274 строк. Модуль Python, 274 строк, явных URL: 0. Назначение модуля: No module docstring detected.\par\medskip
\textbf{\path{tests/test_cli_and_dspy.py}}. Тип: python. Размер: 68622 байт. Форма: 1876 строк. Модуль Python, 1876 строк, явных URL: 2. Назначение модуля: No module docstring detected.\par\medskip
\textbf{\path{tests/test_codex_proxy.py}}. Тип: python. Размер: 57154 байт. Форма: 1571 строк. Модуль Python, 1571 строк, явных URL: 5. Назначение модуля: No module docstring detected.\par\medskip
\textbf{\path{tests/test_dspy_training.py}}. Тип: python. Размер: 66652 байт. Форма: 1653 строк. Модуль Python, 1653 строк, явных URL: 11. Назначение модуля: No module docstring detected.\par\medskip
\textbf{\path{tests/test_exploratorium_translation.py}}. Тип: python. Размер: 5239 байт. Форма: 138 строк. Модуль Python, 138 строк, явных URL: 8. Назначение модуля: No module docstring detected.\par\medskip
\textbf{\path{tests/test_file_summaries.py}}. Тип: python. Размер: 9110 байт. Форма: 230 строк. Модуль Python, 230 строк, явных URL: 0. Назначение модуля: No module docstring detected.\par\medskip
\textbf{\path{tests/test_github_pr_gates.py}}. Тип: python. Размер: 5678 байт. Форма: 183 строк. Модуль Python, 183 строк, явных URL: 1. Назначение модуля: No module docstring detected.\par\medskip
\textbf{\path{tests/test_hushwheel_fixture.py}}. Тип: python. Размер: 4617 байт. Форма: 101 строк. Модуль Python, 101 строк, явных URL: 0. Назначение модуля: No module docstring detected.\par\medskip
\textbf{\path{tests/test_hushwheel_program_surface.py}}. Тип: python. Размер: 4626 байт. Форма: 125 строк. Модуль Python, 125 строк, явных URL: 0. Назначение модуля: No module docstring detected.\par\medskip
\textbf{\path{tests/test_live_azure_integration.py}}. Тип: python. Размер: 2634 байт. Форма: 83 строк. Модуль Python, 83 строк, явных URL: 0. Назначение модуля: No module docstring detected.\par\medskip
\textbf{\path{tests/test_lookup_first.py}}. Тип: python. Размер: 3559 байт. Форма: 108 строк. Модуль Python, 108 строк, явных URL: 0. Назначение модуля: No module docstring detected.\par\medskip
\textbf{\path{tests/test_mcp_branches.py}}. Тип: python. Размер: 1077 байт. Форма: 31 строк. Модуль Python, 31 строк, явных URL: 0. Назначение модуля: No module docstring detected.\par\medskip
\textbf{\path{tests/test_mcp_server.py}}. Тип: python. Размер: 26557 байт. Форма: 799 строк. Модуль Python, 799 строк, явных URL: 0. Назначение модуля: No module docstring detected.\par\medskip
\textbf{\path{tests/test_mcp_stdio.py}}. Тип: python. Размер: 988 байт. Форма: 42 строк. Модуль Python, 42 строк, явных URL: 0. Назначение модуля: No module docstring detected.\par\medskip
\textbf{\path{tests/test_notebook_runner.py}}. Тип: python. Размер: 3088 байт. Форма: 84 строк. Модуль Python, 84 строк, явных URL: 0. Назначение модуля: No module docstring detected.\par\medskip
\textbf{\path{tests/test_pages_site.py}}. Тип: python. Размер: 9045 байт. Форма: 225 строк. Модуль Python, 225 строк, явных URL: 14. Назначение модуля: No module docstring detected.\par\medskip
\textbf{\path{tests/test_population_samples.py}}. Тип: python. Размер: 3144 байт. Форма: 76 строк. Модуль Python, 76 строк, явных URL: 0. Назначение модуля: No module docstring detected.\par\medskip
\textbf{\path{tests/test_project_surfaces.py}}. Тип: python. Размер: 21471 байт. Форма: 484 строк. Модуль Python, 484 строк, явных URL: 1. Назначение модуля: No module docstring detected.\par\medskip
\textbf{\path{tests/test_repository_rag_bdd.py}}. Тип: python. Размер: 879 байт. Форма: 30 строк. Модуль Python, 30 строк, явных URL: 0. Назначение модуля: No module docstring detected.\par\medskip
\textbf{\path{tests/test_retrieval.py}}. Тип: python. Размер: 13362 байт. Форма: 389 строк. Модуль Python, 389 строк, явных URL: 0. Назначение модуля: No module docstring detected.\par\medskip
\textbf{\path{tests/test_runtime_artifacts_azure.py}}. Тип: python. Размер: 58883 байт. Форма: 1554 строк. Модуль Python, 1554 строк, явных URL: 4. Назначение модуля: No module docstring detected.\par\medskip
\textbf{\path{tests/test_rust_lookup.py}}. Тип: python. Размер: 3099 байт. Форма: 95 строк. Модуль Python, 95 строк, явных URL: 0. Назначение модуля: No module docstring detected.\par\medskip
\textbf{\path{tests/test_todo_backlog.py}}. Тип: python. Размер: 3505 байт. Форма: 103 строк. Модуль Python, 103 строк, явных URL: 1. Назначение модуля: No module docstring detected.\par\medskip
\textbf{\path{tests/test_training_samples.py}}. Тип: python. Размер: 84870 байт. Форма: 2152 строк. Модуль Python, 2152 строк, явных URL: 0. Назначение модуля: No module docstring detected.\par\medskip
\textbf{\path{tests/test_utilities.py}}. Тип: python. Размер: 121479 байт. Форма: 3117 строк. Модуль Python, 3117 строк, явных URL: 6. Назначение модуля: No module docstring detected.\par\medskip
\textbf{\path{tests/test_verification.py}}. Тип: python. Размер: 3156 байт. Форма: 66 строк. Модуль Python, 66 строк, явных URL: 0. Назначение модуля: No module docstring detected.\par\medskip
\textbf{\path{tests/test_workflow.py}}. Тип: python. Размер: 3456 байт. Форма: 99 строк. Модуль Python, 99 строк, явных URL: 0. Назначение модуля: No module docstring detected.\par\medskip
\textbf{\path{tests/test_workflow_live.py}}. Тип: python. Размер: 4292 байт. Форма: 128 строк. Модуль Python, 128 строк, явных URL: 0. Назначение модуля: No module docstring detected.\par\medskip
\textbf{\path{todo-backlog.yaml}}. Тип: yaml. Размер: 6116 байт. Форма: 128 строк. Файл YAML с данными или workflow, 128 строк, явных URL: 1. Верхние ключи: repository\_web\_base, items.\par\medskip
\textbf{\path{utilities/README.md}}. Тип: markdown. Размер: 3495 байт. Форма: 54 строк. Документ Markdown, 54 строк, явных URL: 0. Главный заголовок: Repository Utilities.\par\medskip
\textbf{\path{uv.lock}}. Тип: generic. Размер: 784153 байт. Форма: бинарная поверхность. Отслеживаемый файл репозитория, 784153 байт бинарных данных, явных URL: 0. Первая строка: No non-empty preview line.\par\medskip
\subsection*{Русский: Сводки по всем авторским URL}
\textbf{\url{http://127.0.0.1}}. Хост: 127.0.0.1. Упоминаний: 2. Ссылка хранится только как URL. Локальное зеркало в репозитории сейчас не отслеживается. Источники: \path{tests/test_codex_proxy.py}.\par\medskip
\textbf{\url{http://127.0.0.1:40111/openai`}}. Хост: 127.0.0.1:40111. Упоминаний: 1. Ссылка хранится только как URL. Локальное зеркало в репозитории сейчас не отслеживается. Источники: \path{docs/audit/2026-05-02-zz-codex-exec-resume-session-state-stage0.md}.\par\medskip
\textbf{\url{http://127.0.0.1:44973/openai`}}. Хост: 127.0.0.1:44973. Упоминаний: 1. Ссылка хранится только как URL. Локальное зеркало в репозитории сейчас не отслеживается. Источники: \path{docs/audit/2026-05-02-zz-codex-exec-resume-session-state-stage0.md}.\par\medskip
\textbf{\url{https://api.github.com/repos/example/demo/branches/master/protection}}. Хост: api.github.com. Упоминаний: 1. Внешний URL GitHub, упомянутый в 1 файле(ах). Локальные сопроводительные пути: none. Источники: \path{tests/test_github_pr_gates.py}.\par\medskip
\textbf{\url{https://astral.sh/}}. Хост: astral.sh. Упоминаний: 2. Ссылка хранится только как URL. Локальное зеркало в репозитории сейчас не отслеживается. Источники: \path{tests/test_exploratorium_translation.py}\newline \path{tests/test_utilities.py}.\par\medskip
\textbf{\url{https://astral.sh/uv/install.sh}}. Хост: astral.sh. Упоминаний: 1. Ссылка хранится только как URL. Локальное зеркало в репозитории сейчас не отслеживается. Источники: \path{README.md}.\par\medskip
\textbf{\url{https://dspy-helper.openai.azure.com}}. Хост: dspy-helper.openai.azure.com. Упоминаний: 1. Ссылка хранится только как URL. Локальное зеркало в репозитории сейчас не отслеживается. Источники: \path{tests/test_dspy_training.py}.\par\medskip
\textbf{\url{https://dspy-helper.openai.azure.com/}}. Хост: dspy-helper.openai.azure.com. Упоминаний: 1. Ссылка хранится только как URL. Локальное зеркало в репозитории сейчас не отслеживается. Источники: \path{tests/test_dspy_training.py}.\par\medskip
\textbf{\url{https://dspy.ai/tutorials/rag/}}. Хост: dspy.ai. Упоминаний: 1. Ссылка хранится только как URL. Локальное зеркало в репозитории сейчас не отслеживается. Источники: \path{docs/architecture/inspired/dspy-rag-tutorial.md}.\par\medskip
\textbf{\url{https://example.com}}. Хост: example.com. Упоминаний: 2. Ссылка хранится только как URL. Локальное зеркало в репозитории сейчас не отслеживается. Источники: \path{tests/test_pages_site.py}.\par\medskip
\textbf{\url{https://example.com/docs}}. Хост: example.com. Упоминаний: 1. Ссылка хранится только как URL. Локальное зеркало в репозитории сейчас не отслеживается. Источники: \path{tests/test_pages_site.py}.\par\medskip
\textbf{\url{https://example.com/org/repo/blob/main}}. Хост: example.com. Упоминаний: 1. Ссылка хранится только как URL. Локальное зеркало в репозитории сейчас не отслеживается. Источники: \path{tests/test_todo_backlog.py}.\par\medskip
\textbf{\url{https://example.invalid/v1}}. Хост: example.invalid. Упоминаний: 2. Ссылка хранится только как URL. Локальное зеркало в репозитории сейчас не отслеживается. Источники: \path{tests/test_dspy_training.py}.\par\medskip
\textbf{\url{https://example.openai.azure.com}}. Хост: example.openai.azure.com. Упоминаний: 5. Ссылка хранится только как URL. Локальное зеркало в репозитории сейчас не отслеживается. Источники: \path{tests/test_azure_runtime.py}\newline \path{tests/test_dspy_training.py}.\par\medskip
\textbf{\url{https://example.openai.azure.com/}}. Хост: example.openai.azure.com. Упоминаний: 6. Ссылка хранится только как URL. Локальное зеркало в репозитории сейчас не отслеживается. Источники: \path{tests/test_azure_runtime.py}\newline \path{tests/test_dspy_training.py}.\par\medskip
\textbf{\url{https://example.openai.azure.com/openai/deployments/repo-rag-ft}}. Хост: example.openai.azure.com. Упоминаний: 3. Ссылка хранится только как URL. Локальное зеркало в репозитории сейчас не отслеживается. Источники: \path{tests/test_azure_runtime.py}.\par\medskip
\textbf{\url{https://example.openai.azure.com/openai/deployments/repo-rag-ft/}}. Хост: example.openai.azure.com. Упоминаний: 4. Ссылка хранится только как URL. Локальное зеркало в репозитории сейчас не отслеживается. Источники: \path{tests/test_azure_runtime.py}.\par\medskip
\textbf{\url{https://example.openai.azure.com/openai/deployments/repo-rag-ft/chat/completions}}. Хост: example.openai.azure.com. Упоминаний: 1. Ссылка хранится только как URL. Локальное зеркало в репозитории сейчас не отслеживается. Источники: \path{tests/test_dspy_training.py}.\par\medskip
\textbf{\url{https://example.services.ai.azure.com/models}}. Хост: example.services.ai.azure.com. Упоминаний: 6. Ссылка хранится только как URL. Локальное зеркало в репозитории сейчас не отслеживается. Источники: \path{Makefile}\newline \path{docs/architecture/dspy-rag-guide.md}\newline \path{docs/operations/azure-deployment.md}\newline \path{src/repo_rag_lab/utilities.py}\newline \path{tests/test_cli_and_dspy.py}\newline \path{tests/test_project_surfaces.py}.\par\medskip
\textbf{\url{https://from-connection-string}}. Хост: from-connection-string. Упоминаний: 2. Ссылка хранится только как URL. Локальное зеркало в репозитории сейчас не отслеживается. Источники: \path{tests/test_runtime_artifacts_azure.py}.\par\medskip
\textbf{\url{https://github.com/}}. Хост: github.com. Упоминаний: 2. Внешний URL GitHub, упомянутый в 1 файле(ах). Локальные сопроводительные пути: none. Источники: \path{src/repo_rag_lab/pages_site.py}.\par\medskip
\textbf{\url{https://github.com/acme/repo}}. Хост: github.com. Упоминаний: 3. Внешний URL GitHub, упомянутый в 1 файле(ах). Локальные сопроводительные пути: none. Источники: \path{tests/test_codex_proxy.py}.\par\medskip
\textbf{\url{https://github.com/example/demo}}. Хост: github.com. Упоминаний: 10. Внешний URL GitHub, упомянутый в 3 файле(ах). Локальные сопроводительные пути: none. Источники: \path{tests/test_cli_and_dspy.py}\newline \path{tests/test_pages_site.py}\newline \path{tests/test_utilities.py}.\par\medskip
\textbf{\url{https://github.com/example/demo/blob/main/README.md}}. Хост: github.com. Упоминаний: 1. Внешний URL GitHub, упомянутый в 1 файле(ах). Локальные сопроводительные пути: none. Источники: \path{tests/test_pages_site.py}.\par\medskip
\textbf{\url{https://github.com/example/demo/blob/main/assets/logo.png}}. Хост: github.com. Упоминаний: 1. Внешний URL GitHub, упомянутый в 1 файле(ах). Локальные сопроводительные пути: none. Источники: \path{tests/test_pages_site.py}.\par\medskip
\textbf{\url{https://github.com/example/demo/blob/main/src/tool.c}}. Хост: github.com. Упоминаний: 1. Внешний URL GitHub, упомянутый в 1 файле(ах). Локальные сопроводительные пути: none. Источники: \path{tests/test_pages_site.py}.\par\medskip
\textbf{\url{https://github.com/example/demo/tree/main/docs}}. Хост: github.com. Упоминаний: 1. Внешний URL GitHub, упомянутый в 1 файле(ах). Локальные сопроводительные пути: none. Источники: \path{tests/test_pages_site.py}.\par\medskip
\textbf{\url{https://github.com/example/project}}. Хост: github.com. Упоминаний: 2. Внешний URL GitHub, упомянутый в 2 файле(ах). Локальные сопроводительные пути: none. Источники: \path{tests/test_exploratorium_translation.py}\newline \path{tests/test_utilities.py}.\par\medskip
\textbf{\url{https://github.com/example/project/blob/main/README.md}}. Хост: github.com. Упоминаний: 1. Внешний URL GitHub, упомянутый в 1 файле(ах). Локальные сопроводительные пути: none. Источники: \path{tests/test_exploratorium_translation.py}.\par\medskip
\textbf{\url{https://github.com/example/project/blob/main/README.md\#L1-L2}}. Хост: github.com. Упоминаний: 1. Внешний URL GitHub, упомянутый в 1 файле(ах). Локальные сопроводительные пути: none. Источники: \path{tests/test_exploratorium_translation.py}.\par\medskip
\textbf{\url{https://github.com/realagiorganization/dspy_rag_in_repo_docs_and_impl1}}. Хост: github.com. Упоминаний: 1. Внешний URL GitHub, упомянутый в 1 файле(ах). Локальные сопроводительные пути: none. Источники: \path{mkdocs.yml}.\par\medskip
\textbf{\url{https://github.com/realagiorganization/dspy_rag_in_repo_docs_and_impl1/actions/workflows/ci.yml/badge.svg)](https://github.com/realagiorganization/dspy_rag_in_repo_docs_and_impl1/actions/workflows/ci.yml}}. Хост: github.com. Упоминаний: 2. Внешний URL GitHub, упомянутый в 2 файле(ах). Локальные сопроводительные пути: none. Источники: \path{README.md}\newline \path{notebooks/02_agent_workflow_checklist.ipynb}.\par\medskip
\textbf{\url{https://github.com/realagiorganization/dspy_rag_in_repo_docs_and_impl1/actions/workflows/publish.yml/badge.svg)](https://github.com/realagiorganization/dspy_rag_in_repo_docs_and_impl1/actions/workflows/publish.yml}}. Хост: github.com. Упоминаний: 1. Внешний URL GitHub, упомянутый в 1 файле(ах). Локальные сопроводительные пути: none. Источники: \path{README.md}.\par\medskip
\textbf{\url{https://github.com/realagiorganization/dspy_rag_in_repo_docs_and_impl1/blob/master}}. Хост: github.com. Упоминаний: 3. Внешний URL GitHub, упомянутый в 3 файле(ах). Локальные сопроводительные пути: none. Источники: \path{FILES.md}\newline \path{src/repo_rag_lab/todo_backlog.py}\newline \path{todo-backlog.yaml}.\par\medskip
\textbf{\url{https://github.com/realagiorganization/dspy_rag_in_repo_docs_and_impl1/blob/master/.github/workflows/ci.yml}}. Хост: github.com. Упоминаний: 1. Внешний URL GitHub, упомянутый в 1 файле(ах). Локальные сопроводительные пути: none. Источники: \path{publication/todo-backlog-table.tex}.\par\medskip
\textbf{\url{https://github.com/realagiorganization/dspy_rag_in_repo_docs_and_impl1/blob/master/.github/workflows/publish.yml}}. Хост: github.com. Упоминаний: 1. Внешний URL GitHub, упомянутый в 1 файле(ах). Локальные сопроводительные пути: none. Источники: \path{publication/todo-backlog-table.tex}.\par\medskip
\textbf{\url{https://github.com/realagiorganization/dspy_rag_in_repo_docs_and_impl1/blob/master/Makefile\#L142-L143}}. Хост: github.com. Упоминаний: 1. Внешний URL GitHub, упомянутый в 1 файле(ах). Локальные сопроводительные пути: none. Источники: \path{docs/operations/simple-end-to-end-verification-guide.md}.\par\medskip
\textbf{\url{https://github.com/realagiorganization/dspy_rag_in_repo_docs_and_impl1/blob/master/Makefile\#L153-L154}}. Хост: github.com. Упоминаний: 1. Внешний URL GitHub, упомянутый в 1 файле(ах). Локальные сопроводительные пути: none. Источники: \path{docs/operations/simple-end-to-end-verification-guide.md}.\par\medskip
\textbf{\url{https://github.com/realagiorganization/dspy_rag_in_repo_docs_and_impl1/blob/master/Makefile\#L248-L251}}. Хост: github.com. Упоминаний: 1. Внешний URL GitHub, упомянутый в 1 файле(ах). Локальные сопроводительные пути: none. Источники: \path{docs/operations/simple-end-to-end-verification-guide.md}.\par\medskip
\textbf{\url{https://github.com/realagiorganization/dspy_rag_in_repo_docs_and_impl1/blob/master/Makefile\#L65-L66}}. Хост: github.com. Упоминаний: 1. Внешний URL GitHub, упомянутый в 1 файле(ах). Локальные сопроводительные пути: none. Источники: \path{docs/operations/simple-end-to-end-verification-guide.md}.\par\medskip
\textbf{\url{https://github.com/realagiorganization/dspy_rag_in_repo_docs_and_impl1/blob/master/README.md}}. Хост: github.com. Упоминаний: 2. Внешний URL GitHub, упомянутый в 1 файле(ах). Локальные сопроводительные пути: none. Источники: \path{publication/todo-backlog-table.tex}.\par\medskip
\textbf{\url{https://github.com/realagiorganization/dspy_rag_in_repo_docs_and_impl1/blob/master/TODO.MD}}. Хост: github.com. Упоминаний: 1. Внешний URL GitHub, упомянутый в 1 файле(ах). Локальные сопроводительные пути: none. Источники: \path{publication/todo-backlog-table.tex}.\par\medskip
\textbf{\url{https://github.com/realagiorganization/dspy_rag_in_repo_docs_and_impl1/blob/master/docs/architecture/dspy-rag-guide.md}}. Хост: github.com. Упоминаний: 1. Внешний URL GitHub, упомянутый в 1 файле(ах). Локальные сопроводительные пути: none. Источники: \path{publication/todo-backlog-table.tex}.\par\medskip
\textbf{\url{https://github.com/realagiorganization/dspy_rag_in_repo_docs_and_impl1/blob/master/docs/archive/repo-completeness-checklist.md}}. Хост: github.com. Упоминаний: 1. Внешний URL GitHub, упомянутый в 1 файле(ах). Локальные сопроводительные пути: none. Источники: \path{publication/todo-backlog-table.tex}.\par\medskip
\textbf{\url{https://github.com/realagiorganization/dspy_rag_in_repo_docs_and_impl1/blob/master/docs/audit/README.md}}. Хост: github.com. Упоминаний: 3. Внешний URL GitHub, упомянутый в 1 файле(ах). Локальные сопроводительные пути: none. Источники: \path{publication/todo-backlog-table.tex}.\par\medskip
\textbf{\url{https://github.com/realagiorganization/dspy_rag_in_repo_docs_and_impl1/blob/master/docs/operations/azure-deployment.md}}. Хост: github.com. Упоминаний: 1. Внешний URL GitHub, упомянутый в 1 файле(ах). Локальные сопроводительные пути: none. Источники: \path{publication/todo-backlog-table.tex}.\par\medskip
\textbf{\url{https://github.com/realagiorganization/dspy_rag_in_repo_docs_and_impl1/blob/master/docs/planning/dataset-integration-plan.md}}. Хост: github.com. Упоминаний: 1. Внешний URL GitHub, упомянутый в 1 файле(ах). Локальные сопроводительные пути: none. Источники: \path{publication/todo-backlog-table.tex}.\par\medskip
\textbf{\url{https://github.com/realagiorganization/dspy_rag_in_repo_docs_and_impl1/blob/master/docs/planning/repo-hardening-plan.md}}. Хост: github.com. Упоминаний: 1. Внешний URL GitHub, упомянутый в 1 файле(ах). Локальные сопроводительные пути: none. Источники: \path{publication/todo-backlog-table.tex}.\par\medskip
\textbf{\url{https://github.com/realagiorganization/dspy_rag_in_repo_docs_and_impl1/blob/master/notebooks/01_repo_rag_research.ipynb}}. Хост: github.com. Упоминаний: 1. Внешний URL GitHub, упомянутый в 1 файле(ах). Локальные сопроводительные пути: none. Источники: \path{publication/todo-backlog-table.tex}.\par\medskip
\textbf{\url{https://github.com/realagiorganization/dspy_rag_in_repo_docs_and_impl1/blob/master/notebooks/03_dspy_training_lab.ipynb}}. Хост: github.com. Упоминаний: 1. Внешний URL GitHub, упомянутый в 1 файле(ах). Локальные сопроводительные пути: none. Источники: \path{publication/todo-backlog-table.tex}.\par\medskip
\textbf{\url{https://github.com/realagiorganization/dspy_rag_in_repo_docs_and_impl1/blob/master/notebooks/04_sample_population_lab.ipynb}}. Хост: github.com. Упоминаний: 1. Внешний URL GitHub, упомянутый в 1 файле(ах). Локальные сопроводительные пути: none. Источники: \path{publication/todo-backlog-table.tex}.\par\medskip
\textbf{\url{https://github.com/realagiorganization/dspy_rag_in_repo_docs_and_impl1/blob/master/publication/README.md}}. Хост: github.com. Упоминаний: 1. Внешний URL GitHub, упомянутый в 1 файле(ах). Локальные сопроводительные пути: none. Источники: \path{publication/todo-backlog-table.tex}.\par\medskip
\textbf{\url{https://github.com/realagiorganization/dspy_rag_in_repo_docs_and_impl1/blob/master/pyproject.toml}}. Хост: github.com. Упоминаний: 1. Внешний URL GitHub, упомянутый в 1 файле(ах). Локальные сопроводительные пути: none. Источники: \path{publication/todo-backlog-table.tex}.\par\medskip
\textbf{\url{https://github.com/realagiorganization/dspy_rag_in_repo_docs_and_impl1/blob/master/rust-cli/Cargo.lock}}. Хост: github.com. Упоминаний: 1. Внешний URL GitHub, упомянутый в 1 файле(ах). Локальные сопроводительные пути: none. Источники: \path{publication/todo-backlog-table.tex}.\par\medskip
\textbf{\url{https://github.com/realagiorganization/dspy_rag_in_repo_docs_and_impl1/blob/master/rust-cli/Cargo.toml}}. Хост: github.com. Упоминаний: 1. Внешний URL GitHub, упомянутый в 1 файле(ах). Локальные сопроводительные пути: none. Источники: \path{publication/todo-backlog-table.tex}.\par\medskip
\textbf{\url{https://github.com/realagiorganization/dspy_rag_in_repo_docs_and_impl1/blob/master/samples/population/repository_population_candidates.yaml}}. Хост: github.com. Упоминаний: 1. Внешний URL GitHub, упомянутый в 1 файле(ах). Локальные сопроводительные пути: none. Источники: \path{publication/todo-backlog-table.tex}.\par\medskip
\textbf{\url{https://github.com/realagiorganization/dspy_rag_in_repo_docs_and_impl1/blob/master/samples/training/repository_training_examples.yaml}}. Хост: github.com. Упоминаний: 1. Внешний URL GitHub, упомянутый в 1 файле(ах). Локальные сопроводительные пути: none. Источники: \path{publication/todo-backlog-table.tex}.\par\medskip
\textbf{\url{https://github.com/realagiorganization/dspy_rag_in_repo_docs_and_impl1/blob/master/src/repo_rag_lab/benchmarks.py}}. Хост: github.com. Упоминаний: 2. Внешний URL GitHub, упомянутый в 1 файле(ах). Локальные сопроводительные пути: none. Источники: \path{publication/todo-backlog-table.tex}.\par\medskip
\textbf{\url{https://github.com/realagiorganization/dspy_rag_in_repo_docs_and_impl1/blob/master/src/repo_rag_lab/cli.py\#L72-L88}}. Хост: github.com. Упоминаний: 1. Внешний URL GitHub, упомянутый в 1 файле(ах). Локальные сопроводительные пути: none. Источники: \path{docs/operations/simple-end-to-end-verification-guide.md}.\par\medskip
\textbf{\url{https://github.com/realagiorganization/dspy_rag_in_repo_docs_and_impl1/blob/master/src/repo_rag_lab/dspy_workflow.py}}. Хост: github.com. Упоминаний: 1. Внешний URL GitHub, упомянутый в 1 файле(ах). Локальные сопроводительные пути: none. Источники: \path{publication/todo-backlog-table.tex}.\par\medskip
\textbf{\url{https://github.com/realagiorganization/dspy_rag_in_repo_docs_and_impl1/blob/master/src/repo_rag_lab/notebook_scaffolding.py}}. Хост: github.com. Упоминаний: 1. Внешний URL GitHub, упомянутый в 1 файле(ах). Локальные сопроводительные пути: none. Источники: \path{publication/todo-backlog-table.tex}.\par\medskip
\textbf{\url{https://github.com/realagiorganization/dspy_rag_in_repo_docs_and_impl1/blob/master/src/repo_rag_lab/population_samples.py}}. Хост: github.com. Упоминаний: 1. Внешний URL GitHub, упомянутый в 1 файле(ах). Локальные сопроводительные пути: none. Источники: \path{publication/todo-backlog-table.tex}.\par\medskip
\textbf{\url{https://github.com/realagiorganization/dspy_rag_in_repo_docs_and_impl1/blob/master/src/repo_rag_lab/retrieval.py}}. Хост: github.com. Упоминаний: 1. Внешний URL GitHub, упомянутый в 1 файле(ах). Локальные сопроводительные пути: none. Источники: \path{publication/todo-backlog-table.tex}.\par\medskip
\textbf{\url{https://github.com/realagiorganization/dspy_rag_in_repo_docs_and_impl1/blob/master/src/repo_rag_lab/utilities.py\#L108-L124}}. Хост: github.com. Упоминаний: 2. Внешний URL GitHub, упомянутый в 1 файле(ах). Локальные сопроводительные пути: none. Источники: \path{docs/operations/simple-end-to-end-verification-guide.md}.\par\medskip
\textbf{\url{https://github.com/realagiorganization/dspy_rag_in_repo_docs_and_impl1/blob/master/src/repo_rag_lab/utilities.py\#L190-L193}}. Хост: github.com. Упоминаний: 1. Внешний URL GitHub, упомянутый в 1 файле(ах). Локальные сопроводительные пути: none. Источники: \path{docs/operations/simple-end-to-end-verification-guide.md}.\par\medskip
\textbf{\url{https://github.com/realagiorganization/dspy_rag_in_repo_docs_and_impl1/blob/master/src/repo_rag_lab/verification.py\#L11-L51}}. Хост: github.com. Упоминаний: 1. Внешний URL GitHub, упомянутый в 1 файле(ах). Локальные сопроводительные пути: none. Источники: \path{docs/operations/simple-end-to-end-verification-guide.md}.\par\medskip
\textbf{\url{https://github.com/realagiorganization/dspy_rag_in_repo_docs_and_impl1/blob/master/src/repo_rag_lab/verification.py\#L133-L146}}. Хост: github.com. Упоминаний: 1. Внешний URL GitHub, упомянутый в 1 файле(ах). Локальные сопроводительные пути: none. Источники: \path{docs/operations/simple-end-to-end-verification-guide.md}.\par\medskip
\textbf{\url{https://github.com/realagiorganization/dspy_rag_in_repo_docs_and_impl1/blob/master/src/repo_rag_lab/workflow.py}}. Хост: github.com. Упоминаний: 1. Внешний URL GitHub, упомянутый в 1 файле(ах). Локальные сопроводительные пути: none. Источники: \path{publication/todo-backlog-table.tex}.\par\medskip
\textbf{\url{https://github.com/realagiorganization/dspy_rag_in_repo_docs_and_impl1/blob/master/src/repo_rag_lab/workflow.py\#L119-L145}}. Хост: github.com. Упоминаний: 1. Внешний URL GitHub, упомянутый в 1 файле(ах). Локальные сопроводительные пути: none. Источники: \path{docs/operations/simple-end-to-end-verification-guide.md}.\par\medskip
\textbf{\url{https://github.com/realagiorganization/dspy_rag_in_repo_docs_and_impl1/blob/master/src/repo_rag_lab/workflow.py\#L54-L72}}. Хост: github.com. Упоминаний: 1. Внешний URL GitHub, упомянутый в 1 файле(ах). Локальные сопроводительные пути: none. Источники: \path{docs/operations/simple-end-to-end-verification-guide.md}.\par\medskip
\textbf{\url{https://github.com/realagiorganization/dspy_rag_in_repo_docs_and_impl1/blob/master/tests/test_benchmarks_and_notebook_scaffolding.py}}. Хост: github.com. Упоминаний: 1. Внешний URL GitHub, упомянутый в 1 файле(ах). Локальные сопроводительные пути: none. Источники: \path{publication/todo-backlog-table.tex}.\par\medskip
\textbf{\url{https://github.com/realagiorganization/dspy_rag_in_repo_docs_and_impl1/blob/master/tests/test_repository_rag_bdd.py}}. Хост: github.com. Упоминаний: 1. Внешний URL GitHub, упомянутый в 1 файле(ах). Локальные сопроводительные пути: none. Источники: \path{publication/todo-backlog-table.tex}.\par\medskip
\textbf{\url{https://github.com/realagiorganization/dspy_rag_in_repo_docs_and_impl1/blob/master/utilities/README.md}}. Хост: github.com. Упоминаний: 1. Внешний URL GitHub, упомянутый в 1 файле(ах). Локальные сопроводительные пути: none. Источники: \path{publication/todo-backlog-table.tex}.\par\medskip
\textbf{\url{https://github.com/realagiorganization/dspy_rag_in_repo_docs_and_impl1/tree/master/samples/logs/}}. Хост: github.com. Упоминаний: 3. Внешний URL GitHub, упомянутый в 1 файле(ах). Локальные сопроводительные пути: none. Источники: \path{publication/todo-backlog-table.tex}.\par\medskip
\textbf{\url{https://github.com/realagiorganization/dspy_rag_in_repo_docs_and_impl1/tree/master/tests/fixtures/hushwheel_lexiconarium/}}. Хост: github.com. Упоминаний: 1. Внешний URL GitHub, упомянутый в 1 файле(ах). Локальные сопроводительные пути: none. Источники: \path{publication/todo-backlog-table.tex}.\par\medskip
\textbf{\url{https://github.com/realagiorganization/national-debt-relief}}. Хост: github.com. Упоминаний: 1. Внешний URL GitHub, упомянутый в 1 файле(ах). Локальные сопроводительные пути: none. Источники: \path{docs/audit/2026-05-10-aks-run-25632110510-handoff-fixed-runtime-still-heuristic.md}.\par\medskip
\textbf{\url{https://gpt45standard.openai.azure.com/}}. Хост: gpt45standard.openai.azure.com. Упоминаний: 5. Ссылка хранится только как URL. Локальное зеркало в репозитории сейчас не отслеживается. Источники: \path{docs/audit/2026-04-29-live-azure-gpt54-validation.md}.\par\medskip
\textbf{\url{https://gpt45standard.openai.azure.com/`}}. Хост: gpt45standard.openai.azure.com. Упоминаний: 3. Ссылка хранится только как URL. Локальное зеркало в репозитории сейчас не отслеживается. Источники: \path{README.md}\newline \path{docs/audit/2026-04-29-live-azure-gpt54-validation.md}\newline \path{docs/operations/azure-deployment.md}.\par\medskip
\textbf{\url{https://gpt45standard.openai.azure.com/openai/v1}}. Хост: gpt45standard.openai.azure.com. Упоминаний: 1. Ссылка хранится только как URL. Локальное зеркало в репозитории сейчас не отслеживается. Источники: \path{docs/audit/2026-05-01-dataset-azure-responses-api-version-default-fix.md}.\par\medskip
\textbf{\url{https://gpt45standard.openai.azure.com`}}. Хост: gpt45standard.openai.azure.com`. Упоминаний: 1. Ссылка хранится только как URL. Локальное зеркало в репозитории сейчас не отслеживается. Источники: \path{docs/audit/2026-03-18-zzzzzz-dspy-pipeline-and-live-azure-proof.md}.\par\medskip
\textbf{\url{https://img.shields.io/badge/coverage-94.12\%25-brightgreen)](https://github.com/realagiorganization/dspy_rag_in_repo_docs_and_impl1/actions/workflows/ci.yml}}. Хост: img.shields.io. Упоминаний: 1. Ссылка хранится только как URL. Локальное зеркало в репозитории сейчас не отслеживается. Источники: \path{README.md}.\par\medskip
\textbf{\url{https://learn.microsoft.com/azure/ai-foundry/foundry-models/how-to/inference}}. Хост: learn.microsoft.com. Упоминаний: 1. Внешний URL документации Azure. Зеркальная копия upstream не хранится, но есть локальные сопроводительные заметки: docs/operations/azure-deployment.md. Источники: \path{docs/operations/azure-deployment.md}.\par\medskip
\textbf{\url{https://learn.microsoft.com/azure/ai-services/}}. Хост: learn.microsoft.com. Упоминаний: 4. Внешний URL документации Azure. Зеркальная копия upstream не хранится, но есть локальные сопроводительные заметки: docs/operations/azure-deployment.md. Источники: \path{publication/references.bib}\newline \path{tests/test_exploratorium_translation.py}\newline \path{tests/test_utilities.py}.\par\medskip
\textbf{\url{https://learn.microsoft.com/en-us/azure/ai-foundry/openai/how-to/fine-tuning-deploy}}. Хост: learn.microsoft.com. Упоминаний: 1. Внешний URL документации Azure. Зеркальная копия upstream не хранится, но есть локальные сопроводительные заметки: docs/operations/azure-deployment.md. Источники: \path{docs/operations/azure-deployment.md}.\par\medskip
\textbf{\url{https://learn.microsoft.com/en-us/azure/ai-services/openai/how-to/fine-tuning}}. Хост: learn.microsoft.com. Упоминаний: 1. Внешний URL документации Azure. Зеркальная копия upstream не хранится, но есть локальные сопроводительные заметки: docs/operations/azure-deployment.md. Источники: \path{docs/operations/azure-deployment.md}.\par\medskip
\textbf{\url{https://medium.com/@arancibia.juan22/implementing-rag-with-dspy-a-technical-guide-a6ae15f6a455}}. Хост: medium.com. Упоминаний: 1. Ссылка хранится только как URL. Локальное зеркало в репозитории сейчас не отслеживается. Источники: \path{docs/architecture/inspired/implementing-rag-with-dspy-technical-guide.md}.\par\medskip
\textbf{\url{https://modelcontextprotocol.io/}}. Хост: modelcontextprotocol.io. Упоминаний: 4. Внешний URL документации MCP. Зеркальная копия upstream не хранится, но есть локальные сопроводительные заметки: docs/architecture/mcp-discovery.md. Источники: \path{publication/references.bib}\newline \path{tests/test_exploratorium_translation.py}\newline \path{tests/test_utilities.py}.\par\medskip
\textbf{\url{https://realagiorganization.github.io/dspy_rag_in_repo_docs_and_impl1/}}. Хост: realagiorganization.github.io. Упоминаний: 3. Ссылка хранится только как URL. Локальное зеркало в репозитории сейчас не отслеживается. Источники: \path{README.md}\newline \path{docs/audit/2026-03-19-public-pages-and-secret-scan.md}\newline \path{mkdocs.yml}.\par\medskip
\textbf{\url{https://realagiorganization.github.io/dspy_rag_in_repo_docs_and_impl1/`}}. Хост: realagiorganization.github.io. Упоминаний: 1. Ссылка хранится только как URL. Локальное зеркало в репозитории сейчас не отслеживается. Источники: \path{docs/audit/2026-03-19-public-pages-and-secret-scan.md}.\par\medskip
\textbf{\url{https://realagistorage.blob.core.windows.net}}. Хост: realagistorage.blob.core.windows.net. Упоминаний: 1. Ссылка хранится только как URL. Локальное зеркало в репозитории сейчас не отслеживается. Источники: \path{tests/test_runtime_artifacts_azure.py}.\par\medskip
\textbf{\url{https://realagistorage.queue.core.windows.net}}. Хост: realagistorage.queue.core.windows.net. Упоминаний: 1. Ссылка хранится только как URL. Локальное зеркало в репозитории сейчас не отслеживается. Источники: \path{tests/test_runtime_artifacts_azure.py}.\par\medskip


\end{document}
